\documentclass[11pt,a4paper]{article}
\usepackage[T1]{fontenc}
\usepackage{isabelle,isabellesym}

% further packages required for unusual symbols (see also
% isabellesym.sty), use only when needed

\usepackage{amssymb}
  %for \<leadsto>, \<box>, \<diamond>, \<sqsupset>, \<mho>, \<Join>,
  %\<lhd>, \<lesssim>, \<greatersim>, \<lessapprox>, \<greaterapprox>,
  %\<triangleq>, \<yen>, \<lozenge>

%\usepackage{eurosym}
  %for \<euro>

%\usepackage[only,bigsqcap,bigparallel,fatsemi,interleave,sslash]{stmaryrd}
  %for \<Sqinter>, \<Parallel>, \<Zsemi>, \<Parallel>, \<sslash>

%\usepackage{eufrak}
  %for \<AA> ... \<ZZ>, \<aa> ... \<zz> (also included in amssymb)

%\usepackage{textcomp}
  %for \<onequarter>, \<onehalf>, \<threequarters>, \<degree>, \<cent>,
  %\<currency>

% this should be the last package used
\usepackage{pdfsetup}

% urls in roman style, theory text in math-similar italics
\urlstyle{rm}
\isabellestyle{it}

% for uniform font size
%\renewcommand{\isastyle}{\isastyleminor}


\begin{document}

\title{Nominal $\alpha$-Unification}
\author{Guilherme Borges Brand\~ao$^\dag$ \and Thomas Ammer$^\ddag$ \and Daniele Nantes Sobrinho$^\dag$ \and Mauricio Ayala-Rinci\'on$^\dag$ \and  Christian Urban$^\ddag$ \and Maribel Fern\'andez$^\ddag$ \and Mohammad Abdulaziz$^\ddag$\\[2mm]
\mbox{}$^\dag$ Universidade de Bras\'ilia, Bras\'ilia D.F., Brazil \\[1mm]
\mbox{}$^\ddag$King's College London, London, U.K.
 }

\maketitle

\begin{abstract}
This theory verifies a non-deterministic procedure for nominal syntactic unification, i.e., nominal $\alpha$-unification. The procedure is presented as a set of inductive rules, following Urban's seminal Isabelle formalisation in Isabelle \cite{DBLP:journals/tcs/UrbanPG04,DBLP:journals/corr/abs-1012-4890}. 

The formalisation is updated according to the PVS Rocha Oliveira et al. approach \cite{DBLP:journals/entcs/Ayala-RinconFO16,ACRochaOliveira2016}, by proving the properties of symmetry, transitivity, and equivariance of the $\alpha$-equivalence relation separately. This treatment of $\alpha$-equivalence properties was also followed in the Coq formalisation by de Carvalho's et al. \cite{DBLP:journals/tcs/Ayala-RinconSFN19,WLRdeCarvalho2019}. In the PVS formalisation, nominal $\alpha$-equivalence is presented as a functional recursive algorithm, proved sound and complete, from which executable code in Lisp can be extracted. In the Coq formalisation, a non-deterministic rule-based procedure was presented, following Urban's seminal approach, which is proved sound and complete; further, a recursive definition was presented, which is proved equivalent to the inductive procedure, and from which executable code was extracted. 

\end{abstract}
\tableofcontents

% sane default for proof documents
\parindent 0pt\parskip 0.5ex

% generated text of all theories
%
\begin{isabellebody}%
\setisabellecontext{Swap}%
%
\isadelimtheory
%
\endisadelimtheory
%
\isatagtheory
\isacommand{theory}\isamarkupfalse%
\ Swap\ \isanewline
\ \ \isakeyword{imports}\ Main\isanewline
\isakeyword{begin}%
\endisatagtheory
{\isafoldtheory}%
%
\isadelimtheory
\isanewline
%
\endisadelimtheory
\isanewline
\isanewline
\isanewline
\isacommand{fun}\isamarkupfalse%
\ swapa\ {\isacharcolon}{\kern0pt}{\isacharcolon}{\kern0pt}\ {\isachardoublequoteopen}{\isacharparenleft}{\kern0pt}string\ {\isasymtimes}\ string{\isacharparenright}{\kern0pt}\ {\isasymRightarrow}\ string\ {\isasymRightarrow}\ string{\isachardoublequoteclose}\isanewline
\ \ \isakeyword{where}\isanewline
\ \ {\isachardoublequoteopen}swapa\ {\isacharparenleft}{\kern0pt}b{\isacharcomma}{\kern0pt}c{\isacharparenright}{\kern0pt}\ a\ {\isacharequal}{\kern0pt}\ {\isacharparenleft}{\kern0pt}if\ b\ {\isacharequal}{\kern0pt}\ a\ then\ c\ else\ {\isacharparenleft}{\kern0pt}if\ c\ {\isacharequal}{\kern0pt}\ a\ then\ b\ else\ a{\isacharparenright}{\kern0pt}{\isacharparenright}{\kern0pt}{\isachardoublequoteclose}\isanewline
\isanewline
\isacommand{fun}\isamarkupfalse%
\ swapas{\isacharcolon}{\kern0pt}{\isacharcolon}{\kern0pt}\ {\isachardoublequoteopen}{\isacharparenleft}{\kern0pt}string\ {\isasymtimes}\ string{\isacharparenright}{\kern0pt}\ list\ {\isasymRightarrow}\ string\ {\isasymRightarrow}\ string{\isachardoublequoteclose}\isanewline
\ \ \isakeyword{where}\isanewline
\ \ {\isachardoublequoteopen}swapas\ {\isacharbrackleft}{\kern0pt}{\isacharbrackright}{\kern0pt}\ \ \ \ \ a\ {\isacharequal}{\kern0pt}\ a{\isachardoublequoteclose}\isanewline
{\isacharbar}{\kern0pt}\ {\isachardoublequoteopen}swapas\ {\isacharparenleft}{\kern0pt}x{\isacharhash}{\kern0pt}pi{\isacharparenright}{\kern0pt}\ a\ {\isacharequal}{\kern0pt}\ swapa\ x\ {\isacharparenleft}{\kern0pt}swapas\ pi\ a{\isacharparenright}{\kern0pt}{\isachardoublequoteclose}\isanewline
\isanewline
\isacommand{lemma}\isamarkupfalse%
\ {\isacharbrackleft}{\kern0pt}simp{\isacharbrackright}{\kern0pt}{\isacharcolon}{\kern0pt}\ \isanewline
\ \ \isakeyword{shows}\ {\isachardoublequoteopen}swapas\ {\isacharbrackleft}{\kern0pt}{\isacharparenleft}{\kern0pt}a{\isacharcomma}{\kern0pt}a{\isacharparenright}{\kern0pt}{\isacharbrackright}{\kern0pt}\ b{\isacharequal}{\kern0pt}b{\isachardoublequoteclose}\isanewline
%
\isadelimproof
\ \ %
\endisadelimproof
%
\isatagproof
\isacommand{by}\isamarkupfalse%
\ simp%
\endisatagproof
{\isafoldproof}%
%
\isadelimproof
\isanewline
%
\endisadelimproof
\isanewline
\isacommand{lemma}\isamarkupfalse%
\ swapas{\isacharunderscore}{\kern0pt}ab{\isacharunderscore}{\kern0pt}ba{\isacharcolon}{\kern0pt}\isanewline
\ \ \isakeyword{shows}\ {\isachardoublequoteopen}swapas\ {\isacharbrackleft}{\kern0pt}{\isacharparenleft}{\kern0pt}a{\isacharcomma}{\kern0pt}b{\isacharparenright}{\kern0pt}{\isacharbrackright}{\kern0pt}\ {\isacharequal}{\kern0pt}\ swapas\ {\isacharbrackleft}{\kern0pt}{\isacharparenleft}{\kern0pt}b{\isacharcomma}{\kern0pt}a{\isacharparenright}{\kern0pt}{\isacharbrackright}{\kern0pt}{\isachardoublequoteclose}\isanewline
%
\isadelimproof
\ \ %
\endisadelimproof
%
\isatagproof
\isacommand{by}\isamarkupfalse%
\ auto%
\endisatagproof
{\isafoldproof}%
%
\isadelimproof
\isanewline
%
\endisadelimproof
\isanewline
\isacommand{lemma}\isamarkupfalse%
\ swapas{\isacharunderscore}{\kern0pt}append{\isacharcolon}{\kern0pt}\ \isanewline
\ \ \isakeyword{shows}\ {\isachardoublequoteopen}swapas\ {\isacharparenleft}{\kern0pt}pi{\isadigit{1}}{\isacharat}{\kern0pt}pi{\isadigit{2}}{\isacharparenright}{\kern0pt}\ a\ {\isacharequal}{\kern0pt}\ swapas\ pi{\isadigit{1}}\ {\isacharparenleft}{\kern0pt}swapas\ pi{\isadigit{2}}\ a{\isacharparenright}{\kern0pt}{\isachardoublequoteclose}\isanewline
%
\isadelimproof
\ \ %
\endisadelimproof
%
\isatagproof
\isacommand{by}\isamarkupfalse%
\ {\isacharparenleft}{\kern0pt}induct\ pi{\isadigit{1}}{\isacharparenright}{\kern0pt}\ auto%
\endisatagproof
{\isafoldproof}%
%
\isadelimproof
\isanewline
%
\endisadelimproof
\ \isanewline
\isacommand{lemma}\isamarkupfalse%
\ swapas{\isacharunderscore}{\kern0pt}inv\ {\isacharbrackleft}{\kern0pt}simp{\isacharbrackright}{\kern0pt}{\isacharcolon}{\kern0pt}\ \isanewline
\ \ \isakeyword{shows}\ {\isachardoublequoteopen}swapas\ {\isacharparenleft}{\kern0pt}rev\ pi{\isacharparenright}{\kern0pt}\ {\isacharparenleft}{\kern0pt}swapas\ pi\ a{\isacharparenright}{\kern0pt}\ {\isacharequal}{\kern0pt}\ a{\isachardoublequoteclose}\isanewline
%
\isadelimproof
%
\endisadelimproof
%
\isatagproof
\isacommand{apply}\isamarkupfalse%
{\isacharparenleft}{\kern0pt}induct{\isacharunderscore}{\kern0pt}tac\ pi{\isacharparenright}{\kern0pt}\isanewline
\isacommand{apply}\isamarkupfalse%
{\isacharparenleft}{\kern0pt}auto\ simp\ add{\isacharcolon}{\kern0pt}\ swapas{\isacharunderscore}{\kern0pt}append{\isacharparenright}{\kern0pt}\isanewline
\isacommand{done}\isamarkupfalse%
%
\endisatagproof
{\isafoldproof}%
%
\isadelimproof
\isanewline
%
\endisadelimproof
\isanewline
\isacommand{lemma}\isamarkupfalse%
\ swapas{\isacharunderscore}{\kern0pt}rev{\isacharunderscore}{\kern0pt}pi{\isacharunderscore}{\kern0pt}a{\isacharcolon}{\kern0pt}\ \isanewline
\ \ \isakeyword{assumes}\ {\isachardoublequoteopen}swapas\ pi\ a\ {\isacharequal}{\kern0pt}\ b{\isachardoublequoteclose}\isanewline
\ \ \isakeyword{shows}\ {\isachardoublequoteopen}swapas\ {\isacharparenleft}{\kern0pt}rev\ pi{\isacharparenright}{\kern0pt}\ b\ {\isacharequal}{\kern0pt}\ a{\isachardoublequoteclose}\isanewline
%
\isadelimproof
%
\endisadelimproof
%
\isatagproof
\isacommand{using}\isamarkupfalse%
\ assms\ \isacommand{by}\isamarkupfalse%
\ auto%
\endisatagproof
{\isafoldproof}%
%
\isadelimproof
\isanewline
%
\endisadelimproof
\isanewline
\isanewline
\isacommand{lemma}\isamarkupfalse%
\ swapas{\isacharunderscore}{\kern0pt}rev{\isacharunderscore}{\kern0pt}swapas{\isacharunderscore}{\kern0pt}id\ {\isacharbrackleft}{\kern0pt}simp{\isacharbrackright}{\kern0pt}{\isacharcolon}{\kern0pt}\ \isanewline
\ \ \isakeyword{shows}\ {\isachardoublequoteopen}swapas\ pi\ {\isacharparenleft}{\kern0pt}swapas\ {\isacharparenleft}{\kern0pt}rev\ pi{\isacharparenright}{\kern0pt}\ a{\isacharparenright}{\kern0pt}\ {\isacharequal}{\kern0pt}\ a{\isachardoublequoteclose}\isanewline
%
\isadelimproof
\ \ %
\endisadelimproof
%
\isatagproof
\isacommand{by}\isamarkupfalse%
\ {\isacharparenleft}{\kern0pt}metis\ swapas{\isacharunderscore}{\kern0pt}rev{\isacharunderscore}{\kern0pt}pi{\isacharunderscore}{\kern0pt}a{\isacharparenright}{\kern0pt}%
\endisatagproof
{\isafoldproof}%
%
\isadelimproof
\isanewline
%
\endisadelimproof
\isanewline
\isacommand{lemma}\isamarkupfalse%
\ swapas{\isacharunderscore}{\kern0pt}rev{\isacharunderscore}{\kern0pt}pi{\isacharunderscore}{\kern0pt}b{\isacharcolon}{\kern0pt}\ \isanewline
\ \ \isakeyword{assumes}\ {\isachardoublequoteopen}swapas\ {\isacharparenleft}{\kern0pt}rev\ pi{\isacharparenright}{\kern0pt}\ a\ {\isacharequal}{\kern0pt}\ b{\isachardoublequoteclose}\isanewline
\ \ \isakeyword{shows}\ {\isachardoublequoteopen}swapas\ pi\ b\ {\isacharequal}{\kern0pt}\ a{\isachardoublequoteclose}\isanewline
%
\isadelimproof
%
\endisadelimproof
%
\isatagproof
\isacommand{using}\isamarkupfalse%
\ assms\ \isacommand{by}\isamarkupfalse%
\ auto%
\endisatagproof
{\isafoldproof}%
%
\isadelimproof
\isanewline
%
\endisadelimproof
\isanewline
\isacommand{lemma}\isamarkupfalse%
\ swapas{\isacharunderscore}{\kern0pt}comm{\isacharcolon}{\kern0pt}\ \isanewline
\ \ \isakeyword{shows}\ {\isachardoublequoteopen}{\isacharparenleft}{\kern0pt}swapas\ {\isacharparenleft}{\kern0pt}pi{\isacharat}{\kern0pt}{\isacharbrackleft}{\kern0pt}{\isacharparenleft}{\kern0pt}a{\isacharcomma}{\kern0pt}b{\isacharparenright}{\kern0pt}{\isacharbrackright}{\kern0pt}{\isacharparenright}{\kern0pt}\ c{\isacharparenright}{\kern0pt}\ {\isacharequal}{\kern0pt}\ {\isacharparenleft}{\kern0pt}swapas\ {\isacharparenleft}{\kern0pt}{\isacharbrackleft}{\kern0pt}{\isacharparenleft}{\kern0pt}swapas\ pi\ a{\isacharcomma}{\kern0pt}swapas\ pi\ b{\isacharparenright}{\kern0pt}{\isacharbrackright}{\kern0pt}{\isacharat}{\kern0pt}pi{\isacharparenright}{\kern0pt}\ c{\isacharparenright}{\kern0pt}{\isachardoublequoteclose}\isanewline
%
\isadelimproof
\ \ %
\endisadelimproof
%
\isatagproof
\isacommand{by}\isamarkupfalse%
\ {\isacharparenleft}{\kern0pt}metis\ swapa{\isachardot}{\kern0pt}simps\ swapas{\isachardot}{\kern0pt}simps{\isacharparenleft}{\kern0pt}{\isadigit{1}}{\isacharcomma}{\kern0pt}{\isadigit{2}}{\isacharparenright}{\kern0pt}\ swapas{\isacharunderscore}{\kern0pt}append\ swapas{\isacharunderscore}{\kern0pt}rev{\isacharunderscore}{\kern0pt}pi{\isacharunderscore}{\kern0pt}a{\isacharparenright}{\kern0pt}%
\endisatagproof
{\isafoldproof}%
%
\isadelimproof
\isanewline
%
\endisadelimproof
\isanewline
\isanewline
%
\isadelimtheory
\isanewline
%
\endisadelimtheory
%
\isatagtheory
\isacommand{end}\isamarkupfalse%
%
\endisatagtheory
{\isafoldtheory}%
%
\isadelimtheory
%
\endisadelimtheory
%
\end{isabellebody}%
\endinput
%:%file=~/nominal_unification_Isabelle/nominal_alpha_unification/Swap.thy%:%
%:%10=1%:%
%:%11=1%:%
%:%12=2%:%
%:%13=3%:%
%:%18=3%:%
%:%21=4%:%
%:%22=5%:%
%:%23=6%:%
%:%24=7%:%
%:%25=7%:%
%:%26=8%:%
%:%27=9%:%
%:%28=10%:%
%:%29=11%:%
%:%30=11%:%
%:%31=12%:%
%:%32=13%:%
%:%33=14%:%
%:%34=15%:%
%:%35=16%:%
%:%36=16%:%
%:%37=17%:%
%:%40=18%:%
%:%44=18%:%
%:%45=18%:%
%:%50=18%:%
%:%53=19%:%
%:%54=20%:%
%:%55=20%:%
%:%56=21%:%
%:%59=22%:%
%:%63=22%:%
%:%64=22%:%
%:%69=22%:%
%:%72=23%:%
%:%73=24%:%
%:%74=24%:%
%:%75=25%:%
%:%78=26%:%
%:%82=26%:%
%:%83=26%:%
%:%88=26%:%
%:%91=27%:%
%:%92=28%:%
%:%93=28%:%
%:%94=29%:%
%:%101=30%:%
%:%102=30%:%
%:%103=31%:%
%:%104=31%:%
%:%105=32%:%
%:%111=32%:%
%:%114=33%:%
%:%115=34%:%
%:%116=34%:%
%:%117=35%:%
%:%118=36%:%
%:%125=37%:%
%:%126=37%:%
%:%127=37%:%
%:%132=37%:%
%:%135=38%:%
%:%136=39%:%
%:%137=40%:%
%:%138=40%:%
%:%139=41%:%
%:%142=42%:%
%:%146=42%:%
%:%147=42%:%
%:%152=42%:%
%:%155=43%:%
%:%156=44%:%
%:%157=44%:%
%:%158=45%:%
%:%159=46%:%
%:%166=47%:%
%:%167=47%:%
%:%168=47%:%
%:%173=47%:%
%:%176=48%:%
%:%177=49%:%
%:%178=49%:%
%:%179=50%:%
%:%182=51%:%
%:%186=51%:%
%:%187=51%:%
%:%192=51%:%
%:%195=52%:%
%:%196=53%:%
%:%199=54%:%
%:%204=55%:%

%
\begin{isabellebody}%
\setisabellecontext{Terms}%
%
\isadelimtheory
%
\endisadelimtheory
%
\isatagtheory
\isacommand{theory}\isamarkupfalse%
\ Terms\ \isanewline
\ \ \isakeyword{imports}\ Swap\isanewline
\isakeyword{begin}%
\endisatagtheory
{\isafoldtheory}%
%
\isadelimtheory
\isanewline
%
\endisadelimtheory
\isanewline
\isanewline
\isanewline
\isacommand{datatype}\isamarkupfalse%
\ trm\ {\isacharequal}{\kern0pt}\ Abst\ string\ \ trm\ \ \ \isanewline
\ \ \ \ \ \ \ \ \ \ \ \ \ {\isacharbar}{\kern0pt}\ Susp\ {\isachardoublequoteopen}{\isacharparenleft}{\kern0pt}string\ {\isasymtimes}\ string{\isacharparenright}{\kern0pt}\ list{\isachardoublequoteclose}\ string\isanewline
\ \ \ \ \ \ \ \ \ \ \ \ \ {\isacharbar}{\kern0pt}\ Unit\ \ \ \ \ \ \ \ \ \ \ \ \ \ \ \ \ \ \ \ \ \ \ \ \ \ \ \ \ \ \ \ \ \ \ \isanewline
\ \ \ \ \ \ \ \ \ \ \ \ \ {\isacharbar}{\kern0pt}\ Atom\ string\ \ \ \ \ \ \ \ \ \ \ \ \ \ \ \ \ \ \ \ \ \ \ \ \ \ \ \isanewline
\ \ \ \ \ \ \ \ \ \ \ \ \ {\isacharbar}{\kern0pt}\ Paar\ trm\ trm\ \ \ \ \ \ \ \ \ \ \ \ \ \ \ \ \ \ \ \ \ \ \ \ \ \ \isanewline
\ \ \ \ \ \ \ \ \ \ \ \ \ {\isacharbar}{\kern0pt}\ Func\ string\ trm\ \ \ \ \ \ \ \ \ \ \ \ \ \ \ \ \ \ \ \ \ \ \ \isanewline
\isanewline
\isanewline
\isanewline
\isacommand{fun}\isamarkupfalse%
\ swap\ \ {\isacharcolon}{\kern0pt}{\isacharcolon}{\kern0pt}\ {\isachardoublequoteopen}{\isacharparenleft}{\kern0pt}string\ {\isasymtimes}\ string{\isacharparenright}{\kern0pt}\ list\ {\isasymRightarrow}\ trm\ {\isasymRightarrow}\ trm{\isachardoublequoteclose}\isanewline
\ \ \isakeyword{where}\isanewline
\ \ {\isachardoublequoteopen}swap\ \ pi\ {\isacharparenleft}{\kern0pt}Unit{\isacharparenright}{\kern0pt}\ \ \ \ \ \ \ \ {\isacharequal}{\kern0pt}\ Unit{\isachardoublequoteclose}\ \isanewline
{\isacharbar}{\kern0pt}\ {\isachardoublequoteopen}swap\ \ pi\ {\isacharparenleft}{\kern0pt}Atom\ a{\isacharparenright}{\kern0pt}\ \ \ \ \ \ {\isacharequal}{\kern0pt}\ Atom\ {\isacharparenleft}{\kern0pt}swapas\ pi\ a{\isacharparenright}{\kern0pt}{\isachardoublequoteclose}\isanewline
{\isacharbar}{\kern0pt}\ {\isachardoublequoteopen}swap\ \ pi\ {\isacharparenleft}{\kern0pt}Susp\ pi{\isacharprime}{\kern0pt}\ X{\isacharparenright}{\kern0pt}\ \ {\isacharequal}{\kern0pt}\ Susp\ {\isacharparenleft}{\kern0pt}pi\ {\isacharat}{\kern0pt}\ pi{\isacharprime}{\kern0pt}{\isacharparenright}{\kern0pt}\ X{\isachardoublequoteclose}\isanewline
{\isacharbar}{\kern0pt}\ {\isachardoublequoteopen}swap\ \ pi\ {\isacharparenleft}{\kern0pt}Abst\ a\ t{\isacharparenright}{\kern0pt}\ \ \ \ {\isacharequal}{\kern0pt}\ Abst\ {\isacharparenleft}{\kern0pt}swapas\ pi\ a{\isacharparenright}{\kern0pt}\ {\isacharparenleft}{\kern0pt}swap\ pi\ t{\isacharparenright}{\kern0pt}{\isachardoublequoteclose}\isanewline
{\isacharbar}{\kern0pt}\ {\isachardoublequoteopen}swap\ \ pi\ {\isacharparenleft}{\kern0pt}Paar\ t{\isadigit{1}}\ t{\isadigit{2}}{\isacharparenright}{\kern0pt}\ \ {\isacharequal}{\kern0pt}\ Paar\ {\isacharparenleft}{\kern0pt}swap\ pi\ t{\isadigit{1}}{\isacharparenright}{\kern0pt}\ {\isacharparenleft}{\kern0pt}swap\ pi\ t{\isadigit{2}}{\isacharparenright}{\kern0pt}{\isachardoublequoteclose}\isanewline
{\isacharbar}{\kern0pt}\ {\isachardoublequoteopen}swap\ \ pi\ {\isacharparenleft}{\kern0pt}Func\ F\ t{\isacharparenright}{\kern0pt}\ \ \ \ {\isacharequal}{\kern0pt}\ Func\ F\ {\isacharparenleft}{\kern0pt}swap\ pi\ t{\isacharparenright}{\kern0pt}{\isachardoublequoteclose}\isanewline
\isanewline
\isacommand{lemma}\isamarkupfalse%
\ swap{\isacharunderscore}{\kern0pt}id\ {\isacharbrackleft}{\kern0pt}simp{\isacharbrackright}{\kern0pt}{\isacharcolon}{\kern0pt}\ \isanewline
\ \ \isakeyword{shows}\ {\isachardoublequoteopen}swap\ {\isacharbrackleft}{\kern0pt}{\isacharbrackright}{\kern0pt}\ t\ {\isacharequal}{\kern0pt}\ t{\isachardoublequoteclose}\isanewline
%
\isadelimproof
\ \ %
\endisadelimproof
%
\isatagproof
\isacommand{by}\isamarkupfalse%
\ {\isacharparenleft}{\kern0pt}induct{\isacharunderscore}{\kern0pt}tac\ t{\isacharparenright}{\kern0pt}\ {\isacharparenleft}{\kern0pt}auto{\isacharparenright}{\kern0pt}%
\endisatagproof
{\isafoldproof}%
%
\isadelimproof
\isanewline
%
\endisadelimproof
\isanewline
\isanewline
\isacommand{lemma}\isamarkupfalse%
\ swap{\isacharunderscore}{\kern0pt}append{\isacharcolon}{\kern0pt}\ \isanewline
\ \ \isakeyword{shows}\ {\isachardoublequoteopen}swap\ {\isacharparenleft}{\kern0pt}pi{\isadigit{1}}{\isacharat}{\kern0pt}pi{\isadigit{2}}{\isacharparenright}{\kern0pt}\ t\ {\isacharequal}{\kern0pt}\ swap\ pi{\isadigit{1}}\ {\isacharparenleft}{\kern0pt}swap\ pi{\isadigit{2}}\ t{\isacharparenright}{\kern0pt}{\isachardoublequoteclose}\isanewline
%
\isadelimproof
\ \ %
\endisadelimproof
%
\isatagproof
\isacommand{using}\isamarkupfalse%
\ swapas{\isacharunderscore}{\kern0pt}append\ \isacommand{by}\isamarkupfalse%
\ {\isacharparenleft}{\kern0pt}induct\ t\ arbitrary{\isacharcolon}{\kern0pt}\ pi{\isadigit{1}}\ pi{\isadigit{2}}{\isacharparenright}{\kern0pt}\ auto%
\endisatagproof
{\isafoldproof}%
%
\isadelimproof
\isanewline
%
\endisadelimproof
\isanewline
\isanewline
\isacommand{lemma}\isamarkupfalse%
\ swap{\isacharunderscore}{\kern0pt}empty{\isacharcolon}{\kern0pt}\ \isanewline
\ \ \isakeyword{assumes}\ {\isachardoublequoteopen}swap\ pi\ t\ {\isacharequal}{\kern0pt}\ Susp\ {\isacharbrackleft}{\kern0pt}{\isacharbrackright}{\kern0pt}\ X{\isachardoublequoteclose}\ \isanewline
\ \ \isakeyword{shows}\ {\isachardoublequoteopen}pi\ {\isacharequal}{\kern0pt}\ {\isacharbrackleft}{\kern0pt}{\isacharbrackright}{\kern0pt}{\isachardoublequoteclose}\isanewline
%
\isadelimproof
\ \ %
\endisadelimproof
%
\isatagproof
\isacommand{using}\isamarkupfalse%
\ assms\ swap{\isachardot}{\kern0pt}elims\ \isacommand{by}\isamarkupfalse%
\ blast%
\endisatagproof
{\isafoldproof}%
%
\isadelimproof
\isanewline
%
\endisadelimproof
\isanewline
\isanewline
\isanewline
\isanewline
\isacommand{fun}\isamarkupfalse%
\ depth\ {\isacharcolon}{\kern0pt}{\isacharcolon}{\kern0pt}\ {\isachardoublequoteopen}trm\ {\isasymRightarrow}\ nat{\isachardoublequoteclose}\isanewline
\ \ \isakeyword{where}\isanewline
\ \ {\isachardoublequoteopen}depth\ {\isacharparenleft}{\kern0pt}Unit{\isacharparenright}{\kern0pt}\ \ \ \ \ \ {\isacharequal}{\kern0pt}\ {\isadigit{0}}{\isachardoublequoteclose}\isanewline
{\isacharbar}{\kern0pt}\ {\isachardoublequoteopen}depth\ {\isacharparenleft}{\kern0pt}Atom\ a{\isacharparenright}{\kern0pt}\ \ \ \ {\isacharequal}{\kern0pt}\ {\isadigit{0}}{\isachardoublequoteclose}\isanewline
{\isacharbar}{\kern0pt}\ {\isachardoublequoteopen}depth\ {\isacharparenleft}{\kern0pt}Susp\ pi\ X{\isacharparenright}{\kern0pt}\ {\isacharequal}{\kern0pt}\ {\isadigit{0}}{\isachardoublequoteclose}\isanewline
{\isacharbar}{\kern0pt}\ {\isachardoublequoteopen}depth\ {\isacharparenleft}{\kern0pt}Abst\ a\ t{\isacharparenright}{\kern0pt}\ \ {\isacharequal}{\kern0pt}\ {\isadigit{1}}\ {\isacharplus}{\kern0pt}\ depth\ t{\isachardoublequoteclose}\isanewline
{\isacharbar}{\kern0pt}\ {\isachardoublequoteopen}depth\ {\isacharparenleft}{\kern0pt}Func\ F\ t{\isacharparenright}{\kern0pt}\ \ {\isacharequal}{\kern0pt}\ {\isadigit{1}}\ {\isacharplus}{\kern0pt}\ depth\ t{\isachardoublequoteclose}\isanewline
{\isacharbar}{\kern0pt}\ {\isachardoublequoteopen}depth\ {\isacharparenleft}{\kern0pt}Paar\ t\ t{\isacharprime}{\kern0pt}{\isacharparenright}{\kern0pt}\ {\isacharequal}{\kern0pt}\ {\isadigit{1}}\ {\isacharplus}{\kern0pt}\ {\isacharparenleft}{\kern0pt}max\ {\isacharparenleft}{\kern0pt}depth\ t{\isacharparenright}{\kern0pt}\ {\isacharparenleft}{\kern0pt}depth\ t{\isacharprime}{\kern0pt}{\isacharparenright}{\kern0pt}{\isacharparenright}{\kern0pt}{\isachardoublequoteclose}\ \isanewline
\isanewline
\isacommand{lemma}\isamarkupfalse%
\ \isanewline
\ \ Suc{\isacharunderscore}{\kern0pt}max{\isacharunderscore}{\kern0pt}left{\isacharcolon}{\kern0pt}\ \ {\isachardoublequoteopen}Suc{\isacharparenleft}{\kern0pt}max\ x\ y{\isacharparenright}{\kern0pt}\ {\isacharequal}{\kern0pt}\ z\ {\isasymlongrightarrow}\ x\ {\isacharless}{\kern0pt}\ z{\isachardoublequoteclose}\ \isakeyword{and}\isanewline
\ \ Suc{\isacharunderscore}{\kern0pt}max{\isacharunderscore}{\kern0pt}right{\isacharcolon}{\kern0pt}\ {\isachardoublequoteopen}Suc{\isacharparenleft}{\kern0pt}max\ x\ y{\isacharparenright}{\kern0pt}\ {\isacharequal}{\kern0pt}\ z\ {\isasymlongrightarrow}\ y\ {\isacharless}{\kern0pt}\ z{\isachardoublequoteclose}\isanewline
%
\isadelimproof
%
\endisadelimproof
%
\isatagproof
\isacommand{by}\isamarkupfalse%
{\isacharparenleft}{\kern0pt}auto{\isacharparenright}{\kern0pt}%
\endisatagproof
{\isafoldproof}%
%
\isadelimproof
\isanewline
%
\endisadelimproof
\isanewline
\isacommand{lemma}\isamarkupfalse%
\ swap{\isacharunderscore}{\kern0pt}depth\ {\isacharbrackleft}{\kern0pt}simp{\isacharbrackright}{\kern0pt}{\isacharcolon}{\kern0pt}\ \isanewline
\ \ \isakeyword{shows}\ {\isachardoublequoteopen}depth\ {\isacharparenleft}{\kern0pt}swap\ pi\ t{\isacharparenright}{\kern0pt}\ {\isacharequal}{\kern0pt}\ depth\ t{\isachardoublequoteclose}\ \isanewline
%
\isadelimproof
\ \ %
\endisadelimproof
%
\isatagproof
\isacommand{by}\isamarkupfalse%
\ {\isacharparenleft}{\kern0pt}induct\ t{\isacharparenright}{\kern0pt}\ {\isacharparenleft}{\kern0pt}auto{\isacharparenright}{\kern0pt}%
\endisatagproof
{\isafoldproof}%
%
\isadelimproof
\isanewline
%
\endisadelimproof
\isanewline
\isanewline
\isanewline
\isanewline
\isacommand{fun}\isamarkupfalse%
\ occurs\ {\isacharcolon}{\kern0pt}{\isacharcolon}{\kern0pt}\ {\isachardoublequoteopen}string\ {\isasymRightarrow}\ trm\ {\isasymRightarrow}\ bool{\isachardoublequoteclose}\isanewline
\ \ \isakeyword{where}\ \isanewline
\ \ {\isachardoublequoteopen}occurs\ X\ {\isacharparenleft}{\kern0pt}Unit{\isacharparenright}{\kern0pt}\ \ \ \ \ \ \ {\isacharequal}{\kern0pt}\ False{\isachardoublequoteclose}\isanewline
{\isacharbar}{\kern0pt}\ {\isachardoublequoteopen}occurs\ X\ {\isacharparenleft}{\kern0pt}Atom\ a{\isacharparenright}{\kern0pt}\ \ \ \ \ {\isacharequal}{\kern0pt}\ False{\isachardoublequoteclose}\isanewline
{\isacharbar}{\kern0pt}\ {\isachardoublequoteopen}occurs\ X\ {\isacharparenleft}{\kern0pt}Susp\ pi\ Y{\isacharparenright}{\kern0pt}\ \ {\isacharequal}{\kern0pt}\ {\isacharparenleft}{\kern0pt}if\ X\ {\isacharequal}{\kern0pt}\ Y\ then\ True\ else\ False{\isacharparenright}{\kern0pt}{\isachardoublequoteclose}\isanewline
{\isacharbar}{\kern0pt}\ {\isachardoublequoteopen}occurs\ X\ {\isacharparenleft}{\kern0pt}Abst\ a\ t{\isacharparenright}{\kern0pt}\ \ \ {\isacharequal}{\kern0pt}\ occurs\ X\ t{\isachardoublequoteclose}\isanewline
{\isacharbar}{\kern0pt}\ {\isachardoublequoteopen}occurs\ X\ {\isacharparenleft}{\kern0pt}Paar\ t{\isadigit{1}}\ t{\isadigit{2}}{\isacharparenright}{\kern0pt}\ {\isacharequal}{\kern0pt}\ {\isacharparenleft}{\kern0pt}if\ {\isacharparenleft}{\kern0pt}occurs\ X\ t{\isadigit{1}}{\isacharparenright}{\kern0pt}\ then\ True\ else\ {\isacharparenleft}{\kern0pt}occurs\ X\ t{\isadigit{2}}{\isacharparenright}{\kern0pt}{\isacharparenright}{\kern0pt}{\isachardoublequoteclose}\isanewline
{\isacharbar}{\kern0pt}\ {\isachardoublequoteopen}occurs\ X\ {\isacharparenleft}{\kern0pt}Func\ F\ t{\isacharparenright}{\kern0pt}\ \ \ {\isacharequal}{\kern0pt}\ occurs\ X\ t{\isachardoublequoteclose}\isanewline
\isanewline
\isacommand{fun}\isamarkupfalse%
\ vars{\isacharunderscore}{\kern0pt}trm\ {\isacharcolon}{\kern0pt}{\isacharcolon}{\kern0pt}\ {\isachardoublequoteopen}trm\ {\isasymRightarrow}\ string\ set{\isachardoublequoteclose}\isanewline
\ \ \isakeyword{where}\isanewline
\ \ {\isachardoublequoteopen}vars{\isacharunderscore}{\kern0pt}trm\ {\isacharparenleft}{\kern0pt}Unit{\isacharparenright}{\kern0pt}\ \ \ \ \ \ \ {\isacharequal}{\kern0pt}\ {\isacharbraceleft}{\kern0pt}{\isacharbraceright}{\kern0pt}{\isachardoublequoteclose}\isanewline
{\isacharbar}{\kern0pt}\ {\isachardoublequoteopen}vars{\isacharunderscore}{\kern0pt}trm\ {\isacharparenleft}{\kern0pt}Atom\ a{\isacharparenright}{\kern0pt}\ \ \ \ \ {\isacharequal}{\kern0pt}\ {\isacharbraceleft}{\kern0pt}{\isacharbraceright}{\kern0pt}{\isachardoublequoteclose}\isanewline
{\isacharbar}{\kern0pt}\ {\isachardoublequoteopen}vars{\isacharunderscore}{\kern0pt}trm\ {\isacharparenleft}{\kern0pt}Susp\ pi\ X{\isacharparenright}{\kern0pt}\ \ {\isacharequal}{\kern0pt}\ {\isacharbraceleft}{\kern0pt}X{\isacharbraceright}{\kern0pt}{\isachardoublequoteclose}\isanewline
{\isacharbar}{\kern0pt}\ {\isachardoublequoteopen}vars{\isacharunderscore}{\kern0pt}trm\ {\isacharparenleft}{\kern0pt}Paar\ t{\isadigit{1}}\ t{\isadigit{2}}{\isacharparenright}{\kern0pt}\ {\isacharequal}{\kern0pt}\ {\isacharparenleft}{\kern0pt}vars{\isacharunderscore}{\kern0pt}trm\ t{\isadigit{1}}{\isacharparenright}{\kern0pt}{\isasymunion}{\isacharparenleft}{\kern0pt}vars{\isacharunderscore}{\kern0pt}trm\ t{\isadigit{2}}{\isacharparenright}{\kern0pt}{\isachardoublequoteclose}\isanewline
{\isacharbar}{\kern0pt}\ {\isachardoublequoteopen}vars{\isacharunderscore}{\kern0pt}trm\ {\isacharparenleft}{\kern0pt}Abst\ a\ t{\isacharparenright}{\kern0pt}\ \ \ {\isacharequal}{\kern0pt}\ vars{\isacharunderscore}{\kern0pt}trm\ t{\isachardoublequoteclose}\isanewline
{\isacharbar}{\kern0pt}\ {\isachardoublequoteopen}vars{\isacharunderscore}{\kern0pt}trm\ {\isacharparenleft}{\kern0pt}Func\ F\ t{\isacharparenright}{\kern0pt}\ \ \ {\isacharequal}{\kern0pt}\ vars{\isacharunderscore}{\kern0pt}trm\ t{\isachardoublequoteclose}\isanewline
\isanewline
\isacommand{lemma}\isamarkupfalse%
\ vars{\isacharunderscore}{\kern0pt}swap{\isacharcolon}{\kern0pt}\isanewline
\ \ \isakeyword{shows}\ {\isachardoublequoteopen}vars{\isacharunderscore}{\kern0pt}trm\ {\isacharparenleft}{\kern0pt}swap\ pi\ t{\isacharparenright}{\kern0pt}\ {\isacharequal}{\kern0pt}\ vars{\isacharunderscore}{\kern0pt}trm\ t{\isachardoublequoteclose}\isanewline
%
\isadelimproof
\ \ %
\endisadelimproof
%
\isatagproof
\isacommand{by}\isamarkupfalse%
\ {\isacharparenleft}{\kern0pt}induct\ t{\isacharparenright}{\kern0pt}\ auto%
\endisatagproof
{\isafoldproof}%
%
\isadelimproof
\isanewline
%
\endisadelimproof
\isanewline
\isanewline
\isanewline
\isanewline
\isacommand{fun}\isamarkupfalse%
\ sub{\isacharunderscore}{\kern0pt}trms\ {\isacharcolon}{\kern0pt}{\isacharcolon}{\kern0pt}\ {\isachardoublequoteopen}trm\ {\isasymRightarrow}\ trm\ set{\isachardoublequoteclose}\isanewline
\ \ \isakeyword{where}\ \isanewline
\ \ {\isachardoublequoteopen}sub{\isacharunderscore}{\kern0pt}trms\ {\isacharparenleft}{\kern0pt}Unit{\isacharparenright}{\kern0pt}\ \ \ \ \ \ \ {\isacharequal}{\kern0pt}\ {\isacharbraceleft}{\kern0pt}Unit{\isacharbraceright}{\kern0pt}{\isachardoublequoteclose}\isanewline
{\isacharbar}{\kern0pt}\ {\isachardoublequoteopen}sub{\isacharunderscore}{\kern0pt}trms\ {\isacharparenleft}{\kern0pt}Atom\ a{\isacharparenright}{\kern0pt}\ \ \ \ \ {\isacharequal}{\kern0pt}\ {\isacharbraceleft}{\kern0pt}Atom\ a{\isacharbraceright}{\kern0pt}{\isachardoublequoteclose}\isanewline
{\isacharbar}{\kern0pt}\ {\isachardoublequoteopen}sub{\isacharunderscore}{\kern0pt}trms\ {\isacharparenleft}{\kern0pt}Susp\ pi\ Y{\isacharparenright}{\kern0pt}\ \ {\isacharequal}{\kern0pt}\ {\isacharbraceleft}{\kern0pt}Susp\ pi\ Y{\isacharbraceright}{\kern0pt}{\isachardoublequoteclose}\isanewline
{\isacharbar}{\kern0pt}\ {\isachardoublequoteopen}sub{\isacharunderscore}{\kern0pt}trms\ {\isacharparenleft}{\kern0pt}Abst\ a\ t{\isacharparenright}{\kern0pt}\ \ \ {\isacharequal}{\kern0pt}\ {\isacharbraceleft}{\kern0pt}Abst\ a\ t{\isacharbraceright}{\kern0pt}\ {\isasymunion}\ sub{\isacharunderscore}{\kern0pt}trms\ t{\isachardoublequoteclose}\isanewline
{\isacharbar}{\kern0pt}\ {\isachardoublequoteopen}sub{\isacharunderscore}{\kern0pt}trms\ {\isacharparenleft}{\kern0pt}Paar\ t{\isadigit{1}}\ t{\isadigit{2}}{\isacharparenright}{\kern0pt}\ {\isacharequal}{\kern0pt}\ {\isacharbraceleft}{\kern0pt}Paar\ t{\isadigit{1}}\ t{\isadigit{2}}{\isacharbraceright}{\kern0pt}\ {\isasymunion}\ sub{\isacharunderscore}{\kern0pt}trms\ t{\isadigit{1}}\ {\isasymunion}\ sub{\isacharunderscore}{\kern0pt}trms\ t{\isadigit{2}}{\isachardoublequoteclose}\isanewline
{\isacharbar}{\kern0pt}\ {\isachardoublequoteopen}sub{\isacharunderscore}{\kern0pt}trms\ {\isacharparenleft}{\kern0pt}Func\ F\ t{\isacharparenright}{\kern0pt}\ \ \ {\isacharequal}{\kern0pt}\ {\isacharbraceleft}{\kern0pt}Func\ F\ t{\isacharbraceright}{\kern0pt}\ {\isasymunion}\ sub{\isacharunderscore}{\kern0pt}trms\ t{\isachardoublequoteclose}\isanewline
\isanewline
\isacommand{fun}\isamarkupfalse%
\ psub{\isacharunderscore}{\kern0pt}trms\ {\isacharcolon}{\kern0pt}{\isacharcolon}{\kern0pt}\ {\isachardoublequoteopen}trm\ {\isasymRightarrow}\ trm\ set{\isachardoublequoteclose}\isanewline
\ \ \isakeyword{where}\isanewline
\ \ {\isachardoublequoteopen}psub{\isacharunderscore}{\kern0pt}trms\ {\isacharparenleft}{\kern0pt}Unit{\isacharparenright}{\kern0pt}\ \ \ \ \ \ \ {\isacharequal}{\kern0pt}\ {\isacharbraceleft}{\kern0pt}{\isacharbraceright}{\kern0pt}{\isachardoublequoteclose}\isanewline
{\isacharbar}{\kern0pt}\ {\isachardoublequoteopen}psub{\isacharunderscore}{\kern0pt}trms\ {\isacharparenleft}{\kern0pt}Atom\ a{\isacharparenright}{\kern0pt}\ \ \ \ \ {\isacharequal}{\kern0pt}\ {\isacharbraceleft}{\kern0pt}{\isacharbraceright}{\kern0pt}{\isachardoublequoteclose}\isanewline
{\isacharbar}{\kern0pt}\ {\isachardoublequoteopen}psub{\isacharunderscore}{\kern0pt}trms\ {\isacharparenleft}{\kern0pt}Susp\ pi\ X{\isacharparenright}{\kern0pt}\ \ {\isacharequal}{\kern0pt}\ {\isacharbraceleft}{\kern0pt}{\isacharbraceright}{\kern0pt}{\isachardoublequoteclose}\isanewline
{\isacharbar}{\kern0pt}\ {\isachardoublequoteopen}psub{\isacharunderscore}{\kern0pt}trms\ {\isacharparenleft}{\kern0pt}Paar\ t{\isadigit{1}}\ t{\isadigit{2}}{\isacharparenright}{\kern0pt}\ {\isacharequal}{\kern0pt}\ sub{\isacharunderscore}{\kern0pt}trms\ t{\isadigit{1}}\ {\isasymunion}\ sub{\isacharunderscore}{\kern0pt}trms\ t{\isadigit{2}}{\isachardoublequoteclose}\isanewline
{\isacharbar}{\kern0pt}\ {\isachardoublequoteopen}psub{\isacharunderscore}{\kern0pt}trms\ {\isacharparenleft}{\kern0pt}Abst\ a\ t{\isacharparenright}{\kern0pt}\ \ \ {\isacharequal}{\kern0pt}\ sub{\isacharunderscore}{\kern0pt}trms\ t{\isachardoublequoteclose}\isanewline
{\isacharbar}{\kern0pt}\ {\isachardoublequoteopen}psub{\isacharunderscore}{\kern0pt}trms\ {\isacharparenleft}{\kern0pt}Func\ F\ t{\isacharparenright}{\kern0pt}\ \ \ {\isacharequal}{\kern0pt}\ sub{\isacharunderscore}{\kern0pt}trms\ t{\isachardoublequoteclose}\isanewline
\isanewline
\isacommand{lemma}\isamarkupfalse%
\ psub{\isacharunderscore}{\kern0pt}sub{\isacharunderscore}{\kern0pt}trms{\isacharcolon}{\kern0pt}\ \isanewline
\ \ \isakeyword{assumes}\ {\isachardoublequoteopen}t{\isadigit{1}}\ {\isasymin}\ psub{\isacharunderscore}{\kern0pt}trms\ t{\isadigit{2}}{\isachardoublequoteclose}\ \isanewline
\ \ \isakeyword{shows}\ {\isachardoublequoteopen}t{\isadigit{1}}\ {\isasymin}\ sub{\isacharunderscore}{\kern0pt}trms\ t{\isadigit{2}}{\isachardoublequoteclose}\isanewline
%
\isadelimproof
\ \ %
\endisadelimproof
%
\isatagproof
\isacommand{using}\isamarkupfalse%
\ assms\isanewline
\ \ \isacommand{by}\isamarkupfalse%
{\isacharparenleft}{\kern0pt}induct\ t{\isadigit{2}}{\isacharparenright}{\kern0pt}{\isacharparenleft}{\kern0pt}auto{\isacharparenright}{\kern0pt}%
\endisatagproof
{\isafoldproof}%
%
\isadelimproof
\isanewline
%
\endisadelimproof
\isanewline
\isanewline
\isacommand{lemma}\isamarkupfalse%
\ t{\isacharunderscore}{\kern0pt}sub{\isacharunderscore}{\kern0pt}trms{\isacharunderscore}{\kern0pt}t{\isacharcolon}{\kern0pt}\ \isanewline
\ \ \isakeyword{shows}\ {\isachardoublequoteopen}t\ {\isasymin}\ sub{\isacharunderscore}{\kern0pt}trms\ t{\isachardoublequoteclose}\isanewline
%
\isadelimproof
\ \ %
\endisadelimproof
%
\isatagproof
\isacommand{by}\isamarkupfalse%
\ {\isacharparenleft}{\kern0pt}induct\ t{\isacharparenright}{\kern0pt}\ {\isacharparenleft}{\kern0pt}auto{\isacharparenright}{\kern0pt}%
\endisatagproof
{\isafoldproof}%
%
\isadelimproof
\isanewline
%
\endisadelimproof
\isanewline
\isanewline
\isanewline
\isacommand{lemma}\isamarkupfalse%
\ abst{\isacharunderscore}{\kern0pt}psub{\isacharunderscore}{\kern0pt}trms{\isacharcolon}{\kern0pt}\ \isanewline
\ \ \isakeyword{assumes}\ {\isachardoublequoteopen}Abst\ a\ t{\isadigit{1}}\ {\isasymin}\ sub{\isacharunderscore}{\kern0pt}trms\ t{\isadigit{2}}{\isachardoublequoteclose}\isanewline
\ \ \isakeyword{shows}\ {\isachardoublequoteopen}t{\isadigit{1}}\ {\isasymin}\ psub{\isacharunderscore}{\kern0pt}trms\ t{\isadigit{2}}{\isachardoublequoteclose}\isanewline
%
\isadelimproof
\ \ %
\endisadelimproof
%
\isatagproof
\isacommand{using}\isamarkupfalse%
\ assms\isanewline
\isacommand{apply}\isamarkupfalse%
{\isacharparenleft}{\kern0pt}induct\ t{\isadigit{2}}\ arbitrary{\isacharcolon}{\kern0pt}\ t{\isadigit{1}}{\isacharparenright}{\kern0pt}\isanewline
\isacommand{apply}\isamarkupfalse%
{\isacharparenleft}{\kern0pt}auto\ simp\ add{\isacharcolon}{\kern0pt}\ t{\isacharunderscore}{\kern0pt}sub{\isacharunderscore}{\kern0pt}trms{\isacharunderscore}{\kern0pt}t\ intro{\isacharcolon}{\kern0pt}\ psub{\isacharunderscore}{\kern0pt}sub{\isacharunderscore}{\kern0pt}trms{\isacharparenright}{\kern0pt}\isanewline
\isacommand{done}\isamarkupfalse%
%
\endisatagproof
{\isafoldproof}%
%
\isadelimproof
\isanewline
%
\endisadelimproof
\isanewline
\isacommand{lemma}\isamarkupfalse%
\ func{\isacharunderscore}{\kern0pt}psub{\isacharunderscore}{\kern0pt}trms{\isacharcolon}{\kern0pt}\ \isanewline
\ \ \isakeyword{assumes}\ {\isachardoublequoteopen}Func\ F\ t{\isadigit{1}}\ {\isasymin}\ sub{\isacharunderscore}{\kern0pt}trms\ t{\isadigit{2}}{\isachardoublequoteclose}\isanewline
\ \ \isakeyword{shows}\ {\isachardoublequoteopen}t{\isadigit{1}}\ {\isasymin}\ psub{\isacharunderscore}{\kern0pt}trms\ t{\isadigit{2}}{\isachardoublequoteclose}\isanewline
%
\isadelimproof
\ \ %
\endisadelimproof
%
\isatagproof
\isacommand{using}\isamarkupfalse%
\ assms\isanewline
\isacommand{apply}\isamarkupfalse%
{\isacharparenleft}{\kern0pt}induct\ t{\isadigit{2}}{\isacharparenright}{\kern0pt}\isanewline
\isacommand{apply}\isamarkupfalse%
{\isacharparenleft}{\kern0pt}auto\ simp\ add{\isacharcolon}{\kern0pt}\ t{\isacharunderscore}{\kern0pt}sub{\isacharunderscore}{\kern0pt}trms{\isacharunderscore}{\kern0pt}t\ intro{\isacharcolon}{\kern0pt}\ psub{\isacharunderscore}{\kern0pt}sub{\isacharunderscore}{\kern0pt}trms{\isacharparenright}{\kern0pt}\isanewline
\isacommand{done}\isamarkupfalse%
%
\endisatagproof
{\isafoldproof}%
%
\isadelimproof
\isanewline
%
\endisadelimproof
\isanewline
\isacommand{lemma}\isamarkupfalse%
\ paar{\isacharunderscore}{\kern0pt}psub{\isacharunderscore}{\kern0pt}trms{\isacharcolon}{\kern0pt}\ \isanewline
\ \ \isakeyword{assumes}\ {\isachardoublequoteopen}Paar\ t{\isadigit{1}}\ t{\isadigit{2}}\ {\isasymin}\ sub{\isacharunderscore}{\kern0pt}trms\ t{\isadigit{3}}{\isachardoublequoteclose}\isanewline
\ \ \isakeyword{shows}\ {\isachardoublequoteopen}t{\isadigit{1}}\ {\isasymin}\ psub{\isacharunderscore}{\kern0pt}trms\ t{\isadigit{3}}{\isachardoublequoteclose}\ \isakeyword{and}\ {\isachardoublequoteopen}t{\isadigit{2}}\ {\isasymin}\ psub{\isacharunderscore}{\kern0pt}trms\ t{\isadigit{3}}{\isachardoublequoteclose}\isanewline
%
\isadelimproof
\ \ %
\endisadelimproof
%
\isatagproof
\isacommand{using}\isamarkupfalse%
\ assms\isanewline
\isacommand{apply}\isamarkupfalse%
{\isacharparenleft}{\kern0pt}induct\ t{\isadigit{3}}{\isacharparenright}{\kern0pt}\isanewline
\isacommand{apply}\isamarkupfalse%
{\isacharparenleft}{\kern0pt}auto\ simp\ add{\isacharcolon}{\kern0pt}\ t{\isacharunderscore}{\kern0pt}sub{\isacharunderscore}{\kern0pt}trms{\isacharunderscore}{\kern0pt}t\ intro{\isacharcolon}{\kern0pt}\ psub{\isacharunderscore}{\kern0pt}sub{\isacharunderscore}{\kern0pt}trms{\isacharparenright}{\kern0pt}\isanewline
\ \ \isacommand{done}\isamarkupfalse%
%
\endisatagproof
{\isafoldproof}%
%
\isadelimproof
\isanewline
%
\endisadelimproof
\isanewline
\isacommand{lemma}\isamarkupfalse%
\ t{\isacharunderscore}{\kern0pt}not{\isacharunderscore}{\kern0pt}in{\isacharunderscore}{\kern0pt}psub{\isacharunderscore}{\kern0pt}trms{\isacharunderscore}{\kern0pt}t{\isacharcolon}{\kern0pt}\ \isanewline
\ \ \isakeyword{shows}\ {\isachardoublequoteopen}t\ {\isasymnotin}\ psub{\isacharunderscore}{\kern0pt}trms\ t{\isachardoublequoteclose}\isanewline
%
\isadelimproof
%
\endisadelimproof
%
\isatagproof
\isacommand{proof}\isamarkupfalse%
{\isacharparenleft}{\kern0pt}induct\ t{\isacharparenright}{\kern0pt}\isanewline
\ \ \isacommand{case}\isamarkupfalse%
\ {\isacharparenleft}{\kern0pt}Abst\ x{\isadigit{1}}\ t{\isacharparenright}{\kern0pt}\isanewline
\ \ \isacommand{then}\isamarkupfalse%
\ \isacommand{show}\isamarkupfalse%
\ {\isacharquery}{\kern0pt}case\ \isanewline
\ \ \ \ \isacommand{using}\isamarkupfalse%
\ abst{\isacharunderscore}{\kern0pt}psub{\isacharunderscore}{\kern0pt}trms\ \isacommand{by}\isamarkupfalse%
\ auto\isanewline
\isacommand{next}\isamarkupfalse%
\isanewline
\ \ \isacommand{case}\isamarkupfalse%
\ {\isacharparenleft}{\kern0pt}Susp\ x{\isadigit{1}}\ x{\isadigit{2}}{\isacharparenright}{\kern0pt}\isanewline
\ \ \isacommand{then}\isamarkupfalse%
\ \isacommand{show}\isamarkupfalse%
\ {\isacharquery}{\kern0pt}case\ \isacommand{by}\isamarkupfalse%
\ simp\isanewline
\isacommand{next}\isamarkupfalse%
\isanewline
\ \ \isacommand{case}\isamarkupfalse%
\ Unit\isanewline
\ \ \isacommand{then}\isamarkupfalse%
\ \isacommand{show}\isamarkupfalse%
\ {\isacharquery}{\kern0pt}case\ \isacommand{by}\isamarkupfalse%
\ simp\isanewline
\isacommand{next}\isamarkupfalse%
\isanewline
\ \ \isacommand{case}\isamarkupfalse%
\ {\isacharparenleft}{\kern0pt}Atom\ x{\isacharparenright}{\kern0pt}\isanewline
\ \ \isacommand{then}\isamarkupfalse%
\ \isacommand{show}\isamarkupfalse%
\ {\isacharquery}{\kern0pt}case\ \isacommand{by}\isamarkupfalse%
\ simp\isanewline
\isacommand{next}\isamarkupfalse%
\isanewline
\ \ \isacommand{case}\isamarkupfalse%
\ {\isacharparenleft}{\kern0pt}Paar\ t{\isadigit{1}}\ t{\isadigit{2}}{\isacharparenright}{\kern0pt}\isanewline
\ \ \isacommand{then}\isamarkupfalse%
\ \isacommand{have}\isamarkupfalse%
\ {\isachardoublequoteopen}Paar\ t{\isadigit{1}}\ t{\isadigit{2}}\ {\isasymnotin}\ sub{\isacharunderscore}{\kern0pt}trms\ t{\isadigit{1}}{\isachardoublequoteclose}\ {\isachardoublequoteopen}Paar\ t{\isadigit{1}}\ t{\isadigit{2}}\ {\isasymnotin}\ sub{\isacharunderscore}{\kern0pt}trms\ t{\isadigit{2}}{\isachardoublequoteclose}\isanewline
\ \ \ \ \isacommand{using}\isamarkupfalse%
\ paar{\isacharunderscore}{\kern0pt}psub{\isacharunderscore}{\kern0pt}trms\ \isacommand{by}\isamarkupfalse%
\ auto\isanewline
\ \ \isacommand{then}\isamarkupfalse%
\ \isacommand{show}\isamarkupfalse%
\ {\isacharquery}{\kern0pt}case\ \isacommand{by}\isamarkupfalse%
\ simp\isanewline
\isacommand{next}\isamarkupfalse%
\isanewline
\ \ \isacommand{case}\isamarkupfalse%
\ {\isacharparenleft}{\kern0pt}Func\ f\ t{\isacharparenright}{\kern0pt}\isanewline
\ \ \isacommand{then}\isamarkupfalse%
\ \isacommand{show}\isamarkupfalse%
\ {\isacharquery}{\kern0pt}case\isanewline
\ \ \ \ \isacommand{using}\isamarkupfalse%
\ func{\isacharunderscore}{\kern0pt}psub{\isacharunderscore}{\kern0pt}trms\ \isacommand{by}\isamarkupfalse%
\ auto\isanewline
\isacommand{qed}\isamarkupfalse%
%
\endisatagproof
{\isafoldproof}%
%
\isadelimproof
\isanewline
%
\endisadelimproof
\isanewline
\isacommand{lemma}\isamarkupfalse%
\ if{\isacharunderscore}{\kern0pt}sub{\isacharunderscore}{\kern0pt}not{\isacharunderscore}{\kern0pt}eq{\isacharunderscore}{\kern0pt}then{\isacharunderscore}{\kern0pt}psub{\isacharcolon}{\kern0pt}\ \isanewline
\ \ \isakeyword{assumes}\ {\isachardoublequoteopen}t{\isadigit{1}}\ {\isasymin}\ sub{\isacharunderscore}{\kern0pt}trms\ t{\isadigit{2}}{\isachardoublequoteclose}\ {\isachardoublequoteopen}t{\isadigit{1}}\ {\isasymnoteq}\ t{\isadigit{2}}{\isachardoublequoteclose}\isanewline
\ \ \isakeyword{shows}\ {\isachardoublequoteopen}t{\isadigit{1}}\ {\isasymin}\ psub{\isacharunderscore}{\kern0pt}trms\ t{\isadigit{2}}{\isachardoublequoteclose}\isanewline
%
\isadelimproof
\ \ %
\endisadelimproof
%
\isatagproof
\isacommand{using}\isamarkupfalse%
\ assms\isanewline
\ \ \isacommand{by}\isamarkupfalse%
\ {\isacharparenleft}{\kern0pt}induct\ t{\isadigit{2}}{\isacharparenright}{\kern0pt}\ auto%
\endisatagproof
{\isafoldproof}%
%
\isadelimproof
\isanewline
%
\endisadelimproof
\isanewline
\isanewline
\isacommand{lemma}\isamarkupfalse%
\ depth{\isacharunderscore}{\kern0pt}psub{\isacharunderscore}{\kern0pt}trms{\isacharcolon}{\kern0pt}\ \isanewline
\ \ \isakeyword{assumes}\ {\isachardoublequoteopen}t{\isadigit{1}}\ {\isasymin}\ psub{\isacharunderscore}{\kern0pt}trms\ t{\isadigit{2}}{\isachardoublequoteclose}\isanewline
\ \ \isakeyword{shows}\ {\isachardoublequoteopen}depth\ t{\isadigit{1}}\ {\isacharless}{\kern0pt}\ depth\ t{\isadigit{2}}{\isachardoublequoteclose}\isanewline
%
\isadelimproof
\ \ %
\endisadelimproof
%
\isatagproof
\isacommand{using}\isamarkupfalse%
\ assms\isanewline
\isacommand{proof}\isamarkupfalse%
{\isacharparenleft}{\kern0pt}induct\ t{\isadigit{2}}\ arbitrary{\isacharcolon}{\kern0pt}\ t{\isadigit{1}}{\isacharparenright}{\kern0pt}\isanewline
\ \ \isacommand{case}\isamarkupfalse%
\ {\isacharparenleft}{\kern0pt}Abst\ a\ t{\isacharparenright}{\kern0pt}\isanewline
\ \ \isacommand{then}\isamarkupfalse%
\ \isacommand{show}\isamarkupfalse%
\ {\isacharquery}{\kern0pt}case\isanewline
\ \ \ \ \isacommand{using}\isamarkupfalse%
\ depth{\isachardot}{\kern0pt}simps{\isacharparenleft}{\kern0pt}{\isadigit{4}}{\isacharparenright}{\kern0pt}{\isacharbrackleft}{\kern0pt}of\ a\ t{\isacharbrackright}{\kern0pt}\isanewline
\ \ \ \ \ \ psub{\isacharunderscore}{\kern0pt}trms{\isachardot}{\kern0pt}simps{\isacharparenleft}{\kern0pt}{\isadigit{5}}{\isacharparenright}{\kern0pt}{\isacharbrackleft}{\kern0pt}of\ a\ t{\isacharbrackright}{\kern0pt}\ if{\isacharunderscore}{\kern0pt}sub{\isacharunderscore}{\kern0pt}not{\isacharunderscore}{\kern0pt}eq{\isacharunderscore}{\kern0pt}then{\isacharunderscore}{\kern0pt}psub\isanewline
\ \ \ \ \isacommand{by}\isamarkupfalse%
\ fastforce\isanewline
\isacommand{next}\isamarkupfalse%
\isanewline
\ \ \isacommand{case}\isamarkupfalse%
\ {\isacharparenleft}{\kern0pt}Susp\ pi\ X{\isacharparenright}{\kern0pt}\isanewline
\ \ \isacommand{then}\isamarkupfalse%
\ \isacommand{show}\isamarkupfalse%
\ {\isacharquery}{\kern0pt}case\ \isacommand{by}\isamarkupfalse%
\ simp\isanewline
\isacommand{next}\isamarkupfalse%
\isanewline
\ \ \isacommand{case}\isamarkupfalse%
\ Unit\isanewline
\ \ \isacommand{then}\isamarkupfalse%
\ \isacommand{show}\isamarkupfalse%
\ {\isacharquery}{\kern0pt}case\ \isacommand{by}\isamarkupfalse%
\ simp\isanewline
\isacommand{next}\isamarkupfalse%
\isanewline
\ \ \isacommand{case}\isamarkupfalse%
\ {\isacharparenleft}{\kern0pt}Atom\ a{\isacharparenright}{\kern0pt}\isanewline
\ \ \isacommand{then}\isamarkupfalse%
\ \isacommand{show}\isamarkupfalse%
\ {\isacharquery}{\kern0pt}case\ \isacommand{by}\isamarkupfalse%
\ simp\isanewline
\isacommand{next}\isamarkupfalse%
\isanewline
\ \ \isacommand{case}\isamarkupfalse%
\ {\isacharparenleft}{\kern0pt}Paar\ t{\isadigit{1}}{\isacharprime}{\kern0pt}\ t{\isadigit{2}}{\isacharprime}{\kern0pt}{\isacharparenright}{\kern0pt}\isanewline
\ \ \isacommand{then}\isamarkupfalse%
\ \isacommand{show}\isamarkupfalse%
\ {\isacharquery}{\kern0pt}case\ \isanewline
\ \ \ \ \isacommand{using}\isamarkupfalse%
\ depth{\isachardot}{\kern0pt}simps{\isacharparenleft}{\kern0pt}{\isadigit{6}}{\isacharparenright}{\kern0pt}{\isacharbrackleft}{\kern0pt}of\ t{\isadigit{1}}{\isacharprime}{\kern0pt}\ t{\isadigit{2}}{\isacharprime}{\kern0pt}{\isacharbrackright}{\kern0pt}\ \isanewline
\ \ \ \ \ \ psub{\isacharunderscore}{\kern0pt}trms{\isachardot}{\kern0pt}simps{\isacharparenleft}{\kern0pt}{\isadigit{4}}{\isacharparenright}{\kern0pt}\ if{\isacharunderscore}{\kern0pt}sub{\isacharunderscore}{\kern0pt}not{\isacharunderscore}{\kern0pt}eq{\isacharunderscore}{\kern0pt}then{\isacharunderscore}{\kern0pt}psub\ \isacommand{by}\isamarkupfalse%
\ fastforce\isanewline
\isacommand{next}\isamarkupfalse%
\isanewline
\ \ \isacommand{case}\isamarkupfalse%
\ {\isacharparenleft}{\kern0pt}Func\ f\ t{\isadigit{1}}{\isacharprime}{\kern0pt}{\isacharparenright}{\kern0pt}\isanewline
\ \ \isacommand{then}\isamarkupfalse%
\ \isacommand{show}\isamarkupfalse%
\ {\isacharquery}{\kern0pt}case\isanewline
\ \ \ \ \isacommand{using}\isamarkupfalse%
\ depth{\isachardot}{\kern0pt}simps{\isacharparenleft}{\kern0pt}{\isadigit{5}}{\isacharparenright}{\kern0pt}{\isacharbrackleft}{\kern0pt}of\ f\ t{\isadigit{1}}{\isacharprime}{\kern0pt}{\isacharbrackright}{\kern0pt}\isanewline
\ \ \ \ \ \ psub{\isacharunderscore}{\kern0pt}trms{\isachardot}{\kern0pt}simps{\isacharparenleft}{\kern0pt}{\isadigit{6}}{\isacharparenright}{\kern0pt}{\isacharbrackleft}{\kern0pt}of\ f\ t{\isadigit{1}}{\isacharprime}{\kern0pt}{\isacharbrackright}{\kern0pt}\ if{\isacharunderscore}{\kern0pt}sub{\isacharunderscore}{\kern0pt}not{\isacharunderscore}{\kern0pt}eq{\isacharunderscore}{\kern0pt}then{\isacharunderscore}{\kern0pt}psub\isanewline
\ \ \ \ \isacommand{by}\isamarkupfalse%
\ fastforce\isanewline
\isacommand{qed}\isamarkupfalse%
%
\endisatagproof
{\isafoldproof}%
%
\isadelimproof
\isanewline
%
\endisadelimproof
%
\isadelimtheory
\isanewline
%
\endisadelimtheory
%
\isatagtheory
\isacommand{end}\isamarkupfalse%
%
\endisatagtheory
{\isafoldtheory}%
%
\isadelimtheory
%
\endisadelimtheory
%
\end{isabellebody}%
\endinput
%:%file=~/nominal_unification_Isabelle/nominal_alpha_unification/Terms.thy%:%
%:%10=1%:%
%:%11=1%:%
%:%12=2%:%
%:%13=3%:%
%:%18=3%:%
%:%21=4%:%
%:%22=5%:%
%:%23=6%:%
%:%24=7%:%
%:%25=7%:%
%:%26=8%:%
%:%27=9%:%
%:%28=10%:%
%:%29=11%:%
%:%30=12%:%
%:%31=13%:%
%:%32=14%:%
%:%33=15%:%
%:%34=16%:%
%:%35=16%:%
%:%36=17%:%
%:%37=18%:%
%:%38=19%:%
%:%39=20%:%
%:%40=21%:%
%:%41=22%:%
%:%42=23%:%
%:%43=24%:%
%:%44=25%:%
%:%45=25%:%
%:%46=26%:%
%:%49=27%:%
%:%53=27%:%
%:%54=27%:%
%:%59=27%:%
%:%62=28%:%
%:%63=29%:%
%:%64=30%:%
%:%65=30%:%
%:%66=31%:%
%:%69=32%:%
%:%73=32%:%
%:%74=32%:%
%:%75=32%:%
%:%80=32%:%
%:%83=33%:%
%:%84=34%:%
%:%85=35%:%
%:%86=35%:%
%:%87=36%:%
%:%88=37%:%
%:%91=38%:%
%:%95=38%:%
%:%96=38%:%
%:%97=38%:%
%:%102=38%:%
%:%105=39%:%
%:%106=40%:%
%:%107=41%:%
%:%108=42%:%
%:%109=43%:%
%:%110=43%:%
%:%111=44%:%
%:%112=45%:%
%:%113=46%:%
%:%114=47%:%
%:%115=48%:%
%:%116=49%:%
%:%117=50%:%
%:%118=51%:%
%:%119=52%:%
%:%120=52%:%
%:%121=53%:%
%:%122=54%:%
%:%129=55%:%
%:%130=55%:%
%:%135=55%:%
%:%138=56%:%
%:%139=57%:%
%:%140=57%:%
%:%141=58%:%
%:%144=59%:%
%:%148=59%:%
%:%149=59%:%
%:%154=59%:%
%:%157=60%:%
%:%158=61%:%
%:%159=62%:%
%:%160=63%:%
%:%161=64%:%
%:%162=64%:%
%:%163=65%:%
%:%164=66%:%
%:%165=67%:%
%:%166=68%:%
%:%167=69%:%
%:%168=70%:%
%:%169=71%:%
%:%170=72%:%
%:%171=73%:%
%:%172=73%:%
%:%173=74%:%
%:%174=75%:%
%:%175=76%:%
%:%176=77%:%
%:%177=78%:%
%:%178=79%:%
%:%179=80%:%
%:%180=81%:%
%:%181=82%:%
%:%182=82%:%
%:%183=83%:%
%:%186=84%:%
%:%190=84%:%
%:%191=84%:%
%:%196=84%:%
%:%199=85%:%
%:%200=86%:%
%:%201=87%:%
%:%202=88%:%
%:%203=89%:%
%:%204=89%:%
%:%205=90%:%
%:%206=91%:%
%:%207=92%:%
%:%208=93%:%
%:%209=94%:%
%:%210=95%:%
%:%211=96%:%
%:%212=97%:%
%:%213=98%:%
%:%214=98%:%
%:%215=99%:%
%:%216=100%:%
%:%217=101%:%
%:%218=102%:%
%:%219=103%:%
%:%220=104%:%
%:%221=105%:%
%:%222=106%:%
%:%223=107%:%
%:%224=107%:%
%:%225=108%:%
%:%226=109%:%
%:%229=110%:%
%:%233=110%:%
%:%234=110%:%
%:%235=111%:%
%:%236=111%:%
%:%241=111%:%
%:%244=112%:%
%:%245=113%:%
%:%246=114%:%
%:%247=114%:%
%:%248=115%:%
%:%251=116%:%
%:%255=116%:%
%:%256=116%:%
%:%261=116%:%
%:%264=117%:%
%:%265=118%:%
%:%266=119%:%
%:%267=120%:%
%:%268=120%:%
%:%269=121%:%
%:%270=122%:%
%:%273=123%:%
%:%277=123%:%
%:%278=123%:%
%:%279=124%:%
%:%280=124%:%
%:%281=125%:%
%:%282=125%:%
%:%283=126%:%
%:%289=126%:%
%:%292=127%:%
%:%293=128%:%
%:%294=128%:%
%:%295=129%:%
%:%296=130%:%
%:%299=131%:%
%:%303=131%:%
%:%304=131%:%
%:%305=132%:%
%:%306=132%:%
%:%307=133%:%
%:%308=133%:%
%:%309=134%:%
%:%315=134%:%
%:%318=135%:%
%:%319=136%:%
%:%320=136%:%
%:%321=137%:%
%:%322=138%:%
%:%325=139%:%
%:%329=139%:%
%:%330=139%:%
%:%331=140%:%
%:%332=140%:%
%:%333=141%:%
%:%334=141%:%
%:%335=142%:%
%:%341=142%:%
%:%344=143%:%
%:%345=144%:%
%:%346=144%:%
%:%347=145%:%
%:%354=146%:%
%:%355=146%:%
%:%356=147%:%
%:%357=147%:%
%:%358=148%:%
%:%359=148%:%
%:%360=148%:%
%:%361=149%:%
%:%362=149%:%
%:%363=149%:%
%:%364=150%:%
%:%365=150%:%
%:%366=151%:%
%:%367=151%:%
%:%368=152%:%
%:%369=152%:%
%:%370=152%:%
%:%371=152%:%
%:%372=153%:%
%:%373=153%:%
%:%374=154%:%
%:%375=154%:%
%:%376=155%:%
%:%377=155%:%
%:%378=155%:%
%:%379=155%:%
%:%380=156%:%
%:%381=156%:%
%:%382=157%:%
%:%383=157%:%
%:%384=158%:%
%:%385=158%:%
%:%386=158%:%
%:%387=158%:%
%:%388=159%:%
%:%389=159%:%
%:%390=160%:%
%:%391=160%:%
%:%392=161%:%
%:%393=161%:%
%:%394=161%:%
%:%395=162%:%
%:%396=162%:%
%:%397=162%:%
%:%398=163%:%
%:%399=163%:%
%:%400=163%:%
%:%401=163%:%
%:%402=164%:%
%:%403=164%:%
%:%404=165%:%
%:%405=165%:%
%:%406=166%:%
%:%407=166%:%
%:%408=166%:%
%:%409=167%:%
%:%410=167%:%
%:%411=167%:%
%:%412=168%:%
%:%418=168%:%
%:%421=169%:%
%:%422=170%:%
%:%423=170%:%
%:%424=171%:%
%:%425=172%:%
%:%428=173%:%
%:%432=173%:%
%:%433=173%:%
%:%434=174%:%
%:%435=174%:%
%:%440=174%:%
%:%443=175%:%
%:%444=176%:%
%:%445=177%:%
%:%446=177%:%
%:%447=178%:%
%:%448=179%:%
%:%451=180%:%
%:%455=180%:%
%:%456=180%:%
%:%457=181%:%
%:%458=181%:%
%:%459=182%:%
%:%460=182%:%
%:%461=183%:%
%:%462=183%:%
%:%463=183%:%
%:%464=184%:%
%:%465=184%:%
%:%466=185%:%
%:%467=186%:%
%:%468=186%:%
%:%469=187%:%
%:%470=187%:%
%:%471=188%:%
%:%472=188%:%
%:%473=189%:%
%:%474=189%:%
%:%475=189%:%
%:%476=189%:%
%:%477=190%:%
%:%478=190%:%
%:%479=191%:%
%:%480=191%:%
%:%481=192%:%
%:%482=192%:%
%:%483=192%:%
%:%484=192%:%
%:%485=193%:%
%:%486=193%:%
%:%487=194%:%
%:%488=194%:%
%:%489=195%:%
%:%490=195%:%
%:%491=195%:%
%:%492=195%:%
%:%493=196%:%
%:%494=196%:%
%:%495=197%:%
%:%496=197%:%
%:%497=198%:%
%:%498=198%:%
%:%499=198%:%
%:%500=199%:%
%:%501=199%:%
%:%502=200%:%
%:%503=200%:%
%:%504=201%:%
%:%505=201%:%
%:%506=202%:%
%:%507=202%:%
%:%508=203%:%
%:%509=203%:%
%:%510=203%:%
%:%511=204%:%
%:%512=204%:%
%:%513=205%:%
%:%514=206%:%
%:%515=206%:%
%:%516=207%:%
%:%522=207%:%
%:%527=208%:%
%:%532=209%:%

%
\begin{isabellebody}%
\setisabellecontext{Atoms}%
%
\isadelimtheory
\isanewline
%
\endisadelimtheory
%
\isatagtheory
\isacommand{theory}\isamarkupfalse%
\ Atoms\ \isanewline
\isanewline
\isakeyword{imports}\ Terms\isanewline
\isanewline
\isakeyword{begin}%
\endisatagtheory
{\isafoldtheory}%
%
\isadelimtheory
\isanewline
%
\endisadelimtheory
\isanewline
\isacommand{fun}\isamarkupfalse%
\ atms\ {\isacharcolon}{\kern0pt}{\isacharcolon}{\kern0pt}\ {\isachardoublequoteopen}{\isacharparenleft}{\kern0pt}string\ {\isasymtimes}\ string{\isacharparenright}{\kern0pt}\ list\ {\isasymRightarrow}\ string\ set{\isachardoublequoteclose}\ \isakeyword{where}\isanewline
{\isachardoublequoteopen}atms\ {\isacharbrackleft}{\kern0pt}{\isacharbrackright}{\kern0pt}\ {\isacharequal}{\kern0pt}\ {\isacharbraceleft}{\kern0pt}{\isacharbraceright}{\kern0pt}{\isachardoublequoteclose}\ {\isacharbar}{\kern0pt}\ \isanewline
{\isachardoublequoteopen}atms\ {\isacharparenleft}{\kern0pt}x{\isacharhash}{\kern0pt}xs{\isacharparenright}{\kern0pt}\ {\isacharequal}{\kern0pt}\ {\isacharparenleft}{\kern0pt}{\isacharparenleft}{\kern0pt}atms\ xs{\isacharparenright}{\kern0pt}\ {\isasymunion}\ \ {\isacharbraceleft}{\kern0pt}fst{\isacharparenleft}{\kern0pt}x{\isacharparenright}{\kern0pt}{\isacharcomma}{\kern0pt}snd{\isacharparenleft}{\kern0pt}x{\isacharparenright}{\kern0pt}{\isacharbraceright}{\kern0pt}{\isacharparenright}{\kern0pt}{\isachardoublequoteclose}\isanewline
\isanewline
\isanewline
\isacommand{lemma}\isamarkupfalse%
\ {\isacharbrackleft}{\kern0pt}simp{\isacharbrackright}{\kern0pt}{\isacharcolon}{\kern0pt}\ \isanewline
\ \ {\isachardoublequoteopen}atms\ {\isacharparenleft}{\kern0pt}xs{\isacharat}{\kern0pt}ys{\isacharparenright}{\kern0pt}\ {\isacharequal}{\kern0pt}\ atms\ xs\ {\isasymunion}\ atms\ ys{\isachardoublequoteclose}\isanewline
%
\isadelimproof
\ \ %
\endisadelimproof
%
\isatagproof
\isacommand{by}\isamarkupfalse%
\ {\isacharparenleft}{\kern0pt}induct\ xs{\isacharparenright}{\kern0pt}\ auto%
\endisatagproof
{\isafoldproof}%
%
\isadelimproof
\isanewline
%
\endisadelimproof
\isanewline
\isanewline
\isacommand{lemma}\isamarkupfalse%
\ {\isacharbrackleft}{\kern0pt}simp{\isacharbrackright}{\kern0pt}{\isacharcolon}{\kern0pt}\ \isanewline
\ \ {\isachardoublequoteopen}atms\ {\isacharparenleft}{\kern0pt}rev\ pi{\isacharparenright}{\kern0pt}\ {\isacharequal}{\kern0pt}\ atms\ pi{\isachardoublequoteclose}\isanewline
%
\isadelimproof
\ \ %
\endisadelimproof
%
\isatagproof
\isacommand{by}\isamarkupfalse%
\ {\isacharparenleft}{\kern0pt}induct\ pi{\isacharparenright}{\kern0pt}\ auto%
\endisatagproof
{\isafoldproof}%
%
\isadelimproof
\isanewline
%
\endisadelimproof
\isanewline
\isanewline
\isacommand{lemma}\isamarkupfalse%
\ a{\isacharunderscore}{\kern0pt}not{\isacharunderscore}{\kern0pt}in{\isacharunderscore}{\kern0pt}atms{\isacharcolon}{\kern0pt}\ \isanewline
\ \ \isakeyword{assumes}\ {\isachardoublequoteopen}a\ {\isasymnotin}\ atms\ pi{\isachardoublequoteclose}\isanewline
\ \ \isakeyword{shows}\ {\isachardoublequoteopen}a\ {\isacharequal}{\kern0pt}\ swapas\ pi\ a{\isachardoublequoteclose}\isanewline
%
\isadelimproof
\ \ %
\endisadelimproof
%
\isatagproof
\isacommand{using}\isamarkupfalse%
\ assms\ \isacommand{by}\isamarkupfalse%
\ {\isacharparenleft}{\kern0pt}induct\ pi{\isacharparenright}{\kern0pt}\ auto%
\endisatagproof
{\isafoldproof}%
%
\isadelimproof
\isanewline
%
\endisadelimproof
\isanewline
\isanewline
\isacommand{lemma}\isamarkupfalse%
\ swapas{\isacharunderscore}{\kern0pt}pi{\isacharunderscore}{\kern0pt}ineq{\isacharunderscore}{\kern0pt}a{\isacharcolon}{\kern0pt}\ \isanewline
\ \ \isakeyword{assumes}\ \ {\isachardoublequoteopen}swapas\ pi\ a\ {\isasymnoteq}\ a{\isachardoublequoteclose}\isanewline
\ \ \isakeyword{shows}\ {\isachardoublequoteopen}a\ {\isasymin}\ atms\ pi{\isachardoublequoteclose}\isanewline
%
\isadelimproof
\ \ %
\endisadelimproof
%
\isatagproof
\isacommand{using}\isamarkupfalse%
\ a{\isacharunderscore}{\kern0pt}not{\isacharunderscore}{\kern0pt}in{\isacharunderscore}{\kern0pt}atms\ assms\ \isacommand{by}\isamarkupfalse%
\ force%
\endisatagproof
{\isafoldproof}%
%
\isadelimproof
\isanewline
%
\endisadelimproof
\isanewline
\isanewline
\isacommand{lemma}\isamarkupfalse%
\ a{\isacharunderscore}{\kern0pt}ineq{\isacharunderscore}{\kern0pt}swapas{\isacharunderscore}{\kern0pt}pi{\isacharcolon}{\kern0pt}\ \isanewline
\ \ \isakeyword{assumes}\ {\isachardoublequoteopen}a\ {\isasymnoteq}\ swapas\ pi\ a{\isachardoublequoteclose}\isanewline
\ \ \isakeyword{shows}\ {\isachardoublequoteopen}a\ {\isasymin}\ atms\ pi{\isachardoublequoteclose}\isanewline
%
\isadelimproof
\ \ %
\endisadelimproof
%
\isatagproof
\isacommand{using}\isamarkupfalse%
\ a{\isacharunderscore}{\kern0pt}not{\isacharunderscore}{\kern0pt}in{\isacharunderscore}{\kern0pt}atms\ assms\ \isacommand{by}\isamarkupfalse%
\ blast%
\endisatagproof
{\isafoldproof}%
%
\isadelimproof
\isanewline
%
\endisadelimproof
\isanewline
\isanewline
\isacommand{lemma}\isamarkupfalse%
\ swapas{\isacharunderscore}{\kern0pt}pi{\isacharunderscore}{\kern0pt}in{\isacharunderscore}{\kern0pt}atms{\isacharcolon}{\kern0pt}\ \isanewline
\ \ \isakeyword{assumes}\ {\isachardoublequoteopen}a\ {\isasymin}\ atms\ pi{\isachardoublequoteclose}\ \isanewline
\ \ \isakeyword{shows}{\isachardoublequoteopen}swapas\ pi\ a\ {\isasymin}\ atms\ pi{\isachardoublequoteclose}\isanewline
%
\isadelimproof
\ \ %
\endisadelimproof
%
\isatagproof
\isacommand{using}\isamarkupfalse%
\ assms\ \isacommand{by}\isamarkupfalse%
\ {\isacharparenleft}{\kern0pt}metis\ a{\isacharunderscore}{\kern0pt}not{\isacharunderscore}{\kern0pt}in{\isacharunderscore}{\kern0pt}atms\ swapas{\isacharunderscore}{\kern0pt}rev{\isacharunderscore}{\kern0pt}pi{\isacharunderscore}{\kern0pt}a{\isacharparenright}{\kern0pt}%
\endisatagproof
{\isafoldproof}%
%
\isadelimproof
\isanewline
%
\endisadelimproof
\isanewline
%
\isadelimtheory
\isanewline
%
\endisadelimtheory
%
\isatagtheory
\isacommand{end}\isamarkupfalse%
%
\endisatagtheory
{\isafoldtheory}%
%
\isadelimtheory
%
\endisadelimtheory
%
\end{isabellebody}%
\endinput
%:%file=~/nominal_unification_Isabelle/nominal_alpha_unification/Atoms.thy%:%
%:%6=1%:%
%:%11=2%:%
%:%12=2%:%
%:%13=3%:%
%:%14=4%:%
%:%15=5%:%
%:%16=6%:%
%:%21=6%:%
%:%24=7%:%
%:%25=8%:%
%:%26=8%:%
%:%27=9%:%
%:%28=10%:%
%:%29=11%:%
%:%30=12%:%
%:%31=13%:%
%:%32=13%:%
%:%33=14%:%
%:%36=15%:%
%:%40=15%:%
%:%41=15%:%
%:%46=15%:%
%:%49=16%:%
%:%50=17%:%
%:%51=18%:%
%:%52=18%:%
%:%53=19%:%
%:%56=20%:%
%:%60=20%:%
%:%61=20%:%
%:%66=20%:%
%:%69=21%:%
%:%70=22%:%
%:%71=23%:%
%:%72=23%:%
%:%73=24%:%
%:%74=25%:%
%:%77=26%:%
%:%81=26%:%
%:%82=26%:%
%:%83=26%:%
%:%88=26%:%
%:%91=27%:%
%:%92=28%:%
%:%93=29%:%
%:%94=29%:%
%:%95=30%:%
%:%96=31%:%
%:%99=32%:%
%:%103=32%:%
%:%104=32%:%
%:%105=32%:%
%:%110=32%:%
%:%113=33%:%
%:%114=34%:%
%:%115=35%:%
%:%116=35%:%
%:%117=36%:%
%:%118=37%:%
%:%121=38%:%
%:%125=38%:%
%:%126=38%:%
%:%127=38%:%
%:%132=38%:%
%:%135=39%:%
%:%136=40%:%
%:%137=41%:%
%:%138=41%:%
%:%139=42%:%
%:%140=43%:%
%:%143=44%:%
%:%147=44%:%
%:%148=44%:%
%:%149=44%:%
%:%154=44%:%
%:%157=45%:%
%:%160=46%:%
%:%165=47%:%

%
\begin{isabellebody}%
\setisabellecontext{Disagreement}%
%
\isadelimtheory
\isanewline
%
\endisadelimtheory
%
\isatagtheory
\isacommand{theory}\isamarkupfalse%
\ Disagreement\ \isanewline
\isanewline
\isakeyword{imports}\ Atoms\isanewline
\isanewline
\isakeyword{begin}%
\endisatagtheory
{\isafoldtheory}%
%
\isadelimtheory
\isanewline
%
\endisadelimtheory
\isanewline
\isanewline
\isanewline
\isacommand{definition}\isamarkupfalse%
\ \ ds\ {\isacharcolon}{\kern0pt}{\isacharcolon}{\kern0pt}\ {\isachardoublequoteopen}{\isacharparenleft}{\kern0pt}string\ {\isasymtimes}\ string{\isacharparenright}{\kern0pt}\ list\ {\isasymRightarrow}\ {\isacharparenleft}{\kern0pt}string\ {\isasymtimes}\ string{\isacharparenright}{\kern0pt}\ list\ {\isasymRightarrow}\ string\ set{\isachardoublequoteclose}\ \isakeyword{where}\isanewline
\ \ ds{\isacharunderscore}{\kern0pt}def{\isacharcolon}{\kern0pt}\ {\isachardoublequoteopen}ds\ xs\ ys\ \ {\isasymequiv}\ \ {\isacharbraceleft}{\kern0pt}\ a\ {\isachardot}{\kern0pt}\ a\ {\isasymin}\ {\isacharparenleft}{\kern0pt}atms\ xs\ {\isasymunion}\ atms\ ys{\isacharparenright}{\kern0pt}\ {\isasymand}\ {\isacharparenleft}{\kern0pt}swapas\ xs\ a\ {\isasymnoteq}\ swapas\ ys\ a{\isacharparenright}{\kern0pt}\ {\isacharbraceright}{\kern0pt}{\isachardoublequoteclose}\isanewline
\isanewline
\isanewline
\isanewline
\isacommand{lemma}\isamarkupfalse%
\ \isanewline
\ \ ds{\isacharunderscore}{\kern0pt}elem{\isacharcolon}{\kern0pt}\ {\isachardoublequoteopen}{\isasymlbrakk}swapas\ pi\ a{\isasymnoteq}a{\isasymrbrakk}{\isasymLongrightarrow}a{\isasymin}ds\ {\isacharbrackleft}{\kern0pt}{\isacharbrackright}{\kern0pt}\ pi{\isachardoublequoteclose}\isanewline
%
\isadelimproof
\ \ %
\endisadelimproof
%
\isatagproof
\isacommand{using}\isamarkupfalse%
\ ds{\isacharunderscore}{\kern0pt}def\ swapas{\isacharunderscore}{\kern0pt}pi{\isacharunderscore}{\kern0pt}ineq{\isacharunderscore}{\kern0pt}a\ \isacommand{by}\isamarkupfalse%
\ simp%
\endisatagproof
{\isafoldproof}%
%
\isadelimproof
\isanewline
%
\endisadelimproof
\isanewline
\isanewline
\isacommand{corollary}\isamarkupfalse%
\ ds{\isacharunderscore}{\kern0pt}elem{\isacharunderscore}{\kern0pt}cp{\isacharcolon}{\kern0pt}\ {\isachardoublequoteopen}a\ {\isasymnotin}\ ds\ {\isacharbrackleft}{\kern0pt}{\isacharbrackright}{\kern0pt}\ pi\ {\isasymLongrightarrow}\ swapas\ pi\ a\ {\isacharequal}{\kern0pt}\ a{\isachardoublequoteclose}\isanewline
%
\isadelimproof
\ \ %
\endisadelimproof
%
\isatagproof
\isacommand{using}\isamarkupfalse%
\ ds{\isacharunderscore}{\kern0pt}elem\ \isacommand{by}\isamarkupfalse%
\ blast%
\endisatagproof
{\isafoldproof}%
%
\isadelimproof
\isanewline
%
\endisadelimproof
\isanewline
\isacommand{lemma}\isamarkupfalse%
\ \isanewline
\ \ elem{\isacharunderscore}{\kern0pt}ds{\isacharcolon}{\kern0pt}\ {\isachardoublequoteopen}{\isasymlbrakk}a{\isasymin}ds\ {\isacharbrackleft}{\kern0pt}{\isacharbrackright}{\kern0pt}\ pi{\isasymrbrakk}{\isasymLongrightarrow}a{\isasymnoteq}swapas\ pi\ a{\isachardoublequoteclose}\isanewline
%
\isadelimproof
\ \ %
\endisadelimproof
%
\isatagproof
\isacommand{using}\isamarkupfalse%
\ ds{\isacharunderscore}{\kern0pt}def\ \isacommand{by}\isamarkupfalse%
\ simp%
\endisatagproof
{\isafoldproof}%
%
\isadelimproof
\isanewline
%
\endisadelimproof
\isanewline
\isanewline
\isacommand{lemma}\isamarkupfalse%
\ \isanewline
\ \ ds{\isacharunderscore}{\kern0pt}sym{\isacharcolon}{\kern0pt}\ {\isachardoublequoteopen}ds\ pi{\isadigit{1}}\ pi{\isadigit{2}}\ {\isacharequal}{\kern0pt}\ ds\ pi{\isadigit{2}}\ pi{\isadigit{1}}{\isachardoublequoteclose}\isanewline
%
\isadelimproof
\ \ %
\endisadelimproof
%
\isatagproof
\isacommand{using}\isamarkupfalse%
\ ds{\isacharunderscore}{\kern0pt}def\ \isacommand{by}\isamarkupfalse%
\ auto%
\endisatagproof
{\isafoldproof}%
%
\isadelimproof
\isanewline
%
\endisadelimproof
\isanewline
\isanewline
\isacommand{lemma}\isamarkupfalse%
\ \isanewline
\ \ ds{\isacharunderscore}{\kern0pt}trans{\isacharcolon}{\kern0pt}\ {\isachardoublequoteopen}c\ {\isasymin}\ ds\ pi{\isadigit{1}}\ pi{\isadigit{3}}\ {\isasymLongrightarrow}\ {\isacharparenleft}{\kern0pt}c\ {\isasymin}\ ds\ pi{\isadigit{1}}\ pi{\isadigit{2}}\ {\isasymor}\ c\ {\isasymin}\ ds\ pi{\isadigit{2}}\ pi{\isadigit{3}}{\isacharparenright}{\kern0pt}{\isachardoublequoteclose}\isanewline
%
\isadelimproof
%
\endisadelimproof
%
\isatagproof
\isacommand{using}\isamarkupfalse%
\ ds{\isacharunderscore}{\kern0pt}def\ a{\isacharunderscore}{\kern0pt}not{\isacharunderscore}{\kern0pt}in{\isacharunderscore}{\kern0pt}atms\ swapas{\isacharunderscore}{\kern0pt}pi{\isacharunderscore}{\kern0pt}ineq{\isacharunderscore}{\kern0pt}a\ \isacommand{by}\isamarkupfalse%
\ auto%
\endisatagproof
{\isafoldproof}%
%
\isadelimproof
\isanewline
%
\endisadelimproof
\isanewline
\isanewline
\isanewline
\isacommand{lemma}\isamarkupfalse%
\ ds{\isacharunderscore}{\kern0pt}cancel{\isacharunderscore}{\kern0pt}pi{\isacharunderscore}{\kern0pt}left{\isacharcolon}{\kern0pt}\isanewline
\ \ \isakeyword{assumes}\ {\isachardoublequoteopen}{\isacharparenleft}{\kern0pt}c{\isasymin}\ ds\ {\isacharparenleft}{\kern0pt}pi{\isadigit{1}}{\isacharat}{\kern0pt}pi{\isacharparenright}{\kern0pt}\ {\isacharparenleft}{\kern0pt}pi{\isadigit{2}}{\isacharat}{\kern0pt}pi{\isacharparenright}{\kern0pt}{\isacharparenright}{\kern0pt}{\isachardoublequoteclose}\isanewline
\ \ \isakeyword{shows}\ {\isachardoublequoteopen}{\isacharparenleft}{\kern0pt}swapas\ pi\ c\ {\isasymin}\ ds\ pi{\isadigit{1}}\ pi{\isadigit{2}}{\isacharparenright}{\kern0pt}{\isachardoublequoteclose}\isanewline
%
\isadelimproof
\ \ %
\endisadelimproof
%
\isatagproof
\isacommand{using}\isamarkupfalse%
\ assms\ ds{\isacharunderscore}{\kern0pt}def\ swapas{\isacharunderscore}{\kern0pt}append\ a{\isacharunderscore}{\kern0pt}ineq{\isacharunderscore}{\kern0pt}swapas{\isacharunderscore}{\kern0pt}pi\ a{\isacharunderscore}{\kern0pt}not{\isacharunderscore}{\kern0pt}in{\isacharunderscore}{\kern0pt}atms\isanewline
\ \ \isacommand{by}\isamarkupfalse%
\ {\isacharparenleft}{\kern0pt}metis\ {\isacharparenleft}{\kern0pt}mono{\isacharunderscore}{\kern0pt}tags{\isacharcomma}{\kern0pt}\ lifting{\isacharparenright}{\kern0pt}\ Un{\isacharunderscore}{\kern0pt}iff\ mem{\isacharunderscore}{\kern0pt}Collect{\isacharunderscore}{\kern0pt}eq{\isacharparenright}{\kern0pt}%
\endisatagproof
{\isafoldproof}%
%
\isadelimproof
\isanewline
%
\endisadelimproof
\isanewline
\isanewline
\isanewline
\isacommand{lemma}\isamarkupfalse%
\ ds{\isacharunderscore}{\kern0pt}cancel{\isacharunderscore}{\kern0pt}pi{\isacharunderscore}{\kern0pt}right{\isacharcolon}{\kern0pt}\ \isanewline
\ \ {\isachardoublequoteopen}{\isacharparenleft}{\kern0pt}swapas\ pi\ c{\isasymin}\ ds\ pi{\isadigit{1}}\ pi{\isadigit{2}}{\isacharparenright}{\kern0pt}\ {\isasymLongrightarrow}\ {\isacharparenleft}{\kern0pt}c{\isasymin}\ ds\ {\isacharparenleft}{\kern0pt}pi{\isadigit{1}}{\isacharat}{\kern0pt}pi{\isacharparenright}{\kern0pt}\ {\isacharparenleft}{\kern0pt}pi{\isadigit{2}}{\isacharat}{\kern0pt}pi{\isacharparenright}{\kern0pt}{\isacharparenright}{\kern0pt}{\isachardoublequoteclose}\isanewline
%
\isadelimproof
%
\endisadelimproof
%
\isatagproof
\isacommand{apply}\isamarkupfalse%
{\isacharparenleft}{\kern0pt}simp\ only{\isacharcolon}{\kern0pt}\ ds{\isacharunderscore}{\kern0pt}def{\isacharparenright}{\kern0pt}\isanewline
\isacommand{apply}\isamarkupfalse%
{\isacharparenleft}{\kern0pt}auto{\isacharparenright}{\kern0pt}\isanewline
\isacommand{apply}\isamarkupfalse%
{\isacharparenleft}{\kern0pt}simp{\isacharunderscore}{\kern0pt}all\ add{\isacharcolon}{\kern0pt}\ swapas{\isacharunderscore}{\kern0pt}append{\isacharparenright}{\kern0pt}\isanewline
\isacommand{apply}\isamarkupfalse%
{\isacharparenleft}{\kern0pt}rule\ a{\isacharunderscore}{\kern0pt}ineq{\isacharunderscore}{\kern0pt}swapas{\isacharunderscore}{\kern0pt}pi{\isacharcomma}{\kern0pt}clarify{\isacharcomma}{\kern0pt}\isanewline
\ \ \ \ \ \ drule\ a{\isacharunderscore}{\kern0pt}not{\isacharunderscore}{\kern0pt}in{\isacharunderscore}{\kern0pt}atms{\isacharcomma}{\kern0pt}drule\ a{\isacharunderscore}{\kern0pt}not{\isacharunderscore}{\kern0pt}in{\isacharunderscore}{\kern0pt}atms{\isacharcomma}{\kern0pt}simp{\isacharparenright}{\kern0pt}{\isacharplus}{\kern0pt}\isanewline
\ \ \isacommand{done}\isamarkupfalse%
%
\endisatagproof
{\isafoldproof}%
%
\isadelimproof
\isanewline
%
\endisadelimproof
\isanewline
\isanewline
\isanewline
\isacommand{lemma}\isamarkupfalse%
\ ds{\isacharunderscore}{\kern0pt}equality{\isacharcolon}{\kern0pt}\ \isanewline
\ \ {\isachardoublequoteopen}{\isacharparenleft}{\kern0pt}ds\ {\isacharbrackleft}{\kern0pt}{\isacharbrackright}{\kern0pt}\ pi{\isacharparenright}{\kern0pt}{\isacharminus}{\kern0pt}{\isacharbraceleft}{\kern0pt}a{\isacharcomma}{\kern0pt}swapas\ pi\ a{\isacharbraceright}{\kern0pt}\ {\isacharequal}{\kern0pt}\ {\isacharparenleft}{\kern0pt}ds\ {\isacharbrackleft}{\kern0pt}{\isacharbrackright}{\kern0pt}\ {\isacharparenleft}{\kern0pt}{\isacharparenleft}{\kern0pt}a{\isacharcomma}{\kern0pt}swapas\ pi\ a{\isacharparenright}{\kern0pt}{\isacharhash}{\kern0pt}pi{\isacharparenright}{\kern0pt}{\isacharparenright}{\kern0pt}{\isacharminus}{\kern0pt}{\isacharbraceleft}{\kern0pt}swapas\ pi\ a{\isacharbraceright}{\kern0pt}{\isachardoublequoteclose}\isanewline
%
\isadelimproof
\ \ %
\endisadelimproof
%
\isatagproof
\isacommand{using}\isamarkupfalse%
\ ds{\isacharunderscore}{\kern0pt}def\ \isacommand{by}\isamarkupfalse%
\ auto%
\endisatagproof
{\isafoldproof}%
%
\isadelimproof
\isanewline
%
\endisadelimproof
\isanewline
\isanewline
\isacommand{lemma}\isamarkupfalse%
\ ds{\isacharunderscore}{\kern0pt}{\isadigit{7}}{\isacharcolon}{\kern0pt}\ \isanewline
\ \ {\isachardoublequoteopen}{\isasymlbrakk}b{\isasymnoteq}\ swapas\ pi\ b{\isacharsemicolon}{\kern0pt}a{\isasymin}ds\ {\isacharbrackleft}{\kern0pt}{\isacharbrackright}{\kern0pt}\ {\isacharparenleft}{\kern0pt}{\isacharparenleft}{\kern0pt}b{\isacharcomma}{\kern0pt}swapas\ pi\ b{\isacharparenright}{\kern0pt}{\isacharhash}{\kern0pt}pi{\isacharparenright}{\kern0pt}{\isasymrbrakk}{\isasymLongrightarrow}a{\isasymin}ds\ {\isacharbrackleft}{\kern0pt}{\isacharbrackright}{\kern0pt}\ pi{\isachardoublequoteclose}\isanewline
%
\isadelimproof
\ \ %
\endisadelimproof
%
\isatagproof
\isacommand{using}\isamarkupfalse%
\ ds{\isacharunderscore}{\kern0pt}def\ swapas{\isacharunderscore}{\kern0pt}pi{\isacharunderscore}{\kern0pt}in{\isacharunderscore}{\kern0pt}atms\ a{\isacharunderscore}{\kern0pt}ineq{\isacharunderscore}{\kern0pt}swapas{\isacharunderscore}{\kern0pt}pi\ swapas{\isacharunderscore}{\kern0pt}rev{\isacharunderscore}{\kern0pt}pi{\isacharunderscore}{\kern0pt}a\ \isanewline
\ \ \ \ ds{\isacharunderscore}{\kern0pt}elem\ elem{\isacharunderscore}{\kern0pt}ds\ swapa{\isachardot}{\kern0pt}simps\ swapas{\isachardot}{\kern0pt}simps{\isacharparenleft}{\kern0pt}{\isadigit{2}}{\isacharparenright}{\kern0pt}\isanewline
\ \ \isacommand{by}\isamarkupfalse%
\ metis%
\endisatagproof
{\isafoldproof}%
%
\isadelimproof
\isanewline
%
\endisadelimproof
\isanewline
\isanewline
\isanewline
\isacommand{lemma}\isamarkupfalse%
\ ds{\isacharunderscore}{\kern0pt}cancel{\isacharunderscore}{\kern0pt}pi{\isacharunderscore}{\kern0pt}front{\isacharcolon}{\kern0pt}\ \isanewline
\ \ {\isachardoublequoteopen}ds\ {\isacharparenleft}{\kern0pt}pi{\isacharat}{\kern0pt}pi{\isadigit{1}}{\isacharparenright}{\kern0pt}\ {\isacharparenleft}{\kern0pt}pi{\isacharat}{\kern0pt}pi{\isadigit{2}}{\isacharparenright}{\kern0pt}\ {\isacharequal}{\kern0pt}\ ds\ pi{\isadigit{1}}\ pi{\isadigit{2}}{\isachardoublequoteclose}\isanewline
%
\isadelimproof
%
\endisadelimproof
%
\isatagproof
\isacommand{apply}\isamarkupfalse%
{\isacharparenleft}{\kern0pt}simp\ only{\isacharcolon}{\kern0pt}\ ds{\isacharunderscore}{\kern0pt}def{\isacharparenright}{\kern0pt}\isanewline
\isacommand{apply}\isamarkupfalse%
{\isacharparenleft}{\kern0pt}auto{\isacharparenright}{\kern0pt}\isanewline
\isacommand{apply}\isamarkupfalse%
{\isacharparenleft}{\kern0pt}simp{\isacharunderscore}{\kern0pt}all\ add{\isacharcolon}{\kern0pt}\ swapas{\isacharunderscore}{\kern0pt}append{\isacharparenright}{\kern0pt}\isanewline
\isacommand{apply}\isamarkupfalse%
{\isacharparenleft}{\kern0pt}rule\ swapas{\isacharunderscore}{\kern0pt}pi{\isacharunderscore}{\kern0pt}ineq{\isacharunderscore}{\kern0pt}a{\isacharcomma}{\kern0pt}\ clarify{\isacharcomma}{\kern0pt}\ drule\ a{\isacharunderscore}{\kern0pt}not{\isacharunderscore}{\kern0pt}in{\isacharunderscore}{\kern0pt}atms{\isacharcomma}{\kern0pt}\ simp{\isacharparenright}{\kern0pt}{\isacharplus}{\kern0pt}\isanewline
\isacommand{apply}\isamarkupfalse%
{\isacharparenleft}{\kern0pt}drule\ swapas{\isacharunderscore}{\kern0pt}rev{\isacharunderscore}{\kern0pt}pi{\isacharunderscore}{\kern0pt}a{\isacharcomma}{\kern0pt}\ simp{\isacharparenright}{\kern0pt}{\isacharplus}{\kern0pt}\isanewline
\isacommand{done}\isamarkupfalse%
%
\endisatagproof
{\isafoldproof}%
%
\isadelimproof
\isanewline
%
\endisadelimproof
\isanewline
\isacommand{lemma}\isamarkupfalse%
\ ds{\isacharunderscore}{\kern0pt}rev{\isacharunderscore}{\kern0pt}pi{\isacharunderscore}{\kern0pt}pi{\isacharcolon}{\kern0pt}\ \isanewline
\ \ {\isachardoublequoteopen}ds\ {\isacharparenleft}{\kern0pt}{\isacharparenleft}{\kern0pt}rev\ pi{\isadigit{1}}{\isacharparenright}{\kern0pt}{\isacharat}{\kern0pt}pi{\isadigit{1}}{\isacharparenright}{\kern0pt}\ pi{\isadigit{2}}\ {\isacharequal}{\kern0pt}\ ds\ {\isacharbrackleft}{\kern0pt}{\isacharbrackright}{\kern0pt}\ pi{\isadigit{2}}{\isachardoublequoteclose}\isanewline
%
\isadelimproof
%
\endisadelimproof
%
\isatagproof
\isacommand{apply}\isamarkupfalse%
{\isacharparenleft}{\kern0pt}simp\ only{\isacharcolon}{\kern0pt}\ ds{\isacharunderscore}{\kern0pt}def{\isacharparenright}{\kern0pt}\isanewline
\isacommand{apply}\isamarkupfalse%
{\isacharparenleft}{\kern0pt}auto{\isacharparenright}{\kern0pt}\isanewline
\isacommand{apply}\isamarkupfalse%
{\isacharparenleft}{\kern0pt}simp{\isacharunderscore}{\kern0pt}all\ add{\isacharcolon}{\kern0pt}\ swapas{\isacharunderscore}{\kern0pt}append{\isacharparenright}{\kern0pt}\isanewline
\isacommand{apply}\isamarkupfalse%
{\isacharparenleft}{\kern0pt}drule\ a{\isacharunderscore}{\kern0pt}ineq{\isacharunderscore}{\kern0pt}swapas{\isacharunderscore}{\kern0pt}pi{\isacharcomma}{\kern0pt}\ assumption{\isacharparenright}{\kern0pt}{\isacharplus}{\kern0pt}\isanewline
\isacommand{done}\isamarkupfalse%
%
\endisatagproof
{\isafoldproof}%
%
\isadelimproof
\isanewline
%
\endisadelimproof
\isanewline
\isacommand{lemma}\isamarkupfalse%
\ ds{\isacharunderscore}{\kern0pt}rev{\isacharcolon}{\kern0pt}\ \isanewline
\ \ {\isachardoublequoteopen}ds\ {\isacharbrackleft}{\kern0pt}{\isacharbrackright}{\kern0pt}\ {\isacharparenleft}{\kern0pt}{\isacharparenleft}{\kern0pt}rev\ pi{\isadigit{1}}{\isacharparenright}{\kern0pt}{\isacharat}{\kern0pt}pi{\isadigit{2}}{\isacharparenright}{\kern0pt}\ {\isacharequal}{\kern0pt}\ ds\ pi{\isadigit{1}}\ pi{\isadigit{2}}{\isachardoublequoteclose}\isanewline
%
\isadelimproof
\ \ %
\endisadelimproof
%
\isatagproof
\isacommand{using}\isamarkupfalse%
\ ds{\isacharunderscore}{\kern0pt}cancel{\isacharunderscore}{\kern0pt}pi{\isacharunderscore}{\kern0pt}front\ ds{\isacharunderscore}{\kern0pt}rev{\isacharunderscore}{\kern0pt}pi{\isacharunderscore}{\kern0pt}pi\ \isacommand{by}\isamarkupfalse%
\ blast%
\endisatagproof
{\isafoldproof}%
%
\isadelimproof
\isanewline
%
\endisadelimproof
\isanewline
\isacommand{lemma}\isamarkupfalse%
\ ds{\isacharunderscore}{\kern0pt}acabbc{\isacharcolon}{\kern0pt}\ \isanewline
\ \ {\isachardoublequoteopen}{\isasymlbrakk}a{\isasymnoteq}b{\isacharsemicolon}{\kern0pt}b{\isasymnoteq}c{\isacharsemicolon}{\kern0pt}c{\isasymnoteq}a{\isasymrbrakk}{\isasymLongrightarrow}ds\ {\isacharbrackleft}{\kern0pt}{\isacharparenleft}{\kern0pt}a{\isacharcomma}{\kern0pt}\ b{\isacharparenright}{\kern0pt}{\isacharcomma}{\kern0pt}\ {\isacharparenleft}{\kern0pt}b{\isacharcomma}{\kern0pt}\ c{\isacharparenright}{\kern0pt}{\isacharbrackright}{\kern0pt}\ {\isacharbrackleft}{\kern0pt}{\isacharparenleft}{\kern0pt}a{\isacharcomma}{\kern0pt}\ c{\isacharparenright}{\kern0pt}{\isacharbrackright}{\kern0pt}\ {\isacharequal}{\kern0pt}\ {\isacharbraceleft}{\kern0pt}a{\isacharcomma}{\kern0pt}\ b{\isacharbraceright}{\kern0pt}{\isachardoublequoteclose}\isanewline
%
\isadelimproof
\ \ %
\endisadelimproof
%
\isatagproof
\isacommand{using}\isamarkupfalse%
\ ds{\isacharunderscore}{\kern0pt}def\ \isacommand{by}\isamarkupfalse%
\ auto%
\endisatagproof
{\isafoldproof}%
%
\isadelimproof
\isanewline
%
\endisadelimproof
\isanewline
\isacommand{lemma}\isamarkupfalse%
\ ds{\isacharunderscore}{\kern0pt}baab{\isacharcolon}{\kern0pt}\ \isanewline
\ \ {\isachardoublequoteopen}{\isasymlbrakk}a{\isasymnoteq}b{\isasymrbrakk}{\isasymLongrightarrow}ds\ {\isacharbrackleft}{\kern0pt}{\isacharparenleft}{\kern0pt}b{\isacharcomma}{\kern0pt}a{\isacharparenright}{\kern0pt}{\isacharbrackright}{\kern0pt}\ {\isacharbrackleft}{\kern0pt}{\isacharparenleft}{\kern0pt}a{\isacharcomma}{\kern0pt}\ b{\isacharparenright}{\kern0pt}{\isacharbrackright}{\kern0pt}\ {\isacharequal}{\kern0pt}\ {\isacharbraceleft}{\kern0pt}{\isacharbraceright}{\kern0pt}{\isachardoublequoteclose}\isanewline
%
\isadelimproof
\ \ %
\endisadelimproof
%
\isatagproof
\isacommand{using}\isamarkupfalse%
\ ds{\isacharunderscore}{\kern0pt}def\ \isacommand{by}\isamarkupfalse%
\ auto%
\endisatagproof
{\isafoldproof}%
%
\isadelimproof
\isanewline
%
\endisadelimproof
\isanewline
\isacommand{lemma}\isamarkupfalse%
\ ds{\isacharunderscore}{\kern0pt}baab{\isacharunderscore}{\kern0pt}id{\isacharcolon}{\kern0pt}\ \isanewline
{\isachardoublequoteopen}{\isasymlbrakk}a{\isasymnoteq}b{\isasymrbrakk}{\isasymLongrightarrow}ds\ {\isacharparenleft}{\kern0pt}{\isacharbrackleft}{\kern0pt}{\isacharparenleft}{\kern0pt}b{\isacharcomma}{\kern0pt}a{\isacharparenright}{\kern0pt}{\isacharbrackright}{\kern0pt}{\isacharat}{\kern0pt}{\isacharbrackleft}{\kern0pt}{\isacharparenleft}{\kern0pt}a{\isacharcomma}{\kern0pt}\ b{\isacharparenright}{\kern0pt}{\isacharbrackright}{\kern0pt}{\isacharparenright}{\kern0pt}\ {\isacharbrackleft}{\kern0pt}{\isacharbrackright}{\kern0pt}\ {\isacharequal}{\kern0pt}\ {\isacharbraceleft}{\kern0pt}{\isacharbraceright}{\kern0pt}{\isachardoublequoteclose}\isanewline
%
\isadelimproof
\ \ %
\endisadelimproof
%
\isatagproof
\isacommand{using}\isamarkupfalse%
\ ds{\isacharunderscore}{\kern0pt}def\ ds{\isacharunderscore}{\kern0pt}rev\ ds{\isacharunderscore}{\kern0pt}baab\ \isacommand{by}\isamarkupfalse%
\ simp%
\endisatagproof
{\isafoldproof}%
%
\isadelimproof
\isanewline
%
\endisadelimproof
\isanewline
\isacommand{lemma}\isamarkupfalse%
\ ds{\isacharunderscore}{\kern0pt}abab{\isacharcolon}{\kern0pt}\ \isanewline
\ \ {\isachardoublequoteopen}{\isasymlbrakk}a{\isasymnoteq}b{\isasymrbrakk}{\isasymLongrightarrow}ds\ {\isacharbrackleft}{\kern0pt}{\isacharbrackright}{\kern0pt}\ {\isacharbrackleft}{\kern0pt}{\isacharparenleft}{\kern0pt}a{\isacharcomma}{\kern0pt}\ b{\isacharparenright}{\kern0pt}{\isacharcomma}{\kern0pt}\ {\isacharparenleft}{\kern0pt}a{\isacharcomma}{\kern0pt}\ b{\isacharparenright}{\kern0pt}{\isacharbrackright}{\kern0pt}\ {\isacharequal}{\kern0pt}\ {\isacharbraceleft}{\kern0pt}{\isacharbraceright}{\kern0pt}{\isachardoublequoteclose}\isanewline
%
\isadelimproof
\ \ %
\endisadelimproof
%
\isatagproof
\isacommand{using}\isamarkupfalse%
\ ds{\isacharunderscore}{\kern0pt}def\ \isacommand{by}\isamarkupfalse%
\ auto%
\endisatagproof
{\isafoldproof}%
%
\isadelimproof
\isanewline
%
\endisadelimproof
\isanewline
\isacommand{lemma}\isamarkupfalse%
\ ds{\isacharunderscore}{\kern0pt}comm{\isacharcolon}{\kern0pt}\isanewline
{\isachardoublequoteopen}ds\ {\isacharparenleft}{\kern0pt}pi\ {\isacharat}{\kern0pt}\ {\isacharbrackleft}{\kern0pt}{\isacharparenleft}{\kern0pt}a{\isacharcomma}{\kern0pt}b{\isacharparenright}{\kern0pt}{\isacharbrackright}{\kern0pt}{\isacharparenright}{\kern0pt}\ {\isacharparenleft}{\kern0pt}{\isacharbrackleft}{\kern0pt}{\isacharparenleft}{\kern0pt}swapas\ pi\ a{\isacharcomma}{\kern0pt}\ swapas\ pi\ b{\isacharparenright}{\kern0pt}{\isacharbrackright}{\kern0pt}\ {\isacharat}{\kern0pt}\ pi{\isacharparenright}{\kern0pt}\ {\isacharequal}{\kern0pt}\ {\isacharbraceleft}{\kern0pt}{\isacharbraceright}{\kern0pt}{\isachardoublequoteclose}\isanewline
%
\isadelimproof
\ \ %
\endisadelimproof
%
\isatagproof
\isacommand{using}\isamarkupfalse%
\ swapas{\isacharunderscore}{\kern0pt}comm\ ds{\isacharunderscore}{\kern0pt}def\ \isacommand{by}\isamarkupfalse%
\ simp%
\endisatagproof
{\isafoldproof}%
%
\isadelimproof
\isanewline
%
\endisadelimproof
\isanewline
\isacommand{lemma}\isamarkupfalse%
\ ds{\isacharunderscore}{\kern0pt}rev{\isacharunderscore}{\kern0pt}pi{\isacharunderscore}{\kern0pt}id{\isacharcolon}{\kern0pt}\isanewline
{\isachardoublequoteopen}ds\ {\isacharparenleft}{\kern0pt}rev\ pi\ {\isacharat}{\kern0pt}\ pi{\isacharparenright}{\kern0pt}\ {\isacharbrackleft}{\kern0pt}{\isacharbrackright}{\kern0pt}\ {\isacharequal}{\kern0pt}\ {\isacharbraceleft}{\kern0pt}{\isacharbraceright}{\kern0pt}{\isachardoublequoteclose}\isanewline
%
\isadelimproof
\ \ %
\endisadelimproof
%
\isatagproof
\isacommand{using}\isamarkupfalse%
\ ds{\isacharunderscore}{\kern0pt}rev{\isacharunderscore}{\kern0pt}pi{\isacharunderscore}{\kern0pt}pi{\isacharbrackleft}{\kern0pt}of\ pi\ {\isacartoucheopen}{\isacharbrackleft}{\kern0pt}{\isacharbrackright}{\kern0pt}{\isacartoucheclose}{\isacharbrackright}{\kern0pt}\ elem{\isacharunderscore}{\kern0pt}ds\isanewline
\ \ \ \ ds{\isacharunderscore}{\kern0pt}sym{\isacharbrackleft}{\kern0pt}of\ {\isacartoucheopen}{\isacharparenleft}{\kern0pt}rev\ pi\ {\isacharat}{\kern0pt}\ pi{\isacharparenright}{\kern0pt}{\isacartoucheclose}\ {\isacartoucheopen}{\isacharbrackleft}{\kern0pt}{\isacharbrackright}{\kern0pt}{\isacartoucheclose}{\isacharbrackright}{\kern0pt}\ \isacommand{by}\isamarkupfalse%
\ fastforce%
\endisatagproof
{\isafoldproof}%
%
\isadelimproof
\isanewline
%
\endisadelimproof
\isanewline
\isacommand{lemma}\isamarkupfalse%
\ ds{\isacharunderscore}{\kern0pt}pi{\isacharunderscore}{\kern0pt}rev{\isacharunderscore}{\kern0pt}id{\isacharcolon}{\kern0pt}\isanewline
{\isachardoublequoteopen}ds\ {\isacharparenleft}{\kern0pt}pi\ {\isacharat}{\kern0pt}\ rev\ pi{\isacharparenright}{\kern0pt}\ {\isacharbrackleft}{\kern0pt}{\isacharbrackright}{\kern0pt}\ {\isacharequal}{\kern0pt}\ {\isacharbraceleft}{\kern0pt}{\isacharbraceright}{\kern0pt}{\isachardoublequoteclose}\isanewline
%
\isadelimproof
\ \ %
\endisadelimproof
%
\isatagproof
\isacommand{using}\isamarkupfalse%
\ ds{\isacharunderscore}{\kern0pt}rev{\isacharunderscore}{\kern0pt}pi{\isacharunderscore}{\kern0pt}id{\isacharbrackleft}{\kern0pt}of\ {\isacartoucheopen}rev\ pi{\isacartoucheclose}{\isacharbrackright}{\kern0pt}\ rev{\isacharunderscore}{\kern0pt}rev{\isacharunderscore}{\kern0pt}ident{\isacharbrackleft}{\kern0pt}of\ pi{\isacharbrackright}{\kern0pt}\isanewline
\ \ \isacommand{by}\isamarkupfalse%
\ simp%
\endisatagproof
{\isafoldproof}%
%
\isadelimproof
\isanewline
%
\endisadelimproof
\isanewline
\isacommand{lemma}\isamarkupfalse%
\ ds{\isacharunderscore}{\kern0pt}swapas{\isacharunderscore}{\kern0pt}eq{\isacharcolon}{\kern0pt}\isanewline
{\isachardoublequoteopen}ds\ pi{\isadigit{1}}\ pi{\isadigit{2}}\ {\isacharequal}{\kern0pt}\ {\isacharbraceleft}{\kern0pt}{\isacharbraceright}{\kern0pt}\ {\isasymLongrightarrow}\ swapas\ pi{\isadigit{1}}\ a\ {\isacharequal}{\kern0pt}\ swapas\ pi{\isadigit{2}}\ a{\isachardoublequoteclose}\isanewline
%
\isadelimproof
\ \ %
\endisadelimproof
%
\isatagproof
\isacommand{using}\isamarkupfalse%
\ ds{\isacharunderscore}{\kern0pt}elem{\isacharbrackleft}{\kern0pt}of\ {\isacartoucheopen}rev\ pi{\isadigit{1}}\ {\isacharat}{\kern0pt}\ pi{\isadigit{2}}{\isacartoucheclose}{\isacharbrackright}{\kern0pt}\ ds{\isacharunderscore}{\kern0pt}rev{\isacharbrackleft}{\kern0pt}of\ pi{\isadigit{1}}\ pi{\isadigit{2}}{\isacharbrackright}{\kern0pt}\ \isanewline
\ \ empty{\isacharunderscore}{\kern0pt}iff{\isacharbrackleft}{\kern0pt}of\ a{\isacharbrackright}{\kern0pt}\ swapas{\isacharunderscore}{\kern0pt}append{\isacharbrackleft}{\kern0pt}of\ {\isacartoucheopen}rev\ pi{\isadigit{1}}{\isacartoucheclose}\ pi{\isadigit{2}}\ a{\isacharbrackright}{\kern0pt}\ swapas{\isacharunderscore}{\kern0pt}inv\ \isacommand{by}\isamarkupfalse%
\ metis%
\endisatagproof
{\isafoldproof}%
%
\isadelimproof
\isanewline
%
\endisadelimproof
\isanewline
\isanewline
\isanewline
\isanewline
\ \isanewline
\isacommand{fun}\isamarkupfalse%
\ flatten\ {\isacharcolon}{\kern0pt}{\isacharcolon}{\kern0pt}\ {\isachardoublequoteopen}{\isacharparenleft}{\kern0pt}string\ {\isasymtimes}\ string{\isacharparenright}{\kern0pt}list\ {\isasymRightarrow}\ string\ list{\isachardoublequoteclose}\ \isakeyword{where}\isanewline
{\isachardoublequoteopen}flatten\ {\isacharbrackleft}{\kern0pt}{\isacharbrackright}{\kern0pt}\ \ \ \ \ {\isacharequal}{\kern0pt}\ {\isacharbrackleft}{\kern0pt}{\isacharbrackright}{\kern0pt}{\isachardoublequoteclose}\ {\isacharbar}{\kern0pt}\isanewline
{\isachardoublequoteopen}flatten\ {\isacharparenleft}{\kern0pt}x{\isacharhash}{\kern0pt}xs{\isacharparenright}{\kern0pt}\ {\isacharequal}{\kern0pt}\ {\isacharparenleft}{\kern0pt}fst\ x{\isacharparenright}{\kern0pt}{\isacharhash}{\kern0pt}{\isacharparenleft}{\kern0pt}snd\ x{\isacharparenright}{\kern0pt}{\isacharhash}{\kern0pt}{\isacharparenleft}{\kern0pt}flatten\ xs{\isacharparenright}{\kern0pt}{\isachardoublequoteclose}\isanewline
\isanewline
\isacommand{definition}\isamarkupfalse%
\ ds{\isacharunderscore}{\kern0pt}list\ {\isacharcolon}{\kern0pt}{\isacharcolon}{\kern0pt}\ {\isachardoublequoteopen}{\isacharparenleft}{\kern0pt}string\ {\isasymtimes}\ string{\isacharparenright}{\kern0pt}list\ {\isasymRightarrow}\ {\isacharparenleft}{\kern0pt}string\ {\isasymtimes}\ string{\isacharparenright}{\kern0pt}list\ {\isasymRightarrow}\ string\ list{\isachardoublequoteclose}\ \isakeyword{where}\isanewline
\ \ ds{\isacharunderscore}{\kern0pt}list{\isacharunderscore}{\kern0pt}def{\isacharcolon}{\kern0pt}\ {\isachardoublequoteopen}ds{\isacharunderscore}{\kern0pt}list\ pi{\isadigit{1}}\ pi{\isadigit{2}}\ {\isasymequiv}\ remdups\ {\isacharparenleft}{\kern0pt}{\isacharbrackleft}{\kern0pt}x{\isachardot}{\kern0pt}\ x\ {\isacharless}{\kern0pt}{\isacharminus}{\kern0pt}\ {\isacharparenleft}{\kern0pt}flatten\ {\isacharparenleft}{\kern0pt}pi{\isadigit{1}}{\isacharat}{\kern0pt}pi{\isadigit{2}}{\isacharparenright}{\kern0pt}{\isacharparenright}{\kern0pt}{\isacharcomma}{\kern0pt}\ x{\isasymin}ds\ pi{\isadigit{1}}\ pi{\isadigit{2}}{\isacharbrackright}{\kern0pt}{\isacharparenright}{\kern0pt}{\isachardoublequoteclose}\isanewline
\isanewline
\isanewline
\isacommand{lemma}\isamarkupfalse%
\ set{\isacharunderscore}{\kern0pt}flatten{\isacharunderscore}{\kern0pt}eq{\isacharunderscore}{\kern0pt}atms{\isacharcolon}{\kern0pt}\ \isanewline
\ \ {\isachardoublequoteopen}set\ {\isacharparenleft}{\kern0pt}flatten\ pi{\isacharparenright}{\kern0pt}\ {\isacharequal}{\kern0pt}\ atms\ pi{\isachardoublequoteclose}\isanewline
%
\isadelimproof
\ \ %
\endisadelimproof
%
\isatagproof
\isacommand{by}\isamarkupfalse%
\ {\isacharparenleft}{\kern0pt}induct\ pi{\isacharparenright}{\kern0pt}\ auto%
\endisatagproof
{\isafoldproof}%
%
\isadelimproof
\isanewline
%
\endisadelimproof
\isanewline
\isacommand{lemma}\isamarkupfalse%
\ ds{\isacharunderscore}{\kern0pt}list{\isacharunderscore}{\kern0pt}equ{\isacharunderscore}{\kern0pt}ds{\isacharcolon}{\kern0pt}\ \isanewline
\ \ {\isachardoublequoteopen}set\ {\isacharparenleft}{\kern0pt}ds{\isacharunderscore}{\kern0pt}list\ pi{\isadigit{1}}\ pi{\isadigit{2}}{\isacharparenright}{\kern0pt}\ {\isacharequal}{\kern0pt}\ ds\ pi{\isadigit{1}}\ pi{\isadigit{2}}{\isachardoublequoteclose}\isanewline
%
\isadelimproof
\ \ %
\endisadelimproof
%
\isatagproof
\isacommand{using}\isamarkupfalse%
\ ds{\isacharunderscore}{\kern0pt}list{\isacharunderscore}{\kern0pt}def\ ds{\isacharunderscore}{\kern0pt}def\ set{\isacharunderscore}{\kern0pt}flatten{\isacharunderscore}{\kern0pt}eq{\isacharunderscore}{\kern0pt}atms\ \isacommand{by}\isamarkupfalse%
\ auto%
\endisatagproof
{\isafoldproof}%
%
\isadelimproof
\isanewline
%
\endisadelimproof
\isanewline
\isanewline
%
\isadelimtheory
\isanewline
%
\endisadelimtheory
%
\isatagtheory
\isacommand{end}\isamarkupfalse%
%
\endisatagtheory
{\isafoldtheory}%
%
\isadelimtheory
%
\endisadelimtheory
%
\end{isabellebody}%
\endinput
%:%file=~/nominal_unification_Isabelle/nominal_alpha_unification/Disagreement.thy%:%
%:%6=1%:%
%:%11=2%:%
%:%12=2%:%
%:%13=3%:%
%:%14=4%:%
%:%15=5%:%
%:%16=6%:%
%:%21=6%:%
%:%24=7%:%
%:%25=11%:%
%:%26=12%:%
%:%27=13%:%
%:%28=13%:%
%:%29=14%:%
%:%30=15%:%
%:%31=16%:%
%:%32=17%:%
%:%33=18%:%
%:%34=18%:%
%:%35=19%:%
%:%38=20%:%
%:%42=20%:%
%:%43=20%:%
%:%44=20%:%
%:%49=20%:%
%:%52=21%:%
%:%53=22%:%
%:%54=23%:%
%:%55=23%:%
%:%58=24%:%
%:%62=24%:%
%:%63=24%:%
%:%64=24%:%
%:%69=24%:%
%:%72=25%:%
%:%73=26%:%
%:%74=26%:%
%:%75=27%:%
%:%78=28%:%
%:%82=28%:%
%:%83=28%:%
%:%84=28%:%
%:%89=28%:%
%:%92=29%:%
%:%93=30%:%
%:%94=31%:%
%:%95=31%:%
%:%96=32%:%
%:%99=33%:%
%:%103=33%:%
%:%104=33%:%
%:%105=33%:%
%:%110=33%:%
%:%113=34%:%
%:%114=35%:%
%:%115=36%:%
%:%116=36%:%
%:%117=37%:%
%:%124=38%:%
%:%125=38%:%
%:%126=38%:%
%:%131=38%:%
%:%134=39%:%
%:%135=40%:%
%:%136=41%:%
%:%137=42%:%
%:%138=42%:%
%:%139=43%:%
%:%140=44%:%
%:%143=45%:%
%:%147=45%:%
%:%148=45%:%
%:%149=46%:%
%:%150=46%:%
%:%155=46%:%
%:%158=47%:%
%:%159=48%:%
%:%160=49%:%
%:%161=50%:%
%:%162=50%:%
%:%163=51%:%
%:%170=52%:%
%:%171=52%:%
%:%172=53%:%
%:%173=53%:%
%:%174=54%:%
%:%175=54%:%
%:%176=55%:%
%:%177=55%:%
%:%178=56%:%
%:%179=57%:%
%:%185=57%:%
%:%188=58%:%
%:%189=59%:%
%:%190=60%:%
%:%191=61%:%
%:%192=61%:%
%:%193=62%:%
%:%196=63%:%
%:%200=63%:%
%:%201=63%:%
%:%202=63%:%
%:%207=63%:%
%:%210=64%:%
%:%211=65%:%
%:%212=66%:%
%:%213=66%:%
%:%214=67%:%
%:%217=68%:%
%:%221=68%:%
%:%222=68%:%
%:%223=69%:%
%:%224=70%:%
%:%225=70%:%
%:%230=70%:%
%:%233=71%:%
%:%234=72%:%
%:%235=73%:%
%:%236=74%:%
%:%237=74%:%
%:%238=75%:%
%:%245=76%:%
%:%246=76%:%
%:%247=77%:%
%:%248=77%:%
%:%249=78%:%
%:%250=78%:%
%:%251=79%:%
%:%252=79%:%
%:%253=80%:%
%:%254=80%:%
%:%255=81%:%
%:%261=81%:%
%:%264=82%:%
%:%265=83%:%
%:%266=83%:%
%:%267=84%:%
%:%274=85%:%
%:%275=85%:%
%:%276=86%:%
%:%277=86%:%
%:%278=87%:%
%:%279=87%:%
%:%280=88%:%
%:%281=88%:%
%:%282=89%:%
%:%288=89%:%
%:%291=90%:%
%:%292=91%:%
%:%293=91%:%
%:%294=92%:%
%:%297=93%:%
%:%301=93%:%
%:%302=93%:%
%:%303=93%:%
%:%308=93%:%
%:%311=94%:%
%:%312=95%:%
%:%313=95%:%
%:%314=96%:%
%:%317=97%:%
%:%321=97%:%
%:%322=97%:%
%:%323=97%:%
%:%328=97%:%
%:%331=98%:%
%:%332=99%:%
%:%333=99%:%
%:%334=100%:%
%:%337=101%:%
%:%341=101%:%
%:%342=101%:%
%:%343=101%:%
%:%348=101%:%
%:%351=102%:%
%:%352=103%:%
%:%353=103%:%
%:%354=104%:%
%:%357=105%:%
%:%361=105%:%
%:%362=105%:%
%:%363=105%:%
%:%368=105%:%
%:%371=106%:%
%:%372=107%:%
%:%373=107%:%
%:%374=108%:%
%:%377=109%:%
%:%381=109%:%
%:%382=109%:%
%:%383=109%:%
%:%388=109%:%
%:%391=110%:%
%:%392=111%:%
%:%393=111%:%
%:%394=112%:%
%:%397=113%:%
%:%401=113%:%
%:%402=113%:%
%:%403=113%:%
%:%408=113%:%
%:%411=114%:%
%:%412=115%:%
%:%413=115%:%
%:%414=116%:%
%:%417=117%:%
%:%421=117%:%
%:%422=117%:%
%:%423=118%:%
%:%424=118%:%
%:%429=118%:%
%:%432=119%:%
%:%433=120%:%
%:%434=120%:%
%:%435=121%:%
%:%438=122%:%
%:%442=122%:%
%:%443=122%:%
%:%444=123%:%
%:%445=123%:%
%:%450=123%:%
%:%453=124%:%
%:%454=125%:%
%:%455=125%:%
%:%456=126%:%
%:%459=127%:%
%:%463=127%:%
%:%464=127%:%
%:%465=128%:%
%:%466=128%:%
%:%471=128%:%
%:%474=129%:%
%:%475=130%:%
%:%476=131%:%
%:%477=132%:%
%:%478=133%:%
%:%479=134%:%
%:%480=134%:%
%:%481=135%:%
%:%482=136%:%
%:%483=137%:%
%:%484=138%:%
%:%485=138%:%
%:%486=139%:%
%:%487=140%:%
%:%488=141%:%
%:%489=142%:%
%:%490=142%:%
%:%491=143%:%
%:%494=144%:%
%:%498=144%:%
%:%499=144%:%
%:%504=144%:%
%:%507=145%:%
%:%508=146%:%
%:%509=146%:%
%:%510=147%:%
%:%513=148%:%
%:%517=148%:%
%:%518=148%:%
%:%519=148%:%
%:%524=148%:%
%:%527=149%:%
%:%528=150%:%
%:%531=151%:%
%:%536=152%:%

%
\begin{isabellebody}%
\setisabellecontext{Fresh}%
%
\isadelimtheory
%
\endisadelimtheory
%
\isatagtheory
\isacommand{theory}\isamarkupfalse%
\ Fresh\isanewline
\isanewline
\isakeyword{imports}\ Disagreement\isanewline
\isanewline
\isakeyword{begin}%
\endisatagtheory
{\isafoldtheory}%
%
\isadelimtheory
\isanewline
%
\endisadelimtheory
\isanewline
\isacommand{type{\isacharunderscore}{\kern0pt}synonym}\isamarkupfalse%
\ fresh{\isacharunderscore}{\kern0pt}envs\ {\isacharequal}{\kern0pt}\ {\isachardoublequoteopen}{\isacharparenleft}{\kern0pt}string\ {\isasymtimes}\ string{\isacharparenright}{\kern0pt}\ set{\isachardoublequoteclose}\isanewline
\isanewline
\isacommand{inductive}\isamarkupfalse%
\ fresh\ {\isacharcolon}{\kern0pt}{\isacharcolon}{\kern0pt}\ {\isachardoublequoteopen}fresh{\isacharunderscore}{\kern0pt}envs\ {\isasymRightarrow}\ string\ {\isasymRightarrow}\ trm\ {\isasymRightarrow}\ bool{\isachardoublequoteclose}\ {\isacharparenleft}{\kern0pt}{\isachardoublequoteopen}\ {\isacharunderscore}{\kern0pt}\ {\isasymturnstile}\ {\isacharunderscore}{\kern0pt}\ {\isasymsharp}\ {\isacharunderscore}{\kern0pt}{\isachardoublequoteclose}\ {\isacharbrackleft}{\kern0pt}{\isadigit{8}}{\isadigit{0}}{\isacharcomma}{\kern0pt}{\isadigit{8}}{\isadigit{0}}{\isacharcomma}{\kern0pt}{\isadigit{8}}{\isadigit{0}}{\isacharbrackright}{\kern0pt}\ {\isadigit{8}}{\isadigit{0}}{\isacharparenright}{\kern0pt}\ \isanewline
\ \ \isakeyword{where}\isanewline
\ \ fresh{\isacharunderscore}{\kern0pt}abst{\isacharunderscore}{\kern0pt}ab{\isacharbrackleft}{\kern0pt}intro{\isacharbang}{\kern0pt}{\isacharbrackright}{\kern0pt}{\isacharcolon}{\kern0pt}\ {\isachardoublequoteopen}{\isasymlbrakk}nabla\ {\isasymturnstile}\ a\ {\isasymsharp}\ t{\isacharsemicolon}{\kern0pt}\ a{\isasymnoteq}b{\isasymrbrakk}\ {\isasymLongrightarrow}\ nabla\ {\isasymturnstile}\ a\ {\isasymsharp}\ Abst\ b\ t{\isachardoublequoteclose}\ {\isacharbar}{\kern0pt}\isanewline
\ \ fresh{\isacharunderscore}{\kern0pt}abst{\isacharunderscore}{\kern0pt}aa{\isacharbrackleft}{\kern0pt}intro{\isacharbang}{\kern0pt}{\isacharbrackright}{\kern0pt}{\isacharcolon}{\kern0pt}\ {\isachardoublequoteopen}nabla\ {\isasymturnstile}\ a\ {\isasymsharp}\ Abst\ a\ t{\isachardoublequoteclose}\ {\isacharbar}{\kern0pt}\isanewline
\ \ fresh{\isacharunderscore}{\kern0pt}unit{\isacharbrackleft}{\kern0pt}intro{\isacharbang}{\kern0pt}{\isacharbrackright}{\kern0pt}{\isacharcolon}{\kern0pt}\ \ \ \ {\isachardoublequoteopen}nabla\ {\isasymturnstile}\ a\ {\isasymsharp}\ Unit{\isachardoublequoteclose}\ {\isacharbar}{\kern0pt}\isanewline
\ \ fresh{\isacharunderscore}{\kern0pt}atom{\isacharbrackleft}{\kern0pt}intro{\isacharbang}{\kern0pt}{\isacharbrackright}{\kern0pt}{\isacharcolon}{\kern0pt}\ \ \ \ {\isachardoublequoteopen}a\ {\isasymnoteq}\ b\ {\isasymLongrightarrow}\ nabla\ {\isasymturnstile}\ a\ {\isasymsharp}\ Atom\ b{\isachardoublequoteclose}\ {\isacharbar}{\kern0pt}\isanewline
\ \ fresh{\isacharunderscore}{\kern0pt}susp{\isacharbrackleft}{\kern0pt}intro{\isacharbang}{\kern0pt}{\isacharbrackright}{\kern0pt}{\isacharcolon}{\kern0pt}\ \ \ \ {\isachardoublequoteopen}{\isacharparenleft}{\kern0pt}swapas\ {\isacharparenleft}{\kern0pt}rev\ pi{\isacharparenright}{\kern0pt}\ a{\isacharcomma}{\kern0pt}X{\isacharparenright}{\kern0pt}\ {\isasymin}\ nabla\ {\isasymLongrightarrow}\ nabla\ {\isasymturnstile}\ a\ {\isasymsharp}\ Susp\ pi\ X{\isachardoublequoteclose}\ {\isacharbar}{\kern0pt}\isanewline
\ \ fresh{\isacharunderscore}{\kern0pt}paar{\isacharbrackleft}{\kern0pt}intro{\isacharbang}{\kern0pt}{\isacharbrackright}{\kern0pt}{\isacharcolon}{\kern0pt}\ \ \ \ {\isachardoublequoteopen}{\isasymlbrakk}nabla\ {\isasymturnstile}\ a\ {\isasymsharp}\ t{\isadigit{1}}{\isacharsemicolon}{\kern0pt}\ nabla\ {\isasymturnstile}\ a\ {\isasymsharp}\ t{\isadigit{2}}{\isasymrbrakk}\ {\isasymLongrightarrow}\ nabla\ {\isasymturnstile}\ a\ {\isasymsharp}\ Paar\ t{\isadigit{1}}\ t{\isadigit{2}}{\isachardoublequoteclose}\ {\isacharbar}{\kern0pt}\isanewline
\ \ fresh{\isacharunderscore}{\kern0pt}func{\isacharbrackleft}{\kern0pt}intro{\isacharbang}{\kern0pt}{\isacharbrackright}{\kern0pt}{\isacharcolon}{\kern0pt}\ \ \ \ {\isachardoublequoteopen}nabla\ {\isasymturnstile}\ a\ {\isasymsharp}\ t\ {\isasymLongrightarrow}\ nabla\ {\isasymturnstile}\ a\ {\isasymsharp}\ Func\ F\ t{\isachardoublequoteclose}\isanewline
\isanewline
\isacommand{inductive{\isacharunderscore}{\kern0pt}cases}\isamarkupfalse%
\ Fresh{\isacharunderscore}{\kern0pt}elims{\isacharcolon}{\kern0pt}\isanewline
\ \ {\isachardoublequoteopen}nabla\ {\isasymturnstile}\ a\ {\isasymsharp}\ Abst\ b\ t{\isachardoublequoteclose}\isanewline
\ \ {\isachardoublequoteopen}nabla\ {\isasymturnstile}\ a\ {\isasymsharp}\ Unit{\isachardoublequoteclose}\isanewline
\ \ {\isachardoublequoteopen}nabla\ {\isasymturnstile}\ a\ {\isasymsharp}\ Atom\ b{\isachardoublequoteclose}\isanewline
\ \ {\isachardoublequoteopen}nabla\ {\isasymturnstile}\ a\ {\isasymsharp}\ Susp\ pi\ X{\isachardoublequoteclose}\isanewline
\ \ {\isachardoublequoteopen}nabla\ {\isasymturnstile}\ a\ {\isasymsharp}\ Paar\ t{\isadigit{1}}\ t{\isadigit{2}}{\isachardoublequoteclose}\isanewline
\ \ {\isachardoublequoteopen}nabla\ {\isasymturnstile}\ a\ {\isasymsharp}\ Func\ F\ t{\isachardoublequoteclose}\isanewline
\isanewline
\isacommand{lemma}\isamarkupfalse%
\ fresh{\isacharunderscore}{\kern0pt}swap{\isacharunderscore}{\kern0pt}eqvt{\isacharcolon}{\kern0pt}\ \isanewline
\ \ \isakeyword{assumes}\ {\isachardoublequoteopen}nabla\ {\isasymturnstile}\ a\ {\isasymsharp}\ t{\isachardoublequoteclose}\isanewline
\ \ \isakeyword{shows}\ {\isachardoublequoteopen}nabla\ {\isasymturnstile}\ swapas\ pi\ a\ {\isasymsharp}\ swap\ pi\ t{\isachardoublequoteclose}\isanewline
%
\isadelimproof
\ \ %
\endisadelimproof
%
\isatagproof
\isacommand{using}\isamarkupfalse%
\ assms\isanewline
\ \ \isacommand{by}\isamarkupfalse%
\ {\isacharparenleft}{\kern0pt}induct\ arbitrary{\isacharcolon}{\kern0pt}\ pi\ rule{\isacharcolon}{\kern0pt}\ fresh{\isachardot}{\kern0pt}induct{\isacharparenright}{\kern0pt}\isanewline
{\isacharparenleft}{\kern0pt}auto\ simp\ add{\isacharcolon}{\kern0pt}\ swapas{\isacharunderscore}{\kern0pt}append\ dest{\isacharcolon}{\kern0pt}\ swapas{\isacharunderscore}{\kern0pt}rev{\isacharunderscore}{\kern0pt}pi{\isacharunderscore}{\kern0pt}a{\isacharparenright}{\kern0pt}%
\endisatagproof
{\isafoldproof}%
%
\isadelimproof
\isanewline
%
\endisadelimproof
\isanewline
\isanewline
\isacommand{lemma}\isamarkupfalse%
\ fresh{\isacharunderscore}{\kern0pt}swap{\isacharunderscore}{\kern0pt}left{\isacharcolon}{\kern0pt}\ \isanewline
\ \ \isakeyword{assumes}\ {\isachardoublequoteopen}nabla\ {\isasymturnstile}\ a\ {\isasymsharp}\ swap\ pi\ t{\isachardoublequoteclose}\isanewline
\ \ \isakeyword{shows}\ {\isachardoublequoteopen}nabla\ {\isasymturnstile}\ swapas\ {\isacharparenleft}{\kern0pt}rev\ pi{\isacharparenright}{\kern0pt}\ a\ {\isasymsharp}\ t{\isachardoublequoteclose}\isanewline
%
\isadelimproof
\ \ %
\endisadelimproof
%
\isatagproof
\isacommand{using}\isamarkupfalse%
\ assms\isanewline
\isacommand{proof}\isamarkupfalse%
{\isacharparenleft}{\kern0pt}induct\ t\ arbitrary{\isacharcolon}{\kern0pt}\ a\ pi{\isacharparenright}{\kern0pt}\isanewline
\ \ \isacommand{case}\isamarkupfalse%
\ {\isacharparenleft}{\kern0pt}Abst\ x{\isadigit{1}}\ t{\isacharparenright}{\kern0pt}\isanewline
\ \ \isacommand{have}\isamarkupfalse%
\ {\isachardoublequoteopen}nabla\ {\isasymturnstile}\ a\ {\isasymsharp}\ swap\ pi\ {\isacharparenleft}{\kern0pt}Abst\ x{\isadigit{1}}\ t{\isacharparenright}{\kern0pt}{\isachardoublequoteclose}\ \isacommand{by}\isamarkupfalse%
\ fact\isanewline
\ \ \isacommand{then}\isamarkupfalse%
\ \isacommand{have}\isamarkupfalse%
\ {\isachardoublequoteopen}{\isacharparenleft}{\kern0pt}nabla\ {\isasymturnstile}\ a\ {\isasymsharp}\ swap\ pi\ t\ {\isasymand}\ a\ {\isasymnoteq}\ swapas\ pi\ x{\isadigit{1}}{\isacharparenright}{\kern0pt}\ {\isasymor}\ {\isacharparenleft}{\kern0pt}a\ {\isacharequal}{\kern0pt}\ swapas\ pi\ x{\isadigit{1}}{\isacharparenright}{\kern0pt}{\isachardoublequoteclose}\isanewline
\ \ \ \ \isacommand{by}\isamarkupfalse%
\ {\isacharparenleft}{\kern0pt}metis\ Fresh{\isacharunderscore}{\kern0pt}elims{\isacharparenleft}{\kern0pt}{\isadigit{1}}{\isacharparenright}{\kern0pt}\ swap{\isachardot}{\kern0pt}simps{\isacharparenleft}{\kern0pt}{\isadigit{4}}{\isacharparenright}{\kern0pt}{\isacharparenright}{\kern0pt}\isanewline
\ \ \isacommand{moreover}\isamarkupfalse%
\ \isacommand{{\isacharbraceleft}{\kern0pt}}\isamarkupfalse%
\isanewline
\ \ \ \ \isacommand{assume}\isamarkupfalse%
\ as{\isacharcolon}{\kern0pt}\ {\isachardoublequoteopen}nabla\ {\isasymturnstile}\ a\ {\isasymsharp}\ swap\ pi\ t\ {\isasymand}\ a\ {\isasymnoteq}\ swapas\ pi\ x{\isadigit{1}}{\isachardoublequoteclose}\ \isanewline
\ \ \ \ \isacommand{have}\isamarkupfalse%
\ {\isachardoublequoteopen}nabla\ {\isasymturnstile}\ swapas\ {\isacharparenleft}{\kern0pt}rev\ pi{\isacharparenright}{\kern0pt}\ a\ {\isasymsharp}\ Abst\ x{\isadigit{1}}\ t{\isachardoublequoteclose}\isanewline
\ \ \ \ \ \ \isacommand{using}\isamarkupfalse%
\ Abst{\isachardot}{\kern0pt}hyps\ as\ \isacommand{by}\isamarkupfalse%
\ auto\ \isanewline
\ \ \isacommand{{\isacharbraceright}{\kern0pt}}\isamarkupfalse%
\isanewline
\ \ \isacommand{moreover}\isamarkupfalse%
\ \isacommand{{\isacharbraceleft}{\kern0pt}}\isamarkupfalse%
\isanewline
\ \ \ \ \isacommand{assume}\isamarkupfalse%
\ as{\isacharcolon}{\kern0pt}\ {\isachardoublequoteopen}a\ {\isacharequal}{\kern0pt}\ swapas\ pi\ x{\isadigit{1}}{\isachardoublequoteclose}\isanewline
\ \ \ \ \isacommand{then}\isamarkupfalse%
\ \isacommand{have}\isamarkupfalse%
\ {\isachardoublequoteopen}nabla\ {\isasymturnstile}\ swapas\ {\isacharparenleft}{\kern0pt}rev\ pi{\isacharparenright}{\kern0pt}\ a\ {\isasymsharp}\ Abst\ x{\isadigit{1}}\ t{\isachardoublequoteclose}\isanewline
\ \ \ \ \ \ \isacommand{by}\isamarkupfalse%
\ {\isacharparenleft}{\kern0pt}simp\ add{\isacharcolon}{\kern0pt}\ fresh{\isacharunderscore}{\kern0pt}abst{\isacharunderscore}{\kern0pt}aa{\isacharparenright}{\kern0pt}\ \ \ \ \ \ \isanewline
\ \ \isacommand{{\isacharbraceright}{\kern0pt}}\isamarkupfalse%
\ \isanewline
\ \ \isacommand{ultimately}\isamarkupfalse%
\ \isacommand{show}\isamarkupfalse%
\ {\isachardoublequoteopen}nabla\ {\isasymturnstile}\ swapas\ {\isacharparenleft}{\kern0pt}rev\ pi{\isacharparenright}{\kern0pt}\ a\ {\isasymsharp}\ Abst\ x{\isadigit{1}}\ t{\isachardoublequoteclose}\ \isanewline
\ \ \ \ \isacommand{by}\isamarkupfalse%
\ argo\isanewline
\isacommand{next}\isamarkupfalse%
\isanewline
\ \ \isacommand{case}\isamarkupfalse%
\ {\isacharparenleft}{\kern0pt}Susp\ x{\isadigit{1}}\ x{\isadigit{2}}{\isacharparenright}{\kern0pt}\isanewline
\ \ \isacommand{have}\isamarkupfalse%
\ {\isachardoublequoteopen}nabla\ {\isasymturnstile}\ a\ {\isasymsharp}\ swap\ pi\ {\isacharparenleft}{\kern0pt}Susp\ x{\isadigit{1}}\ x{\isadigit{2}}{\isacharparenright}{\kern0pt}{\isachardoublequoteclose}\ \isacommand{by}\isamarkupfalse%
\ fact\isanewline
\ \ \isacommand{have}\isamarkupfalse%
\ {\isachardoublequoteopen}{\isacharparenleft}{\kern0pt}swapas\ {\isacharparenleft}{\kern0pt}rev\ x{\isadigit{1}}{\isacharparenright}{\kern0pt}\ {\isacharparenleft}{\kern0pt}swapas\ {\isacharparenleft}{\kern0pt}rev\ pi{\isacharparenright}{\kern0pt}\ a{\isacharparenright}{\kern0pt}{\isacharcomma}{\kern0pt}\ x{\isadigit{2}}{\isacharparenright}{\kern0pt}\ {\isasymin}\ nabla{\isachardoublequoteclose}\isanewline
\ \ \ \ \isacommand{by}\isamarkupfalse%
\ {\isacharparenleft}{\kern0pt}metis\ Fresh{\isacharunderscore}{\kern0pt}elims{\isacharparenleft}{\kern0pt}{\isadigit{4}}{\isacharparenright}{\kern0pt}\ Susp\ rev{\isacharunderscore}{\kern0pt}append\ swap{\isachardot}{\kern0pt}simps{\isacharparenleft}{\kern0pt}{\isadigit{3}}{\isacharparenright}{\kern0pt}\ swapas{\isacharunderscore}{\kern0pt}append{\isacharparenright}{\kern0pt}\ \isanewline
\ \ \isacommand{then}\isamarkupfalse%
\ \isacommand{show}\isamarkupfalse%
\ {\isachardoublequoteopen}nabla\ {\isasymturnstile}\ swapas\ {\isacharparenleft}{\kern0pt}rev\ pi{\isacharparenright}{\kern0pt}\ a\ {\isasymsharp}\ Susp\ x{\isadigit{1}}\ x{\isadigit{2}}{\isachardoublequoteclose}\isanewline
\ \ \ \ \isacommand{by}\isamarkupfalse%
\ {\isacharparenleft}{\kern0pt}simp\ add{\isacharcolon}{\kern0pt}\ fresh{\isacharunderscore}{\kern0pt}susp{\isacharparenright}{\kern0pt}\isanewline
\isacommand{next}\isamarkupfalse%
\isanewline
\ \ \isacommand{case}\isamarkupfalse%
\ {\isacharparenleft}{\kern0pt}Atom\ x{\isacharparenright}{\kern0pt}\isanewline
\ \ \isacommand{have}\isamarkupfalse%
\ {\isachardoublequoteopen}nabla\ {\isasymturnstile}\ a\ {\isasymsharp}\ swap\ pi\ {\isacharparenleft}{\kern0pt}Atom\ x{\isacharparenright}{\kern0pt}{\isachardoublequoteclose}\ \isacommand{by}\isamarkupfalse%
\ fact\isanewline
\ \ \isacommand{then}\isamarkupfalse%
\ \isacommand{have}\isamarkupfalse%
\ {\isachardoublequoteopen}swapas\ pi\ x\ {\isasymnoteq}\ a{\isachardoublequoteclose}\ \isacommand{by}\isamarkupfalse%
\ {\isacharparenleft}{\kern0pt}auto\ elim{\isacharcolon}{\kern0pt}\ Fresh{\isacharunderscore}{\kern0pt}elims{\isacharparenright}{\kern0pt}\isanewline
\ \ \isacommand{then}\isamarkupfalse%
\ \isacommand{have}\isamarkupfalse%
\ {\isachardoublequoteopen}x\ {\isasymnoteq}\ swapas\ {\isacharparenleft}{\kern0pt}rev\ pi{\isacharparenright}{\kern0pt}\ a{\isachardoublequoteclose}\ \isacommand{using}\isamarkupfalse%
\ swapas{\isacharunderscore}{\kern0pt}rev{\isacharunderscore}{\kern0pt}pi{\isacharunderscore}{\kern0pt}b\ \isacommand{by}\isamarkupfalse%
\ blast\ \ \isanewline
\ \ \isacommand{then}\isamarkupfalse%
\ \isacommand{show}\isamarkupfalse%
\ {\isachardoublequoteopen}nabla\ {\isasymturnstile}\ swapas\ {\isacharparenleft}{\kern0pt}rev\ pi{\isacharparenright}{\kern0pt}\ a\ {\isasymsharp}\ Atom\ x{\isachardoublequoteclose}\isanewline
\ \ \ \ \isacommand{by}\isamarkupfalse%
\ {\isacharparenleft}{\kern0pt}simp\ add{\isacharcolon}{\kern0pt}\ fresh{\isacharunderscore}{\kern0pt}atom{\isacharparenright}{\kern0pt}\ \isanewline
\isacommand{qed}\isamarkupfalse%
\ {\isacharparenleft}{\kern0pt}auto\ elim{\isacharcolon}{\kern0pt}\ Fresh{\isacharunderscore}{\kern0pt}elims{\isacharparenright}{\kern0pt}%
\endisatagproof
{\isafoldproof}%
%
\isadelimproof
\isanewline
%
\endisadelimproof
\isanewline
\isanewline
\isanewline
\isacommand{lemma}\isamarkupfalse%
\ fresh{\isacharunderscore}{\kern0pt}swap{\isacharunderscore}{\kern0pt}right{\isacharcolon}{\kern0pt}\ \isanewline
\ \ \isakeyword{assumes}\ {\isachardoublequoteopen}nabla\ {\isasymturnstile}\ swapas\ {\isacharparenleft}{\kern0pt}rev\ pi{\isacharparenright}{\kern0pt}\ a\ {\isasymsharp}\ t{\isachardoublequoteclose}\isanewline
\ \ \isakeyword{shows}\ {\isachardoublequoteopen}nabla\ {\isasymturnstile}\ a\ {\isasymsharp}\ swap\ pi\ t{\isachardoublequoteclose}\isanewline
%
\isadelimproof
\ \ %
\endisadelimproof
%
\isatagproof
\isacommand{using}\isamarkupfalse%
\ assms\ fresh{\isacharunderscore}{\kern0pt}swap{\isacharunderscore}{\kern0pt}eqvt{\isacharbrackleft}{\kern0pt}of\ nabla\ {\isacartoucheopen}swapas\ {\isacharparenleft}{\kern0pt}rev\ pi{\isacharparenright}{\kern0pt}\ a{\isacartoucheclose}\ t\ pi{\isacharbrackright}{\kern0pt}\ \isanewline
\ \ \ \ swapas{\isacharunderscore}{\kern0pt}rev{\isacharunderscore}{\kern0pt}swapas{\isacharunderscore}{\kern0pt}id{\isacharbrackleft}{\kern0pt}of\ pi\ a{\isacharbrackright}{\kern0pt}\ \isacommand{by}\isamarkupfalse%
\ simp%
\endisatagproof
{\isafoldproof}%
%
\isadelimproof
\isanewline
%
\endisadelimproof
\isanewline
\isacommand{lemma}\isamarkupfalse%
\ fresh{\isacharunderscore}{\kern0pt}weak{\isacharcolon}{\kern0pt}\ \isanewline
\ \ \isakeyword{assumes}\ {\isachardoublequoteopen}nabla{\isadigit{1}}\ {\isasymturnstile}\ a\ {\isasymsharp}\ t{\isachardoublequoteclose}\isanewline
\ \ \isakeyword{shows}\ {\isachardoublequoteopen}{\isacharparenleft}{\kern0pt}nabla{\isadigit{1}}\ {\isasymunion}\ nabla{\isadigit{2}}{\isacharparenright}{\kern0pt}\ {\isasymturnstile}\ a\ {\isasymsharp}\ t{\isachardoublequoteclose}\isanewline
%
\isadelimproof
\ \ %
\endisadelimproof
%
\isatagproof
\isacommand{using}\isamarkupfalse%
\ assms\isanewline
\ \ \isacommand{by}\isamarkupfalse%
{\isacharparenleft}{\kern0pt}induct\ rule{\isacharcolon}{\kern0pt}\ fresh{\isachardot}{\kern0pt}induct{\isacharparenright}{\kern0pt}{\isacharparenleft}{\kern0pt}auto{\isacharparenright}{\kern0pt}%
\endisatagproof
{\isafoldproof}%
%
\isadelimproof
\isanewline
%
\endisadelimproof
\isanewline
\isacommand{lemma}\isamarkupfalse%
\ ds{\isacharunderscore}{\kern0pt}empty{\isacharunderscore}{\kern0pt}fresh{\isacharunderscore}{\kern0pt}{\isadigit{1}}{\isacharcolon}{\kern0pt}\isanewline
\ \ \isakeyword{assumes}\ {\isachardoublequoteopen}ds\ pi{\isadigit{1}}\ pi{\isadigit{2}}\ {\isacharequal}{\kern0pt}\ {\isacharbraceleft}{\kern0pt}{\isacharbraceright}{\kern0pt}{\isachardoublequoteclose}\isanewline
\ \ \isakeyword{shows}\ {\isachardoublequoteopen}nabla\ {\isasymturnstile}\ swapas\ pi{\isadigit{1}}\ a\ {\isasymsharp}\ t\ {\isasymLongrightarrow}\ nabla\ {\isasymturnstile}\ swapas\ pi{\isadigit{2}}\ a\ {\isasymsharp}\ t{\isachardoublequoteclose}\isanewline
%
\isadelimproof
\ \ %
\endisadelimproof
%
\isatagproof
\isacommand{using}\isamarkupfalse%
\ assms\ ds{\isacharunderscore}{\kern0pt}swapas{\isacharunderscore}{\kern0pt}eq\ \isacommand{by}\isamarkupfalse%
\ simp%
\endisatagproof
{\isafoldproof}%
%
\isadelimproof
\isanewline
%
\endisadelimproof
\isanewline
\isacommand{lemma}\isamarkupfalse%
\ ds{\isacharunderscore}{\kern0pt}empty{\isacharunderscore}{\kern0pt}fresh{\isacharunderscore}{\kern0pt}{\isadigit{2}}{\isacharcolon}{\kern0pt}\isanewline
\ \ \isakeyword{assumes}\ {\isachardoublequoteopen}ds\ pi{\isadigit{1}}\ pi{\isadigit{2}}\ {\isacharequal}{\kern0pt}\ {\isacharbraceleft}{\kern0pt}{\isacharbraceright}{\kern0pt}{\isachardoublequoteclose}\isanewline
\ \ \isakeyword{shows}\ {\isachardoublequoteopen}nabla\ {\isasymturnstile}\ a\ {\isasymsharp}\ swap\ pi{\isadigit{1}}\ t\ {\isasymLongrightarrow}\ nabla\ {\isasymturnstile}\ a\ {\isasymsharp}\ swap\ pi{\isadigit{2}}\ t{\isachardoublequoteclose}\isanewline
%
\isadelimproof
\ \ %
\endisadelimproof
%
\isatagproof
\isacommand{using}\isamarkupfalse%
\ assms\ \isanewline
\isacommand{proof}\isamarkupfalse%
{\isacharminus}{\kern0pt}\isanewline
\ \ \isacommand{assume}\isamarkupfalse%
\ assm{\isadigit{1}}{\isacharcolon}{\kern0pt}\ {\isachardoublequoteopen}nabla\ {\isasymturnstile}\ a\ {\isasymsharp}\ swap\ pi{\isadigit{1}}\ t{\isachardoublequoteclose}\isanewline
\ \ \isacommand{have}\isamarkupfalse%
\ i{\isacharcolon}{\kern0pt}\ {\isachardoublequoteopen}swapas\ {\isacharparenleft}{\kern0pt}rev\ pi{\isadigit{1}}{\isacharparenright}{\kern0pt}\ a\ {\isacharequal}{\kern0pt}\ swapas\ {\isacharparenleft}{\kern0pt}rev\ pi{\isadigit{2}}{\isacharparenright}{\kern0pt}\ a{\isachardoublequoteclose}\isanewline
\ \ \ \ \isacommand{using}\isamarkupfalse%
\ assms\ ds{\isacharunderscore}{\kern0pt}swapas{\isacharunderscore}{\kern0pt}eq\ swapas{\isacharunderscore}{\kern0pt}rev{\isacharunderscore}{\kern0pt}pi{\isacharunderscore}{\kern0pt}a\ \isacommand{by}\isamarkupfalse%
\ metis\isanewline
\ \ \isacommand{with}\isamarkupfalse%
\ assms\ \isacommand{have}\isamarkupfalse%
\ ii{\isacharcolon}{\kern0pt}\ {\isachardoublequoteopen}nabla\ {\isasymturnstile}\ swapas\ {\isacharparenleft}{\kern0pt}rev\ pi{\isadigit{2}}{\isacharparenright}{\kern0pt}\ a\ {\isasymsharp}\ t{\isachardoublequoteclose}\isanewline
\ \ \ \ \isacommand{using}\isamarkupfalse%
\ fresh{\isacharunderscore}{\kern0pt}swap{\isacharunderscore}{\kern0pt}left{\isacharbrackleft}{\kern0pt}OF\ assm{\isadigit{1}}{\isacharbrackright}{\kern0pt}\ i\ \isacommand{by}\isamarkupfalse%
\ simp\isanewline
\ \ \isacommand{show}\isamarkupfalse%
\ {\isacharquery}{\kern0pt}thesis\ \isanewline
\ \ \ \ \isacommand{using}\isamarkupfalse%
\ fresh{\isacharunderscore}{\kern0pt}swap{\isacharunderscore}{\kern0pt}right{\isacharbrackleft}{\kern0pt}OF\ ii{\isacharbrackright}{\kern0pt}\ \isacommand{by}\isamarkupfalse%
\ simp\isanewline
\isacommand{qed}\isamarkupfalse%
%
\endisatagproof
{\isafoldproof}%
%
\isadelimproof
\isanewline
%
\endisadelimproof
%
\isadelimtheory
\isanewline
%
\endisadelimtheory
%
\isatagtheory
\isacommand{end}\isamarkupfalse%
%
\endisatagtheory
{\isafoldtheory}%
%
\isadelimtheory
%
\endisadelimtheory
%
\end{isabellebody}%
\endinput
%:%file=~/nominal_unification_Isabelle/nominal_alpha_unification/Fresh.thy%:%
%:%10=1%:%
%:%11=1%:%
%:%12=2%:%
%:%13=3%:%
%:%14=4%:%
%:%15=5%:%
%:%20=5%:%
%:%23=6%:%
%:%24=7%:%
%:%25=7%:%
%:%26=8%:%
%:%27=9%:%
%:%28=9%:%
%:%29=10%:%
%:%30=11%:%
%:%31=12%:%
%:%32=13%:%
%:%33=14%:%
%:%34=15%:%
%:%35=16%:%
%:%36=17%:%
%:%37=18%:%
%:%38=19%:%
%:%39=19%:%
%:%40=20%:%
%:%41=21%:%
%:%42=22%:%
%:%43=23%:%
%:%44=24%:%
%:%45=25%:%
%:%46=26%:%
%:%47=27%:%
%:%48=27%:%
%:%49=28%:%
%:%50=29%:%
%:%53=30%:%
%:%57=30%:%
%:%58=30%:%
%:%59=31%:%
%:%60=31%:%
%:%61=32%:%
%:%66=32%:%
%:%69=33%:%
%:%70=34%:%
%:%71=35%:%
%:%72=35%:%
%:%73=36%:%
%:%74=37%:%
%:%77=38%:%
%:%81=38%:%
%:%82=38%:%
%:%83=39%:%
%:%84=39%:%
%:%85=40%:%
%:%86=40%:%
%:%87=41%:%
%:%88=41%:%
%:%89=41%:%
%:%90=42%:%
%:%91=42%:%
%:%92=42%:%
%:%93=43%:%
%:%94=43%:%
%:%95=44%:%
%:%96=44%:%
%:%97=44%:%
%:%98=45%:%
%:%99=45%:%
%:%100=46%:%
%:%101=46%:%
%:%102=47%:%
%:%103=47%:%
%:%104=47%:%
%:%105=48%:%
%:%106=48%:%
%:%107=49%:%
%:%108=49%:%
%:%109=49%:%
%:%110=50%:%
%:%111=50%:%
%:%112=51%:%
%:%113=51%:%
%:%114=51%:%
%:%115=52%:%
%:%116=52%:%
%:%117=53%:%
%:%118=53%:%
%:%119=54%:%
%:%120=54%:%
%:%121=54%:%
%:%122=55%:%
%:%123=55%:%
%:%124=56%:%
%:%125=56%:%
%:%126=57%:%
%:%127=57%:%
%:%128=58%:%
%:%129=58%:%
%:%130=58%:%
%:%131=59%:%
%:%132=59%:%
%:%133=60%:%
%:%134=60%:%
%:%135=61%:%
%:%136=61%:%
%:%137=61%:%
%:%138=62%:%
%:%139=62%:%
%:%140=63%:%
%:%141=63%:%
%:%142=64%:%
%:%143=64%:%
%:%144=65%:%
%:%145=65%:%
%:%146=65%:%
%:%147=66%:%
%:%148=66%:%
%:%149=66%:%
%:%150=66%:%
%:%151=67%:%
%:%152=67%:%
%:%153=67%:%
%:%154=67%:%
%:%155=67%:%
%:%156=68%:%
%:%157=68%:%
%:%158=68%:%
%:%159=69%:%
%:%160=69%:%
%:%161=70%:%
%:%162=70%:%
%:%167=70%:%
%:%170=71%:%
%:%171=72%:%
%:%172=73%:%
%:%173=74%:%
%:%174=74%:%
%:%175=75%:%
%:%176=76%:%
%:%179=77%:%
%:%183=77%:%
%:%184=77%:%
%:%185=78%:%
%:%186=78%:%
%:%191=78%:%
%:%194=79%:%
%:%195=80%:%
%:%196=80%:%
%:%197=81%:%
%:%198=82%:%
%:%201=83%:%
%:%205=83%:%
%:%206=83%:%
%:%207=84%:%
%:%208=84%:%
%:%213=84%:%
%:%216=85%:%
%:%217=86%:%
%:%218=86%:%
%:%219=87%:%
%:%220=88%:%
%:%223=89%:%
%:%227=89%:%
%:%228=89%:%
%:%229=89%:%
%:%234=89%:%
%:%237=90%:%
%:%238=91%:%
%:%239=91%:%
%:%240=92%:%
%:%241=93%:%
%:%244=94%:%
%:%248=94%:%
%:%249=94%:%
%:%250=95%:%
%:%251=95%:%
%:%252=96%:%
%:%253=96%:%
%:%254=97%:%
%:%255=97%:%
%:%256=98%:%
%:%257=98%:%
%:%258=98%:%
%:%259=99%:%
%:%260=99%:%
%:%261=99%:%
%:%262=100%:%
%:%263=100%:%
%:%264=100%:%
%:%265=101%:%
%:%266=101%:%
%:%267=102%:%
%:%268=102%:%
%:%269=102%:%
%:%270=103%:%
%:%276=103%:%
%:%281=104%:%
%:%286=105%:%

%
\begin{isabellebody}%
\setisabellecontext{PreEqu}%
%
\isadelimtheory
%
\endisadelimtheory
%
\isatagtheory
\isacommand{theory}\isamarkupfalse%
\ PreEqu\isanewline
\isanewline
\isakeyword{imports}\ Fresh\isanewline
\isanewline
\isakeyword{begin}%
\endisatagtheory
{\isafoldtheory}%
%
\isadelimtheory
\isanewline
%
\endisadelimtheory
\isanewline
\isanewline
\isacommand{inductive}\isamarkupfalse%
\ equ\ {\isacharcolon}{\kern0pt}{\isacharcolon}{\kern0pt}\ {\isachardoublequoteopen}fresh{\isacharunderscore}{\kern0pt}envs\ {\isasymRightarrow}\ trm\ {\isasymRightarrow}\ trm\ {\isasymRightarrow}\ bool{\isachardoublequoteclose}\ {\isacharparenleft}{\kern0pt}{\isachardoublequoteopen}\ {\isacharunderscore}{\kern0pt}\ {\isasymturnstile}\ {\isacharunderscore}{\kern0pt}\ {\isasymapprox}\ {\isacharunderscore}{\kern0pt}{\isachardoublequoteclose}\ {\isacharbrackleft}{\kern0pt}{\isadigit{8}}{\isadigit{0}}{\isacharcomma}{\kern0pt}{\isadigit{8}}{\isadigit{0}}{\isacharcomma}{\kern0pt}{\isadigit{8}}{\isadigit{0}}{\isacharbrackright}{\kern0pt}\ {\isadigit{8}}{\isadigit{0}}{\isacharparenright}{\kern0pt}\ \isakeyword{where}\isanewline
equ{\isacharunderscore}{\kern0pt}abst{\isacharunderscore}{\kern0pt}ab{\isacharbrackleft}{\kern0pt}intro{\isacharbang}{\kern0pt}{\isacharbrackright}{\kern0pt}{\isacharcolon}{\kern0pt}\ {\isachardoublequoteopen}{\isasymlbrakk}a{\isasymnoteq}b{\isacharsemicolon}{\kern0pt}{\isacharparenleft}{\kern0pt}nabla\ {\isasymturnstile}\ a\ {\isasymsharp}\ t{\isadigit{2}}{\isacharparenright}{\kern0pt}{\isacharsemicolon}{\kern0pt}{\isacharparenleft}{\kern0pt}nabla\ {\isasymturnstile}\ t{\isadigit{1}}\ {\isasymapprox}\ {\isacharparenleft}{\kern0pt}swap\ {\isacharbrackleft}{\kern0pt}{\isacharparenleft}{\kern0pt}a{\isacharcomma}{\kern0pt}b{\isacharparenright}{\kern0pt}{\isacharbrackright}{\kern0pt}\ t{\isadigit{2}}{\isacharparenright}{\kern0pt}{\isacharparenright}{\kern0pt}{\isasymrbrakk}\ \isanewline
\ \ \ \ \ \ \ \ \ \ \ \ \ \ \ \ \ \ \ \ \ \ {\isasymLongrightarrow}\ {\isacharparenleft}{\kern0pt}nabla\ {\isasymturnstile}\ Abst\ a\ t{\isadigit{1}}\ {\isasymapprox}\ Abst\ b\ t{\isadigit{2}}{\isacharparenright}{\kern0pt}{\isachardoublequoteclose}\ {\isacharbar}{\kern0pt}\isanewline
equ{\isacharunderscore}{\kern0pt}abst{\isacharunderscore}{\kern0pt}aa{\isacharbrackleft}{\kern0pt}intro{\isacharbang}{\kern0pt}{\isacharbrackright}{\kern0pt}{\isacharcolon}{\kern0pt}\ {\isachardoublequoteopen}{\isacharparenleft}{\kern0pt}nabla\ {\isasymturnstile}\ t{\isadigit{1}}\ {\isasymapprox}\ t{\isadigit{2}}{\isacharparenright}{\kern0pt}\ {\isasymLongrightarrow}\ {\isacharparenleft}{\kern0pt}nabla\ {\isasymturnstile}\ Abst\ a\ t{\isadigit{1}}\ {\isasymapprox}\ Abst\ a\ t{\isadigit{2}}{\isacharparenright}{\kern0pt}{\isachardoublequoteclose}\ {\isacharbar}{\kern0pt}\isanewline
equ{\isacharunderscore}{\kern0pt}unit{\isacharbrackleft}{\kern0pt}intro{\isacharbang}{\kern0pt}{\isacharbrackright}{\kern0pt}{\isacharcolon}{\kern0pt}\ \ \ \ {\isachardoublequoteopen}{\isacharparenleft}{\kern0pt}nabla\ {\isasymturnstile}\ Unit\ {\isasymapprox}\ Unit{\isacharparenright}{\kern0pt}{\isachardoublequoteclose}\ {\isacharbar}{\kern0pt}\isanewline
equ{\isacharunderscore}{\kern0pt}atom{\isacharbrackleft}{\kern0pt}intro{\isacharbang}{\kern0pt}{\isacharbrackright}{\kern0pt}{\isacharcolon}{\kern0pt}\ \ \ \ {\isachardoublequoteopen}a{\isacharequal}{\kern0pt}b{\isasymLongrightarrow}nabla\ {\isasymturnstile}\ Atom\ a\ {\isasymapprox}\ Atom\ b{\isachardoublequoteclose}\ {\isacharbar}{\kern0pt}\isanewline
equ{\isacharunderscore}{\kern0pt}susp{\isacharbrackleft}{\kern0pt}intro{\isacharbang}{\kern0pt}{\isacharbrackright}{\kern0pt}{\isacharcolon}{\kern0pt}\ \ \ \ {\isachardoublequoteopen}{\isacharparenleft}{\kern0pt}{\isasymforall}\ c\ {\isasymin}\ ds\ pi{\isadigit{1}}\ pi{\isadigit{2}}{\isachardot}{\kern0pt}\ {\isacharparenleft}{\kern0pt}c{\isacharcomma}{\kern0pt}X{\isacharparenright}{\kern0pt}\ {\isasymin}\ nabla{\isacharparenright}{\kern0pt}\ {\isasymLongrightarrow}\ {\isacharparenleft}{\kern0pt}nabla\ {\isasymturnstile}\ Susp\ pi{\isadigit{1}}\ X\ {\isasymapprox}\ Susp\ pi{\isadigit{2}}\ X{\isacharparenright}{\kern0pt}{\isachardoublequoteclose}\ {\isacharbar}{\kern0pt}\isanewline
equ{\isacharunderscore}{\kern0pt}paar{\isacharbrackleft}{\kern0pt}intro{\isacharbang}{\kern0pt}{\isacharbrackright}{\kern0pt}{\isacharcolon}{\kern0pt}\ \ \ \ {\isachardoublequoteopen}{\isasymlbrakk}{\isacharparenleft}{\kern0pt}nabla\ {\isasymturnstile}\ t{\isadigit{1}}{\isasymapprox}t{\isadigit{2}}{\isacharparenright}{\kern0pt}{\isacharsemicolon}{\kern0pt}{\isacharparenleft}{\kern0pt}nabla\ {\isasymturnstile}\ s{\isadigit{1}}{\isasymapprox}s{\isadigit{2}}{\isacharparenright}{\kern0pt}{\isasymrbrakk}\ {\isasymLongrightarrow}\ {\isacharparenleft}{\kern0pt}nabla\ {\isasymturnstile}\ Paar\ t{\isadigit{1}}\ s{\isadigit{1}}\ {\isasymapprox}\ Paar\ t{\isadigit{2}}\ s{\isadigit{2}}{\isacharparenright}{\kern0pt}{\isachardoublequoteclose}\ {\isacharbar}{\kern0pt}\isanewline
equ{\isacharunderscore}{\kern0pt}func{\isacharbrackleft}{\kern0pt}intro{\isacharbang}{\kern0pt}{\isacharbrackright}{\kern0pt}{\isacharcolon}{\kern0pt}\ \ \ \ {\isachardoublequoteopen}{\isacharparenleft}{\kern0pt}nabla\ {\isasymturnstile}\ t{\isadigit{1}}\ {\isasymapprox}\ t{\isadigit{2}}{\isacharparenright}{\kern0pt}\ {\isasymLongrightarrow}\ {\isacharparenleft}{\kern0pt}nabla\ {\isasymturnstile}\ Func\ F\ t{\isadigit{1}}\ {\isasymapprox}\ Func\ F\ t{\isadigit{2}}{\isacharparenright}{\kern0pt}{\isachardoublequoteclose}\isanewline
\isanewline
\isacommand{inductive{\isacharunderscore}{\kern0pt}cases}\isamarkupfalse%
\ Equ{\isacharunderscore}{\kern0pt}elims{\isacharcolon}{\kern0pt}\isanewline
{\isachardoublequoteopen}nabla\ {\isasymturnstile}\ Atom\ a\ {\isasymapprox}\ Atom\ b{\isachardoublequoteclose}\isanewline
{\isachardoublequoteopen}nabla\ {\isasymturnstile}\ Unit\ {\isasymapprox}\ Unit{\isachardoublequoteclose}\isanewline
{\isachardoublequoteopen}nabla\ {\isasymturnstile}\ Susp\ pi{\isadigit{1}}\ X\ {\isasymapprox}\ Susp\ pi{\isadigit{2}}\ X{\isachardoublequoteclose}\isanewline
{\isachardoublequoteopen}nabla\ {\isasymturnstile}\ Paar\ s{\isadigit{1}}\ t{\isadigit{1}}\ {\isasymapprox}\ Paar\ s{\isadigit{2}}\ t{\isadigit{2}}{\isachardoublequoteclose}\isanewline
{\isachardoublequoteopen}nabla\ {\isasymturnstile}\ Func\ F\ t{\isadigit{1}}\ {\isasymapprox}\ Func\ F\ t{\isadigit{2}}{\isachardoublequoteclose}\isanewline
{\isachardoublequoteopen}nabla\ {\isasymturnstile}\ Abst\ a\ t{\isadigit{1}}\ {\isasymapprox}\ Abst\ a\ t{\isadigit{2}}{\isachardoublequoteclose}\isanewline
{\isachardoublequoteopen}nabla\ {\isasymturnstile}\ Abst\ a\ t{\isadigit{1}}\ {\isasymapprox}\ Abst\ b\ t{\isadigit{2}}{\isachardoublequoteclose}\isanewline
\isanewline
\isanewline
\isacommand{lemma}\isamarkupfalse%
\ equ{\isacharunderscore}{\kern0pt}depth{\isacharcolon}{\kern0pt}\ \isanewline
\ \ \isakeyword{assumes}\ {\isachardoublequoteopen}nabla\ {\isasymturnstile}\ t{\isadigit{1}}\ {\isasymapprox}\ t{\isadigit{2}}{\isachardoublequoteclose}\isanewline
\ \ \isakeyword{shows}\ {\isachardoublequoteopen}depth\ t{\isadigit{1}}\ {\isacharequal}{\kern0pt}\ depth\ t{\isadigit{2}}{\isachardoublequoteclose}\isanewline
%
\isadelimproof
\ \ %
\endisadelimproof
%
\isatagproof
\isacommand{using}\isamarkupfalse%
\ assms\ \isacommand{by}\isamarkupfalse%
\ {\isacharparenleft}{\kern0pt}rule\ equ{\isachardot}{\kern0pt}induct{\isacharcomma}{\kern0pt}\ simp{\isacharunderscore}{\kern0pt}all{\isacharparenright}{\kern0pt}%
\endisatagproof
{\isafoldproof}%
%
\isadelimproof
\isanewline
%
\endisadelimproof
\isanewline
\isanewline
\isacommand{lemma}\isamarkupfalse%
\ rev{\isacharunderscore}{\kern0pt}pi{\isacharunderscore}{\kern0pt}pi{\isacharunderscore}{\kern0pt}equ{\isacharcolon}{\kern0pt}\ {\isachardoublequoteopen}{\isacharparenleft}{\kern0pt}nabla\ {\isasymturnstile}\ swap\ {\isacharparenleft}{\kern0pt}rev\ pi{\isacharparenright}{\kern0pt}\ {\isacharparenleft}{\kern0pt}swap\ pi\ t{\isacharparenright}{\kern0pt}\ {\isasymapprox}\ t{\isacharparenright}{\kern0pt}{\isachardoublequoteclose}\isanewline
%
\isadelimproof
%
\endisadelimproof
%
\isatagproof
\isacommand{proof}\isamarkupfalse%
{\isacharparenleft}{\kern0pt}induct\ t{\isacharparenright}{\kern0pt}\isanewline
\ \ \isacommand{case}\isamarkupfalse%
\ {\isacharparenleft}{\kern0pt}Susp\ pi{\isacharprime}{\kern0pt}\ X{\isacharparenright}{\kern0pt}\isanewline
\ \ \isacommand{then}\isamarkupfalse%
\ \isacommand{show}\isamarkupfalse%
\ {\isacharquery}{\kern0pt}case\ \isanewline
\ \ \ \ \isacommand{using}\isamarkupfalse%
\ ds{\isacharunderscore}{\kern0pt}cancel{\isacharunderscore}{\kern0pt}pi{\isacharunderscore}{\kern0pt}left{\isacharbrackleft}{\kern0pt}of\ {\isacharunderscore}{\kern0pt}\ {\isacartoucheopen}rev\ pi\ {\isacharat}{\kern0pt}\ pi{\isacartoucheclose}\ {\isacharunderscore}{\kern0pt}\ {\isacartoucheopen}{\isacharbrackleft}{\kern0pt}{\isacharbrackright}{\kern0pt}{\isacartoucheclose}{\isacharbrackright}{\kern0pt}\isanewline
\ \ \ \ ds{\isacharunderscore}{\kern0pt}rev{\isacharunderscore}{\kern0pt}pi{\isacharunderscore}{\kern0pt}pi\ elem{\isacharunderscore}{\kern0pt}ds\ ds{\isacharunderscore}{\kern0pt}rev{\isacharunderscore}{\kern0pt}pi{\isacharunderscore}{\kern0pt}id\ \isacommand{by}\isamarkupfalse%
\ auto\isanewline
\isacommand{qed}\isamarkupfalse%
\ {\isacharparenleft}{\kern0pt}auto{\isacharparenright}{\kern0pt}%
\endisatagproof
{\isafoldproof}%
%
\isadelimproof
\isanewline
%
\endisadelimproof
\ \ \isanewline
\isanewline
\isacommand{lemma}\isamarkupfalse%
\ equ{\isacharunderscore}{\kern0pt}pi{\isacharunderscore}{\kern0pt}right{\isacharcolon}{\kern0pt}\ \isanewline
\ \ \isakeyword{assumes}\ {\isachardoublequoteopen}{\isasymforall}a\ {\isasymin}\ ds\ {\isacharbrackleft}{\kern0pt}{\isacharbrackright}{\kern0pt}\ pi{\isachardot}{\kern0pt}\ nabla\ {\isasymturnstile}\ a\ {\isasymsharp}\ t{\isachardoublequoteclose}\isanewline
\ \ \isakeyword{shows}\ {\isachardoublequoteopen}nabla\ {\isasymturnstile}\ t\ {\isasymapprox}\ swap\ pi\ t{\isachardoublequoteclose}\isanewline
%
\isadelimproof
\ \ %
\endisadelimproof
%
\isatagproof
\isacommand{using}\isamarkupfalse%
\ assms\isanewline
\isacommand{proof}\isamarkupfalse%
{\isacharparenleft}{\kern0pt}induct\ t\ arbitrary{\isacharcolon}{\kern0pt}\ pi{\isacharparenright}{\kern0pt}\isanewline
\ \ \ \isacommand{case}\isamarkupfalse%
\ {\isacharparenleft}{\kern0pt}Abst\ b\ t{\isacharprime}{\kern0pt}{\isacharparenright}{\kern0pt}\isanewline
\ \ \ \isacommand{have}\isamarkupfalse%
\ {\isachardoublequoteopen}swapas\ pi\ b\ {\isacharequal}{\kern0pt}\ b\ {\isasymor}\ swapas\ pi\ b\ {\isasymnoteq}\ b{\isachardoublequoteclose}\ \isacommand{by}\isamarkupfalse%
\ simp\isanewline
\ \ \isacommand{moreover}\isamarkupfalse%
\ \isanewline
\ \ \isacommand{{\isacharbraceleft}{\kern0pt}}\isamarkupfalse%
\ \isacommand{assume}\isamarkupfalse%
\ eq{\isacharcolon}{\kern0pt}\ {\isachardoublequoteopen}swapas\ pi\ b\ {\isacharequal}{\kern0pt}\ b{\isachardoublequoteclose}\isanewline
\ \ \ \ \isacommand{have}\isamarkupfalse%
\ {\isachardoublequoteopen}nabla\ {\isasymturnstile}\ Abst\ b\ t{\isacharprime}{\kern0pt}\ {\isasymapprox}\ Abst\ b\ {\isacharparenleft}{\kern0pt}swap\ pi\ t{\isacharprime}{\kern0pt}{\isacharparenright}{\kern0pt}{\isachardoublequoteclose}\isanewline
\ \ \ \ \ \ \isacommand{apply}\isamarkupfalse%
\ {\isacharparenleft}{\kern0pt}rule\ equ{\isacharunderscore}{\kern0pt}abst{\isacharunderscore}{\kern0pt}aa{\isacharparenright}{\kern0pt}\isanewline
\ \ \ \ \ \ \isacommand{apply}\isamarkupfalse%
\ {\isacharparenleft}{\kern0pt}rule\ Abst{\isachardot}{\kern0pt}hyps{\isacharparenright}{\kern0pt}\isanewline
\ \ \ \ \ \ \isacommand{apply}\isamarkupfalse%
\ {\isacharparenleft}{\kern0pt}rule\ ballI{\isacharparenright}{\kern0pt}\isanewline
\ \ \ \ \ \ \isacommand{subgoal}\isamarkupfalse%
\ \isakeyword{for}\ a\ \isanewline
\ \ \ \ \ \ \ \ \isacommand{apply}\isamarkupfalse%
\ {\isacharparenleft}{\kern0pt}rule\ Fresh{\isacharunderscore}{\kern0pt}elims{\isacharparenleft}{\kern0pt}{\isadigit{1}}{\isacharparenright}{\kern0pt}{\isacharbrackleft}{\kern0pt}of\ nabla\ a\ b\ t{\isacharprime}{\kern0pt}{\isacharbrackright}{\kern0pt}{\isacharparenright}{\kern0pt}\isanewline
\ \ \ \ \ \ \ \ \isacommand{using}\isamarkupfalse%
\ Abst{\isachardot}{\kern0pt}prems\ elem{\isacharunderscore}{\kern0pt}ds{\isacharbrackleft}{\kern0pt}of\ a\ pi{\isacharbrackright}{\kern0pt}\ eq\ \isacommand{by}\isamarkupfalse%
\ auto\isanewline
\ \ \ \ \ \ \isacommand{done}\isamarkupfalse%
\isanewline
\ \ \ \ \isacommand{then}\isamarkupfalse%
\ \isacommand{have}\isamarkupfalse%
\ {\isachardoublequoteopen}nabla\ {\isasymturnstile}\ Abst\ b\ t{\isacharprime}{\kern0pt}\ {\isasymapprox}\ swap\ pi\ {\isacharparenleft}{\kern0pt}Abst\ b\ t{\isacharprime}{\kern0pt}{\isacharparenright}{\kern0pt}{\isachardoublequoteclose}\ \isacommand{using}\isamarkupfalse%
\ eq\ \isacommand{by}\isamarkupfalse%
\ simp\isanewline
\ \ \isacommand{{\isacharbraceright}{\kern0pt}}\isamarkupfalse%
\isanewline
\ \ \isacommand{moreover}\isamarkupfalse%
\isanewline
\ \ \isacommand{{\isacharbraceleft}{\kern0pt}}\isamarkupfalse%
\ \isacommand{assume}\isamarkupfalse%
\ neq{\isacharcolon}{\kern0pt}\ {\isachardoublequoteopen}b\ {\isasymnoteq}\ swapas\ pi\ b{\isachardoublequoteclose}\ \ \isanewline
\ \ \ \ \isacommand{obtain}\isamarkupfalse%
\ c\ \isakeyword{where}\ c{\isacharunderscore}{\kern0pt}def{\isacharcolon}{\kern0pt}\ {\isachardoublequoteopen}c\ {\isacharequal}{\kern0pt}\ swapas\ pi\ b{\isachardoublequoteclose}\ \isacommand{by}\isamarkupfalse%
\ simp\isanewline
\ \ \ \ \isacommand{have}\isamarkupfalse%
\ {\isachardoublequoteopen}c\ {\isasymin}\ ds\ {\isacharbrackleft}{\kern0pt}{\isacharbrackright}{\kern0pt}\ pi{\isachardoublequoteclose}\isanewline
\ \ \ \ \ \ \isacommand{by}\isamarkupfalse%
\ {\isacharparenleft}{\kern0pt}metis\ append{\isachardot}{\kern0pt}right{\isacharunderscore}{\kern0pt}neutral\ c{\isacharunderscore}{\kern0pt}def\ ds{\isacharunderscore}{\kern0pt}cancel{\isacharunderscore}{\kern0pt}pi{\isacharunderscore}{\kern0pt}front\ ds{\isacharunderscore}{\kern0pt}cancel{\isacharunderscore}{\kern0pt}pi{\isacharunderscore}{\kern0pt}left\ ds{\isacharunderscore}{\kern0pt}elem\ neq{\isacharparenright}{\kern0pt}\isanewline
\ \ \ \ \isacommand{then}\isamarkupfalse%
\ \isacommand{have}\isamarkupfalse%
\ one{\isacharcolon}{\kern0pt}\ {\isachardoublequoteopen}nabla\ {\isasymturnstile}\ c\ {\isasymsharp}\ Abst\ b\ t{\isacharprime}{\kern0pt}{\isachardoublequoteclose}\ \isacommand{using}\isamarkupfalse%
\ assms\ Abst{\isachardot}{\kern0pt}prems\ \isacommand{by}\isamarkupfalse%
\ auto\isanewline
\ \ \ \ \isacommand{have}\isamarkupfalse%
\ two{\isacharcolon}{\kern0pt}\ {\isachardoublequoteopen}b\ {\isasymnoteq}\ c{\isachardoublequoteclose}\ \isacommand{using}\isamarkupfalse%
\ neq\ c{\isacharunderscore}{\kern0pt}def\ \isacommand{by}\isamarkupfalse%
\ blast\isanewline
\ \ \ \ \isacommand{have}\isamarkupfalse%
\ {\isachardoublequoteopen}nabla\ {\isasymturnstile}\ c\ {\isasymsharp}\ t{\isacharprime}{\kern0pt}{\isachardoublequoteclose}\ \isacommand{using}\isamarkupfalse%
\ one\ two\ \isacommand{by}\isamarkupfalse%
\ cases\ auto\isanewline
\ \ \ \ \isacommand{have}\isamarkupfalse%
\ {\isachardoublequoteopen}nabla\ {\isasymturnstile}\ Abst\ b\ t{\isacharprime}{\kern0pt}\ {\isasymapprox}\ Abst\ c\ {\isacharparenleft}{\kern0pt}swap\ pi\ t{\isacharprime}{\kern0pt}{\isacharparenright}{\kern0pt}{\isachardoublequoteclose}\ \isanewline
\ \ \ \ \isacommand{proof}\isamarkupfalse%
{\isacharparenleft}{\kern0pt}rule\ equ{\isacharunderscore}{\kern0pt}abst{\isacharunderscore}{\kern0pt}ab{\isacharparenright}{\kern0pt}\isanewline
\ \ \ \ \ \ \isacommand{show}\isamarkupfalse%
\ {\isachardoublequoteopen}b\ {\isasymnoteq}\ c{\isachardoublequoteclose}\ \isacommand{using}\isamarkupfalse%
\ neq\ c{\isacharunderscore}{\kern0pt}def\ \isacommand{by}\isamarkupfalse%
\ blast\isanewline
\ \ \ \ \ \ \isacommand{have}\isamarkupfalse%
\ b{\isacharunderscore}{\kern0pt}is{\isacharunderscore}{\kern0pt}swap{\isacharcolon}{\kern0pt}\ {\isachardoublequoteopen}b\ {\isacharequal}{\kern0pt}\ swapas\ {\isacharparenleft}{\kern0pt}rev\ pi{\isacharparenright}{\kern0pt}\ c{\isachardoublequoteclose}\ \isacommand{using}\isamarkupfalse%
\ c{\isacharunderscore}{\kern0pt}def\ \isacommand{by}\isamarkupfalse%
\ simp\isanewline
\ \ \ \ \ \ \isacommand{show}\isamarkupfalse%
\ {\isachardoublequoteopen}nabla\ {\isasymturnstile}\ b\ {\isasymsharp}\ swap\ pi\ t{\isacharprime}{\kern0pt}{\isachardoublequoteclose}\ \isanewline
\ \ \ \ \ \ \ \ \isacommand{using}\isamarkupfalse%
\ Abst{\isachardot}{\kern0pt}prems\ Fresh{\isacharunderscore}{\kern0pt}elims{\isacharparenleft}{\kern0pt}{\isadigit{1}}{\isacharparenright}{\kern0pt}\ ds{\isacharunderscore}{\kern0pt}elem\ fresh{\isacharunderscore}{\kern0pt}swap{\isacharunderscore}{\kern0pt}right\ neq\ swapas{\isacharunderscore}{\kern0pt}rev{\isacharunderscore}{\kern0pt}pi{\isacharunderscore}{\kern0pt}a\ \isacommand{by}\isamarkupfalse%
\ metis\isanewline
\ \ \ \ \ \ \isacommand{show}\isamarkupfalse%
\ {\isachardoublequoteopen}nabla\ {\isasymturnstile}\ t{\isacharprime}{\kern0pt}\ {\isasymapprox}\ swap\ {\isacharbrackleft}{\kern0pt}{\isacharparenleft}{\kern0pt}b{\isacharcomma}{\kern0pt}\ c{\isacharparenright}{\kern0pt}{\isacharbrackright}{\kern0pt}\ {\isacharparenleft}{\kern0pt}swap\ pi\ t{\isacharprime}{\kern0pt}{\isacharparenright}{\kern0pt}{\isachardoublequoteclose}\isanewline
\ \ \ \ \ \ \isacommand{proof}\isamarkupfalse%
\ {\isacharminus}{\kern0pt}\isanewline
\ \ \ \ \ \ \ \ \ \isacommand{have}\isamarkupfalse%
\ fresh{\isacharunderscore}{\kern0pt}ext{\isacharcolon}{\kern0pt}\isanewline
\ \ \ \ \ \ \ \ \ \ {\isachardoublequoteopen}{\isasymforall}a\ {\isasymin}\ ds\ {\isacharbrackleft}{\kern0pt}{\isacharbrackright}{\kern0pt}\ {\isacharparenleft}{\kern0pt}{\isacharparenleft}{\kern0pt}b{\isacharcomma}{\kern0pt}c{\isacharparenright}{\kern0pt}\ {\isacharhash}{\kern0pt}\ pi{\isacharparenright}{\kern0pt}{\isachardot}{\kern0pt}\ nabla\ {\isasymturnstile}\ a\ {\isasymsharp}\ t{\isacharprime}{\kern0pt}{\isachardoublequoteclose}\isanewline
\ \ \ \ \ \ \ \ \isacommand{proof}\isamarkupfalse%
\ {\isacharparenleft}{\kern0pt}rule\ ballI{\isacharparenright}{\kern0pt}\isanewline
\ \ \ \ \ \ \ \ \ \ \isacommand{fix}\isamarkupfalse%
\ a\ \isacommand{assume}\isamarkupfalse%
\ A{\isacharcolon}{\kern0pt}\ {\isachardoublequoteopen}a\ {\isasymin}\ ds\ {\isacharbrackleft}{\kern0pt}{\isacharbrackright}{\kern0pt}\ {\isacharparenleft}{\kern0pt}{\isacharparenleft}{\kern0pt}b{\isacharcomma}{\kern0pt}c{\isacharparenright}{\kern0pt}{\isacharhash}{\kern0pt}pi{\isacharparenright}{\kern0pt}{\isachardoublequoteclose}\isanewline
\ \ \ \ \ \ \ \ \ \ \isacommand{then}\isamarkupfalse%
\ \isacommand{have}\isamarkupfalse%
\ {\isachardoublequoteopen}a\ {\isacharequal}{\kern0pt}\ c\ {\isasymor}\ a\ {\isasymin}\ ds\ {\isacharbrackleft}{\kern0pt}{\isacharbrackright}{\kern0pt}\ pi{\isachardoublequoteclose}\isanewline
\ \ \ \ \ \ \ \ \ \ \ \ \isacommand{using}\isamarkupfalse%
\ c{\isacharunderscore}{\kern0pt}def\ ds{\isacharunderscore}{\kern0pt}{\isadigit{7}}\ neq\ \isacommand{by}\isamarkupfalse%
\ auto\ \ \isanewline
\ \ \ \ \ \ \ \ \ \ \isacommand{then}\isamarkupfalse%
\ \isacommand{show}\isamarkupfalse%
\ {\isachardoublequoteopen}nabla\ {\isasymturnstile}\ a\ {\isasymsharp}\ t{\isacharprime}{\kern0pt}{\isachardoublequoteclose}\isanewline
\ \ \ \ \ \ \ \ \ \ \isacommand{proof}\isamarkupfalse%
\isanewline
\ \ \ \ \ \ \ \ \ \ \ \ \isacommand{assume}\isamarkupfalse%
\ {\isachardoublequoteopen}a\ {\isacharequal}{\kern0pt}\ c{\isachardoublequoteclose}\isanewline
\ \ \ \ \ \ \ \ \ \ \ \ \isacommand{with}\isamarkupfalse%
\ {\isacartoucheopen}nabla\ {\isasymturnstile}\ c\ {\isasymsharp}\ t{\isacharprime}{\kern0pt}{\isacartoucheclose}\ \isacommand{show}\isamarkupfalse%
\ {\isacharquery}{\kern0pt}thesis\ \isacommand{by}\isamarkupfalse%
\ simp\isanewline
\ \ \ \ \ \ \ \ \ \ \isacommand{next}\isamarkupfalse%
\isanewline
\ \ \ \ \ \ \ \ \ \ \ \ \isacommand{assume}\isamarkupfalse%
\ {\isachardoublequoteopen}a\ {\isasymin}\ ds\ {\isacharbrackleft}{\kern0pt}{\isacharbrackright}{\kern0pt}\ pi{\isachardoublequoteclose}\isanewline
\ \ \ \ \ \ \ \ \ \ \ \ \isacommand{show}\isamarkupfalse%
\ {\isacharquery}{\kern0pt}thesis\ \isanewline
\ \ \ \ \ \ \ \ \ \ \ \ \isacommand{proof}\isamarkupfalse%
{\isacharparenleft}{\kern0pt}rule\ Fresh{\isacharunderscore}{\kern0pt}elims{\isacharparenleft}{\kern0pt}{\isadigit{1}}{\isacharparenright}{\kern0pt}{\isacharbrackleft}{\kern0pt}of\ nabla\ a\ b\ t{\isacharprime}{\kern0pt}{\isacharbrackright}{\kern0pt}{\isacharcomma}{\kern0pt}\ goal{\isacharunderscore}{\kern0pt}cases{\isacharparenright}{\kern0pt}\isanewline
\ \ \ \ \ \ \ \ \ \ \ \ \ \ \isacommand{case}\isamarkupfalse%
\ {\isadigit{1}}\isanewline
\ \ \ \ \ \ \ \ \ \ \ \ \ \ \isacommand{then}\isamarkupfalse%
\ \isacommand{show}\isamarkupfalse%
\ {\isacharquery}{\kern0pt}case\isanewline
\ \ \ \ \ \ \ \ \ \ \ \ \ \ \ \ \isacommand{using}\isamarkupfalse%
\ Abst{\isachardot}{\kern0pt}prems\ {\isacartoucheopen}a\ {\isasymin}\ ds\ {\isacharbrackleft}{\kern0pt}{\isacharbrackright}{\kern0pt}\ pi{\isacartoucheclose}\isanewline
\ \ \ \ \ \ \ \ \ \ \ \ \ \ \isacommand{by}\isamarkupfalse%
\ simp\isanewline
\ \ \ \ \ \ \ \ \ \ \ \ \isacommand{next}\isamarkupfalse%
\isanewline
\ \ \ \ \ \ \ \ \ \ \ \ \ \ \isacommand{case}\isamarkupfalse%
\ {\isadigit{2}}\isanewline
\ \ \ \ \ \ \ \ \ \ \ \ \ \ \isacommand{then}\isamarkupfalse%
\ \isacommand{show}\isamarkupfalse%
\ {\isacharquery}{\kern0pt}case\ \isacommand{by}\isamarkupfalse%
\ simp\isanewline
\ \ \ \ \ \ \ \ \ \ \ \ \isacommand{next}\isamarkupfalse%
\isanewline
\ \ \ \ \ \ \ \ \ \ \ \ \ \ \isacommand{case}\isamarkupfalse%
\ {\isadigit{3}}\isanewline
\ \ \ \ \ \ \ \ \ \ \ \ \ \ \isacommand{have}\isamarkupfalse%
\ {\isachardoublequoteopen}a\ {\isasymnoteq}\ swapas\ {\isacharparenleft}{\kern0pt}{\isacharparenleft}{\kern0pt}b{\isacharcomma}{\kern0pt}\ c{\isacharparenright}{\kern0pt}\ {\isacharhash}{\kern0pt}\ pi{\isacharparenright}{\kern0pt}\ a{\isachardoublequoteclose}\ \isanewline
\ \ \ \ \ \ \ \ \ \ \ \ \ \ \ \ \isacommand{using}\isamarkupfalse%
\ A\ elem{\isacharunderscore}{\kern0pt}ds\isanewline
\ \ \ \ \ \ \ \ \ \ \ \ \ \ \ \ \isacommand{by}\isamarkupfalse%
\ blast\isanewline
\ \ \ \ \ \ \ \ \ \ \ \ \ \ \isacommand{hence}\isamarkupfalse%
\ {\isachardoublequoteopen}a\ {\isasymnoteq}\ b{\isachardoublequoteclose}\ \isacommand{using}\isamarkupfalse%
\ c{\isacharunderscore}{\kern0pt}def\ A\ \isanewline
\ \ \ \ \ \ \ \ \ \ \ \ \ \ \ \ \isacommand{by}\isamarkupfalse%
\ {\isacharparenleft}{\kern0pt}auto\ simp\ add{\isacharcolon}{\kern0pt}\ if{\isacharunderscore}{\kern0pt}split{\isacharparenright}{\kern0pt}\isanewline
\ \ \ \ \ \ \ \ \ \ \ \ \ \ \isacommand{hence}\isamarkupfalse%
\ False\ \isacommand{using}\isamarkupfalse%
\ {\isadigit{3}}\ \isacommand{by}\isamarkupfalse%
\ simp\isanewline
\ \ \ \ \ \ \ \ \ \ \ \ \ \ \isacommand{then}\isamarkupfalse%
\ \isacommand{show}\isamarkupfalse%
\ {\isacharquery}{\kern0pt}case\ \isacommand{by}\isamarkupfalse%
\ simp\isanewline
\ \ \ \ \ \ \ \ \ \ \isacommand{qed}\isamarkupfalse%
\isanewline
\ \ \ \ \ \ \ \ \isacommand{qed}\isamarkupfalse%
\isanewline
\ \ \ \ \ \ \isacommand{qed}\isamarkupfalse%
\isanewline
\ \ \ \ \ \ \ \ \isacommand{from}\isamarkupfalse%
\ Abst{\isachardot}{\kern0pt}hyps{\isacharbrackleft}{\kern0pt}OF\ fresh{\isacharunderscore}{\kern0pt}ext{\isacharbrackright}{\kern0pt}\isanewline
\ \ \ \ \ \ \ \ \isacommand{have}\isamarkupfalse%
\ {\isachardoublequoteopen}nabla\ {\isasymturnstile}\ t{\isacharprime}{\kern0pt}\ {\isasymapprox}\ swap\ {\isacharparenleft}{\kern0pt}{\isacharbrackleft}{\kern0pt}{\isacharparenleft}{\kern0pt}b{\isacharcomma}{\kern0pt}c{\isacharparenright}{\kern0pt}{\isacharbrackright}{\kern0pt}{\isacharat}{\kern0pt}pi{\isacharparenright}{\kern0pt}\ t{\isacharprime}{\kern0pt}{\isachardoublequoteclose}\isanewline
\ \ \ \ \ \ \ \ \ \ \isacommand{by}\isamarkupfalse%
\ simp\isanewline
\ \ \ \ \ \ \ \ \isacommand{moreover}\isamarkupfalse%
\ \isacommand{have}\isamarkupfalse%
\ {\isachardoublequoteopen}swap\ {\isacharparenleft}{\kern0pt}{\isacharbrackleft}{\kern0pt}{\isacharparenleft}{\kern0pt}b{\isacharcomma}{\kern0pt}c{\isacharparenright}{\kern0pt}{\isacharbrackright}{\kern0pt}{\isacharat}{\kern0pt}pi{\isacharparenright}{\kern0pt}\ t{\isacharprime}{\kern0pt}\ {\isacharequal}{\kern0pt}\ swap\ {\isacharbrackleft}{\kern0pt}{\isacharparenleft}{\kern0pt}b{\isacharcomma}{\kern0pt}c{\isacharparenright}{\kern0pt}{\isacharbrackright}{\kern0pt}\ {\isacharparenleft}{\kern0pt}swap\ pi\ t{\isacharprime}{\kern0pt}{\isacharparenright}{\kern0pt}{\isachardoublequoteclose}\isanewline
\ \ \ \ \ \ \ \ \ \ \isacommand{using}\isamarkupfalse%
\ swap{\isacharunderscore}{\kern0pt}append\ \isacommand{by}\isamarkupfalse%
\ blast\isanewline
\ \ \ \ \ \ \ \ \isacommand{ultimately}\isamarkupfalse%
\ \isacommand{show}\isamarkupfalse%
\ {\isacharquery}{\kern0pt}thesis\ \isacommand{by}\isamarkupfalse%
\ simp\isanewline
\ \ \ \ \ \ \isacommand{qed}\isamarkupfalse%
\isanewline
\ \ \ \ \isacommand{qed}\isamarkupfalse%
\isanewline
\ \ \ \isacommand{then}\isamarkupfalse%
\ \isacommand{have}\isamarkupfalse%
\ {\isachardoublequoteopen}nabla\ {\isasymturnstile}\ Abst\ b\ t{\isacharprime}{\kern0pt}\ {\isasymapprox}\ swap\ pi\ {\isacharparenleft}{\kern0pt}Abst\ b\ t{\isacharprime}{\kern0pt}{\isacharparenright}{\kern0pt}{\isachardoublequoteclose}\ \isacommand{using}\isamarkupfalse%
\ neq\ c{\isacharunderscore}{\kern0pt}def\ \isacommand{by}\isamarkupfalse%
\ force\isanewline
\ \isacommand{{\isacharbraceright}{\kern0pt}}\isamarkupfalse%
\isanewline
\isanewline
\ \ \isacommand{ultimately}\isamarkupfalse%
\ \isacommand{show}\isamarkupfalse%
\ {\isacharquery}{\kern0pt}case\ \isacommand{by}\isamarkupfalse%
\ argo\isanewline
\isanewline
\isacommand{next}\isamarkupfalse%
\isanewline
\ \ \isacommand{case}\isamarkupfalse%
\ {\isacharparenleft}{\kern0pt}Susp\ pi{\isacharprime}{\kern0pt}\ X{\isacharparenright}{\kern0pt}\isanewline
\ \ \isacommand{have}\isamarkupfalse%
\ {\isachardoublequoteopen}{\isasymforall}a{\isasymin}ds\ {\isacharbrackleft}{\kern0pt}{\isacharbrackright}{\kern0pt}\ pi{\isachardot}{\kern0pt}\ nabla\ {\isasymturnstile}\ a\ {\isasymsharp}\ Susp\ pi{\isacharprime}{\kern0pt}\ X{\isachardoublequoteclose}\ \isacommand{by}\isamarkupfalse%
\ fact\isanewline
\ \ \isacommand{then}\isamarkupfalse%
\ \isacommand{have}\isamarkupfalse%
\ {\isachardoublequoteopen}nabla\ {\isasymturnstile}\ Susp\ pi{\isacharprime}{\kern0pt}\ X\ {\isasymapprox}\ Susp\ {\isacharparenleft}{\kern0pt}pi\ {\isacharat}{\kern0pt}\ pi{\isacharprime}{\kern0pt}{\isacharparenright}{\kern0pt}\ X{\isachardoublequoteclose}\isanewline
\ \ \ \isacommand{using}\isamarkupfalse%
\ Fresh{\isacharunderscore}{\kern0pt}elims{\isacharparenleft}{\kern0pt}{\isadigit{4}}{\isacharparenright}{\kern0pt}\ equ{\isacharunderscore}{\kern0pt}susp\ ds{\isacharunderscore}{\kern0pt}def\ ds{\isacharunderscore}{\kern0pt}cancel{\isacharunderscore}{\kern0pt}pi{\isacharunderscore}{\kern0pt}left\isanewline
\ \ \ \isacommand{by}\isamarkupfalse%
\ {\isacharparenleft}{\kern0pt}metis\ {\isacharparenleft}{\kern0pt}lifting{\isacharparenright}{\kern0pt}\ append{\isacharunderscore}{\kern0pt}self{\isacharunderscore}{\kern0pt}conv{\isadigit{2}}\ swapas{\isacharunderscore}{\kern0pt}rev{\isacharunderscore}{\kern0pt}pi{\isacharunderscore}{\kern0pt}a{\isacharparenright}{\kern0pt}\isanewline
\ \ \isacommand{then}\isamarkupfalse%
\ \isacommand{show}\isamarkupfalse%
\ {\isacharquery}{\kern0pt}case\ \isacommand{by}\isamarkupfalse%
\ simp\isanewline
\isacommand{next}\isamarkupfalse%
\isanewline
\ \ \isacommand{case}\isamarkupfalse%
\ Unit\isanewline
\ \ \isacommand{then}\isamarkupfalse%
\ \isacommand{show}\isamarkupfalse%
\ {\isacharquery}{\kern0pt}case\ \isanewline
\ \ \ \ \isacommand{using}\isamarkupfalse%
\ equ{\isacharunderscore}{\kern0pt}unit\ swap{\isachardot}{\kern0pt}simps{\isacharparenleft}{\kern0pt}{\isadigit{2}}{\isacharparenright}{\kern0pt}\ \isacommand{by}\isamarkupfalse%
\ force\isanewline
\isacommand{next}\isamarkupfalse%
\isanewline
\ \ \isacommand{case}\isamarkupfalse%
\ {\isacharparenleft}{\kern0pt}Atom\ b{\isacharparenright}{\kern0pt}\isanewline
\ \ \ \ \isacommand{then}\isamarkupfalse%
\ \isacommand{show}\isamarkupfalse%
\ {\isacharquery}{\kern0pt}case\isanewline
\ \ \ \ \ \ \isacommand{apply}\isamarkupfalse%
\ simp\isanewline
\ \ \ \ \isacommand{apply}\isamarkupfalse%
\ {\isacharparenleft}{\kern0pt}rule\ equ{\isacharunderscore}{\kern0pt}atom{\isacharparenright}{\kern0pt}\isanewline
\ \ \ \ \isacommand{using}\isamarkupfalse%
\ Fresh{\isacharunderscore}{\kern0pt}elims{\isacharparenleft}{\kern0pt}{\isadigit{3}}{\isacharparenright}{\kern0pt}\ ds{\isacharunderscore}{\kern0pt}elem{\isacharunderscore}{\kern0pt}cp\isanewline
\ \ \ \ \isacommand{by}\isamarkupfalse%
\ metis\isanewline
\isacommand{next}\isamarkupfalse%
\isanewline
\ \ \isacommand{case}\isamarkupfalse%
\ {\isacharparenleft}{\kern0pt}Paar\ t{\isadigit{1}}\ t{\isadigit{2}}{\isacharparenright}{\kern0pt}\isanewline
\ \ \isacommand{then}\isamarkupfalse%
\ \isacommand{have}\isamarkupfalse%
\ hypsp{\isadigit{1}}{\isacharcolon}{\kern0pt}\ {\isachardoublequoteopen}{\isacharparenleft}{\kern0pt}{\isasymforall}a{\isasymin}ds\ {\isacharbrackleft}{\kern0pt}{\isacharbrackright}{\kern0pt}\ pi{\isachardot}{\kern0pt}\ \ nabla\ {\isasymturnstile}\ a\ {\isasymsharp}\ t{\isadigit{1}}{\isacharparenright}{\kern0pt}{\isachardoublequoteclose}\isanewline
\ \ \ \ \isakeyword{and}\ hypsp{\isadigit{2}}{\isacharcolon}{\kern0pt}\ {\isachardoublequoteopen}{\isacharparenleft}{\kern0pt}{\isasymforall}a{\isasymin}ds\ {\isacharbrackleft}{\kern0pt}{\isacharbrackright}{\kern0pt}\ pi{\isachardot}{\kern0pt}\ \ nabla\ {\isasymturnstile}\ a\ {\isasymsharp}\ t{\isadigit{2}}{\isacharparenright}{\kern0pt}{\isachardoublequoteclose}\isanewline
\ \ \ \ \isacommand{using}\isamarkupfalse%
\ Paar{\isachardot}{\kern0pt}prems\ Fresh{\isacharunderscore}{\kern0pt}elims{\isacharparenleft}{\kern0pt}{\isadigit{5}}{\isacharparenright}{\kern0pt}\ \isanewline
\ \ \ \ \ \isacommand{apply}\isamarkupfalse%
\ meson\isanewline
\ \ \ \ \isacommand{using}\isamarkupfalse%
\ Paar{\isachardot}{\kern0pt}prems\ Fresh{\isacharunderscore}{\kern0pt}elims{\isacharparenleft}{\kern0pt}{\isadigit{5}}{\isacharparenright}{\kern0pt}\ \isacommand{by}\isamarkupfalse%
\ blast\isanewline
\ \ \isacommand{have}\isamarkupfalse%
\ {\isachardoublequoteopen}nabla\ {\isasymturnstile}\ Paar\ t{\isadigit{1}}\ t{\isadigit{2}}\ {\isasymapprox}\ Paar\ {\isacharparenleft}{\kern0pt}swap\ pi\ t{\isadigit{1}}{\isacharparenright}{\kern0pt}\ {\isacharparenleft}{\kern0pt}swap\ pi\ t{\isadigit{2}}{\isacharparenright}{\kern0pt}{\isachardoublequoteclose}\isanewline
\ \ \ \ \isacommand{apply}\isamarkupfalse%
{\isacharparenleft}{\kern0pt}rule\ equ{\isacharunderscore}{\kern0pt}paar{\isacharparenright}{\kern0pt}\isanewline
\ \ \ \ \ \isacommand{apply}\isamarkupfalse%
{\isacharparenleft}{\kern0pt}rule\ Paar{\isachardot}{\kern0pt}hyps{\isacharparenright}{\kern0pt}\isanewline
\ \ \ \ \ \isacommand{apply}\isamarkupfalse%
\ {\isacharparenleft}{\kern0pt}simp\ add{\isacharcolon}{\kern0pt}\ hypsp{\isadigit{1}}{\isacharparenright}{\kern0pt}\isanewline
\ \ \ \ \isacommand{by}\isamarkupfalse%
\ {\isacharparenleft}{\kern0pt}simp\ add{\isacharcolon}{\kern0pt}\ Paar{\isachardot}{\kern0pt}hyps{\isacharparenleft}{\kern0pt}{\isadigit{2}}{\isacharparenright}{\kern0pt}\ hypsp{\isadigit{2}}{\isacharparenright}{\kern0pt}\isanewline
\ \ \isacommand{then}\isamarkupfalse%
\ \isacommand{show}\isamarkupfalse%
\ {\isacharquery}{\kern0pt}case\ \isacommand{by}\isamarkupfalse%
\ simp\isanewline
\isacommand{next}\isamarkupfalse%
\isanewline
\ \ \isacommand{case}\isamarkupfalse%
\ {\isacharparenleft}{\kern0pt}Func\ f\ t{\isacharprime}{\kern0pt}{\isacharparenright}{\kern0pt}\isanewline
\ \ \isacommand{then}\isamarkupfalse%
\ \isacommand{have}\isamarkupfalse%
\ hypsf{\isacharcolon}{\kern0pt}\ {\isachardoublequoteopen}{\isasymforall}a{\isasymin}ds\ {\isacharbrackleft}{\kern0pt}{\isacharbrackright}{\kern0pt}\ pi{\isachardot}{\kern0pt}\ nabla\ {\isasymturnstile}\ a\ {\isasymsharp}\ t{\isacharprime}{\kern0pt}{\isachardoublequoteclose}\ \isacommand{using}\isamarkupfalse%
\ fresh{\isacharunderscore}{\kern0pt}func\isanewline
\ \ \ \ \isacommand{by}\isamarkupfalse%
\ {\isacharparenleft}{\kern0pt}metis\ Fresh{\isacharunderscore}{\kern0pt}elims{\isacharparenleft}{\kern0pt}{\isadigit{6}}{\isacharparenright}{\kern0pt}{\isacharparenright}{\kern0pt}\isanewline
\ \ \isacommand{have}\isamarkupfalse%
\ {\isachardoublequoteopen}nabla\ {\isasymturnstile}\ Func\ f\ t{\isacharprime}{\kern0pt}\ {\isasymapprox}\ Func\ f\ {\isacharparenleft}{\kern0pt}swap\ pi\ t{\isacharprime}{\kern0pt}{\isacharparenright}{\kern0pt}{\isachardoublequoteclose}\isanewline
\ \ \ \ \isacommand{apply}\isamarkupfalse%
{\isacharparenleft}{\kern0pt}rule\ equ{\isacharunderscore}{\kern0pt}func{\isacharparenright}{\kern0pt}\isanewline
\ \ \ \ \isacommand{apply}\isamarkupfalse%
{\isacharparenleft}{\kern0pt}rule\ Func{\isachardot}{\kern0pt}hyps{\isacharparenright}{\kern0pt}\isanewline
\ \ \ \ \isacommand{using}\isamarkupfalse%
\ hypsf\ \isacommand{by}\isamarkupfalse%
\ auto\isanewline
\ \ \isacommand{then}\isamarkupfalse%
\ \isacommand{show}\isamarkupfalse%
\ {\isacharquery}{\kern0pt}case\isanewline
\ \ \ \ \isacommand{by}\isamarkupfalse%
\ simp\isanewline
\isacommand{qed}\isamarkupfalse%
%
\endisatagproof
{\isafoldproof}%
%
\isadelimproof
\isanewline
%
\endisadelimproof
\isanewline
\isacommand{lemma}\isamarkupfalse%
\ pi{\isacharunderscore}{\kern0pt}comm{\isacharcolon}{\kern0pt}\ {\isachardoublequoteopen}nabla\ {\isasymturnstile}\ {\isacharparenleft}{\kern0pt}swap\ {\isacharparenleft}{\kern0pt}pi\ {\isacharat}{\kern0pt}\ {\isacharbrackleft}{\kern0pt}{\isacharparenleft}{\kern0pt}a{\isacharcomma}{\kern0pt}b{\isacharparenright}{\kern0pt}{\isacharbrackright}{\kern0pt}{\isacharparenright}{\kern0pt}\ t{\isacharparenright}{\kern0pt}\ {\isasymapprox}\ {\isacharparenleft}{\kern0pt}swap\ {\isacharparenleft}{\kern0pt}{\isacharbrackleft}{\kern0pt}{\isacharparenleft}{\kern0pt}swapas\ pi\ a{\isacharcomma}{\kern0pt}\ swapas\ pi\ b{\isacharparenright}{\kern0pt}{\isacharbrackright}{\kern0pt}\ {\isacharat}{\kern0pt}\ pi{\isacharparenright}{\kern0pt}\ t{\isacharparenright}{\kern0pt}{\isachardoublequoteclose}\isanewline
%
\isadelimproof
%
\endisadelimproof
%
\isatagproof
\isacommand{proof}\isamarkupfalse%
{\isacharparenleft}{\kern0pt}induct\ t{\isacharparenright}{\kern0pt}\isanewline
\ \ \isacommand{case}\isamarkupfalse%
\ {\isacharparenleft}{\kern0pt}Abst\ c\ t{\isacharparenright}{\kern0pt}\isanewline
\ \ \isacommand{then}\isamarkupfalse%
\ \isacommand{show}\isamarkupfalse%
\ {\isacharquery}{\kern0pt}case\ \isacommand{using}\isamarkupfalse%
\ swapas{\isacharunderscore}{\kern0pt}comm\ \isacommand{by}\isamarkupfalse%
\ force\isanewline
\isacommand{next}\isamarkupfalse%
\isanewline
\ \ \isacommand{case}\isamarkupfalse%
\ {\isacharparenleft}{\kern0pt}Susp\ {\isasympi}{\isacharprime}{\kern0pt}\ X{\isacharparenright}{\kern0pt}\isanewline
\ \ \isacommand{have}\isamarkupfalse%
\ rew{\isadigit{1}}{\isacharcolon}{\kern0pt}\isanewline
\ \ \ \ {\isachardoublequoteopen}swap\ {\isacharparenleft}{\kern0pt}pi\ {\isacharat}{\kern0pt}\ {\isacharbrackleft}{\kern0pt}{\isacharparenleft}{\kern0pt}a{\isacharcomma}{\kern0pt}b{\isacharparenright}{\kern0pt}{\isacharbrackright}{\kern0pt}{\isacharparenright}{\kern0pt}\ {\isacharparenleft}{\kern0pt}Susp\ {\isasympi}{\isacharprime}{\kern0pt}\ X{\isacharparenright}{\kern0pt}\ {\isacharequal}{\kern0pt}\ Susp\ {\isacharparenleft}{\kern0pt}pi\ {\isacharat}{\kern0pt}\ {\isacharbrackleft}{\kern0pt}{\isacharparenleft}{\kern0pt}a{\isacharcomma}{\kern0pt}b{\isacharparenright}{\kern0pt}{\isacharbrackright}{\kern0pt}\ {\isacharat}{\kern0pt}\ {\isasympi}{\isacharprime}{\kern0pt}{\isacharparenright}{\kern0pt}\ X{\isachardoublequoteclose}\isanewline
\ \ \ \ \isacommand{by}\isamarkupfalse%
\ simp\isanewline
\ \ \isacommand{have}\isamarkupfalse%
\ rew{\isadigit{2}}{\isacharcolon}{\kern0pt}\isanewline
\ \ \ \ {\isachardoublequoteopen}swap\ {\isacharparenleft}{\kern0pt}{\isacharbrackleft}{\kern0pt}{\isacharparenleft}{\kern0pt}swapas\ pi\ a{\isacharcomma}{\kern0pt}\ swapas\ pi\ b{\isacharparenright}{\kern0pt}{\isacharbrackright}{\kern0pt}\ {\isacharat}{\kern0pt}\ pi{\isacharparenright}{\kern0pt}\ {\isacharparenleft}{\kern0pt}Susp\ {\isasympi}{\isacharprime}{\kern0pt}\ X{\isacharparenright}{\kern0pt}\ \isanewline
\ \ \ \ {\isacharequal}{\kern0pt}\ Susp\ {\isacharparenleft}{\kern0pt}{\isacharbrackleft}{\kern0pt}{\isacharparenleft}{\kern0pt}swapas\ pi\ a{\isacharcomma}{\kern0pt}\ swapas\ pi\ b{\isacharparenright}{\kern0pt}{\isacharbrackright}{\kern0pt}\ {\isacharat}{\kern0pt}\ pi\ {\isacharat}{\kern0pt}\ {\isasympi}{\isacharprime}{\kern0pt}{\isacharparenright}{\kern0pt}\ X{\isachardoublequoteclose}\isanewline
\ \ \ \ \isacommand{by}\isamarkupfalse%
\ simp\isanewline
\ \ \isacommand{have}\isamarkupfalse%
\ fresh{\isacharcolon}{\kern0pt}\isanewline
\ \ {\isachardoublequoteopen}{\isasymforall}c\ {\isasymin}\ ds\ {\isacharparenleft}{\kern0pt}pi\ {\isacharat}{\kern0pt}\ {\isacharbrackleft}{\kern0pt}{\isacharparenleft}{\kern0pt}a{\isacharcomma}{\kern0pt}b{\isacharparenright}{\kern0pt}{\isacharbrackright}{\kern0pt}\ {\isacharat}{\kern0pt}\ {\isasympi}{\isacharprime}{\kern0pt}{\isacharparenright}{\kern0pt}\ {\isacharparenleft}{\kern0pt}{\isacharbrackleft}{\kern0pt}{\isacharparenleft}{\kern0pt}swapas\ pi\ a{\isacharcomma}{\kern0pt}\ swapas\ pi\ b{\isacharparenright}{\kern0pt}{\isacharbrackright}{\kern0pt}\ {\isacharat}{\kern0pt}\ pi\ {\isacharat}{\kern0pt}\ {\isasympi}{\isacharprime}{\kern0pt}{\isacharparenright}{\kern0pt}{\isachardot}{\kern0pt}\ {\isacharparenleft}{\kern0pt}c{\isacharcomma}{\kern0pt}\ X{\isacharparenright}{\kern0pt}\ {\isasymin}\ nabla{\isachardoublequoteclose}\isanewline
\ \ \ \ \isacommand{using}\isamarkupfalse%
\ swapas{\isacharunderscore}{\kern0pt}append\ swapas{\isacharunderscore}{\kern0pt}comm\ ds{\isacharunderscore}{\kern0pt}def\ \isacommand{by}\isamarkupfalse%
\ simp\isanewline
\ \ \isacommand{from}\isamarkupfalse%
\ rew{\isadigit{1}}\ rew{\isadigit{2}}\ fresh\isanewline
\ \ \isacommand{show}\isamarkupfalse%
\ {\isacharquery}{\kern0pt}case\isanewline
\ \ \ \ \isacommand{using}\isamarkupfalse%
\ equ{\isacharunderscore}{\kern0pt}susp\ \isacommand{by}\isamarkupfalse%
\ simp\isanewline
\isacommand{next}\isamarkupfalse%
\isanewline
\ \ \isacommand{case}\isamarkupfalse%
\ Unit\isanewline
\ \ \isacommand{then}\isamarkupfalse%
\ \isacommand{show}\isamarkupfalse%
\ {\isacharquery}{\kern0pt}case\ \isacommand{by}\isamarkupfalse%
\ force\isanewline
\isacommand{next}\isamarkupfalse%
\isanewline
\ \ \isacommand{case}\isamarkupfalse%
\ {\isacharparenleft}{\kern0pt}Atom\ c{\isacharparenright}{\kern0pt}\isanewline
\ \ \isacommand{then}\isamarkupfalse%
\ \isacommand{show}\isamarkupfalse%
\ {\isacharquery}{\kern0pt}case\ \isanewline
\ \ \ \ \isacommand{using}\isamarkupfalse%
\ equ{\isacharunderscore}{\kern0pt}atom\ swapas{\isacharunderscore}{\kern0pt}comm\ \isacommand{by}\isamarkupfalse%
\ auto\isanewline
\isacommand{next}\isamarkupfalse%
\isanewline
\ \ \isacommand{case}\isamarkupfalse%
\ {\isacharparenleft}{\kern0pt}Paar\ t{\isadigit{1}}\ t{\isadigit{2}}{\isacharparenright}{\kern0pt}\isanewline
\ \ \isacommand{then}\isamarkupfalse%
\ \isacommand{show}\isamarkupfalse%
\ {\isacharquery}{\kern0pt}case\ \isacommand{by}\isamarkupfalse%
\ force\isanewline
\isacommand{next}\isamarkupfalse%
\isanewline
\ \ \isacommand{case}\isamarkupfalse%
\ {\isacharparenleft}{\kern0pt}Func\ f\ t{\isacharparenright}{\kern0pt}\isanewline
\ \ \isacommand{then}\isamarkupfalse%
\ \isacommand{show}\isamarkupfalse%
\ {\isacharquery}{\kern0pt}case\ \isacommand{by}\isamarkupfalse%
\ force\isanewline
\isacommand{qed}\isamarkupfalse%
%
\endisatagproof
{\isafoldproof}%
%
\isadelimproof
\isanewline
%
\endisadelimproof
\isanewline
\isanewline
\isacommand{lemma}\isamarkupfalse%
\ l{\isadigit{3}}{\isacharunderscore}{\kern0pt}jud{\isacharcolon}{\kern0pt}\ \isanewline
\ \ \isakeyword{assumes}\ {\isachardoublequoteopen}{\isacharparenleft}{\kern0pt}nabla\ {\isasymturnstile}\ t{\isadigit{1}}\ {\isasymapprox}\ t{\isadigit{2}}{\isacharparenright}{\kern0pt}{\isachardoublequoteclose}\isanewline
\ \ \isakeyword{shows}\ {\isachardoublequoteopen}{\isacharparenleft}{\kern0pt}nabla\ {\isasymturnstile}\ a\ {\isasymsharp}\ t{\isadigit{1}}{\isacharparenright}{\kern0pt}\ {\isasymlongrightarrow}\ {\isacharparenleft}{\kern0pt}nabla\ {\isasymturnstile}\ a\ {\isasymsharp}\ t{\isadigit{2}}{\isacharparenright}{\kern0pt}{\isachardoublequoteclose}\isanewline
%
\isadelimproof
\ \ %
\endisadelimproof
%
\isatagproof
\isacommand{using}\isamarkupfalse%
\ assms\isanewline
\isacommand{proof}\isamarkupfalse%
{\isacharparenleft}{\kern0pt}induction\ rule{\isacharcolon}{\kern0pt}\ equ{\isachardot}{\kern0pt}induct{\isacharparenright}{\kern0pt}\isanewline
\ \ \isacommand{case}\isamarkupfalse%
\ {\isacharparenleft}{\kern0pt}equ{\isacharunderscore}{\kern0pt}abst{\isacharunderscore}{\kern0pt}ab\ b\ c\ nabla\ t{\isadigit{2}}\ t{\isadigit{1}}{\isacharparenright}{\kern0pt}\isanewline
\ \ \isacommand{then}\isamarkupfalse%
\ \isacommand{show}\isamarkupfalse%
\ {\isacharquery}{\kern0pt}case\ \isanewline
\ \ \ \ \isacommand{using}\isamarkupfalse%
\ fresh{\isacharunderscore}{\kern0pt}swap{\isacharunderscore}{\kern0pt}left\ Fresh{\isacharunderscore}{\kern0pt}elims{\isacharparenleft}{\kern0pt}{\isadigit{1}}{\isacharparenright}{\kern0pt}\ fresh{\isacharunderscore}{\kern0pt}abst{\isacharunderscore}{\kern0pt}aa\ fresh{\isacharunderscore}{\kern0pt}abst{\isacharunderscore}{\kern0pt}ab\ rev{\isacharunderscore}{\kern0pt}singleton{\isacharunderscore}{\kern0pt}conv\ swapa{\isachardot}{\kern0pt}simps\ swapas{\isachardot}{\kern0pt}simps{\isacharparenleft}{\kern0pt}{\isadigit{1}}{\isacharcomma}{\kern0pt}{\isadigit{2}}{\isacharparenright}{\kern0pt}\isanewline
\ \ \ \ \isacommand{by}\isamarkupfalse%
\ metis\isanewline
\isacommand{next}\isamarkupfalse%
\isanewline
\ \ \isacommand{case}\isamarkupfalse%
\ {\isacharparenleft}{\kern0pt}equ{\isacharunderscore}{\kern0pt}abst{\isacharunderscore}{\kern0pt}aa\ nabla\ t{\isadigit{1}}\ t{\isadigit{2}}\ b{\isacharparenright}{\kern0pt}\isanewline
\ \ \isacommand{then}\isamarkupfalse%
\ \isacommand{show}\isamarkupfalse%
\ {\isacharquery}{\kern0pt}case\ \isanewline
\ \ \ \ \isacommand{using}\isamarkupfalse%
\ fresh{\isacharunderscore}{\kern0pt}swap{\isacharunderscore}{\kern0pt}left\ Fresh{\isacharunderscore}{\kern0pt}elims{\isacharparenleft}{\kern0pt}{\isadigit{1}}{\isacharparenright}{\kern0pt}\ fresh{\isacharunderscore}{\kern0pt}abst{\isacharunderscore}{\kern0pt}aa\ fresh{\isacharunderscore}{\kern0pt}abst{\isacharunderscore}{\kern0pt}ab\ \isanewline
\ \ \ \ \isacommand{by}\isamarkupfalse%
\ metis\isanewline
\isacommand{next}\isamarkupfalse%
\isanewline
\ \ \isacommand{case}\isamarkupfalse%
\ {\isacharparenleft}{\kern0pt}equ{\isacharunderscore}{\kern0pt}unit\ nabla{\isacharparenright}{\kern0pt}\isanewline
\ \ \isacommand{then}\isamarkupfalse%
\ \isacommand{show}\isamarkupfalse%
\ {\isacharquery}{\kern0pt}case\isanewline
\ \ \ \ \isacommand{by}\isamarkupfalse%
\ blast\isanewline
\isacommand{next}\isamarkupfalse%
\isanewline
\ \ \isacommand{case}\isamarkupfalse%
\ {\isacharparenleft}{\kern0pt}equ{\isacharunderscore}{\kern0pt}atom\ b\ c\ nabla{\isacharparenright}{\kern0pt}\isanewline
\ \ \isacommand{then}\isamarkupfalse%
\ \isacommand{show}\isamarkupfalse%
\ {\isacharquery}{\kern0pt}case\ \isanewline
\ \ \ \ \isacommand{by}\isamarkupfalse%
\ simp\isanewline
\isacommand{next}\isamarkupfalse%
\isanewline
\ \ \isacommand{case}\isamarkupfalse%
\ {\isacharparenleft}{\kern0pt}equ{\isacharunderscore}{\kern0pt}susp\ pi{\isadigit{1}}\ pi{\isadigit{2}}\ X\ nabla{\isacharparenright}{\kern0pt}\isanewline
\ \ \isacommand{then}\isamarkupfalse%
\ \isacommand{show}\isamarkupfalse%
\ {\isacharquery}{\kern0pt}case\ \isanewline
\ \ \ \ \isacommand{using}\isamarkupfalse%
\ Fresh{\isacharunderscore}{\kern0pt}elims{\isacharparenleft}{\kern0pt}{\isadigit{4}}{\isacharparenright}{\kern0pt}\ ds{\isacharunderscore}{\kern0pt}elem{\isacharunderscore}{\kern0pt}cp\ ds{\isacharunderscore}{\kern0pt}rev\ fresh{\isacharunderscore}{\kern0pt}susp\ swapas{\isacharunderscore}{\kern0pt}append\ swapas{\isacharunderscore}{\kern0pt}rev{\isacharunderscore}{\kern0pt}pi{\isacharunderscore}{\kern0pt}a\isanewline
\ \ \ \ \isacommand{by}\isamarkupfalse%
\ metis\isanewline
\isacommand{next}\isamarkupfalse%
\isanewline
\ \ \isacommand{case}\isamarkupfalse%
\ {\isacharparenleft}{\kern0pt}equ{\isacharunderscore}{\kern0pt}paar\ nabla\ t{\isadigit{1}}\ t{\isadigit{2}}\ s{\isadigit{1}}\ s{\isadigit{2}}{\isacharparenright}{\kern0pt}\isanewline
\ \ \isacommand{then}\isamarkupfalse%
\ \isacommand{show}\isamarkupfalse%
\ {\isacharquery}{\kern0pt}case\isanewline
\ \ \ \ \isacommand{using}\isamarkupfalse%
\ Fresh{\isacharunderscore}{\kern0pt}elims{\isacharparenleft}{\kern0pt}{\isadigit{5}}{\isacharparenright}{\kern0pt}\ \isacommand{by}\isamarkupfalse%
\ blast\isanewline
\isacommand{next}\isamarkupfalse%
\isanewline
\ \ \isacommand{case}\isamarkupfalse%
\ {\isacharparenleft}{\kern0pt}equ{\isacharunderscore}{\kern0pt}func\ nabla\ t{\isadigit{1}}\ t{\isadigit{2}}\ f{\isacharparenright}{\kern0pt}\isanewline
\ \ \isacommand{then}\isamarkupfalse%
\ \isacommand{show}\isamarkupfalse%
\ {\isacharquery}{\kern0pt}case\ \isanewline
\ \ \ \ \isacommand{using}\isamarkupfalse%
\ Fresh{\isacharunderscore}{\kern0pt}elims{\isacharparenleft}{\kern0pt}{\isadigit{6}}{\isacharparenright}{\kern0pt}\ \isacommand{by}\isamarkupfalse%
\ blast\isanewline
\isacommand{qed}\isamarkupfalse%
%
\endisatagproof
{\isafoldproof}%
%
\isadelimproof
\isanewline
%
\endisadelimproof
\isanewline
\isanewline
\isanewline
\isanewline
\isacommand{lemma}\isamarkupfalse%
\ ds{\isacharunderscore}{\kern0pt}empty{\isacharunderscore}{\kern0pt}equiv{\isacharunderscore}{\kern0pt}{\isadigit{1}}{\isacharcolon}{\kern0pt}\isanewline
\ \ \isakeyword{assumes}\ {\isachardoublequoteopen}ds\ pi{\isadigit{1}}\ pi{\isadigit{2}}\ {\isacharequal}{\kern0pt}\ {\isacharbraceleft}{\kern0pt}{\isacharbraceright}{\kern0pt}{\isachardoublequoteclose}\isanewline
\ \ \isakeyword{shows}\ {\isachardoublequoteopen}nabla\ {\isasymturnstile}\ swap\ pi{\isadigit{1}}\ t{\isadigit{1}}\ {\isasymapprox}\ t{\isadigit{2}}\ {\isasymLongrightarrow}\ nabla\ {\isasymturnstile}\ swap\ pi{\isadigit{2}}\ t{\isadigit{1}}\ {\isasymapprox}\ t{\isadigit{2}}{\isachardoublequoteclose}\isanewline
%
\isadelimproof
\ \ %
\endisadelimproof
%
\isatagproof
\isacommand{using}\isamarkupfalse%
\ assms\isanewline
\isacommand{proof}\isamarkupfalse%
{\isacharparenleft}{\kern0pt}induction\ t{\isadigit{1}}\ arbitrary{\isacharcolon}{\kern0pt}\ t{\isadigit{2}}{\isacharparenright}{\kern0pt}\isanewline
\ \ \isacommand{case}\isamarkupfalse%
\ {\isacharparenleft}{\kern0pt}Abst\ a\ t{\isadigit{1}}{\isacharprime}{\kern0pt}{\isacharparenright}{\kern0pt}\isanewline
\ \ \isacommand{then}\isamarkupfalse%
\ \isacommand{obtain}\isamarkupfalse%
\ b\ t{\isadigit{2}}{\isacharprime}{\kern0pt}\isanewline
\ \ \ \ \isakeyword{where}\ t{\isadigit{2}}{\isacharcolon}{\kern0pt}\ {\isachardoublequoteopen}t{\isadigit{2}}\ {\isacharequal}{\kern0pt}\ Abst\ b\ t{\isadigit{2}}{\isacharprime}{\kern0pt}{\isachardoublequoteclose}\isanewline
\ \ \ \ \isacommand{by}\isamarkupfalse%
{\isacharparenleft}{\kern0pt}auto\ elim{\isacharcolon}{\kern0pt}\ equ{\isachardot}{\kern0pt}cases{\isacharparenright}{\kern0pt}\isanewline
\ \ \isacommand{with}\isamarkupfalse%
\ Abst\ \isacommand{have}\isamarkupfalse%
\ i{\isacharcolon}{\kern0pt}\ {\isachardoublequoteopen}nabla\ {\isasymturnstile}\ Abst\ {\isacharparenleft}{\kern0pt}swapas\ pi{\isadigit{1}}\ a{\isacharparenright}{\kern0pt}\ {\isacharparenleft}{\kern0pt}swap\ pi{\isadigit{1}}\ t{\isadigit{1}}{\isacharprime}{\kern0pt}{\isacharparenright}{\kern0pt}\ {\isasymapprox}\ Abst\ b\ t{\isadigit{2}}{\isacharprime}{\kern0pt}{\isachardoublequoteclose}\isanewline
\ \ \ \ \isacommand{using}\isamarkupfalse%
\ swap{\isachardot}{\kern0pt}simps{\isacharparenleft}{\kern0pt}{\isadigit{4}}{\isacharparenright}{\kern0pt}\ \isacommand{by}\isamarkupfalse%
\ simp\isanewline
\ \ \isacommand{then}\isamarkupfalse%
\ \isacommand{show}\isamarkupfalse%
\ {\isacharquery}{\kern0pt}case\isanewline
\ \ \isacommand{proof}\isamarkupfalse%
{\isacharparenleft}{\kern0pt}cases\ {\isacartoucheopen}swapas\ pi{\isadigit{1}}\ a\ {\isacharequal}{\kern0pt}\ b{\isacartoucheclose}{\isacharparenright}{\kern0pt}\isanewline
\ \ \ \ \isacommand{case}\isamarkupfalse%
\ True\isanewline
\ \ \ \ \isacommand{then}\isamarkupfalse%
\ \isacommand{have}\isamarkupfalse%
\ {\isachardoublequoteopen}nabla\ {\isasymturnstile}\ swap\ pi{\isadigit{1}}\ t{\isadigit{1}}{\isacharprime}{\kern0pt}\ {\isasymapprox}\ t{\isadigit{2}}{\isacharprime}{\kern0pt}{\isachardoublequoteclose}\ \isanewline
\ \ \ \ \ \ \isacommand{using}\isamarkupfalse%
\ Abst\ i\ Equ{\isacharunderscore}{\kern0pt}elims{\isacharparenleft}{\kern0pt}{\isadigit{6}}{\isacharparenright}{\kern0pt}\ \isacommand{by}\isamarkupfalse%
\ blast\isanewline
\ \ \ \ \isacommand{then}\isamarkupfalse%
\ \isacommand{have}\isamarkupfalse%
\ ii{\isacharcolon}{\kern0pt}\ {\isachardoublequoteopen}nabla\ {\isasymturnstile}\ swap\ pi{\isadigit{2}}\ t{\isadigit{1}}{\isacharprime}{\kern0pt}\ {\isasymapprox}\ t{\isadigit{2}}{\isacharprime}{\kern0pt}{\isachardoublequoteclose}\isanewline
\ \ \ \ \ \ \isacommand{using}\isamarkupfalse%
\ Abst\ \isacommand{by}\isamarkupfalse%
\ simp\isanewline
\ \ \ \ \isacommand{then}\isamarkupfalse%
\ \isacommand{have}\isamarkupfalse%
\ {\isachardoublequoteopen}nabla\ {\isasymturnstile}\ Abst\ {\isacharparenleft}{\kern0pt}swapas\ pi{\isadigit{2}}\ a{\isacharparenright}{\kern0pt}\ {\isacharparenleft}{\kern0pt}swap\ pi{\isadigit{2}}\ t{\isadigit{1}}{\isacharprime}{\kern0pt}{\isacharparenright}{\kern0pt}\ {\isasymapprox}\ Abst\ b\ t{\isadigit{2}}{\isacharprime}{\kern0pt}{\isachardoublequoteclose}\isanewline
\ \ \ \ \ \ \isacommand{using}\isamarkupfalse%
\ True\ ds{\isacharunderscore}{\kern0pt}swapas{\isacharunderscore}{\kern0pt}eq{\isacharbrackleft}{\kern0pt}OF\ Abst{\isacharparenleft}{\kern0pt}{\isadigit{3}}{\isacharparenright}{\kern0pt}{\isacharcomma}{\kern0pt}\ of\ a{\isacharbrackright}{\kern0pt}\ equ{\isacharunderscore}{\kern0pt}abst{\isacharunderscore}{\kern0pt}aa{\isacharbrackleft}{\kern0pt}OF\ ii{\isacharcomma}{\kern0pt}\ of\ {\isacartoucheopen}swapas\ pi{\isadigit{2}}\ a{\isacartoucheclose}{\isacharbrackright}{\kern0pt}\ \isanewline
\ \ \ \ \ \ \isacommand{by}\isamarkupfalse%
\ simp\isanewline
\ \ \ \ \isacommand{then}\isamarkupfalse%
\ \isacommand{show}\isamarkupfalse%
\ {\isacharquery}{\kern0pt}thesis\ \isanewline
\ \ \ \ \ \ \isacommand{using}\isamarkupfalse%
\ t{\isadigit{2}}\ \isacommand{by}\isamarkupfalse%
\ simp\isanewline
\ \ \isacommand{next}\isamarkupfalse%
\isanewline
\ \ \ \ \isacommand{case}\isamarkupfalse%
\ False\isanewline
\ \ \ \ \isacommand{then}\isamarkupfalse%
\ \isacommand{have}\isamarkupfalse%
\ equ{\isacharunderscore}{\kern0pt}ab{\isacharunderscore}{\kern0pt}pi{\isadigit{1}}{\isacharcolon}{\kern0pt}\ {\isachardoublequoteopen}nabla\ {\isasymturnstile}\ swap\ pi{\isadigit{1}}\ t{\isadigit{1}}{\isacharprime}{\kern0pt}\ {\isasymapprox}\ swap\ {\isacharbrackleft}{\kern0pt}{\isacharparenleft}{\kern0pt}swapas\ pi{\isadigit{1}}\ a{\isacharcomma}{\kern0pt}\ b{\isacharparenright}{\kern0pt}{\isacharbrackright}{\kern0pt}\ t{\isadigit{2}}{\isacharprime}{\kern0pt}{\isachardoublequoteclose}\ \isanewline
\ \ \ \ \ \ \isakeyword{and}\ fresh{\isacharunderscore}{\kern0pt}ab{\isacharunderscore}{\kern0pt}pi{\isadigit{1}}{\isacharcolon}{\kern0pt}\ {\isachardoublequoteopen}nabla\ {\isasymturnstile}\ swapas\ pi{\isadigit{1}}\ a\ {\isasymsharp}\ t{\isadigit{2}}{\isacharprime}{\kern0pt}{\isachardoublequoteclose}\isanewline
\ \ \ \ \ \ \isacommand{using}\isamarkupfalse%
\ i\ Equ{\isacharunderscore}{\kern0pt}elims{\isacharparenleft}{\kern0pt}{\isadigit{7}}{\isacharparenright}{\kern0pt}{\isacharbrackleft}{\kern0pt}of\ nabla\ {\isacartoucheopen}swapas\ pi{\isadigit{1}}\ a{\isacartoucheclose}\ {\isacartoucheopen}swap\ pi{\isadigit{1}}\ t{\isadigit{1}}{\isacharprime}{\kern0pt}{\isacartoucheclose}\ b\ t{\isadigit{2}}{\isacharprime}{\kern0pt}{\isacharbrackright}{\kern0pt}\ \isanewline
\ \ \ \ \ \ \isacommand{by}\isamarkupfalse%
\ blast{\isacharplus}{\kern0pt}\isanewline
\ \ \ \ \isacommand{have}\isamarkupfalse%
\ equ{\isacharunderscore}{\kern0pt}ab{\isacharunderscore}{\kern0pt}pi{\isadigit{2}}{\isacharcolon}{\kern0pt}\ {\isachardoublequoteopen}nabla\ {\isasymturnstile}\ swap\ pi{\isadigit{2}}\ t{\isadigit{1}}{\isacharprime}{\kern0pt}\ {\isasymapprox}\ swap\ {\isacharbrackleft}{\kern0pt}{\isacharparenleft}{\kern0pt}swapas\ pi{\isadigit{2}}\ a{\isacharcomma}{\kern0pt}\ b{\isacharparenright}{\kern0pt}{\isacharbrackright}{\kern0pt}\ t{\isadigit{2}}{\isacharprime}{\kern0pt}{\isachardoublequoteclose}\isanewline
\ \ \ \ \ \ \isacommand{using}\isamarkupfalse%
\ equ{\isacharunderscore}{\kern0pt}ab{\isacharunderscore}{\kern0pt}pi{\isadigit{1}}\ ds{\isacharunderscore}{\kern0pt}swapas{\isacharunderscore}{\kern0pt}eq{\isacharbrackleft}{\kern0pt}OF\ Abst{\isacharparenleft}{\kern0pt}{\isadigit{3}}{\isacharparenright}{\kern0pt}{\isacharcomma}{\kern0pt}\ of\ a{\isacharbrackright}{\kern0pt}\ \isanewline
\ \ \ \ \ \ \ \ Abst{\isacharparenleft}{\kern0pt}{\isadigit{3}}{\isacharparenright}{\kern0pt}\ Abst{\isacharparenleft}{\kern0pt}{\isadigit{1}}{\isacharparenright}{\kern0pt}{\isacharbrackleft}{\kern0pt}of\ {\isacartoucheopen}swap\ {\isacharbrackleft}{\kern0pt}{\isacharparenleft}{\kern0pt}swapas\ pi{\isadigit{1}}\ a{\isacharcomma}{\kern0pt}\ b{\isacharparenright}{\kern0pt}{\isacharbrackright}{\kern0pt}\ t{\isadigit{2}}{\isacharprime}{\kern0pt}{\isacartoucheclose}{\isacharbrackright}{\kern0pt}\ \isacommand{by}\isamarkupfalse%
\ auto\isanewline
\ \ \ \ \isacommand{have}\isamarkupfalse%
\ fresh{\isacharunderscore}{\kern0pt}ab{\isacharunderscore}{\kern0pt}pi{\isadigit{2}}{\isacharcolon}{\kern0pt}\ {\isachardoublequoteopen}nabla\ {\isasymturnstile}\ swapas\ pi{\isadigit{2}}\ a\ {\isasymsharp}\ t{\isadigit{2}}{\isacharprime}{\kern0pt}{\isachardoublequoteclose}\isanewline
\ \ \ \ \ \ \isacommand{using}\isamarkupfalse%
\ fresh{\isacharunderscore}{\kern0pt}ab{\isacharunderscore}{\kern0pt}pi{\isadigit{1}}\ ds{\isacharunderscore}{\kern0pt}swapas{\isacharunderscore}{\kern0pt}eq{\isacharbrackleft}{\kern0pt}OF\ Abst{\isacharparenleft}{\kern0pt}{\isadigit{3}}{\isacharparenright}{\kern0pt}{\isacharcomma}{\kern0pt}\ of\ a{\isacharbrackright}{\kern0pt}\ \isacommand{by}\isamarkupfalse%
\ simp\isanewline
\ \ \ \ \isacommand{with}\isamarkupfalse%
\ equ{\isacharunderscore}{\kern0pt}ab{\isacharunderscore}{\kern0pt}pi{\isadigit{2}}\ \isacommand{have}\isamarkupfalse%
\ {\isachardoublequoteopen}nabla\ {\isasymturnstile}\ Abst\ {\isacharparenleft}{\kern0pt}swapas\ pi{\isadigit{2}}\ a{\isacharparenright}{\kern0pt}\ {\isacharparenleft}{\kern0pt}swap\ pi{\isadigit{2}}\ t{\isadigit{1}}{\isacharprime}{\kern0pt}{\isacharparenright}{\kern0pt}\ {\isasymapprox}\ Abst\ b\ t{\isadigit{2}}{\isacharprime}{\kern0pt}{\isachardoublequoteclose}\isanewline
\ \ \ \ \ \ \isacommand{using}\isamarkupfalse%
\ equ{\isacharunderscore}{\kern0pt}abst{\isacharunderscore}{\kern0pt}ab{\isacharbrackleft}{\kern0pt}OF\ False{\isacharcomma}{\kern0pt}\ of\ nabla\ t{\isadigit{2}}{\isacharprime}{\kern0pt}{\isacharbrackright}{\kern0pt}\ ds{\isacharunderscore}{\kern0pt}swapas{\isacharunderscore}{\kern0pt}eq{\isacharbrackleft}{\kern0pt}OF\ Abst{\isacharparenleft}{\kern0pt}{\isadigit{3}}{\isacharparenright}{\kern0pt}{\isacharcomma}{\kern0pt}\ of\ a{\isacharbrackright}{\kern0pt}\ \isacommand{by}\isamarkupfalse%
\ simp\isanewline
\ \ \ \ \isacommand{then}\isamarkupfalse%
\ \isacommand{show}\isamarkupfalse%
\ {\isacharquery}{\kern0pt}thesis\ \isacommand{using}\isamarkupfalse%
\ t{\isadigit{2}}\ \isacommand{by}\isamarkupfalse%
\ simp\isanewline
\ \ \isacommand{qed}\isamarkupfalse%
\isanewline
\isacommand{next}\isamarkupfalse%
\isanewline
\ \ \isacommand{case}\isamarkupfalse%
\ {\isacharparenleft}{\kern0pt}Susp\ pi\ X{\isacharparenright}{\kern0pt}\isanewline
\ \ \isacommand{then}\isamarkupfalse%
\ \isacommand{obtain}\isamarkupfalse%
\ pi{\isacharprime}{\kern0pt}\isanewline
\ \ \ \ \isakeyword{where}\ t{\isadigit{2}}{\isacharcolon}{\kern0pt}\ {\isachardoublequoteopen}t{\isadigit{2}}\ {\isacharequal}{\kern0pt}\ Susp\ pi{\isacharprime}{\kern0pt}\ X{\isachardoublequoteclose}\isanewline
\ \ \ \ \isacommand{by}\isamarkupfalse%
\ {\isacharparenleft}{\kern0pt}auto\ elim{\isacharcolon}{\kern0pt}equ{\isachardot}{\kern0pt}cases{\isacharparenright}{\kern0pt}\isanewline
\ \ \isacommand{with}\isamarkupfalse%
\ Susp\ \isacommand{have}\isamarkupfalse%
\ i{\isacharcolon}{\kern0pt}\ {\isachardoublequoteopen}nabla\ {\isasymturnstile}\ Susp\ {\isacharparenleft}{\kern0pt}pi{\isadigit{1}}\ {\isacharat}{\kern0pt}\ pi{\isacharparenright}{\kern0pt}\ X\ {\isasymapprox}\ Susp\ pi{\isacharprime}{\kern0pt}\ X{\isachardoublequoteclose}\ \isanewline
\ \ \ \ \isacommand{by}\isamarkupfalse%
\ simp\isanewline
\ \ \isacommand{then}\isamarkupfalse%
\ \isacommand{have}\isamarkupfalse%
\ {\isachardoublequoteopen}{\isasymforall}\ c\ {\isasymin}\ ds\ {\isacharparenleft}{\kern0pt}pi{\isadigit{1}}{\isacharat}{\kern0pt}pi{\isacharparenright}{\kern0pt}\ pi{\isacharprime}{\kern0pt}{\isachardot}{\kern0pt}\ {\isacharparenleft}{\kern0pt}c{\isacharcomma}{\kern0pt}X{\isacharparenright}{\kern0pt}\ {\isasymin}\ nabla{\isachardoublequoteclose}\isanewline
\ \ \ \ \isacommand{using}\isamarkupfalse%
\ Equ{\isacharunderscore}{\kern0pt}elims{\isacharparenleft}{\kern0pt}{\isadigit{3}}{\isacharparenright}{\kern0pt}{\isacharbrackleft}{\kern0pt}OF\ i{\isacharbrackright}{\kern0pt}\ \isacommand{by}\isamarkupfalse%
\ auto\isanewline
\ \ \isacommand{with}\isamarkupfalse%
\ Susp\ \isacommand{have}\isamarkupfalse%
\ {\isachardoublequoteopen}{\isasymforall}\ c\ {\isasymin}\ ds\ {\isacharparenleft}{\kern0pt}pi{\isadigit{2}}{\isacharat}{\kern0pt}pi{\isacharparenright}{\kern0pt}\ pi{\isacharprime}{\kern0pt}{\isachardot}{\kern0pt}\ {\isacharparenleft}{\kern0pt}c{\isacharcomma}{\kern0pt}X{\isacharparenright}{\kern0pt}\ {\isasymin}\ nabla{\isachardoublequoteclose}\isanewline
\ \ \ \ \isacommand{using}\isamarkupfalse%
\ ds{\isacharunderscore}{\kern0pt}cancel{\isacharunderscore}{\kern0pt}pi{\isacharunderscore}{\kern0pt}left\ ds{\isacharunderscore}{\kern0pt}sym\ ds{\isacharunderscore}{\kern0pt}trans\isanewline
\ \ \ \ \isacommand{by}\isamarkupfalse%
\ blast\isanewline
\ \ \isacommand{then}\isamarkupfalse%
\ \isacommand{have}\isamarkupfalse%
\ {\isachardoublequoteopen}nabla\ {\isasymturnstile}\ Susp\ {\isacharparenleft}{\kern0pt}pi{\isadigit{2}}\ {\isacharat}{\kern0pt}\ pi{\isacharparenright}{\kern0pt}\ X\ {\isasymapprox}\ Susp\ pi{\isacharprime}{\kern0pt}\ X{\isachardoublequoteclose}\isanewline
\ \ \ \ \isacommand{using}\isamarkupfalse%
\ equ{\isacharunderscore}{\kern0pt}susp\ \isacommand{by}\isamarkupfalse%
\ simp\isanewline
\ \ \isacommand{then}\isamarkupfalse%
\ \isacommand{show}\isamarkupfalse%
\ {\isacharquery}{\kern0pt}case\ \isacommand{using}\isamarkupfalse%
\ t{\isadigit{2}}\ swap{\isachardot}{\kern0pt}simps{\isacharparenleft}{\kern0pt}{\isadigit{3}}{\isacharparenright}{\kern0pt}\ \isacommand{by}\isamarkupfalse%
\ simp\isanewline
\isacommand{next}\isamarkupfalse%
\isanewline
\ \ \isacommand{case}\isamarkupfalse%
\ {\isacharparenleft}{\kern0pt}Atom\ a{\isacharparenright}{\kern0pt}\isanewline
\ \ \isacommand{then}\isamarkupfalse%
\ \isacommand{show}\isamarkupfalse%
\ {\isacharquery}{\kern0pt}case\isanewline
\ \ \ \ \isacommand{using}\isamarkupfalse%
\ ds{\isacharunderscore}{\kern0pt}swapas{\isacharunderscore}{\kern0pt}eq{\isacharbrackleft}{\kern0pt}OF\ Atom{\isacharparenleft}{\kern0pt}{\isadigit{2}}{\isacharparenright}{\kern0pt}{\isacharcomma}{\kern0pt}\ of\ a{\isacharbrackright}{\kern0pt}\isanewline
\ \ \ \ \isacommand{by}\isamarkupfalse%
\ {\isacharparenleft}{\kern0pt}auto\ elim{\isacharcolon}{\kern0pt}\ equ{\isachardot}{\kern0pt}cases{\isacharparenright}{\kern0pt}\isanewline
\isacommand{qed}\isamarkupfalse%
\ {\isacharparenleft}{\kern0pt}auto\ elim{\isacharcolon}{\kern0pt}\ equ{\isachardot}{\kern0pt}cases{\isacharparenright}{\kern0pt}%
\endisatagproof
{\isafoldproof}%
%
\isadelimproof
\isanewline
%
\endisadelimproof
\isanewline
\isanewline
\isacommand{lemma}\isamarkupfalse%
\ ds{\isacharunderscore}{\kern0pt}empty{\isacharunderscore}{\kern0pt}equiv{\isacharunderscore}{\kern0pt}{\isadigit{2}}{\isacharcolon}{\kern0pt}\isanewline
\ \ \isakeyword{assumes}\ {\isachardoublequoteopen}ds\ pi{\isadigit{1}}\ pi{\isadigit{2}}\ {\isacharequal}{\kern0pt}\ {\isacharbraceleft}{\kern0pt}{\isacharbraceright}{\kern0pt}{\isachardoublequoteclose}\isanewline
\ \ \isakeyword{shows}\ {\isachardoublequoteopen}nabla\ {\isasymturnstile}\ t{\isadigit{1}}\ {\isasymapprox}\ swap\ pi{\isadigit{1}}\ t{\isadigit{2}}\ {\isasymLongrightarrow}\ nabla\ {\isasymturnstile}\ t{\isadigit{1}}\ {\isasymapprox}\ swap\ pi{\isadigit{2}}\ t{\isadigit{2}}{\isachardoublequoteclose}\isanewline
%
\isadelimproof
%
\endisadelimproof
%
\isatagproof
\isacommand{using}\isamarkupfalse%
\ assms\isanewline
\isacommand{proof}\isamarkupfalse%
{\isacharparenleft}{\kern0pt}induction\ t{\isadigit{2}}\ arbitrary{\isacharcolon}{\kern0pt}\ pi{\isadigit{1}}\ pi{\isadigit{2}}\ t{\isadigit{1}}{\isacharparenright}{\kern0pt}\isanewline
\ \ \isacommand{case}\isamarkupfalse%
\ {\isacharparenleft}{\kern0pt}Abst\ b\ t{\isadigit{2}}{\isacharprime}{\kern0pt}{\isacharparenright}{\kern0pt}\isanewline
\ \ \isacommand{then}\isamarkupfalse%
\ \isacommand{obtain}\isamarkupfalse%
\ a\ t{\isadigit{1}}{\isacharprime}{\kern0pt}\isanewline
\ \ \ \ \isakeyword{where}\ t{\isadigit{1}}{\isacharcolon}{\kern0pt}\ {\isachardoublequoteopen}t{\isadigit{1}}\ {\isacharequal}{\kern0pt}\ Abst\ a\ t{\isadigit{1}}{\isacharprime}{\kern0pt}{\isachardoublequoteclose}\isanewline
\ \ \ \ \isacommand{by}\isamarkupfalse%
{\isacharparenleft}{\kern0pt}auto\ elim{\isacharcolon}{\kern0pt}\ equ{\isachardot}{\kern0pt}cases{\isacharparenright}{\kern0pt}\isanewline
\ \ \isacommand{with}\isamarkupfalse%
\ Abst\ \isacommand{have}\isamarkupfalse%
\ i{\isacharcolon}{\kern0pt}\ {\isachardoublequoteopen}nabla\ {\isasymturnstile}\ Abst\ a\ t{\isadigit{1}}{\isacharprime}{\kern0pt}\ {\isasymapprox}\ Abst\ {\isacharparenleft}{\kern0pt}swapas\ pi{\isadigit{1}}\ b{\isacharparenright}{\kern0pt}\ {\isacharparenleft}{\kern0pt}swap\ pi{\isadigit{1}}\ t{\isadigit{2}}{\isacharprime}{\kern0pt}{\isacharparenright}{\kern0pt}{\isachardoublequoteclose}\isanewline
\ \ \ \ \isacommand{using}\isamarkupfalse%
\ swap{\isachardot}{\kern0pt}simps{\isacharparenleft}{\kern0pt}{\isadigit{4}}{\isacharparenright}{\kern0pt}\ \isacommand{by}\isamarkupfalse%
\ simp\isanewline
\ \ \isacommand{then}\isamarkupfalse%
\ \isacommand{show}\isamarkupfalse%
\ {\isacharquery}{\kern0pt}case\isanewline
\ \ \isacommand{proof}\isamarkupfalse%
{\isacharparenleft}{\kern0pt}cases\ {\isacartoucheopen}swapas\ pi{\isadigit{1}}\ b\ {\isacharequal}{\kern0pt}\ a{\isacartoucheclose}{\isacharparenright}{\kern0pt}\isanewline
\ \ \ \ \isacommand{case}\isamarkupfalse%
\ True\isanewline
\ \ \ \ \isacommand{then}\isamarkupfalse%
\ \isacommand{have}\isamarkupfalse%
\ {\isachardoublequoteopen}nabla\ {\isasymturnstile}\ t{\isadigit{1}}{\isacharprime}{\kern0pt}\ {\isasymapprox}\ swap\ pi{\isadigit{1}}\ t{\isadigit{2}}{\isacharprime}{\kern0pt}{\isachardoublequoteclose}\ \isanewline
\ \ \ \ \ \ \isacommand{using}\isamarkupfalse%
\ Abst\ i\ Equ{\isacharunderscore}{\kern0pt}elims{\isacharparenleft}{\kern0pt}{\isadigit{6}}{\isacharparenright}{\kern0pt}\ \isacommand{by}\isamarkupfalse%
\ blast\isanewline
\ \ \ \ \isacommand{then}\isamarkupfalse%
\ \isacommand{have}\isamarkupfalse%
\ ii{\isacharcolon}{\kern0pt}\ {\isachardoublequoteopen}nabla\ {\isasymturnstile}\ t{\isadigit{1}}{\isacharprime}{\kern0pt}\ {\isasymapprox}\ swap\ pi{\isadigit{2}}\ t{\isadigit{2}}{\isacharprime}{\kern0pt}{\isachardoublequoteclose}\isanewline
\ \ \ \ \ \ \isacommand{using}\isamarkupfalse%
\ Abst\ \isacommand{by}\isamarkupfalse%
\ simp\isanewline
\ \ \ \ \isacommand{then}\isamarkupfalse%
\ \isacommand{have}\isamarkupfalse%
\ {\isachardoublequoteopen}nabla\ {\isasymturnstile}\ Abst\ a\ t{\isadigit{1}}{\isacharprime}{\kern0pt}\ {\isasymapprox}\ Abst\ {\isacharparenleft}{\kern0pt}swapas\ pi{\isadigit{2}}\ b{\isacharparenright}{\kern0pt}\ {\isacharparenleft}{\kern0pt}swap\ pi{\isadigit{2}}\ t{\isadigit{2}}{\isacharprime}{\kern0pt}{\isacharparenright}{\kern0pt}{\isachardoublequoteclose}\isanewline
\ \ \ \ \ \ \isacommand{using}\isamarkupfalse%
\ True\ ds{\isacharunderscore}{\kern0pt}swapas{\isacharunderscore}{\kern0pt}eq{\isacharbrackleft}{\kern0pt}OF\ Abst{\isacharparenleft}{\kern0pt}{\isadigit{3}}{\isacharparenright}{\kern0pt}{\isacharcomma}{\kern0pt}\ of\ b{\isacharbrackright}{\kern0pt}\ equ{\isacharunderscore}{\kern0pt}abst{\isacharunderscore}{\kern0pt}aa{\isacharbrackleft}{\kern0pt}OF\ ii{\isacharcomma}{\kern0pt}\ of\ {\isacartoucheopen}swapas\ pi{\isadigit{2}}\ b{\isacartoucheclose}{\isacharbrackright}{\kern0pt}\ \isanewline
\ \ \ \ \ \ \isacommand{by}\isamarkupfalse%
\ simp\isanewline
\ \ \ \ \isacommand{then}\isamarkupfalse%
\ \isacommand{show}\isamarkupfalse%
\ {\isacharquery}{\kern0pt}thesis\ \isanewline
\ \ \ \ \ \ \isacommand{using}\isamarkupfalse%
\ t{\isadigit{1}}\ \isacommand{by}\isamarkupfalse%
\ simp\isanewline
\ \ \isacommand{next}\isamarkupfalse%
\isanewline
\ \ \ \ \isacommand{case}\isamarkupfalse%
\ False\isanewline
\ \ \ \ \isacommand{then}\isamarkupfalse%
\ \isacommand{have}\isamarkupfalse%
\ equ{\isacharunderscore}{\kern0pt}ab{\isacharcolon}{\kern0pt}\ {\isachardoublequoteopen}nabla\ {\isasymturnstile}\ t{\isadigit{1}}{\isacharprime}{\kern0pt}\ {\isasymapprox}\ swap\ {\isacharbrackleft}{\kern0pt}{\isacharparenleft}{\kern0pt}a{\isacharcomma}{\kern0pt}\ swapas\ pi{\isadigit{1}}\ b{\isacharparenright}{\kern0pt}{\isacharbrackright}{\kern0pt}\ {\isacharparenleft}{\kern0pt}swap\ pi{\isadigit{1}}\ t{\isadigit{2}}{\isacharprime}{\kern0pt}{\isacharparenright}{\kern0pt}{\isachardoublequoteclose}\ \isanewline
\ \ \ \ \ \ \isakeyword{and}\ fresh{\isacharunderscore}{\kern0pt}ab{\isacharunderscore}{\kern0pt}pi{\isadigit{1}}{\isacharcolon}{\kern0pt}\ {\isachardoublequoteopen}nabla\ {\isasymturnstile}\ a\ {\isasymsharp}\ swap\ pi{\isadigit{1}}\ t{\isadigit{2}}{\isacharprime}{\kern0pt}{\isachardoublequoteclose}\isanewline
\ \ \ \ \ \ \isacommand{using}\isamarkupfalse%
\ i\ Equ{\isacharunderscore}{\kern0pt}elims{\isacharparenleft}{\kern0pt}{\isadigit{7}}{\isacharparenright}{\kern0pt}\ \isacommand{by}\isamarkupfalse%
\ blast{\isacharplus}{\kern0pt}\isanewline
\isanewline
\ \ \ \ \isacommand{then}\isamarkupfalse%
\ \isacommand{have}\isamarkupfalse%
\ equ{\isacharunderscore}{\kern0pt}ab{\isacharunderscore}{\kern0pt}pi{\isadigit{1}}{\isacharcolon}{\kern0pt}\ {\isachardoublequoteopen}nabla\ {\isasymturnstile}\ t{\isadigit{1}}{\isacharprime}{\kern0pt}\ {\isasymapprox}\ swap\ {\isacharparenleft}{\kern0pt}{\isacharbrackleft}{\kern0pt}{\isacharparenleft}{\kern0pt}a{\isacharcomma}{\kern0pt}\ swapas\ pi{\isadigit{1}}\ b{\isacharparenright}{\kern0pt}{\isacharbrackright}{\kern0pt}{\isacharat}{\kern0pt}\ pi{\isadigit{1}}{\isacharparenright}{\kern0pt}\ t{\isadigit{2}}{\isacharprime}{\kern0pt}{\isachardoublequoteclose}\isanewline
\ \ \ \ \ \ \isacommand{using}\isamarkupfalse%
\ swap{\isacharunderscore}{\kern0pt}append{\isacharbrackleft}{\kern0pt}of\ {\isacartoucheopen}{\isacharbrackleft}{\kern0pt}{\isacharparenleft}{\kern0pt}a{\isacharcomma}{\kern0pt}\ swapas\ pi{\isadigit{1}}\ b{\isacharparenright}{\kern0pt}{\isacharbrackright}{\kern0pt}{\isacartoucheclose}\ pi{\isadigit{1}}\ t{\isadigit{2}}{\isacharprime}{\kern0pt}{\isacharbrackright}{\kern0pt}\ \isacommand{by}\isamarkupfalse%
\ simp\isanewline
\isanewline
\ \ \ \ \isacommand{have}\isamarkupfalse%
\ ds{\isacharunderscore}{\kern0pt}empty{\isacharcolon}{\kern0pt}\ {\isachardoublequoteopen}ds\ {\isacharparenleft}{\kern0pt}{\isacharbrackleft}{\kern0pt}{\isacharparenleft}{\kern0pt}a{\isacharcomma}{\kern0pt}\ swapas\ pi{\isadigit{1}}\ b{\isacharparenright}{\kern0pt}{\isacharbrackright}{\kern0pt}\ {\isacharat}{\kern0pt}\ pi{\isadigit{1}}{\isacharparenright}{\kern0pt}\ {\isacharparenleft}{\kern0pt}{\isacharbrackleft}{\kern0pt}{\isacharparenleft}{\kern0pt}a{\isacharcomma}{\kern0pt}\ swapas\ pi{\isadigit{2}}\ b{\isacharparenright}{\kern0pt}{\isacharbrackright}{\kern0pt}\ {\isacharat}{\kern0pt}\ pi{\isadigit{2}}{\isacharparenright}{\kern0pt}\ {\isacharequal}{\kern0pt}\ {\isacharbraceleft}{\kern0pt}{\isacharbraceright}{\kern0pt}{\isachardoublequoteclose}\isanewline
\ \ \ \ \ \ \isacommand{using}\isamarkupfalse%
\ Abst{\isacharparenleft}{\kern0pt}{\isadigit{3}}{\isacharparenright}{\kern0pt}\ ds{\isacharunderscore}{\kern0pt}swapas{\isacharunderscore}{\kern0pt}eq{\isacharbrackleft}{\kern0pt}OF\ Abst{\isacharparenleft}{\kern0pt}{\isadigit{3}}{\isacharparenright}{\kern0pt}{\isacharcomma}{\kern0pt}\ of\ b{\isacharbrackright}{\kern0pt}\ \isanewline
\ \ \ \ \ \ \ \ ds{\isacharunderscore}{\kern0pt}cancel{\isacharunderscore}{\kern0pt}pi{\isacharunderscore}{\kern0pt}front{\isacharbrackleft}{\kern0pt}of\ {\isacartoucheopen}{\isacharbrackleft}{\kern0pt}{\isacharparenleft}{\kern0pt}a{\isacharcomma}{\kern0pt}\ swapas\ pi{\isadigit{1}}\ b{\isacharparenright}{\kern0pt}{\isacharbrackright}{\kern0pt}{\isacartoucheclose}\ pi{\isadigit{1}}\ pi{\isadigit{2}}{\isacharbrackright}{\kern0pt}\ \isacommand{by}\isamarkupfalse%
\ simp\isanewline
\ \ \ \ \isanewline
\ \ \ \ \isacommand{hence}\isamarkupfalse%
\ {\isachardoublequoteopen}nabla\ {\isasymturnstile}\ t{\isadigit{1}}{\isacharprime}{\kern0pt}\ {\isasymapprox}\ swap\ {\isacharparenleft}{\kern0pt}{\isacharbrackleft}{\kern0pt}{\isacharparenleft}{\kern0pt}a{\isacharcomma}{\kern0pt}\ swapas\ pi{\isadigit{2}}\ b{\isacharparenright}{\kern0pt}{\isacharbrackright}{\kern0pt}{\isacharat}{\kern0pt}\ pi{\isadigit{2}}{\isacharparenright}{\kern0pt}\ t{\isadigit{2}}{\isacharprime}{\kern0pt}{\isachardoublequoteclose}\isanewline
\ \ \ \ \ \ \isacommand{using}\isamarkupfalse%
\ Abst{\isacharparenleft}{\kern0pt}{\isadigit{1}}{\isacharparenright}{\kern0pt}{\isacharbrackleft}{\kern0pt}OF\ equ{\isacharunderscore}{\kern0pt}ab{\isacharunderscore}{\kern0pt}pi{\isadigit{1}}\ ds{\isacharunderscore}{\kern0pt}empty{\isacharbrackright}{\kern0pt}\ \isacommand{by}\isamarkupfalse%
\ simp\isanewline
\ \ \ \ \isacommand{hence}\isamarkupfalse%
\ equ{\isacharunderscore}{\kern0pt}ab{\isacharunderscore}{\kern0pt}pi{\isadigit{2}}{\isacharcolon}{\kern0pt}\ {\isachardoublequoteopen}nabla\ {\isasymturnstile}\ t{\isadigit{1}}{\isacharprime}{\kern0pt}\ {\isasymapprox}\ swap\ {\isacharbrackleft}{\kern0pt}{\isacharparenleft}{\kern0pt}a{\isacharcomma}{\kern0pt}\ swapas\ pi{\isadigit{2}}\ b{\isacharparenright}{\kern0pt}{\isacharbrackright}{\kern0pt}\ {\isacharparenleft}{\kern0pt}swap\ pi{\isadigit{2}}\ t{\isadigit{2}}{\isacharprime}{\kern0pt}{\isacharparenright}{\kern0pt}{\isachardoublequoteclose}\isanewline
\ \ \ \ \ \ \isacommand{using}\isamarkupfalse%
\ swap{\isacharunderscore}{\kern0pt}append{\isacharbrackleft}{\kern0pt}of\ {\isacartoucheopen}{\isacharbrackleft}{\kern0pt}{\isacharparenleft}{\kern0pt}a{\isacharcomma}{\kern0pt}\ swapas\ pi{\isadigit{2}}\ b{\isacharparenright}{\kern0pt}{\isacharbrackright}{\kern0pt}{\isacartoucheclose}\ pi{\isadigit{2}}\ t{\isadigit{2}}{\isacharprime}{\kern0pt}{\isacharbrackright}{\kern0pt}\ \isacommand{by}\isamarkupfalse%
\ simp\isanewline
\isanewline
\ \ \ \ \isacommand{have}\isamarkupfalse%
\ fresh{\isacharunderscore}{\kern0pt}ab{\isacharunderscore}{\kern0pt}pi{\isadigit{2}}{\isacharcolon}{\kern0pt}\ {\isachardoublequoteopen}nabla\ {\isasymturnstile}\ a\ {\isasymsharp}\ swap\ pi{\isadigit{2}}\ t{\isadigit{2}}{\isacharprime}{\kern0pt}{\isachardoublequoteclose}\isanewline
\ \ \ \ \ \ \isacommand{using}\isamarkupfalse%
\ ds{\isacharunderscore}{\kern0pt}empty{\isacharunderscore}{\kern0pt}fresh{\isacharunderscore}{\kern0pt}{\isadigit{2}}\ Abst{\isacharparenleft}{\kern0pt}{\isadigit{3}}{\isacharparenright}{\kern0pt}\ fresh{\isacharunderscore}{\kern0pt}ab{\isacharunderscore}{\kern0pt}pi{\isadigit{1}}\ \isacommand{by}\isamarkupfalse%
\ auto\isanewline
\ \ \ \ \isacommand{with}\isamarkupfalse%
\ equ{\isacharunderscore}{\kern0pt}ab{\isacharunderscore}{\kern0pt}pi{\isadigit{2}}\ \isacommand{have}\isamarkupfalse%
\ {\isachardoublequoteopen}nabla\ {\isasymturnstile}\ Abst\ a\ t{\isadigit{1}}{\isacharprime}{\kern0pt}\ {\isasymapprox}\ Abst\ {\isacharparenleft}{\kern0pt}swapas\ pi{\isadigit{2}}\ b{\isacharparenright}{\kern0pt}\ {\isacharparenleft}{\kern0pt}swap\ pi{\isadigit{2}}\ t{\isadigit{2}}{\isacharprime}{\kern0pt}{\isacharparenright}{\kern0pt}{\isachardoublequoteclose}\isanewline
\ \ \ \ \ \ \isacommand{using}\isamarkupfalse%
\ equ{\isacharunderscore}{\kern0pt}abst{\isacharunderscore}{\kern0pt}ab\ ds{\isacharunderscore}{\kern0pt}swapas{\isacharunderscore}{\kern0pt}eq\ False\ assms\ Abst{\isacharparenleft}{\kern0pt}{\isadigit{3}}{\isacharparenright}{\kern0pt}\ \isacommand{by}\isamarkupfalse%
\ auto\isanewline
\ \ \ \ \isacommand{then}\isamarkupfalse%
\ \isacommand{show}\isamarkupfalse%
\ {\isacharquery}{\kern0pt}thesis\ \isacommand{using}\isamarkupfalse%
\ t{\isadigit{1}}\ \isacommand{by}\isamarkupfalse%
\ simp\isanewline
\ \ \isacommand{qed}\isamarkupfalse%
\isanewline
\isacommand{next}\isamarkupfalse%
\isanewline
\ \ \isacommand{case}\isamarkupfalse%
\ {\isacharparenleft}{\kern0pt}Susp\ pi{\isacharprime}{\kern0pt}\ X{\isacharparenright}{\kern0pt}\isanewline
\ \ \isacommand{then}\isamarkupfalse%
\ \isacommand{obtain}\isamarkupfalse%
\ pi\isanewline
\ \ \ \ \isakeyword{where}\ t{\isadigit{1}}{\isacharcolon}{\kern0pt}\ {\isachardoublequoteopen}t{\isadigit{1}}\ {\isacharequal}{\kern0pt}\ Susp\ pi\ X{\isachardoublequoteclose}\isanewline
\ \ \ \ \isacommand{by}\isamarkupfalse%
\ {\isacharparenleft}{\kern0pt}auto\ elim{\isacharcolon}{\kern0pt}equ{\isachardot}{\kern0pt}cases{\isacharparenright}{\kern0pt}\isanewline
\ \ \isacommand{with}\isamarkupfalse%
\ Susp\ \isacommand{have}\isamarkupfalse%
\ i{\isacharcolon}{\kern0pt}\ {\isachardoublequoteopen}nabla\ {\isasymturnstile}\ Susp\ pi\ X\ {\isasymapprox}\ Susp\ {\isacharparenleft}{\kern0pt}pi{\isadigit{1}}\ {\isacharat}{\kern0pt}\ pi{\isacharprime}{\kern0pt}{\isacharparenright}{\kern0pt}\ X{\isachardoublequoteclose}\ \isanewline
\ \ \ \ \isacommand{by}\isamarkupfalse%
\ simp\isanewline
\ \ \isacommand{then}\isamarkupfalse%
\ \isacommand{have}\isamarkupfalse%
\ {\isachardoublequoteopen}{\isasymforall}\ c\ {\isasymin}\ ds\ pi\ {\isacharparenleft}{\kern0pt}pi{\isadigit{1}}\ {\isacharat}{\kern0pt}\ pi{\isacharprime}{\kern0pt}{\isacharparenright}{\kern0pt}{\isachardot}{\kern0pt}\ {\isacharparenleft}{\kern0pt}c{\isacharcomma}{\kern0pt}X{\isacharparenright}{\kern0pt}\ {\isasymin}\ nabla{\isachardoublequoteclose}\isanewline
\ \ \ \ \isacommand{using}\isamarkupfalse%
\ Equ{\isacharunderscore}{\kern0pt}elims{\isacharparenleft}{\kern0pt}{\isadigit{3}}{\isacharparenright}{\kern0pt}{\isacharbrackleft}{\kern0pt}OF\ i{\isacharbrackright}{\kern0pt}\ \isacommand{by}\isamarkupfalse%
\ auto\isanewline
\ \ \isacommand{with}\isamarkupfalse%
\ Susp\ \isacommand{have}\isamarkupfalse%
\ {\isachardoublequoteopen}{\isasymforall}\ c\ {\isasymin}\ ds\ pi\ {\isacharparenleft}{\kern0pt}pi{\isadigit{2}}\ {\isacharat}{\kern0pt}\ pi{\isacharprime}{\kern0pt}{\isacharparenright}{\kern0pt}{\isachardot}{\kern0pt}\ {\isacharparenleft}{\kern0pt}c{\isacharcomma}{\kern0pt}X{\isacharparenright}{\kern0pt}\ {\isasymin}\ nabla{\isachardoublequoteclose}\isanewline
\ \ \ \ \isacommand{using}\isamarkupfalse%
\ ds{\isacharunderscore}{\kern0pt}cancel{\isacharunderscore}{\kern0pt}pi{\isacharunderscore}{\kern0pt}left\ ds{\isacharunderscore}{\kern0pt}sym\ ds{\isacharunderscore}{\kern0pt}trans\isanewline
\ \ \ \ \isacommand{by}\isamarkupfalse%
\ blast\isanewline
\ \ \isacommand{then}\isamarkupfalse%
\ \isacommand{have}\isamarkupfalse%
\ {\isachardoublequoteopen}nabla\ {\isasymturnstile}\ Susp\ pi\ X\ {\isasymapprox}\ Susp\ {\isacharparenleft}{\kern0pt}pi{\isadigit{2}}\ {\isacharat}{\kern0pt}\ pi{\isacharprime}{\kern0pt}{\isacharparenright}{\kern0pt}\ X{\isachardoublequoteclose}\isanewline
\ \ \ \ \isacommand{using}\isamarkupfalse%
\ equ{\isacharunderscore}{\kern0pt}susp\ \isacommand{by}\isamarkupfalse%
\ simp\isanewline
\ \ \isacommand{then}\isamarkupfalse%
\ \isacommand{show}\isamarkupfalse%
\ {\isacharquery}{\kern0pt}case\ \isacommand{using}\isamarkupfalse%
\ t{\isadigit{1}}\ swap{\isachardot}{\kern0pt}simps{\isacharparenleft}{\kern0pt}{\isadigit{3}}{\isacharparenright}{\kern0pt}\ \isacommand{by}\isamarkupfalse%
\ simp\isanewline
\isacommand{next}\isamarkupfalse%
\isanewline
\ \ \isacommand{case}\isamarkupfalse%
\ {\isacharparenleft}{\kern0pt}Atom\ a{\isacharparenright}{\kern0pt}\isanewline
\ \ \isacommand{then}\isamarkupfalse%
\ \isacommand{show}\isamarkupfalse%
\ {\isacharquery}{\kern0pt}case\isanewline
\ \ \ \ \isacommand{using}\isamarkupfalse%
\ ds{\isacharunderscore}{\kern0pt}swapas{\isacharunderscore}{\kern0pt}eq\ \isacommand{by}\isamarkupfalse%
\ {\isacharparenleft}{\kern0pt}auto\ elim{\isacharcolon}{\kern0pt}\ equ{\isachardot}{\kern0pt}cases{\isacharparenright}{\kern0pt}\isanewline
\isacommand{qed}\isamarkupfalse%
\ {\isacharparenleft}{\kern0pt}auto\ elim{\isacharcolon}{\kern0pt}\ equ{\isachardot}{\kern0pt}cases{\isacharparenright}{\kern0pt}%
\endisatagproof
{\isafoldproof}%
%
\isadelimproof
\isanewline
%
\endisadelimproof
\isanewline
\isanewline
\isacommand{lemma}\isamarkupfalse%
\ equ{\isacharunderscore}{\kern0pt}equivariance{\isacharcolon}{\kern0pt}\isanewline
\ \ \isakeyword{assumes}\ {\isachardoublequoteopen}nabla\ {\isasymturnstile}\ t{\isadigit{1}}\ {\isasymapprox}\ t{\isadigit{2}}{\isachardoublequoteclose}\isanewline
\ \ \isakeyword{shows}\ {\isachardoublequoteopen}nabla\ {\isasymturnstile}\ swap\ pi\ t{\isadigit{1}}\ {\isasymapprox}\ swap\ pi\ t{\isadigit{2}}{\isachardoublequoteclose}\isanewline
%
\isadelimproof
\ \ %
\endisadelimproof
%
\isatagproof
\isacommand{using}\isamarkupfalse%
\ assms\isanewline
\isacommand{proof}\isamarkupfalse%
{\isacharparenleft}{\kern0pt}induct\ rule{\isacharcolon}{\kern0pt}\ equ{\isachardot}{\kern0pt}induct{\isacharparenright}{\kern0pt}\isanewline
\ \ \isacommand{case}\isamarkupfalse%
\ {\isacharparenleft}{\kern0pt}equ{\isacharunderscore}{\kern0pt}abst{\isacharunderscore}{\kern0pt}ab\ a\ b\ nabla\ t{\isadigit{2}}{\isacharprime}{\kern0pt}\ t{\isadigit{1}}{\isacharprime}{\kern0pt}{\isacharparenright}{\kern0pt}\isanewline
\ \ \isacommand{have}\isamarkupfalse%
\ i{\isacharcolon}{\kern0pt}\ {\isachardoublequoteopen}nabla\ {\isasymturnstile}\ swapas\ pi\ a\ {\isasymsharp}\ swap\ pi\ t{\isadigit{2}}{\isacharprime}{\kern0pt}{\isachardoublequoteclose}\isanewline
\ \ \ \ \isacommand{using}\isamarkupfalse%
\ fresh{\isacharunderscore}{\kern0pt}swap{\isacharunderscore}{\kern0pt}eqvt{\isacharbrackleft}{\kern0pt}OF\ equ{\isacharunderscore}{\kern0pt}abst{\isacharunderscore}{\kern0pt}ab{\isacharparenleft}{\kern0pt}{\isadigit{2}}{\isacharparenright}{\kern0pt}{\isacharcomma}{\kern0pt}\ of\ pi{\isacharbrackright}{\kern0pt}\ \isacommand{by}\isamarkupfalse%
\ simp\isanewline
\isanewline
\ \ \isacommand{have}\isamarkupfalse%
\ {\isachardoublequoteopen}nabla\ {\isasymturnstile}\ swap\ pi\ t{\isadigit{1}}{\isacharprime}{\kern0pt}\ {\isasymapprox}\ swap\ {\isacharparenleft}{\kern0pt}pi\ {\isacharat}{\kern0pt}\ {\isacharbrackleft}{\kern0pt}{\isacharparenleft}{\kern0pt}a{\isacharcomma}{\kern0pt}b{\isacharparenright}{\kern0pt}{\isacharbrackright}{\kern0pt}{\isacharparenright}{\kern0pt}\ t{\isadigit{2}}{\isacharprime}{\kern0pt}{\isachardoublequoteclose}\isanewline
\ \ \ \ \isacommand{using}\isamarkupfalse%
\ equ{\isacharunderscore}{\kern0pt}abst{\isacharunderscore}{\kern0pt}ab{\isacharparenleft}{\kern0pt}{\isadigit{4}}{\isacharparenright}{\kern0pt}\ swap{\isacharunderscore}{\kern0pt}append{\isacharbrackleft}{\kern0pt}of\ pi\ {\isacartoucheopen}{\isacharbrackleft}{\kern0pt}{\isacharparenleft}{\kern0pt}a{\isacharcomma}{\kern0pt}b{\isacharparenright}{\kern0pt}{\isacharbrackright}{\kern0pt}{\isacartoucheclose}\ t{\isadigit{2}}{\isacharprime}{\kern0pt}{\isacharbrackright}{\kern0pt}\ \isacommand{by}\isamarkupfalse%
\ simp\isanewline
\ \ \isacommand{then}\isamarkupfalse%
\ \isacommand{have}\isamarkupfalse%
\ {\isachardoublequoteopen}nabla\ {\isasymturnstile}\ swap\ pi\ t{\isadigit{1}}{\isacharprime}{\kern0pt}\ {\isasymapprox}\ swap\ {\isacharparenleft}{\kern0pt}{\isacharbrackleft}{\kern0pt}{\isacharparenleft}{\kern0pt}swapas\ pi\ a{\isacharcomma}{\kern0pt}\ swapas\ pi\ b{\isacharparenright}{\kern0pt}{\isacharbrackright}{\kern0pt}\ {\isacharat}{\kern0pt}\ pi{\isacharparenright}{\kern0pt}\ t{\isadigit{2}}{\isacharprime}{\kern0pt}{\isachardoublequoteclose}\isanewline
\ \ \ \ \isacommand{using}\isamarkupfalse%
\ ds{\isacharunderscore}{\kern0pt}empty{\isacharunderscore}{\kern0pt}equiv{\isacharunderscore}{\kern0pt}{\isadigit{2}}{\isacharbrackleft}{\kern0pt}OF\ ds{\isacharunderscore}{\kern0pt}comm{\isacharbrackleft}{\kern0pt}of\ pi\ a\ b{\isacharbrackright}{\kern0pt}{\isacharbrackright}{\kern0pt}\ \isacommand{by}\isamarkupfalse%
\ auto\isanewline
\ \ \isacommand{then}\isamarkupfalse%
\ \isacommand{have}\isamarkupfalse%
\ ii{\isacharcolon}{\kern0pt}\ {\isachardoublequoteopen}nabla\ {\isasymturnstile}\ swap\ pi\ t{\isadigit{1}}{\isacharprime}{\kern0pt}\ {\isasymapprox}\ swap\ {\isacharbrackleft}{\kern0pt}{\isacharparenleft}{\kern0pt}swapas\ pi\ a{\isacharcomma}{\kern0pt}\ swapas\ pi\ b{\isacharparenright}{\kern0pt}{\isacharbrackright}{\kern0pt}\ {\isacharparenleft}{\kern0pt}swap\ pi\ t{\isadigit{2}}{\isacharprime}{\kern0pt}{\isacharparenright}{\kern0pt}{\isachardoublequoteclose}\isanewline
\ \ \ \ \isacommand{using}\isamarkupfalse%
\ ds{\isacharunderscore}{\kern0pt}empty{\isacharunderscore}{\kern0pt}equiv{\isacharunderscore}{\kern0pt}{\isadigit{2}}{\isacharbrackleft}{\kern0pt}OF\ ds{\isacharunderscore}{\kern0pt}comm{\isacharbrackleft}{\kern0pt}of\ pi\ a\ b{\isacharbrackright}{\kern0pt}{\isacharbrackright}{\kern0pt}\ \isanewline
\ \ \ \ \ \ swap{\isacharunderscore}{\kern0pt}append{\isacharbrackleft}{\kern0pt}of\ {\isacartoucheopen}{\isacharbrackleft}{\kern0pt}{\isacharparenleft}{\kern0pt}swapas\ pi\ a{\isacharcomma}{\kern0pt}\ swapas\ pi\ b{\isacharparenright}{\kern0pt}{\isacharbrackright}{\kern0pt}{\isacartoucheclose}\ pi\ t{\isadigit{2}}{\isacharprime}{\kern0pt}{\isacharbrackright}{\kern0pt}\ \isacommand{by}\isamarkupfalse%
\ simp\isanewline
\isanewline
\ \ \isacommand{have}\isamarkupfalse%
\ iii{\isacharcolon}{\kern0pt}\ {\isachardoublequoteopen}swapas\ pi\ a\ {\isasymnoteq}\ swapas\ pi\ b{\isachardoublequoteclose}\ \isanewline
\ \ \ \ \isacommand{using}\isamarkupfalse%
\ equ{\isacharunderscore}{\kern0pt}abst{\isacharunderscore}{\kern0pt}ab{\isachardot}{\kern0pt}hyps{\isacharparenleft}{\kern0pt}{\isadigit{1}}{\isacharparenright}{\kern0pt}\ swapas{\isacharunderscore}{\kern0pt}rev{\isacharunderscore}{\kern0pt}pi{\isacharunderscore}{\kern0pt}a\ \isacommand{by}\isamarkupfalse%
\ blast\isanewline
\ \ \isanewline
\ \ \isacommand{with}\isamarkupfalse%
\ i\ ii\isanewline
\ \ \isacommand{have}\isamarkupfalse%
\ {\isachardoublequoteopen}nabla\ {\isasymturnstile}\ Abst\ {\isacharparenleft}{\kern0pt}swapas\ pi\ a{\isacharparenright}{\kern0pt}\ {\isacharparenleft}{\kern0pt}swap\ pi\ t{\isadigit{1}}{\isacharprime}{\kern0pt}{\isacharparenright}{\kern0pt}\ {\isasymapprox}\ Abst\ {\isacharparenleft}{\kern0pt}swapas\ pi\ b{\isacharparenright}{\kern0pt}\ {\isacharparenleft}{\kern0pt}swap\ pi\ t{\isadigit{2}}{\isacharprime}{\kern0pt}{\isacharparenright}{\kern0pt}{\isachardoublequoteclose}\isanewline
\ \ \ \ \isacommand{using}\isamarkupfalse%
\ equ{\isachardot}{\kern0pt}equ{\isacharunderscore}{\kern0pt}abst{\isacharunderscore}{\kern0pt}ab\ \isacommand{by}\isamarkupfalse%
\ auto\isanewline
\ \ \isacommand{hence}\isamarkupfalse%
\ {\isachardoublequoteopen}nabla\ {\isasymturnstile}\ swap\ pi\ {\isacharparenleft}{\kern0pt}Abst\ a\ t{\isadigit{1}}{\isacharprime}{\kern0pt}{\isacharparenright}{\kern0pt}\ {\isasymapprox}\ swap\ pi\ {\isacharparenleft}{\kern0pt}Abst\ b\ t{\isadigit{2}}{\isacharprime}{\kern0pt}{\isacharparenright}{\kern0pt}{\isachardoublequoteclose}\isanewline
\ \ \ \ \isacommand{using}\isamarkupfalse%
\ swap{\isachardot}{\kern0pt}simps{\isacharparenleft}{\kern0pt}{\isadigit{4}}{\isacharparenright}{\kern0pt}\ \isacommand{by}\isamarkupfalse%
\ simp\isanewline
\ \ \isacommand{then}\isamarkupfalse%
\ \isacommand{show}\isamarkupfalse%
\ {\isacharquery}{\kern0pt}case\ \isacommand{by}\isamarkupfalse%
\ simp\isanewline
\isacommand{next}\isamarkupfalse%
\isanewline
\ \ \isacommand{case}\isamarkupfalse%
\ {\isacharparenleft}{\kern0pt}equ{\isacharunderscore}{\kern0pt}abst{\isacharunderscore}{\kern0pt}aa\ nabla\ t{\isadigit{1}}{\isacharprime}{\kern0pt}\ t{\isadigit{2}}{\isacharprime}{\kern0pt}\ a{\isacharparenright}{\kern0pt}\isanewline
\ \ \isacommand{then}\isamarkupfalse%
\ \isacommand{show}\isamarkupfalse%
\ {\isacharquery}{\kern0pt}case\isanewline
\ \ \ \ \isacommand{using}\isamarkupfalse%
\ equ{\isachardot}{\kern0pt}equ{\isacharunderscore}{\kern0pt}abst{\isacharunderscore}{\kern0pt}aa\ \isacommand{by}\isamarkupfalse%
\ simp\isanewline
\isacommand{next}\isamarkupfalse%
\isanewline
\ \ \isacommand{case}\isamarkupfalse%
\ {\isacharparenleft}{\kern0pt}equ{\isacharunderscore}{\kern0pt}susp\ pi{\isadigit{1}}\ pi{\isadigit{2}}\ X\ nabla{\isacharparenright}{\kern0pt}\isanewline
\ \ \isacommand{then}\isamarkupfalse%
\ \isacommand{show}\isamarkupfalse%
\ {\isacharquery}{\kern0pt}case\ \isacommand{using}\isamarkupfalse%
\ ds{\isacharunderscore}{\kern0pt}cancel{\isacharunderscore}{\kern0pt}pi{\isacharunderscore}{\kern0pt}front\isanewline
\ \ \ \ \isacommand{by}\isamarkupfalse%
\ auto\isanewline
\isacommand{qed}\isamarkupfalse%
\ {\isacharparenleft}{\kern0pt}auto{\isacharparenright}{\kern0pt}%
\endisatagproof
{\isafoldproof}%
%
\isadelimproof
\isanewline
%
\endisadelimproof
\isanewline
\isacommand{lemma}\isamarkupfalse%
\ swap{\isacharunderscore}{\kern0pt}inv{\isacharunderscore}{\kern0pt}side{\isacharcolon}{\kern0pt}\ \isanewline
\ \ \isakeyword{shows}\ {\isachardoublequoteopen}nabla\ {\isasymturnstile}\ swap\ pi\ t{\isadigit{1}}\ {\isasymapprox}\ t{\isadigit{2}}\ {\isacharequal}{\kern0pt}\ nabla\ {\isasymturnstile}\ t{\isadigit{1}}\ {\isasymapprox}\ swap\ {\isacharparenleft}{\kern0pt}rev\ pi{\isacharparenright}{\kern0pt}\ t{\isadigit{2}}{\isachardoublequoteclose}\isanewline
%
\isadelimproof
%
\endisadelimproof
%
\isatagproof
\isacommand{proof}\isamarkupfalse%
\isanewline
\ \ \isanewline
\ \ \isacommand{{\isacharbraceleft}{\kern0pt}}\isamarkupfalse%
\isacommand{assume}\isamarkupfalse%
\ assm{\isadigit{1}}{\isacharcolon}{\kern0pt}\ {\isachardoublequoteopen}nabla\ {\isasymturnstile}\ swap\ pi\ t{\isadigit{1}}\ {\isasymapprox}\ t{\isadigit{2}}{\isachardoublequoteclose}\isanewline
\ \ \isacommand{from}\isamarkupfalse%
\ equ{\isacharunderscore}{\kern0pt}equivariance{\isacharbrackleft}{\kern0pt}OF\ assm{\isadigit{1}}{\isacharcomma}{\kern0pt}\ of\ {\isacartoucheopen}rev\ pi{\isacartoucheclose}{\isacharbrackright}{\kern0pt}\isanewline
\ \ \isacommand{have}\isamarkupfalse%
\ {\isachardoublequoteopen}nabla\ {\isasymturnstile}\ swap\ {\isacharparenleft}{\kern0pt}rev\ pi\ {\isacharat}{\kern0pt}\ pi{\isacharparenright}{\kern0pt}\ t{\isadigit{1}}\ \ {\isasymapprox}\ swap\ {\isacharparenleft}{\kern0pt}rev\ pi{\isacharparenright}{\kern0pt}\ t{\isadigit{2}}{\isachardoublequoteclose}\isanewline
\ \ \ \ \isacommand{using}\isamarkupfalse%
\ swap{\isacharunderscore}{\kern0pt}append{\isacharbrackleft}{\kern0pt}of\ {\isacartoucheopen}rev\ pi{\isacartoucheclose}\ pi\ t{\isadigit{1}}{\isacharbrackright}{\kern0pt}\ \isacommand{by}\isamarkupfalse%
\ simp\isanewline
\ \ \isacommand{then}\isamarkupfalse%
\ \isacommand{show}\isamarkupfalse%
\ {\isachardoublequoteopen}nabla\ {\isasymturnstile}\ t{\isadigit{1}}\ {\isasymapprox}\ swap\ {\isacharparenleft}{\kern0pt}rev\ pi{\isacharparenright}{\kern0pt}\ t{\isadigit{2}}{\isachardoublequoteclose}\isanewline
\ \ \ \ \isacommand{using}\isamarkupfalse%
\ ds{\isacharunderscore}{\kern0pt}empty{\isacharunderscore}{\kern0pt}equiv{\isacharunderscore}{\kern0pt}{\isadigit{1}}{\isacharbrackleft}{\kern0pt}OF\ ds{\isacharunderscore}{\kern0pt}rev{\isacharunderscore}{\kern0pt}pi{\isacharunderscore}{\kern0pt}id{\isacharbrackleft}{\kern0pt}of\ pi{\isacharbrackright}{\kern0pt}{\isacharcomma}{\kern0pt}\ of\ nabla\ t{\isadigit{1}}\ {\isacartoucheopen}swap\ {\isacharparenleft}{\kern0pt}rev\ pi{\isacharparenright}{\kern0pt}\ t{\isadigit{2}}{\isacartoucheclose}{\isacharbrackright}{\kern0pt}\ \isanewline
\ \ \ \ \ \ swap{\isacharunderscore}{\kern0pt}id{\isacharbrackleft}{\kern0pt}of\ t{\isadigit{1}}{\isacharbrackright}{\kern0pt}\ \isacommand{by}\isamarkupfalse%
\ simp\isacommand{{\isacharbraceright}{\kern0pt}}\isamarkupfalse%
\isanewline
\isanewline
\ \ \isacommand{{\isacharbraceleft}{\kern0pt}}\isamarkupfalse%
\isacommand{assume}\isamarkupfalse%
\ assm{\isadigit{2}}{\isacharcolon}{\kern0pt}\ {\isachardoublequoteopen}nabla\ {\isasymturnstile}\ t{\isadigit{1}}\ {\isasymapprox}\ swap\ {\isacharparenleft}{\kern0pt}rev\ pi{\isacharparenright}{\kern0pt}\ t{\isadigit{2}}{\isachardoublequoteclose}\isanewline
\ \ \isacommand{from}\isamarkupfalse%
\ equ{\isacharunderscore}{\kern0pt}equivariance{\isacharbrackleft}{\kern0pt}OF\ assm{\isadigit{2}}{\isacharcomma}{\kern0pt}\ of\ pi{\isacharbrackright}{\kern0pt}\isanewline
\ \ \isacommand{have}\isamarkupfalse%
\ {\isachardoublequoteopen}nabla\ {\isasymturnstile}\ swap\ pi\ t{\isadigit{1}}\ {\isasymapprox}\ swap\ {\isacharparenleft}{\kern0pt}pi\ {\isacharat}{\kern0pt}\ rev\ pi{\isacharparenright}{\kern0pt}\ t{\isadigit{2}}{\isachardoublequoteclose}\isanewline
\ \ \ \ \isacommand{using}\isamarkupfalse%
\ swap{\isacharunderscore}{\kern0pt}append{\isacharbrackleft}{\kern0pt}of\ pi\ {\isacartoucheopen}rev\ pi{\isacartoucheclose}\ t{\isadigit{2}}{\isacharbrackright}{\kern0pt}\ \isacommand{by}\isamarkupfalse%
\ simp\isanewline
\ \ \isacommand{then}\isamarkupfalse%
\ \isacommand{show}\isamarkupfalse%
\ {\isachardoublequoteopen}nabla\ {\isasymturnstile}\ swap\ pi\ t{\isadigit{1}}\ {\isasymapprox}\ t{\isadigit{2}}{\isachardoublequoteclose}\isanewline
\ \ \ \ \isacommand{using}\isamarkupfalse%
\ ds{\isacharunderscore}{\kern0pt}empty{\isacharunderscore}{\kern0pt}equiv{\isacharunderscore}{\kern0pt}{\isadigit{2}}{\isacharbrackleft}{\kern0pt}OF\ ds{\isacharunderscore}{\kern0pt}pi{\isacharunderscore}{\kern0pt}rev{\isacharunderscore}{\kern0pt}id{\isacharbrackleft}{\kern0pt}of\ pi{\isacharbrackright}{\kern0pt}{\isacharcomma}{\kern0pt}\ of\ nabla\ {\isacartoucheopen}swap\ pi\ t{\isadigit{1}}{\isacartoucheclose}\ t{\isadigit{2}}{\isacharbrackright}{\kern0pt}\ \isanewline
\ \ \ \ \ \ swap{\isacharunderscore}{\kern0pt}id{\isacharbrackleft}{\kern0pt}of\ t{\isadigit{2}}{\isacharbrackright}{\kern0pt}\ \isacommand{by}\isamarkupfalse%
\ simp\isacommand{{\isacharbraceright}{\kern0pt}}\isamarkupfalse%
\isanewline
\isacommand{qed}\isamarkupfalse%
%
\endisatagproof
{\isafoldproof}%
%
\isadelimproof
\isanewline
%
\endisadelimproof
\isanewline
\isacommand{lemma}\isamarkupfalse%
\ equ{\isacharunderscore}{\kern0pt}swap{\isacharunderscore}{\kern0pt}abba{\isacharcolon}{\kern0pt}\isanewline
\ \ \isakeyword{assumes}\ {\isachardoublequoteopen}n\ {\isacharequal}{\kern0pt}\ depth\ t{\isadigit{1}}{\isachardoublequoteclose}\isanewline
\ \ \isakeyword{shows}\ {\isachardoublequoteopen}nabla\ {\isasymturnstile}\ swap\ {\isacharbrackleft}{\kern0pt}{\isacharparenleft}{\kern0pt}a{\isacharcomma}{\kern0pt}b{\isacharparenright}{\kern0pt}{\isacharbrackright}{\kern0pt}\ t{\isadigit{1}}\ {\isasymapprox}\ t{\isadigit{2}}\ {\isasymLongrightarrow}\ nabla\ {\isasymturnstile}\ t{\isadigit{1}}\ {\isasymapprox}\ swap\ {\isacharbrackleft}{\kern0pt}{\isacharparenleft}{\kern0pt}b{\isacharcomma}{\kern0pt}a{\isacharparenright}{\kern0pt}{\isacharbrackright}{\kern0pt}\ t{\isadigit{2}}{\isachardoublequoteclose}\isanewline
%
\isadelimproof
%
\endisadelimproof
%
\isatagproof
\isacommand{proof}\isamarkupfalse%
{\isacharminus}{\kern0pt}\isanewline
\ \ \isacommand{assume}\isamarkupfalse%
\ {\isachardoublequoteopen}nabla\ {\isasymturnstile}\ swap\ {\isacharbrackleft}{\kern0pt}{\isacharparenleft}{\kern0pt}a{\isacharcomma}{\kern0pt}b{\isacharparenright}{\kern0pt}{\isacharbrackright}{\kern0pt}\ t{\isadigit{1}}\ {\isasymapprox}\ t{\isadigit{2}}{\isachardoublequoteclose}\isanewline
\ \ \isacommand{with}\isamarkupfalse%
\ ds{\isacharunderscore}{\kern0pt}baab{\isacharbrackleft}{\kern0pt}of\ b\ a{\isacharbrackright}{\kern0pt}\ \isanewline
\ \ \isacommand{have}\isamarkupfalse%
\ {\isachardoublequoteopen}nabla\ {\isasymturnstile}\ swap\ {\isacharbrackleft}{\kern0pt}{\isacharparenleft}{\kern0pt}b{\isacharcomma}{\kern0pt}a{\isacharparenright}{\kern0pt}{\isacharbrackright}{\kern0pt}\ t{\isadigit{1}}\ {\isasymapprox}\ t{\isadigit{2}}{\isachardoublequoteclose}\isanewline
\ \ \ \ \isacommand{using}\isamarkupfalse%
\ ds{\isacharunderscore}{\kern0pt}empty{\isacharunderscore}{\kern0pt}equiv{\isacharunderscore}{\kern0pt}{\isadigit{1}}\ \isacommand{by}\isamarkupfalse%
\ blast\isanewline
\ \ \ \ \isacommand{then}\isamarkupfalse%
\ \isacommand{show}\isamarkupfalse%
\ {\isacharquery}{\kern0pt}thesis\isanewline
\ \ \ \ \isacommand{using}\isamarkupfalse%
\ swap{\isacharunderscore}{\kern0pt}inv{\isacharunderscore}{\kern0pt}side{\isacharbrackleft}{\kern0pt}of\ nabla\ {\isacartoucheopen}{\isacharbrackleft}{\kern0pt}{\isacharparenleft}{\kern0pt}b{\isacharcomma}{\kern0pt}a{\isacharparenright}{\kern0pt}{\isacharbrackright}{\kern0pt}{\isacartoucheclose}{\isacharbrackright}{\kern0pt}\ \isacommand{by}\isamarkupfalse%
\ simp\isanewline
\isacommand{qed}\isamarkupfalse%
%
\endisatagproof
{\isafoldproof}%
%
\isadelimproof
\isanewline
%
\endisadelimproof
\isanewline
\isacommand{lemma}\isamarkupfalse%
\ equ{\isacharunderscore}{\kern0pt}equiv{\isacharunderscore}{\kern0pt}pi{\isacharcolon}{\kern0pt}\ \isanewline
\ \ \isakeyword{assumes}\ {\isachardoublequoteopen}{\isasymforall}a\ {\isasymin}\ ds\ pi{\isadigit{1}}\ pi{\isadigit{2}}{\isachardot}{\kern0pt}\ nabla\ {\isasymturnstile}\ a\ {\isasymsharp}\ t{\isachardoublequoteclose}\isanewline
\ \ \isakeyword{shows}\ {\isachardoublequoteopen}nabla\ {\isasymturnstile}\ swap\ pi{\isadigit{1}}\ t\ {\isasymapprox}\ swap\ pi{\isadigit{2}}\ t{\isachardoublequoteclose}\isanewline
%
\isadelimproof
\ \ %
\endisadelimproof
%
\isatagproof
\isacommand{using}\isamarkupfalse%
\ assms\ equ{\isacharunderscore}{\kern0pt}pi{\isacharunderscore}{\kern0pt}right{\isacharbrackleft}{\kern0pt}of\ {\isacartoucheopen}rev\ pi{\isadigit{1}}\ {\isacharat}{\kern0pt}\ pi{\isadigit{2}}{\isacartoucheclose}\ nabla\ t{\isacharbrackright}{\kern0pt}\isanewline
\ \ \ \ swap{\isacharunderscore}{\kern0pt}inv{\isacharunderscore}{\kern0pt}side{\isacharbrackleft}{\kern0pt}of\ nabla\ pi{\isadigit{1}}\ t\ {\isacartoucheopen}swap\ pi{\isadigit{2}}\ t{\isacartoucheclose}{\isacharbrackright}{\kern0pt}\ ds{\isacharunderscore}{\kern0pt}rev\ swap{\isacharunderscore}{\kern0pt}append\ \isacommand{by}\isamarkupfalse%
\ simp%
\endisatagproof
{\isafoldproof}%
%
\isadelimproof
\isanewline
%
\endisadelimproof
\isanewline
\isanewline
\isacommand{lemma}\isamarkupfalse%
\ equ{\isacharunderscore}{\kern0pt}symm{\isacharcolon}{\kern0pt}\isanewline
\ \ \isakeyword{shows}\ {\isachardoublequoteopen}{\isacharparenleft}{\kern0pt}nabla\ {\isasymturnstile}\ t{\isadigit{1}}\ {\isasymapprox}\ t{\isadigit{2}}{\isacharparenright}{\kern0pt}\ {\isasymLongrightarrow}\ {\isacharparenleft}{\kern0pt}nabla\ {\isasymturnstile}\ t{\isadigit{2}}\ {\isasymapprox}\ t{\isadigit{1}}{\isacharparenright}{\kern0pt}{\isachardoublequoteclose}\isanewline
%
\isadelimproof
%
\endisadelimproof
%
\isatagproof
\isacommand{proof}\isamarkupfalse%
{\isacharparenleft}{\kern0pt}induction\ rule{\isacharcolon}{\kern0pt}\ equ{\isachardot}{\kern0pt}induct{\isacharparenright}{\kern0pt}\isanewline
\ \ \isacommand{case}\isamarkupfalse%
\ {\isacharparenleft}{\kern0pt}equ{\isacharunderscore}{\kern0pt}abst{\isacharunderscore}{\kern0pt}ab\ a\ b\ nabla\ t{\isadigit{2}}{\isacharprime}{\kern0pt}\ t{\isadigit{1}}{\isacharprime}{\kern0pt}{\isacharparenright}{\kern0pt}\isanewline
\ \ \isacommand{have}\isamarkupfalse%
\ i{\isacharcolon}{\kern0pt}\ {\isachardoublequoteopen}nabla\ {\isasymturnstile}\ b\ {\isasymsharp}\ swap\ {\isacharbrackleft}{\kern0pt}{\isacharparenleft}{\kern0pt}a{\isacharcomma}{\kern0pt}\ b{\isacharparenright}{\kern0pt}{\isacharbrackright}{\kern0pt}\ t{\isadigit{2}}{\isacharprime}{\kern0pt}{\isachardoublequoteclose}\ \isanewline
\ \ \ \ \isacommand{using}\isamarkupfalse%
\ fresh{\isacharunderscore}{\kern0pt}swap{\isacharunderscore}{\kern0pt}eqvt{\isacharbrackleft}{\kern0pt}of\ nabla\ a\ t{\isadigit{2}}{\isacharprime}{\kern0pt}\ {\isachardoublequoteopen}{\isacharbrackleft}{\kern0pt}{\isacharparenleft}{\kern0pt}a{\isacharcomma}{\kern0pt}b{\isacharparenright}{\kern0pt}{\isacharbrackright}{\kern0pt}{\isachardoublequoteclose}{\isacharbrackright}{\kern0pt}\ equ{\isacharunderscore}{\kern0pt}abst{\isacharunderscore}{\kern0pt}ab{\isacharparenleft}{\kern0pt}{\isadigit{2}}{\isacharparenright}{\kern0pt}\ \isacommand{by}\isamarkupfalse%
\ auto\isanewline
\ \ \isacommand{with}\isamarkupfalse%
\ equ{\isacharunderscore}{\kern0pt}abst{\isacharunderscore}{\kern0pt}ab{\isacharparenleft}{\kern0pt}{\isadigit{4}}{\isacharparenright}{\kern0pt}\ \isanewline
\ \ \isacommand{have}\isamarkupfalse%
\ b{\isacharunderscore}{\kern0pt}fresh{\isacharcolon}{\kern0pt}\ {\isachardoublequoteopen}nabla\ {\isasymturnstile}\ b\ {\isasymsharp}\ t{\isadigit{1}}{\isacharprime}{\kern0pt}{\isachardoublequoteclose}\isanewline
\ \ \ \ \isacommand{using}\isamarkupfalse%
\ l{\isadigit{3}}{\isacharunderscore}{\kern0pt}jud\ \isacommand{by}\isamarkupfalse%
\ blast\isanewline
\ \ \isacommand{from}\isamarkupfalse%
\ equ{\isacharunderscore}{\kern0pt}abst{\isacharunderscore}{\kern0pt}ab{\isacharparenleft}{\kern0pt}{\isadigit{4}}{\isacharparenright}{\kern0pt}\ \isanewline
\ \ \isacommand{have}\isamarkupfalse%
\ ii{\isacharcolon}{\kern0pt}\ {\isachardoublequoteopen}nabla\ {\isasymturnstile}\ t{\isadigit{2}}{\isacharprime}{\kern0pt}\ {\isasymapprox}\ swap\ {\isacharbrackleft}{\kern0pt}{\isacharparenleft}{\kern0pt}b{\isacharcomma}{\kern0pt}\ a{\isacharparenright}{\kern0pt}{\isacharbrackright}{\kern0pt}\ t{\isadigit{1}}{\isacharprime}{\kern0pt}{\isachardoublequoteclose}\isanewline
\ \ \ \ \isacommand{using}\isamarkupfalse%
\ equ{\isacharunderscore}{\kern0pt}swap{\isacharunderscore}{\kern0pt}abba\ \isacommand{by}\isamarkupfalse%
\ simp\isanewline
\ \ \isacommand{show}\isamarkupfalse%
\ {\isacharquery}{\kern0pt}case\ \isanewline
\ \ \ \ \isacommand{using}\isamarkupfalse%
\ equ{\isachardot}{\kern0pt}equ{\isacharunderscore}{\kern0pt}abst{\isacharunderscore}{\kern0pt}ab{\isacharbrackleft}{\kern0pt}OF\ equ{\isacharunderscore}{\kern0pt}abst{\isacharunderscore}{\kern0pt}ab{\isacharparenleft}{\kern0pt}{\isadigit{1}}{\isacharparenright}{\kern0pt}{\isacharbrackleft}{\kern0pt}symmetric{\isacharbrackright}{\kern0pt}\ b{\isacharunderscore}{\kern0pt}fresh\ ii{\isacharbrackright}{\kern0pt}\ \isacommand{by}\isamarkupfalse%
\ simp\isanewline
\isacommand{next}\isamarkupfalse%
\isanewline
\ \ \isacommand{case}\isamarkupfalse%
\ {\isacharparenleft}{\kern0pt}equ{\isacharunderscore}{\kern0pt}abst{\isacharunderscore}{\kern0pt}aa\ nabla\ t{\isadigit{1}}{\isacharprime}{\kern0pt}\ t{\isadigit{2}}{\isacharprime}{\kern0pt}\ a{\isacharparenright}{\kern0pt}\isanewline
\ \ \isacommand{then}\isamarkupfalse%
\ \isacommand{show}\isamarkupfalse%
\ {\isacharquery}{\kern0pt}case\ \isanewline
\ \ \ \ \isacommand{using}\isamarkupfalse%
\ equ{\isachardot}{\kern0pt}equ{\isacharunderscore}{\kern0pt}abst{\isacharunderscore}{\kern0pt}aa{\isacharbrackleft}{\kern0pt}OF\ equ{\isacharunderscore}{\kern0pt}abst{\isacharunderscore}{\kern0pt}aa{\isacharparenleft}{\kern0pt}{\isadigit{2}}{\isacharparenright}{\kern0pt}{\isacharcomma}{\kern0pt}\ of\ a{\isacharbrackright}{\kern0pt}\ \isacommand{by}\isamarkupfalse%
\ simp\isanewline
\isacommand{next}\isamarkupfalse%
\isanewline
\ \ \isacommand{case}\isamarkupfalse%
\ {\isacharparenleft}{\kern0pt}equ{\isacharunderscore}{\kern0pt}unit\ nabla{\isacharparenright}{\kern0pt}\isanewline
\ \ \isacommand{then}\isamarkupfalse%
\ \isacommand{show}\isamarkupfalse%
\ {\isacharquery}{\kern0pt}case\ \isacommand{by}\isamarkupfalse%
\ auto\isanewline
\isacommand{next}\isamarkupfalse%
\isanewline
\ \ \isacommand{case}\isamarkupfalse%
\ {\isacharparenleft}{\kern0pt}equ{\isacharunderscore}{\kern0pt}atom\ a\ b\ nabla{\isacharparenright}{\kern0pt}\isanewline
\ \ \isacommand{then}\isamarkupfalse%
\ \isacommand{show}\isamarkupfalse%
\ {\isacharquery}{\kern0pt}case\ \isacommand{by}\isamarkupfalse%
\ auto\isanewline
\isacommand{next}\isamarkupfalse%
\isanewline
\ \ \isacommand{case}\isamarkupfalse%
\ {\isacharparenleft}{\kern0pt}equ{\isacharunderscore}{\kern0pt}susp\ pi{\isadigit{1}}\ pi{\isadigit{2}}\ X\ nabla{\isacharparenright}{\kern0pt}\isanewline
\ \ \isacommand{then}\isamarkupfalse%
\ \isacommand{show}\isamarkupfalse%
\ {\isacharquery}{\kern0pt}case\ \isanewline
\ \ \ \ \isacommand{using}\isamarkupfalse%
\ ds{\isacharunderscore}{\kern0pt}sym\ \isacommand{by}\isamarkupfalse%
\ auto\isanewline
\isacommand{next}\isamarkupfalse%
\isanewline
\ \ \isacommand{case}\isamarkupfalse%
\ {\isacharparenleft}{\kern0pt}equ{\isacharunderscore}{\kern0pt}paar\ nabla\ t{\isadigit{1}}{\isacharprime}{\kern0pt}\ t{\isadigit{2}}{\isacharprime}{\kern0pt}\ s{\isadigit{1}}{\isacharprime}{\kern0pt}\ s{\isadigit{2}}{\isacharprime}{\kern0pt}{\isacharparenright}{\kern0pt}\isanewline
\ \ \isacommand{then}\isamarkupfalse%
\ \isacommand{show}\isamarkupfalse%
\ {\isacharquery}{\kern0pt}case\ \isacommand{by}\isamarkupfalse%
\ auto\isanewline
\isacommand{next}\isamarkupfalse%
\isanewline
\ \ \isacommand{case}\isamarkupfalse%
\ {\isacharparenleft}{\kern0pt}equ{\isacharunderscore}{\kern0pt}func\ nabla\ t{\isadigit{1}}{\isacharprime}{\kern0pt}\ t{\isadigit{2}}{\isacharprime}{\kern0pt}\ f{\isacharparenright}{\kern0pt}\isanewline
\ \ \isacommand{then}\isamarkupfalse%
\ \isacommand{show}\isamarkupfalse%
\ {\isacharquery}{\kern0pt}case\ \isanewline
\ \ \ \ \isacommand{using}\isamarkupfalse%
\ equ{\isachardot}{\kern0pt}equ{\isacharunderscore}{\kern0pt}func{\isacharbrackleft}{\kern0pt}OF\ equ{\isacharunderscore}{\kern0pt}func{\isacharparenleft}{\kern0pt}{\isadigit{2}}{\isacharparenright}{\kern0pt}{\isacharcomma}{\kern0pt}\ of\ f{\isacharbrackright}{\kern0pt}\ \isacommand{by}\isamarkupfalse%
\ simp\isanewline
\isacommand{qed}\isamarkupfalse%
%
\endisatagproof
{\isafoldproof}%
%
\isadelimproof
\isanewline
%
\endisadelimproof
\isanewline
\isanewline
\isacommand{lemma}\isamarkupfalse%
\ equ{\isacharunderscore}{\kern0pt}trans{\isacharcolon}{\kern0pt}\isanewline
\ \ \isakeyword{assumes}\ {\isachardoublequoteopen}nabla\ {\isasymturnstile}\ t{\isadigit{1}}\ {\isasymapprox}\ t{\isadigit{2}}{\isachardoublequoteclose}\ {\isachardoublequoteopen}nabla\ {\isasymturnstile}\ t{\isadigit{2}}\ {\isasymapprox}\ t{\isadigit{3}}{\isachardoublequoteclose}\isanewline
\ \ \isakeyword{shows}\ {\isachardoublequoteopen}nabla\ {\isasymturnstile}\ t{\isadigit{1}}\ {\isasymapprox}\ t{\isadigit{3}}{\isachardoublequoteclose}\isanewline
%
\isadelimproof
\ \ %
\endisadelimproof
%
\isatagproof
\isacommand{using}\isamarkupfalse%
\ assms\isanewline
\isacommand{proof}\isamarkupfalse%
{\isacharparenleft}{\kern0pt}induction\ {\isacartoucheopen}depth\ t{\isadigit{1}}{\isacartoucheclose}\ arbitrary{\isacharcolon}{\kern0pt}\ t{\isadigit{1}}\ t{\isadigit{2}}\ t{\isadigit{3}}\ rule{\isacharcolon}{\kern0pt}\ nat{\isacharunderscore}{\kern0pt}less{\isacharunderscore}{\kern0pt}induct{\isacharparenright}{\kern0pt}\isanewline
\ \ \isacommand{case}\isamarkupfalse%
\ {\isadigit{1}}\isanewline
\ \ \isacommand{note}\isamarkupfalse%
\ IH\ {\isacharequal}{\kern0pt}\ this\isanewline
\ \ \isacommand{have}\isamarkupfalse%
\ IH{\isacharunderscore}{\kern0pt}usable{\isacharcolon}{\kern0pt}\ \isanewline
\ \ \ \ {\isachardoublequoteopen}depth\ t{\isadigit{1}}{\isacharprime}{\kern0pt}\ {\isacharless}{\kern0pt}\ depth\ t{\isadigit{1}}\ {\isasymLongrightarrow}\ nabla\ {\isasymturnstile}\ t{\isadigit{1}}{\isacharprime}{\kern0pt}\ {\isasymapprox}\ t{\isadigit{2}}{\isacharprime}{\kern0pt}\ {\isasymLongrightarrow}\ \ nabla\ {\isasymturnstile}\ t{\isadigit{2}}{\isacharprime}{\kern0pt}\ {\isasymapprox}\ t{\isadigit{3}}{\isacharprime}{\kern0pt}\ {\isasymLongrightarrow}\ nabla\ {\isasymturnstile}\ t{\isadigit{1}}{\isacharprime}{\kern0pt}\ {\isasymapprox}\ t{\isadigit{3}}{\isacharprime}{\kern0pt}{\isachardoublequoteclose}\isanewline
\ \ \ \ \isakeyword{for}\ t{\isadigit{1}}{\isacharprime}{\kern0pt}\ t{\isadigit{2}}{\isacharprime}{\kern0pt}\ t{\isadigit{3}}{\isacharprime}{\kern0pt}\ \isacommand{using}\isamarkupfalse%
\ IH\ \isacommand{by}\isamarkupfalse%
\ blast\isanewline
\ \ \isacommand{have}\isamarkupfalse%
\ {\isachardoublequoteopen}nabla\ {\isasymturnstile}\ t{\isadigit{1}}\ {\isasymapprox}\ t{\isadigit{2}}\ {\isasymLongrightarrow}\ nabla\ {\isasymturnstile}\ t{\isadigit{2}}\ {\isasymapprox}\ t{\isadigit{3}}\ {\isasymLongrightarrow}\ \ nabla\ {\isasymturnstile}\ t{\isadigit{1}}\ {\isasymapprox}\ t{\isadigit{3}}{\isachardoublequoteclose}\isanewline
\ \ \isacommand{proof}\isamarkupfalse%
{\isacharminus}{\kern0pt}\isanewline
\ \ \ \ \isacommand{assume}\isamarkupfalse%
\ t{\isadigit{1}}{\isadigit{2}}{\isacharcolon}{\kern0pt}\ {\isachardoublequoteopen}nabla\ {\isasymturnstile}\ t{\isadigit{1}}\ {\isasymapprox}\ t{\isadigit{2}}{\isachardoublequoteclose}\ \isakeyword{and}\ t{\isadigit{2}}{\isadigit{3}}{\isacharcolon}{\kern0pt}\ {\isachardoublequoteopen}nabla\ {\isasymturnstile}\ t{\isadigit{2}}\ {\isasymapprox}\ t{\isadigit{3}}{\isachardoublequoteclose}\isanewline
\ \ \ \ \isacommand{show}\isamarkupfalse%
\ {\isacharquery}{\kern0pt}thesis\isanewline
\ \ \ \ \ \isacommand{proof}\isamarkupfalse%
{\isacharparenleft}{\kern0pt}cases\ rule{\isacharcolon}{\kern0pt}\ equ{\isachardot}{\kern0pt}cases{\isacharbrackleft}{\kern0pt}OF\ {\isacartoucheopen}nabla\ {\isasymturnstile}\ t{\isadigit{1}}\ {\isasymapprox}\ t{\isadigit{2}}{\isacartoucheclose}{\isacharbrackright}{\kern0pt}{\isacharparenright}{\kern0pt}\isanewline
\ \ \ \ \ \ \ \isacommand{case}\isamarkupfalse%
\ {\isacharparenleft}{\kern0pt}{\isadigit{1}}\ a\ b\ nabla\ t{\isadigit{2}}{\isacharprime}{\kern0pt}\ t{\isadigit{1}}{\isacharprime}{\kern0pt}{\isacharparenright}{\kern0pt}\isanewline
\ \ \ \ \ \ \ \isacommand{from}\isamarkupfalse%
\ {\isadigit{1}}{\isacharparenleft}{\kern0pt}{\isadigit{2}}{\isacharparenright}{\kern0pt}\ \isacommand{have}\isamarkupfalse%
\ deptht{\isadigit{1}}{\isacharcolon}{\kern0pt}\ {\isachardoublequoteopen}depth\ t{\isadigit{1}}{\isacharprime}{\kern0pt}\ {\isacharless}{\kern0pt}\ depth\ t{\isadigit{1}}{\isachardoublequoteclose}\isanewline
\ \ \ \ \ \ \ \ \ \isacommand{using}\isamarkupfalse%
\ depth{\isachardot}{\kern0pt}simps{\isacharparenleft}{\kern0pt}{\isadigit{4}}{\isacharparenright}{\kern0pt}\ \isacommand{by}\isamarkupfalse%
\ auto\isanewline
\ \ \ \ \ \ \ \isacommand{from}\isamarkupfalse%
\ t{\isadigit{2}}{\isadigit{3}}\ \isacommand{show}\isamarkupfalse%
\ {\isacharquery}{\kern0pt}thesis\isanewline
\ \ \ \ \ \ \ \isacommand{proof}\isamarkupfalse%
{\isacharparenleft}{\kern0pt}cases\ rule{\isacharcolon}{\kern0pt}\ equ{\isachardot}{\kern0pt}cases{\isacharparenright}{\kern0pt}\isanewline
\ \ \ \ \ \ \ \ \ \isacommand{case}\isamarkupfalse%
\ {\isacharparenleft}{\kern0pt}equ{\isacharunderscore}{\kern0pt}abst{\isacharunderscore}{\kern0pt}ab\ b\ c\ t{\isadigit{3}}{\isacharprime}{\kern0pt}\ t{\isadigit{2}}{\isacharprime}{\kern0pt}{\isacharparenright}{\kern0pt}\isanewline
\ \ \ \ \ \ \ \ \ \isacommand{then}\isamarkupfalse%
\ \isacommand{show}\isamarkupfalse%
\ {\isacharquery}{\kern0pt}thesis\ \isanewline
\ \ \ \ \ \ \ \ \ \isacommand{proof}\isamarkupfalse%
{\isacharparenleft}{\kern0pt}cases\ {\isachardoublequoteopen}a\ {\isacharequal}{\kern0pt}\ c{\isachardoublequoteclose}{\isacharparenright}{\kern0pt}\isanewline
\ \ \ \ \ \ \ \ \ \ \ \isacommand{case}\isamarkupfalse%
\ True\isanewline
\ \ \ \ \ \ \ \ \ \ \ \isacommand{have}\isamarkupfalse%
\ i{\isacharcolon}{\kern0pt}\ {\isachardoublequoteopen}nabla\ {\isasymturnstile}\ swap\ {\isacharbrackleft}{\kern0pt}{\isacharparenleft}{\kern0pt}c{\isacharcomma}{\kern0pt}b{\isacharparenright}{\kern0pt}{\isacharbrackright}{\kern0pt}\ t{\isadigit{2}}{\isacharprime}{\kern0pt}\ {\isasymapprox}\ swap\ {\isacharparenleft}{\kern0pt}{\isacharbrackleft}{\kern0pt}{\isacharparenleft}{\kern0pt}c{\isacharcomma}{\kern0pt}b{\isacharparenright}{\kern0pt}{\isacharbrackright}{\kern0pt}{\isacharat}{\kern0pt}{\isacharbrackleft}{\kern0pt}{\isacharparenleft}{\kern0pt}b{\isacharcomma}{\kern0pt}c{\isacharparenright}{\kern0pt}{\isacharbrackright}{\kern0pt}{\isacharparenright}{\kern0pt}\ t{\isadigit{3}}{\isacharprime}{\kern0pt}{\isachardoublequoteclose}\isanewline
\ \ \ \ \ \ \ \ \ \ \ \ \ \isacommand{using}\isamarkupfalse%
\ equ{\isacharunderscore}{\kern0pt}equivariance{\isacharbrackleft}{\kern0pt}OF\ equ{\isacharunderscore}{\kern0pt}abst{\isacharunderscore}{\kern0pt}ab{\isacharparenleft}{\kern0pt}{\isadigit{5}}{\isacharparenright}{\kern0pt}{\isacharcomma}{\kern0pt}\ of\ {\isacartoucheopen}{\isacharbrackleft}{\kern0pt}{\isacharparenleft}{\kern0pt}c{\isacharcomma}{\kern0pt}b{\isacharparenright}{\kern0pt}{\isacharbrackright}{\kern0pt}{\isacartoucheclose}{\isacharbrackright}{\kern0pt}\ {\isadigit{1}}{\isacharparenleft}{\kern0pt}{\isadigit{1}}{\isacharparenright}{\kern0pt}\isanewline
\ \ \ \ \ \ \ \ \ \ \ \ \ \ \ swap{\isacharunderscore}{\kern0pt}append{\isacharbrackleft}{\kern0pt}of\ {\isacartoucheopen}{\isacharbrackleft}{\kern0pt}{\isacharparenleft}{\kern0pt}c{\isacharcomma}{\kern0pt}b{\isacharparenright}{\kern0pt}{\isacharbrackright}{\kern0pt}{\isacartoucheclose}\ {\isacartoucheopen}{\isacharbrackleft}{\kern0pt}{\isacharparenleft}{\kern0pt}b{\isacharcomma}{\kern0pt}c{\isacharparenright}{\kern0pt}{\isacharbrackright}{\kern0pt}{\isacartoucheclose}\ t{\isadigit{3}}{\isacharprime}{\kern0pt}{\isacharbrackright}{\kern0pt}\ \isacommand{by}\isamarkupfalse%
\ simp\isanewline
\ \ \ \ \ \ \ \ \ \ \ \isacommand{have}\isamarkupfalse%
\ ii{\isacharcolon}{\kern0pt}\ {\isachardoublequoteopen}nabla\ {\isasymturnstile}\ swap\ {\isacharbrackleft}{\kern0pt}{\isacharparenleft}{\kern0pt}c{\isacharcomma}{\kern0pt}b{\isacharparenright}{\kern0pt}{\isacharbrackright}{\kern0pt}\ t{\isadigit{2}}{\isacharprime}{\kern0pt}\ {\isasymapprox}\ t{\isadigit{3}}{\isacharprime}{\kern0pt}{\isachardoublequoteclose}\isanewline
\ \ \ \ \ \ \ \ \ \ \ \ \ \isacommand{using}\isamarkupfalse%
\ ds{\isacharunderscore}{\kern0pt}empty{\isacharunderscore}{\kern0pt}equiv{\isacharunderscore}{\kern0pt}{\isadigit{2}}{\isacharbrackleft}{\kern0pt}OF\ ds{\isacharunderscore}{\kern0pt}baab{\isacharunderscore}{\kern0pt}id{\isacharbrackleft}{\kern0pt}OF\ equ{\isacharunderscore}{\kern0pt}abst{\isacharunderscore}{\kern0pt}ab{\isacharparenleft}{\kern0pt}{\isadigit{3}}{\isacharparenright}{\kern0pt}{\isacharbrackright}{\kern0pt}\ i{\isacharbrackright}{\kern0pt}\ \isacommand{by}\isamarkupfalse%
\ simp\isanewline
\ \ \ \ \ \ \ \ \ \ \ \isacommand{with}\isamarkupfalse%
\ equ{\isacharunderscore}{\kern0pt}abst{\isacharunderscore}{\kern0pt}ab\ {\isadigit{1}}\ True\isanewline
\ \ \ \ \ \ \ \ \ \ \ \isacommand{have}\isamarkupfalse%
\ iii{\isacharcolon}{\kern0pt}\ {\isachardoublequoteopen}nabla\ {\isasymturnstile}\ t{\isadigit{1}}{\isacharprime}{\kern0pt}\ {\isasymapprox}\ swap\ {\isacharbrackleft}{\kern0pt}{\isacharparenleft}{\kern0pt}c{\isacharcomma}{\kern0pt}\ b{\isacharparenright}{\kern0pt}{\isacharbrackright}{\kern0pt}\ t{\isadigit{2}}{\isacharprime}{\kern0pt}{\isachardoublequoteclose}\ \isacommand{by}\isamarkupfalse%
\ blast\isanewline
\ \ \ \ \ \ \ \ \ \ \ \isacommand{with}\isamarkupfalse%
\ ii\ \isacommand{have}\isamarkupfalse%
\ {\isachardoublequoteopen}nabla\ {\isasymturnstile}\ t{\isadigit{1}}{\isacharprime}{\kern0pt}\ {\isasymapprox}\ t{\isadigit{3}}{\isacharprime}{\kern0pt}{\isachardoublequoteclose}\isanewline
\ \ \ \ \ \ \ \ \ \ \ \ \ \isacommand{using}\isamarkupfalse%
\ IH{\isacharunderscore}{\kern0pt}usable{\isacharbrackleft}{\kern0pt}OF\ deptht{\isadigit{1}}{\isacharcomma}{\kern0pt}\ of\ {\isacartoucheopen}swap\ {\isacharbrackleft}{\kern0pt}{\isacharparenleft}{\kern0pt}c{\isacharcomma}{\kern0pt}\ b{\isacharparenright}{\kern0pt}{\isacharbrackright}{\kern0pt}\ t{\isadigit{2}}{\isacharprime}{\kern0pt}{\isacartoucheclose}\ t{\isadigit{3}}{\isacharprime}{\kern0pt}{\isacharbrackright}{\kern0pt}\ {\isadigit{1}}{\isacharparenleft}{\kern0pt}{\isadigit{1}}{\isacharparenright}{\kern0pt}\isanewline
\ \ \ \ \ \ \ \ \ \ \ \ \ \isacommand{by}\isamarkupfalse%
\ simp\isanewline
\ \ \ \ \ \ \ \ \ \ \ \isacommand{then}\isamarkupfalse%
\ \isacommand{show}\isamarkupfalse%
\ {\isacharquery}{\kern0pt}thesis\ \isacommand{using}\isamarkupfalse%
\ True\ equ{\isacharunderscore}{\kern0pt}abst{\isacharunderscore}{\kern0pt}ab{\isacharparenleft}{\kern0pt}{\isadigit{2}}{\isacharparenright}{\kern0pt}\ {\isadigit{1}}{\isacharparenleft}{\kern0pt}{\isadigit{1}}{\isacharcomma}{\kern0pt}{\isadigit{2}}{\isacharparenright}{\kern0pt}\ \isacommand{by}\isamarkupfalse%
\ auto\isanewline
\ \ \ \ \ \ \ \ \ \isacommand{next}\isamarkupfalse%
\isanewline
\ \ \ \ \ \ \ \ \ \ \ \isacommand{case}\isamarkupfalse%
\ False\isanewline
\isanewline
\ \ \ \ \ \ \ \ \ \ \ \isacommand{have}\isamarkupfalse%
\ deptht{\isadigit{2}}{\isacharcolon}{\kern0pt}\ {\isachardoublequoteopen}depth\ {\isacharparenleft}{\kern0pt}swap\ {\isacharbrackleft}{\kern0pt}{\isacharparenleft}{\kern0pt}a{\isacharcomma}{\kern0pt}b{\isacharparenright}{\kern0pt}{\isacharbrackright}{\kern0pt}\ t{\isadigit{2}}{\isacharprime}{\kern0pt}{\isacharparenright}{\kern0pt}\ {\isacharless}{\kern0pt}\ depth\ t{\isadigit{1}}{\isachardoublequoteclose}\isanewline
\ \ \ \ \ \ \ \ \ \ \ \ \ \ \isacommand{using}\isamarkupfalse%
\ depth{\isachardot}{\kern0pt}simps{\isacharparenleft}{\kern0pt}{\isadigit{4}}{\isacharparenright}{\kern0pt}\ equ{\isacharunderscore}{\kern0pt}depth{\isacharbrackleft}{\kern0pt}OF\ {\isadigit{1}}{\isacharparenleft}{\kern0pt}{\isadigit{6}}{\isacharparenright}{\kern0pt}{\isacharbrackright}{\kern0pt}\ swap{\isacharunderscore}{\kern0pt}depth\ {\isadigit{1}}{\isacharparenleft}{\kern0pt}{\isadigit{3}}{\isacharparenright}{\kern0pt}\ deptht{\isadigit{1}}\isanewline
\ \ \ \ \ \ \ \ \ \ \ \ \ \ equ{\isacharunderscore}{\kern0pt}abst{\isacharunderscore}{\kern0pt}ab{\isacharparenleft}{\kern0pt}{\isadigit{1}}{\isacharparenright}{\kern0pt}\ \isacommand{by}\isamarkupfalse%
\ simp\isanewline
\isanewline
\ \ \ \ \ \ \ \ \ \ \ \ \isacommand{from}\isamarkupfalse%
\ equ{\isacharunderscore}{\kern0pt}abst{\isacharunderscore}{\kern0pt}ab{\isacharparenleft}{\kern0pt}{\isadigit{1}}{\isacharcomma}{\kern0pt}{\isadigit{4}}{\isacharcomma}{\kern0pt}{\isadigit{5}}{\isacharparenright}{\kern0pt}\ {\isadigit{1}}{\isacharparenleft}{\kern0pt}{\isadigit{1}}{\isacharcomma}{\kern0pt}{\isadigit{3}}{\isacharcomma}{\kern0pt}{\isadigit{5}}{\isacharparenright}{\kern0pt}\isanewline
\ \ \ \ \ \ \ \ \ \ \ \ \ \ \isacommand{have}\isamarkupfalse%
\ ii{\isacharcolon}{\kern0pt}\ {\isachardoublequoteopen}nabla\ {\isasymturnstile}\ a\ {\isasymsharp}\ swap\ {\isacharbrackleft}{\kern0pt}{\isacharparenleft}{\kern0pt}b{\isacharcomma}{\kern0pt}c{\isacharparenright}{\kern0pt}{\isacharbrackright}{\kern0pt}\ t{\isadigit{3}}{\isacharprime}{\kern0pt}{\isachardoublequoteclose}\isanewline
\ \ \ \ \ \ \ \ \ \ \ \ \ \ \ \ \isacommand{using}\isamarkupfalse%
\ l{\isadigit{3}}{\isacharunderscore}{\kern0pt}jud\ \isacommand{by}\isamarkupfalse%
\ simp\isanewline
\isanewline
\ \ \ \ \ \ \ \ \ \ \ \ \isacommand{then}\isamarkupfalse%
\ \isacommand{have}\isamarkupfalse%
\ a{\isacharunderscore}{\kern0pt}fresh{\isacharunderscore}{\kern0pt}t{\isadigit{3}}{\isacharcolon}{\kern0pt}\ {\isachardoublequoteopen}nabla\ {\isasymturnstile}\ a\ {\isasymsharp}\ t{\isadigit{3}}{\isacharprime}{\kern0pt}{\isachardoublequoteclose}\isanewline
\ \ \ \ \ \ \ \ \ \ \ \ \ \ \isacommand{using}\isamarkupfalse%
\ fresh{\isacharunderscore}{\kern0pt}swap{\isacharunderscore}{\kern0pt}left{\isacharbrackleft}{\kern0pt}OF\ ii{\isacharbrackright}{\kern0pt}\ False\ equ{\isacharunderscore}{\kern0pt}abst{\isacharunderscore}{\kern0pt}ab\ {\isadigit{1}}\ \isacommand{by}\isamarkupfalse%
\ simp\isanewline
\ \ \ \ \ \ \ \ \ \ \ \ \isacommand{have}\isamarkupfalse%
\ b{\isacharunderscore}{\kern0pt}fresh{\isacharunderscore}{\kern0pt}t{\isadigit{3}}{\isacharcolon}{\kern0pt}\ {\isachardoublequoteopen}nabla\ {\isasymturnstile}\ b\ {\isasymsharp}\ t{\isadigit{3}}{\isacharprime}{\kern0pt}{\isachardoublequoteclose}\ \isanewline
\ \ \ \ \ \ \ \ \ \ \ \ \ \ \isacommand{using}\isamarkupfalse%
\ equ{\isacharunderscore}{\kern0pt}abst{\isacharunderscore}{\kern0pt}ab{\isacharparenleft}{\kern0pt}{\isadigit{4}}{\isacharparenright}{\kern0pt}\ {\isadigit{1}}{\isacharparenleft}{\kern0pt}{\isadigit{1}}{\isacharparenright}{\kern0pt}\ \isacommand{by}\isamarkupfalse%
\ simp\isanewline
\isanewline
\ \ \ \ \ \ \ \ \ \ \ \ \isacommand{from}\isamarkupfalse%
\ {\isadigit{1}}{\isacharparenleft}{\kern0pt}{\isadigit{1}}{\isacharparenright}{\kern0pt}\ equ{\isacharunderscore}{\kern0pt}abst{\isacharunderscore}{\kern0pt}ab{\isacharparenleft}{\kern0pt}{\isadigit{5}}{\isacharparenright}{\kern0pt}\isanewline
\ \ \ \ \ \ \ \ \ \ \ \ \isacommand{have}\isamarkupfalse%
\ {\isachardoublequoteopen}nabla\ {\isasymturnstile}\ swap\ {\isacharbrackleft}{\kern0pt}{\isacharparenleft}{\kern0pt}a{\isacharcomma}{\kern0pt}b{\isacharparenright}{\kern0pt}{\isacharbrackright}{\kern0pt}\ t{\isadigit{2}}{\isacharprime}{\kern0pt}\ {\isasymapprox}\ swap\ {\isacharparenleft}{\kern0pt}{\isacharbrackleft}{\kern0pt}{\isacharparenleft}{\kern0pt}a{\isacharcomma}{\kern0pt}b{\isacharparenright}{\kern0pt}{\isacharbrackright}{\kern0pt}{\isacharat}{\kern0pt}{\isacharbrackleft}{\kern0pt}{\isacharparenleft}{\kern0pt}b{\isacharcomma}{\kern0pt}c{\isacharparenright}{\kern0pt}{\isacharbrackright}{\kern0pt}{\isacharparenright}{\kern0pt}\ t{\isadigit{3}}{\isacharprime}{\kern0pt}{\isachardoublequoteclose}\isanewline
\ \ \ \ \ \ \ \ \ \ \ \ \ \ \isacommand{using}\isamarkupfalse%
\ equ{\isacharunderscore}{\kern0pt}equivariance{\isacharbrackleft}{\kern0pt}OF\ equ{\isacharunderscore}{\kern0pt}abst{\isacharunderscore}{\kern0pt}ab{\isacharparenleft}{\kern0pt}{\isadigit{5}}{\isacharparenright}{\kern0pt}{\isacharbrackright}{\kern0pt}\isanewline
\ \ \ \ \ \ \ \ \ \ \ \ \ swap{\isacharunderscore}{\kern0pt}append{\isacharbrackleft}{\kern0pt}of\ {\isacartoucheopen}{\isacharbrackleft}{\kern0pt}{\isacharparenleft}{\kern0pt}a{\isacharcomma}{\kern0pt}\ b{\isacharparenright}{\kern0pt}{\isacharbrackright}{\kern0pt}{\isacartoucheclose}\ {\isacartoucheopen}{\isacharbrackleft}{\kern0pt}{\isacharparenleft}{\kern0pt}b{\isacharcomma}{\kern0pt}c{\isacharparenright}{\kern0pt}{\isacharbrackright}{\kern0pt}{\isacartoucheclose}\ t{\isadigit{3}}{\isacharprime}{\kern0pt}{\isacharbrackright}{\kern0pt}\ \isacommand{by}\isamarkupfalse%
\ simp\isanewline
\ \ \ \ \ \ \ \ \ \ \ \ \isacommand{then}\isamarkupfalse%
\ \isacommand{have}\isamarkupfalse%
\ iii{\isacharcolon}{\kern0pt}\ {\isachardoublequoteopen}nabla\ {\isasymturnstile}\ swap\ {\isacharbrackleft}{\kern0pt}{\isacharparenleft}{\kern0pt}a{\isacharcomma}{\kern0pt}b{\isacharparenright}{\kern0pt}{\isacharbrackright}{\kern0pt}\ t{\isadigit{2}}{\isacharprime}{\kern0pt}\ {\isasymapprox}\ swap\ {\isacharparenleft}{\kern0pt}{\isacharbrackleft}{\kern0pt}{\isacharparenleft}{\kern0pt}a{\isacharcomma}{\kern0pt}b{\isacharparenright}{\kern0pt}{\isacharcomma}{\kern0pt}{\isacharparenleft}{\kern0pt}b{\isacharcomma}{\kern0pt}c{\isacharparenright}{\kern0pt}{\isacharbrackright}{\kern0pt}{\isacharparenright}{\kern0pt}\ t{\isadigit{3}}{\isacharprime}{\kern0pt}{\isachardoublequoteclose}\isanewline
\ \ \ \ \ \ \ \ \ \ \ \ \ \ \isacommand{by}\isamarkupfalse%
\ simp\isanewline
\ \ \ \ \ \ \ \ \ \ \ \ \isanewline
\ \ \ \ \ \ \ \ \ \ \ \ \isacommand{have}\isamarkupfalse%
\ {\isachardoublequoteopen}ds\ {\isacharbrackleft}{\kern0pt}{\isacharparenleft}{\kern0pt}a{\isacharcomma}{\kern0pt}b{\isacharparenright}{\kern0pt}{\isacharcomma}{\kern0pt}\ {\isacharparenleft}{\kern0pt}b{\isacharcomma}{\kern0pt}c{\isacharparenright}{\kern0pt}{\isacharbrackright}{\kern0pt}\ {\isacharbrackleft}{\kern0pt}{\isacharparenleft}{\kern0pt}a{\isacharcomma}{\kern0pt}c{\isacharparenright}{\kern0pt}{\isacharbrackright}{\kern0pt}\ {\isacharequal}{\kern0pt}\ {\isacharbraceleft}{\kern0pt}a{\isacharcomma}{\kern0pt}b{\isacharbraceright}{\kern0pt}{\isachardoublequoteclose}\ \isanewline
\ \ \ \ \ \ \ \ \ \ \ \ \ \ \isacommand{using}\isamarkupfalse%
\ ds{\isacharunderscore}{\kern0pt}acabbc\ equ{\isacharunderscore}{\kern0pt}abst{\isacharunderscore}{\kern0pt}ab\ {\isadigit{1}}\ False\ \isacommand{by}\isamarkupfalse%
\ simp\isanewline
\ \ \ \ \ \ \ \ \ \ \ \ \isacommand{with}\isamarkupfalse%
\ a{\isacharunderscore}{\kern0pt}fresh{\isacharunderscore}{\kern0pt}t{\isadigit{3}}\ b{\isacharunderscore}{\kern0pt}fresh{\isacharunderscore}{\kern0pt}t{\isadigit{3}}\ \isanewline
\ \ \ \ \ \ \ \ \ \ \ \ \isacommand{have}\isamarkupfalse%
\ iv{\isacharcolon}{\kern0pt}\ {\isachardoublequoteopen}nabla\ {\isasymturnstile}\ swap\ {\isacharbrackleft}{\kern0pt}{\isacharparenleft}{\kern0pt}a{\isacharcomma}{\kern0pt}b{\isacharparenright}{\kern0pt}{\isacharcomma}{\kern0pt}{\isacharparenleft}{\kern0pt}b{\isacharcomma}{\kern0pt}c{\isacharparenright}{\kern0pt}{\isacharbrackright}{\kern0pt}\ t{\isadigit{3}}{\isacharprime}{\kern0pt}\ {\isasymapprox}\ swap\ {\isacharbrackleft}{\kern0pt}{\isacharparenleft}{\kern0pt}a{\isacharcomma}{\kern0pt}c{\isacharparenright}{\kern0pt}{\isacharbrackright}{\kern0pt}\ t{\isadigit{3}}{\isacharprime}{\kern0pt}{\isachardoublequoteclose}\isanewline
\ \ \ \ \ \ \ \ \ \ \ \ \ \ \isacommand{using}\isamarkupfalse%
\ equ{\isacharunderscore}{\kern0pt}equiv{\isacharunderscore}{\kern0pt}pi\ \isacommand{by}\isamarkupfalse%
\ simp\isanewline
\isanewline
\ \ \ \ \ \ \ \ \ \ \ \ \isacommand{with}\isamarkupfalse%
\ {\isadigit{1}}{\isacharparenleft}{\kern0pt}{\isadigit{1}}{\isacharparenright}{\kern0pt}\ iii\ \isacommand{have}\isamarkupfalse%
\ {\isachardoublequoteopen}nabla\ {\isasymturnstile}\ swap\ {\isacharbrackleft}{\kern0pt}{\isacharparenleft}{\kern0pt}a{\isacharcomma}{\kern0pt}b{\isacharparenright}{\kern0pt}{\isacharbrackright}{\kern0pt}\ t{\isadigit{2}}{\isacharprime}{\kern0pt}\ {\isasymapprox}\ swap\ {\isacharbrackleft}{\kern0pt}{\isacharparenleft}{\kern0pt}a{\isacharcomma}{\kern0pt}c{\isacharparenright}{\kern0pt}{\isacharbrackright}{\kern0pt}\ t{\isadigit{3}}{\isacharprime}{\kern0pt}{\isachardoublequoteclose}\isanewline
\ \ \ \ \ \ \ \ \ \ \ \ \ \ \isacommand{using}\isamarkupfalse%
\ IH{\isacharunderscore}{\kern0pt}usable{\isacharbrackleft}{\kern0pt}OF\ deptht{\isadigit{2}}{\isacharcomma}{\kern0pt}\ of\ {\isacartoucheopen}swap\ {\isacharparenleft}{\kern0pt}{\isacharbrackleft}{\kern0pt}{\isacharparenleft}{\kern0pt}a{\isacharcomma}{\kern0pt}b{\isacharparenright}{\kern0pt}{\isacharcomma}{\kern0pt}{\isacharparenleft}{\kern0pt}b{\isacharcomma}{\kern0pt}c{\isacharparenright}{\kern0pt}{\isacharbrackright}{\kern0pt}{\isacharparenright}{\kern0pt}\ t{\isadigit{3}}{\isacharprime}{\kern0pt}{\isacartoucheclose}\ {\isacartoucheopen}swap\ {\isacharbrackleft}{\kern0pt}{\isacharparenleft}{\kern0pt}a{\isacharcomma}{\kern0pt}c{\isacharparenright}{\kern0pt}{\isacharbrackright}{\kern0pt}\ t{\isadigit{3}}{\isacharprime}{\kern0pt}{\isacartoucheclose}{\isacharbrackright}{\kern0pt}\isanewline
\ \ \ \ \ \ \ \ \ \ \ \ \ \ \isacommand{by}\isamarkupfalse%
\ simp\isanewline
\isanewline
\ \ \ \ \ \ \ \ \ \ \ \ \isacommand{with}\isamarkupfalse%
\ {\isadigit{1}}{\isacharparenleft}{\kern0pt}{\isadigit{1}}{\isacharcomma}{\kern0pt}{\isadigit{3}}{\isacharcomma}{\kern0pt}{\isadigit{6}}{\isacharparenright}{\kern0pt}\ \isacommand{have}\isamarkupfalse%
\ {\isachardoublequoteopen}nabla\ {\isasymturnstile}\ t{\isadigit{1}}{\isacharprime}{\kern0pt}\ {\isasymapprox}\ swap\ {\isacharbrackleft}{\kern0pt}{\isacharparenleft}{\kern0pt}a{\isacharcomma}{\kern0pt}c{\isacharparenright}{\kern0pt}{\isacharbrackright}{\kern0pt}\ t{\isadigit{3}}{\isacharprime}{\kern0pt}{\isachardoublequoteclose}\isanewline
\ \ \ \ \ \ \ \ \ \ \ \ \ \ \isacommand{using}\isamarkupfalse%
\ IH{\isacharunderscore}{\kern0pt}usable{\isacharbrackleft}{\kern0pt}OF\ deptht{\isadigit{1}}{\isacharcomma}{\kern0pt}\ of\ {\isacartoucheopen}swap\ {\isacharbrackleft}{\kern0pt}{\isacharparenleft}{\kern0pt}a{\isacharcomma}{\kern0pt}\ b{\isacharparenright}{\kern0pt}{\isacharbrackright}{\kern0pt}\ t{\isadigit{2}}{\isacharprime}{\kern0pt}{\isacartoucheclose}\ {\isacartoucheopen}swap\ {\isacharbrackleft}{\kern0pt}{\isacharparenleft}{\kern0pt}a{\isacharcomma}{\kern0pt}c{\isacharparenright}{\kern0pt}{\isacharbrackright}{\kern0pt}\ t{\isadigit{3}}{\isacharprime}{\kern0pt}{\isacartoucheclose}{\isacharbrackright}{\kern0pt}\isanewline
\ \ \ \ \ \ \ \ \ \ \ \ \ \ equ{\isacharunderscore}{\kern0pt}abst{\isacharunderscore}{\kern0pt}ab{\isacharparenleft}{\kern0pt}{\isadigit{1}}{\isacharparenright}{\kern0pt}\ \isacommand{by}\isamarkupfalse%
\ simp\isanewline
\ \ \ \ \ \ \ \ \ \ \ \ \isanewline
\ \ \ \ \ \ \ \ \ \ \ \ \isacommand{with}\isamarkupfalse%
\ {\isadigit{1}}{\isacharparenleft}{\kern0pt}{\isadigit{1}}{\isacharcomma}{\kern0pt}{\isadigit{2}}{\isacharparenright}{\kern0pt}\ equ{\isacharunderscore}{\kern0pt}abst{\isacharunderscore}{\kern0pt}ab{\isacharparenleft}{\kern0pt}{\isadigit{2}}{\isacharparenright}{\kern0pt}\ \isanewline
\ \ \ \ \ \ \ \ \ \ \ \ \isacommand{show}\isamarkupfalse%
\ {\isacharquery}{\kern0pt}thesis\ \isacommand{using}\isamarkupfalse%
\ equ{\isachardot}{\kern0pt}equ{\isacharunderscore}{\kern0pt}abst{\isacharunderscore}{\kern0pt}ab{\isacharbrackleft}{\kern0pt}OF\ False\ a{\isacharunderscore}{\kern0pt}fresh{\isacharunderscore}{\kern0pt}t{\isadigit{3}}{\isacharbrackright}{\kern0pt}\ \isacommand{by}\isamarkupfalse%
\ simp\isanewline
\ \ \ \ \ \ \ \ \ \isacommand{qed}\isamarkupfalse%
\isanewline
\ \ \ \ \ \ \ \isacommand{next}\isamarkupfalse%
\isanewline
\ \ \ \ \ \ \ \ \ \isacommand{case}\isamarkupfalse%
\ {\isacharparenleft}{\kern0pt}equ{\isacharunderscore}{\kern0pt}abst{\isacharunderscore}{\kern0pt}aa\ t{\isadigit{2}}{\isacharprime}{\kern0pt}\ t{\isadigit{3}}{\isacharprime}{\kern0pt}\ b{\isacharparenright}{\kern0pt}\isanewline
\ \ \ \ \ \ \ \ \ \isacommand{from}\isamarkupfalse%
\ this{\isacharparenleft}{\kern0pt}{\isadigit{3}}{\isacharparenright}{\kern0pt}\ {\isadigit{1}}{\isacharparenleft}{\kern0pt}{\isadigit{1}}{\isacharparenright}{\kern0pt}\isanewline
\ \ \ \ \ \ \ \ \ \isacommand{have}\isamarkupfalse%
\ i{\isacharcolon}{\kern0pt}\ {\isachardoublequoteopen}nabla\ {\isasymturnstile}\ swap\ {\isacharbrackleft}{\kern0pt}{\isacharparenleft}{\kern0pt}a{\isacharcomma}{\kern0pt}b{\isacharparenright}{\kern0pt}{\isacharbrackright}{\kern0pt}\ t{\isadigit{2}}{\isacharprime}{\kern0pt}\ {\isasymapprox}\ swap\ {\isacharbrackleft}{\kern0pt}{\isacharparenleft}{\kern0pt}a{\isacharcomma}{\kern0pt}b{\isacharparenright}{\kern0pt}{\isacharbrackright}{\kern0pt}\ t{\isadigit{3}}{\isacharprime}{\kern0pt}{\isachardoublequoteclose}\isanewline
\ \ \ \ \ \ \ \ \ \ \ \isacommand{using}\isamarkupfalse%
\ equ{\isacharunderscore}{\kern0pt}equivariance{\isacharbrackleft}{\kern0pt}of\ nabla\ t{\isadigit{2}}{\isacharprime}{\kern0pt}\ t{\isadigit{3}}{\isacharprime}{\kern0pt}\ {\isacartoucheopen}{\isacharbrackleft}{\kern0pt}{\isacharparenleft}{\kern0pt}a{\isacharcomma}{\kern0pt}b{\isacharparenright}{\kern0pt}{\isacharbrackright}{\kern0pt}{\isacartoucheclose}{\isacharbrackright}{\kern0pt}\ \isacommand{by}\isamarkupfalse%
\ simp\isanewline
\isanewline
\ \ \ \ \ \ \ \ \ \isacommand{have}\isamarkupfalse%
\ {\isachardoublequoteopen}nabla\ {\isasymturnstile}\ b\ {\isasymsharp}\ swap\ {\isacharbrackleft}{\kern0pt}{\isacharparenleft}{\kern0pt}a{\isacharcomma}{\kern0pt}b{\isacharparenright}{\kern0pt}{\isacharbrackright}{\kern0pt}\ t{\isadigit{2}}{\isacharprime}{\kern0pt}{\isachardoublequoteclose}\ \isanewline
\ \ \ \ \ \ \ \ \ \ \ \isacommand{using}\isamarkupfalse%
\ {\isadigit{1}}{\isacharparenleft}{\kern0pt}{\isadigit{1}}{\isacharcomma}{\kern0pt}{\isadigit{3}}{\isacharcomma}{\kern0pt}{\isadigit{4}}{\isacharcomma}{\kern0pt}{\isadigit{5}}{\isacharparenright}{\kern0pt}\ equ{\isacharunderscore}{\kern0pt}abst{\isacharunderscore}{\kern0pt}aa{\isacharparenleft}{\kern0pt}{\isadigit{1}}{\isacharparenright}{\kern0pt}\ fresh{\isacharunderscore}{\kern0pt}swap{\isacharunderscore}{\kern0pt}eqvt{\isacharbrackleft}{\kern0pt}of\ nabla\ a\ t{\isadigit{2}}{\isacharprime}{\kern0pt}\ {\isacartoucheopen}{\isacharbrackleft}{\kern0pt}{\isacharparenleft}{\kern0pt}a{\isacharcomma}{\kern0pt}b{\isacharparenright}{\kern0pt}{\isacharbrackright}{\kern0pt}{\isacartoucheclose}{\isacharbrackright}{\kern0pt}\isanewline
\ \ \ \ \ \ \ \ \ \ \ \isacommand{by}\isamarkupfalse%
\ auto\isanewline
\ \ \ \ \ \ \ \ \ \isacommand{then}\isamarkupfalse%
\ \isacommand{have}\isamarkupfalse%
\ ii{\isacharcolon}{\kern0pt}\ {\isachardoublequoteopen}nabla\ {\isasymturnstile}\ b\ {\isasymsharp}\ swap\ {\isacharbrackleft}{\kern0pt}{\isacharparenleft}{\kern0pt}a{\isacharcomma}{\kern0pt}b{\isacharparenright}{\kern0pt}{\isacharbrackright}{\kern0pt}\ t{\isadigit{3}}{\isacharprime}{\kern0pt}{\isachardoublequoteclose}\isanewline
\ \ \ \ \ \ \ \ \ \ \ \isacommand{using}\isamarkupfalse%
\ l{\isadigit{3}}{\isacharunderscore}{\kern0pt}jud{\isacharbrackleft}{\kern0pt}OF\ i{\isacharcomma}{\kern0pt}\ of\ b{\isacharbrackright}{\kern0pt}\ \isacommand{by}\isamarkupfalse%
\ simp\isanewline
\ \ \ \ \ \ \ \ \ \isacommand{have}\isamarkupfalse%
\ {\isachardoublequoteopen}nabla\ {\isasymturnstile}\ swapas\ {\isacharbrackleft}{\kern0pt}{\isacharparenleft}{\kern0pt}a{\isacharcomma}{\kern0pt}b{\isacharparenright}{\kern0pt}{\isacharbrackright}{\kern0pt}\ b\ {\isasymsharp}\ t{\isadigit{3}}{\isacharprime}{\kern0pt}{\isachardoublequoteclose}\isanewline
\ \ \ \ \ \ \ \ \ \ \ \isacommand{using}\isamarkupfalse%
\ fresh{\isacharunderscore}{\kern0pt}swap{\isacharunderscore}{\kern0pt}left{\isacharbrackleft}{\kern0pt}OF\ ii{\isacharbrackright}{\kern0pt}\ \isacommand{by}\isamarkupfalse%
\ simp\isanewline
\ \ \ \ \ \ \ \ \ \isacommand{then}\isamarkupfalse%
\ \isacommand{have}\isamarkupfalse%
\ a{\isacharunderscore}{\kern0pt}fresh{\isacharunderscore}{\kern0pt}t{\isadigit{3}}{\isacharcolon}{\kern0pt}\ {\isachardoublequoteopen}nabla\ {\isasymturnstile}\ a\ {\isasymsharp}\ t{\isadigit{3}}{\isacharprime}{\kern0pt}{\isachardoublequoteclose}\isanewline
\ \ \ \ \ \ \ \ \ \ \ \isacommand{using}\isamarkupfalse%
\ {\isadigit{1}}{\isacharparenleft}{\kern0pt}{\isadigit{3}}{\isacharcomma}{\kern0pt}{\isadigit{4}}{\isacharparenright}{\kern0pt}\ equ{\isacharunderscore}{\kern0pt}abst{\isacharunderscore}{\kern0pt}aa{\isacharparenleft}{\kern0pt}{\isadigit{1}}{\isacharparenright}{\kern0pt}\ \isacommand{by}\isamarkupfalse%
\ simp\isanewline
\isanewline
\ \ \ \ \ \ \ \ \ \isacommand{have}\isamarkupfalse%
\ is{\isacharunderscore}{\kern0pt}equ{\isacharcolon}{\kern0pt}\ {\isachardoublequoteopen}nabla\ {\isasymturnstile}\ t{\isadigit{1}}{\isacharprime}{\kern0pt}\ {\isasymapprox}\ swap\ {\isacharbrackleft}{\kern0pt}{\isacharparenleft}{\kern0pt}a{\isacharcomma}{\kern0pt}b{\isacharparenright}{\kern0pt}{\isacharbrackright}{\kern0pt}\ t{\isadigit{3}}{\isacharprime}{\kern0pt}{\isachardoublequoteclose}\isanewline
\ \ \ \ \ \ \ \ \ \ \ \isacommand{using}\isamarkupfalse%
\ {\isadigit{1}}\ i\ IH{\isacharunderscore}{\kern0pt}usable{\isacharbrackleft}{\kern0pt}OF\ deptht{\isadigit{1}}{\isacharcomma}{\kern0pt}\ of\ {\isacartoucheopen}swap\ {\isacharbrackleft}{\kern0pt}{\isacharparenleft}{\kern0pt}a{\isacharcomma}{\kern0pt}b{\isacharparenright}{\kern0pt}{\isacharbrackright}{\kern0pt}\ t{\isadigit{2}}{\isacharprime}{\kern0pt}{\isacartoucheclose}\ {\isacartoucheopen}swap\ {\isacharbrackleft}{\kern0pt}{\isacharparenleft}{\kern0pt}a{\isacharcomma}{\kern0pt}b{\isacharparenright}{\kern0pt}{\isacharbrackright}{\kern0pt}\ t{\isadigit{3}}{\isacharprime}{\kern0pt}{\isacartoucheclose}{\isacharbrackright}{\kern0pt}\isanewline
\ \ \ \ \ \ \ \ \ \ \ equ{\isacharunderscore}{\kern0pt}abst{\isacharunderscore}{\kern0pt}aa{\isacharparenleft}{\kern0pt}{\isadigit{1}}{\isacharparenright}{\kern0pt}\ \isacommand{by}\isamarkupfalse%
\ auto\isanewline
\ \ \ \ \ \ \ \ \ \ \isanewline
\ \ \ \ \ \ \ \ \ \isacommand{from}\isamarkupfalse%
\ a{\isacharunderscore}{\kern0pt}fresh{\isacharunderscore}{\kern0pt}t{\isadigit{3}}\ is{\isacharunderscore}{\kern0pt}equ\ \isacommand{show}\isamarkupfalse%
\ {\isacharquery}{\kern0pt}thesis\isanewline
\ \ \ \ \ \ \ \ \ \ \ \isacommand{using}\isamarkupfalse%
\ {\isadigit{1}}\ equ{\isacharunderscore}{\kern0pt}abst{\isacharunderscore}{\kern0pt}aa{\isacharparenleft}{\kern0pt}{\isadigit{1}}{\isacharcomma}{\kern0pt}{\isadigit{2}}{\isacharparenright}{\kern0pt}\ equ{\isachardot}{\kern0pt}equ{\isacharunderscore}{\kern0pt}abst{\isacharunderscore}{\kern0pt}ab\ \isacommand{by}\isamarkupfalse%
\ simp\isanewline
\ \ \ \ \ \ \ \isacommand{qed}\isamarkupfalse%
\ {\isacharparenleft}{\kern0pt}auto\ simp\ add{\isacharcolon}{\kern0pt}\ {\isadigit{1}}\ t{\isadigit{1}}{\isadigit{2}}{\isacharparenright}{\kern0pt}\isanewline
\ \ \ \ \ \isacommand{next}\isamarkupfalse%
\isanewline
\ \ \ \ \ \ \ \isacommand{case}\isamarkupfalse%
\ {\isacharparenleft}{\kern0pt}{\isadigit{2}}\ nabla\ t{\isadigit{1}}{\isacharprime}{\kern0pt}\ t{\isadigit{2}}{\isacharprime}{\kern0pt}\ a{\isacharparenright}{\kern0pt}\isanewline
\ \ \ \ \ \ \ \isacommand{have}\isamarkupfalse%
\ deptht{\isadigit{1}}{\isacharcolon}{\kern0pt}\ {\isachardoublequoteopen}depth\ t{\isadigit{1}}{\isacharprime}{\kern0pt}\ {\isacharless}{\kern0pt}\ depth\ t{\isadigit{1}}{\isachardoublequoteclose}\isanewline
\ \ \ \ \ \ \ \ \ \isacommand{using}\isamarkupfalse%
\ depth{\isachardot}{\kern0pt}simps{\isacharparenleft}{\kern0pt}{\isadigit{4}}{\isacharparenright}{\kern0pt}\ {\isadigit{2}}{\isacharparenleft}{\kern0pt}{\isadigit{2}}{\isacharparenright}{\kern0pt}\ \isacommand{by}\isamarkupfalse%
\ simp\isanewline
\ \ \ \ \ \ \ \ \ \isacommand{from}\isamarkupfalse%
\ t{\isadigit{2}}{\isadigit{3}}\ \isacommand{show}\isamarkupfalse%
\ {\isacharquery}{\kern0pt}thesis\ \isanewline
\ \ \ \ \ \ \ \ \ \isacommand{proof}\isamarkupfalse%
{\isacharparenleft}{\kern0pt}cases\ rule{\isacharcolon}{\kern0pt}\ equ{\isachardot}{\kern0pt}cases{\isacharparenright}{\kern0pt}\isanewline
\ \ \ \ \ \ \ \ \ \ \ \isacommand{case}\isamarkupfalse%
\ {\isacharparenleft}{\kern0pt}equ{\isacharunderscore}{\kern0pt}abst{\isacharunderscore}{\kern0pt}ab\ a\ c\ t{\isadigit{3}}{\isacharprime}{\kern0pt}\ t{\isadigit{2}}{\isacharprime}{\kern0pt}{\isacharparenright}{\kern0pt}\isanewline
\ \ \ \ \ \ \ \ \ \ \ \isacommand{have}\isamarkupfalse%
\ {\isachardoublequoteopen}nabla\ {\isasymturnstile}\ t{\isadigit{1}}{\isacharprime}{\kern0pt}\ {\isasymapprox}\ swap\ {\isacharbrackleft}{\kern0pt}{\isacharparenleft}{\kern0pt}a{\isacharcomma}{\kern0pt}c{\isacharparenright}{\kern0pt}{\isacharbrackright}{\kern0pt}\ t{\isadigit{3}}{\isacharprime}{\kern0pt}{\isachardoublequoteclose}\isanewline
\ \ \ \ \ \ \ \ \ \ \ \ \ \isacommand{using}\isamarkupfalse%
\ {\isadigit{2}}{\isacharparenleft}{\kern0pt}{\isadigit{1}}{\isacharcomma}{\kern0pt}{\isadigit{2}}{\isacharcomma}{\kern0pt}{\isadigit{3}}{\isacharcomma}{\kern0pt}{\isadigit{4}}{\isacharparenright}{\kern0pt}\ IH{\isacharunderscore}{\kern0pt}usable{\isacharbrackleft}{\kern0pt}OF\ deptht{\isadigit{1}}{\isacharcomma}{\kern0pt}\ of\ {\isacartoucheopen}t{\isadigit{2}}{\isacharprime}{\kern0pt}{\isacartoucheclose}\ {\isacartoucheopen}swap\ {\isacharbrackleft}{\kern0pt}{\isacharparenleft}{\kern0pt}a{\isacharcomma}{\kern0pt}c{\isacharparenright}{\kern0pt}{\isacharbrackright}{\kern0pt}\ t{\isadigit{3}}{\isacharprime}{\kern0pt}{\isacartoucheclose}{\isacharbrackright}{\kern0pt}\ \isanewline
\ \ \ \ \ \ \ \ \ \ \ \ equ{\isacharunderscore}{\kern0pt}abst{\isacharunderscore}{\kern0pt}ab{\isacharparenleft}{\kern0pt}{\isadigit{1}}{\isacharcomma}{\kern0pt}{\isadigit{4}}{\isacharcomma}{\kern0pt}{\isadigit{5}}{\isacharparenright}{\kern0pt}\ \isacommand{by}\isamarkupfalse%
\ simp\isanewline
\ \ \ \ \ \ \ \ \ \ \ \isacommand{then}\isamarkupfalse%
\ \isacommand{show}\isamarkupfalse%
\ {\isacharquery}{\kern0pt}thesis\isanewline
\ \ \ \ \ \ \ \ \ \ \ \ \ \isacommand{using}\isamarkupfalse%
\ equ{\isachardot}{\kern0pt}equ{\isacharunderscore}{\kern0pt}abst{\isacharunderscore}{\kern0pt}ab{\isacharbrackleft}{\kern0pt}OF\ equ{\isacharunderscore}{\kern0pt}abst{\isacharunderscore}{\kern0pt}ab{\isacharparenleft}{\kern0pt}{\isadigit{3}}{\isacharparenright}{\kern0pt}{\isacharbrackright}{\kern0pt}\isanewline
\ \ \ \ \ \ \ \ \ \ \ \ \ \ \ {\isadigit{2}}{\isacharparenleft}{\kern0pt}{\isadigit{1}}{\isacharcomma}{\kern0pt}{\isadigit{2}}{\isacharcomma}{\kern0pt}{\isadigit{3}}{\isacharparenright}{\kern0pt}\ equ{\isacharunderscore}{\kern0pt}abst{\isacharunderscore}{\kern0pt}ab{\isacharparenleft}{\kern0pt}{\isadigit{1}}{\isacharcomma}{\kern0pt}{\isadigit{2}}{\isacharcomma}{\kern0pt}{\isadigit{4}}{\isacharparenright}{\kern0pt}\ \isacommand{by}\isamarkupfalse%
\ auto\isanewline
\ \ \ \ \ \ \ \ \ \isacommand{next}\isamarkupfalse%
\isanewline
\ \ \ \ \ \ \ \ \ \ \ \isacommand{case}\isamarkupfalse%
\ {\isacharparenleft}{\kern0pt}equ{\isacharunderscore}{\kern0pt}abst{\isacharunderscore}{\kern0pt}aa\ t{\isadigit{2}}{\isacharprime}{\kern0pt}\ t{\isadigit{3}}{\isacharprime}{\kern0pt}\ a{\isacharparenright}{\kern0pt}\isanewline
\ \ \ \ \ \ \ \ \ \ \ \isacommand{then}\isamarkupfalse%
\ \isacommand{show}\isamarkupfalse%
\ {\isacharquery}{\kern0pt}thesis\ \isanewline
\ \ \ \ \ \ \ \ \ \ \ \ \ \isacommand{using}\isamarkupfalse%
\ {\isadigit{1}}\ {\isadigit{2}}{\isacharparenleft}{\kern0pt}{\isadigit{1}}{\isacharcomma}{\kern0pt}{\isadigit{2}}{\isacharcomma}{\kern0pt}{\isadigit{3}}{\isacharcomma}{\kern0pt}{\isadigit{4}}{\isacharparenright}{\kern0pt}\ \isacommand{by}\isamarkupfalse%
\ auto\isanewline
\ \ \ \ \ \ \ \ \ \isacommand{qed}\isamarkupfalse%
\ {\isacharparenleft}{\kern0pt}auto\ simp\ add{\isacharcolon}{\kern0pt}\ {\isadigit{2}}{\isacharparenright}{\kern0pt}\isanewline
\ \ \ \ \ \isacommand{next}\isamarkupfalse%
\isanewline
\ \ \ \ \ \ \ \isacommand{case}\isamarkupfalse%
\ {\isacharparenleft}{\kern0pt}{\isadigit{3}}\ nabla{\isacharparenright}{\kern0pt}\isanewline
\ \ \ \ \ \ \ \ \isacommand{from}\isamarkupfalse%
\ t{\isadigit{2}}{\isadigit{3}}\ \isacommand{show}\isamarkupfalse%
\ {\isacharquery}{\kern0pt}thesis\isanewline
\ \ \ \ \ \ \ \ \isacommand{proof}\isamarkupfalse%
{\isacharparenleft}{\kern0pt}cases\ rule{\isacharcolon}{\kern0pt}\ equ{\isachardot}{\kern0pt}cases{\isacharparenright}{\kern0pt}\isanewline
\ \ \ \ \ \ \ \ \ \ \isacommand{case}\isamarkupfalse%
\ equ{\isacharunderscore}{\kern0pt}unit\isanewline
\ \ \ \ \ \ \ \ \ \ \isacommand{then}\isamarkupfalse%
\ \isacommand{show}\isamarkupfalse%
\ {\isacharquery}{\kern0pt}thesis\ \isacommand{using}\isamarkupfalse%
\ t{\isadigit{1}}{\isadigit{2}}\ \isacommand{by}\isamarkupfalse%
\ auto\isanewline
\ \ \ \ \ \ \ \ \isacommand{qed}\isamarkupfalse%
\ {\isacharparenleft}{\kern0pt}auto\ simp\ add{\isacharcolon}{\kern0pt}\ {\isadigit{3}}{\isacharparenright}{\kern0pt}\isanewline
\ \ \ \ \ \isacommand{next}\isamarkupfalse%
\isanewline
\ \ \ \ \ \ \ \isacommand{case}\isamarkupfalse%
\ {\isacharparenleft}{\kern0pt}{\isadigit{4}}\ a\ b\ nabla{\isacharparenright}{\kern0pt}\isanewline
\ \ \ \ \ \ \ \ \isacommand{from}\isamarkupfalse%
\ t{\isadigit{2}}{\isadigit{3}}\ \isacommand{show}\isamarkupfalse%
\ {\isacharquery}{\kern0pt}thesis\ \isanewline
\ \ \ \ \ \ \ \ \isacommand{proof}\isamarkupfalse%
{\isacharparenleft}{\kern0pt}cases\ rule{\isacharcolon}{\kern0pt}\ equ{\isachardot}{\kern0pt}cases{\isacharparenright}{\kern0pt}\isanewline
\ \ \ \ \ \ \ \ \ \ \isacommand{case}\isamarkupfalse%
\ {\isacharparenleft}{\kern0pt}equ{\isacharunderscore}{\kern0pt}atom\ a\ b{\isacharparenright}{\kern0pt}\isanewline
\ \ \ \ \ \ \ \ \ \ \isacommand{then}\isamarkupfalse%
\ \isacommand{show}\isamarkupfalse%
\ {\isacharquery}{\kern0pt}thesis\isanewline
\ \ \ \ \ \ \ \ \ \ \ \ \isacommand{using}\isamarkupfalse%
\ t{\isadigit{1}}{\isadigit{2}}\ \isacommand{by}\isamarkupfalse%
\ auto\isanewline
\ \ \ \ \ \ \ \ \isacommand{qed}\isamarkupfalse%
\ {\isacharparenleft}{\kern0pt}auto\ simp\ add{\isacharcolon}{\kern0pt}\ {\isadigit{4}}{\isacharparenright}{\kern0pt}\isanewline
\ \ \ \ \ \isacommand{next}\isamarkupfalse%
\isanewline
\ \ \ \ \ \ \ \isacommand{case}\isamarkupfalse%
\ {\isacharparenleft}{\kern0pt}{\isadigit{5}}\ pi{\isadigit{1}}\ pi{\isadigit{2}}\ X\ nabla{\isacharparenright}{\kern0pt}\isanewline
\ \ \ \ \ \ \ \ \isacommand{from}\isamarkupfalse%
\ t{\isadigit{2}}{\isadigit{3}}\ \isacommand{show}\isamarkupfalse%
\ {\isacharquery}{\kern0pt}thesis\ \isanewline
\ \ \ \ \ \ \ \ \isacommand{proof}\isamarkupfalse%
{\isacharparenleft}{\kern0pt}cases\ rule{\isacharcolon}{\kern0pt}\ equ{\isachardot}{\kern0pt}cases{\isacharparenright}{\kern0pt}\isanewline
\ \ \ \ \ \ \ \ \ \ \isacommand{case}\isamarkupfalse%
\ {\isacharparenleft}{\kern0pt}equ{\isacharunderscore}{\kern0pt}susp\ pi{\isadigit{1}}\ pi{\isadigit{2}}\ X{\isacharparenright}{\kern0pt}\isanewline
\ \ \ \ \ \ \ \ \ \ \isacommand{then}\isamarkupfalse%
\ \isacommand{show}\isamarkupfalse%
\ {\isacharquery}{\kern0pt}thesis\ \isanewline
\ \ \ \ \ \ \ \ \ \ \ \ \isacommand{using}\isamarkupfalse%
\ ds{\isacharunderscore}{\kern0pt}trans\ {\isadigit{5}}\ \isacommand{by}\isamarkupfalse%
\ blast\isanewline
\ \ \ \ \ \ \ \ \isacommand{qed}\isamarkupfalse%
\ {\isacharparenleft}{\kern0pt}auto\ simp\ add{\isacharcolon}{\kern0pt}\ {\isadigit{5}}{\isacharparenright}{\kern0pt}\isanewline
\ \ \ \ \ \isacommand{next}\isamarkupfalse%
\isanewline
\ \ \ \ \ \ \ \isacommand{case}\isamarkupfalse%
\ {\isacharparenleft}{\kern0pt}{\isadigit{6}}\ nabla\ t{\isadigit{1}}{\isacharprime}{\kern0pt}\ t{\isadigit{2}}{\isacharprime}{\kern0pt}\ s{\isadigit{1}}{\isacharprime}{\kern0pt}\ s{\isadigit{2}}{\isacharprime}{\kern0pt}{\isacharparenright}{\kern0pt}\isanewline
\ \ \ \ \ \ \ \ \isacommand{from}\isamarkupfalse%
\ t{\isadigit{2}}{\isadigit{3}}\ \isacommand{show}\isamarkupfalse%
\ {\isacharquery}{\kern0pt}thesis\ \isanewline
\ \ \ \ \ \ \ \ \isacommand{proof}\isamarkupfalse%
{\isacharparenleft}{\kern0pt}cases\ rule{\isacharcolon}{\kern0pt}\ equ{\isachardot}{\kern0pt}cases{\isacharparenright}{\kern0pt}\isanewline
\ \ \ \ \ \ \ \ \ \ \isacommand{case}\isamarkupfalse%
\ {\isacharparenleft}{\kern0pt}equ{\isacharunderscore}{\kern0pt}paar\ t{\isadigit{2}}{\isacharprime}{\kern0pt}\ t{\isadigit{3}}{\isacharprime}{\kern0pt}\ s{\isadigit{2}}{\isacharprime}{\kern0pt}\ s{\isadigit{3}}{\isacharprime}{\kern0pt}{\isacharparenright}{\kern0pt}\isanewline
\ \ \ \ \ \ \ \ \ \ \isacommand{have}\isamarkupfalse%
\ deptht{\isadigit{1}}{\isacharcolon}{\kern0pt}\ {\isachardoublequoteopen}depth\ t{\isadigit{1}}{\isacharprime}{\kern0pt}\ {\isacharless}{\kern0pt}\ depth\ t{\isadigit{1}}{\isachardoublequoteclose}\ \isanewline
\ \ \ \ \ \ \ \ \ \ \ \ \isacommand{using}\isamarkupfalse%
\ {\isadigit{6}}{\isacharparenleft}{\kern0pt}{\isadigit{2}}{\isacharparenright}{\kern0pt}\ \isacommand{by}\isamarkupfalse%
\ simp\isanewline
\ \ \ \ \ \ \ \ \ \isacommand{have}\isamarkupfalse%
\ depths{\isadigit{1}}{\isacharcolon}{\kern0pt}\ {\isachardoublequoteopen}depth\ s{\isadigit{1}}{\isacharprime}{\kern0pt}\ {\isacharless}{\kern0pt}\ depth\ t{\isadigit{1}}{\isachardoublequoteclose}\ \isanewline
\ \ \ \ \ \ \ \ \ \ \ \isacommand{using}\isamarkupfalse%
\ {\isadigit{6}}{\isacharparenleft}{\kern0pt}{\isadigit{2}}{\isacharparenright}{\kern0pt}\ \isacommand{by}\isamarkupfalse%
\ simp\isanewline
\ \ \ \ \ \ \ \ \ \isacommand{have}\isamarkupfalse%
\ a{\isacharcolon}{\kern0pt}\ {\isachardoublequoteopen}nabla\ {\isasymturnstile}\ t{\isadigit{1}}{\isacharprime}{\kern0pt}\ {\isasymapprox}\ t{\isadigit{2}}{\isacharprime}{\kern0pt}{\isachardoublequoteclose}\ \isanewline
\ \ \ \ \ \ \ \ \ \ \ \isacommand{using}\isamarkupfalse%
\ {\isachardoublequoteopen}{\isadigit{6}}{\isachardoublequoteclose}{\isacharparenleft}{\kern0pt}{\isadigit{3}}{\isacharcomma}{\kern0pt}{\isadigit{4}}{\isacharparenright}{\kern0pt}\ equ{\isacharunderscore}{\kern0pt}paar\ \isacommand{by}\isamarkupfalse%
\ auto\isanewline
\ \ \ \ \ \ \ \ \ \isacommand{have}\isamarkupfalse%
\ b{\isacharcolon}{\kern0pt}\ {\isachardoublequoteopen}nabla\ {\isasymturnstile}\ s{\isadigit{1}}{\isacharprime}{\kern0pt}\ {\isasymapprox}\ s{\isadigit{2}}{\isacharprime}{\kern0pt}{\isachardoublequoteclose}\ \isanewline
\ \ \ \ \ \ \ \ \ \ \ \isacommand{using}\isamarkupfalse%
\ {\isachardoublequoteopen}{\isadigit{6}}{\isachardoublequoteclose}\ equ{\isacharunderscore}{\kern0pt}paar\ \isacommand{by}\isamarkupfalse%
\ auto\isanewline
\ \ \ \ \ \ \ \ \ \isacommand{from}\isamarkupfalse%
\ a\ \isacommand{have}\isamarkupfalse%
\ t{\isadigit{1}}{\isacharunderscore}{\kern0pt}equal{\isacharunderscore}{\kern0pt}t{\isadigit{3}}{\isacharcolon}{\kern0pt}\ {\isachardoublequoteopen}nabla\ {\isasymturnstile}\ t{\isadigit{1}}{\isacharprime}{\kern0pt}\ {\isasymapprox}\ t{\isadigit{3}}{\isacharprime}{\kern0pt}{\isachardoublequoteclose}\ \isanewline
\ \ \ \ \ \ \ \ \ \ \ \isacommand{using}\isamarkupfalse%
\ IH{\isacharunderscore}{\kern0pt}usable{\isacharbrackleft}{\kern0pt}of\ t{\isadigit{1}}{\isacharprime}{\kern0pt}\ t{\isadigit{2}}{\isacharprime}{\kern0pt}\ t{\isadigit{3}}{\isacharprime}{\kern0pt}{\isacharbrackright}{\kern0pt}\ equ{\isacharunderscore}{\kern0pt}paar{\isacharparenleft}{\kern0pt}{\isadigit{3}}{\isacharparenright}{\kern0pt}\ {\isadigit{6}}{\isacharparenleft}{\kern0pt}{\isadigit{1}}{\isacharparenright}{\kern0pt}\ deptht{\isadigit{1}}\isanewline
\ \ \ \ \ \ \ \ \ \ \ \isacommand{by}\isamarkupfalse%
\ simp\isanewline
\ \ \ \ \ \ \ \ \ \isacommand{from}\isamarkupfalse%
\ b\ \isacommand{have}\isamarkupfalse%
\ s{\isadigit{1}}{\isacharunderscore}{\kern0pt}equal{\isacharunderscore}{\kern0pt}s{\isadigit{3}}{\isacharcolon}{\kern0pt}\ {\isachardoublequoteopen}nabla\ {\isasymturnstile}\ s{\isadigit{1}}{\isacharprime}{\kern0pt}\ {\isasymapprox}\ s{\isadigit{3}}{\isacharprime}{\kern0pt}{\isachardoublequoteclose}\ \isanewline
\ \ \ \ \ \ \ \ \ \ \ \isacommand{using}\isamarkupfalse%
\ IH{\isacharunderscore}{\kern0pt}usable{\isacharbrackleft}{\kern0pt}of\ s{\isadigit{1}}{\isacharprime}{\kern0pt}\ s{\isadigit{2}}{\isacharprime}{\kern0pt}\ s{\isadigit{3}}{\isacharprime}{\kern0pt}{\isacharbrackright}{\kern0pt}\ equ{\isacharunderscore}{\kern0pt}paar{\isacharparenleft}{\kern0pt}{\isadigit{4}}{\isacharparenright}{\kern0pt}\ {\isadigit{6}}{\isacharparenleft}{\kern0pt}{\isadigit{1}}{\isacharparenright}{\kern0pt}\ depths{\isadigit{1}}\isanewline
\ \ \ \ \ \ \ \ \ \ \ \isacommand{by}\isamarkupfalse%
\ simp\isanewline
\ \ \ \ \ \ \ \ \ \isacommand{from}\isamarkupfalse%
\ t{\isadigit{1}}{\isacharunderscore}{\kern0pt}equal{\isacharunderscore}{\kern0pt}t{\isadigit{3}}\ s{\isadigit{1}}{\isacharunderscore}{\kern0pt}equal{\isacharunderscore}{\kern0pt}s{\isadigit{3}}\ \isacommand{show}\isamarkupfalse%
\ {\isacharquery}{\kern0pt}thesis\ \isanewline
\ \ \ \ \ \ \ \ \ \ \ \isacommand{using}\isamarkupfalse%
\ equ{\isachardot}{\kern0pt}equ{\isacharunderscore}{\kern0pt}paar{\isacharbrackleft}{\kern0pt}of\ nabla\ t{\isadigit{1}}{\isacharprime}{\kern0pt}\ t{\isadigit{3}}{\isacharprime}{\kern0pt}\ s{\isadigit{1}}{\isacharprime}{\kern0pt}\ s{\isadigit{3}}{\isacharprime}{\kern0pt}{\isacharbrackright}{\kern0pt}\ {\isadigit{6}}{\isacharparenleft}{\kern0pt}{\isadigit{1}}{\isacharcomma}{\kern0pt}{\isadigit{2}}{\isacharparenright}{\kern0pt}\ equ{\isacharunderscore}{\kern0pt}paar{\isacharparenleft}{\kern0pt}{\isadigit{2}}{\isacharparenright}{\kern0pt}\isanewline
\ \ \ \ \ \ \ \ \ \ \ \isacommand{by}\isamarkupfalse%
\ simp\isanewline
\ \ \ \ \ \ \ \ \isacommand{qed}\isamarkupfalse%
\ {\isacharparenleft}{\kern0pt}auto\ simp\ add{\isacharcolon}{\kern0pt}\ {\isadigit{6}}{\isacharparenright}{\kern0pt}\isanewline
\ \ \ \ \ \isacommand{next}\isamarkupfalse%
\isanewline
\ \ \ \ \ \ \ \isacommand{case}\isamarkupfalse%
\ {\isacharparenleft}{\kern0pt}{\isadigit{7}}\ nabla\ t{\isadigit{1}}\ t{\isadigit{2}}\ F{\isacharparenright}{\kern0pt}\isanewline
\ \ \ \ \ \ \ \ \isacommand{from}\isamarkupfalse%
\ t{\isadigit{2}}{\isadigit{3}}\ \isacommand{show}\isamarkupfalse%
\ {\isacharquery}{\kern0pt}thesis\ \isanewline
\ \ \ \ \ \ \ \isacommand{proof}\isamarkupfalse%
{\isacharparenleft}{\kern0pt}cases\ rule{\isacharcolon}{\kern0pt}\ equ{\isachardot}{\kern0pt}cases{\isacharparenright}{\kern0pt}\isanewline
\ \ \ \ \ \ \ \ \ \isacommand{case}\isamarkupfalse%
\ {\isacharparenleft}{\kern0pt}equ{\isacharunderscore}{\kern0pt}func\ t{\isadigit{1}}\ t{\isadigit{2}}\ F{\isacharparenright}{\kern0pt}\isanewline
\ \ \ \ \ \ \ \ \ \isacommand{then}\isamarkupfalse%
\ \isacommand{show}\isamarkupfalse%
\ {\isacharquery}{\kern0pt}thesis\ \isanewline
\ \ \ \ \ \ \ \ \ \ \ \isacommand{using}\isamarkupfalse%
\ {\isadigit{1}}\ {\isadigit{7}}{\isacharparenleft}{\kern0pt}{\isadigit{1}}{\isacharcomma}{\kern0pt}{\isadigit{2}}{\isacharcomma}{\kern0pt}{\isadigit{3}}{\isacharcomma}{\kern0pt}{\isadigit{4}}{\isacharparenright}{\kern0pt}\ \isacommand{by}\isamarkupfalse%
\ auto\isanewline
\ \ \ \ \ \ \ \isacommand{qed}\isamarkupfalse%
\ {\isacharparenleft}{\kern0pt}auto\ simp\ add{\isacharcolon}{\kern0pt}\ {\isadigit{7}}{\isacharparenright}{\kern0pt}\isanewline
\ \ \ \ \ \isacommand{qed}\isamarkupfalse%
\isanewline
\ \ \ \isacommand{qed}\isamarkupfalse%
\isanewline
\ \ \ \isacommand{then}\isamarkupfalse%
\ \isacommand{show}\isamarkupfalse%
\ {\isacharquery}{\kern0pt}case\ \isanewline
\ \ \ \ \ \isacommand{using}\isamarkupfalse%
\ IH{\isacharparenleft}{\kern0pt}{\isadigit{2}}{\isacharcomma}{\kern0pt}{\isadigit{3}}{\isacharparenright}{\kern0pt}\ \isacommand{by}\isamarkupfalse%
\ blast\isanewline
\ \isacommand{qed}\isamarkupfalse%
%
\endisatagproof
{\isafoldproof}%
%
\isadelimproof
\isanewline
%
\endisadelimproof
\isanewline
\isanewline
\isanewline
\isacommand{lemma}\isamarkupfalse%
\ pi{\isacharunderscore}{\kern0pt}right{\isacharunderscore}{\kern0pt}equ{\isacharunderscore}{\kern0pt}help{\isacharcolon}{\kern0pt}\isanewline
\ \ \isakeyword{assumes}\ {\isachardoublequoteopen}{\isacharparenleft}{\kern0pt}n\ {\isacharequal}{\kern0pt}\ depth\ t{\isacharparenright}{\kern0pt}{\isachardoublequoteclose}\isanewline
\ \ \isakeyword{shows}\ {\isachardoublequoteopen}nabla\ {\isasymturnstile}\ t\ {\isasymapprox}\ swap\ pi\ t\ {\isasymLongrightarrow}\ {\isasymforall}\ a\ {\isasymin}\ ds\ {\isacharbrackleft}{\kern0pt}{\isacharbrackright}{\kern0pt}\ pi{\isachardot}{\kern0pt}\ nabla\ {\isasymturnstile}\ a\ {\isasymsharp}\ t{\isachardoublequoteclose}\isanewline
%
\isadelimproof
\ \ %
\endisadelimproof
%
\isatagproof
\isacommand{using}\isamarkupfalse%
\ assms\ \isanewline
\isacommand{proof}\isamarkupfalse%
{\isacharparenleft}{\kern0pt}induction\ n\ arbitrary{\isacharcolon}{\kern0pt}\ t\ pi\ rule{\isacharcolon}{\kern0pt}\ nat{\isacharunderscore}{\kern0pt}less{\isacharunderscore}{\kern0pt}induct{\isacharparenright}{\kern0pt}\isanewline
\ \ \isacommand{case}\isamarkupfalse%
\ {\isacharparenleft}{\kern0pt}{\isadigit{1}}\ n{\isacharparenright}{\kern0pt}\isanewline
\ \ \isacommand{note}\isamarkupfalse%
\ IH\ {\isacharequal}{\kern0pt}\ this\isanewline
\ \ \isacommand{then}\isamarkupfalse%
\ \isacommand{show}\isamarkupfalse%
\ {\isacharquery}{\kern0pt}case\ \isanewline
\ \ \ \ \isacommand{proof}\isamarkupfalse%
{\isacharparenleft}{\kern0pt}cases\ rule{\isacharcolon}{\kern0pt}\ equ{\isachardot}{\kern0pt}cases{\isacharbrackleft}{\kern0pt}OF\ {\isacartoucheopen}nabla\ {\isasymturnstile}\ t\ {\isasymapprox}\ swap\ pi\ t{\isacartoucheclose}{\isacharbrackright}{\kern0pt}{\isacharparenright}{\kern0pt}\isanewline
\ \ \ \ \ \ \isacommand{case}\isamarkupfalse%
\ {\isacharparenleft}{\kern0pt}{\isadigit{1}}\ b\ c\ nabla\ t{\isadigit{2}}\ t{\isadigit{1}}{\isacharparenright}{\kern0pt}\isanewline
\ \ \ \ \ \ \ \ \isacommand{have}\isamarkupfalse%
\ deptht{\isadigit{1}}{\isacharcolon}{\kern0pt}\ {\isachardoublequoteopen}depth\ t{\isadigit{1}}\ {\isacharless}{\kern0pt}\ n{\isachardoublequoteclose}\isanewline
\ \ \ \ \ \ \ \ \ \ \isacommand{using}\isamarkupfalse%
\ {\isadigit{1}}\ depth{\isachardot}{\kern0pt}simps{\isacharparenleft}{\kern0pt}{\isadigit{4}}{\isacharparenright}{\kern0pt}\ IH\ \isacommand{by}\isamarkupfalse%
\ simp\isanewline
\ \ \ \ \ \ \ \ \isacommand{have}\isamarkupfalse%
\ swap{\isacharunderscore}{\kern0pt}pi{\isacharunderscore}{\kern0pt}t{\isadigit{1}}{\isacharunderscore}{\kern0pt}is{\isacharunderscore}{\kern0pt}t{\isadigit{2}}{\isacharcolon}{\kern0pt}\ {\isachardoublequoteopen}swap\ pi\ t{\isadigit{1}}\ {\isacharequal}{\kern0pt}\ t{\isadigit{2}}{\isachardoublequoteclose}\isanewline
\ \ \ \ \ \ \ \ \ \ \isacommand{using}\isamarkupfalse%
\ {\isadigit{1}}{\isacharparenleft}{\kern0pt}{\isadigit{2}}{\isacharcomma}{\kern0pt}{\isadigit{3}}{\isacharparenright}{\kern0pt}\ \isacommand{by}\isamarkupfalse%
\ auto\isanewline
\ \ \ \ \ \ \ \ \isacommand{have}\isamarkupfalse%
\ swap{\isacharunderscore}{\kern0pt}pi{\isacharunderscore}{\kern0pt}b{\isacharunderscore}{\kern0pt}is{\isacharunderscore}{\kern0pt}c{\isacharcolon}{\kern0pt}\ {\isachardoublequoteopen}swapas\ pi\ b\ {\isacharequal}{\kern0pt}\ c{\isachardoublequoteclose}\isanewline
\ \ \ \ \ \ \ \ \ \ \isacommand{using}\isamarkupfalse%
\ {\isadigit{1}}{\isacharparenleft}{\kern0pt}{\isadigit{2}}{\isacharcomma}{\kern0pt}{\isadigit{3}}{\isacharparenright}{\kern0pt}\ \isacommand{by}\isamarkupfalse%
\ auto\isanewline
\ \ \ \ \ \ \ \ \isacommand{with}\isamarkupfalse%
\ swap{\isacharunderscore}{\kern0pt}pi{\isacharunderscore}{\kern0pt}t{\isadigit{1}}{\isacharunderscore}{\kern0pt}is{\isacharunderscore}{\kern0pt}t{\isadigit{2}}\ \isacommand{have}\isamarkupfalse%
\ {\isachardoublequoteopen}nabla\ {\isasymturnstile}\ t{\isadigit{1}}\ {\isasymapprox}\ swap\ {\isacharbrackleft}{\kern0pt}{\isacharparenleft}{\kern0pt}b{\isacharcomma}{\kern0pt}\ swapas\ pi\ b{\isacharparenright}{\kern0pt}{\isacharbrackright}{\kern0pt}\ {\isacharparenleft}{\kern0pt}swap\ pi\ t{\isadigit{1}}{\isacharparenright}{\kern0pt}{\isachardoublequoteclose}\isanewline
\ \ \ \ \ \ \ \ \ \ \isacommand{using}\isamarkupfalse%
\ {\isadigit{1}}{\isacharparenleft}{\kern0pt}{\isadigit{1}}{\isacharcomma}{\kern0pt}{\isadigit{6}}{\isacharparenright}{\kern0pt}\ \isacommand{by}\isamarkupfalse%
\ simp\isanewline
\ \ \ \ \ \ \ \ \isacommand{then}\isamarkupfalse%
\ \isacommand{have}\isamarkupfalse%
\ {\isachardoublequoteopen}nabla\ {\isasymturnstile}\ t{\isadigit{1}}\ {\isasymapprox}\ swap\ {\isacharparenleft}{\kern0pt}{\isacharbrackleft}{\kern0pt}{\isacharparenleft}{\kern0pt}b{\isacharcomma}{\kern0pt}\ swapas\ pi\ b{\isacharparenright}{\kern0pt}{\isacharbrackright}{\kern0pt}\ {\isacharat}{\kern0pt}\ pi{\isacharparenright}{\kern0pt}\ t{\isadigit{1}}{\isachardoublequoteclose}\isanewline
\ \ \ \ \ \ \ \ \ \ \isacommand{using}\isamarkupfalse%
\ swap{\isacharunderscore}{\kern0pt}append{\isacharbrackleft}{\kern0pt}of\ {\isacartoucheopen}{\isacharbrackleft}{\kern0pt}{\isacharparenleft}{\kern0pt}b{\isacharcomma}{\kern0pt}\ swapas\ pi\ b{\isacharparenright}{\kern0pt}{\isacharbrackright}{\kern0pt}{\isacartoucheclose}\ pi\ t{\isadigit{1}}{\isacharbrackright}{\kern0pt}\ \isacommand{by}\isamarkupfalse%
\ simp\isanewline
\ \ \ \ \ \ \ \ \isacommand{with}\isamarkupfalse%
\ deptht{\isadigit{1}}\ \isacommand{have}\isamarkupfalse%
\ {\isachardoublequoteopen}{\isasymforall}\ a\ {\isasymin}\ ds\ {\isacharbrackleft}{\kern0pt}{\isacharbrackright}{\kern0pt}\ {\isacharparenleft}{\kern0pt}{\isacharparenleft}{\kern0pt}b{\isacharcomma}{\kern0pt}\ swapas\ pi\ b{\isacharparenright}{\kern0pt}{\isacharhash}{\kern0pt}pi{\isacharparenright}{\kern0pt}{\isachardot}{\kern0pt}\ nabla\ {\isasymturnstile}\ a\ {\isasymsharp}\ t{\isadigit{1}}{\isachardoublequoteclose}\isanewline
\ \ \ \ \ \ \ \ \ \ \isacommand{using}\isamarkupfalse%
\ {\isachardoublequoteopen}{\isadigit{1}}{\isachardot}{\kern0pt}IH{\isachardoublequoteclose}\ {\isadigit{1}}\ \isacommand{by}\isamarkupfalse%
\ auto\isanewline
\ \ \ \ \ \ \ \ \isacommand{then}\isamarkupfalse%
\ \isacommand{have}\isamarkupfalse%
\ ds{\isacharunderscore}{\kern0pt}minus{\isacharunderscore}{\kern0pt}bc{\isacharcolon}{\kern0pt}\ {\isachardoublequoteopen}{\isasymforall}\ a\ {\isasymin}\ ds\ {\isacharbrackleft}{\kern0pt}{\isacharbrackright}{\kern0pt}\ pi\ {\isacharminus}{\kern0pt}\ {\isacharbraceleft}{\kern0pt}b{\isacharcomma}{\kern0pt}\ c{\isacharbraceright}{\kern0pt}{\isachardot}{\kern0pt}\ nabla\ {\isasymturnstile}\ a\ {\isasymsharp}\ t{\isadigit{1}}{\isachardoublequoteclose}\isanewline
\ \ \ \ \ \ \ \ \ \ \isacommand{using}\isamarkupfalse%
\ ds{\isacharunderscore}{\kern0pt}equality\ swap{\isacharunderscore}{\kern0pt}pi{\isacharunderscore}{\kern0pt}b{\isacharunderscore}{\kern0pt}is{\isacharunderscore}{\kern0pt}c\ \isacommand{by}\isamarkupfalse%
\ auto\isanewline
\ \ \ \ \ \ \ \ \isacommand{have}\isamarkupfalse%
\ c{\isacharunderscore}{\kern0pt}fresh{\isacharunderscore}{\kern0pt}t{\isadigit{1}}{\isacharcolon}{\kern0pt}\ {\isachardoublequoteopen}nabla\ {\isasymturnstile}\ c\ {\isasymsharp}\ t{\isadigit{1}}{\isachardoublequoteclose}\isanewline
\ \ \ \ \ \ \ \ \ \ \isacommand{using}\isamarkupfalse%
\ {\isadigit{1}}\ equ{\isacharunderscore}{\kern0pt}symm\ l{\isadigit{3}}{\isacharunderscore}{\kern0pt}jud\ fresh{\isacharunderscore}{\kern0pt}swap{\isacharunderscore}{\kern0pt}eqvt\ Equ{\isacharunderscore}{\kern0pt}elims\ equ{\isacharunderscore}{\kern0pt}abst{\isacharunderscore}{\kern0pt}ab\ \isacommand{by}\isamarkupfalse%
\ meson\isanewline
\ \ \ \ \ \ \ \ \isacommand{with}\isamarkupfalse%
\ ds{\isacharunderscore}{\kern0pt}minus{\isacharunderscore}{\kern0pt}bc\ \isacommand{have}\isamarkupfalse%
\ {\isachardoublequoteopen}{\isasymforall}\ a\ {\isasymin}\ ds\ {\isacharbrackleft}{\kern0pt}{\isacharbrackright}{\kern0pt}\ pi\ {\isacharminus}{\kern0pt}\ {\isacharbraceleft}{\kern0pt}b{\isacharbraceright}{\kern0pt}{\isachardot}{\kern0pt}\ nabla\ {\isasymturnstile}\ a\ {\isasymsharp}\ t{\isadigit{1}}{\isachardoublequoteclose}\isanewline
\ \ \ \ \ \ \ \ \ \ \isacommand{by}\isamarkupfalse%
\ auto\isanewline
\ \ \ \ \ \ \ \ \isacommand{then}\isamarkupfalse%
\ \isacommand{have}\isamarkupfalse%
\ ds{\isacharunderscore}{\kern0pt}minus{\isacharunderscore}{\kern0pt}b{\isacharcolon}{\kern0pt}\ {\isachardoublequoteopen}{\isasymforall}\ a\ {\isasymin}\ ds\ {\isacharbrackleft}{\kern0pt}{\isacharbrackright}{\kern0pt}\ pi\ {\isacharminus}{\kern0pt}\ {\isacharbraceleft}{\kern0pt}b{\isacharbraceright}{\kern0pt}{\isachardot}{\kern0pt}\ nabla\ {\isasymturnstile}\ a\ {\isasymsharp}\ Abst\ b\ t{\isadigit{1}}{\isachardoublequoteclose}\isanewline
\ \ \ \ \ \ \ \ \ \ \isacommand{using}\isamarkupfalse%
\ fresh{\isacharunderscore}{\kern0pt}abst{\isacharunderscore}{\kern0pt}ab\ \isacommand{by}\isamarkupfalse%
\ blast\isanewline
\ \ \ \ \ \ \ \ \isacommand{have}\isamarkupfalse%
\ b{\isacharunderscore}{\kern0pt}fresh{\isacharunderscore}{\kern0pt}abst{\isacharunderscore}{\kern0pt}t{\isadigit{1}}{\isacharcolon}{\kern0pt}\ {\isachardoublequoteopen}nabla\ {\isasymturnstile}\ b\ {\isasymsharp}\ Abst\ b\ t{\isadigit{1}}{\isachardoublequoteclose}\isanewline
\ \ \ \ \ \ \ \ \ \ \isacommand{using}\isamarkupfalse%
\ fresh{\isacharunderscore}{\kern0pt}abst{\isacharunderscore}{\kern0pt}aa\ \isacommand{by}\isamarkupfalse%
\ simp\isanewline
\ \ \ \ \ \ \ \ \isacommand{have}\isamarkupfalse%
\ {\isachardoublequoteopen}{\isasymforall}\ a\ {\isasymin}\ ds\ {\isacharbrackleft}{\kern0pt}{\isacharbrackright}{\kern0pt}\ pi{\isachardot}{\kern0pt}\ nabla\ {\isasymturnstile}\ a\ {\isasymsharp}\ Abst\ b\ t{\isadigit{1}}{\isachardoublequoteclose}\isanewline
\ \ \ \ \ \ \ \ \isacommand{proof}\isamarkupfalse%
{\isacharparenleft}{\kern0pt}rule{\isacharcomma}{\kern0pt}\ goal{\isacharunderscore}{\kern0pt}cases{\isacharparenright}{\kern0pt}\isanewline
\ \ \ \ \ \ \ \ \ \ \isacommand{case}\isamarkupfalse%
\ {\isacharparenleft}{\kern0pt}{\isadigit{1}}\ a{\isacharparenright}{\kern0pt}\isanewline
\ \ \ \ \ \ \ \ \ \ \isacommand{then}\isamarkupfalse%
\ \isacommand{show}\isamarkupfalse%
\ {\isacharquery}{\kern0pt}case\isanewline
\ \ \ \ \ \ \ \ \ \ \ \ \isacommand{using}\isamarkupfalse%
\ b{\isacharunderscore}{\kern0pt}fresh{\isacharunderscore}{\kern0pt}abst{\isacharunderscore}{\kern0pt}t{\isadigit{1}}\ ds{\isacharunderscore}{\kern0pt}minus{\isacharunderscore}{\kern0pt}b\isanewline
\ \ \ \ \ \ \ \ \ \ \ \ \isacommand{by}\isamarkupfalse%
\ {\isacharparenleft}{\kern0pt}cases\ {\isachardoublequoteopen}a{\isacharequal}{\kern0pt}b{\isachardoublequoteclose}{\isacharcomma}{\kern0pt}\ auto{\isacharparenright}{\kern0pt}\isanewline
\ \ \ \ \ \ \ \ \isacommand{qed}\isamarkupfalse%
\isanewline
\ \ \ \ \ \ \ \ \isacommand{with}\isamarkupfalse%
\ {\isadigit{1}}\ \isacommand{show}\isamarkupfalse%
\ {\isacharquery}{\kern0pt}thesis\isanewline
\ \ \ \ \ \ \ \ \ \ \isacommand{by}\isamarkupfalse%
\ auto\isanewline
\ \ \ \ \isacommand{next}\isamarkupfalse%
\isanewline
\ \ \ \ \ \ \isacommand{case}\isamarkupfalse%
\ {\isacharparenleft}{\kern0pt}{\isadigit{2}}\ nabla\ t{\isadigit{1}}\ t{\isadigit{2}}\ b{\isacharparenright}{\kern0pt}\isanewline
\ \ \ \ \ \ \isacommand{have}\isamarkupfalse%
\ deptht{\isadigit{1}}{\isacharcolon}{\kern0pt}\ {\isachardoublequoteopen}depth\ t{\isadigit{1}}\ {\isacharless}{\kern0pt}\ n{\isachardoublequoteclose}\isanewline
\ \ \ \ \ \ \ \ \isacommand{using}\isamarkupfalse%
\ {\isadigit{2}}\ depth{\isachardot}{\kern0pt}simps{\isacharparenleft}{\kern0pt}{\isadigit{4}}{\isacharparenright}{\kern0pt}\ IH\ \isacommand{by}\isamarkupfalse%
\ simp\isanewline
\ \ \ \ \ \ \isacommand{have}\isamarkupfalse%
\ swap{\isacharunderscore}{\kern0pt}pi{\isacharunderscore}{\kern0pt}t{\isadigit{1}}{\isacharunderscore}{\kern0pt}is{\isacharunderscore}{\kern0pt}t{\isadigit{2}}{\isacharcolon}{\kern0pt}\ {\isachardoublequoteopen}swap\ pi\ t{\isadigit{1}}\ {\isacharequal}{\kern0pt}\ t{\isadigit{2}}{\isachardoublequoteclose}\isanewline
\ \ \ \ \ \ \ \ \isacommand{using}\isamarkupfalse%
\ {\isadigit{2}}{\isacharparenleft}{\kern0pt}{\isadigit{2}}{\isacharcomma}{\kern0pt}{\isadigit{3}}{\isacharparenright}{\kern0pt}\ \isacommand{by}\isamarkupfalse%
\ auto\isanewline
\ \ \ \ \ \ \isacommand{from}\isamarkupfalse%
\ swap{\isacharunderscore}{\kern0pt}pi{\isacharunderscore}{\kern0pt}t{\isadigit{1}}{\isacharunderscore}{\kern0pt}is{\isacharunderscore}{\kern0pt}t{\isadigit{2}}\ deptht{\isadigit{1}}\ \isanewline
\ \ \ \ \ \ \isacommand{have}\isamarkupfalse%
\ {\isachardoublequoteopen}{\isasymforall}\ a\ {\isasymin}\ ds\ {\isacharbrackleft}{\kern0pt}{\isacharbrackright}{\kern0pt}\ pi{\isachardot}{\kern0pt}\ nabla\ {\isasymturnstile}\ a\ {\isasymsharp}\ t{\isadigit{1}}{\isachardoublequoteclose}\isanewline
\ \ \ \ \ \ \ \ \isacommand{using}\isamarkupfalse%
\ {\isachardoublequoteopen}{\isadigit{1}}{\isachardot}{\kern0pt}IH{\isachardoublequoteclose}\ {\isadigit{2}}{\isacharparenleft}{\kern0pt}{\isadigit{1}}{\isacharcomma}{\kern0pt}{\isadigit{4}}{\isacharparenright}{\kern0pt}\ \isacommand{by}\isamarkupfalse%
\ auto\isanewline
\ \ \ \ \ \ \isacommand{then}\isamarkupfalse%
\ \isacommand{show}\isamarkupfalse%
\ {\isacharquery}{\kern0pt}thesis\ \isanewline
\ \ \ \ \ \ \ \ \isacommand{using}\isamarkupfalse%
\ {\isadigit{2}}{\isacharparenleft}{\kern0pt}{\isadigit{1}}{\isacharcomma}{\kern0pt}{\isadigit{2}}{\isacharcomma}{\kern0pt}{\isadigit{3}}{\isacharparenright}{\kern0pt}\ elem{\isacharunderscore}{\kern0pt}ds\ \isacommand{by}\isamarkupfalse%
\ fastforce\isanewline
\ \ \ \ \isacommand{next}\isamarkupfalse%
\isanewline
\ \ \ \ \ \ \isacommand{case}\isamarkupfalse%
\ {\isacharparenleft}{\kern0pt}{\isadigit{3}}\ nabla{\isacharparenright}{\kern0pt}\isanewline
\ \ \ \ \ \ \isacommand{then}\isamarkupfalse%
\ \isacommand{show}\isamarkupfalse%
\ {\isacharquery}{\kern0pt}thesis\ \isanewline
\ \ \ \ \ \ \ \ \isacommand{by}\isamarkupfalse%
\ blast\isanewline
\ \ \ \ \isacommand{next}\isamarkupfalse%
\isanewline
\ \ \ \ \ \ \isacommand{case}\isamarkupfalse%
\ {\isacharparenleft}{\kern0pt}{\isadigit{4}}\ a\ b\ nabla{\isacharparenright}{\kern0pt}\isanewline
\ \ \ \ \ \ \isacommand{then}\isamarkupfalse%
\ \isacommand{show}\isamarkupfalse%
\ {\isacharquery}{\kern0pt}thesis\isanewline
\ \ \ \ \ \ \ \ \isacommand{using}\isamarkupfalse%
\ elem{\isacharunderscore}{\kern0pt}ds\ \isacommand{by}\isamarkupfalse%
\ force\isanewline
\ \ \ \ \isacommand{next}\isamarkupfalse%
\isanewline
\ \ \ \ \ \ \isacommand{case}\isamarkupfalse%
\ {\isacharparenleft}{\kern0pt}{\isadigit{5}}\ pi{\isadigit{1}}\ pi{\isadigit{2}}\ X\ nabla{\isacharparenright}{\kern0pt}\isanewline
\ \ \ \ \ \ \isacommand{have}\isamarkupfalse%
\ {\isachardoublequoteopen}swap\ pi\ t\ {\isacharequal}{\kern0pt}\ Susp\ {\isacharparenleft}{\kern0pt}pi{\isacharat}{\kern0pt}pi{\isadigit{1}}{\isacharparenright}{\kern0pt}\ X{\isachardoublequoteclose}\ \isanewline
\ \ \ \ \ \ \ \ \isacommand{using}\isamarkupfalse%
\ {\isachardoublequoteopen}{\isadigit{5}}{\isachardoublequoteclose}{\isacharparenleft}{\kern0pt}{\isadigit{2}}{\isacharparenright}{\kern0pt}\ \isacommand{by}\isamarkupfalse%
\ auto\isanewline
\ \ \ \ \ \ \isacommand{also}\isamarkupfalse%
\ \isacommand{have}\isamarkupfalse%
\ {\isachardoublequoteopen}{\isachardot}{\kern0pt}{\isachardot}{\kern0pt}{\isachardot}{\kern0pt}\ {\isacharequal}{\kern0pt}\ Susp\ pi{\isadigit{2}}\ X{\isachardoublequoteclose}\isanewline
\ \ \ \ \ \ \ \ \isacommand{using}\isamarkupfalse%
\ {\isadigit{5}}{\isacharparenleft}{\kern0pt}{\isadigit{3}}{\isacharparenright}{\kern0pt}\ calculation\ \isacommand{by}\isamarkupfalse%
\ simp\isanewline
\ \ \ \ \ \ \isacommand{then}\isamarkupfalse%
\ \isacommand{have}\isamarkupfalse%
\ {\isachardoublequoteopen}pi{\isacharat}{\kern0pt}pi{\isadigit{1}}\ {\isacharequal}{\kern0pt}\ pi{\isadigit{2}}{\isachardoublequoteclose}\ \isacommand{by}\isamarkupfalse%
\ simp\isanewline
\ \ \ \ \ \ \isacommand{hence}\isamarkupfalse%
\ {\isachardoublequoteopen}{\isasymforall}c{\isasymin}ds\ pi{\isadigit{1}}\ {\isacharparenleft}{\kern0pt}pi{\isacharat}{\kern0pt}pi{\isadigit{1}}{\isacharparenright}{\kern0pt}{\isachardot}{\kern0pt}\ {\isacharparenleft}{\kern0pt}c{\isacharcomma}{\kern0pt}\ X{\isacharparenright}{\kern0pt}\ {\isasymin}\ nabla{\isachardoublequoteclose}\isanewline
\ \ \ \ \ \ \ \ \isacommand{using}\isamarkupfalse%
\ {\isadigit{5}}{\isacharparenleft}{\kern0pt}{\isadigit{4}}{\isacharparenright}{\kern0pt}\ \isacommand{by}\isamarkupfalse%
\ simp\isanewline
\ \ \ \ \ \ \isacommand{then}\isamarkupfalse%
\ \isacommand{show}\isamarkupfalse%
\ {\isacharquery}{\kern0pt}thesis\ \isanewline
\ \ \ \ \ \ \ \ \isacommand{using}\isamarkupfalse%
\ {\isachardoublequoteopen}{\isadigit{5}}{\isachardoublequoteclose}{\isacharparenleft}{\kern0pt}{\isadigit{1}}{\isacharcomma}{\kern0pt}{\isadigit{2}}{\isacharparenright}{\kern0pt}\ fresh{\isacharunderscore}{\kern0pt}susp\ ds{\isacharunderscore}{\kern0pt}cancel{\isacharunderscore}{\kern0pt}pi{\isacharunderscore}{\kern0pt}right{\isacharbrackleft}{\kern0pt}of\ {\isacharunderscore}{\kern0pt}\ {\isacharunderscore}{\kern0pt}\ {\isachardoublequoteopen}{\isacharbrackleft}{\kern0pt}{\isacharbrackright}{\kern0pt}{\isachardoublequoteclose}\ {\isacharunderscore}{\kern0pt}{\isacharbrackright}{\kern0pt}\isanewline
\ \ \ \ \ \ \ \ \isacommand{by}\isamarkupfalse%
\ simp\isanewline
\ \ \ \ \isacommand{next}\isamarkupfalse%
\isanewline
\ \ \ \ \ \ \isacommand{case}\isamarkupfalse%
\ {\isacharparenleft}{\kern0pt}{\isadigit{6}}\ nabla\ t{\isadigit{1}}\ t{\isadigit{2}}\ s{\isadigit{1}}\ s{\isadigit{2}}{\isacharparenright}{\kern0pt}\isanewline
\ \ \ \ \ \ \isacommand{have}\isamarkupfalse%
\ deptht{\isadigit{1}}{\isacharcolon}{\kern0pt}\ {\isachardoublequoteopen}depth\ t{\isadigit{1}}\ {\isacharless}{\kern0pt}\ n{\isachardoublequoteclose}\isanewline
\ \ \ \ \ \ \ \ \isacommand{using}\isamarkupfalse%
\ {\isadigit{6}}\ depth{\isachardot}{\kern0pt}simps{\isacharparenleft}{\kern0pt}{\isadigit{5}}{\isacharparenright}{\kern0pt}\ IH\isanewline
\ \ \ \ \ \ \ \ \isacommand{by}\isamarkupfalse%
\ simp\isanewline
\ \ \ \ \ \ \isacommand{have}\isamarkupfalse%
\ depths{\isadigit{1}}{\isacharcolon}{\kern0pt}\ {\isachardoublequoteopen}depth\ s{\isadigit{1}}\ {\isacharless}{\kern0pt}\ n{\isachardoublequoteclose}\isanewline
\ \ \ \ \ \ \ \ \isacommand{using}\isamarkupfalse%
\ {\isadigit{6}}\ depth{\isachardot}{\kern0pt}simps{\isacharparenleft}{\kern0pt}{\isadigit{5}}{\isacharparenright}{\kern0pt}\ IH\isanewline
\ \ \ \ \ \ \ \ \isacommand{by}\isamarkupfalse%
\ simp\isanewline
\ \ \ \ \ \ \isacommand{from}\isamarkupfalse%
\ deptht{\isadigit{1}}\ depths{\isadigit{1}}\ \isacommand{show}\isamarkupfalse%
\ {\isacharquery}{\kern0pt}thesis\ \isanewline
\ \ \ \ \ \ \ \ \isacommand{using}\isamarkupfalse%
\ {\isachardoublequoteopen}{\isadigit{1}}{\isachardot}{\kern0pt}IH{\isachardoublequoteclose}\ {\isadigit{6}}\ \isacommand{by}\isamarkupfalse%
\ auto\isanewline
\ \ \ \ \isacommand{next}\isamarkupfalse%
\isanewline
\ \ \ \ \ \ \isacommand{case}\isamarkupfalse%
\ {\isacharparenleft}{\kern0pt}{\isadigit{7}}\ nabla\ t{\isadigit{1}}\ t{\isadigit{2}}\ f{\isacharparenright}{\kern0pt}\isanewline
\ \ \ \ \ \ \isacommand{have}\isamarkupfalse%
\ deptht{\isadigit{1}}{\isacharcolon}{\kern0pt}\ {\isachardoublequoteopen}depth\ t{\isadigit{1}}\ {\isacharless}{\kern0pt}\ n{\isachardoublequoteclose}\isanewline
\ \ \ \ \ \ \ \ \isacommand{using}\isamarkupfalse%
\ {\isadigit{7}}{\isacharparenleft}{\kern0pt}{\isadigit{2}}{\isacharparenright}{\kern0pt}\ depth{\isachardot}{\kern0pt}simps{\isacharparenleft}{\kern0pt}{\isadigit{5}}{\isacharparenright}{\kern0pt}\ IH\ \isanewline
\ \ \ \ \ \ \ \ \isacommand{by}\isamarkupfalse%
\ auto\isanewline
\ \ \ \ \ \ \isacommand{have}\isamarkupfalse%
\ {\isachardoublequoteopen}nabla\ {\isasymturnstile}\ t{\isadigit{1}}\ {\isasymapprox}\ swap\ pi\ t{\isadigit{1}}{\isachardoublequoteclose}\isanewline
\ \ \ \ \ \ \ \ \isacommand{using}\isamarkupfalse%
\ {\isadigit{7}}\ \isacommand{by}\isamarkupfalse%
\ simp\isanewline
\ \ \ \ \ \ \isacommand{then}\isamarkupfalse%
\ \isacommand{have}\isamarkupfalse%
\ {\isachardoublequoteopen}{\isasymforall}\ a\ {\isasymin}\ ds\ {\isacharbrackleft}{\kern0pt}{\isacharbrackright}{\kern0pt}\ pi{\isachardot}{\kern0pt}\ nabla\ {\isasymturnstile}\ a\ {\isasymsharp}\ t{\isadigit{1}}{\isachardoublequoteclose}\isanewline
\ \ \ \ \ \ \ \ \isacommand{using}\isamarkupfalse%
\ {\isadigit{7}}{\isacharparenleft}{\kern0pt}{\isadigit{1}}{\isacharparenright}{\kern0pt}\ IH\ deptht{\isadigit{1}}\ \isacommand{by}\isamarkupfalse%
\ simp\isanewline
\ \ \ \ \ \ \isacommand{then}\isamarkupfalse%
\ \isacommand{show}\isamarkupfalse%
\ {\isacharquery}{\kern0pt}thesis\ \isanewline
\ \ \ \ \ \ \ \ \isacommand{using}\isamarkupfalse%
\ {\isadigit{7}}{\isacharparenleft}{\kern0pt}{\isadigit{1}}{\isacharcomma}{\kern0pt}{\isadigit{2}}{\isacharparenright}{\kern0pt}\ fresh{\isacharunderscore}{\kern0pt}func{\isacharbrackleft}{\kern0pt}of\ nabla\ {\isacharunderscore}{\kern0pt}\ t{\isadigit{1}}\ f{\isacharbrackright}{\kern0pt}\ \isanewline
\ \ \ \ \ \ \ \ \isacommand{by}\isamarkupfalse%
\ simp\isanewline
\ \ \ \ \isacommand{qed}\isamarkupfalse%
\isanewline
\ \ \isacommand{qed}\isamarkupfalse%
%
\endisatagproof
{\isafoldproof}%
%
\isadelimproof
\isanewline
%
\endisadelimproof
\isanewline
%
\isadelimtheory
\isanewline
%
\endisadelimtheory
%
\isatagtheory
\isacommand{end}\isamarkupfalse%
%
\endisatagtheory
{\isafoldtheory}%
%
\isadelimtheory
%
\endisadelimtheory
%
\end{isabellebody}%
\endinput
%:%file=~/nominal_unification_Isabelle/nominal_alpha_unification/PreEqu.thy%:%
%:%10=1%:%
%:%11=1%:%
%:%12=2%:%
%:%13=3%:%
%:%14=4%:%
%:%15=5%:%
%:%20=5%:%
%:%23=6%:%
%:%24=7%:%
%:%25=8%:%
%:%26=8%:%
%:%27=9%:%
%:%28=10%:%
%:%29=11%:%
%:%30=12%:%
%:%31=13%:%
%:%32=14%:%
%:%33=15%:%
%:%34=16%:%
%:%35=17%:%
%:%36=18%:%
%:%37=18%:%
%:%38=19%:%
%:%39=20%:%
%:%40=21%:%
%:%41=22%:%
%:%42=23%:%
%:%43=24%:%
%:%44=25%:%
%:%45=26%:%
%:%46=27%:%
%:%47=28%:%
%:%48=28%:%
%:%49=29%:%
%:%50=30%:%
%:%53=31%:%
%:%57=31%:%
%:%58=31%:%
%:%59=31%:%
%:%64=31%:%
%:%67=32%:%
%:%68=33%:%
%:%69=34%:%
%:%70=34%:%
%:%77=35%:%
%:%78=35%:%
%:%79=36%:%
%:%80=36%:%
%:%81=37%:%
%:%82=37%:%
%:%83=37%:%
%:%84=38%:%
%:%85=38%:%
%:%86=39%:%
%:%87=39%:%
%:%88=40%:%
%:%89=40%:%
%:%94=40%:%
%:%97=41%:%
%:%98=42%:%
%:%99=43%:%
%:%100=43%:%
%:%101=44%:%
%:%102=45%:%
%:%105=46%:%
%:%109=46%:%
%:%110=46%:%
%:%111=47%:%
%:%112=47%:%
%:%113=48%:%
%:%114=48%:%
%:%115=49%:%
%:%116=49%:%
%:%117=49%:%
%:%118=50%:%
%:%119=50%:%
%:%120=51%:%
%:%121=51%:%
%:%122=51%:%
%:%123=52%:%
%:%124=52%:%
%:%125=53%:%
%:%126=53%:%
%:%127=54%:%
%:%128=54%:%
%:%129=55%:%
%:%130=55%:%
%:%131=56%:%
%:%132=56%:%
%:%133=57%:%
%:%134=57%:%
%:%135=58%:%
%:%136=58%:%
%:%137=58%:%
%:%138=59%:%
%:%139=59%:%
%:%140=60%:%
%:%141=60%:%
%:%142=60%:%
%:%143=60%:%
%:%144=60%:%
%:%145=61%:%
%:%146=61%:%
%:%147=62%:%
%:%148=62%:%
%:%149=63%:%
%:%150=63%:%
%:%151=63%:%
%:%152=64%:%
%:%153=64%:%
%:%154=64%:%
%:%155=65%:%
%:%156=65%:%
%:%157=66%:%
%:%158=66%:%
%:%159=67%:%
%:%160=67%:%
%:%161=67%:%
%:%162=67%:%
%:%163=67%:%
%:%164=68%:%
%:%165=68%:%
%:%166=68%:%
%:%167=68%:%
%:%168=69%:%
%:%169=69%:%
%:%170=69%:%
%:%171=69%:%
%:%172=70%:%
%:%173=70%:%
%:%174=71%:%
%:%175=71%:%
%:%176=72%:%
%:%177=72%:%
%:%178=72%:%
%:%179=72%:%
%:%180=73%:%
%:%181=73%:%
%:%182=73%:%
%:%183=73%:%
%:%184=74%:%
%:%185=74%:%
%:%186=75%:%
%:%187=75%:%
%:%188=75%:%
%:%189=76%:%
%:%190=76%:%
%:%191=77%:%
%:%192=77%:%
%:%193=78%:%
%:%194=78%:%
%:%195=79%:%
%:%196=80%:%
%:%197=80%:%
%:%198=81%:%
%:%199=81%:%
%:%200=81%:%
%:%201=82%:%
%:%202=82%:%
%:%203=82%:%
%:%204=83%:%
%:%205=83%:%
%:%206=83%:%
%:%207=84%:%
%:%208=84%:%
%:%209=84%:%
%:%210=85%:%
%:%211=85%:%
%:%212=86%:%
%:%213=86%:%
%:%214=87%:%
%:%215=87%:%
%:%216=87%:%
%:%217=87%:%
%:%218=88%:%
%:%219=88%:%
%:%220=89%:%
%:%221=89%:%
%:%222=90%:%
%:%223=90%:%
%:%224=91%:%
%:%225=91%:%
%:%226=92%:%
%:%227=92%:%
%:%228=93%:%
%:%229=93%:%
%:%230=93%:%
%:%231=94%:%
%:%232=94%:%
%:%233=95%:%
%:%234=95%:%
%:%235=96%:%
%:%236=96%:%
%:%237=97%:%
%:%238=97%:%
%:%239=98%:%
%:%240=98%:%
%:%241=98%:%
%:%242=98%:%
%:%243=99%:%
%:%244=99%:%
%:%245=100%:%
%:%246=100%:%
%:%247=101%:%
%:%248=101%:%
%:%249=102%:%
%:%250=102%:%
%:%251=103%:%
%:%252=103%:%
%:%253=104%:%
%:%254=104%:%
%:%255=104%:%
%:%256=105%:%
%:%257=105%:%
%:%258=106%:%
%:%259=106%:%
%:%260=106%:%
%:%261=106%:%
%:%262=107%:%
%:%263=107%:%
%:%264=107%:%
%:%265=107%:%
%:%266=108%:%
%:%267=108%:%
%:%268=109%:%
%:%269=109%:%
%:%270=110%:%
%:%271=110%:%
%:%272=111%:%
%:%273=111%:%
%:%274=112%:%
%:%275=112%:%
%:%276=113%:%
%:%277=113%:%
%:%278=114%:%
%:%279=114%:%
%:%280=114%:%
%:%281=115%:%
%:%282=115%:%
%:%283=115%:%
%:%284=116%:%
%:%285=116%:%
%:%286=116%:%
%:%287=116%:%
%:%288=117%:%
%:%289=117%:%
%:%290=118%:%
%:%291=118%:%
%:%292=119%:%
%:%293=119%:%
%:%294=119%:%
%:%295=119%:%
%:%296=119%:%
%:%297=120%:%
%:%298=120%:%
%:%299=121%:%
%:%300=122%:%
%:%301=122%:%
%:%302=122%:%
%:%303=122%:%
%:%304=123%:%
%:%305=124%:%
%:%306=124%:%
%:%307=125%:%
%:%308=125%:%
%:%309=126%:%
%:%310=126%:%
%:%311=126%:%
%:%312=127%:%
%:%313=127%:%
%:%314=127%:%
%:%315=128%:%
%:%316=128%:%
%:%317=129%:%
%:%318=129%:%
%:%319=130%:%
%:%320=130%:%
%:%321=130%:%
%:%322=130%:%
%:%323=131%:%
%:%324=131%:%
%:%325=132%:%
%:%326=132%:%
%:%327=133%:%
%:%328=133%:%
%:%329=133%:%
%:%330=134%:%
%:%331=134%:%
%:%332=134%:%
%:%333=135%:%
%:%334=135%:%
%:%335=136%:%
%:%336=136%:%
%:%337=137%:%
%:%338=137%:%
%:%339=137%:%
%:%340=138%:%
%:%341=138%:%
%:%342=139%:%
%:%343=139%:%
%:%344=140%:%
%:%345=140%:%
%:%346=141%:%
%:%347=141%:%
%:%348=142%:%
%:%349=142%:%
%:%350=143%:%
%:%351=143%:%
%:%352=144%:%
%:%353=144%:%
%:%354=144%:%
%:%355=145%:%
%:%356=146%:%
%:%357=146%:%
%:%358=147%:%
%:%359=147%:%
%:%360=148%:%
%:%361=148%:%
%:%362=148%:%
%:%363=149%:%
%:%364=149%:%
%:%365=150%:%
%:%366=150%:%
%:%367=151%:%
%:%368=151%:%
%:%369=152%:%
%:%370=152%:%
%:%371=153%:%
%:%372=153%:%
%:%373=154%:%
%:%374=154%:%
%:%375=154%:%
%:%376=154%:%
%:%377=155%:%
%:%378=155%:%
%:%379=156%:%
%:%380=156%:%
%:%381=157%:%
%:%382=157%:%
%:%383=157%:%
%:%384=157%:%
%:%385=158%:%
%:%386=158%:%
%:%387=159%:%
%:%388=159%:%
%:%389=160%:%
%:%390=160%:%
%:%391=161%:%
%:%392=161%:%
%:%393=162%:%
%:%394=162%:%
%:%395=162%:%
%:%396=163%:%
%:%397=163%:%
%:%398=163%:%
%:%399=164%:%
%:%400=164%:%
%:%401=165%:%
%:%407=165%:%
%:%410=166%:%
%:%411=167%:%
%:%412=167%:%
%:%419=168%:%
%:%420=168%:%
%:%421=169%:%
%:%422=169%:%
%:%423=170%:%
%:%424=170%:%
%:%425=170%:%
%:%426=170%:%
%:%427=170%:%
%:%428=171%:%
%:%429=171%:%
%:%430=172%:%
%:%431=172%:%
%:%432=173%:%
%:%433=173%:%
%:%434=174%:%
%:%435=175%:%
%:%436=175%:%
%:%437=176%:%
%:%438=176%:%
%:%439=177%:%
%:%440=178%:%
%:%441=179%:%
%:%442=179%:%
%:%443=180%:%
%:%444=180%:%
%:%445=181%:%
%:%446=182%:%
%:%447=182%:%
%:%448=182%:%
%:%449=183%:%
%:%450=183%:%
%:%451=184%:%
%:%452=184%:%
%:%453=185%:%
%:%454=185%:%
%:%455=185%:%
%:%456=186%:%
%:%457=186%:%
%:%458=187%:%
%:%459=187%:%
%:%460=188%:%
%:%461=188%:%
%:%462=188%:%
%:%463=188%:%
%:%464=189%:%
%:%465=189%:%
%:%466=190%:%
%:%467=190%:%
%:%468=191%:%
%:%469=191%:%
%:%470=191%:%
%:%471=192%:%
%:%472=192%:%
%:%473=192%:%
%:%474=193%:%
%:%475=193%:%
%:%476=194%:%
%:%477=194%:%
%:%478=195%:%
%:%479=195%:%
%:%480=195%:%
%:%481=195%:%
%:%482=196%:%
%:%483=196%:%
%:%484=197%:%
%:%485=197%:%
%:%486=198%:%
%:%487=198%:%
%:%488=198%:%
%:%489=198%:%
%:%490=199%:%
%:%496=199%:%
%:%499=200%:%
%:%500=201%:%
%:%501=202%:%
%:%502=202%:%
%:%503=203%:%
%:%504=204%:%
%:%507=205%:%
%:%511=205%:%
%:%512=205%:%
%:%513=206%:%
%:%514=206%:%
%:%515=207%:%
%:%516=207%:%
%:%517=208%:%
%:%518=208%:%
%:%519=208%:%
%:%520=209%:%
%:%521=209%:%
%:%522=210%:%
%:%523=210%:%
%:%524=211%:%
%:%525=211%:%
%:%526=212%:%
%:%527=212%:%
%:%528=213%:%
%:%529=213%:%
%:%530=213%:%
%:%531=214%:%
%:%532=214%:%
%:%533=215%:%
%:%534=215%:%
%:%535=216%:%
%:%536=216%:%
%:%537=217%:%
%:%538=217%:%
%:%539=218%:%
%:%540=218%:%
%:%541=218%:%
%:%542=219%:%
%:%543=219%:%
%:%544=220%:%
%:%545=220%:%
%:%546=221%:%
%:%547=221%:%
%:%548=222%:%
%:%549=222%:%
%:%550=222%:%
%:%551=223%:%
%:%552=223%:%
%:%553=224%:%
%:%554=224%:%
%:%555=225%:%
%:%556=225%:%
%:%557=226%:%
%:%558=226%:%
%:%559=226%:%
%:%560=227%:%
%:%561=227%:%
%:%562=228%:%
%:%563=228%:%
%:%564=229%:%
%:%565=229%:%
%:%566=230%:%
%:%567=230%:%
%:%568=231%:%
%:%569=231%:%
%:%570=231%:%
%:%571=232%:%
%:%572=232%:%
%:%573=232%:%
%:%574=233%:%
%:%575=233%:%
%:%576=234%:%
%:%577=234%:%
%:%578=235%:%
%:%579=235%:%
%:%580=235%:%
%:%581=236%:%
%:%582=236%:%
%:%583=236%:%
%:%584=237%:%
%:%590=237%:%
%:%593=238%:%
%:%594=239%:%
%:%595=240%:%
%:%596=241%:%
%:%597=242%:%
%:%598=242%:%
%:%599=243%:%
%:%600=244%:%
%:%603=245%:%
%:%607=245%:%
%:%608=245%:%
%:%609=246%:%
%:%610=246%:%
%:%611=247%:%
%:%612=247%:%
%:%613=248%:%
%:%614=248%:%
%:%615=248%:%
%:%616=249%:%
%:%617=250%:%
%:%618=250%:%
%:%619=251%:%
%:%620=251%:%
%:%621=251%:%
%:%622=252%:%
%:%623=252%:%
%:%624=252%:%
%:%625=253%:%
%:%626=253%:%
%:%627=253%:%
%:%628=254%:%
%:%629=254%:%
%:%630=255%:%
%:%631=255%:%
%:%632=256%:%
%:%633=256%:%
%:%634=256%:%
%:%635=257%:%
%:%636=257%:%
%:%637=257%:%
%:%638=258%:%
%:%639=258%:%
%:%640=258%:%
%:%641=259%:%
%:%642=259%:%
%:%643=259%:%
%:%644=260%:%
%:%645=260%:%
%:%646=260%:%
%:%647=261%:%
%:%648=261%:%
%:%649=262%:%
%:%650=262%:%
%:%651=263%:%
%:%652=263%:%
%:%653=263%:%
%:%654=264%:%
%:%655=264%:%
%:%656=264%:%
%:%657=265%:%
%:%658=265%:%
%:%659=266%:%
%:%660=266%:%
%:%661=267%:%
%:%662=267%:%
%:%663=267%:%
%:%664=268%:%
%:%665=269%:%
%:%666=269%:%
%:%667=270%:%
%:%668=270%:%
%:%669=271%:%
%:%670=271%:%
%:%671=272%:%
%:%672=272%:%
%:%673=273%:%
%:%674=273%:%
%:%675=274%:%
%:%676=274%:%
%:%677=275%:%
%:%678=275%:%
%:%679=275%:%
%:%680=276%:%
%:%681=276%:%
%:%682=276%:%
%:%683=277%:%
%:%684=277%:%
%:%685=277%:%
%:%686=278%:%
%:%687=278%:%
%:%688=278%:%
%:%689=278%:%
%:%690=278%:%
%:%691=279%:%
%:%692=279%:%
%:%693=280%:%
%:%694=280%:%
%:%695=281%:%
%:%696=281%:%
%:%697=282%:%
%:%698=282%:%
%:%699=282%:%
%:%700=283%:%
%:%701=284%:%
%:%702=284%:%
%:%703=285%:%
%:%704=285%:%
%:%705=285%:%
%:%706=286%:%
%:%707=286%:%
%:%708=287%:%
%:%709=287%:%
%:%710=287%:%
%:%711=288%:%
%:%712=288%:%
%:%713=288%:%
%:%714=289%:%
%:%715=289%:%
%:%716=289%:%
%:%717=290%:%
%:%718=290%:%
%:%719=291%:%
%:%720=291%:%
%:%721=292%:%
%:%722=292%:%
%:%723=292%:%
%:%724=293%:%
%:%725=293%:%
%:%726=293%:%
%:%727=294%:%
%:%728=294%:%
%:%729=294%:%
%:%730=294%:%
%:%731=294%:%
%:%732=295%:%
%:%733=295%:%
%:%734=296%:%
%:%735=296%:%
%:%736=297%:%
%:%737=297%:%
%:%738=297%:%
%:%739=298%:%
%:%740=298%:%
%:%741=299%:%
%:%742=299%:%
%:%743=300%:%
%:%744=300%:%
%:%749=300%:%
%:%752=301%:%
%:%753=302%:%
%:%754=303%:%
%:%755=303%:%
%:%756=304%:%
%:%757=305%:%
%:%764=306%:%
%:%765=306%:%
%:%766=307%:%
%:%767=307%:%
%:%768=308%:%
%:%769=308%:%
%:%770=309%:%
%:%771=309%:%
%:%772=309%:%
%:%773=310%:%
%:%774=311%:%
%:%775=311%:%
%:%776=312%:%
%:%777=312%:%
%:%778=312%:%
%:%779=313%:%
%:%780=313%:%
%:%781=313%:%
%:%782=314%:%
%:%783=314%:%
%:%784=314%:%
%:%785=315%:%
%:%786=315%:%
%:%787=316%:%
%:%788=316%:%
%:%789=317%:%
%:%790=317%:%
%:%791=317%:%
%:%792=318%:%
%:%793=318%:%
%:%794=318%:%
%:%795=319%:%
%:%796=319%:%
%:%797=319%:%
%:%798=320%:%
%:%799=320%:%
%:%800=320%:%
%:%801=321%:%
%:%802=321%:%
%:%803=321%:%
%:%804=322%:%
%:%805=322%:%
%:%806=323%:%
%:%807=323%:%
%:%808=324%:%
%:%809=324%:%
%:%810=324%:%
%:%811=325%:%
%:%812=325%:%
%:%813=325%:%
%:%814=326%:%
%:%815=326%:%
%:%816=327%:%
%:%817=327%:%
%:%818=328%:%
%:%819=328%:%
%:%820=328%:%
%:%821=329%:%
%:%822=330%:%
%:%823=330%:%
%:%824=330%:%
%:%825=331%:%
%:%826=332%:%
%:%827=332%:%
%:%828=332%:%
%:%829=333%:%
%:%830=333%:%
%:%831=333%:%
%:%832=334%:%
%:%833=335%:%
%:%834=335%:%
%:%835=336%:%
%:%836=336%:%
%:%837=337%:%
%:%838=337%:%
%:%839=338%:%
%:%840=339%:%
%:%841=339%:%
%:%842=340%:%
%:%843=340%:%
%:%844=340%:%
%:%845=341%:%
%:%846=341%:%
%:%847=342%:%
%:%848=342%:%
%:%849=342%:%
%:%850=343%:%
%:%851=344%:%
%:%852=344%:%
%:%853=345%:%
%:%854=345%:%
%:%855=345%:%
%:%856=346%:%
%:%857=346%:%
%:%858=346%:%
%:%859=347%:%
%:%860=347%:%
%:%861=347%:%
%:%862=348%:%
%:%863=348%:%
%:%864=348%:%
%:%865=348%:%
%:%866=348%:%
%:%867=349%:%
%:%868=349%:%
%:%869=350%:%
%:%870=350%:%
%:%871=351%:%
%:%872=351%:%
%:%873=352%:%
%:%874=352%:%
%:%875=352%:%
%:%876=353%:%
%:%877=354%:%
%:%878=354%:%
%:%879=355%:%
%:%880=355%:%
%:%881=355%:%
%:%882=356%:%
%:%883=356%:%
%:%884=357%:%
%:%885=357%:%
%:%886=357%:%
%:%887=358%:%
%:%888=358%:%
%:%889=358%:%
%:%890=359%:%
%:%891=359%:%
%:%892=359%:%
%:%893=360%:%
%:%894=360%:%
%:%895=361%:%
%:%896=361%:%
%:%897=362%:%
%:%898=362%:%
%:%899=362%:%
%:%900=363%:%
%:%901=363%:%
%:%902=363%:%
%:%903=364%:%
%:%904=364%:%
%:%905=364%:%
%:%906=364%:%
%:%907=364%:%
%:%908=365%:%
%:%909=365%:%
%:%910=366%:%
%:%911=366%:%
%:%912=367%:%
%:%913=367%:%
%:%914=367%:%
%:%915=368%:%
%:%916=368%:%
%:%917=368%:%
%:%918=369%:%
%:%919=369%:%
%:%924=369%:%
%:%927=370%:%
%:%928=371%:%
%:%929=372%:%
%:%930=372%:%
%:%931=373%:%
%:%932=374%:%
%:%935=375%:%
%:%939=375%:%
%:%940=375%:%
%:%941=376%:%
%:%942=376%:%
%:%943=377%:%
%:%944=377%:%
%:%945=378%:%
%:%946=378%:%
%:%947=379%:%
%:%948=379%:%
%:%949=379%:%
%:%950=380%:%
%:%951=381%:%
%:%952=381%:%
%:%953=382%:%
%:%954=382%:%
%:%955=382%:%
%:%956=383%:%
%:%957=383%:%
%:%958=383%:%
%:%959=384%:%
%:%960=384%:%
%:%961=384%:%
%:%962=385%:%
%:%963=385%:%
%:%964=385%:%
%:%965=386%:%
%:%966=386%:%
%:%967=387%:%
%:%968=387%:%
%:%969=388%:%
%:%970=389%:%
%:%971=389%:%
%:%972=390%:%
%:%973=390%:%
%:%974=390%:%
%:%975=391%:%
%:%976=392%:%
%:%977=392%:%
%:%978=393%:%
%:%979=393%:%
%:%980=394%:%
%:%981=394%:%
%:%982=394%:%
%:%983=395%:%
%:%984=395%:%
%:%985=396%:%
%:%986=396%:%
%:%987=396%:%
%:%988=397%:%
%:%989=397%:%
%:%990=397%:%
%:%991=397%:%
%:%992=398%:%
%:%993=398%:%
%:%994=399%:%
%:%995=399%:%
%:%996=400%:%
%:%997=400%:%
%:%998=400%:%
%:%999=401%:%
%:%1000=401%:%
%:%1001=401%:%
%:%1002=402%:%
%:%1003=402%:%
%:%1004=403%:%
%:%1005=403%:%
%:%1006=404%:%
%:%1007=404%:%
%:%1008=404%:%
%:%1009=404%:%
%:%1010=405%:%
%:%1011=405%:%
%:%1012=406%:%
%:%1013=406%:%
%:%1018=406%:%
%:%1021=407%:%
%:%1022=408%:%
%:%1023=408%:%
%:%1024=409%:%
%:%1031=410%:%
%:%1032=410%:%
%:%1033=411%:%
%:%1034=412%:%
%:%1035=412%:%
%:%1036=412%:%
%:%1037=413%:%
%:%1038=413%:%
%:%1039=414%:%
%:%1040=414%:%
%:%1041=415%:%
%:%1042=415%:%
%:%1043=415%:%
%:%1044=416%:%
%:%1045=416%:%
%:%1046=416%:%
%:%1047=417%:%
%:%1048=417%:%
%:%1049=418%:%
%:%1050=418%:%
%:%1051=418%:%
%:%1052=419%:%
%:%1053=420%:%
%:%1054=420%:%
%:%1055=420%:%
%:%1056=421%:%
%:%1057=421%:%
%:%1058=422%:%
%:%1059=422%:%
%:%1060=423%:%
%:%1061=423%:%
%:%1062=423%:%
%:%1063=424%:%
%:%1064=424%:%
%:%1065=424%:%
%:%1066=425%:%
%:%1067=425%:%
%:%1068=426%:%
%:%1069=426%:%
%:%1070=426%:%
%:%1071=427%:%
%:%1077=427%:%
%:%1080=428%:%
%:%1081=429%:%
%:%1082=429%:%
%:%1083=430%:%
%:%1084=431%:%
%:%1091=432%:%
%:%1092=432%:%
%:%1093=433%:%
%:%1094=433%:%
%:%1095=434%:%
%:%1096=434%:%
%:%1097=435%:%
%:%1098=435%:%
%:%1099=436%:%
%:%1100=436%:%
%:%1101=436%:%
%:%1102=437%:%
%:%1103=437%:%
%:%1104=437%:%
%:%1105=438%:%
%:%1106=438%:%
%:%1107=438%:%
%:%1108=439%:%
%:%1114=439%:%
%:%1117=440%:%
%:%1118=441%:%
%:%1119=441%:%
%:%1120=442%:%
%:%1121=443%:%
%:%1124=444%:%
%:%1128=444%:%
%:%1129=444%:%
%:%1130=445%:%
%:%1131=445%:%
%:%1136=445%:%
%:%1139=446%:%
%:%1140=447%:%
%:%1141=448%:%
%:%1142=448%:%
%:%1143=449%:%
%:%1150=450%:%
%:%1151=450%:%
%:%1152=451%:%
%:%1153=451%:%
%:%1154=452%:%
%:%1155=452%:%
%:%1156=453%:%
%:%1157=453%:%
%:%1158=453%:%
%:%1159=454%:%
%:%1160=454%:%
%:%1161=455%:%
%:%1162=455%:%
%:%1163=456%:%
%:%1164=456%:%
%:%1165=456%:%
%:%1166=457%:%
%:%1167=457%:%
%:%1168=458%:%
%:%1169=458%:%
%:%1170=459%:%
%:%1171=459%:%
%:%1172=459%:%
%:%1173=460%:%
%:%1174=460%:%
%:%1175=461%:%
%:%1176=461%:%
%:%1177=461%:%
%:%1178=462%:%
%:%1179=462%:%
%:%1180=463%:%
%:%1181=463%:%
%:%1182=464%:%
%:%1183=464%:%
%:%1184=464%:%
%:%1185=465%:%
%:%1186=465%:%
%:%1187=465%:%
%:%1188=466%:%
%:%1189=466%:%
%:%1190=467%:%
%:%1191=467%:%
%:%1192=468%:%
%:%1193=468%:%
%:%1194=468%:%
%:%1195=468%:%
%:%1196=469%:%
%:%1197=469%:%
%:%1198=470%:%
%:%1199=470%:%
%:%1200=471%:%
%:%1201=471%:%
%:%1202=471%:%
%:%1203=471%:%
%:%1204=472%:%
%:%1205=472%:%
%:%1206=473%:%
%:%1207=473%:%
%:%1208=474%:%
%:%1209=474%:%
%:%1210=474%:%
%:%1211=475%:%
%:%1212=475%:%
%:%1213=475%:%
%:%1214=476%:%
%:%1215=476%:%
%:%1216=477%:%
%:%1217=477%:%
%:%1218=478%:%
%:%1219=478%:%
%:%1220=478%:%
%:%1221=478%:%
%:%1222=479%:%
%:%1223=479%:%
%:%1224=480%:%
%:%1225=480%:%
%:%1226=481%:%
%:%1227=481%:%
%:%1228=481%:%
%:%1229=482%:%
%:%1230=482%:%
%:%1231=482%:%
%:%1232=483%:%
%:%1238=483%:%
%:%1241=484%:%
%:%1242=485%:%
%:%1243=486%:%
%:%1244=486%:%
%:%1245=487%:%
%:%1246=488%:%
%:%1249=489%:%
%:%1253=489%:%
%:%1254=489%:%
%:%1255=490%:%
%:%1256=490%:%
%:%1257=491%:%
%:%1258=491%:%
%:%1259=492%:%
%:%1260=492%:%
%:%1261=493%:%
%:%1262=493%:%
%:%1263=494%:%
%:%1264=495%:%
%:%1265=495%:%
%:%1266=495%:%
%:%1267=496%:%
%:%1268=496%:%
%:%1269=497%:%
%:%1270=497%:%
%:%1271=498%:%
%:%1272=498%:%
%:%1273=499%:%
%:%1274=499%:%
%:%1275=500%:%
%:%1276=500%:%
%:%1277=501%:%
%:%1278=501%:%
%:%1279=502%:%
%:%1280=502%:%
%:%1281=502%:%
%:%1282=503%:%
%:%1283=503%:%
%:%1284=503%:%
%:%1285=504%:%
%:%1286=504%:%
%:%1287=504%:%
%:%1288=505%:%
%:%1289=505%:%
%:%1290=506%:%
%:%1291=506%:%
%:%1292=507%:%
%:%1293=507%:%
%:%1294=507%:%
%:%1295=508%:%
%:%1296=508%:%
%:%1297=509%:%
%:%1298=509%:%
%:%1299=510%:%
%:%1300=510%:%
%:%1301=511%:%
%:%1302=511%:%
%:%1303=512%:%
%:%1304=512%:%
%:%1305=513%:%
%:%1306=513%:%
%:%1307=514%:%
%:%1308=514%:%
%:%1309=514%:%
%:%1310=515%:%
%:%1311=515%:%
%:%1312=516%:%
%:%1313=516%:%
%:%1314=516%:%
%:%1315=517%:%
%:%1316=517%:%
%:%1317=517%:%
%:%1318=518%:%
%:%1319=518%:%
%:%1320=519%:%
%:%1321=519%:%
%:%1322=520%:%
%:%1323=520%:%
%:%1324=520%:%
%:%1325=520%:%
%:%1326=520%:%
%:%1327=521%:%
%:%1328=521%:%
%:%1329=522%:%
%:%1330=522%:%
%:%1331=523%:%
%:%1332=524%:%
%:%1333=524%:%
%:%1334=525%:%
%:%1335=525%:%
%:%1336=526%:%
%:%1337=526%:%
%:%1338=527%:%
%:%1339=528%:%
%:%1340=528%:%
%:%1341=529%:%
%:%1342=529%:%
%:%1343=530%:%
%:%1344=530%:%
%:%1345=530%:%
%:%1346=531%:%
%:%1347=532%:%
%:%1348=532%:%
%:%1349=532%:%
%:%1350=533%:%
%:%1351=533%:%
%:%1352=533%:%
%:%1353=534%:%
%:%1354=534%:%
%:%1355=535%:%
%:%1356=535%:%
%:%1357=535%:%
%:%1358=536%:%
%:%1359=537%:%
%:%1360=537%:%
%:%1361=538%:%
%:%1362=538%:%
%:%1363=539%:%
%:%1364=539%:%
%:%1365=540%:%
%:%1366=540%:%
%:%1367=541%:%
%:%1368=541%:%
%:%1369=541%:%
%:%1370=542%:%
%:%1371=542%:%
%:%1372=543%:%
%:%1373=544%:%
%:%1374=544%:%
%:%1375=545%:%
%:%1376=545%:%
%:%1377=545%:%
%:%1378=546%:%
%:%1379=546%:%
%:%1380=547%:%
%:%1381=547%:%
%:%1382=548%:%
%:%1383=548%:%
%:%1384=548%:%
%:%1385=549%:%
%:%1386=550%:%
%:%1387=550%:%
%:%1388=550%:%
%:%1389=551%:%
%:%1390=551%:%
%:%1391=552%:%
%:%1392=552%:%
%:%1393=553%:%
%:%1394=554%:%
%:%1395=554%:%
%:%1396=554%:%
%:%1397=555%:%
%:%1398=555%:%
%:%1399=556%:%
%:%1400=556%:%
%:%1401=557%:%
%:%1402=558%:%
%:%1403=558%:%
%:%1404=559%:%
%:%1405=559%:%
%:%1406=559%:%
%:%1407=559%:%
%:%1408=560%:%
%:%1409=560%:%
%:%1410=561%:%
%:%1411=561%:%
%:%1412=562%:%
%:%1413=562%:%
%:%1414=563%:%
%:%1415=563%:%
%:%1416=564%:%
%:%1417=564%:%
%:%1418=565%:%
%:%1419=565%:%
%:%1420=565%:%
%:%1421=566%:%
%:%1422=567%:%
%:%1423=567%:%
%:%1424=568%:%
%:%1425=568%:%
%:%1426=569%:%
%:%1427=569%:%
%:%1428=570%:%
%:%1429=570%:%
%:%1430=570%:%
%:%1431=571%:%
%:%1432=571%:%
%:%1433=571%:%
%:%1434=572%:%
%:%1435=572%:%
%:%1436=573%:%
%:%1437=573%:%
%:%1438=573%:%
%:%1439=574%:%
%:%1440=574%:%
%:%1441=574%:%
%:%1442=575%:%
%:%1443=575%:%
%:%1444=575%:%
%:%1445=576%:%
%:%1446=577%:%
%:%1447=577%:%
%:%1448=578%:%
%:%1449=578%:%
%:%1450=579%:%
%:%1451=579%:%
%:%1452=580%:%
%:%1453=581%:%
%:%1454=581%:%
%:%1455=581%:%
%:%1456=582%:%
%:%1457=582%:%
%:%1458=582%:%
%:%1459=583%:%
%:%1460=583%:%
%:%1461=584%:%
%:%1462=584%:%
%:%1463=585%:%
%:%1464=585%:%
%:%1465=586%:%
%:%1466=586%:%
%:%1467=587%:%
%:%1468=587%:%
%:%1469=587%:%
%:%1470=588%:%
%:%1471=588%:%
%:%1472=588%:%
%:%1473=589%:%
%:%1474=589%:%
%:%1475=590%:%
%:%1476=590%:%
%:%1477=591%:%
%:%1478=591%:%
%:%1479=592%:%
%:%1480=592%:%
%:%1481=593%:%
%:%1482=593%:%
%:%1483=594%:%
%:%1484=594%:%
%:%1485=594%:%
%:%1486=595%:%
%:%1487=595%:%
%:%1488=596%:%
%:%1489=596%:%
%:%1490=597%:%
%:%1491=597%:%
%:%1492=598%:%
%:%1493=598%:%
%:%1494=599%:%
%:%1495=599%:%
%:%1496=599%:%
%:%1497=600%:%
%:%1498=600%:%
%:%1499=600%:%
%:%1500=601%:%
%:%1501=601%:%
%:%1502=602%:%
%:%1503=602%:%
%:%1504=603%:%
%:%1505=603%:%
%:%1506=604%:%
%:%1507=604%:%
%:%1508=604%:%
%:%1509=605%:%
%:%1510=605%:%
%:%1511=606%:%
%:%1512=606%:%
%:%1513=607%:%
%:%1514=607%:%
%:%1515=607%:%
%:%1516=607%:%
%:%1517=607%:%
%:%1518=608%:%
%:%1519=608%:%
%:%1520=609%:%
%:%1521=609%:%
%:%1522=610%:%
%:%1523=610%:%
%:%1524=611%:%
%:%1525=611%:%
%:%1526=611%:%
%:%1527=612%:%
%:%1528=612%:%
%:%1529=613%:%
%:%1530=613%:%
%:%1531=614%:%
%:%1532=614%:%
%:%1533=614%:%
%:%1534=615%:%
%:%1535=615%:%
%:%1536=615%:%
%:%1537=616%:%
%:%1538=616%:%
%:%1539=617%:%
%:%1540=617%:%
%:%1541=618%:%
%:%1542=618%:%
%:%1543=619%:%
%:%1544=619%:%
%:%1545=619%:%
%:%1546=620%:%
%:%1547=620%:%
%:%1548=621%:%
%:%1549=621%:%
%:%1550=622%:%
%:%1551=622%:%
%:%1552=622%:%
%:%1553=623%:%
%:%1554=623%:%
%:%1555=623%:%
%:%1556=624%:%
%:%1557=624%:%
%:%1558=625%:%
%:%1559=625%:%
%:%1560=626%:%
%:%1561=626%:%
%:%1562=627%:%
%:%1563=627%:%
%:%1564=627%:%
%:%1565=628%:%
%:%1566=628%:%
%:%1567=629%:%
%:%1568=629%:%
%:%1569=630%:%
%:%1570=630%:%
%:%1571=631%:%
%:%1572=631%:%
%:%1573=631%:%
%:%1574=632%:%
%:%1575=632%:%
%:%1576=633%:%
%:%1577=633%:%
%:%1578=633%:%
%:%1579=634%:%
%:%1580=634%:%
%:%1581=635%:%
%:%1582=635%:%
%:%1583=635%:%
%:%1584=636%:%
%:%1585=636%:%
%:%1586=637%:%
%:%1587=637%:%
%:%1588=637%:%
%:%1589=638%:%
%:%1590=638%:%
%:%1591=638%:%
%:%1592=639%:%
%:%1593=639%:%
%:%1594=640%:%
%:%1595=640%:%
%:%1596=641%:%
%:%1597=641%:%
%:%1598=641%:%
%:%1599=642%:%
%:%1600=642%:%
%:%1601=643%:%
%:%1602=643%:%
%:%1603=644%:%
%:%1604=644%:%
%:%1605=644%:%
%:%1606=645%:%
%:%1607=645%:%
%:%1608=646%:%
%:%1609=646%:%
%:%1610=647%:%
%:%1611=647%:%
%:%1612=648%:%
%:%1613=648%:%
%:%1614=649%:%
%:%1615=649%:%
%:%1616=650%:%
%:%1617=650%:%
%:%1618=650%:%
%:%1619=651%:%
%:%1620=651%:%
%:%1621=652%:%
%:%1622=652%:%
%:%1623=653%:%
%:%1624=653%:%
%:%1625=653%:%
%:%1626=654%:%
%:%1627=654%:%
%:%1628=654%:%
%:%1629=655%:%
%:%1630=655%:%
%:%1631=656%:%
%:%1632=656%:%
%:%1633=657%:%
%:%1634=657%:%
%:%1635=658%:%
%:%1636=658%:%
%:%1637=658%:%
%:%1638=659%:%
%:%1639=659%:%
%:%1640=659%:%
%:%1641=660%:%
%:%1647=660%:%
%:%1650=661%:%
%:%1651=662%:%
%:%1652=663%:%
%:%1653=664%:%
%:%1654=664%:%
%:%1655=665%:%
%:%1656=666%:%
%:%1659=667%:%
%:%1663=667%:%
%:%1664=667%:%
%:%1665=668%:%
%:%1666=668%:%
%:%1667=669%:%
%:%1668=669%:%
%:%1669=670%:%
%:%1670=670%:%
%:%1671=671%:%
%:%1672=671%:%
%:%1673=671%:%
%:%1674=672%:%
%:%1675=672%:%
%:%1676=673%:%
%:%1677=673%:%
%:%1678=674%:%
%:%1679=674%:%
%:%1680=675%:%
%:%1681=675%:%
%:%1682=675%:%
%:%1683=676%:%
%:%1684=676%:%
%:%1685=677%:%
%:%1686=677%:%
%:%1687=677%:%
%:%1688=678%:%
%:%1689=678%:%
%:%1690=679%:%
%:%1691=679%:%
%:%1692=679%:%
%:%1693=680%:%
%:%1694=680%:%
%:%1695=680%:%
%:%1696=681%:%
%:%1697=681%:%
%:%1698=681%:%
%:%1699=682%:%
%:%1700=682%:%
%:%1701=682%:%
%:%1702=683%:%
%:%1703=683%:%
%:%1704=683%:%
%:%1705=684%:%
%:%1706=684%:%
%:%1707=684%:%
%:%1708=685%:%
%:%1709=685%:%
%:%1710=685%:%
%:%1711=686%:%
%:%1712=686%:%
%:%1713=686%:%
%:%1714=687%:%
%:%1715=687%:%
%:%1716=687%:%
%:%1717=688%:%
%:%1718=688%:%
%:%1719=689%:%
%:%1720=689%:%
%:%1721=689%:%
%:%1722=690%:%
%:%1723=690%:%
%:%1724=690%:%
%:%1725=691%:%
%:%1726=691%:%
%:%1727=692%:%
%:%1728=692%:%
%:%1729=692%:%
%:%1730=693%:%
%:%1731=693%:%
%:%1732=693%:%
%:%1733=694%:%
%:%1734=694%:%
%:%1735=695%:%
%:%1736=695%:%
%:%1737=695%:%
%:%1738=696%:%
%:%1739=696%:%
%:%1740=697%:%
%:%1741=697%:%
%:%1742=698%:%
%:%1743=698%:%
%:%1744=699%:%
%:%1745=699%:%
%:%1746=699%:%
%:%1747=700%:%
%:%1748=700%:%
%:%1749=701%:%
%:%1750=701%:%
%:%1751=702%:%
%:%1752=702%:%
%:%1753=703%:%
%:%1754=703%:%
%:%1755=703%:%
%:%1756=704%:%
%:%1757=704%:%
%:%1758=705%:%
%:%1759=705%:%
%:%1760=706%:%
%:%1761=706%:%
%:%1762=707%:%
%:%1763=707%:%
%:%1764=708%:%
%:%1765=708%:%
%:%1766=708%:%
%:%1767=709%:%
%:%1768=709%:%
%:%1769=710%:%
%:%1770=710%:%
%:%1771=710%:%
%:%1772=711%:%
%:%1773=711%:%
%:%1774=712%:%
%:%1775=712%:%
%:%1776=713%:%
%:%1777=713%:%
%:%1778=713%:%
%:%1779=714%:%
%:%1780=714%:%
%:%1781=714%:%
%:%1782=715%:%
%:%1783=715%:%
%:%1784=715%:%
%:%1785=716%:%
%:%1786=716%:%
%:%1787=717%:%
%:%1788=717%:%
%:%1789=718%:%
%:%1790=718%:%
%:%1791=718%:%
%:%1792=719%:%
%:%1793=719%:%
%:%1794=720%:%
%:%1795=720%:%
%:%1796=721%:%
%:%1797=721%:%
%:%1798=722%:%
%:%1799=722%:%
%:%1800=722%:%
%:%1801=723%:%
%:%1802=723%:%
%:%1803=723%:%
%:%1804=724%:%
%:%1805=724%:%
%:%1806=725%:%
%:%1807=725%:%
%:%1808=726%:%
%:%1809=726%:%
%:%1810=727%:%
%:%1811=727%:%
%:%1812=727%:%
%:%1813=728%:%
%:%1814=728%:%
%:%1815=728%:%
%:%1816=729%:%
%:%1817=729%:%
%:%1818=729%:%
%:%1819=730%:%
%:%1820=730%:%
%:%1821=730%:%
%:%1822=730%:%
%:%1823=731%:%
%:%1824=731%:%
%:%1825=732%:%
%:%1826=732%:%
%:%1827=732%:%
%:%1828=733%:%
%:%1829=733%:%
%:%1830=733%:%
%:%1831=734%:%
%:%1832=734%:%
%:%1833=735%:%
%:%1834=735%:%
%:%1835=736%:%
%:%1836=736%:%
%:%1837=737%:%
%:%1838=737%:%
%:%1839=738%:%
%:%1840=738%:%
%:%1841=739%:%
%:%1842=739%:%
%:%1843=740%:%
%:%1844=740%:%
%:%1845=741%:%
%:%1846=741%:%
%:%1847=742%:%
%:%1848=742%:%
%:%1849=743%:%
%:%1850=743%:%
%:%1851=744%:%
%:%1852=744%:%
%:%1853=744%:%
%:%1854=745%:%
%:%1855=745%:%
%:%1856=745%:%
%:%1857=746%:%
%:%1858=746%:%
%:%1859=747%:%
%:%1860=747%:%
%:%1861=748%:%
%:%1862=748%:%
%:%1863=749%:%
%:%1864=749%:%
%:%1865=750%:%
%:%1866=750%:%
%:%1867=751%:%
%:%1868=751%:%
%:%1869=752%:%
%:%1870=752%:%
%:%1871=752%:%
%:%1872=753%:%
%:%1873=753%:%
%:%1874=753%:%
%:%1875=754%:%
%:%1876=754%:%
%:%1877=754%:%
%:%1878=755%:%
%:%1879=755%:%
%:%1880=755%:%
%:%1881=756%:%
%:%1882=756%:%
%:%1883=757%:%
%:%1884=757%:%
%:%1885=758%:%
%:%1886=758%:%
%:%1887=759%:%
%:%1893=759%:%
%:%1896=760%:%
%:%1899=761%:%
%:%1904=762%:%

%
\begin{isabellebody}%
\setisabellecontext{Equ}%
%
\isadelimtheory
%
\endisadelimtheory
%
\isatagtheory
\isacommand{theory}\isamarkupfalse%
\ Equ\isanewline
\isanewline
\isakeyword{imports}\ PreEqu\isanewline
\isanewline
\isakeyword{begin}%
\endisatagtheory
{\isafoldtheory}%
%
\isadelimtheory
\isanewline
%
\endisadelimtheory
\isanewline
\isacommand{lemma}\isamarkupfalse%
\ equ{\isacharunderscore}{\kern0pt}refl{\isacharcolon}{\kern0pt}\ \isanewline
\ \ {\isachardoublequoteopen}nabla\ {\isasymturnstile}\ t\ {\isasymapprox}\ t{\isachardoublequoteclose}\isanewline
%
\isadelimproof
\ \ %
\endisadelimproof
%
\isatagproof
\isacommand{by}\isamarkupfalse%
{\isacharparenleft}{\kern0pt}induct\ t{\isacharcomma}{\kern0pt}\ auto\ simp\ add{\isacharcolon}{\kern0pt}\ ds{\isacharunderscore}{\kern0pt}def{\isacharparenright}{\kern0pt}%
\endisatagproof
{\isafoldproof}%
%
\isadelimproof
\isanewline
%
\endisadelimproof
\isanewline
\isanewline
\isacommand{lemma}\isamarkupfalse%
\ equ{\isacharunderscore}{\kern0pt}dec{\isacharunderscore}{\kern0pt}pi{\isacharcolon}{\kern0pt}\isanewline
\ \ {\isachardoublequoteopen}nabla\ {\isasymturnstile}\ swap\ pi\ t{\isadigit{1}}\ {\isasymapprox}\ swap\ pi\ t{\isadigit{2}}\ {\isasymLongrightarrow}\ nabla\ {\isasymturnstile}\ t{\isadigit{1}}\ {\isasymapprox}\ t{\isadigit{2}}{\isachardoublequoteclose}\isanewline
%
\isadelimproof
%
\endisadelimproof
%
\isatagproof
\isacommand{proof}\isamarkupfalse%
{\isacharminus}{\kern0pt}\isanewline
\ \ \isacommand{have}\isamarkupfalse%
\ i{\isacharcolon}{\kern0pt}\ {\isachardoublequoteopen}nabla\ {\isasymturnstile}\ swap\ {\isacharparenleft}{\kern0pt}rev\ pi{\isacharparenright}{\kern0pt}\ {\isacharparenleft}{\kern0pt}swap\ pi\ t{\isadigit{1}}{\isacharparenright}{\kern0pt}\ {\isasymapprox}\ t{\isadigit{1}}{\isachardoublequoteclose}\isanewline
\ \ \ \ {\isachardoublequoteopen}nabla\ {\isasymturnstile}\ swap\ {\isacharparenleft}{\kern0pt}rev\ pi{\isacharparenright}{\kern0pt}\ {\isacharparenleft}{\kern0pt}swap\ pi\ t{\isadigit{2}}{\isacharparenright}{\kern0pt}\ {\isasymapprox}\ t{\isadigit{2}}{\isachardoublequoteclose}\isanewline
\ \ \ \ \isacommand{using}\isamarkupfalse%
\ rev{\isacharunderscore}{\kern0pt}pi{\isacharunderscore}{\kern0pt}pi{\isacharunderscore}{\kern0pt}equ\ \isacommand{by}\isamarkupfalse%
\ auto\isanewline
\ \ \isacommand{assume}\isamarkupfalse%
\ {\isachardoublequoteopen}nabla\ {\isasymturnstile}\ swap\ pi\ t{\isadigit{1}}\ {\isasymapprox}\ swap\ pi\ t{\isadigit{2}}{\isachardoublequoteclose}\isanewline
\ \ \isacommand{then}\isamarkupfalse%
\ \isacommand{have}\isamarkupfalse%
\ {\isachardoublequoteopen}nabla\ {\isasymturnstile}\ swap\ {\isacharparenleft}{\kern0pt}rev\ pi{\isacharparenright}{\kern0pt}\ {\isacharparenleft}{\kern0pt}swap\ pi\ t{\isadigit{1}}{\isacharparenright}{\kern0pt}\ {\isasymapprox}\ swap\ {\isacharparenleft}{\kern0pt}rev\ pi{\isacharparenright}{\kern0pt}\ {\isacharparenleft}{\kern0pt}swap\ pi\ t{\isadigit{2}}{\isacharparenright}{\kern0pt}{\isachardoublequoteclose}\isanewline
\ \ \ \ \isacommand{using}\isamarkupfalse%
\ equ{\isacharunderscore}{\kern0pt}equivariance\ \isacommand{by}\isamarkupfalse%
\ simp\isanewline
\ \ \isacommand{then}\isamarkupfalse%
\ \isacommand{show}\isamarkupfalse%
\ {\isacharquery}{\kern0pt}thesis\ \isacommand{using}\isamarkupfalse%
\ i\ equ{\isacharunderscore}{\kern0pt}symm\ equ{\isacharunderscore}{\kern0pt}trans\ \isacommand{by}\isamarkupfalse%
\ meson\ \isanewline
\isacommand{qed}\isamarkupfalse%
%
\endisatagproof
{\isafoldproof}%
%
\isadelimproof
\isanewline
%
\endisadelimproof
\isanewline
\isanewline
\isacommand{lemma}\isamarkupfalse%
\ equ{\isacharunderscore}{\kern0pt}involutive{\isacharunderscore}{\kern0pt}left{\isacharcolon}{\kern0pt}\ \isanewline
\ \ {\isachardoublequoteopen}nabla\ {\isasymturnstile}\ swap\ {\isacharparenleft}{\kern0pt}rev\ pi{\isacharparenright}{\kern0pt}\ {\isacharparenleft}{\kern0pt}swap\ pi\ t{\isadigit{1}}{\isacharparenright}{\kern0pt}\ {\isasymapprox}\ t{\isadigit{2}}\ {\isacharequal}{\kern0pt}\ nabla\ {\isasymturnstile}\ t{\isadigit{1}}\ {\isasymapprox}\ t{\isadigit{2}}{\isachardoublequoteclose}\isanewline
%
\isadelimproof
%
\endisadelimproof
%
\isatagproof
\isacommand{proof}\isamarkupfalse%
{\isacharparenleft}{\kern0pt}auto{\isacharparenright}{\kern0pt}\isanewline
\ \ \isacommand{have}\isamarkupfalse%
\ i{\isacharcolon}{\kern0pt}\ {\isachardoublequoteopen}nabla\ {\isasymturnstile}\ t{\isadigit{1}}\ {\isasymapprox}\ swap\ {\isacharparenleft}{\kern0pt}rev\ pi{\isacharparenright}{\kern0pt}\ {\isacharparenleft}{\kern0pt}swap\ pi\ t{\isadigit{1}}{\isacharparenright}{\kern0pt}{\isachardoublequoteclose}\isanewline
\ \ \ \ \isacommand{using}\isamarkupfalse%
\ rev{\isacharunderscore}{\kern0pt}pi{\isacharunderscore}{\kern0pt}pi{\isacharunderscore}{\kern0pt}equ\ equ{\isacharunderscore}{\kern0pt}symm\ \isacommand{by}\isamarkupfalse%
\ blast\isanewline
\ \ \isacommand{show}\isamarkupfalse%
\ {\isachardoublequoteopen}nabla\ {\isasymturnstile}\ swap\ {\isacharparenleft}{\kern0pt}rev\ pi{\isacharparenright}{\kern0pt}\ {\isacharparenleft}{\kern0pt}swap\ pi\ t{\isadigit{1}}{\isacharparenright}{\kern0pt}\ {\isasymapprox}\ t{\isadigit{2}}\ {\isasymLongrightarrow}\ nabla\ {\isasymturnstile}\ t{\isadigit{1}}\ {\isasymapprox}\ t{\isadigit{2}}{\isachardoublequoteclose}\isanewline
\ \ \ \ \isacommand{using}\isamarkupfalse%
\ i\ equ{\isacharunderscore}{\kern0pt}trans\ \isacommand{by}\isamarkupfalse%
\ blast\isanewline
\ \ \isacommand{show}\isamarkupfalse%
\ {\isachardoublequoteopen}nabla\ {\isasymturnstile}\ t{\isadigit{1}}\ {\isasymapprox}\ t{\isadigit{2}}\ {\isasymLongrightarrow}\ nabla\ {\isasymturnstile}\ swap\ {\isacharparenleft}{\kern0pt}rev\ pi{\isacharparenright}{\kern0pt}\ {\isacharparenleft}{\kern0pt}swap\ pi\ t{\isadigit{1}}{\isacharparenright}{\kern0pt}\ {\isasymapprox}\ t{\isadigit{2}}{\isachardoublequoteclose}\isanewline
\ \ \ \ \isacommand{using}\isamarkupfalse%
\ i\ equ{\isacharunderscore}{\kern0pt}trans\ equ{\isacharunderscore}{\kern0pt}symm\ \isacommand{by}\isamarkupfalse%
\ blast\isanewline
\isacommand{qed}\isamarkupfalse%
%
\endisatagproof
{\isafoldproof}%
%
\isadelimproof
\isanewline
%
\endisadelimproof
\isanewline
\isanewline
\isacommand{lemma}\isamarkupfalse%
\ equ{\isacharunderscore}{\kern0pt}pi{\isacharunderscore}{\kern0pt}to{\isacharunderscore}{\kern0pt}left{\isacharcolon}{\kern0pt}\ \isanewline
\ \ {\isachardoublequoteopen}nabla\ {\isasymturnstile}\ swap\ {\isacharparenleft}{\kern0pt}rev\ pi{\isacharparenright}{\kern0pt}\ t{\isadigit{1}}\ {\isasymapprox}\ t{\isadigit{2}}\ {\isacharequal}{\kern0pt}\ nabla\ {\isasymturnstile}\ t{\isadigit{1}}\ {\isasymapprox}\ swap\ pi\ t{\isadigit{2}}{\isachardoublequoteclose}\isanewline
%
\isadelimproof
%
\endisadelimproof
%
\isatagproof
\isacommand{proof}\isamarkupfalse%
\isanewline
\isanewline
\ \ \isacommand{{\isacharbraceleft}{\kern0pt}}\isamarkupfalse%
\isacommand{assume}\isamarkupfalse%
\ i{\isacharcolon}{\kern0pt}\ {\isachardoublequoteopen}nabla\ {\isasymturnstile}\ swap\ {\isacharparenleft}{\kern0pt}rev\ pi{\isacharparenright}{\kern0pt}\ t{\isadigit{1}}\ {\isasymapprox}\ t{\isadigit{2}}{\isachardoublequoteclose}\isanewline
\ \ \isacommand{have}\isamarkupfalse%
\ {\isachardoublequoteopen}nabla\ {\isasymturnstile}\ swap\ pi\ {\isacharparenleft}{\kern0pt}swap\ {\isacharparenleft}{\kern0pt}rev\ pi{\isacharparenright}{\kern0pt}\ t{\isadigit{1}}{\isacharparenright}{\kern0pt}\ {\isasymapprox}\ swap\ pi\ t{\isadigit{2}}{\isachardoublequoteclose}\isanewline
\ \ \ \ \isacommand{using}\isamarkupfalse%
\ equ{\isacharunderscore}{\kern0pt}equivariance{\isacharbrackleft}{\kern0pt}OF\ i{\isacharcomma}{\kern0pt}\ of\ pi{\isacharbrackright}{\kern0pt}\ \isacommand{by}\isamarkupfalse%
\ simp\isanewline
\ \ \isacommand{then}\isamarkupfalse%
\ \isacommand{show}\isamarkupfalse%
\ {\isachardoublequoteopen}nabla\ {\isasymturnstile}\ t{\isadigit{1}}\ {\isasymapprox}\ swap\ pi\ t{\isadigit{2}}{\isachardoublequoteclose}\isanewline
\ \ \ \ \isacommand{using}\isamarkupfalse%
\ equ{\isacharunderscore}{\kern0pt}involutive{\isacharunderscore}{\kern0pt}left{\isacharbrackleft}{\kern0pt}of\ nabla\ {\isacartoucheopen}rev\ pi{\isacartoucheclose}\ t{\isadigit{1}}\ {\isacartoucheopen}swap\ pi\ t{\isadigit{2}}{\isacartoucheclose}{\isacharbrackright}{\kern0pt}\ rev{\isacharunderscore}{\kern0pt}rev{\isacharunderscore}{\kern0pt}ident{\isacharbrackleft}{\kern0pt}of\ pi{\isacharbrackright}{\kern0pt}\isanewline
\ \ \ \ \isacommand{by}\isamarkupfalse%
\ simp\isacommand{{\isacharbraceright}{\kern0pt}}\isamarkupfalse%
\isanewline
\isanewline
\ \ \isacommand{{\isacharbraceleft}{\kern0pt}}\isamarkupfalse%
\isacommand{assume}\isamarkupfalse%
\ i{\isacharcolon}{\kern0pt}\ {\isachardoublequoteopen}nabla\ {\isasymturnstile}\ t{\isadigit{1}}\ {\isasymapprox}\ swap\ pi\ t{\isadigit{2}}{\isachardoublequoteclose}\isanewline
\ \ \isacommand{have}\isamarkupfalse%
\ ii{\isacharcolon}{\kern0pt}\ {\isachardoublequoteopen}nabla\ {\isasymturnstile}\ swap\ {\isacharparenleft}{\kern0pt}rev\ pi{\isacharparenright}{\kern0pt}\ t{\isadigit{1}}\ {\isasymapprox}\ swap\ {\isacharparenleft}{\kern0pt}rev\ pi{\isacharparenright}{\kern0pt}\ {\isacharparenleft}{\kern0pt}swap\ pi\ t{\isadigit{2}}{\isacharparenright}{\kern0pt}{\isachardoublequoteclose}\isanewline
\ \ \ \ \isacommand{using}\isamarkupfalse%
\ equ{\isacharunderscore}{\kern0pt}equivariance{\isacharbrackleft}{\kern0pt}OF\ i{\isacharcomma}{\kern0pt}\ of\ {\isacartoucheopen}rev\ pi{\isacartoucheclose}{\isacharbrackright}{\kern0pt}\ \isacommand{by}\isamarkupfalse%
\ simp\isanewline
\ \ \isacommand{then}\isamarkupfalse%
\ \isacommand{have}\isamarkupfalse%
\ iii{\isacharcolon}{\kern0pt}\ {\isachardoublequoteopen}nabla\ {\isasymturnstile}\ swap\ {\isacharparenleft}{\kern0pt}rev\ pi{\isacharparenright}{\kern0pt}\ {\isacharparenleft}{\kern0pt}swap\ pi\ t{\isadigit{2}}{\isacharparenright}{\kern0pt}\ {\isasymapprox}\ swap\ {\isacharparenleft}{\kern0pt}rev\ pi{\isacharparenright}{\kern0pt}\ t{\isadigit{1}}{\isachardoublequoteclose}\isanewline
\ \ \ \ \isacommand{using}\isamarkupfalse%
\ equ{\isacharunderscore}{\kern0pt}symm{\isacharbrackleft}{\kern0pt}OF\ ii{\isacharbrackright}{\kern0pt}\ \isacommand{by}\isamarkupfalse%
\ simp\isanewline
\ \ \isacommand{then}\isamarkupfalse%
\ \isacommand{have}\isamarkupfalse%
\ iv{\isacharcolon}{\kern0pt}\ {\isachardoublequoteopen}nabla\ {\isasymturnstile}\ t{\isadigit{2}}\ {\isasymapprox}\ swap\ {\isacharparenleft}{\kern0pt}rev\ pi{\isacharparenright}{\kern0pt}\ t{\isadigit{1}}{\isachardoublequoteclose}\isanewline
\ \ \ \ \isacommand{using}\isamarkupfalse%
\ equ{\isacharunderscore}{\kern0pt}involutive{\isacharunderscore}{\kern0pt}left{\isacharbrackleft}{\kern0pt}of\ nabla\ pi\ t{\isadigit{2}}\ {\isacartoucheopen}swap\ {\isacharparenleft}{\kern0pt}rev\ pi{\isacharparenright}{\kern0pt}\ t{\isadigit{1}}{\isacartoucheclose}{\isacharbrackright}{\kern0pt}\ \isacommand{by}\isamarkupfalse%
\ simp\isanewline
\ \ \isacommand{then}\isamarkupfalse%
\ \isacommand{show}\isamarkupfalse%
\ {\isachardoublequoteopen}nabla\ {\isasymturnstile}\ swap\ {\isacharparenleft}{\kern0pt}rev\ pi{\isacharparenright}{\kern0pt}\ t{\isadigit{1}}\ {\isasymapprox}\ t{\isadigit{2}}{\isachardoublequoteclose}\isanewline
\ \ \ \ \isacommand{using}\isamarkupfalse%
\ equ{\isacharunderscore}{\kern0pt}symm{\isacharbrackleft}{\kern0pt}OF\ iv{\isacharbrackright}{\kern0pt}\ \isacommand{by}\isamarkupfalse%
\ simp\isacommand{{\isacharbraceright}{\kern0pt}}\isamarkupfalse%
\isanewline
\isanewline
\isacommand{qed}\isamarkupfalse%
%
\endisatagproof
{\isafoldproof}%
%
\isadelimproof
\isanewline
%
\endisadelimproof
\ \ \ \ \isanewline
\isanewline
\isacommand{lemma}\isamarkupfalse%
\ equ{\isacharunderscore}{\kern0pt}pi{\isacharunderscore}{\kern0pt}to{\isacharunderscore}{\kern0pt}right{\isacharcolon}{\kern0pt}\ \isanewline
\ \ {\isachardoublequoteopen}nabla{\isasymturnstile}t{\isadigit{1}}\ {\isasymapprox}\ swap\ {\isacharparenleft}{\kern0pt}rev\ pi{\isacharparenright}{\kern0pt}\ t{\isadigit{2}}\ {\isacharequal}{\kern0pt}\ nabla{\isasymturnstile}swap\ pi\ t{\isadigit{1}}{\isasymapprox}t{\isadigit{2}}{\isachardoublequoteclose}\isanewline
%
\isadelimproof
%
\endisadelimproof
%
\isatagproof
\isacommand{proof}\isamarkupfalse%
\isanewline
\ \ \isacommand{{\isacharbraceleft}{\kern0pt}}\isamarkupfalse%
\isacommand{assume}\isamarkupfalse%
\ i{\isacharcolon}{\kern0pt}\ {\isachardoublequoteopen}nabla\ {\isasymturnstile}\ t{\isadigit{1}}\ {\isasymapprox}\ swap\ {\isacharparenleft}{\kern0pt}rev\ pi{\isacharparenright}{\kern0pt}\ t{\isadigit{2}}{\isachardoublequoteclose}\isanewline
\ \ \ \ \isacommand{then}\isamarkupfalse%
\ \isacommand{show}\isamarkupfalse%
\ {\isachardoublequoteopen}nabla\ {\isasymturnstile}\ swap\ pi\ t{\isadigit{1}}\ {\isasymapprox}\ \ t{\isadigit{2}}{\isachardoublequoteclose}\isanewline
\ \ \ \ \ \ \isacommand{using}\isamarkupfalse%
\ equ{\isacharunderscore}{\kern0pt}involutive{\isacharunderscore}{\kern0pt}left\ equ{\isacharunderscore}{\kern0pt}dec{\isacharunderscore}{\kern0pt}pi\ \isacommand{by}\isamarkupfalse%
\ blast\isacommand{{\isacharbraceright}{\kern0pt}}\isamarkupfalse%
\isanewline
\ \ \isacommand{{\isacharbraceleft}{\kern0pt}}\isamarkupfalse%
\isacommand{assume}\isamarkupfalse%
\ ii{\isacharcolon}{\kern0pt}\ {\isachardoublequoteopen}nabla\ {\isasymturnstile}\ swap\ pi\ t{\isadigit{1}}\ {\isasymapprox}\ t{\isadigit{2}}{\isachardoublequoteclose}\isanewline
\ \ \ \ \isacommand{then}\isamarkupfalse%
\ \isacommand{show}\isamarkupfalse%
\ {\isachardoublequoteopen}nabla\ {\isasymturnstile}\ t{\isadigit{1}}\ {\isasymapprox}\ swap\ {\isacharparenleft}{\kern0pt}rev\ pi{\isacharparenright}{\kern0pt}\ t{\isadigit{2}}{\isachardoublequoteclose}\isanewline
\ \ \ \ \ \ \isacommand{using}\isamarkupfalse%
\ equ{\isacharunderscore}{\kern0pt}involutive{\isacharunderscore}{\kern0pt}left\ equ{\isacharunderscore}{\kern0pt}equivariance\ \isacommand{by}\isamarkupfalse%
\ blast\isacommand{{\isacharbraceright}{\kern0pt}}\isamarkupfalse%
\isanewline
\isacommand{qed}\isamarkupfalse%
%
\endisatagproof
{\isafoldproof}%
%
\isadelimproof
\isanewline
%
\endisadelimproof
\isanewline
\isanewline
\isacommand{lemma}\isamarkupfalse%
\ equ{\isacharunderscore}{\kern0pt}involutive{\isacharunderscore}{\kern0pt}right{\isacharcolon}{\kern0pt}\ \isanewline
\ \ {\isachardoublequoteopen}nabla\ {\isasymturnstile}\ t{\isadigit{1}}\ {\isasymapprox}\ swap\ {\isacharparenleft}{\kern0pt}rev\ pi{\isacharparenright}{\kern0pt}\ {\isacharparenleft}{\kern0pt}swap\ pi\ t{\isadigit{2}}{\isacharparenright}{\kern0pt}\ {\isacharequal}{\kern0pt}\ nabla\ {\isasymturnstile}\ t{\isadigit{1}}\ {\isasymapprox}\ t{\isadigit{2}}{\isachardoublequoteclose}\isanewline
%
\isadelimproof
\ \ %
\endisadelimproof
%
\isatagproof
\isacommand{using}\isamarkupfalse%
\ equ{\isacharunderscore}{\kern0pt}dec{\isacharunderscore}{\kern0pt}pi{\isacharbrackleft}{\kern0pt}of\ nabla\ pi\ t{\isadigit{1}}\ t{\isadigit{2}}{\isacharbrackright}{\kern0pt}\ equ{\isacharunderscore}{\kern0pt}equivariance{\isacharbrackleft}{\kern0pt}of\ nabla\ t{\isadigit{1}}\ t{\isadigit{2}}\ pi{\isacharbrackright}{\kern0pt}\ \isanewline
\ \ \ \ equ{\isacharunderscore}{\kern0pt}pi{\isacharunderscore}{\kern0pt}to{\isacharunderscore}{\kern0pt}right{\isacharbrackleft}{\kern0pt}of\ nabla\ t{\isadigit{1}}\ pi\ {\isacartoucheopen}swap\ pi\ t{\isadigit{2}}{\isacartoucheclose}{\isacharbrackright}{\kern0pt}\isanewline
\ \ \isacommand{by}\isamarkupfalse%
\ auto%
\endisatagproof
{\isafoldproof}%
%
\isadelimproof
\isanewline
%
\endisadelimproof
\isanewline
\isacommand{lemma}\isamarkupfalse%
\ equ{\isacharunderscore}{\kern0pt}pi{\isadigit{1}}{\isacharunderscore}{\kern0pt}pi{\isadigit{2}}{\isacharunderscore}{\kern0pt}add{\isacharcolon}{\kern0pt}\ \isanewline
\ \ {\isachardoublequoteopen}{\isacharparenleft}{\kern0pt}{\isasymforall}a{\isasymin}\ ds\ pi{\isadigit{1}}\ pi{\isadigit{2}}{\isachardot}{\kern0pt}\ nabla{\isasymturnstile}a{\isasymsharp}t{\isacharparenright}{\kern0pt}\ {\isasymLongrightarrow}\ {\isacharparenleft}{\kern0pt}nabla{\isasymturnstile}swap\ pi{\isadigit{1}}\ t\ {\isasymapprox}\ swap\ pi{\isadigit{2}}\ t{\isacharparenright}{\kern0pt}{\isachardoublequoteclose}\isanewline
%
\isadelimproof
%
\endisadelimproof
%
\isatagproof
\isacommand{proof}\isamarkupfalse%
{\isacharminus}{\kern0pt}\isanewline
\ \ \isacommand{assume}\isamarkupfalse%
\ {\isachardoublequoteopen}{\isacharparenleft}{\kern0pt}{\isasymforall}a{\isasymin}\ ds\ pi{\isadigit{1}}\ pi{\isadigit{2}}{\isachardot}{\kern0pt}\ nabla{\isasymturnstile}a{\isasymsharp}t{\isacharparenright}{\kern0pt}{\isachardoublequoteclose}\isanewline
\ \ \isacommand{hence}\isamarkupfalse%
\ {\isachardoublequoteopen}nabla\ {\isasymturnstile}\ t\ {\isasymapprox}\ swap\ {\isacharparenleft}{\kern0pt}rev\ pi{\isadigit{1}}\ {\isacharat}{\kern0pt}\ pi{\isadigit{2}}{\isacharparenright}{\kern0pt}\ t{\isachardoublequoteclose}\ \isanewline
\ \ \ \ \isacommand{using}\isamarkupfalse%
\ ds{\isacharunderscore}{\kern0pt}rev\ equ{\isacharunderscore}{\kern0pt}pi{\isacharunderscore}{\kern0pt}right\ \isacommand{by}\isamarkupfalse%
\ simp\isanewline
\ \ \isacommand{hence}\isamarkupfalse%
\ {\isachardoublequoteopen}nabla\ {\isasymturnstile}\ t\ {\isasymapprox}\ swap\ {\isacharparenleft}{\kern0pt}rev\ pi{\isadigit{1}}{\isacharparenright}{\kern0pt}\ {\isacharparenleft}{\kern0pt}swap\ pi{\isadigit{2}}\ t{\isacharparenright}{\kern0pt}{\isachardoublequoteclose}\isanewline
\ \ \ \ \isacommand{using}\isamarkupfalse%
\ swap{\isacharunderscore}{\kern0pt}append\ \isacommand{by}\isamarkupfalse%
\ auto\isanewline
\ \ \isacommand{then}\isamarkupfalse%
\ \isacommand{show}\isamarkupfalse%
\ {\isachardoublequoteopen}nabla{\isasymturnstile}swap\ pi{\isadigit{1}}\ t\ {\isasymapprox}\ swap\ pi{\isadigit{2}}\ t{\isachardoublequoteclose}\isanewline
\ \ \ \ \isacommand{using}\isamarkupfalse%
\ equ{\isacharunderscore}{\kern0pt}pi{\isacharunderscore}{\kern0pt}to{\isacharunderscore}{\kern0pt}right\ \isacommand{by}\isamarkupfalse%
\ simp\isanewline
\isacommand{qed}\isamarkupfalse%
%
\endisatagproof
{\isafoldproof}%
%
\isadelimproof
\isanewline
%
\endisadelimproof
\isanewline
\isanewline
\isanewline
\isacommand{lemma}\isamarkupfalse%
\ pi{\isacharunderscore}{\kern0pt}right{\isacharunderscore}{\kern0pt}equ{\isacharcolon}{\kern0pt}\ {\isachardoublequoteopen}{\isacharparenleft}{\kern0pt}nabla\ {\isasymturnstile}\ t\ {\isasymapprox}\ swap\ pi\ t{\isacharparenright}{\kern0pt}\ {\isasymLongrightarrow}\ {\isacharparenleft}{\kern0pt}{\isasymforall}a{\isasymin}\ ds\ {\isacharbrackleft}{\kern0pt}{\isacharbrackright}{\kern0pt}\ pi{\isachardot}{\kern0pt}\ nabla\ {\isasymturnstile}\ a\ {\isasymsharp}\ t{\isacharparenright}{\kern0pt}{\isachardoublequoteclose}\isanewline
%
\isadelimproof
\ \ %
\endisadelimproof
%
\isatagproof
\isacommand{using}\isamarkupfalse%
\ pi{\isacharunderscore}{\kern0pt}right{\isacharunderscore}{\kern0pt}equ{\isacharunderscore}{\kern0pt}help\ \isacommand{by}\isamarkupfalse%
\ blast%
\endisatagproof
{\isafoldproof}%
%
\isadelimproof
\isanewline
%
\endisadelimproof
\isanewline
\isanewline
\isacommand{lemma}\isamarkupfalse%
\ equ{\isacharunderscore}{\kern0pt}pi{\isadigit{1}}{\isacharunderscore}{\kern0pt}pi{\isadigit{2}}{\isacharunderscore}{\kern0pt}dec{\isacharcolon}{\kern0pt}\ \ \isanewline
\ \ {\isachardoublequoteopen}{\isacharparenleft}{\kern0pt}nabla\ {\isasymturnstile}\ swap\ pi{\isadigit{1}}\ t\ {\isasymapprox}\ swap\ pi{\isadigit{2}}\ t{\isacharparenright}{\kern0pt}\ {\isasymLongrightarrow}\ {\isacharparenleft}{\kern0pt}{\isasymforall}\ a\ {\isasymin}\ ds\ pi{\isadigit{1}}\ pi{\isadigit{2}}{\isachardot}{\kern0pt}\ nabla{\isasymturnstile}a\ {\isasymsharp}\ t{\isacharparenright}{\kern0pt}{\isachardoublequoteclose}\isanewline
%
\isadelimproof
%
\endisadelimproof
%
\isatagproof
\isacommand{proof}\isamarkupfalse%
{\isacharminus}{\kern0pt}\isanewline
\ \ \isacommand{assume}\isamarkupfalse%
\ assm{\isacharcolon}{\kern0pt}\ {\isachardoublequoteopen}nabla\ {\isasymturnstile}\ swap\ pi{\isadigit{1}}\ t\ {\isasymapprox}\ swap\ pi{\isadigit{2}}\ t{\isachardoublequoteclose}\isanewline
\ \ \isacommand{then}\isamarkupfalse%
\ \isacommand{have}\isamarkupfalse%
\ {\isachardoublequoteopen}nabla\ {\isasymturnstile}\ t\ {\isasymapprox}\ swap\ {\isacharparenleft}{\kern0pt}{\isacharparenleft}{\kern0pt}rev\ pi{\isadigit{1}}{\isacharparenright}{\kern0pt}\ {\isacharat}{\kern0pt}\ pi{\isadigit{2}}{\isacharparenright}{\kern0pt}\ t{\isachardoublequoteclose}\isanewline
\ \ \ \ \isacommand{using}\isamarkupfalse%
\ equ{\isacharunderscore}{\kern0pt}pi{\isacharunderscore}{\kern0pt}to{\isacharunderscore}{\kern0pt}right\ swap{\isacharunderscore}{\kern0pt}append\ \isacommand{by}\isamarkupfalse%
\ simp\isanewline
\ \ \isacommand{then}\isamarkupfalse%
\ \isacommand{show}\isamarkupfalse%
\ {\isachardoublequoteopen}{\isasymforall}\ a\ {\isasymin}\ ds\ pi{\isadigit{1}}\ pi{\isadigit{2}}{\isachardot}{\kern0pt}\ nabla{\isasymturnstile}a\ {\isasymsharp}\ t{\isachardoublequoteclose}\isanewline
\ \ \ \ \isacommand{using}\isamarkupfalse%
\ pi{\isacharunderscore}{\kern0pt}right{\isacharunderscore}{\kern0pt}equ\ ds{\isacharunderscore}{\kern0pt}rev{\isacharbrackleft}{\kern0pt}of\ pi{\isadigit{1}}\ pi{\isadigit{2}}{\isacharbrackright}{\kern0pt}\ \isacommand{by}\isamarkupfalse%
\ simp\isanewline
\isacommand{qed}\isamarkupfalse%
%
\endisatagproof
{\isafoldproof}%
%
\isadelimproof
\isanewline
%
\endisadelimproof
\isanewline
\isacommand{lemma}\isamarkupfalse%
\ equ{\isacharunderscore}{\kern0pt}weak{\isacharcolon}{\kern0pt}\ \isanewline
\ \ {\isachardoublequoteopen}nabla{\isadigit{1}}\ {\isasymturnstile}\ t{\isadigit{1}}\ {\isasymapprox}\ t{\isadigit{2}}\ {\isasymLongrightarrow}\ {\isacharparenleft}{\kern0pt}nabla{\isadigit{1}}\ {\isasymunion}\ nabla{\isadigit{2}}{\isacharparenright}{\kern0pt}\ {\isasymturnstile}\ t{\isadigit{1}}\ {\isasymapprox}\ t{\isadigit{2}}{\isachardoublequoteclose}\isanewline
%
\isadelimproof
\ \ %
\endisadelimproof
%
\isatagproof
\isacommand{by}\isamarkupfalse%
{\isacharparenleft}{\kern0pt}erule\ equ{\isachardot}{\kern0pt}induct{\isacharcomma}{\kern0pt}\ auto\ simp\ add{\isacharcolon}{\kern0pt}\ fresh{\isacharunderscore}{\kern0pt}weak{\isacharparenright}{\kern0pt}%
\endisatagproof
{\isafoldproof}%
%
\isadelimproof
\isanewline
%
\endisadelimproof
\isanewline
\isanewline
\isanewline
\isanewline
\isanewline
\isanewline
\isacommand{lemma}\isamarkupfalse%
\ psub{\isacharunderscore}{\kern0pt}trm{\isacharunderscore}{\kern0pt}not{\isacharunderscore}{\kern0pt}equ{\isacharcolon}{\kern0pt}\ \isanewline
\ \ {\isachardoublequoteopen}{\isasymforall}\ t{\isadigit{2}}\ {\isasymin}\ psub{\isacharunderscore}{\kern0pt}trms\ t{\isadigit{1}}{\isachardot}{\kern0pt}\ {\isacharparenleft}{\kern0pt}{\isasymnot}{\isacharparenleft}{\kern0pt}{\isasymexists}\ pi{\isachardot}{\kern0pt}\ {\isacharparenleft}{\kern0pt}nabla\ {\isasymturnstile}\ t{\isadigit{1}}\ {\isasymapprox}\ swap\ pi\ t{\isadigit{2}}{\isacharparenright}{\kern0pt}{\isacharparenright}{\kern0pt}{\isacharparenright}{\kern0pt}{\isachardoublequoteclose}\isanewline
%
\isadelimproof
%
\endisadelimproof
%
\isatagproof
\isacommand{proof}\isamarkupfalse%
\isanewline
\ \ \isacommand{fix}\isamarkupfalse%
\ t{\isadigit{2}}\isanewline
\ \ \isacommand{assume}\isamarkupfalse%
\ i{\isacharcolon}{\kern0pt}\ {\isachardoublequoteopen}t{\isadigit{2}}\ {\isasymin}\ psub{\isacharunderscore}{\kern0pt}trms\ t{\isadigit{1}}{\isachardoublequoteclose}\isanewline
\ \ \isacommand{show}\isamarkupfalse%
\ {\isachardoublequoteopen}{\isasymnot}\ {\isacharparenleft}{\kern0pt}{\isasymexists}pi{\isachardot}{\kern0pt}\ nabla\ {\isasymturnstile}\ t{\isadigit{1}}\ {\isasymapprox}\ swap\ pi\ t{\isadigit{2}}{\isacharparenright}{\kern0pt}{\isachardoublequoteclose}\isanewline
\ \ \isacommand{proof}\isamarkupfalse%
\isanewline
\ \ \ \ \isacommand{assume}\isamarkupfalse%
\ {\isachardoublequoteopen}{\isasymexists}pi{\isachardot}{\kern0pt}\ nabla\ {\isasymturnstile}\ t{\isadigit{1}}\ {\isasymapprox}\ swap\ pi\ t{\isadigit{2}}{\isachardoublequoteclose}\isanewline
\ \ \ \ \isacommand{then}\isamarkupfalse%
\ \isacommand{obtain}\isamarkupfalse%
\ pi\ \isakeyword{where}\ H{\isacharcolon}{\kern0pt}\isanewline
\ \ \ \ \ \ {\isachardoublequoteopen}nabla\ {\isasymturnstile}\ t{\isadigit{1}}\ {\isasymapprox}\ swap\ pi\ t{\isadigit{2}}{\isachardoublequoteclose}\ \isacommand{by}\isamarkupfalse%
\ blast\isanewline
\isanewline
\ \ \ \ \isacommand{from}\isamarkupfalse%
\ equ{\isacharunderscore}{\kern0pt}depth{\isacharbrackleft}{\kern0pt}OF\ H{\isacharbrackright}{\kern0pt}\isanewline
\ \ \ \ \isacommand{have}\isamarkupfalse%
\ {\isachardoublequoteopen}depth\ t{\isadigit{1}}\ {\isacharequal}{\kern0pt}\ depth\ {\isacharparenleft}{\kern0pt}swap\ pi\ t{\isadigit{2}}{\isacharparenright}{\kern0pt}{\isachardoublequoteclose}\isanewline
\ \ \ \ \ \ \isacommand{by}\isamarkupfalse%
\ simp\isanewline
\ \ \ \ \isacommand{hence}\isamarkupfalse%
\ {\isachardoublequoteopen}depth\ t{\isadigit{1}}\ {\isacharequal}{\kern0pt}\ depth\ t{\isadigit{2}}{\isachardoublequoteclose}\ \isanewline
\ \ \ \ \ \ \isacommand{by}\isamarkupfalse%
\ simp\isanewline
\ \ \ \ \isacommand{moreover}\isamarkupfalse%
\ \isacommand{have}\isamarkupfalse%
\ {\isachardoublequoteopen}depth\ t{\isadigit{2}}\ {\isacharless}{\kern0pt}\ depth\ t{\isadigit{1}}{\isachardoublequoteclose}\isanewline
\ \ \ \ \ \ \isacommand{using}\isamarkupfalse%
\ i\ depth{\isacharunderscore}{\kern0pt}psub{\isacharunderscore}{\kern0pt}trms\ \isanewline
\ \ \ \ \ \ \isacommand{by}\isamarkupfalse%
\ simp\isanewline
\isanewline
\ \ \ \ \isacommand{ultimately}\isamarkupfalse%
\ \isacommand{show}\isamarkupfalse%
\ False\ \isacommand{by}\isamarkupfalse%
\ auto\isanewline
\ \ \isacommand{qed}\isamarkupfalse%
\isanewline
\isacommand{qed}\isamarkupfalse%
%
\endisatagproof
{\isafoldproof}%
%
\isadelimproof
\isanewline
%
\endisadelimproof
%
\isadelimtheory
\isanewline
%
\endisadelimtheory
%
\isatagtheory
\isacommand{end}\isamarkupfalse%
%
\endisatagtheory
{\isafoldtheory}%
%
\isadelimtheory
%
\endisadelimtheory
%
\end{isabellebody}%
\endinput
%:%file=~/nominal_unification_Isabelle/nominal_alpha_unification/Equ.thy%:%
%:%10=1%:%
%:%11=1%:%
%:%12=2%:%
%:%13=3%:%
%:%14=4%:%
%:%15=5%:%
%:%20=5%:%
%:%23=6%:%
%:%24=7%:%
%:%25=7%:%
%:%26=8%:%
%:%29=9%:%
%:%33=9%:%
%:%34=9%:%
%:%39=9%:%
%:%42=10%:%
%:%43=11%:%
%:%44=12%:%
%:%45=12%:%
%:%46=13%:%
%:%53=14%:%
%:%54=14%:%
%:%55=15%:%
%:%56=15%:%
%:%57=16%:%
%:%58=17%:%
%:%59=17%:%
%:%60=17%:%
%:%61=18%:%
%:%62=18%:%
%:%63=19%:%
%:%64=19%:%
%:%65=19%:%
%:%66=20%:%
%:%67=20%:%
%:%68=20%:%
%:%69=21%:%
%:%70=21%:%
%:%71=21%:%
%:%72=21%:%
%:%73=21%:%
%:%74=22%:%
%:%80=22%:%
%:%83=23%:%
%:%84=24%:%
%:%85=25%:%
%:%86=25%:%
%:%87=26%:%
%:%94=27%:%
%:%95=27%:%
%:%96=28%:%
%:%97=28%:%
%:%98=29%:%
%:%99=29%:%
%:%100=29%:%
%:%101=30%:%
%:%102=30%:%
%:%103=31%:%
%:%104=31%:%
%:%105=31%:%
%:%106=32%:%
%:%107=32%:%
%:%108=33%:%
%:%109=33%:%
%:%110=33%:%
%:%111=34%:%
%:%117=34%:%
%:%120=35%:%
%:%121=36%:%
%:%122=37%:%
%:%123=37%:%
%:%124=38%:%
%:%131=39%:%
%:%132=39%:%
%:%133=40%:%
%:%134=41%:%
%:%135=41%:%
%:%136=41%:%
%:%137=42%:%
%:%138=42%:%
%:%139=43%:%
%:%140=43%:%
%:%141=43%:%
%:%142=44%:%
%:%143=44%:%
%:%144=44%:%
%:%145=45%:%
%:%146=45%:%
%:%147=46%:%
%:%148=46%:%
%:%149=46%:%
%:%150=47%:%
%:%151=48%:%
%:%152=48%:%
%:%153=48%:%
%:%154=49%:%
%:%155=49%:%
%:%156=50%:%
%:%157=50%:%
%:%158=50%:%
%:%159=51%:%
%:%160=51%:%
%:%161=51%:%
%:%162=52%:%
%:%163=52%:%
%:%164=52%:%
%:%165=53%:%
%:%166=53%:%
%:%167=53%:%
%:%168=54%:%
%:%169=54%:%
%:%170=54%:%
%:%171=55%:%
%:%172=55%:%
%:%173=55%:%
%:%174=56%:%
%:%175=56%:%
%:%176=56%:%
%:%177=56%:%
%:%178=57%:%
%:%179=58%:%
%:%185=58%:%
%:%188=59%:%
%:%189=60%:%
%:%190=61%:%
%:%191=61%:%
%:%192=62%:%
%:%199=63%:%
%:%200=63%:%
%:%201=64%:%
%:%202=64%:%
%:%203=64%:%
%:%204=65%:%
%:%205=65%:%
%:%206=65%:%
%:%207=66%:%
%:%208=66%:%
%:%209=66%:%
%:%210=66%:%
%:%211=67%:%
%:%212=67%:%
%:%213=67%:%
%:%214=68%:%
%:%215=68%:%
%:%216=68%:%
%:%217=69%:%
%:%218=69%:%
%:%219=69%:%
%:%220=69%:%
%:%221=70%:%
%:%227=70%:%
%:%230=71%:%
%:%231=72%:%
%:%232=73%:%
%:%233=73%:%
%:%234=74%:%
%:%237=75%:%
%:%241=75%:%
%:%242=75%:%
%:%243=76%:%
%:%244=77%:%
%:%245=77%:%
%:%250=77%:%
%:%253=78%:%
%:%254=79%:%
%:%255=79%:%
%:%256=80%:%
%:%263=81%:%
%:%264=81%:%
%:%265=82%:%
%:%266=82%:%
%:%267=83%:%
%:%268=83%:%
%:%269=84%:%
%:%270=84%:%
%:%271=84%:%
%:%272=85%:%
%:%273=85%:%
%:%274=86%:%
%:%275=86%:%
%:%276=86%:%
%:%277=87%:%
%:%278=87%:%
%:%279=87%:%
%:%280=88%:%
%:%281=88%:%
%:%282=88%:%
%:%283=89%:%
%:%289=89%:%
%:%292=90%:%
%:%293=91%:%
%:%294=92%:%
%:%295=93%:%
%:%296=93%:%
%:%299=94%:%
%:%303=94%:%
%:%304=94%:%
%:%305=94%:%
%:%310=94%:%
%:%313=95%:%
%:%314=96%:%
%:%315=97%:%
%:%316=97%:%
%:%317=98%:%
%:%324=99%:%
%:%325=99%:%
%:%326=100%:%
%:%327=100%:%
%:%328=101%:%
%:%329=101%:%
%:%330=101%:%
%:%331=102%:%
%:%332=102%:%
%:%333=102%:%
%:%334=103%:%
%:%335=103%:%
%:%336=103%:%
%:%337=104%:%
%:%338=104%:%
%:%339=104%:%
%:%340=105%:%
%:%346=105%:%
%:%349=106%:%
%:%350=107%:%
%:%351=107%:%
%:%352=108%:%
%:%355=109%:%
%:%359=109%:%
%:%360=109%:%
%:%365=109%:%
%:%368=110%:%
%:%369=111%:%
%:%370=112%:%
%:%371=113%:%
%:%372=114%:%
%:%373=115%:%
%:%374=116%:%
%:%375=116%:%
%:%376=117%:%
%:%383=118%:%
%:%384=118%:%
%:%385=119%:%
%:%386=119%:%
%:%387=120%:%
%:%388=120%:%
%:%389=121%:%
%:%390=121%:%
%:%391=122%:%
%:%392=122%:%
%:%393=123%:%
%:%394=123%:%
%:%395=124%:%
%:%396=124%:%
%:%397=124%:%
%:%398=125%:%
%:%399=125%:%
%:%400=126%:%
%:%401=127%:%
%:%402=127%:%
%:%403=128%:%
%:%404=128%:%
%:%405=129%:%
%:%406=129%:%
%:%407=130%:%
%:%408=130%:%
%:%409=131%:%
%:%410=131%:%
%:%411=132%:%
%:%412=132%:%
%:%413=132%:%
%:%414=133%:%
%:%415=133%:%
%:%416=134%:%
%:%417=134%:%
%:%418=135%:%
%:%419=136%:%
%:%420=136%:%
%:%421=136%:%
%:%422=136%:%
%:%423=137%:%
%:%424=137%:%
%:%425=138%:%
%:%431=138%:%
%:%436=139%:%
%:%441=140%:%

%
\begin{isabellebody}%
\setisabellecontext{Substs}%
%
\isadelimtheory
%
\endisadelimtheory
%
\isatagtheory
\isacommand{theory}\isamarkupfalse%
\ Substs\isanewline
\isanewline
\isakeyword{imports}\ Equ\isanewline
\isanewline
\isakeyword{begin}%
\endisatagtheory
{\isafoldtheory}%
%
\isadelimtheory
\isanewline
%
\endisadelimtheory
\isanewline
\isanewline
\isanewline
\isacommand{type{\isacharunderscore}{\kern0pt}synonym}\isamarkupfalse%
\ substs\ {\isacharequal}{\kern0pt}\ {\isachardoublequoteopen}{\isacharparenleft}{\kern0pt}string\ {\isasymtimes}\ trm{\isacharparenright}{\kern0pt}\ list{\isachardoublequoteclose}\isanewline
\isanewline
\isacommand{fun}\isamarkupfalse%
\ look{\isacharunderscore}{\kern0pt}up\ {\isacharcolon}{\kern0pt}{\isacharcolon}{\kern0pt}\ {\isachardoublequoteopen}string\ {\isasymRightarrow}\ substs\ {\isasymRightarrow}\ trm{\isachardoublequoteclose}\ \isakeyword{where}\isanewline
\ \ {\isachardoublequoteopen}look{\isacharunderscore}{\kern0pt}up\ X\ {\isacharbrackleft}{\kern0pt}{\isacharbrackright}{\kern0pt}\ \ \ \ \ {\isacharequal}{\kern0pt}\ Susp\ {\isacharbrackleft}{\kern0pt}{\isacharbrackright}{\kern0pt}\ X{\isachardoublequoteclose}\ {\isacharbar}{\kern0pt}\isanewline
\ \ {\isachardoublequoteopen}look{\isacharunderscore}{\kern0pt}up\ X\ {\isacharparenleft}{\kern0pt}x{\isacharhash}{\kern0pt}xs{\isacharparenright}{\kern0pt}\ {\isacharequal}{\kern0pt}\ {\isacharparenleft}{\kern0pt}if\ {\isacharparenleft}{\kern0pt}X\ {\isacharequal}{\kern0pt}\ fst\ x{\isacharparenright}{\kern0pt}\ then\ {\isacharparenleft}{\kern0pt}snd\ x{\isacharparenright}{\kern0pt}\ else\ look{\isacharunderscore}{\kern0pt}up\ X\ xs{\isacharparenright}{\kern0pt}{\isachardoublequoteclose}\isanewline
\isanewline
\isacommand{fun}\isamarkupfalse%
\ subst\ {\isacharcolon}{\kern0pt}{\isacharcolon}{\kern0pt}\ {\isachardoublequoteopen}substs\ {\isasymRightarrow}\ trm\ {\isasymRightarrow}\ trm{\isachardoublequoteclose}\ \isakeyword{where}\isanewline
\ \ subst{\isacharunderscore}{\kern0pt}unit{\isacharcolon}{\kern0pt}\ {\isachardoublequoteopen}subst\ s\ {\isacharparenleft}{\kern0pt}Unit{\isacharparenright}{\kern0pt}\ \ \ \ \ \ \ \ \ \ {\isacharequal}{\kern0pt}\ Unit{\isachardoublequoteclose}\ {\isacharbar}{\kern0pt}\isanewline
\ \ subst{\isacharunderscore}{\kern0pt}susp{\isacharcolon}{\kern0pt}\ {\isachardoublequoteopen}subst\ s\ {\isacharparenleft}{\kern0pt}Susp\ pi\ X{\isacharparenright}{\kern0pt}\ \ \ \ \ {\isacharequal}{\kern0pt}\ swap\ pi\ {\isacharparenleft}{\kern0pt}look{\isacharunderscore}{\kern0pt}up\ X\ s{\isacharparenright}{\kern0pt}{\isachardoublequoteclose}\ {\isacharbar}{\kern0pt}\isanewline
\ \ subst{\isacharunderscore}{\kern0pt}atom{\isacharcolon}{\kern0pt}\ {\isachardoublequoteopen}subst\ s\ {\isacharparenleft}{\kern0pt}Atom\ a{\isacharparenright}{\kern0pt}\ \ \ \ \ \ \ \ {\isacharequal}{\kern0pt}\ Atom\ a{\isachardoublequoteclose}\ {\isacharbar}{\kern0pt}\isanewline
\ \ subst{\isacharunderscore}{\kern0pt}abst{\isacharcolon}{\kern0pt}\ {\isachardoublequoteopen}subst\ s\ {\isacharparenleft}{\kern0pt}Abst\ a\ t{\isacharparenright}{\kern0pt}\ \ \ \ \ \ {\isacharequal}{\kern0pt}\ Abst\ a\ {\isacharparenleft}{\kern0pt}subst\ s\ t{\isacharparenright}{\kern0pt}{\isachardoublequoteclose}\ {\isacharbar}{\kern0pt}\isanewline
\ \ subst{\isacharunderscore}{\kern0pt}paar{\isacharcolon}{\kern0pt}\ {\isachardoublequoteopen}subst\ s\ {\isacharparenleft}{\kern0pt}Paar\ t{\isadigit{1}}\ t{\isadigit{2}}{\isacharparenright}{\kern0pt}\ \ \ \ {\isacharequal}{\kern0pt}\ Paar\ {\isacharparenleft}{\kern0pt}subst\ s\ t{\isadigit{1}}{\isacharparenright}{\kern0pt}\ {\isacharparenleft}{\kern0pt}subst\ s\ t{\isadigit{2}}{\isacharparenright}{\kern0pt}{\isachardoublequoteclose}\ {\isacharbar}{\kern0pt}\isanewline
\ \ subst{\isacharunderscore}{\kern0pt}func{\isacharcolon}{\kern0pt}\ {\isachardoublequoteopen}subst\ s\ {\isacharparenleft}{\kern0pt}Func\ F\ t{\isacharparenright}{\kern0pt}\ \ \ \ \ \ {\isacharequal}{\kern0pt}\ Func\ F\ {\isacharparenleft}{\kern0pt}subst\ s\ t{\isacharparenright}{\kern0pt}{\isachardoublequoteclose}\isanewline
\isanewline
\isacommand{declare}\isamarkupfalse%
\ subst{\isacharunderscore}{\kern0pt}susp\ {\isacharbrackleft}{\kern0pt}simp\ del{\isacharbrackright}{\kern0pt}\isanewline
\isanewline
\isanewline
\isanewline
\isacommand{fun}\isamarkupfalse%
\ alist{\isacharunderscore}{\kern0pt}rec\ {\isacharcolon}{\kern0pt}{\isacharcolon}{\kern0pt}\ {\isachardoublequoteopen}substs\ {\isasymRightarrow}\ substs\ {\isasymRightarrow}\ {\isacharparenleft}{\kern0pt}string{\isasymRightarrow}trm{\isasymRightarrow}substs{\isasymRightarrow}substs{\isasymRightarrow}substs{\isacharparenright}{\kern0pt}\ {\isasymRightarrow}\ substs{\isachardoublequoteclose}\ \isakeyword{where}\isanewline
\ \ {\isachardoublequoteopen}alist{\isacharunderscore}{\kern0pt}rec\ {\isacharbrackleft}{\kern0pt}{\isacharbrackright}{\kern0pt}\ c\ d\ {\isacharequal}{\kern0pt}\ c{\isachardoublequoteclose}\ {\isacharbar}{\kern0pt}\isanewline
\ \ {\isachardoublequoteopen}alist{\isacharunderscore}{\kern0pt}rec\ {\isacharparenleft}{\kern0pt}p\ {\isacharhash}{\kern0pt}\ al{\isacharparenright}{\kern0pt}\ c\ d\ {\isacharequal}{\kern0pt}\ d\ {\isacharparenleft}{\kern0pt}fst\ p{\isacharparenright}{\kern0pt}\ {\isacharparenleft}{\kern0pt}snd\ p{\isacharparenright}{\kern0pt}\ al\ {\isacharparenleft}{\kern0pt}alist{\isacharunderscore}{\kern0pt}rec\ al\ c\ d{\isacharparenright}{\kern0pt}{\isachardoublequoteclose}\isanewline
\isanewline
\isacommand{definition}\isamarkupfalse%
\ comp\ {\isacharcolon}{\kern0pt}{\isacharcolon}{\kern0pt}\ \ {\isachardoublequoteopen}substs\ {\isasymRightarrow}\ substs\ {\isasymRightarrow}\ substs{\isachardoublequoteclose}\ {\isacharparenleft}{\kern0pt}\isakeyword{infixl}\ {\isacartoucheopen}{\isasymbullet}{\isacartoucheclose}\ {\isadigit{8}}{\isadigit{1}}{\isacharparenright}{\kern0pt}\ \isakeyword{where}\isanewline
\ {\isachardoublequoteopen}s{\isadigit{1}}\ {\isasymbullet}\ s{\isadigit{2}}\ {\isasymequiv}\ alist{\isacharunderscore}{\kern0pt}rec\ s{\isadigit{2}}\ s{\isadigit{1}}\ {\isacharparenleft}{\kern0pt}{\isasymlambda}\ x\ y\ xs\ g{\isachardot}{\kern0pt}\ {\isacharparenleft}{\kern0pt}x{\isacharcomma}{\kern0pt}subst\ s{\isadigit{1}}\ y{\isacharparenright}{\kern0pt}\ {\isacharhash}{\kern0pt}\ g{\isacharparenright}{\kern0pt}{\isachardoublequoteclose}\isanewline
\isanewline
\isanewline
\isanewline
\isacommand{definition}\isamarkupfalse%
\ domn\ {\isacharcolon}{\kern0pt}{\isacharcolon}{\kern0pt}\ {\isachardoublequoteopen}{\isacharparenleft}{\kern0pt}trm\ {\isasymRightarrow}\ trm{\isacharparenright}{\kern0pt}\ {\isasymRightarrow}\ string\ set{\isachardoublequoteclose}\ \isakeyword{where}\isanewline
\ \ {\isachardoublequoteopen}domn\ s\ {\isasymequiv}\ {\isacharbraceleft}{\kern0pt}X{\isachardot}{\kern0pt}\ {\isacharparenleft}{\kern0pt}s\ {\isacharparenleft}{\kern0pt}Susp\ {\isacharbrackleft}{\kern0pt}{\isacharbrackright}{\kern0pt}\ X{\isacharparenright}{\kern0pt}{\isacharparenright}{\kern0pt}\ {\isasymnoteq}\ {\isacharparenleft}{\kern0pt}Susp\ {\isacharbrackleft}{\kern0pt}{\isacharbrackright}{\kern0pt}\ X{\isacharparenright}{\kern0pt}{\isacharbraceright}{\kern0pt}{\isachardoublequoteclose}\ \isanewline
\isanewline
\isanewline
\isanewline
\isacommand{definition}\isamarkupfalse%
\ ext{\isacharunderscore}{\kern0pt}subst\ {\isacharcolon}{\kern0pt}{\isacharcolon}{\kern0pt}\ {\isachardoublequoteopen}fresh{\isacharunderscore}{\kern0pt}envs\ {\isasymRightarrow}\ {\isacharparenleft}{\kern0pt}trm\ {\isasymRightarrow}\ trm{\isacharparenright}{\kern0pt}\ {\isasymRightarrow}\ fresh{\isacharunderscore}{\kern0pt}envs\ {\isasymRightarrow}\ bool{\isachardoublequoteclose}\ {\isacharparenleft}{\kern0pt}{\isachardoublequoteopen}\ {\isacharunderscore}{\kern0pt}\ {\isasymTurnstile}\ {\isacharunderscore}{\kern0pt}\ {\isacharunderscore}{\kern0pt}\ {\isachardoublequoteclose}\ {\isacharbrackleft}{\kern0pt}{\isadigit{8}}{\isadigit{0}}{\isacharcomma}{\kern0pt}{\isadigit{8}}{\isadigit{0}}{\isacharcomma}{\kern0pt}{\isadigit{8}}{\isadigit{0}}{\isacharbrackright}{\kern0pt}\ {\isadigit{8}}{\isadigit{0}}{\isacharparenright}{\kern0pt}\ \isakeyword{where}\isanewline
\ \ {\isachardoublequoteopen}nabla{\isadigit{1}}\ {\isasymTurnstile}\ s\ nabla{\isadigit{2}}\ {\isasymequiv}\ {\isacharparenleft}{\kern0pt}{\isasymforall}{\isacharparenleft}{\kern0pt}a{\isacharcomma}{\kern0pt}X{\isacharparenright}{\kern0pt}{\isasymin}nabla{\isadigit{2}}{\isachardot}{\kern0pt}\ nabla{\isadigit{1}}{\isasymturnstile}a{\isasymsharp}\ s\ {\isacharparenleft}{\kern0pt}Susp\ {\isacharbrackleft}{\kern0pt}{\isacharbrackright}{\kern0pt}\ X{\isacharparenright}{\kern0pt}{\isacharparenright}{\kern0pt}{\isachardoublequoteclose}\isanewline
\isanewline
\isanewline
\isanewline
\isacommand{definition}\isamarkupfalse%
\ subst{\isacharunderscore}{\kern0pt}equ\ {\isacharcolon}{\kern0pt}{\isacharcolon}{\kern0pt}\ {\isachardoublequoteopen}fresh{\isacharunderscore}{\kern0pt}envs\ {\isasymRightarrow}\ {\isacharparenleft}{\kern0pt}trm{\isasymRightarrow}trm{\isacharparenright}{\kern0pt}\ {\isasymRightarrow}\ {\isacharparenleft}{\kern0pt}trm{\isasymRightarrow}trm{\isacharparenright}{\kern0pt}\ {\isasymRightarrow}\ bool{\isachardoublequoteclose}\ {\isacharparenleft}{\kern0pt}{\isachardoublequoteopen}\ {\isacharunderscore}{\kern0pt}\ {\isasymTurnstile}\ {\isacharunderscore}{\kern0pt}\ {\isasymapprox}\ {\isacharunderscore}{\kern0pt}{\isachardoublequoteclose}\ {\isacharbrackleft}{\kern0pt}{\isadigit{8}}{\isadigit{0}}{\isacharcomma}{\kern0pt}{\isadigit{8}}{\isadigit{0}}{\isacharcomma}{\kern0pt}{\isadigit{8}}{\isadigit{0}}{\isacharbrackright}{\kern0pt}\ {\isadigit{8}}{\isadigit{0}}{\isacharparenright}{\kern0pt}\isanewline
\ \ \isakeyword{where}\isanewline
\ \ {\isachardoublequoteopen}nabla\ {\isasymTurnstile}\ s{\isadigit{1}}\ {\isasymapprox}\ s{\isadigit{2}}\ {\isasymequiv}\ \ {\isasymforall}X{\isasymin}{\isacharparenleft}{\kern0pt}domn\ s{\isadigit{1}}{\isasymunion}domn\ s{\isadigit{2}}{\isacharparenright}{\kern0pt}{\isachardot}{\kern0pt}\ {\isacharparenleft}{\kern0pt}nabla\ {\isasymturnstile}\ s{\isadigit{1}}\ {\isacharparenleft}{\kern0pt}Susp\ {\isacharbrackleft}{\kern0pt}{\isacharbrackright}{\kern0pt}\ X{\isacharparenright}{\kern0pt}\ {\isasymapprox}\ s{\isadigit{2}}\ {\isacharparenleft}{\kern0pt}Susp\ {\isacharbrackleft}{\kern0pt}{\isacharbrackright}{\kern0pt}\ X{\isacharparenright}{\kern0pt}{\isacharparenright}{\kern0pt}{\isachardoublequoteclose}\isanewline
\isanewline
\isacommand{lemma}\isamarkupfalse%
\ subst{\isacharunderscore}{\kern0pt}equ{\isacharunderscore}{\kern0pt}symm{\isacharcolon}{\kern0pt}\ \isanewline
\ \ \isakeyword{assumes}\ {\isachardoublequoteopen}nabla\ {\isasymTurnstile}\ s{\isadigit{1}}\ {\isasymapprox}\ s{\isadigit{2}}{\isachardoublequoteclose}\isanewline
\ \ \isakeyword{shows}\ {\isachardoublequoteopen}nabla\ {\isasymTurnstile}\ s{\isadigit{2}}\ {\isasymapprox}\ s{\isadigit{1}}{\isachardoublequoteclose}\isanewline
%
\isadelimproof
\ \ %
\endisadelimproof
%
\isatagproof
\isacommand{using}\isamarkupfalse%
\ assms\ equ{\isacharunderscore}{\kern0pt}symm\ subst{\isacharunderscore}{\kern0pt}equ{\isacharunderscore}{\kern0pt}def{\isacharbrackleft}{\kern0pt}of\ nabla\ s{\isadigit{1}}\ s{\isadigit{2}}{\isacharbrackright}{\kern0pt}\isanewline
\ \ \ \ subst{\isacharunderscore}{\kern0pt}equ{\isacharunderscore}{\kern0pt}def{\isacharbrackleft}{\kern0pt}of\ nabla\ s{\isadigit{2}}\ s{\isadigit{1}}{\isacharbrackright}{\kern0pt}\ Un{\isacharunderscore}{\kern0pt}commute{\isacharbrackleft}{\kern0pt}of\ {\isacartoucheopen}domn\ s{\isadigit{1}}{\isacartoucheclose}\ {\isacartoucheopen}domn\ s{\isadigit{2}}{\isacartoucheclose}{\isacharbrackright}{\kern0pt}\ \isacommand{by}\isamarkupfalse%
\ blast%
\endisatagproof
{\isafoldproof}%
%
\isadelimproof
\isanewline
%
\endisadelimproof
\isanewline
\isacommand{lemma}\isamarkupfalse%
\ subst{\isacharunderscore}{\kern0pt}equ{\isacharunderscore}{\kern0pt}refl{\isacharcolon}{\kern0pt}\isanewline
\ \ \isakeyword{shows}\ {\isachardoublequoteopen}nabla\ {\isasymTurnstile}\ s\ {\isasymapprox}\ s{\isachardoublequoteclose}\isanewline
%
\isadelimproof
\ \ %
\endisadelimproof
%
\isatagproof
\isacommand{using}\isamarkupfalse%
\ equ{\isacharunderscore}{\kern0pt}refl\ subst{\isacharunderscore}{\kern0pt}equ{\isacharunderscore}{\kern0pt}def\ Un{\isacharunderscore}{\kern0pt}absorb\ \isacommand{by}\isamarkupfalse%
\ simp%
\endisatagproof
{\isafoldproof}%
%
\isadelimproof
\isanewline
%
\endisadelimproof
\isanewline
\isacommand{lemma}\isamarkupfalse%
\ not{\isacharunderscore}{\kern0pt}in{\isacharunderscore}{\kern0pt}domn{\isacharcolon}{\kern0pt}\ {\isachardoublequoteopen}X\ {\isasymnotin}\ {\isacharparenleft}{\kern0pt}domn\ s{\isacharparenright}{\kern0pt}{\isasymLongrightarrow}\ {\isacharparenleft}{\kern0pt}s\ {\isacharparenleft}{\kern0pt}Susp\ {\isacharbrackleft}{\kern0pt}{\isacharbrackright}{\kern0pt}\ X{\isacharparenright}{\kern0pt}{\isacharparenright}{\kern0pt}\ {\isacharequal}{\kern0pt}\ {\isacharparenleft}{\kern0pt}Susp\ {\isacharbrackleft}{\kern0pt}{\isacharbrackright}{\kern0pt}\ X{\isacharparenright}{\kern0pt}{\isachardoublequoteclose}\isanewline
%
\isadelimproof
\ \ %
\endisadelimproof
%
\isatagproof
\isacommand{using}\isamarkupfalse%
\ domn{\isacharunderscore}{\kern0pt}def\ \isacommand{by}\isamarkupfalse%
\ simp%
\endisatagproof
{\isafoldproof}%
%
\isadelimproof
\isanewline
%
\endisadelimproof
\isanewline
\isacommand{lemma}\isamarkupfalse%
\ subst{\isacharunderscore}{\kern0pt}swap{\isacharunderscore}{\kern0pt}comm{\isacharcolon}{\kern0pt}\ {\isachardoublequoteopen}subst\ s\ {\isacharparenleft}{\kern0pt}swap\ pi\ t{\isacharparenright}{\kern0pt}\ {\isacharequal}{\kern0pt}\ swap\ pi\ {\isacharparenleft}{\kern0pt}subst\ s\ t{\isacharparenright}{\kern0pt}{\isachardoublequoteclose}\isanewline
%
\isadelimproof
\ \ %
\endisadelimproof
%
\isatagproof
\isacommand{using}\isamarkupfalse%
\ swap{\isacharunderscore}{\kern0pt}append\ subst{\isacharunderscore}{\kern0pt}susp\ \isacommand{by}\isamarkupfalse%
\ {\isacharparenleft}{\kern0pt}induct\ t{\isacharparenright}{\kern0pt}\ auto%
\endisatagproof
{\isafoldproof}%
%
\isadelimproof
\isanewline
%
\endisadelimproof
\isanewline
\isacommand{lemma}\isamarkupfalse%
\ subst{\isacharunderscore}{\kern0pt}not{\isacharunderscore}{\kern0pt}occurs{\isacharcolon}{\kern0pt}\ \isanewline
\ {\isachardoublequoteopen}{\isasymnot}{\isacharparenleft}{\kern0pt}occurs\ X\ t{\isacharparenright}{\kern0pt}\ {\isasymLongrightarrow}\ subst\ {\isacharbrackleft}{\kern0pt}{\isacharparenleft}{\kern0pt}X{\isacharcomma}{\kern0pt}t{\isadigit{2}}{\isacharparenright}{\kern0pt}{\isacharbrackright}{\kern0pt}\ t\ {\isacharequal}{\kern0pt}\ t{\isachardoublequoteclose}\isanewline
%
\isadelimproof
%
\endisadelimproof
%
\isatagproof
\isacommand{proof}\isamarkupfalse%
{\isacharminus}{\kern0pt}\isanewline
\ \ \isacommand{have}\isamarkupfalse%
\ i{\isacharcolon}{\kern0pt}\ {\isachardoublequoteopen}{\isasymnot}{\isacharparenleft}{\kern0pt}occurs\ X\ t{\isacharparenright}{\kern0pt}\ {\isasymlongrightarrow}\ subst\ {\isacharbrackleft}{\kern0pt}{\isacharparenleft}{\kern0pt}X{\isacharcomma}{\kern0pt}t{\isadigit{2}}{\isacharparenright}{\kern0pt}{\isacharbrackright}{\kern0pt}\ t\ {\isacharequal}{\kern0pt}\ t{\isachardoublequoteclose}\isanewline
\ \ \ \ \isacommand{using}\isamarkupfalse%
\ subst{\isacharunderscore}{\kern0pt}susp\ \isacommand{by}\isamarkupfalse%
\ {\isacharparenleft}{\kern0pt}induct\ t{\isacharparenright}{\kern0pt}\ auto\isanewline
\ \ \isacommand{assume}\isamarkupfalse%
\ {\isachardoublequoteopen}{\isasymnot}{\isacharparenleft}{\kern0pt}occurs\ X\ t{\isacharparenright}{\kern0pt}{\isachardoublequoteclose}\isanewline
\ \ \isacommand{with}\isamarkupfalse%
\ i\ \isacommand{show}\isamarkupfalse%
\ {\isacharquery}{\kern0pt}thesis\ \isacommand{by}\isamarkupfalse%
\ simp\isanewline
\isacommand{qed}\isamarkupfalse%
%
\endisatagproof
{\isafoldproof}%
%
\isadelimproof
\isanewline
%
\endisadelimproof
\isanewline
\isanewline
\isacommand{lemma}\isamarkupfalse%
\ id{\isacharunderscore}{\kern0pt}subst\ {\isacharbrackleft}{\kern0pt}simp{\isacharbrackright}{\kern0pt}{\isacharcolon}{\kern0pt}\ {\isachardoublequoteopen}subst\ {\isacharbrackleft}{\kern0pt}{\isacharbrackright}{\kern0pt}\ t\ {\isacharequal}{\kern0pt}\ t{\isachardoublequoteclose}\isanewline
%
\isadelimproof
\ \ %
\endisadelimproof
%
\isatagproof
\isacommand{using}\isamarkupfalse%
\ subst{\isacharunderscore}{\kern0pt}susp\ \isacommand{by}\isamarkupfalse%
\ {\isacharparenleft}{\kern0pt}induct\ t{\isacharparenright}{\kern0pt}\ auto%
\endisatagproof
{\isafoldproof}%
%
\isadelimproof
\isanewline
%
\endisadelimproof
\isanewline
\isacommand{lemma}\isamarkupfalse%
\ subst{\isacharunderscore}{\kern0pt}comp{\isacharunderscore}{\kern0pt}id\ {\isacharbrackleft}{\kern0pt}simp{\isacharbrackright}{\kern0pt}{\isacharcolon}{\kern0pt}\ {\isachardoublequoteopen}subst\ {\isacharparenleft}{\kern0pt}s\ {\isasymbullet}\ {\isacharbrackleft}{\kern0pt}{\isacharbrackright}{\kern0pt}{\isacharparenright}{\kern0pt}\ {\isacharequal}{\kern0pt}\ subst\ s{\isachardoublequoteclose}\isanewline
%
\isadelimproof
\ \ %
\endisadelimproof
%
\isatagproof
\isacommand{using}\isamarkupfalse%
\ comp{\isacharunderscore}{\kern0pt}def\ \isacommand{by}\isamarkupfalse%
\ simp%
\endisatagproof
{\isafoldproof}%
%
\isadelimproof
\isanewline
%
\endisadelimproof
\isanewline
\isacommand{lemma}\isamarkupfalse%
\ id{\isacharunderscore}{\kern0pt}comp{\isacharunderscore}{\kern0pt}subst\ {\isacharbrackleft}{\kern0pt}simp{\isacharbrackright}{\kern0pt}{\isacharcolon}{\kern0pt}\ {\isachardoublequoteopen}subst\ {\isacharparenleft}{\kern0pt}{\isacharbrackleft}{\kern0pt}{\isacharbrackright}{\kern0pt}\ {\isasymbullet}\ s{\isacharparenright}{\kern0pt}\ {\isacharequal}{\kern0pt}\ subst\ s{\isachardoublequoteclose}\isanewline
%
\isadelimproof
%
\endisadelimproof
%
\isatagproof
\isacommand{proof}\isamarkupfalse%
{\isacharparenleft}{\kern0pt}rule\ ext{\isacharparenright}{\kern0pt}\isanewline
\ \ \isacommand{fix}\isamarkupfalse%
\ t\isanewline
\ \ \isacommand{show}\isamarkupfalse%
\ {\isachardoublequoteopen}subst\ {\isacharparenleft}{\kern0pt}{\isacharbrackleft}{\kern0pt}{\isacharbrackright}{\kern0pt}\ {\isasymbullet}\ s{\isacharparenright}{\kern0pt}\ t\ {\isacharequal}{\kern0pt}\ subst\ s\ t{\isachardoublequoteclose}\isanewline
\ \ \isacommand{proof}\isamarkupfalse%
{\isacharparenleft}{\kern0pt}induction\ t\ arbitrary{\isacharcolon}{\kern0pt}\ s{\isacharparenright}{\kern0pt}\isanewline
\ \ \ \ \isacommand{case}\isamarkupfalse%
\ {\isacharparenleft}{\kern0pt}Susp\ pi\ X{\isacharparenright}{\kern0pt}\isanewline
\ \ \ \ \isacommand{then}\isamarkupfalse%
\ \isacommand{show}\isamarkupfalse%
\ {\isacharquery}{\kern0pt}case\ \isacommand{using}\isamarkupfalse%
\ comp{\isacharunderscore}{\kern0pt}def\ subst{\isacharunderscore}{\kern0pt}susp\isanewline
\ \ \ \ \isacommand{proof}\isamarkupfalse%
\ {\isacharparenleft}{\kern0pt}induction\ s{\isacharparenright}{\kern0pt}\isanewline
\ \ \ \ \ \ \isacommand{case}\isamarkupfalse%
\ Nil\isanewline
\ \ \ \ \ \ \isacommand{then}\isamarkupfalse%
\ \isacommand{show}\isamarkupfalse%
\ {\isacharquery}{\kern0pt}case\isanewline
\ \ \ \ \ \ \ \ \isacommand{using}\isamarkupfalse%
\ comp{\isacharunderscore}{\kern0pt}def\ subst{\isacharunderscore}{\kern0pt}susp\ \isacommand{by}\isamarkupfalse%
\ simp\isanewline
\ \ \ \ \isacommand{next}\isamarkupfalse%
\isanewline
\ \ \ \ \ \ \isacommand{case}\isamarkupfalse%
\ {\isacharparenleft}{\kern0pt}Cons\ a\ s{\isacharparenright}{\kern0pt}\isanewline
\ \ \ \ \ \ \isacommand{then}\isamarkupfalse%
\ \isacommand{show}\isamarkupfalse%
\ {\isacharquery}{\kern0pt}case\ \isanewline
\ \ \ \ \ \ \ \ \isacommand{using}\isamarkupfalse%
\ comp{\isacharunderscore}{\kern0pt}def\ subst{\isacharunderscore}{\kern0pt}susp\ \isacommand{by}\isamarkupfalse%
\ simp\isanewline
\ \ \ \ \isacommand{qed}\isamarkupfalse%
\isanewline
\ \ \isacommand{qed}\isamarkupfalse%
\ {\isacharparenleft}{\kern0pt}simp{\isacharunderscore}{\kern0pt}all{\isacharparenright}{\kern0pt}\isanewline
\isacommand{qed}\isamarkupfalse%
%
\endisatagproof
{\isafoldproof}%
%
\isadelimproof
\isanewline
%
\endisadelimproof
\isanewline
\isanewline
\isacommand{lemma}\isamarkupfalse%
\ subst{\isacharunderscore}{\kern0pt}comp{\isacharunderscore}{\kern0pt}expand{\isacharcolon}{\kern0pt}\ {\isachardoublequoteopen}subst\ {\isacharparenleft}{\kern0pt}s{\isadigit{1}}\ {\isasymbullet}\ s{\isadigit{2}}{\isacharparenright}{\kern0pt}\ t\ {\isacharequal}{\kern0pt}\ subst\ s{\isadigit{1}}\ {\isacharparenleft}{\kern0pt}subst\ s{\isadigit{2}}\ t{\isacharparenright}{\kern0pt}{\isachardoublequoteclose}\isanewline
%
\isadelimproof
%
\endisadelimproof
%
\isatagproof
\isacommand{proof}\isamarkupfalse%
{\isacharparenleft}{\kern0pt}induct\ t{\isacharparenright}{\kern0pt}\isanewline
\ \ \isacommand{case}\isamarkupfalse%
\ {\isacharparenleft}{\kern0pt}Susp\ x{\isadigit{1}}\ x{\isadigit{2}}{\isacharparenright}{\kern0pt}\isanewline
\ \ \isacommand{then}\isamarkupfalse%
\ \isacommand{show}\isamarkupfalse%
\ {\isacharquery}{\kern0pt}case\ \isanewline
\ \ \ \isacommand{proof}\isamarkupfalse%
{\isacharparenleft}{\kern0pt}induct\ s{\isadigit{2}}{\isacharparenright}{\kern0pt}\isanewline
\ \ \ \ \ \ \isacommand{case}\isamarkupfalse%
\ Nil\isanewline
\ \ \ \ \ \ \isacommand{then}\isamarkupfalse%
\ \isacommand{show}\isamarkupfalse%
\ {\isacharquery}{\kern0pt}case\ \isacommand{using}\isamarkupfalse%
\ comp{\isacharunderscore}{\kern0pt}def\ subst{\isacharunderscore}{\kern0pt}susp\ subst{\isacharunderscore}{\kern0pt}swap{\isacharunderscore}{\kern0pt}comm\ \isacommand{by}\isamarkupfalse%
\ simp\isanewline
\ \ \ \isacommand{next}\isamarkupfalse%
\isanewline
\ \ \ \ \isacommand{case}\isamarkupfalse%
\ {\isacharparenleft}{\kern0pt}Cons\ a\ s{\isadigit{2}}{\isacharparenright}{\kern0pt}\isanewline
\ \ \ \ \isacommand{then}\isamarkupfalse%
\ \isacommand{show}\isamarkupfalse%
\ {\isacharquery}{\kern0pt}case\ \isacommand{using}\isamarkupfalse%
\ comp{\isacharunderscore}{\kern0pt}def\ subst{\isacharunderscore}{\kern0pt}susp\ subst{\isacharunderscore}{\kern0pt}swap{\isacharunderscore}{\kern0pt}comm\ \isacommand{by}\isamarkupfalse%
\ simp\isanewline
\ \ \isacommand{qed}\isamarkupfalse%
\isanewline
\isacommand{qed}\isamarkupfalse%
\ {\isacharparenleft}{\kern0pt}simp{\isacharunderscore}{\kern0pt}all{\isacharparenright}{\kern0pt}%
\endisatagproof
{\isafoldproof}%
%
\isadelimproof
\isanewline
%
\endisadelimproof
\isanewline
\isanewline
\isanewline
\isacommand{lemma}\isamarkupfalse%
\ subst{\isacharunderscore}{\kern0pt}assoc{\isacharcolon}{\kern0pt}\ {\isachardoublequoteopen}subst\ {\isacharparenleft}{\kern0pt}s{\isadigit{1}}\ {\isasymbullet}\ {\isacharparenleft}{\kern0pt}s{\isadigit{2}}\ {\isasymbullet}\ s{\isadigit{3}}{\isacharparenright}{\kern0pt}{\isacharparenright}{\kern0pt}{\isacharequal}{\kern0pt}subst\ {\isacharparenleft}{\kern0pt}{\isacharparenleft}{\kern0pt}s{\isadigit{1}}\ {\isasymbullet}\ s{\isadigit{2}}{\isacharparenright}{\kern0pt}\ {\isasymbullet}\ s{\isadigit{3}}{\isacharparenright}{\kern0pt}{\isachardoublequoteclose}\isanewline
%
\isadelimproof
%
\endisadelimproof
%
\isatagproof
\isacommand{proof}\isamarkupfalse%
{\isacharparenleft}{\kern0pt}rule\ ext{\isacharparenright}{\kern0pt}\isanewline
\ \ \isacommand{fix}\isamarkupfalse%
\ t\ \isanewline
\ \ \isacommand{show}\isamarkupfalse%
\ {\isachardoublequoteopen}subst\ {\isacharparenleft}{\kern0pt}s{\isadigit{1}}\ {\isasymbullet}\ {\isacharparenleft}{\kern0pt}s{\isadigit{2}}\ {\isasymbullet}\ s{\isadigit{3}}{\isacharparenright}{\kern0pt}{\isacharparenright}{\kern0pt}\ t\ {\isacharequal}{\kern0pt}\ subst\ {\isacharparenleft}{\kern0pt}s{\isadigit{1}}\ {\isasymbullet}\ s{\isadigit{2}}\ {\isasymbullet}\ s{\isadigit{3}}{\isacharparenright}{\kern0pt}\ t{\isachardoublequoteclose}\isanewline
\ \ \isacommand{proof}\isamarkupfalse%
{\isacharparenleft}{\kern0pt}induct\ t{\isacharparenright}{\kern0pt}\isanewline
\ \ \ \ \isacommand{case}\isamarkupfalse%
\ {\isacharparenleft}{\kern0pt}Susp\ pi\ X{\isacharparenright}{\kern0pt}\isanewline
\ \ \ \ \isacommand{then}\isamarkupfalse%
\ \isacommand{show}\isamarkupfalse%
\ {\isacharquery}{\kern0pt}case\ \isacommand{using}\isamarkupfalse%
\ subst{\isacharunderscore}{\kern0pt}comp{\isacharunderscore}{\kern0pt}expand\ \isacommand{by}\isamarkupfalse%
\ simp\isanewline
\ \ \isacommand{qed}\isamarkupfalse%
\ {\isacharparenleft}{\kern0pt}simp{\isacharunderscore}{\kern0pt}all{\isacharparenright}{\kern0pt}\isanewline
\isacommand{qed}\isamarkupfalse%
%
\endisatagproof
{\isafoldproof}%
%
\isadelimproof
\isanewline
%
\endisadelimproof
\isanewline
\isanewline
\isacommand{lemma}\isamarkupfalse%
\ fresh{\isacharunderscore}{\kern0pt}subst{\isacharcolon}{\kern0pt}\isanewline
\ \ \isakeyword{assumes}\ {\isachardoublequoteopen}nabla{\isadigit{1}}\ {\isasymturnstile}\ a\ {\isasymsharp}\ t{\isachardoublequoteclose}\isanewline
\ \ \ \ \isakeyword{and}\ {\isachardoublequoteopen}nabla{\isadigit{2}}\ {\isasymTurnstile}\ {\isacharparenleft}{\kern0pt}subst\ s{\isacharparenright}{\kern0pt}\ nabla{\isadigit{1}}{\isachardoublequoteclose}\isanewline
\ \ \isakeyword{shows}\ {\isachardoublequoteopen}nabla{\isadigit{2}}\ {\isasymturnstile}\ a\ {\isasymsharp}\ subst\ s\ t{\isachardoublequoteclose}\isanewline
%
\isadelimproof
\ \ %
\endisadelimproof
%
\isatagproof
\isacommand{using}\isamarkupfalse%
\ assms\isanewline
\isacommand{proof}\isamarkupfalse%
{\isacharparenleft}{\kern0pt}induction\ rule{\isacharcolon}{\kern0pt}\ fresh{\isachardot}{\kern0pt}induct{\isacharparenright}{\kern0pt}\isanewline
\ \ \isacommand{case}\isamarkupfalse%
\ {\isacharparenleft}{\kern0pt}fresh{\isacharunderscore}{\kern0pt}susp\ pi\ a\ X\ nabla{\isacharparenright}{\kern0pt}\isanewline
\ \ \isacommand{then}\isamarkupfalse%
\ \isacommand{show}\isamarkupfalse%
\ {\isacharquery}{\kern0pt}case\ \isanewline
\ \ \ \ \isacommand{using}\isamarkupfalse%
\ ext{\isacharunderscore}{\kern0pt}subst{\isacharunderscore}{\kern0pt}def\ subst{\isacharunderscore}{\kern0pt}susp\ fresh{\isacharunderscore}{\kern0pt}swap{\isacharunderscore}{\kern0pt}right\ \isanewline
\ \ \ \ \isacommand{by}\isamarkupfalse%
\ auto\isanewline
\isacommand{qed}\isamarkupfalse%
\ {\isacharparenleft}{\kern0pt}auto{\isacharparenright}{\kern0pt}%
\endisatagproof
{\isafoldproof}%
%
\isadelimproof
\isanewline
%
\endisadelimproof
\isanewline
\isanewline
\isacommand{lemma}\isamarkupfalse%
\ equ{\isacharunderscore}{\kern0pt}subst{\isacharcolon}{\kern0pt}\isanewline
\ \ \isakeyword{assumes}\ {\isachardoublequoteopen}nabla{\isadigit{1}}{\isasymturnstile}\ t{\isadigit{1}}\ {\isasymapprox}t{\isadigit{2}}{\isachardoublequoteclose}\isanewline
\ \ \ \ \isakeyword{and}\ {\isachardoublequoteopen}nabla{\isadigit{2}}\ {\isasymTurnstile}\ {\isacharparenleft}{\kern0pt}subst\ s{\isacharparenright}{\kern0pt}\ nabla{\isadigit{1}}{\isachardoublequoteclose}\isanewline
\ \ \isakeyword{shows}\ {\isachardoublequoteopen}nabla{\isadigit{2}}\ {\isasymturnstile}{\isacharparenleft}{\kern0pt}subst\ s\ t{\isadigit{1}}{\isacharparenright}{\kern0pt}{\isasymapprox}{\isacharparenleft}{\kern0pt}subst\ s\ t{\isadigit{2}}{\isacharparenright}{\kern0pt}{\isachardoublequoteclose}\isanewline
%
\isadelimproof
\ \ %
\endisadelimproof
%
\isatagproof
\isacommand{using}\isamarkupfalse%
\ assms\isanewline
\isacommand{proof}\isamarkupfalse%
{\isacharparenleft}{\kern0pt}induction\ rule{\isacharcolon}{\kern0pt}\ equ{\isachardot}{\kern0pt}induct{\isacharparenright}{\kern0pt}\isanewline
\ \ \isacommand{case}\isamarkupfalse%
\ {\isacharparenleft}{\kern0pt}equ{\isacharunderscore}{\kern0pt}abst{\isacharunderscore}{\kern0pt}ab\ a\ b\ nabla\ t{\isadigit{2}}\ t{\isadigit{1}}{\isacharparenright}{\kern0pt}\isanewline
\ \ \isacommand{then}\isamarkupfalse%
\ \isacommand{show}\isamarkupfalse%
\ {\isacharquery}{\kern0pt}case\ \isanewline
\ \ \ \ \isacommand{using}\isamarkupfalse%
\ subst{\isacharunderscore}{\kern0pt}swap{\isacharunderscore}{\kern0pt}comm\ fresh{\isacharunderscore}{\kern0pt}subst\ equ{\isachardot}{\kern0pt}equ{\isacharunderscore}{\kern0pt}abst{\isacharunderscore}{\kern0pt}ab\ \isacommand{by}\isamarkupfalse%
\ simp\ \isanewline
\isacommand{next}\isamarkupfalse%
\isanewline
\ \ \isacommand{case}\isamarkupfalse%
\ {\isacharparenleft}{\kern0pt}equ{\isacharunderscore}{\kern0pt}susp\ pi{\isadigit{1}}\ pi{\isadigit{2}}\ X\ nabla{\isacharparenright}{\kern0pt}\isanewline
\ \ \isacommand{then}\isamarkupfalse%
\ \isacommand{show}\isamarkupfalse%
\ {\isacharquery}{\kern0pt}case\ \isanewline
\ \ \ \ \isacommand{using}\isamarkupfalse%
\ subst{\isacharunderscore}{\kern0pt}susp\ equ{\isacharunderscore}{\kern0pt}pi{\isadigit{1}}{\isacharunderscore}{\kern0pt}pi{\isadigit{2}}{\isacharunderscore}{\kern0pt}add\ ext{\isacharunderscore}{\kern0pt}subst{\isacharunderscore}{\kern0pt}def\ subst{\isacharunderscore}{\kern0pt}susp\isanewline
\ \ \ \ \isacommand{by}\isamarkupfalse%
\ fastforce\isanewline
\isacommand{qed}\isamarkupfalse%
\ {\isacharparenleft}{\kern0pt}auto{\isacharparenright}{\kern0pt}%
\endisatagproof
{\isafoldproof}%
%
\isadelimproof
\isanewline
%
\endisadelimproof
\isanewline
\isacommand{lemma}\isamarkupfalse%
\ unif{\isacharunderscore}{\kern0pt}{\isadigit{1}}{\isacharcolon}{\kern0pt}\ \isanewline
\ \ \isakeyword{assumes}\ {\isachardoublequoteopen}nabla\ {\isasymturnstile}\ subst\ s\ {\isacharparenleft}{\kern0pt}Susp\ pi\ X{\isacharparenright}{\kern0pt}\ {\isasymapprox}\ subst\ s\ t{\isachardoublequoteclose}\isanewline
\ \ \isakeyword{shows}\ {\isachardoublequoteopen}nabla\ {\isasymTurnstile}\ subst\ {\isacharparenleft}{\kern0pt}s\ {\isasymbullet}\ {\isacharbrackleft}{\kern0pt}{\isacharparenleft}{\kern0pt}X{\isacharcomma}{\kern0pt}\ swap\ {\isacharparenleft}{\kern0pt}rev\ pi{\isacharparenright}{\kern0pt}\ t{\isacharparenright}{\kern0pt}{\isacharbrackright}{\kern0pt}{\isacharparenright}{\kern0pt}\ {\isasymapprox}\ subst\ s{\isachardoublequoteclose}\isanewline
%
\isadelimproof
\ \ %
\endisadelimproof
%
\isatagproof
\isacommand{using}\isamarkupfalse%
\ assms\isanewline
\ \ \isacommand{apply}\isamarkupfalse%
{\isacharparenleft}{\kern0pt}simp\ only{\isacharcolon}{\kern0pt}\ subst{\isacharunderscore}{\kern0pt}equ{\isacharunderscore}{\kern0pt}def{\isacharparenright}{\kern0pt}\isanewline
\isacommand{proof}\isamarkupfalse%
\isanewline
\ \ \isacommand{fix}\isamarkupfalse%
\ Xa\isanewline
\ \ \isacommand{assume}\isamarkupfalse%
\ hyps{\isacharcolon}{\kern0pt}\ {\isachardoublequoteopen}nabla\ {\isasymturnstile}\ subst\ s\ {\isacharparenleft}{\kern0pt}Susp\ pi\ X{\isacharparenright}{\kern0pt}\ {\isasymapprox}\ subst\ s\ t{\isachardoublequoteclose}\isanewline
\ \ \ {\isachardoublequoteopen}Xa\ {\isasymin}\ domn\ {\isacharparenleft}{\kern0pt}subst\ {\isacharparenleft}{\kern0pt}s\ {\isasymbullet}\ {\isacharbrackleft}{\kern0pt}{\isacharparenleft}{\kern0pt}X{\isacharcomma}{\kern0pt}\ swap\ {\isacharparenleft}{\kern0pt}rev\ pi{\isacharparenright}{\kern0pt}\ t{\isacharparenright}{\kern0pt}{\isacharbrackright}{\kern0pt}{\isacharparenright}{\kern0pt}{\isacharparenright}{\kern0pt}\ {\isasymunion}\ domn\ {\isacharparenleft}{\kern0pt}subst\ s{\isacharparenright}{\kern0pt}{\isachardoublequoteclose}\isanewline
\ \ \isacommand{show}\isamarkupfalse%
\ {\isachardoublequoteopen}nabla\ {\isasymturnstile}\ subst\ {\isacharparenleft}{\kern0pt}s\ {\isasymbullet}\ {\isacharbrackleft}{\kern0pt}{\isacharparenleft}{\kern0pt}X{\isacharcomma}{\kern0pt}\ swap\ {\isacharparenleft}{\kern0pt}rev\ pi{\isacharparenright}{\kern0pt}\ t{\isacharparenright}{\kern0pt}{\isacharbrackright}{\kern0pt}{\isacharparenright}{\kern0pt}\ {\isacharparenleft}{\kern0pt}Susp\ {\isacharbrackleft}{\kern0pt}{\isacharbrackright}{\kern0pt}\ Xa{\isacharparenright}{\kern0pt}\ {\isasymapprox}\ subst\ s\ {\isacharparenleft}{\kern0pt}Susp\ {\isacharbrackleft}{\kern0pt}{\isacharbrackright}{\kern0pt}\ Xa{\isacharparenright}{\kern0pt}{\isachardoublequoteclose}\isanewline
\ \ \isacommand{proof}\isamarkupfalse%
{\isacharminus}{\kern0pt}\isanewline
\ \ \ \ \isacommand{have}\isamarkupfalse%
\ one{\isacharcolon}{\kern0pt}\ \isanewline
\ \ \ \ \ \ {\isachardoublequoteopen}nabla\ {\isasymturnstile}\ subst\ {\isacharparenleft}{\kern0pt}s\ {\isasymbullet}\ {\isacharbrackleft}{\kern0pt}{\isacharparenleft}{\kern0pt}X{\isacharcomma}{\kern0pt}\ swap\ {\isacharparenleft}{\kern0pt}rev\ pi{\isacharparenright}{\kern0pt}\ t{\isacharparenright}{\kern0pt}{\isacharbrackright}{\kern0pt}{\isacharparenright}{\kern0pt}\ {\isacharparenleft}{\kern0pt}Susp\ pi\ Xa{\isacharparenright}{\kern0pt}\ {\isasymapprox}\ \isanewline
\ \ \ \ \ \ \ subst\ s\ {\isacharparenleft}{\kern0pt}subst\ {\isacharbrackleft}{\kern0pt}{\isacharparenleft}{\kern0pt}X{\isacharcomma}{\kern0pt}\ swap\ {\isacharparenleft}{\kern0pt}rev\ pi{\isacharparenright}{\kern0pt}\ t{\isacharparenright}{\kern0pt}{\isacharbrackright}{\kern0pt}\ {\isacharparenleft}{\kern0pt}Susp\ pi\ Xa{\isacharparenright}{\kern0pt}{\isacharparenright}{\kern0pt}{\isachardoublequoteclose}\isanewline
\ \ \ \ \ \ \isacommand{using}\isamarkupfalse%
\ subst{\isacharunderscore}{\kern0pt}comp{\isacharunderscore}{\kern0pt}expand\ equ{\isacharunderscore}{\kern0pt}refl\ \isacommand{by}\isamarkupfalse%
\ simp\isanewline
\ \ \ \ \isacommand{have}\isamarkupfalse%
\ two{\isacharcolon}{\kern0pt}\ \isanewline
\ \ \ \ \ \ {\isachardoublequoteopen}nabla\ {\isasymturnstile}\ subst\ s\ {\isacharparenleft}{\kern0pt}subst\ {\isacharbrackleft}{\kern0pt}{\isacharparenleft}{\kern0pt}X{\isacharcomma}{\kern0pt}\ swap\ {\isacharparenleft}{\kern0pt}rev\ pi{\isacharparenright}{\kern0pt}\ t{\isacharparenright}{\kern0pt}{\isacharbrackright}{\kern0pt}\ {\isacharparenleft}{\kern0pt}Susp\ pi\ Xa{\isacharparenright}{\kern0pt}{\isacharparenright}{\kern0pt}\ {\isasymapprox}\ \isanewline
\ \ \ \ \ \ subst\ s\ {\isacharparenleft}{\kern0pt}swap\ pi\ {\isacharparenleft}{\kern0pt}look{\isacharunderscore}{\kern0pt}up\ Xa\ {\isacharbrackleft}{\kern0pt}{\isacharparenleft}{\kern0pt}X{\isacharcomma}{\kern0pt}\ swap\ {\isacharparenleft}{\kern0pt}rev\ pi{\isacharparenright}{\kern0pt}\ t{\isacharparenright}{\kern0pt}{\isacharbrackright}{\kern0pt}{\isacharparenright}{\kern0pt}{\isacharparenright}{\kern0pt}{\isachardoublequoteclose}\isanewline
\ \ \ \ \ \ \isacommand{using}\isamarkupfalse%
\ equ{\isacharunderscore}{\kern0pt}refl\ subst{\isacharunderscore}{\kern0pt}susp{\isacharbrackleft}{\kern0pt}of\ {\isacartoucheopen}{\isacharbrackleft}{\kern0pt}{\isacharparenleft}{\kern0pt}X{\isacharcomma}{\kern0pt}\ swap\ {\isacharparenleft}{\kern0pt}rev\ pi{\isacharparenright}{\kern0pt}t{\isacharparenright}{\kern0pt}{\isacharbrackright}{\kern0pt}{\isacartoucheclose}\ pi\ Xa{\isacharbrackright}{\kern0pt}\ \isanewline
\ \ \ \ \ \ \isacommand{by}\isamarkupfalse%
\ simp\isanewline
\ \ \ \ \isacommand{have}\isamarkupfalse%
\ {\isachardoublequoteopen}nabla\ {\isasymturnstile}\ subst\ {\isacharparenleft}{\kern0pt}s\ {\isasymbullet}\ {\isacharbrackleft}{\kern0pt}{\isacharparenleft}{\kern0pt}X{\isacharcomma}{\kern0pt}\ swap\ {\isacharparenleft}{\kern0pt}rev\ pi{\isacharparenright}{\kern0pt}\ t{\isacharparenright}{\kern0pt}{\isacharbrackright}{\kern0pt}{\isacharparenright}{\kern0pt}\ {\isacharparenleft}{\kern0pt}Susp\ pi\ Xa{\isacharparenright}{\kern0pt}\ {\isasymapprox}\ \isanewline
\ \ \ \ subst\ s\ {\isacharparenleft}{\kern0pt}swap\ pi\ {\isacharparenleft}{\kern0pt}look{\isacharunderscore}{\kern0pt}up\ Xa\ {\isacharbrackleft}{\kern0pt}{\isacharparenleft}{\kern0pt}X{\isacharcomma}{\kern0pt}\ swap\ {\isacharparenleft}{\kern0pt}rev\ pi{\isacharparenright}{\kern0pt}\ t{\isacharparenright}{\kern0pt}{\isacharbrackright}{\kern0pt}{\isacharparenright}{\kern0pt}{\isacharparenright}{\kern0pt}{\isachardoublequoteclose}\isanewline
\ \ \ \ \ \ \isacommand{using}\isamarkupfalse%
\ equ{\isacharunderscore}{\kern0pt}trans{\isacharbrackleft}{\kern0pt}OF\ one\ two{\isacharbrackright}{\kern0pt}\ \isacommand{by}\isamarkupfalse%
\ simp\isanewline
\ \ \ \ \isacommand{hence}\isamarkupfalse%
\ \ {\isachardoublequoteopen}nabla\ {\isasymturnstile}\ swap\ pi\ {\isacharparenleft}{\kern0pt}look{\isacharunderscore}{\kern0pt}up\ Xa\ {\isacharparenleft}{\kern0pt}s\ {\isasymbullet}\ {\isacharbrackleft}{\kern0pt}{\isacharparenleft}{\kern0pt}X{\isacharcomma}{\kern0pt}\ swap\ {\isacharparenleft}{\kern0pt}rev\ pi{\isacharparenright}{\kern0pt}\ t{\isacharparenright}{\kern0pt}{\isacharbrackright}{\kern0pt}{\isacharparenright}{\kern0pt}{\isacharparenright}{\kern0pt}\ {\isasymapprox}\ \isanewline
\ \ \ \ subst\ s\ {\isacharparenleft}{\kern0pt}swap\ pi\ {\isacharparenleft}{\kern0pt}look{\isacharunderscore}{\kern0pt}up\ Xa\ {\isacharbrackleft}{\kern0pt}{\isacharparenleft}{\kern0pt}X{\isacharcomma}{\kern0pt}\ swap\ {\isacharparenleft}{\kern0pt}rev\ pi{\isacharparenright}{\kern0pt}\ t{\isacharparenright}{\kern0pt}{\isacharbrackright}{\kern0pt}{\isacharparenright}{\kern0pt}{\isacharparenright}{\kern0pt}{\isachardoublequoteclose}\isanewline
\ \ \ \ \ \ \isacommand{using}\isamarkupfalse%
\ subst{\isacharunderscore}{\kern0pt}susp{\isacharbrackleft}{\kern0pt}of\ {\isacartoucheopen}{\isacharparenleft}{\kern0pt}s\ {\isasymbullet}\ {\isacharbrackleft}{\kern0pt}{\isacharparenleft}{\kern0pt}X{\isacharcomma}{\kern0pt}\ swap\ {\isacharparenleft}{\kern0pt}rev\ pi{\isacharparenright}{\kern0pt}\ t{\isacharparenright}{\kern0pt}{\isacharbrackright}{\kern0pt}{\isacharparenright}{\kern0pt}{\isacartoucheclose}{\isacharbrackright}{\kern0pt}\ \isacommand{by}\isamarkupfalse%
\ simp\isanewline
\ \ \ \ \isacommand{hence}\isamarkupfalse%
\ swap{\isacharunderscore}{\kern0pt}pi{\isacharunderscore}{\kern0pt}outside{\isacharcolon}{\kern0pt}\ \ {\isachardoublequoteopen}nabla\ {\isasymturnstile}\ swap\ pi\ {\isacharparenleft}{\kern0pt}look{\isacharunderscore}{\kern0pt}up\ Xa\ {\isacharparenleft}{\kern0pt}s\ {\isasymbullet}\ {\isacharbrackleft}{\kern0pt}{\isacharparenleft}{\kern0pt}X{\isacharcomma}{\kern0pt}\ swap\ {\isacharparenleft}{\kern0pt}rev\ pi{\isacharparenright}{\kern0pt}\ t{\isacharparenright}{\kern0pt}{\isacharbrackright}{\kern0pt}{\isacharparenright}{\kern0pt}{\isacharparenright}{\kern0pt}\ {\isasymapprox}\ \isanewline
\ \ \ \ swap\ pi\ {\isacharparenleft}{\kern0pt}subst\ s\ {\isacharparenleft}{\kern0pt}look{\isacharunderscore}{\kern0pt}up\ Xa\ {\isacharbrackleft}{\kern0pt}{\isacharparenleft}{\kern0pt}X{\isacharcomma}{\kern0pt}\ swap\ {\isacharparenleft}{\kern0pt}rev\ pi{\isacharparenright}{\kern0pt}\ t{\isacharparenright}{\kern0pt}{\isacharbrackright}{\kern0pt}{\isacharparenright}{\kern0pt}{\isacharparenright}{\kern0pt}{\isachardoublequoteclose}\isanewline
\ \ \ \ \ \ \isacommand{using}\isamarkupfalse%
\ subst{\isacharunderscore}{\kern0pt}swap{\isacharunderscore}{\kern0pt}comm\ \isacommand{by}\isamarkupfalse%
\ simp\isanewline
\ \ \ \ \isacommand{hence}\isamarkupfalse%
\ {\isachardoublequoteopen}nabla\ {\isasymturnstile}\ {\isacharparenleft}{\kern0pt}look{\isacharunderscore}{\kern0pt}up\ Xa\ {\isacharparenleft}{\kern0pt}s\ {\isasymbullet}\ {\isacharbrackleft}{\kern0pt}{\isacharparenleft}{\kern0pt}X{\isacharcomma}{\kern0pt}\ swap\ {\isacharparenleft}{\kern0pt}rev\ pi{\isacharparenright}{\kern0pt}\ t{\isacharparenright}{\kern0pt}{\isacharbrackright}{\kern0pt}{\isacharparenright}{\kern0pt}{\isacharparenright}{\kern0pt}\ {\isasymapprox}\ \isanewline
\ \ \ \ {\isacharparenleft}{\kern0pt}subst\ s\ {\isacharparenleft}{\kern0pt}look{\isacharunderscore}{\kern0pt}up\ Xa\ {\isacharbrackleft}{\kern0pt}{\isacharparenleft}{\kern0pt}X{\isacharcomma}{\kern0pt}\ swap\ {\isacharparenleft}{\kern0pt}rev\ pi{\isacharparenright}{\kern0pt}\ t{\isacharparenright}{\kern0pt}{\isacharbrackright}{\kern0pt}{\isacharparenright}{\kern0pt}{\isacharparenright}{\kern0pt}{\isachardoublequoteclose}\isanewline
\ \ \ \ \ \ \isacommand{using}\isamarkupfalse%
\ equ{\isacharunderscore}{\kern0pt}dec{\isacharunderscore}{\kern0pt}pi{\isacharbrackleft}{\kern0pt}OF\ swap{\isacharunderscore}{\kern0pt}pi{\isacharunderscore}{\kern0pt}outside{\isacharbrackright}{\kern0pt}\ \isacommand{by}\isamarkupfalse%
\ simp\isanewline
\ \ \ \ \isacommand{hence}\isamarkupfalse%
\ three{\isacharcolon}{\kern0pt}\ \isanewline
\ \ \ \ {\isachardoublequoteopen}nabla\ {\isasymturnstile}\ subst\ {\isacharparenleft}{\kern0pt}s\ {\isasymbullet}\ {\isacharbrackleft}{\kern0pt}{\isacharparenleft}{\kern0pt}X{\isacharcomma}{\kern0pt}\ swap\ {\isacharparenleft}{\kern0pt}rev\ pi{\isacharparenright}{\kern0pt}\ t{\isacharparenright}{\kern0pt}{\isacharbrackright}{\kern0pt}{\isacharparenright}{\kern0pt}\ {\isacharparenleft}{\kern0pt}Susp\ {\isacharbrackleft}{\kern0pt}{\isacharbrackright}{\kern0pt}\ Xa{\isacharparenright}{\kern0pt}\ {\isasymapprox}\ \isanewline
\ \ \ \ {\isacharparenleft}{\kern0pt}subst\ s\ {\isacharparenleft}{\kern0pt}look{\isacharunderscore}{\kern0pt}up\ Xa\ {\isacharbrackleft}{\kern0pt}{\isacharparenleft}{\kern0pt}X{\isacharcomma}{\kern0pt}\ swap\ {\isacharparenleft}{\kern0pt}rev\ pi{\isacharparenright}{\kern0pt}\ t{\isacharparenright}{\kern0pt}{\isacharbrackright}{\kern0pt}{\isacharparenright}{\kern0pt}{\isacharparenright}{\kern0pt}{\isachardoublequoteclose}\ \isanewline
\ \ \ \ \ \ \isacommand{using}\isamarkupfalse%
\ subst{\isacharunderscore}{\kern0pt}susp\ \isacommand{by}\isamarkupfalse%
\ simp\isanewline
\ \ \ \ \isacommand{then}\isamarkupfalse%
\ \isacommand{show}\isamarkupfalse%
\ {\isacharquery}{\kern0pt}thesis\isanewline
\ \ \ \ \ \ \isacommand{proof}\isamarkupfalse%
{\isacharparenleft}{\kern0pt}cases\ {\isachardoublequoteopen}Xa\ {\isacharequal}{\kern0pt}\ X{\isachardoublequoteclose}{\isacharparenright}{\kern0pt}\isanewline
\ \ \ \ \ \ \ \ \isacommand{case}\isamarkupfalse%
\ True\isanewline
\ \ \ \ \ \ \ \ \isacommand{hence}\isamarkupfalse%
\ {\isachardoublequoteopen}subst\ s\ {\isacharparenleft}{\kern0pt}look{\isacharunderscore}{\kern0pt}up\ Xa\ {\isacharbrackleft}{\kern0pt}{\isacharparenleft}{\kern0pt}X{\isacharcomma}{\kern0pt}\ swap\ {\isacharparenleft}{\kern0pt}rev\ pi{\isacharparenright}{\kern0pt}\ t{\isacharparenright}{\kern0pt}{\isacharbrackright}{\kern0pt}{\isacharparenright}{\kern0pt}\ {\isacharequal}{\kern0pt}\ subst\ s\ {\isacharparenleft}{\kern0pt}swap\ {\isacharparenleft}{\kern0pt}rev\ pi{\isacharparenright}{\kern0pt}\ t{\isacharparenright}{\kern0pt}{\isachardoublequoteclose}\isanewline
\ \ \ \ \ \ \ \ \ \ \isacommand{by}\isamarkupfalse%
\ simp\isanewline
\ \ \ \ \ \ \ \ \isacommand{with}\isamarkupfalse%
\ three\ \isacommand{have}\isamarkupfalse%
\isanewline
\ \ \ \ \ \ \ \ \ \ {\isachardoublequoteopen}nabla\ {\isasymturnstile}\ subst\ {\isacharparenleft}{\kern0pt}s\ {\isasymbullet}\ {\isacharbrackleft}{\kern0pt}{\isacharparenleft}{\kern0pt}X{\isacharcomma}{\kern0pt}\ swap\ {\isacharparenleft}{\kern0pt}rev\ pi{\isacharparenright}{\kern0pt}\ t{\isacharparenright}{\kern0pt}{\isacharbrackright}{\kern0pt}{\isacharparenright}{\kern0pt}\ {\isacharparenleft}{\kern0pt}Susp\ {\isacharbrackleft}{\kern0pt}{\isacharbrackright}{\kern0pt}\ Xa{\isacharparenright}{\kern0pt}\ {\isasymapprox}\ swap\ {\isacharparenleft}{\kern0pt}rev\ pi{\isacharparenright}{\kern0pt}\ {\isacharparenleft}{\kern0pt}subst\ s\ t{\isacharparenright}{\kern0pt}{\isachardoublequoteclose}\isanewline
\ \ \ \ \ \ \ \ \ \ \isacommand{using}\isamarkupfalse%
\ subst{\isacharunderscore}{\kern0pt}swap{\isacharunderscore}{\kern0pt}comm\ \isacommand{by}\isamarkupfalse%
\ auto\isanewline
\ \ \ \ \ \ \ \ \isacommand{hence}\isamarkupfalse%
\ i{\isacharcolon}{\kern0pt}\ {\isachardoublequoteopen}nabla\ {\isasymturnstile}\ swap\ pi\ {\isacharparenleft}{\kern0pt}subst\ {\isacharparenleft}{\kern0pt}s\ {\isasymbullet}\ {\isacharbrackleft}{\kern0pt}{\isacharparenleft}{\kern0pt}X{\isacharcomma}{\kern0pt}\ swap\ {\isacharparenleft}{\kern0pt}rev\ pi{\isacharparenright}{\kern0pt}\ t{\isacharparenright}{\kern0pt}{\isacharbrackright}{\kern0pt}{\isacharparenright}{\kern0pt}\ {\isacharparenleft}{\kern0pt}Susp\ {\isacharbrackleft}{\kern0pt}{\isacharbrackright}{\kern0pt}\ Xa{\isacharparenright}{\kern0pt}{\isacharparenright}{\kern0pt}\isanewline
\ \ \ \ \ \ \ \ \ \ {\isasymapprox}\ subst\ s\ t{\isachardoublequoteclose}\ \isacommand{using}\isamarkupfalse%
\ swap{\isacharunderscore}{\kern0pt}inv{\isacharunderscore}{\kern0pt}side\ \isacommand{by}\isamarkupfalse%
\ simp\isanewline
\ \ \ \ \ \ \ \ \isacommand{hence}\isamarkupfalse%
\ {\isachardoublequoteopen}nabla\ {\isasymturnstile}\ swap\ pi\ {\isacharparenleft}{\kern0pt}subst\ {\isacharparenleft}{\kern0pt}s\ {\isasymbullet}\ {\isacharbrackleft}{\kern0pt}{\isacharparenleft}{\kern0pt}X{\isacharcomma}{\kern0pt}\ swap\ {\isacharparenleft}{\kern0pt}rev\ pi{\isacharparenright}{\kern0pt}\ t{\isacharparenright}{\kern0pt}{\isacharbrackright}{\kern0pt}{\isacharparenright}{\kern0pt}\ {\isacharparenleft}{\kern0pt}Susp\ {\isacharbrackleft}{\kern0pt}{\isacharbrackright}{\kern0pt}\ Xa{\isacharparenright}{\kern0pt}{\isacharparenright}{\kern0pt}\isanewline
\ \ \ \ \ \ \ \ \ \ {\isasymapprox}\ subst\ s\ {\isacharparenleft}{\kern0pt}Susp\ pi\ Xa{\isacharparenright}{\kern0pt}{\isachardoublequoteclose}\isanewline
\ \ \ \ \ \ \ \ \ \ \isacommand{using}\isamarkupfalse%
\ True\ equ{\isacharunderscore}{\kern0pt}trans{\isacharbrackleft}{\kern0pt}OF\ i\ equ{\isacharunderscore}{\kern0pt}symm{\isacharbrackleft}{\kern0pt}OF\ assms{\isacharparenleft}{\kern0pt}{\isadigit{1}}{\isacharparenright}{\kern0pt}{\isacharbrackright}{\kern0pt}{\isacharbrackright}{\kern0pt}\ \isacommand{by}\isamarkupfalse%
\ simp\isanewline
\ \ \ \ \ \ \ \ \isacommand{hence}\isamarkupfalse%
\ {\isachardoublequoteopen}nabla\ {\isasymturnstile}\ swap\ pi\ {\isacharparenleft}{\kern0pt}subst\ {\isacharparenleft}{\kern0pt}s\ {\isasymbullet}\ {\isacharbrackleft}{\kern0pt}{\isacharparenleft}{\kern0pt}X{\isacharcomma}{\kern0pt}\ swap\ {\isacharparenleft}{\kern0pt}rev\ pi{\isacharparenright}{\kern0pt}\ t{\isacharparenright}{\kern0pt}{\isacharbrackright}{\kern0pt}{\isacharparenright}{\kern0pt}\ {\isacharparenleft}{\kern0pt}Susp\ {\isacharbrackleft}{\kern0pt}{\isacharbrackright}{\kern0pt}\ Xa{\isacharparenright}{\kern0pt}{\isacharparenright}{\kern0pt}\isanewline
\ \ \ \ \ \ \ \ \ \ {\isasymapprox}\ swap\ pi\ {\isacharparenleft}{\kern0pt}look{\isacharunderscore}{\kern0pt}up\ Xa\ s{\isacharparenright}{\kern0pt}{\isachardoublequoteclose}\isanewline
\ \ \ \ \ \ \ \ \ \ \isacommand{using}\isamarkupfalse%
\ subst{\isacharunderscore}{\kern0pt}susp\ \isacommand{by}\isamarkupfalse%
\ simp\isanewline
\ \ \ \ \ \ \ \ \isacommand{then}\isamarkupfalse%
\ \isacommand{show}\isamarkupfalse%
\ {\isachardoublequoteopen}nabla\ {\isasymturnstile}\ {\isacharparenleft}{\kern0pt}subst\ {\isacharparenleft}{\kern0pt}s\ {\isasymbullet}\ {\isacharbrackleft}{\kern0pt}{\isacharparenleft}{\kern0pt}X{\isacharcomma}{\kern0pt}\ swap\ {\isacharparenleft}{\kern0pt}rev\ pi{\isacharparenright}{\kern0pt}\ t{\isacharparenright}{\kern0pt}{\isacharbrackright}{\kern0pt}{\isacharparenright}{\kern0pt}\ {\isacharparenleft}{\kern0pt}Susp\ {\isacharbrackleft}{\kern0pt}{\isacharbrackright}{\kern0pt}\ Xa{\isacharparenright}{\kern0pt}{\isacharparenright}{\kern0pt}\isanewline
\ \ \ \ \ \ \ \ \ \ {\isasymapprox}\ subst\ s\ {\isacharparenleft}{\kern0pt}Susp\ {\isacharbrackleft}{\kern0pt}{\isacharbrackright}{\kern0pt}\ Xa{\isacharparenright}{\kern0pt}{\isachardoublequoteclose}\isanewline
\ \ \ \ \ \ \ \ \ \ \isacommand{using}\isamarkupfalse%
\ equ{\isacharunderscore}{\kern0pt}dec{\isacharunderscore}{\kern0pt}pi\ subst{\isacharunderscore}{\kern0pt}susp\ \isacommand{by}\isamarkupfalse%
\ simp\isanewline
\ \ \ \ \ \ \isacommand{next}\isamarkupfalse%
\isanewline
\ \ \ \ \ \ \ \ \isacommand{case}\isamarkupfalse%
\ False\isanewline
\ \ \ \ \ \ \ \ \isacommand{hence}\isamarkupfalse%
\ {\isachardoublequoteopen}subst\ s\ {\isacharparenleft}{\kern0pt}look{\isacharunderscore}{\kern0pt}up\ Xa\ {\isacharbrackleft}{\kern0pt}{\isacharparenleft}{\kern0pt}X{\isacharcomma}{\kern0pt}\ swap\ {\isacharparenleft}{\kern0pt}rev\ pi{\isacharparenright}{\kern0pt}\ t{\isacharparenright}{\kern0pt}{\isacharbrackright}{\kern0pt}{\isacharparenright}{\kern0pt}\ {\isacharequal}{\kern0pt}\ subst\ s\ {\isacharparenleft}{\kern0pt}Susp\ {\isacharbrackleft}{\kern0pt}{\isacharbrackright}{\kern0pt}\ Xa{\isacharparenright}{\kern0pt}{\isachardoublequoteclose}\isanewline
\ \ \ \ \ \ \ \ \ \ \isacommand{by}\isamarkupfalse%
\ simp\isanewline
\ \ \ \ \ \ \ \ \isacommand{with}\isamarkupfalse%
\ three\ \isacommand{show}\isamarkupfalse%
\ {\isacharquery}{\kern0pt}thesis\ \isacommand{by}\isamarkupfalse%
\ auto\isanewline
\ \ \ \ \ \ \isacommand{qed}\isamarkupfalse%
\isanewline
\ \ \ \ \isacommand{qed}\isamarkupfalse%
\isanewline
\ \ \isacommand{qed}\isamarkupfalse%
%
\endisatagproof
{\isafoldproof}%
%
\isadelimproof
\isanewline
%
\endisadelimproof
\isanewline
\isacommand{lemma}\isamarkupfalse%
\ subst{\isacharunderscore}{\kern0pt}equ{\isacharunderscore}{\kern0pt}a{\isacharcolon}{\kern0pt}\isanewline
\ \ \isakeyword{assumes}\ {\isachardoublequoteopen}nabla\ {\isasymTurnstile}\ {\isacharparenleft}{\kern0pt}subst\ s{\isadigit{1}}{\isacharparenright}{\kern0pt}\ {\isasymapprox}\ {\isacharparenleft}{\kern0pt}subst\ s{\isadigit{2}}{\isacharparenright}{\kern0pt}{\isachardoublequoteclose}\isanewline
\ \ \ \ \isakeyword{and}\ {\isachardoublequoteopen}nabla\ {\isasymturnstile}\ {\isacharparenleft}{\kern0pt}subst\ s{\isadigit{2}}\ t{\isadigit{1}}{\isacharparenright}{\kern0pt}\ {\isasymapprox}\ t{\isadigit{2}}{\isachardoublequoteclose}\isanewline
\ \ \isakeyword{shows}\ {\isachardoublequoteopen}nabla\ {\isasymturnstile}\ {\isacharparenleft}{\kern0pt}subst\ s{\isadigit{1}}\ t{\isadigit{1}}{\isacharparenright}{\kern0pt}\ {\isasymapprox}\ t{\isadigit{2}}{\isachardoublequoteclose}\isanewline
%
\isadelimproof
\ \ %
\endisadelimproof
%
\isatagproof
\isacommand{using}\isamarkupfalse%
\ assms\isanewline
\isacommand{proof}\isamarkupfalse%
{\isacharparenleft}{\kern0pt}induct\ t{\isadigit{1}}\ arbitrary{\isacharcolon}{\kern0pt}\ t{\isadigit{2}}{\isacharparenright}{\kern0pt}\isanewline
\ \ \isacommand{case}\isamarkupfalse%
\ {\isacharparenleft}{\kern0pt}Abst\ a\ t{\isadigit{1}}{\isacharprime}{\kern0pt}{\isacharparenright}{\kern0pt}\isanewline
\ \ \isacommand{then}\isamarkupfalse%
\ \isacommand{obtain}\isamarkupfalse%
\ t{\isadigit{2}}{\isacharprime}{\kern0pt}\ b\ \isanewline
\ \ \ \ \isakeyword{where}\ t{\isadigit{2}}{\isacharcolon}{\kern0pt}\ {\isachardoublequoteopen}t{\isadigit{2}}\ {\isacharequal}{\kern0pt}\ Abst\ b\ t{\isadigit{2}}{\isacharprime}{\kern0pt}{\isachardoublequoteclose}\isanewline
\ \ \ \ \isacommand{by}\isamarkupfalse%
{\isacharparenleft}{\kern0pt}auto\ elim{\isacharcolon}{\kern0pt}\ equ{\isachardot}{\kern0pt}cases{\isacharparenright}{\kern0pt}\isanewline
\ \ \isacommand{with}\isamarkupfalse%
\ Abst{\isacharparenleft}{\kern0pt}{\isadigit{3}}{\isacharparenright}{\kern0pt}\ \isacommand{have}\isamarkupfalse%
\isanewline
\ \ \ \ i{\isacharcolon}{\kern0pt}\ {\isachardoublequoteopen}nabla\ {\isasymturnstile}\ Abst\ a\ {\isacharparenleft}{\kern0pt}subst\ s{\isadigit{2}}\ t{\isadigit{1}}{\isacharprime}{\kern0pt}{\isacharparenright}{\kern0pt}\ {\isasymapprox}\ Abst\ b\ t{\isadigit{2}}{\isacharprime}{\kern0pt}{\isachardoublequoteclose}\isanewline
\ \ \ \ \isacommand{using}\isamarkupfalse%
\ subst{\isacharunderscore}{\kern0pt}abst\ \isacommand{by}\isamarkupfalse%
\ simp\isanewline
\ \ \isacommand{then}\isamarkupfalse%
\ \isacommand{show}\isamarkupfalse%
\ {\isacharquery}{\kern0pt}case\ \isanewline
\ \ \isacommand{proof}\isamarkupfalse%
{\isacharparenleft}{\kern0pt}cases\ {\isachardoublequoteopen}a{\isacharequal}{\kern0pt}b{\isachardoublequoteclose}{\isacharparenright}{\kern0pt}\isanewline
\ \ \ \ \isacommand{case}\isamarkupfalse%
\ True\isanewline
\ \ \ \ \isacommand{hence}\isamarkupfalse%
\ {\isachardoublequoteopen}nabla\ {\isasymturnstile}\ {\isacharparenleft}{\kern0pt}subst\ s{\isadigit{2}}\ t{\isadigit{1}}{\isacharprime}{\kern0pt}{\isacharparenright}{\kern0pt}\ {\isasymapprox}\ t{\isadigit{2}}{\isacharprime}{\kern0pt}{\isachardoublequoteclose}\ \isanewline
\ \ \ \ \ \ \isacommand{using}\isamarkupfalse%
\ Equ{\isacharunderscore}{\kern0pt}elims{\isacharparenleft}{\kern0pt}{\isadigit{7}}{\isacharparenright}{\kern0pt}{\isacharbrackleft}{\kern0pt}OF\ i{\isacharbrackright}{\kern0pt}\ \isacommand{by}\isamarkupfalse%
\ auto\isanewline
\ \ \ \ \isacommand{with}\isamarkupfalse%
\ Abst{\isacharparenleft}{\kern0pt}{\isadigit{1}}{\isacharparenright}{\kern0pt}{\isacharbrackleft}{\kern0pt}OF\ Abst{\isacharparenleft}{\kern0pt}{\isadigit{2}}{\isacharparenright}{\kern0pt}{\isacharcomma}{\kern0pt}\ of\ t{\isadigit{2}}{\isacharprime}{\kern0pt}{\isacharbrackright}{\kern0pt}\isanewline
\ \ \ \ \isacommand{have}\isamarkupfalse%
\ {\isachardoublequoteopen}nabla\ {\isasymturnstile}\ subst\ s{\isadigit{1}}\ t{\isadigit{1}}{\isacharprime}{\kern0pt}\ {\isasymapprox}\ t{\isadigit{2}}{\isacharprime}{\kern0pt}{\isachardoublequoteclose}\ \isanewline
\ \ \ \ \ \ \isacommand{by}\isamarkupfalse%
\ simp\isanewline
\ \ \ \ \isacommand{hence}\isamarkupfalse%
\ {\isachardoublequoteopen}nabla\ {\isasymturnstile}\ Abst\ a\ {\isacharparenleft}{\kern0pt}subst\ s{\isadigit{1}}\ t{\isadigit{1}}{\isacharprime}{\kern0pt}{\isacharparenright}{\kern0pt}\ {\isasymapprox}\ Abst\ b\ t{\isadigit{2}}{\isacharprime}{\kern0pt}{\isachardoublequoteclose}\isanewline
\ \ \ \ \ \ \isacommand{using}\isamarkupfalse%
\ True\ equ{\isacharunderscore}{\kern0pt}abst{\isacharunderscore}{\kern0pt}aa\ \isacommand{by}\isamarkupfalse%
\ simp\isanewline
\ \ \ \ \isacommand{with}\isamarkupfalse%
\ t{\isadigit{2}}\ \isacommand{show}\isamarkupfalse%
\ {\isacharquery}{\kern0pt}thesis\isanewline
\ \ \ \ \ \ \isacommand{using}\isamarkupfalse%
\ subst{\isacharunderscore}{\kern0pt}abst{\isacharbrackleft}{\kern0pt}of\ s{\isadigit{1}}\ a\ t{\isadigit{1}}{\isacharprime}{\kern0pt}{\isacharbrackright}{\kern0pt}\ \isacommand{by}\isamarkupfalse%
\ simp\isanewline
\ \ \isacommand{next}\isamarkupfalse%
\isanewline
\ \ \ \ \isacommand{case}\isamarkupfalse%
\ False\isanewline
\ \ \ \ \isacommand{hence}\isamarkupfalse%
\ {\isachardoublequoteopen}nabla\ {\isasymturnstile}\ {\isacharparenleft}{\kern0pt}subst\ s{\isadigit{2}}\ t{\isadigit{1}}{\isacharprime}{\kern0pt}{\isacharparenright}{\kern0pt}\ {\isasymapprox}\ swap\ {\isacharbrackleft}{\kern0pt}{\isacharparenleft}{\kern0pt}a{\isacharcomma}{\kern0pt}b{\isacharparenright}{\kern0pt}{\isacharbrackright}{\kern0pt}\ t{\isadigit{2}}{\isacharprime}{\kern0pt}{\isachardoublequoteclose}\isanewline
\ \ \ \ \ \ \ \ \ \ \isakeyword{and}\ fresh{\isacharcolon}{\kern0pt}\ {\isachardoublequoteopen}nabla\ {\isasymturnstile}\ a\ {\isasymsharp}\ t{\isadigit{2}}{\isacharprime}{\kern0pt}{\isachardoublequoteclose}\isanewline
\ \ \ \ \ \ \isacommand{using}\isamarkupfalse%
\ Equ{\isacharunderscore}{\kern0pt}elims{\isacharparenleft}{\kern0pt}{\isadigit{7}}{\isacharparenright}{\kern0pt}{\isacharbrackleft}{\kern0pt}OF\ i{\isacharbrackright}{\kern0pt}\ \isacommand{by}\isamarkupfalse%
\ auto\isanewline
\ \ \ \ \ \isacommand{with}\isamarkupfalse%
\ Abst{\isacharparenleft}{\kern0pt}{\isadigit{1}}{\isacharparenright}{\kern0pt}{\isacharbrackleft}{\kern0pt}OF\ Abst{\isacharparenleft}{\kern0pt}{\isadigit{2}}{\isacharparenright}{\kern0pt}{\isacharcomma}{\kern0pt}\ of\ {\isacartoucheopen}swap\ {\isacharbrackleft}{\kern0pt}{\isacharparenleft}{\kern0pt}a{\isacharcomma}{\kern0pt}b{\isacharparenright}{\kern0pt}{\isacharbrackright}{\kern0pt}\ t{\isadigit{2}}{\isacharprime}{\kern0pt}{\isacartoucheclose}{\isacharbrackright}{\kern0pt}\isanewline
\ \ \ \ \isacommand{have}\isamarkupfalse%
\ {\isachardoublequoteopen}nabla\ {\isasymturnstile}\ subst\ s{\isadigit{1}}\ t{\isadigit{1}}{\isacharprime}{\kern0pt}\ {\isasymapprox}\ \ swap\ {\isacharbrackleft}{\kern0pt}{\isacharparenleft}{\kern0pt}a{\isacharcomma}{\kern0pt}b{\isacharparenright}{\kern0pt}{\isacharbrackright}{\kern0pt}\ t{\isadigit{2}}{\isacharprime}{\kern0pt}{\isachardoublequoteclose}\ \isanewline
\ \ \ \ \ \ \isacommand{by}\isamarkupfalse%
\ simp\isanewline
\ \ \ \ \isacommand{hence}\isamarkupfalse%
\ {\isachardoublequoteopen}nabla\ {\isasymturnstile}\ Abst\ a\ {\isacharparenleft}{\kern0pt}subst\ s{\isadigit{1}}\ t{\isadigit{1}}{\isacharprime}{\kern0pt}{\isacharparenright}{\kern0pt}\ {\isasymapprox}\ Abst\ b\ t{\isadigit{2}}{\isacharprime}{\kern0pt}{\isachardoublequoteclose}\isanewline
\ \ \ \ \ \ \isacommand{using}\isamarkupfalse%
\ equ{\isacharunderscore}{\kern0pt}abst{\isacharunderscore}{\kern0pt}ab{\isacharbrackleft}{\kern0pt}OF\ False\ fresh{\isacharbrackright}{\kern0pt}\ \isacommand{by}\isamarkupfalse%
\ simp\isanewline
\ \ \ \ \isacommand{with}\isamarkupfalse%
\ t{\isadigit{2}}\ \isacommand{show}\isamarkupfalse%
\ {\isacharquery}{\kern0pt}thesis\isanewline
\ \ \ \ \ \ \isacommand{using}\isamarkupfalse%
\ subst{\isacharunderscore}{\kern0pt}abst{\isacharbrackleft}{\kern0pt}of\ s{\isadigit{1}}\ a\ t{\isadigit{1}}{\isacharprime}{\kern0pt}{\isacharbrackright}{\kern0pt}\ \isacommand{by}\isamarkupfalse%
\ simp\isanewline
\ \ \isacommand{qed}\isamarkupfalse%
\isanewline
\isacommand{next}\isamarkupfalse%
\isanewline
\ \ \isacommand{case}\isamarkupfalse%
\ {\isacharparenleft}{\kern0pt}Susp\ pi\ X{\isacharparenright}{\kern0pt}\isanewline
\ \ \isacommand{hence}\isamarkupfalse%
\ {\isachardoublequoteopen}nabla\ {\isasymturnstile}\ swap\ pi\ {\isacharparenleft}{\kern0pt}look{\isacharunderscore}{\kern0pt}up\ X\ s{\isadigit{2}}{\isacharparenright}{\kern0pt}\ {\isasymapprox}\ t{\isadigit{2}}{\isachardoublequoteclose}\isanewline
\ \ \ \ \isacommand{using}\isamarkupfalse%
\ subst{\isacharunderscore}{\kern0pt}susp\ \isacommand{by}\isamarkupfalse%
\ simp\isanewline
\ \ \isacommand{hence}\isamarkupfalse%
\ {\isachardoublequoteopen}nabla\ {\isasymturnstile}\ {\isacharparenleft}{\kern0pt}look{\isacharunderscore}{\kern0pt}up\ X\ s{\isadigit{2}}{\isacharparenright}{\kern0pt}\ {\isasymapprox}\ swap\ {\isacharparenleft}{\kern0pt}rev\ pi{\isacharparenright}{\kern0pt}\ t{\isadigit{2}}{\isachardoublequoteclose}\isanewline
\ \ \ \ \isacommand{using}\isamarkupfalse%
\ subst{\isacharunderscore}{\kern0pt}susp\ swap{\isacharunderscore}{\kern0pt}inv{\isacharunderscore}{\kern0pt}side\ \isacommand{by}\isamarkupfalse%
\ simp\isanewline
\ \ \isacommand{hence}\isamarkupfalse%
\ i{\isacharcolon}{\kern0pt}\ {\isachardoublequoteopen}nabla\ {\isasymturnstile}\ subst\ s{\isadigit{2}}\ {\isacharparenleft}{\kern0pt}Susp\ {\isacharbrackleft}{\kern0pt}{\isacharbrackright}{\kern0pt}\ X{\isacharparenright}{\kern0pt}\ {\isasymapprox}\ swap\ {\isacharparenleft}{\kern0pt}rev\ pi{\isacharparenright}{\kern0pt}\ t{\isadigit{2}}{\isachardoublequoteclose}\isanewline
\ \ \ \ \isacommand{using}\isamarkupfalse%
\ subst{\isacharunderscore}{\kern0pt}susp{\isacharbrackleft}{\kern0pt}of\ s{\isadigit{2}}\ {\isacartoucheopen}{\isacharbrackleft}{\kern0pt}{\isacharbrackright}{\kern0pt}{\isacartoucheclose}{\isacharbrackright}{\kern0pt}\ \isacommand{by}\isamarkupfalse%
\ simp\isanewline
\isanewline
\ \ \isacommand{with}\isamarkupfalse%
\ Susp{\isacharparenleft}{\kern0pt}{\isadigit{1}}{\isacharparenright}{\kern0pt}\ subst{\isacharunderscore}{\kern0pt}equ{\isacharunderscore}{\kern0pt}def{\isacharbrackleft}{\kern0pt}of\ nabla\ {\isacartoucheopen}subst\ s{\isadigit{1}}{\isacartoucheclose}\ {\isacartoucheopen}subst\ s{\isadigit{2}}{\isacartoucheclose}{\isacharbrackright}{\kern0pt}\isanewline
\ \ \isacommand{have}\isamarkupfalse%
\ {\isachardoublequoteopen}nabla\ {\isasymturnstile}\ subst\ s{\isadigit{1}}\ {\isacharparenleft}{\kern0pt}Susp\ {\isacharbrackleft}{\kern0pt}{\isacharbrackright}{\kern0pt}\ X{\isacharparenright}{\kern0pt}\ {\isasymapprox}\ subst\ s{\isadigit{2}}\ {\isacharparenleft}{\kern0pt}Susp\ {\isacharbrackleft}{\kern0pt}{\isacharbrackright}{\kern0pt}\ X{\isacharparenright}{\kern0pt}{\isachardoublequoteclose}\isanewline
\ \ \ \ \isacommand{using}\isamarkupfalse%
\ UnCI\ equ{\isacharunderscore}{\kern0pt}refl\ not{\isacharunderscore}{\kern0pt}in{\isacharunderscore}{\kern0pt}domn\ \isacommand{by}\isamarkupfalse%
\ metis\isanewline
\isanewline
\ \ \isacommand{with}\isamarkupfalse%
\ i\ \isacommand{have}\isamarkupfalse%
\ {\isachardoublequoteopen}nabla\ {\isasymturnstile}\ subst\ s{\isadigit{1}}\ {\isacharparenleft}{\kern0pt}Susp\ {\isacharbrackleft}{\kern0pt}{\isacharbrackright}{\kern0pt}\ X{\isacharparenright}{\kern0pt}\ {\isasymapprox}\ swap\ {\isacharparenleft}{\kern0pt}rev\ pi{\isacharparenright}{\kern0pt}\ t{\isadigit{2}}{\isachardoublequoteclose}\isanewline
\ \ \ \ \isacommand{using}\isamarkupfalse%
\ equ{\isacharunderscore}{\kern0pt}trans\ \isacommand{by}\isamarkupfalse%
\ blast\isanewline
\ \ \isacommand{hence}\isamarkupfalse%
\ {\isachardoublequoteopen}nabla\ {\isasymturnstile}\ {\isacharparenleft}{\kern0pt}look{\isacharunderscore}{\kern0pt}up\ X\ s{\isadigit{1}}{\isacharparenright}{\kern0pt}\ {\isasymapprox}\ swap\ {\isacharparenleft}{\kern0pt}rev\ pi{\isacharparenright}{\kern0pt}\ t{\isadigit{2}}{\isachardoublequoteclose}\isanewline
\ \ \ \ \isacommand{using}\isamarkupfalse%
\ subst{\isacharunderscore}{\kern0pt}susp\ \isacommand{by}\isamarkupfalse%
\ simp\isanewline
\ \ \isacommand{hence}\isamarkupfalse%
\ {\isachardoublequoteopen}nabla\ {\isasymturnstile}\ swap\ pi\ {\isacharparenleft}{\kern0pt}look{\isacharunderscore}{\kern0pt}up\ X\ s{\isadigit{1}}{\isacharparenright}{\kern0pt}\ {\isasymapprox}\ \ t{\isadigit{2}}{\isachardoublequoteclose}\isanewline
\ \ \ \ \isacommand{using}\isamarkupfalse%
\ swap{\isacharunderscore}{\kern0pt}inv{\isacharunderscore}{\kern0pt}side\ \isacommand{by}\isamarkupfalse%
\ simp\isanewline
\ \ \ \ \isacommand{then}\isamarkupfalse%
\ \isacommand{show}\isamarkupfalse%
\ {\isacharquery}{\kern0pt}case\ \isanewline
\ \ \ \ \ \ \isacommand{using}\isamarkupfalse%
\ subst{\isacharunderscore}{\kern0pt}susp\ \isacommand{by}\isamarkupfalse%
\ simp\isanewline
\isacommand{next}\isamarkupfalse%
\isanewline
\ \ \isacommand{case}\isamarkupfalse%
\ {\isacharparenleft}{\kern0pt}Paar\ t{\isadigit{1}}{\isadigit{1}}\ t{\isadigit{1}}{\isadigit{2}}{\isacharparenright}{\kern0pt}\isanewline
\ \ \isacommand{then}\isamarkupfalse%
\ \isacommand{obtain}\isamarkupfalse%
\ t{\isadigit{2}}{\isadigit{1}}\ t{\isadigit{2}}{\isadigit{2}}\isanewline
\ \ \ \ \isakeyword{where}\ t{\isadigit{2}}{\isacharcolon}{\kern0pt}\ {\isachardoublequoteopen}t{\isadigit{2}}\ {\isacharequal}{\kern0pt}\ Paar\ t{\isadigit{2}}{\isadigit{1}}\ t{\isadigit{2}}{\isadigit{2}}{\isachardoublequoteclose}\isanewline
\ \ \ \ \isacommand{by}\isamarkupfalse%
{\isacharparenleft}{\kern0pt}auto\ elim{\isacharcolon}{\kern0pt}\ equ{\isachardot}{\kern0pt}cases{\isacharparenright}{\kern0pt}\isanewline
\ \ \isacommand{with}\isamarkupfalse%
\ Paar{\isacharparenleft}{\kern0pt}{\isadigit{4}}{\isacharparenright}{\kern0pt}\ \isanewline
\ \ \isacommand{have}\isamarkupfalse%
\ paar{\isacharunderscore}{\kern0pt}i{\isacharcolon}{\kern0pt}\ {\isachardoublequoteopen}nabla\ {\isasymturnstile}\ subst\ s{\isadigit{2}}\ t{\isadigit{1}}{\isadigit{1}}\ {\isasymapprox}\ t{\isadigit{2}}{\isadigit{1}}{\isachardoublequoteclose}\isanewline
\ \ \ \ \isakeyword{and}\ paar{\isacharunderscore}{\kern0pt}ii{\isacharcolon}{\kern0pt}\ {\isachardoublequoteopen}nabla\ {\isasymturnstile}\ subst\ s{\isadigit{2}}\ t{\isadigit{1}}{\isadigit{2}}\ {\isasymapprox}\ t{\isadigit{2}}{\isadigit{2}}{\isachardoublequoteclose}\isanewline
\ \ \ \ \isacommand{by}\isamarkupfalse%
\ {\isacharparenleft}{\kern0pt}auto\ elim{\isacharcolon}{\kern0pt}\ equ{\isachardot}{\kern0pt}cases{\isacharparenright}{\kern0pt}\isanewline
\ \ \isacommand{from}\isamarkupfalse%
\ t{\isadigit{2}}\ \isacommand{show}\isamarkupfalse%
\ {\isacharquery}{\kern0pt}case\isanewline
\ \ \ \ \isacommand{using}\isamarkupfalse%
\ equ{\isacharunderscore}{\kern0pt}paar{\isacharbrackleft}{\kern0pt}OF\ Paar{\isacharparenleft}{\kern0pt}{\isadigit{1}}{\isacharparenright}{\kern0pt}{\isacharbrackleft}{\kern0pt}OF\ Paar{\isacharparenleft}{\kern0pt}{\isadigit{3}}{\isacharparenright}{\kern0pt}\ paar{\isacharunderscore}{\kern0pt}i{\isacharbrackright}{\kern0pt}\ Paar{\isacharparenleft}{\kern0pt}{\isadigit{2}}{\isacharparenright}{\kern0pt}{\isacharbrackleft}{\kern0pt}OF\ Paar{\isacharparenleft}{\kern0pt}{\isadigit{3}}{\isacharparenright}{\kern0pt}\ paar{\isacharunderscore}{\kern0pt}ii{\isacharbrackright}{\kern0pt}{\isacharbrackright}{\kern0pt}\isanewline
\ \ \ \ subst{\isacharunderscore}{\kern0pt}paar\ \isacommand{by}\isamarkupfalse%
\ simp\isanewline
\isacommand{next}\isamarkupfalse%
\isanewline
\ \ \isacommand{case}\isamarkupfalse%
\ {\isacharparenleft}{\kern0pt}Func\ f\ t{\isadigit{1}}{\isacharprime}{\kern0pt}{\isacharparenright}{\kern0pt}\isanewline
\ \ \isacommand{then}\isamarkupfalse%
\ \isacommand{obtain}\isamarkupfalse%
\ t{\isadigit{2}}{\isacharprime}{\kern0pt}\ \isanewline
\ \ \ \ \isakeyword{where}\ t{\isadigit{2}}{\isacharcolon}{\kern0pt}\ {\isachardoublequoteopen}t{\isadigit{2}}\ {\isacharequal}{\kern0pt}\ Func\ f\ t{\isadigit{2}}{\isacharprime}{\kern0pt}{\isachardoublequoteclose}\isanewline
\ \ \ \ \isacommand{by}\isamarkupfalse%
{\isacharparenleft}{\kern0pt}auto\ elim{\isacharcolon}{\kern0pt}\ equ{\isachardot}{\kern0pt}cases{\isacharparenright}{\kern0pt}\isanewline
\ \ \ \isacommand{with}\isamarkupfalse%
\ Func{\isacharparenleft}{\kern0pt}{\isadigit{3}}{\isacharparenright}{\kern0pt}\ \isanewline
\ \ \ \isacommand{have}\isamarkupfalse%
\ func{\isacharcolon}{\kern0pt}\ {\isachardoublequoteopen}nabla\ {\isasymturnstile}\ subst\ s{\isadigit{2}}\ t{\isadigit{1}}{\isacharprime}{\kern0pt}\ {\isasymapprox}\ t{\isadigit{2}}{\isacharprime}{\kern0pt}{\isachardoublequoteclose}\isanewline
\ \ \ \isacommand{by}\isamarkupfalse%
\ {\isacharparenleft}{\kern0pt}auto\ elim{\isacharcolon}{\kern0pt}\ equ{\isachardot}{\kern0pt}cases{\isacharparenright}{\kern0pt}\isanewline
\ \ \isacommand{from}\isamarkupfalse%
\ t{\isadigit{2}}\ \isacommand{show}\isamarkupfalse%
\ {\isacharquery}{\kern0pt}case\isanewline
\ \ \ \ \isacommand{using}\isamarkupfalse%
\ equ{\isacharunderscore}{\kern0pt}func{\isacharbrackleft}{\kern0pt}OF\ Func{\isacharparenleft}{\kern0pt}{\isadigit{1}}{\isacharparenright}{\kern0pt}{\isacharbrackleft}{\kern0pt}OF\ Func{\isacharparenleft}{\kern0pt}{\isadigit{2}}{\isacharparenright}{\kern0pt}\ func{\isacharbrackright}{\kern0pt}{\isacharcomma}{\kern0pt}\ of\ f{\isacharbrackright}{\kern0pt}\ \isacommand{by}\isamarkupfalse%
\ simp\isanewline
\isacommand{qed}\isamarkupfalse%
\ {\isacharparenleft}{\kern0pt}simp{\isacharunderscore}{\kern0pt}all{\isacharparenright}{\kern0pt}%
\endisatagproof
{\isafoldproof}%
%
\isadelimproof
\isanewline
%
\endisadelimproof
\isanewline
\isanewline
\isanewline
\isacommand{lemma}\isamarkupfalse%
\ unif{\isacharunderscore}{\kern0pt}{\isadigit{2}}a{\isacharcolon}{\kern0pt}\ \isanewline
\ \ \isakeyword{assumes}\ {\isachardoublequoteopen}nabla{\isasymTurnstile}subst\ s{\isadigit{1}}{\isasymapprox}subst\ s{\isadigit{2}}{\isachardoublequoteclose}\isanewline
\ \ \isakeyword{and}\ {\isachardoublequoteopen}{\isacharparenleft}{\kern0pt}nabla{\isasymturnstile}subst\ s{\isadigit{2}}\ t{\isadigit{1}}\ {\isasymapprox}\ subst\ s{\isadigit{2}}\ t{\isadigit{2}}{\isacharparenright}{\kern0pt}{\isachardoublequoteclose}\ \ \ \ \ \ \ \ \ \ \ \ \ \isanewline
\isakeyword{shows}{\isachardoublequoteopen}{\isacharparenleft}{\kern0pt}nabla{\isasymturnstile}subst\ s{\isadigit{1}}\ t{\isadigit{1}}\ {\isasymapprox}\ subst\ s{\isadigit{1}}\ t{\isadigit{2}}{\isacharparenright}{\kern0pt}{\isachardoublequoteclose}\isanewline
%
\isadelimproof
%
\endisadelimproof
%
\isatagproof
\isacommand{proof}\isamarkupfalse%
{\isacharminus}{\kern0pt}\isanewline
\ \ \isacommand{have}\isamarkupfalse%
\ i{\isacharcolon}{\kern0pt}\ {\isachardoublequoteopen}{\isacharparenleft}{\kern0pt}nabla{\isasymturnstile}subst\ s{\isadigit{1}}\ t{\isadigit{1}}\ {\isasymapprox}\ subst\ s{\isadigit{2}}\ t{\isadigit{2}}{\isacharparenright}{\kern0pt}{\isachardoublequoteclose}\isanewline
\ \ \ \ \isacommand{using}\isamarkupfalse%
\ subst{\isacharunderscore}{\kern0pt}equ{\isacharunderscore}{\kern0pt}a{\isacharbrackleft}{\kern0pt}OF\ assms{\isacharparenleft}{\kern0pt}{\isadigit{1}}{\isacharcomma}{\kern0pt}{\isadigit{2}}{\isacharparenright}{\kern0pt}{\isacharbrackright}{\kern0pt}\ \isacommand{by}\isamarkupfalse%
\ simp\isanewline
\ \ \isacommand{show}\isamarkupfalse%
\ {\isachardoublequoteopen}{\isacharparenleft}{\kern0pt}nabla{\isasymturnstile}subst\ s{\isadigit{1}}\ t{\isadigit{1}}\ {\isasymapprox}\ subst\ s{\isadigit{1}}\ t{\isadigit{2}}{\isacharparenright}{\kern0pt}{\isachardoublequoteclose}\isanewline
\ \ \ \ \isacommand{using}\isamarkupfalse%
\ subst{\isacharunderscore}{\kern0pt}equ{\isacharunderscore}{\kern0pt}a{\isacharbrackleft}{\kern0pt}OF\ assms{\isacharparenleft}{\kern0pt}{\isadigit{1}}{\isacharparenright}{\kern0pt}\ equ{\isacharunderscore}{\kern0pt}symm{\isacharbrackleft}{\kern0pt}OF\ i{\isacharbrackright}{\kern0pt}{\isacharbrackright}{\kern0pt}\ equ{\isacharunderscore}{\kern0pt}symm\ \isacommand{by}\isamarkupfalse%
\ auto\isanewline
\isacommand{qed}\isamarkupfalse%
%
\endisatagproof
{\isafoldproof}%
%
\isadelimproof
\isanewline
%
\endisadelimproof
\isanewline
\isacommand{lemma}\isamarkupfalse%
\ unif{\isacharunderscore}{\kern0pt}{\isadigit{2}}b{\isacharcolon}{\kern0pt}\isanewline
\ \ \isakeyword{assumes}\ {\isachardoublequoteopen}nabla\ {\isasymTurnstile}\ subst\ s{\isadigit{1}}\ {\isasymapprox}\ subst\ s{\isadigit{2}}{\isachardoublequoteclose}\ \isanewline
\ \ \isakeyword{and}\ {\isachardoublequoteopen}nabla\ {\isasymturnstile}\ a\ {\isasymsharp}\ subst\ s{\isadigit{2}}\ t{\isachardoublequoteclose}\isanewline
\isakeyword{shows}\ {\isachardoublequoteopen}nabla\ {\isasymturnstile}\ a\ {\isasymsharp}\ subst\ s{\isadigit{1}}\ t{\isachardoublequoteclose}\ \isanewline
%
\isadelimproof
\ \ %
\endisadelimproof
%
\isatagproof
\isacommand{using}\isamarkupfalse%
\ assms\isanewline
\isacommand{proof}\isamarkupfalse%
{\isacharparenleft}{\kern0pt}induct\ t{\isacharparenright}{\kern0pt}\isanewline
\ \ \isacommand{case}\isamarkupfalse%
\ {\isacharparenleft}{\kern0pt}Abst\ b\ t{\isacharprime}{\kern0pt}{\isacharparenright}{\kern0pt}\isanewline
\ \ \isacommand{then}\isamarkupfalse%
\ \isacommand{show}\isamarkupfalse%
\ {\isacharquery}{\kern0pt}case\isanewline
\ \ \isacommand{proof}\isamarkupfalse%
{\isacharparenleft}{\kern0pt}cases\ {\isachardoublequoteopen}a{\isacharequal}{\kern0pt}b{\isachardoublequoteclose}{\isacharparenright}{\kern0pt}\isanewline
\ \ \ \ \isacommand{case}\isamarkupfalse%
\ True\isanewline
\ \ \ \ \isacommand{then}\isamarkupfalse%
\ \isacommand{show}\isamarkupfalse%
\ {\isacharquery}{\kern0pt}thesis\ \isanewline
\ \ \ \ \ \ \isacommand{using}\isamarkupfalse%
\ fresh{\isacharunderscore}{\kern0pt}abst{\isacharunderscore}{\kern0pt}aa\ \isacommand{by}\isamarkupfalse%
\ simp\ \isanewline
\ \ \isacommand{next}\isamarkupfalse%
\isanewline
\ \ \ \ \isacommand{case}\isamarkupfalse%
\ False\isanewline
\ \ \ \ \isacommand{with}\isamarkupfalse%
\ Abst\ \isacommand{show}\isamarkupfalse%
\ {\isacharquery}{\kern0pt}thesis\ \isanewline
\ \ \ \ \ \ \isacommand{using}\isamarkupfalse%
\ Fresh{\isacharunderscore}{\kern0pt}elims{\isacharparenleft}{\kern0pt}{\isadigit{1}}{\isacharparenright}{\kern0pt}\ Substs{\isachardot}{\kern0pt}subst{\isacharunderscore}{\kern0pt}abst\ fresh{\isacharunderscore}{\kern0pt}abst{\isacharunderscore}{\kern0pt}ab\isanewline
\ \ \ \ \ \ \isacommand{by}\isamarkupfalse%
\ metis\isanewline
\ \ \isacommand{qed}\isamarkupfalse%
\isanewline
\isacommand{next}\isamarkupfalse%
\isanewline
\ \ \isacommand{case}\isamarkupfalse%
\ {\isacharparenleft}{\kern0pt}Susp\ pi\ X{\isacharparenright}{\kern0pt}\isanewline
\ \ \isacommand{then}\isamarkupfalse%
\ \isacommand{show}\isamarkupfalse%
\ {\isacharquery}{\kern0pt}case\isanewline
\ \ \ \ \isacommand{using}\isamarkupfalse%
\ equ{\isacharunderscore}{\kern0pt}refl\ equ{\isacharunderscore}{\kern0pt}symm\ l{\isadigit{3}}{\isacharunderscore}{\kern0pt}jud\ subst{\isacharunderscore}{\kern0pt}equ{\isacharunderscore}{\kern0pt}a\ \isacommand{by}\isamarkupfalse%
\ meson\isanewline
\isacommand{next}\isamarkupfalse%
\isanewline
\ \ \isacommand{case}\isamarkupfalse%
\ {\isacharparenleft}{\kern0pt}Paar\ t{\isadigit{1}}\ t{\isadigit{2}}{\isacharparenright}{\kern0pt}\isanewline
\ \ \isacommand{then}\isamarkupfalse%
\ \isacommand{show}\isamarkupfalse%
\ {\isacharquery}{\kern0pt}case\ \isanewline
\ \ \ \ \isacommand{using}\isamarkupfalse%
\ Fresh{\isacharunderscore}{\kern0pt}elims{\isacharparenleft}{\kern0pt}{\isadigit{5}}{\isacharparenright}{\kern0pt}\ subst{\isacharunderscore}{\kern0pt}paar\ fresh{\isacharunderscore}{\kern0pt}paar\ \isanewline
\ \ \ \ \isacommand{by}\isamarkupfalse%
\ metis\isanewline
\isacommand{next}\isamarkupfalse%
\isanewline
\ \ \isacommand{case}\isamarkupfalse%
\ {\isacharparenleft}{\kern0pt}Func\ f\ t{\isacharprime}{\kern0pt}{\isacharparenright}{\kern0pt}\isanewline
\ \ \isacommand{then}\isamarkupfalse%
\ \isacommand{show}\isamarkupfalse%
\ {\isacharquery}{\kern0pt}case\isanewline
\ \ \ \ \isacommand{using}\isamarkupfalse%
\ Fresh{\isacharunderscore}{\kern0pt}elims{\isacharparenleft}{\kern0pt}{\isadigit{6}}{\isacharparenright}{\kern0pt}\ subst{\isacharunderscore}{\kern0pt}func\ fresh{\isacharunderscore}{\kern0pt}func\isanewline
\ \ \ \ \isacommand{by}\isamarkupfalse%
\ metis\isanewline
\isacommand{qed}\isamarkupfalse%
\ {\isacharparenleft}{\kern0pt}simp{\isacharunderscore}{\kern0pt}all{\isacharparenright}{\kern0pt}%
\endisatagproof
{\isafoldproof}%
%
\isadelimproof
\isanewline
%
\endisadelimproof
\isanewline
\isacommand{lemma}\isamarkupfalse%
\ subst{\isacharunderscore}{\kern0pt}equ{\isacharunderscore}{\kern0pt}to{\isacharunderscore}{\kern0pt}trm{\isacharcolon}{\kern0pt}\ {\isachardoublequoteopen}nabla\ {\isasymTurnstile}\ subst\ s{\isadigit{1}}\ {\isasymapprox}\ subst\ s{\isadigit{2}}\ {\isasymLongrightarrow}\ nabla{\isasymturnstile}subst\ s{\isadigit{1}}\ t\ {\isasymapprox}\ subst\ s{\isadigit{2}}\ t{\isachardoublequoteclose}\isanewline
%
\isadelimproof
\ \ %
\endisadelimproof
%
\isatagproof
\isacommand{using}\isamarkupfalse%
\ equ{\isacharunderscore}{\kern0pt}refl\ subst{\isacharunderscore}{\kern0pt}equ{\isacharunderscore}{\kern0pt}a\ \isacommand{by}\isamarkupfalse%
\ simp%
\endisatagproof
{\isafoldproof}%
%
\isadelimproof
\isanewline
%
\endisadelimproof
\isanewline
\isanewline
\isacommand{lemma}\isamarkupfalse%
\ subst{\isacharunderscore}{\kern0pt}cancel{\isacharunderscore}{\kern0pt}right{\isacharcolon}{\kern0pt}\ \isanewline
{\isachardoublequoteopen}{\isasymlbrakk}nabla{\isasymTurnstile}{\isacharparenleft}{\kern0pt}subst\ s{\isadigit{1}}{\isacharparenright}{\kern0pt}{\isasymapprox}{\isacharparenleft}{\kern0pt}subst\ s{\isadigit{2}}{\isacharparenright}{\kern0pt}{\isasymrbrakk}{\isasymLongrightarrow}nabla{\isasymTurnstile}{\isacharparenleft}{\kern0pt}subst\ {\isacharparenleft}{\kern0pt}s{\isadigit{1}}{\isasymbullet}s{\isacharparenright}{\kern0pt}{\isacharparenright}{\kern0pt}{\isasymapprox}{\isacharparenleft}{\kern0pt}subst\ {\isacharparenleft}{\kern0pt}s{\isadigit{2}}{\isasymbullet}s{\isacharparenright}{\kern0pt}{\isacharparenright}{\kern0pt}{\isachardoublequoteclose}\ \isanewline
%
\isadelimproof
\ \ %
\endisadelimproof
%
\isatagproof
\isacommand{using}\isamarkupfalse%
\ subst{\isacharunderscore}{\kern0pt}comp{\isacharunderscore}{\kern0pt}expand\ subst{\isacharunderscore}{\kern0pt}equ{\isacharunderscore}{\kern0pt}def\ subst{\isacharunderscore}{\kern0pt}equ{\isacharunderscore}{\kern0pt}to{\isacharunderscore}{\kern0pt}trm\isanewline
\ \ \isacommand{by}\isamarkupfalse%
\ simp%
\endisatagproof
{\isafoldproof}%
%
\isadelimproof
\isanewline
%
\endisadelimproof
\isanewline
\isacommand{lemma}\isamarkupfalse%
\ subst{\isacharunderscore}{\kern0pt}trans{\isacharcolon}{\kern0pt}\ {\isachardoublequoteopen}{\isasymlbrakk}nabla{\isasymTurnstile}subst\ s{\isadigit{1}}\ {\isasymapprox}\ subst\ s{\isadigit{2}}{\isacharsemicolon}{\kern0pt}\ nabla\ {\isasymTurnstile}\ subst\ s{\isadigit{2}}\ {\isasymapprox}\ subst\ s{\isadigit{3}}{\isasymrbrakk}{\isasymLongrightarrow}nabla{\isasymTurnstile}subst\ s{\isadigit{1}}{\isasymapprox}subst\ s{\isadigit{3}}{\isachardoublequoteclose}\isanewline
%
\isadelimproof
\ \ %
\endisadelimproof
%
\isatagproof
\isacommand{using}\isamarkupfalse%
\ equ{\isacharunderscore}{\kern0pt}trans\ subst{\isacharunderscore}{\kern0pt}equ{\isacharunderscore}{\kern0pt}def\ subst{\isacharunderscore}{\kern0pt}equ{\isacharunderscore}{\kern0pt}to{\isacharunderscore}{\kern0pt}trm\ \isacommand{by}\isamarkupfalse%
\ meson%
\endisatagproof
{\isafoldproof}%
%
\isadelimproof
\isanewline
%
\endisadelimproof
\isanewline
\isanewline
\isanewline
\isanewline
\isacommand{lemma}\isamarkupfalse%
\ occurs{\isacharunderscore}{\kern0pt}sub{\isacharunderscore}{\kern0pt}trm{\isacharunderscore}{\kern0pt}equ{\isacharcolon}{\kern0pt}\ \isanewline
\ \ {\isachardoublequoteopen}occurs\ X\ t{\isadigit{1}}\ {\isasymLongrightarrow}\ {\isacharparenleft}{\kern0pt}{\isasymexists}\ t{\isadigit{2}}\ {\isasymin}\ sub{\isacharunderscore}{\kern0pt}trms\ {\isacharparenleft}{\kern0pt}subst\ s\ t{\isadigit{1}}{\isacharparenright}{\kern0pt}{\isachardot}{\kern0pt}\ {\isacharparenleft}{\kern0pt}{\isasymexists}pi{\isachardot}{\kern0pt}{\isacharparenleft}{\kern0pt}nabla{\isasymturnstile}subst\ s\ {\isacharparenleft}{\kern0pt}Susp\ pi{\isadigit{1}}\ X{\isacharparenright}{\kern0pt}{\isasymapprox}{\isacharparenleft}{\kern0pt}swap\ pi\ t{\isadigit{2}}{\isacharparenright}{\kern0pt}{\isacharparenright}{\kern0pt}{\isacharparenright}{\kern0pt}{\isacharparenright}{\kern0pt}{\isachardoublequoteclose}\ \isanewline
%
\isadelimproof
%
\endisadelimproof
%
\isatagproof
\isacommand{proof}\isamarkupfalse%
{\isacharparenleft}{\kern0pt}induct\ t{\isadigit{1}}{\isacharparenright}{\kern0pt}\isanewline
\ \ \isacommand{case}\isamarkupfalse%
\ {\isacharparenleft}{\kern0pt}Susp\ pi{\isacharprime}{\kern0pt}\ X{\isacharparenright}{\kern0pt}\isanewline
\ \ \isacommand{then}\isamarkupfalse%
\ \isacommand{show}\isamarkupfalse%
\ {\isacharquery}{\kern0pt}case\ \isanewline
\ \ \ \ \isacommand{using}\isamarkupfalse%
\ equ{\isacharunderscore}{\kern0pt}pi{\isacharunderscore}{\kern0pt}to{\isacharunderscore}{\kern0pt}left\ equ{\isacharunderscore}{\kern0pt}involutive{\isacharunderscore}{\kern0pt}left\ \isanewline
\ \ \ \ \ \ equ{\isacharunderscore}{\kern0pt}refl\ subst{\isacharunderscore}{\kern0pt}susp\ swap{\isacharunderscore}{\kern0pt}append\ t{\isacharunderscore}{\kern0pt}sub{\isacharunderscore}{\kern0pt}trms{\isacharunderscore}{\kern0pt}t\isanewline
\ \ \ \ occurs{\isachardot}{\kern0pt}simps{\isacharparenleft}{\kern0pt}{\isadigit{3}}{\isacharparenright}{\kern0pt}\ \isacommand{by}\isamarkupfalse%
\ metis\isanewline
\isacommand{next}\isamarkupfalse%
\isanewline
\ \ \isacommand{case}\isamarkupfalse%
\ {\isacharparenleft}{\kern0pt}Paar\ t{\isadigit{1}}{\isadigit{1}}\ t{\isadigit{1}}{\isadigit{2}}{\isacharparenright}{\kern0pt}\isanewline
\ \ \isacommand{then}\isamarkupfalse%
\ \isacommand{show}\isamarkupfalse%
\ {\isacharquery}{\kern0pt}case\ \isanewline
\ \ \ \ \isacommand{using}\isamarkupfalse%
\ Substs{\isachardot}{\kern0pt}subst{\isacharunderscore}{\kern0pt}paar\ Un{\isacharunderscore}{\kern0pt}iff\ occurs{\isachardot}{\kern0pt}simps{\isacharparenleft}{\kern0pt}{\isadigit{5}}{\isacharparenright}{\kern0pt}\ psub{\isacharunderscore}{\kern0pt}sub{\isacharunderscore}{\kern0pt}trms\isanewline
\ \ \ \ \ \ \ \ psub{\isacharunderscore}{\kern0pt}trms{\isachardot}{\kern0pt}simps{\isacharparenleft}{\kern0pt}{\isadigit{4}}{\isacharparenright}{\kern0pt}\ \isacommand{by}\isamarkupfalse%
\ {\isacharparenleft}{\kern0pt}metis{\isacharparenright}{\kern0pt}\isanewline
\isacommand{qed}\isamarkupfalse%
{\isacharparenleft}{\kern0pt}auto{\isacharparenright}{\kern0pt}%
\endisatagproof
{\isafoldproof}%
%
\isadelimproof
\isanewline
%
\endisadelimproof
\isanewline
\isacommand{lemma}\isamarkupfalse%
\ ext{\isacharunderscore}{\kern0pt}subst{\isacharunderscore}{\kern0pt}strong{\isacharcolon}{\kern0pt}\isanewline
\ \ \isakeyword{assumes}\ {\isachardoublequoteopen}nabla{\isadigit{1}}\ {\isasymTurnstile}\ {\isacharparenleft}{\kern0pt}subst\ s{\isacharparenright}{\kern0pt}\ {\isacharparenleft}{\kern0pt}nabla{\isadigit{2}}\ {\isasymunion}\ nabla{\isadigit{3}}{\isacharparenright}{\kern0pt}{\isachardoublequoteclose}\isanewline
\ \ \isakeyword{shows}\ {\isachardoublequoteopen}nabla{\isadigit{1}}\ {\isasymTurnstile}\ {\isacharparenleft}{\kern0pt}subst\ s{\isacharparenright}{\kern0pt}\ {\isacharparenleft}{\kern0pt}nabla{\isadigit{2}}{\isacharparenright}{\kern0pt}{\isachardoublequoteclose}\ \isakeyword{and}\ {\isachardoublequoteopen}nabla{\isadigit{1}}\ {\isasymTurnstile}\ {\isacharparenleft}{\kern0pt}subst\ s{\isacharparenright}{\kern0pt}\ {\isacharparenleft}{\kern0pt}nabla{\isadigit{3}}{\isacharparenright}{\kern0pt}{\isachardoublequoteclose}\isanewline
%
\isadelimproof
\ \ %
\endisadelimproof
%
\isatagproof
\isacommand{using}\isamarkupfalse%
\ assms\isanewline
\ \ \isacommand{unfolding}\isamarkupfalse%
\ ext{\isacharunderscore}{\kern0pt}subst{\isacharunderscore}{\kern0pt}def\ \isacommand{by}\isamarkupfalse%
\ auto%
\endisatagproof
{\isafoldproof}%
%
\isadelimproof
\isanewline
%
\endisadelimproof
\isanewline
\isacommand{lemma}\isamarkupfalse%
\ ext{\isacharunderscore}{\kern0pt}subst{\isacharunderscore}{\kern0pt}id{\isacharcolon}{\kern0pt}\isanewline
\ \ \isakeyword{shows}\ {\isachardoublequoteopen}nabla\ {\isasymTurnstile}\ {\isacharparenleft}{\kern0pt}subst\ {\isacharbrackleft}{\kern0pt}{\isacharbrackright}{\kern0pt}{\isacharparenright}{\kern0pt}\ nabla{\isachardoublequoteclose}\isanewline
%
\isadelimproof
\ \ %
\endisadelimproof
%
\isatagproof
\isacommand{unfolding}\isamarkupfalse%
\ ext{\isacharunderscore}{\kern0pt}subst{\isacharunderscore}{\kern0pt}def\ id{\isacharunderscore}{\kern0pt}subst\ \isacommand{by}\isamarkupfalse%
\ auto%
\endisatagproof
{\isafoldproof}%
%
\isadelimproof
\isanewline
%
\endisadelimproof
%
\isadelimtheory
\isanewline
%
\endisadelimtheory
%
\isatagtheory
\isacommand{end}\isamarkupfalse%
%
\endisatagtheory
{\isafoldtheory}%
%
\isadelimtheory
%
\endisadelimtheory
%
\end{isabellebody}%
\endinput
%:%file=~/nominal_unification_Isabelle/nominal_alpha_unification/Substs.thy%:%
%:%10=1%:%
%:%11=1%:%
%:%12=2%:%
%:%13=3%:%
%:%14=4%:%
%:%15=5%:%
%:%20=5%:%
%:%23=6%:%
%:%24=7%:%
%:%25=8%:%
%:%26=9%:%
%:%27=9%:%
%:%28=10%:%
%:%29=11%:%
%:%30=11%:%
%:%31=12%:%
%:%32=13%:%
%:%33=14%:%
%:%34=15%:%
%:%35=15%:%
%:%36=16%:%
%:%37=17%:%
%:%38=18%:%
%:%39=19%:%
%:%40=20%:%
%:%41=21%:%
%:%42=22%:%
%:%43=23%:%
%:%44=23%:%
%:%45=24%:%
%:%46=25%:%
%:%47=26%:%
%:%48=27%:%
%:%49=27%:%
%:%50=28%:%
%:%51=29%:%
%:%52=30%:%
%:%53=31%:%
%:%54=31%:%
%:%55=32%:%
%:%56=33%:%
%:%57=34%:%
%:%58=35%:%
%:%59=36%:%
%:%60=36%:%
%:%61=37%:%
%:%62=38%:%
%:%63=39%:%
%:%64=40%:%
%:%65=41%:%
%:%66=41%:%
%:%67=42%:%
%:%68=43%:%
%:%69=44%:%
%:%70=45%:%
%:%71=46%:%
%:%72=46%:%
%:%73=47%:%
%:%74=48%:%
%:%75=49%:%
%:%76=50%:%
%:%77=50%:%
%:%78=51%:%
%:%79=52%:%
%:%82=53%:%
%:%86=53%:%
%:%87=53%:%
%:%88=54%:%
%:%89=54%:%
%:%94=54%:%
%:%97=55%:%
%:%98=56%:%
%:%99=56%:%
%:%100=57%:%
%:%103=58%:%
%:%107=58%:%
%:%108=58%:%
%:%109=58%:%
%:%114=58%:%
%:%117=59%:%
%:%118=60%:%
%:%119=60%:%
%:%122=61%:%
%:%126=61%:%
%:%127=61%:%
%:%128=61%:%
%:%133=61%:%
%:%136=62%:%
%:%137=63%:%
%:%138=63%:%
%:%141=64%:%
%:%145=64%:%
%:%146=64%:%
%:%147=64%:%
%:%152=64%:%
%:%155=65%:%
%:%156=66%:%
%:%157=66%:%
%:%158=67%:%
%:%165=68%:%
%:%166=68%:%
%:%167=69%:%
%:%168=69%:%
%:%169=70%:%
%:%170=70%:%
%:%171=70%:%
%:%172=71%:%
%:%173=71%:%
%:%174=72%:%
%:%175=72%:%
%:%176=72%:%
%:%177=72%:%
%:%178=73%:%
%:%184=73%:%
%:%187=74%:%
%:%188=75%:%
%:%189=76%:%
%:%190=76%:%
%:%193=77%:%
%:%197=77%:%
%:%198=77%:%
%:%199=77%:%
%:%204=77%:%
%:%207=78%:%
%:%208=79%:%
%:%209=79%:%
%:%212=80%:%
%:%216=80%:%
%:%217=80%:%
%:%218=80%:%
%:%223=80%:%
%:%226=81%:%
%:%227=82%:%
%:%228=82%:%
%:%235=83%:%
%:%236=83%:%
%:%237=84%:%
%:%238=84%:%
%:%239=85%:%
%:%240=85%:%
%:%241=86%:%
%:%242=86%:%
%:%243=87%:%
%:%244=87%:%
%:%245=88%:%
%:%246=88%:%
%:%247=88%:%
%:%248=88%:%
%:%249=89%:%
%:%250=89%:%
%:%251=90%:%
%:%252=90%:%
%:%253=91%:%
%:%254=91%:%
%:%255=91%:%
%:%256=92%:%
%:%257=92%:%
%:%258=92%:%
%:%259=93%:%
%:%260=93%:%
%:%261=94%:%
%:%262=94%:%
%:%263=95%:%
%:%264=95%:%
%:%265=95%:%
%:%266=96%:%
%:%267=96%:%
%:%268=96%:%
%:%269=97%:%
%:%270=97%:%
%:%271=98%:%
%:%272=98%:%
%:%273=99%:%
%:%279=99%:%
%:%282=100%:%
%:%283=101%:%
%:%284=102%:%
%:%285=102%:%
%:%292=103%:%
%:%293=103%:%
%:%294=104%:%
%:%295=104%:%
%:%296=105%:%
%:%297=105%:%
%:%298=105%:%
%:%299=106%:%
%:%300=106%:%
%:%301=107%:%
%:%302=107%:%
%:%303=108%:%
%:%304=108%:%
%:%305=108%:%
%:%306=108%:%
%:%307=108%:%
%:%308=109%:%
%:%309=109%:%
%:%310=110%:%
%:%311=110%:%
%:%312=111%:%
%:%313=111%:%
%:%314=111%:%
%:%315=111%:%
%:%316=111%:%
%:%317=112%:%
%:%318=112%:%
%:%319=113%:%
%:%320=113%:%
%:%325=113%:%
%:%328=114%:%
%:%329=115%:%
%:%330=116%:%
%:%331=117%:%
%:%332=117%:%
%:%339=118%:%
%:%340=118%:%
%:%341=119%:%
%:%342=119%:%
%:%343=120%:%
%:%344=120%:%
%:%345=121%:%
%:%346=121%:%
%:%347=122%:%
%:%348=122%:%
%:%349=123%:%
%:%350=123%:%
%:%351=123%:%
%:%352=123%:%
%:%353=123%:%
%:%354=124%:%
%:%355=124%:%
%:%356=125%:%
%:%362=125%:%
%:%365=126%:%
%:%366=127%:%
%:%367=128%:%
%:%368=128%:%
%:%369=129%:%
%:%370=130%:%
%:%371=131%:%
%:%374=132%:%
%:%378=132%:%
%:%379=132%:%
%:%380=133%:%
%:%381=133%:%
%:%382=134%:%
%:%383=134%:%
%:%384=135%:%
%:%385=135%:%
%:%386=135%:%
%:%387=136%:%
%:%388=136%:%
%:%389=137%:%
%:%390=137%:%
%:%391=138%:%
%:%392=138%:%
%:%397=138%:%
%:%400=139%:%
%:%401=140%:%
%:%402=141%:%
%:%403=141%:%
%:%404=142%:%
%:%405=143%:%
%:%406=144%:%
%:%409=145%:%
%:%413=145%:%
%:%414=145%:%
%:%415=146%:%
%:%416=146%:%
%:%417=147%:%
%:%418=147%:%
%:%419=148%:%
%:%420=148%:%
%:%421=148%:%
%:%422=149%:%
%:%423=149%:%
%:%424=149%:%
%:%425=150%:%
%:%426=150%:%
%:%427=151%:%
%:%428=151%:%
%:%429=152%:%
%:%430=152%:%
%:%431=152%:%
%:%432=153%:%
%:%433=153%:%
%:%434=154%:%
%:%435=154%:%
%:%436=155%:%
%:%437=155%:%
%:%442=155%:%
%:%445=156%:%
%:%446=157%:%
%:%447=157%:%
%:%448=158%:%
%:%449=159%:%
%:%452=160%:%
%:%456=160%:%
%:%457=160%:%
%:%458=161%:%
%:%459=161%:%
%:%460=162%:%
%:%461=162%:%
%:%462=163%:%
%:%463=163%:%
%:%464=164%:%
%:%465=164%:%
%:%466=165%:%
%:%467=166%:%
%:%468=166%:%
%:%469=167%:%
%:%470=167%:%
%:%471=168%:%
%:%472=168%:%
%:%473=169%:%
%:%474=170%:%
%:%475=171%:%
%:%476=171%:%
%:%477=171%:%
%:%478=172%:%
%:%479=172%:%
%:%480=173%:%
%:%481=174%:%
%:%482=175%:%
%:%483=175%:%
%:%484=176%:%
%:%485=176%:%
%:%486=177%:%
%:%487=177%:%
%:%488=178%:%
%:%489=179%:%
%:%490=179%:%
%:%491=179%:%
%:%492=180%:%
%:%493=180%:%
%:%494=181%:%
%:%495=182%:%
%:%496=182%:%
%:%497=182%:%
%:%498=183%:%
%:%499=183%:%
%:%500=184%:%
%:%501=185%:%
%:%502=185%:%
%:%503=185%:%
%:%504=186%:%
%:%505=186%:%
%:%506=187%:%
%:%507=188%:%
%:%508=188%:%
%:%509=188%:%
%:%510=189%:%
%:%511=189%:%
%:%512=190%:%
%:%513=191%:%
%:%514=192%:%
%:%515=192%:%
%:%516=192%:%
%:%517=193%:%
%:%518=193%:%
%:%519=193%:%
%:%520=194%:%
%:%521=194%:%
%:%522=195%:%
%:%523=195%:%
%:%524=196%:%
%:%525=196%:%
%:%526=197%:%
%:%527=197%:%
%:%528=198%:%
%:%529=198%:%
%:%530=198%:%
%:%531=199%:%
%:%532=200%:%
%:%533=200%:%
%:%534=200%:%
%:%535=201%:%
%:%536=201%:%
%:%537=202%:%
%:%538=202%:%
%:%539=202%:%
%:%540=203%:%
%:%541=203%:%
%:%542=204%:%
%:%543=205%:%
%:%544=205%:%
%:%545=205%:%
%:%546=206%:%
%:%547=206%:%
%:%548=207%:%
%:%549=208%:%
%:%550=208%:%
%:%551=208%:%
%:%552=209%:%
%:%553=209%:%
%:%554=209%:%
%:%555=210%:%
%:%556=211%:%
%:%557=211%:%
%:%558=211%:%
%:%559=212%:%
%:%560=212%:%
%:%561=213%:%
%:%562=213%:%
%:%563=214%:%
%:%564=214%:%
%:%565=215%:%
%:%566=215%:%
%:%567=216%:%
%:%568=216%:%
%:%569=216%:%
%:%570=216%:%
%:%571=217%:%
%:%572=217%:%
%:%573=218%:%
%:%574=218%:%
%:%575=219%:%
%:%581=219%:%
%:%584=220%:%
%:%585=221%:%
%:%586=221%:%
%:%587=222%:%
%:%588=223%:%
%:%589=224%:%
%:%592=225%:%
%:%596=225%:%
%:%597=225%:%
%:%598=226%:%
%:%599=226%:%
%:%600=227%:%
%:%601=227%:%
%:%602=228%:%
%:%603=228%:%
%:%604=228%:%
%:%605=229%:%
%:%606=230%:%
%:%607=230%:%
%:%608=231%:%
%:%609=231%:%
%:%610=231%:%
%:%611=232%:%
%:%612=233%:%
%:%613=233%:%
%:%614=233%:%
%:%615=234%:%
%:%616=234%:%
%:%617=234%:%
%:%618=235%:%
%:%619=235%:%
%:%620=236%:%
%:%621=236%:%
%:%622=237%:%
%:%623=237%:%
%:%624=238%:%
%:%625=238%:%
%:%626=238%:%
%:%627=239%:%
%:%628=239%:%
%:%629=240%:%
%:%630=240%:%
%:%631=241%:%
%:%632=241%:%
%:%633=242%:%
%:%634=242%:%
%:%635=243%:%
%:%636=243%:%
%:%637=243%:%
%:%638=244%:%
%:%639=244%:%
%:%640=244%:%
%:%641=245%:%
%:%642=245%:%
%:%643=245%:%
%:%644=246%:%
%:%645=246%:%
%:%646=247%:%
%:%647=247%:%
%:%648=248%:%
%:%649=248%:%
%:%650=249%:%
%:%651=250%:%
%:%652=250%:%
%:%653=250%:%
%:%654=251%:%
%:%655=251%:%
%:%656=252%:%
%:%657=252%:%
%:%658=253%:%
%:%659=253%:%
%:%660=254%:%
%:%661=254%:%
%:%662=255%:%
%:%663=255%:%
%:%664=255%:%
%:%665=256%:%
%:%666=256%:%
%:%667=256%:%
%:%668=257%:%
%:%669=257%:%
%:%670=257%:%
%:%671=258%:%
%:%672=258%:%
%:%673=259%:%
%:%674=259%:%
%:%675=260%:%
%:%676=260%:%
%:%677=261%:%
%:%678=261%:%
%:%679=262%:%
%:%680=262%:%
%:%681=262%:%
%:%682=263%:%
%:%683=263%:%
%:%684=264%:%
%:%685=264%:%
%:%686=264%:%
%:%687=265%:%
%:%688=265%:%
%:%689=266%:%
%:%690=266%:%
%:%691=266%:%
%:%692=267%:%
%:%693=268%:%
%:%694=268%:%
%:%695=269%:%
%:%696=269%:%
%:%697=270%:%
%:%698=270%:%
%:%699=270%:%
%:%700=271%:%
%:%701=272%:%
%:%702=272%:%
%:%703=272%:%
%:%704=273%:%
%:%705=273%:%
%:%706=273%:%
%:%707=274%:%
%:%708=274%:%
%:%709=275%:%
%:%710=275%:%
%:%711=275%:%
%:%712=276%:%
%:%713=276%:%
%:%714=277%:%
%:%715=277%:%
%:%716=277%:%
%:%717=278%:%
%:%718=278%:%
%:%719=278%:%
%:%720=279%:%
%:%721=279%:%
%:%722=279%:%
%:%723=280%:%
%:%724=280%:%
%:%725=281%:%
%:%726=281%:%
%:%727=282%:%
%:%728=282%:%
%:%729=282%:%
%:%730=283%:%
%:%731=284%:%
%:%732=284%:%
%:%733=285%:%
%:%734=285%:%
%:%735=286%:%
%:%736=286%:%
%:%737=287%:%
%:%738=288%:%
%:%739=288%:%
%:%740=289%:%
%:%741=289%:%
%:%742=289%:%
%:%743=290%:%
%:%744=290%:%
%:%745=291%:%
%:%746=291%:%
%:%747=292%:%
%:%748=292%:%
%:%749=293%:%
%:%750=293%:%
%:%751=294%:%
%:%752=294%:%
%:%753=294%:%
%:%754=295%:%
%:%755=296%:%
%:%756=296%:%
%:%757=297%:%
%:%758=297%:%
%:%759=298%:%
%:%760=298%:%
%:%761=299%:%
%:%762=299%:%
%:%763=300%:%
%:%764=300%:%
%:%765=300%:%
%:%766=301%:%
%:%767=301%:%
%:%768=301%:%
%:%769=302%:%
%:%770=302%:%
%:%775=302%:%
%:%778=303%:%
%:%779=304%:%
%:%780=305%:%
%:%781=306%:%
%:%782=306%:%
%:%783=307%:%
%:%784=308%:%
%:%785=309%:%
%:%792=310%:%
%:%793=310%:%
%:%794=311%:%
%:%795=311%:%
%:%796=312%:%
%:%797=312%:%
%:%798=312%:%
%:%799=313%:%
%:%800=313%:%
%:%801=314%:%
%:%802=314%:%
%:%803=314%:%
%:%804=315%:%
%:%810=315%:%
%:%813=316%:%
%:%814=317%:%
%:%815=317%:%
%:%816=318%:%
%:%817=319%:%
%:%818=320%:%
%:%821=321%:%
%:%825=321%:%
%:%826=321%:%
%:%827=322%:%
%:%828=322%:%
%:%829=323%:%
%:%830=323%:%
%:%831=324%:%
%:%832=324%:%
%:%833=324%:%
%:%834=325%:%
%:%835=325%:%
%:%836=326%:%
%:%837=326%:%
%:%838=327%:%
%:%839=327%:%
%:%840=327%:%
%:%841=328%:%
%:%842=328%:%
%:%843=328%:%
%:%844=329%:%
%:%845=329%:%
%:%846=330%:%
%:%847=330%:%
%:%848=331%:%
%:%849=331%:%
%:%850=331%:%
%:%851=332%:%
%:%852=332%:%
%:%853=333%:%
%:%854=333%:%
%:%855=334%:%
%:%856=334%:%
%:%857=335%:%
%:%858=335%:%
%:%859=336%:%
%:%860=336%:%
%:%861=337%:%
%:%862=337%:%
%:%863=337%:%
%:%864=338%:%
%:%865=338%:%
%:%866=338%:%
%:%867=339%:%
%:%868=339%:%
%:%869=340%:%
%:%870=340%:%
%:%871=341%:%
%:%872=341%:%
%:%873=341%:%
%:%874=342%:%
%:%875=342%:%
%:%876=343%:%
%:%877=343%:%
%:%878=344%:%
%:%879=344%:%
%:%880=345%:%
%:%881=345%:%
%:%882=346%:%
%:%883=346%:%
%:%884=346%:%
%:%885=347%:%
%:%886=347%:%
%:%887=348%:%
%:%888=348%:%
%:%889=349%:%
%:%890=349%:%
%:%895=349%:%
%:%898=350%:%
%:%899=351%:%
%:%900=351%:%
%:%903=352%:%
%:%907=352%:%
%:%908=352%:%
%:%909=352%:%
%:%914=352%:%
%:%917=353%:%
%:%918=354%:%
%:%919=355%:%
%:%920=355%:%
%:%921=356%:%
%:%924=357%:%
%:%928=357%:%
%:%929=357%:%
%:%930=358%:%
%:%931=358%:%
%:%936=358%:%
%:%939=359%:%
%:%940=360%:%
%:%941=360%:%
%:%944=361%:%
%:%948=361%:%
%:%949=361%:%
%:%950=361%:%
%:%955=361%:%
%:%958=362%:%
%:%959=363%:%
%:%960=364%:%
%:%961=365%:%
%:%962=366%:%
%:%963=366%:%
%:%964=367%:%
%:%971=368%:%
%:%972=368%:%
%:%973=369%:%
%:%974=369%:%
%:%975=370%:%
%:%976=370%:%
%:%977=370%:%
%:%978=371%:%
%:%979=371%:%
%:%980=372%:%
%:%981=373%:%
%:%982=373%:%
%:%983=374%:%
%:%984=374%:%
%:%985=375%:%
%:%986=375%:%
%:%987=376%:%
%:%988=376%:%
%:%989=376%:%
%:%990=377%:%
%:%991=377%:%
%:%992=378%:%
%:%993=378%:%
%:%994=379%:%
%:%995=379%:%
%:%1000=379%:%
%:%1003=380%:%
%:%1004=381%:%
%:%1005=381%:%
%:%1006=382%:%
%:%1007=383%:%
%:%1010=384%:%
%:%1014=384%:%
%:%1015=384%:%
%:%1016=385%:%
%:%1017=385%:%
%:%1018=385%:%
%:%1023=385%:%
%:%1026=386%:%
%:%1027=387%:%
%:%1028=387%:%
%:%1029=388%:%
%:%1032=389%:%
%:%1036=389%:%
%:%1037=389%:%
%:%1038=389%:%
%:%1043=389%:%
%:%1048=390%:%
%:%1053=391%:%

%
\begin{isabellebody}%
\setisabellecontext{Mgu}%
%
\isadelimtheory
%
\endisadelimtheory
%
\isatagtheory
\isacommand{theory}\isamarkupfalse%
\ Mgu\isanewline
\isanewline
\isakeyword{imports}\ Substs\isanewline
\isanewline
\isakeyword{begin}%
\endisatagtheory
{\isafoldtheory}%
%
\isadelimtheory
\isanewline
%
\endisadelimtheory
\isanewline
\isanewline
\isanewline
\isacommand{syntax}\isamarkupfalse%
\ equ{\isacharunderscore}{\kern0pt}prob\ {\isacharcolon}{\kern0pt}{\isacharcolon}{\kern0pt}\ {\isachardoublequoteopen}trm\ {\isasymRightarrow}\ trm\ {\isasymRightarrow}\ {\isacharparenleft}{\kern0pt}trm\ {\isasymtimes}\ trm{\isacharparenright}{\kern0pt}{\isachardoublequoteclose}\ \ {\isacharparenleft}{\kern0pt}\isakeyword{infix}\ {\isachardoublequoteopen}{\isasymapprox}{\isacharquery}{\kern0pt}{\isachardoublequoteclose}\ {\isadigit{8}}{\isadigit{1}}{\isacharparenright}{\kern0pt}\isanewline
\ \ fresh{\isacharunderscore}{\kern0pt}prob\ {\isacharcolon}{\kern0pt}{\isacharcolon}{\kern0pt}\ {\isachardoublequoteopen}string\ {\isasymRightarrow}\ trm\ {\isasymRightarrow}\ {\isacharparenleft}{\kern0pt}string\ {\isasymtimes}\ trm{\isacharparenright}{\kern0pt}{\isachardoublequoteclose}\ {\isacharparenleft}{\kern0pt}\isakeyword{infix}\ {\isachardoublequoteopen}{\isasymsharp}{\isacharquery}{\kern0pt}{\isachardoublequoteclose}\ {\isadigit{8}}{\isadigit{1}}{\isacharparenright}{\kern0pt}\isanewline
\isacommand{translations}\isamarkupfalse%
\isanewline
\ {\isachardoublequoteopen}t{\isadigit{1}}\ {\isasymapprox}{\isacharquery}{\kern0pt}\ t{\isadigit{2}}{\isachardoublequoteclose}\ {\isasymrightharpoonup}\ {\isachardoublequoteopen}{\isacharparenleft}{\kern0pt}t{\isadigit{1}}{\isacharcomma}{\kern0pt}\ t{\isadigit{2}}{\isacharparenright}{\kern0pt}{\isachardoublequoteclose}\isanewline
\ {\isachardoublequoteopen}a\ {\isasymsharp}{\isacharquery}{\kern0pt}\ t{\isachardoublequoteclose}\ \ {\isasymrightharpoonup}\ {\isachardoublequoteopen}{\isacharparenleft}{\kern0pt}a{\isacharcomma}{\kern0pt}t{\isacharparenright}{\kern0pt}{\isachardoublequoteclose}\isanewline
\isanewline
\isanewline
\isanewline
\isanewline
\isanewline
\isacommand{type{\isacharunderscore}{\kern0pt}synonym}\isamarkupfalse%
\ problem{\isacharunderscore}{\kern0pt}type\ {\isacharequal}{\kern0pt}\ {\isachardoublequoteopen}{\isacharparenleft}{\kern0pt}{\isacharparenleft}{\kern0pt}trm\ {\isasymtimes}\ trm{\isacharparenright}{\kern0pt}\ list{\isacharparenright}{\kern0pt}\ {\isasymtimes}\ {\isacharparenleft}{\kern0pt}{\isacharparenleft}{\kern0pt}string\ {\isasymtimes}\ trm{\isacharparenright}{\kern0pt}\ list{\isacharparenright}{\kern0pt}{\isachardoublequoteclose}\isanewline
\isanewline
\isacommand{type{\isacharunderscore}{\kern0pt}synonym}\isamarkupfalse%
\ unifier{\isacharunderscore}{\kern0pt}type\ {\isacharequal}{\kern0pt}\ {\isachardoublequoteopen}fresh{\isacharunderscore}{\kern0pt}envs\ {\isasymtimes}\ substs{\isachardoublequoteclose}\isanewline
\isanewline
\isacommand{definition}\isamarkupfalse%
\ U{\isacharcolon}{\kern0pt}{\isacharcolon}{\kern0pt}\ {\isachardoublequoteopen}problem{\isacharunderscore}{\kern0pt}type\ {\isasymRightarrow}\ {\isacharparenleft}{\kern0pt}unifier{\isacharunderscore}{\kern0pt}type\ set{\isacharparenright}{\kern0pt}{\isachardoublequoteclose}\isanewline
\ \ \isakeyword{where}\ all{\isacharunderscore}{\kern0pt}solutions{\isacharunderscore}{\kern0pt}def{\isacharcolon}{\kern0pt}\ \isanewline
\ \ \ \ {\isachardoublequoteopen}U\ P\ \ {\isasymequiv}\ {\isacharbraceleft}{\kern0pt}{\isacharparenleft}{\kern0pt}nabla{\isacharcomma}{\kern0pt}\ s{\isacharparenright}{\kern0pt}{\isachardot}{\kern0pt}\ \isanewline
\ \ \ \ \ \ \ \ \ \ \ \ \ {\isacharparenleft}{\kern0pt}{\isasymforall}\ {\isacharparenleft}{\kern0pt}t{\isadigit{1}}{\isacharcomma}{\kern0pt}t{\isadigit{2}}{\isacharparenright}{\kern0pt}\ {\isasymin}\ set\ {\isacharparenleft}{\kern0pt}fst\ P{\isacharparenright}{\kern0pt}{\isachardot}{\kern0pt}\ nabla\ {\isasymturnstile}\ subst\ s\ t{\isadigit{1}}\ {\isasymapprox}\ subst\ s\ t{\isadigit{2}}{\isacharparenright}{\kern0pt}\ {\isasymand}\ \isanewline
\ \ \ \ \ \ \ \ \ \ \ \ \ {\isacharparenleft}{\kern0pt}{\isasymforall}\ {\isacharparenleft}{\kern0pt}a{\isacharcomma}{\kern0pt}t{\isacharparenright}{\kern0pt}\ {\isasymin}\ set\ {\isacharparenleft}{\kern0pt}snd\ P{\isacharparenright}{\kern0pt}{\isachardot}{\kern0pt}\ nabla\ {\isasymturnstile}\ a\ {\isasymsharp}\ subst\ s\ t{\isacharparenright}{\kern0pt}{\isacharbraceright}{\kern0pt}{\isachardoublequoteclose}\isanewline
\isanewline
\isanewline
\isanewline
\isanewline
\isacommand{type{\isacharunderscore}{\kern0pt}synonym}\isamarkupfalse%
\ eprobs\ {\isacharequal}{\kern0pt}\ {\isachardoublequoteopen}{\isacharparenleft}{\kern0pt}{\isacharparenleft}{\kern0pt}trm\ {\isasymtimes}\ trm{\isacharparenright}{\kern0pt}\ list{\isacharparenright}{\kern0pt}{\isachardoublequoteclose}\isanewline
\isacommand{type{\isacharunderscore}{\kern0pt}synonym}\isamarkupfalse%
\ fprobs\ {\isacharequal}{\kern0pt}\ {\isachardoublequoteopen}{\isacharparenleft}{\kern0pt}{\isacharparenleft}{\kern0pt}string\ {\isasymtimes}\ trm{\isacharparenright}{\kern0pt}\ list{\isacharparenright}{\kern0pt}{\isachardoublequoteclose}\isanewline
\isacommand{type{\isacharunderscore}{\kern0pt}synonym}\isamarkupfalse%
\ probs\ {\isacharequal}{\kern0pt}\ {\isachardoublequoteopen}eprobs\ {\isasymtimes}\ fprobs{\isachardoublequoteclose}\isanewline
\isanewline
\isanewline
\isacommand{fun}\isamarkupfalse%
\ vars{\isacharunderscore}{\kern0pt}fprobs{\isacharcolon}{\kern0pt}{\isacharcolon}{\kern0pt}\ {\isachardoublequoteopen}{\isacharparenleft}{\kern0pt}{\isacharparenleft}{\kern0pt}string\ {\isasymtimes}\ trm{\isacharparenright}{\kern0pt}\ list{\isacharparenright}{\kern0pt}\ {\isasymRightarrow}\ {\isacharparenleft}{\kern0pt}string\ set{\isacharparenright}{\kern0pt}{\isachardoublequoteclose}\isanewline
\ \ \isakeyword{where}\isanewline
\ \ {\isachardoublequoteopen}vars{\isacharunderscore}{\kern0pt}fprobs\ {\isacharbrackleft}{\kern0pt}{\isacharbrackright}{\kern0pt}\ {\isacharequal}{\kern0pt}\ {\isacharbraceleft}{\kern0pt}{\isacharbraceright}{\kern0pt}{\isachardoublequoteclose}\ {\isacharbar}{\kern0pt}\isanewline
\ \ {\isachardoublequoteopen}vars{\isacharunderscore}{\kern0pt}fprobs\ {\isacharparenleft}{\kern0pt}x{\isacharhash}{\kern0pt}xs{\isacharparenright}{\kern0pt}\ {\isacharequal}{\kern0pt}\ {\isacharparenleft}{\kern0pt}vars{\isacharunderscore}{\kern0pt}trm\ {\isacharparenleft}{\kern0pt}snd\ x{\isacharparenright}{\kern0pt}{\isacharparenright}{\kern0pt}{\isasymunion}{\isacharparenleft}{\kern0pt}vars{\isacharunderscore}{\kern0pt}fprobs\ xs{\isacharparenright}{\kern0pt}{\isachardoublequoteclose}\isanewline
\isanewline
\isacommand{fun}\isamarkupfalse%
\ \ vars{\isacharunderscore}{\kern0pt}eprobs\ {\isacharcolon}{\kern0pt}{\isacharcolon}{\kern0pt}\ {\isachardoublequoteopen}{\isacharparenleft}{\kern0pt}{\isacharparenleft}{\kern0pt}trm\ {\isasymtimes}\ trm{\isacharparenright}{\kern0pt}\ list{\isacharparenright}{\kern0pt}\ {\isasymRightarrow}\ {\isacharparenleft}{\kern0pt}string\ set{\isacharparenright}{\kern0pt}{\isachardoublequoteclose}\isanewline
\ \ \isakeyword{where}\isanewline
\ \ \ {\isachardoublequoteopen}vars{\isacharunderscore}{\kern0pt}eprobs\ {\isacharbrackleft}{\kern0pt}{\isacharbrackright}{\kern0pt}\ {\isacharequal}{\kern0pt}\ {\isacharbraceleft}{\kern0pt}{\isacharbraceright}{\kern0pt}{\isachardoublequoteclose}\ {\isacharbar}{\kern0pt}\isanewline
\ \ \ {\isachardoublequoteopen}vars{\isacharunderscore}{\kern0pt}eprobs\ {\isacharparenleft}{\kern0pt}x{\isacharhash}{\kern0pt}xs{\isacharparenright}{\kern0pt}\ {\isacharequal}{\kern0pt}\ {\isacharparenleft}{\kern0pt}vars{\isacharunderscore}{\kern0pt}trm\ {\isacharparenleft}{\kern0pt}snd\ x{\isacharparenright}{\kern0pt}{\isacharparenright}{\kern0pt}{\isasymunion}{\isacharparenleft}{\kern0pt}vars{\isacharunderscore}{\kern0pt}trm\ {\isacharparenleft}{\kern0pt}fst\ x{\isacharparenright}{\kern0pt}{\isacharparenright}{\kern0pt}{\isasymunion}{\isacharparenleft}{\kern0pt}vars{\isacharunderscore}{\kern0pt}eprobs\ xs{\isacharparenright}{\kern0pt}{\isachardoublequoteclose}\isanewline
\isanewline
\isacommand{definition}\isamarkupfalse%
\ vars{\isacharunderscore}{\kern0pt}probs\ {\isacharcolon}{\kern0pt}{\isacharcolon}{\kern0pt}\ {\isachardoublequoteopen}problem{\isacharunderscore}{\kern0pt}type\ {\isasymRightarrow}\ nat{\isachardoublequoteclose}\isanewline
\ \ \isakeyword{where}\ {\isachardoublequoteopen}vars{\isacharunderscore}{\kern0pt}probs\ P\ {\isasymequiv}\ card{\isacharparenleft}{\kern0pt}{\isacharparenleft}{\kern0pt}vars{\isacharunderscore}{\kern0pt}fprobs\ {\isacharparenleft}{\kern0pt}snd\ P{\isacharparenright}{\kern0pt}{\isacharparenright}{\kern0pt}{\isasymunion}{\isacharparenleft}{\kern0pt}vars{\isacharunderscore}{\kern0pt}eprobs\ {\isacharparenleft}{\kern0pt}fst\ P{\isacharparenright}{\kern0pt}{\isacharparenright}{\kern0pt}{\isacharparenright}{\kern0pt}{\isachardoublequoteclose}\isanewline
\isanewline
\isanewline
\isanewline
\isanewline
\isacommand{definition}\isamarkupfalse%
\ mgu\ {\isacharcolon}{\kern0pt}{\isacharcolon}{\kern0pt}\ {\isachardoublequoteopen}problem{\isacharunderscore}{\kern0pt}type\ {\isasymRightarrow}\ unifier{\isacharunderscore}{\kern0pt}type\ {\isasymRightarrow}\ bool{\isachardoublequoteclose}\isanewline
\ \ \isakeyword{where}\ {\isachardoublequoteopen}mgu\ P\ unif\ {\isasymequiv}\ \isanewline
\ \ \ {\isasymforall}\ {\isacharparenleft}{\kern0pt}nabla{\isacharcomma}{\kern0pt}s{\isadigit{1}}{\isacharparenright}{\kern0pt}{\isasymin}\ U\ P{\isachardot}{\kern0pt}\ {\isacharparenleft}{\kern0pt}{\isasymexists}\ s{\isadigit{2}}{\isachardot}{\kern0pt}\ {\isacharparenleft}{\kern0pt}nabla\ {\isasymTurnstile}\ {\isacharparenleft}{\kern0pt}subst\ s{\isadigit{2}}{\isacharparenright}{\kern0pt}\ {\isacharparenleft}{\kern0pt}fst\ unif{\isacharparenright}{\kern0pt}{\isacharparenright}{\kern0pt}\ {\isasymand}\ \isanewline
\ \ \ \ \ \ \ \ \ \ \ \ \ \ \ \ \ \ \ \ \ \ \ \ \ \ \ \ \ \ {\isacharparenleft}{\kern0pt}nabla\ {\isasymTurnstile}\ subst\ {\isacharparenleft}{\kern0pt}s{\isadigit{2}}\ {\isasymbullet}{\isacharparenleft}{\kern0pt}snd\ unif{\isacharparenright}{\kern0pt}{\isacharparenright}{\kern0pt}\ {\isasymapprox}\ subst\ s{\isadigit{1}}{\isacharparenright}{\kern0pt}{\isacharparenright}{\kern0pt}{\isachardoublequoteclose}\isanewline
\isanewline
\isanewline
\isanewline
\isanewline
\isacommand{definition}\isamarkupfalse%
\ idem\ {\isacharcolon}{\kern0pt}{\isacharcolon}{\kern0pt}\ {\isachardoublequoteopen}unifier{\isacharunderscore}{\kern0pt}type\ {\isasymRightarrow}\ bool{\isachardoublequoteclose}\isanewline
\ \ \isakeyword{where}\ {\isachardoublequoteopen}idem\ unif\ {\isasymequiv}\ {\isacharparenleft}{\kern0pt}fst\ unif{\isacharparenright}{\kern0pt}{\isasymTurnstile}\ subst\ {\isacharparenleft}{\kern0pt}{\isacharparenleft}{\kern0pt}snd\ unif{\isacharparenright}{\kern0pt}{\isasymbullet}{\isacharparenleft}{\kern0pt}snd\ unif{\isacharparenright}{\kern0pt}{\isacharparenright}{\kern0pt}\ {\isasymapprox}\ subst\ {\isacharparenleft}{\kern0pt}snd\ unif{\isacharparenright}{\kern0pt}{\isachardoublequoteclose}\isanewline
\isanewline
\isanewline
\isanewline
\isanewline
\isanewline
\isacommand{definition}\isamarkupfalse%
\ apply{\isacharunderscore}{\kern0pt}subst{\isacharunderscore}{\kern0pt}eprobs\ {\isacharcolon}{\kern0pt}{\isacharcolon}{\kern0pt}\ {\isachardoublequoteopen}substs\ {\isasymRightarrow}\ eprobs\ {\isasymRightarrow}\ eprobs{\isachardoublequoteclose}\isanewline
\ \ \isakeyword{where}\ {\isachardoublequoteopen}apply{\isacharunderscore}{\kern0pt}subst{\isacharunderscore}{\kern0pt}eprobs\ s\ P\ {\isasymequiv}\ map\ {\isacharparenleft}{\kern0pt}{\isasymlambda}{\isacharparenleft}{\kern0pt}t{\isadigit{1}}{\isacharcomma}{\kern0pt}\ t{\isadigit{2}}{\isacharparenright}{\kern0pt}{\isachardot}{\kern0pt}\ {\isacharparenleft}{\kern0pt}subst\ s\ t{\isadigit{1}}\ {\isasymapprox}{\isacharquery}{\kern0pt}\ subst\ s\ t{\isadigit{2}}{\isacharparenright}{\kern0pt}{\isacharparenright}{\kern0pt}\ P{\isachardoublequoteclose}\isanewline
\isanewline
\isacommand{definition}\isamarkupfalse%
\ apply{\isacharunderscore}{\kern0pt}subst{\isacharunderscore}{\kern0pt}fprobs\ {\isacharcolon}{\kern0pt}{\isacharcolon}{\kern0pt}\ {\isachardoublequoteopen}substs\ {\isasymRightarrow}\ fprobs\ {\isasymRightarrow}\ fprobs{\isachardoublequoteclose}\isanewline
\ \ \isakeyword{where}\ {\isachardoublequoteopen}apply{\isacharunderscore}{\kern0pt}subst{\isacharunderscore}{\kern0pt}fprobs\ s\ P\ {\isasymequiv}\ map\ {\isacharparenleft}{\kern0pt}{\isasymlambda}{\isacharparenleft}{\kern0pt}a{\isacharcomma}{\kern0pt}\ t{\isacharparenright}{\kern0pt}{\isachardot}{\kern0pt}\ {\isacharparenleft}{\kern0pt}a\ {\isasymsharp}{\isacharquery}{\kern0pt}\ subst\ s\ t{\isacharparenright}{\kern0pt}{\isacharparenright}{\kern0pt}\ P{\isachardoublequoteclose}\isanewline
\isanewline
\isanewline
\isacommand{definition}\isamarkupfalse%
\ apply{\isacharunderscore}{\kern0pt}subst\ {\isacharcolon}{\kern0pt}{\isacharcolon}{\kern0pt}\ {\isachardoublequoteopen}substs\ {\isasymRightarrow}\ problem{\isacharunderscore}{\kern0pt}type\ {\isasymRightarrow}\ problem{\isacharunderscore}{\kern0pt}type{\isachardoublequoteclose}\isanewline
\ \ \isakeyword{where}\ {\isachardoublequoteopen}apply{\isacharunderscore}{\kern0pt}subst\ s\ P\ {\isasymequiv}\ {\isacharparenleft}{\kern0pt}map\ {\isacharparenleft}{\kern0pt}{\isasymlambda}{\isacharparenleft}{\kern0pt}t{\isadigit{1}}{\isacharcomma}{\kern0pt}\ t{\isadigit{2}}{\isacharparenright}{\kern0pt}{\isachardot}{\kern0pt}\ {\isacharparenleft}{\kern0pt}subst\ s\ t{\isadigit{1}}\ {\isasymapprox}{\isacharquery}{\kern0pt}\ subst\ s\ t{\isadigit{2}}{\isacharparenright}{\kern0pt}{\isacharparenright}{\kern0pt}\ {\isacharparenleft}{\kern0pt}fst\ P{\isacharparenright}{\kern0pt}{\isacharcomma}{\kern0pt}\isanewline
\ \ \ \ \ \ \ \ \ \ \ \ \ \ \ \ \ \ \ \ \ \ map\ {\isacharparenleft}{\kern0pt}{\isasymlambda}{\isacharparenleft}{\kern0pt}a{\isacharcomma}{\kern0pt}\ t{\isacharparenright}{\kern0pt}{\isachardot}{\kern0pt}\ {\isacharparenleft}{\kern0pt}a\ {\isasymsharp}{\isacharquery}{\kern0pt}\ {\isacharparenleft}{\kern0pt}subst\ s\ t{\isacharparenright}{\kern0pt}{\isacharparenright}{\kern0pt}{\isacharparenright}{\kern0pt}\ {\isacharparenleft}{\kern0pt}snd\ P{\isacharparenright}{\kern0pt}{\isacharparenright}{\kern0pt}{\isachardoublequoteclose}\isanewline
\isanewline
\isanewline
\isanewline
\isanewline
\isacommand{inductive}\isamarkupfalse%
\ s{\isacharunderscore}{\kern0pt}red\ {\isacharcolon}{\kern0pt}{\isacharcolon}{\kern0pt}\ {\isachardoublequoteopen}problem{\isacharunderscore}{\kern0pt}type\ {\isasymRightarrow}\ substs\ {\isasymRightarrow}\ problem{\isacharunderscore}{\kern0pt}type\ {\isasymRightarrow}\ bool{\isachardoublequoteclose}\ \ {\isacharparenleft}{\kern0pt}{\isachardoublequoteopen}{\isacharunderscore}{\kern0pt}\ {\isasymturnstile}\ {\isacharunderscore}{\kern0pt}\ {\isasymleadsto}\ {\isacharunderscore}{\kern0pt}\ {\isachardoublequoteclose}\ {\isacharbrackleft}{\kern0pt}{\isadigit{8}}{\isadigit{0}}{\isacharcomma}{\kern0pt}{\isadigit{8}}{\isadigit{0}}{\isacharcomma}{\kern0pt}{\isadigit{8}}{\isadigit{0}}{\isacharbrackright}{\kern0pt}\ {\isadigit{8}}{\isadigit{0}}{\isacharparenright}{\kern0pt}\isanewline
\ \ \isakeyword{where}\isanewline
\ \ unit{\isacharunderscore}{\kern0pt}sred{\isacharbrackleft}{\kern0pt}intro{\isacharbang}{\kern0pt}{\isacharbrackright}{\kern0pt}{\isacharcolon}{\kern0pt}\ \ \ \ {\isachardoublequoteopen}{\isacharparenleft}{\kern0pt}{\isacharparenleft}{\kern0pt}Unit\ {\isasymapprox}{\isacharquery}{\kern0pt}\ Unit{\isacharparenright}{\kern0pt}\ {\isacharhash}{\kern0pt}\ xs{\isacharcomma}{\kern0pt}\ ys{\isacharparenright}{\kern0pt}\ {\isasymturnstile}\ {\isacharbrackleft}{\kern0pt}{\isacharbrackright}{\kern0pt}\ {\isasymleadsto}\ {\isacharparenleft}{\kern0pt}xs{\isacharcomma}{\kern0pt}\ ys{\isacharparenright}{\kern0pt}{\isachardoublequoteclose}\ {\isacharbar}{\kern0pt}\isanewline
\ \ paar{\isacharunderscore}{\kern0pt}sred{\isacharbrackleft}{\kern0pt}intro{\isacharbang}{\kern0pt}{\isacharbrackright}{\kern0pt}{\isacharcolon}{\kern0pt}\ \ \ \ {\isachardoublequoteopen}{\isacharparenleft}{\kern0pt}{\isacharparenleft}{\kern0pt}Paar\ t{\isadigit{1}}\ t{\isadigit{2}}\ {\isasymapprox}{\isacharquery}{\kern0pt}\ Paar\ s{\isadigit{1}}\ s{\isadigit{2}}{\isacharparenright}{\kern0pt}{\isacharhash}{\kern0pt}xs{\isacharcomma}{\kern0pt}ys{\isacharparenright}{\kern0pt}\ {\isasymturnstile}{\isacharbrackleft}{\kern0pt}{\isacharbrackright}{\kern0pt}{\isasymleadsto}\ {\isacharparenleft}{\kern0pt}{\isacharparenleft}{\kern0pt}t{\isadigit{1}}{\isasymapprox}{\isacharquery}{\kern0pt}s{\isadigit{1}}{\isacharparenright}{\kern0pt}{\isacharhash}{\kern0pt}{\isacharparenleft}{\kern0pt}t{\isadigit{2}}{\isasymapprox}{\isacharquery}{\kern0pt}s{\isadigit{2}}{\isacharparenright}{\kern0pt}{\isacharhash}{\kern0pt}xs{\isacharcomma}{\kern0pt}ys{\isacharparenright}{\kern0pt}{\isachardoublequoteclose}\ {\isacharbar}{\kern0pt}\isanewline
\ \ func{\isacharunderscore}{\kern0pt}sred{\isacharbrackleft}{\kern0pt}intro{\isacharbang}{\kern0pt}{\isacharbrackright}{\kern0pt}{\isacharcolon}{\kern0pt}\ \ \ \ {\isachardoublequoteopen}{\isacharparenleft}{\kern0pt}{\isacharparenleft}{\kern0pt}Func\ F\ t{\isadigit{1}}\ {\isasymapprox}{\isacharquery}{\kern0pt}\ Func\ F\ t{\isadigit{2}}{\isacharparenright}{\kern0pt}{\isacharhash}{\kern0pt}xs{\isacharcomma}{\kern0pt}ys{\isacharparenright}{\kern0pt}\ {\isasymturnstile}{\isacharbrackleft}{\kern0pt}{\isacharbrackright}{\kern0pt}{\isasymleadsto}\ {\isacharparenleft}{\kern0pt}{\isacharparenleft}{\kern0pt}t{\isadigit{1}}{\isasymapprox}{\isacharquery}{\kern0pt}t{\isadigit{2}}{\isacharparenright}{\kern0pt}{\isacharhash}{\kern0pt}xs{\isacharcomma}{\kern0pt}ys{\isacharparenright}{\kern0pt}{\isachardoublequoteclose}\ {\isacharbar}{\kern0pt}\isanewline
\ \ abst{\isacharunderscore}{\kern0pt}aa{\isacharunderscore}{\kern0pt}sred{\isacharbrackleft}{\kern0pt}intro{\isacharbang}{\kern0pt}{\isacharbrackright}{\kern0pt}{\isacharcolon}{\kern0pt}\ {\isachardoublequoteopen}{\isacharparenleft}{\kern0pt}{\isacharparenleft}{\kern0pt}Abst\ a\ t{\isadigit{1}}\ {\isasymapprox}{\isacharquery}{\kern0pt}\ Abst\ a\ t{\isadigit{2}}{\isacharparenright}{\kern0pt}{\isacharhash}{\kern0pt}xs{\isacharcomma}{\kern0pt}ys{\isacharparenright}{\kern0pt}\ {\isasymturnstile}{\isacharbrackleft}{\kern0pt}{\isacharbrackright}{\kern0pt}{\isasymleadsto}\ {\isacharparenleft}{\kern0pt}{\isacharparenleft}{\kern0pt}t{\isadigit{1}}{\isasymapprox}{\isacharquery}{\kern0pt}t{\isadigit{2}}{\isacharparenright}{\kern0pt}{\isacharhash}{\kern0pt}xs{\isacharcomma}{\kern0pt}ys{\isacharparenright}{\kern0pt}{\isachardoublequoteclose}\ {\isacharbar}{\kern0pt}\isanewline
\ \ abst{\isacharunderscore}{\kern0pt}ab{\isacharunderscore}{\kern0pt}sred{\isacharbrackleft}{\kern0pt}intro{\isacharbang}{\kern0pt}{\isacharbrackright}{\kern0pt}{\isacharcolon}{\kern0pt}\ {\isachardoublequoteopen}a{\isasymnoteq}b\ {\isasymLongrightarrow}\ \isanewline
\ \ \ \ \ \ \ \ \ \ \ \ \ \ \ \ \ \ \ \ \ \ \ {\isacharparenleft}{\kern0pt}{\isacharparenleft}{\kern0pt}Abst\ a\ t{\isadigit{1}}{\isasymapprox}{\isacharquery}{\kern0pt}Abst\ b\ t{\isadigit{2}}{\isacharparenright}{\kern0pt}{\isacharhash}{\kern0pt}xs{\isacharcomma}{\kern0pt}ys{\isacharparenright}{\kern0pt}\ {\isasymturnstile}{\isacharbrackleft}{\kern0pt}{\isacharbrackright}{\kern0pt}{\isasymleadsto}\ {\isacharparenleft}{\kern0pt}{\isacharparenleft}{\kern0pt}t{\isadigit{1}}{\isasymapprox}{\isacharquery}{\kern0pt}swap\ {\isacharbrackleft}{\kern0pt}{\isacharparenleft}{\kern0pt}a{\isacharcomma}{\kern0pt}b{\isacharparenright}{\kern0pt}{\isacharbrackright}{\kern0pt}\ t{\isadigit{2}}{\isacharparenright}{\kern0pt}{\isacharhash}{\kern0pt}xs{\isacharcomma}{\kern0pt}{\isacharparenleft}{\kern0pt}a{\isasymsharp}{\isacharquery}{\kern0pt}t{\isadigit{2}}{\isacharparenright}{\kern0pt}{\isacharhash}{\kern0pt}ys{\isacharparenright}{\kern0pt}{\isachardoublequoteclose}\ {\isacharbar}{\kern0pt}\isanewline
\ \ atom{\isacharunderscore}{\kern0pt}sred{\isacharbrackleft}{\kern0pt}intro{\isacharbang}{\kern0pt}{\isacharbrackright}{\kern0pt}{\isacharcolon}{\kern0pt}\ \ \ \ {\isachardoublequoteopen}{\isacharparenleft}{\kern0pt}{\isacharparenleft}{\kern0pt}Atom\ a{\isasymapprox}{\isacharquery}{\kern0pt}Atom\ a{\isacharparenright}{\kern0pt}{\isacharhash}{\kern0pt}xs{\isacharcomma}{\kern0pt}ys{\isacharparenright}{\kern0pt}\ {\isasymturnstile}{\isacharbrackleft}{\kern0pt}{\isacharbrackright}{\kern0pt}{\isasymleadsto}\ {\isacharparenleft}{\kern0pt}xs{\isacharcomma}{\kern0pt}ys{\isacharparenright}{\kern0pt}{\isachardoublequoteclose}\ {\isacharbar}{\kern0pt}\isanewline
\ \ susp{\isacharunderscore}{\kern0pt}sred{\isacharbrackleft}{\kern0pt}intro{\isacharbang}{\kern0pt}{\isacharbrackright}{\kern0pt}{\isacharcolon}{\kern0pt}\ \ \ \ {\isachardoublequoteopen}{\isacharparenleft}{\kern0pt}{\isacharparenleft}{\kern0pt}Susp\ pi{\isadigit{1}}\ X{\isasymapprox}{\isacharquery}{\kern0pt}Susp\ pi{\isadigit{2}}\ X{\isacharparenright}{\kern0pt}{\isacharhash}{\kern0pt}xs{\isacharcomma}{\kern0pt}ys{\isacharparenright}{\kern0pt}\ \isanewline
\ \ \ \ \ \ \ \ \ \ \ \ \ \ \ \ \ \ \ \ \ \ \ \ \ \ \ \ \ \ \ \ {\isasymturnstile}{\isacharbrackleft}{\kern0pt}{\isacharbrackright}{\kern0pt}{\isasymleadsto}\ {\isacharparenleft}{\kern0pt}xs{\isacharcomma}{\kern0pt}{\isacharparenleft}{\kern0pt}map\ {\isacharparenleft}{\kern0pt}{\isasymlambda}a{\isachardot}{\kern0pt}\ a{\isasymsharp}{\isacharquery}{\kern0pt}\ Susp\ {\isacharbrackleft}{\kern0pt}{\isacharbrackright}{\kern0pt}\ X{\isacharparenright}{\kern0pt}\ {\isacharparenleft}{\kern0pt}ds{\isacharunderscore}{\kern0pt}list\ pi{\isadigit{1}}\ pi{\isadigit{2}}{\isacharparenright}{\kern0pt}{\isacharparenright}{\kern0pt}{\isacharat}{\kern0pt}ys{\isacharparenright}{\kern0pt}{\isachardoublequoteclose}\ {\isacharbar}{\kern0pt}\ \isanewline
\ \ var{\isacharunderscore}{\kern0pt}{\isadigit{1}}{\isacharunderscore}{\kern0pt}sred{\isacharbrackleft}{\kern0pt}intro{\isacharbang}{\kern0pt}{\isacharbrackright}{\kern0pt}{\isacharcolon}{\kern0pt}\ \ \ {\isachardoublequoteopen}{\isasymnot}{\isacharparenleft}{\kern0pt}occurs\ X\ t{\isacharparenright}{\kern0pt}\ {\isasymLongrightarrow}\ {\isacharparenleft}{\kern0pt}{\isacharparenleft}{\kern0pt}Susp\ pi\ X{\isasymapprox}{\isacharquery}{\kern0pt}t{\isacharparenright}{\kern0pt}{\isacharhash}{\kern0pt}xs{\isacharcomma}{\kern0pt}ys{\isacharparenright}{\kern0pt}\ \isanewline
\ \ \ \ \ \ \ \ \ \ \ \ \ \ \ \ \ \ \ \ \ \ \ \ \ \ \ \ \ \ \ {\isasymturnstile}{\isacharbrackleft}{\kern0pt}{\isacharparenleft}{\kern0pt}X{\isacharcomma}{\kern0pt}swap\ {\isacharparenleft}{\kern0pt}rev\ pi{\isacharparenright}{\kern0pt}\ t{\isacharparenright}{\kern0pt}{\isacharbrackright}{\kern0pt}{\isasymleadsto}\ apply{\isacharunderscore}{\kern0pt}subst\ {\isacharbrackleft}{\kern0pt}{\isacharparenleft}{\kern0pt}X{\isacharcomma}{\kern0pt}swap\ {\isacharparenleft}{\kern0pt}rev\ pi{\isacharparenright}{\kern0pt}\ t{\isacharparenright}{\kern0pt}{\isacharbrackright}{\kern0pt}\ {\isacharparenleft}{\kern0pt}xs{\isacharcomma}{\kern0pt}ys{\isacharparenright}{\kern0pt}{\isachardoublequoteclose}\ {\isacharbar}{\kern0pt}\isanewline
\ \ var{\isacharunderscore}{\kern0pt}{\isadigit{2}}{\isacharunderscore}{\kern0pt}sred{\isacharbrackleft}{\kern0pt}intro{\isacharbang}{\kern0pt}{\isacharbrackright}{\kern0pt}{\isacharcolon}{\kern0pt}\ \ \ {\isachardoublequoteopen}{\isasymnot}{\isacharparenleft}{\kern0pt}occurs\ X\ t{\isacharparenright}{\kern0pt}\ {\isasymLongrightarrow}\ {\isacharparenleft}{\kern0pt}{\isacharparenleft}{\kern0pt}t{\isasymapprox}{\isacharquery}{\kern0pt}Susp\ pi\ X{\isacharparenright}{\kern0pt}{\isacharhash}{\kern0pt}xs{\isacharcomma}{\kern0pt}ys{\isacharparenright}{\kern0pt}\ \isanewline
\ \ \ \ \ \ \ \ \ \ \ \ \ \ \ \ \ \ \ \ \ \ \ \ \ \ \ \ \ \ \ {\isasymturnstile}{\isacharbrackleft}{\kern0pt}{\isacharparenleft}{\kern0pt}X{\isacharcomma}{\kern0pt}swap\ {\isacharparenleft}{\kern0pt}rev\ pi{\isacharparenright}{\kern0pt}\ t{\isacharparenright}{\kern0pt}{\isacharbrackright}{\kern0pt}{\isasymleadsto}\ apply{\isacharunderscore}{\kern0pt}subst\ {\isacharbrackleft}{\kern0pt}{\isacharparenleft}{\kern0pt}X{\isacharcomma}{\kern0pt}swap\ {\isacharparenleft}{\kern0pt}rev\ pi{\isacharparenright}{\kern0pt}\ t{\isacharparenright}{\kern0pt}{\isacharbrackright}{\kern0pt}\ {\isacharparenleft}{\kern0pt}xs{\isacharcomma}{\kern0pt}ys{\isacharparenright}{\kern0pt}{\isachardoublequoteclose}\isanewline
\isanewline
\isacommand{inductive{\isacharunderscore}{\kern0pt}cases}\isamarkupfalse%
\ s{\isacharunderscore}{\kern0pt}red{\isacharunderscore}{\kern0pt}elims{\isacharcolon}{\kern0pt}\isanewline
{\isachardoublequoteopen}{\isacharparenleft}{\kern0pt}{\isacharparenleft}{\kern0pt}Unit\ {\isasymapprox}{\isacharquery}{\kern0pt}\ Unit{\isacharparenright}{\kern0pt}\ {\isacharhash}{\kern0pt}\ xs{\isacharcomma}{\kern0pt}\ ys{\isacharparenright}{\kern0pt}\ {\isasymturnstile}{\isacharbrackleft}{\kern0pt}{\isacharbrackright}{\kern0pt}\ {\isasymleadsto}\ {\isacharparenleft}{\kern0pt}xs{\isacharcomma}{\kern0pt}\ ys{\isacharparenright}{\kern0pt}{\isachardoublequoteclose}\isanewline
{\isachardoublequoteopen}{\isacharparenleft}{\kern0pt}{\isacharparenleft}{\kern0pt}Paar\ t{\isadigit{1}}\ t{\isadigit{2}}\ {\isasymapprox}{\isacharquery}{\kern0pt}\ Paar\ s{\isadigit{1}}\ s{\isadigit{2}}{\isacharparenright}{\kern0pt}{\isacharhash}{\kern0pt}xs{\isacharcomma}{\kern0pt}ys{\isacharparenright}{\kern0pt}\ {\isasymturnstile}{\isacharbrackleft}{\kern0pt}{\isacharbrackright}{\kern0pt}{\isasymleadsto}\ {\isacharparenleft}{\kern0pt}{\isacharparenleft}{\kern0pt}t{\isadigit{1}}{\isasymapprox}{\isacharquery}{\kern0pt}s{\isadigit{1}}{\isacharparenright}{\kern0pt}{\isacharhash}{\kern0pt}{\isacharparenleft}{\kern0pt}t{\isadigit{2}}{\isasymapprox}{\isacharquery}{\kern0pt}s{\isadigit{2}}{\isacharparenright}{\kern0pt}{\isacharhash}{\kern0pt}xs{\isacharcomma}{\kern0pt}ys{\isacharparenright}{\kern0pt}{\isachardoublequoteclose}\isanewline
{\isachardoublequoteopen}{\isacharparenleft}{\kern0pt}{\isacharparenleft}{\kern0pt}Func\ F\ t{\isadigit{1}}\ {\isasymapprox}{\isacharquery}{\kern0pt}\ Func\ F\ t{\isadigit{2}}{\isacharparenright}{\kern0pt}{\isacharhash}{\kern0pt}xs{\isacharcomma}{\kern0pt}ys{\isacharparenright}{\kern0pt}\ {\isasymturnstile}{\isacharbrackleft}{\kern0pt}{\isacharbrackright}{\kern0pt}{\isasymleadsto}\ {\isacharparenleft}{\kern0pt}{\isacharparenleft}{\kern0pt}t{\isadigit{1}}{\isasymapprox}{\isacharquery}{\kern0pt}t{\isadigit{2}}{\isacharparenright}{\kern0pt}{\isacharhash}{\kern0pt}xs{\isacharcomma}{\kern0pt}ys{\isacharparenright}{\kern0pt}{\isachardoublequoteclose}\isanewline
{\isachardoublequoteopen}{\isacharparenleft}{\kern0pt}{\isacharparenleft}{\kern0pt}Abst\ a\ t{\isadigit{1}}\ {\isasymapprox}{\isacharquery}{\kern0pt}\ Abst\ a\ t{\isadigit{2}}{\isacharparenright}{\kern0pt}{\isacharhash}{\kern0pt}xs{\isacharcomma}{\kern0pt}ys{\isacharparenright}{\kern0pt}\ {\isasymturnstile}{\isacharbrackleft}{\kern0pt}{\isacharbrackright}{\kern0pt}{\isasymleadsto}\ {\isacharparenleft}{\kern0pt}{\isacharparenleft}{\kern0pt}t{\isadigit{1}}{\isasymapprox}{\isacharquery}{\kern0pt}t{\isadigit{2}}{\isacharparenright}{\kern0pt}{\isacharhash}{\kern0pt}xs{\isacharcomma}{\kern0pt}ys{\isacharparenright}{\kern0pt}{\isachardoublequoteclose}\isanewline
{\isachardoublequoteopen}{\isacharparenleft}{\kern0pt}{\isacharparenleft}{\kern0pt}Abst\ a\ t{\isadigit{1}}{\isasymapprox}{\isacharquery}{\kern0pt}Abst\ b\ t{\isadigit{2}}{\isacharparenright}{\kern0pt}{\isacharhash}{\kern0pt}xs{\isacharcomma}{\kern0pt}ys{\isacharparenright}{\kern0pt}\ {\isasymturnstile}{\isacharbrackleft}{\kern0pt}{\isacharbrackright}{\kern0pt}{\isasymleadsto}\ {\isacharparenleft}{\kern0pt}{\isacharparenleft}{\kern0pt}t{\isadigit{1}}{\isasymapprox}{\isacharquery}{\kern0pt}swap\ {\isacharbrackleft}{\kern0pt}{\isacharparenleft}{\kern0pt}a{\isacharcomma}{\kern0pt}b{\isacharparenright}{\kern0pt}{\isacharbrackright}{\kern0pt}\ t{\isadigit{2}}{\isacharparenright}{\kern0pt}{\isacharhash}{\kern0pt}xs{\isacharcomma}{\kern0pt}{\isacharparenleft}{\kern0pt}a{\isasymsharp}{\isacharquery}{\kern0pt}t{\isadigit{2}}{\isacharparenright}{\kern0pt}{\isacharhash}{\kern0pt}ys{\isacharparenright}{\kern0pt}{\isachardoublequoteclose}\isanewline
{\isachardoublequoteopen}{\isacharparenleft}{\kern0pt}{\isacharparenleft}{\kern0pt}Atom\ a{\isasymapprox}{\isacharquery}{\kern0pt}Atom\ a{\isacharparenright}{\kern0pt}{\isacharhash}{\kern0pt}xs{\isacharcomma}{\kern0pt}ys{\isacharparenright}{\kern0pt}\ {\isasymturnstile}{\isacharbrackleft}{\kern0pt}{\isacharbrackright}{\kern0pt}{\isasymleadsto}\ {\isacharparenleft}{\kern0pt}xs{\isacharcomma}{\kern0pt}ys{\isacharparenright}{\kern0pt}{\isachardoublequoteclose}\isanewline
{\isachardoublequoteopen}{\isacharparenleft}{\kern0pt}{\isacharparenleft}{\kern0pt}Susp\ pi{\isadigit{1}}\ X{\isasymapprox}{\isacharquery}{\kern0pt}Susp\ pi{\isadigit{2}}\ X{\isacharparenright}{\kern0pt}{\isacharhash}{\kern0pt}xs{\isacharcomma}{\kern0pt}ys{\isacharparenright}{\kern0pt}\ {\isasymturnstile}{\isacharbrackleft}{\kern0pt}{\isacharbrackright}{\kern0pt}{\isasymleadsto}\ {\isacharparenleft}{\kern0pt}xs{\isacharcomma}{\kern0pt}{\isacharparenleft}{\kern0pt}map\ {\isacharparenleft}{\kern0pt}{\isasymlambda}a{\isachardot}{\kern0pt}\ a{\isasymsharp}{\isacharquery}{\kern0pt}\ Susp\ {\isacharbrackleft}{\kern0pt}{\isacharbrackright}{\kern0pt}\ X{\isacharparenright}{\kern0pt}\ {\isacharparenleft}{\kern0pt}ds{\isacharunderscore}{\kern0pt}list\ pi{\isadigit{1}}\ pi{\isadigit{2}}{\isacharparenright}{\kern0pt}{\isacharparenright}{\kern0pt}{\isacharat}{\kern0pt}ys{\isacharparenright}{\kern0pt}{\isachardoublequoteclose}\isanewline
{\isachardoublequoteopen}{\isacharparenleft}{\kern0pt}{\isacharparenleft}{\kern0pt}Susp\ pi\ X{\isasymapprox}{\isacharquery}{\kern0pt}t{\isacharparenright}{\kern0pt}{\isacharhash}{\kern0pt}xs{\isacharcomma}{\kern0pt}ys{\isacharparenright}{\kern0pt}\ {\isasymturnstile}{\isacharbrackleft}{\kern0pt}{\isacharparenleft}{\kern0pt}X{\isacharcomma}{\kern0pt}swap\ {\isacharparenleft}{\kern0pt}rev\ pi{\isacharparenright}{\kern0pt}\ t{\isacharparenright}{\kern0pt}{\isacharbrackright}{\kern0pt}{\isasymleadsto}\ apply{\isacharunderscore}{\kern0pt}subst\ {\isacharbrackleft}{\kern0pt}{\isacharparenleft}{\kern0pt}X{\isacharcomma}{\kern0pt}swap\ {\isacharparenleft}{\kern0pt}rev\ pi{\isacharparenright}{\kern0pt}\ t{\isacharparenright}{\kern0pt}{\isacharbrackright}{\kern0pt}\ {\isacharparenleft}{\kern0pt}xs{\isacharcomma}{\kern0pt}ys{\isacharparenright}{\kern0pt}{\isachardoublequoteclose}\isanewline
{\isachardoublequoteopen}{\isacharparenleft}{\kern0pt}{\isacharparenleft}{\kern0pt}t{\isasymapprox}{\isacharquery}{\kern0pt}Susp\ pi\ X{\isacharparenright}{\kern0pt}{\isacharhash}{\kern0pt}xs{\isacharcomma}{\kern0pt}ys{\isacharparenright}{\kern0pt}\ {\isasymturnstile}{\isacharbrackleft}{\kern0pt}{\isacharparenleft}{\kern0pt}X{\isacharcomma}{\kern0pt}swap\ {\isacharparenleft}{\kern0pt}rev\ pi{\isacharparenright}{\kern0pt}\ t{\isacharparenright}{\kern0pt}{\isacharbrackright}{\kern0pt}{\isasymleadsto}\ apply{\isacharunderscore}{\kern0pt}subst\ {\isacharbrackleft}{\kern0pt}{\isacharparenleft}{\kern0pt}X{\isacharcomma}{\kern0pt}swap\ {\isacharparenleft}{\kern0pt}rev\ pi{\isacharparenright}{\kern0pt}\ t{\isacharparenright}{\kern0pt}{\isacharbrackright}{\kern0pt}\ {\isacharparenleft}{\kern0pt}xs{\isacharcomma}{\kern0pt}ys{\isacharparenright}{\kern0pt}{\isachardoublequoteclose}\isanewline
\isanewline
\isanewline
\isanewline
\isanewline
\isacommand{lemma}\isamarkupfalse%
\ solution{\isacharunderscore}{\kern0pt}weak{\isacharcolon}{\kern0pt}\isanewline
\ \ \isakeyword{assumes}\ {\isachardoublequoteopen}{\isacharparenleft}{\kern0pt}nabla{\isadigit{1}}{\isacharcomma}{\kern0pt}\ s{\isacharparenright}{\kern0pt}\ {\isasymin}\ U\ P{\isachardoublequoteclose}\isanewline
\ \ \isakeyword{shows}\ {\isachardoublequoteopen}{\isacharparenleft}{\kern0pt}nabla{\isadigit{1}}\ {\isasymunion}\ nabla{\isadigit{3}}{\isacharcomma}{\kern0pt}\ s{\isacharparenright}{\kern0pt}\ {\isasymin}\ U\ P{\isachardoublequoteclose}\isanewline
%
\isadelimproof
\ \ %
\endisadelimproof
%
\isatagproof
\isacommand{using}\isamarkupfalse%
\ assms\isanewline
\ \ \isacommand{unfolding}\isamarkupfalse%
\ all{\isacharunderscore}{\kern0pt}solutions{\isacharunderscore}{\kern0pt}def\ \isacommand{using}\isamarkupfalse%
\ fresh{\isacharunderscore}{\kern0pt}weak\ equ{\isacharunderscore}{\kern0pt}weak\ \isacommand{by}\isamarkupfalse%
\ auto%
\endisatagproof
{\isafoldproof}%
%
\isadelimproof
\isanewline
%
\endisadelimproof
\isanewline
\isacommand{lemma}\isamarkupfalse%
\ solution{\isacharunderscore}{\kern0pt}comp{\isacharunderscore}{\kern0pt}id{\isacharcolon}{\kern0pt}\ \isanewline
\ \ \isakeyword{shows}\ {\isachardoublequoteopen}{\isacharparenleft}{\kern0pt}{\isacharparenleft}{\kern0pt}nabla{\isacharcomma}{\kern0pt}\ s{\isacharparenright}{\kern0pt}\ {\isasymin}\ U\ P{\isacharparenright}{\kern0pt}\ {\isacharequal}{\kern0pt}\ {\isacharparenleft}{\kern0pt}{\isacharparenleft}{\kern0pt}nabla{\isacharcomma}{\kern0pt}\ s\ {\isasymbullet}\ {\isacharbrackleft}{\kern0pt}{\isacharbrackright}{\kern0pt}{\isacharparenright}{\kern0pt}\ {\isasymin}\ U\ P{\isacharparenright}{\kern0pt}{\isachardoublequoteclose}\ \isakeyword{and}\ \isanewline
\ \ \ \ {\isachardoublequoteopen}{\isacharparenleft}{\kern0pt}{\isacharparenleft}{\kern0pt}nabla{\isacharcomma}{\kern0pt}\ s{\isacharparenright}{\kern0pt}\ {\isasymin}\ U\ P{\isacharparenright}{\kern0pt}\ {\isacharequal}{\kern0pt}\ {\isacharparenleft}{\kern0pt}{\isacharparenleft}{\kern0pt}nabla{\isacharcomma}{\kern0pt}\ {\isacharbrackleft}{\kern0pt}{\isacharbrackright}{\kern0pt}\ {\isasymbullet}\ s{\isacharparenright}{\kern0pt}\ {\isasymin}\ U\ P{\isacharparenright}{\kern0pt}{\isachardoublequoteclose}\isanewline
%
\isadelimproof
\ \ %
\endisadelimproof
%
\isatagproof
\isacommand{unfolding}\isamarkupfalse%
\ all{\isacharunderscore}{\kern0pt}solutions{\isacharunderscore}{\kern0pt}def\ \isacommand{by}\isamarkupfalse%
\ auto%
\endisatagproof
{\isafoldproof}%
%
\isadelimproof
\isanewline
%
\endisadelimproof
\isanewline
\isacommand{lemma}\isamarkupfalse%
\ solutions{\isacharunderscore}{\kern0pt}subst{\isacharunderscore}{\kern0pt}assoc{\isacharcolon}{\kern0pt}\isanewline
\ \ \ {\isachardoublequoteopen}{\isacharparenleft}{\kern0pt}{\isacharparenleft}{\kern0pt}nabla{\isacharcomma}{\kern0pt}\ s{\isadigit{1}}\ {\isasymbullet}\ {\isacharparenleft}{\kern0pt}s{\isadigit{2}}\ {\isasymbullet}\ s{\isadigit{3}}{\isacharparenright}{\kern0pt}{\isacharparenright}{\kern0pt}\ {\isasymin}\ U\ P{\isacharparenright}{\kern0pt}\ {\isacharequal}{\kern0pt}\ {\isacharparenleft}{\kern0pt}{\isacharparenleft}{\kern0pt}nabla{\isacharcomma}{\kern0pt}\ {\isacharparenleft}{\kern0pt}s{\isadigit{1}}\ {\isasymbullet}\ s{\isadigit{2}}{\isacharparenright}{\kern0pt}\ {\isasymbullet}\ s{\isadigit{3}}{\isacharparenright}{\kern0pt}\ {\isasymin}\ U\ P{\isacharparenright}{\kern0pt}{\isachardoublequoteclose}\isanewline
%
\isadelimproof
\ \ %
\endisadelimproof
%
\isatagproof
\isacommand{unfolding}\isamarkupfalse%
\ all{\isacharunderscore}{\kern0pt}solutions{\isacharunderscore}{\kern0pt}def\ \isacommand{using}\isamarkupfalse%
\ subst{\isacharunderscore}{\kern0pt}assoc\ \isacommand{by}\isamarkupfalse%
\ simp%
\endisatagproof
{\isafoldproof}%
%
\isadelimproof
\isanewline
%
\endisadelimproof
\isanewline
\isanewline
\isanewline
\isacommand{inductive}\isamarkupfalse%
\ c{\isacharunderscore}{\kern0pt}red\ {\isacharcolon}{\kern0pt}{\isacharcolon}{\kern0pt}\ {\isachardoublequoteopen}problem{\isacharunderscore}{\kern0pt}type\ {\isasymRightarrow}\ fresh{\isacharunderscore}{\kern0pt}envs\ {\isasymRightarrow}\ problem{\isacharunderscore}{\kern0pt}type\ {\isasymRightarrow}\ bool{\isachardoublequoteclose}\ {\isacharparenleft}{\kern0pt}{\isachardoublequoteopen}{\isacharunderscore}{\kern0pt}\ {\isasymturnstile}\ {\isacharunderscore}{\kern0pt}\ {\isasymrightarrow}\ {\isacharunderscore}{\kern0pt}\ {\isachardoublequoteclose}\ {\isacharbrackleft}{\kern0pt}{\isadigit{8}}{\isadigit{0}}{\isacharcomma}{\kern0pt}{\isadigit{8}}{\isadigit{0}}{\isacharcomma}{\kern0pt}{\isadigit{8}}{\isadigit{0}}{\isacharbrackright}{\kern0pt}\ {\isadigit{8}}{\isadigit{0}}{\isacharparenright}{\kern0pt}\isanewline
\ \ \isakeyword{where}\isanewline
\ \ unit{\isacharunderscore}{\kern0pt}cred{\isacharbrackleft}{\kern0pt}intro{\isacharbang}{\kern0pt}{\isacharbrackright}{\kern0pt}{\isacharcolon}{\kern0pt}\ \ \ \ {\isachardoublequoteopen}{\isacharparenleft}{\kern0pt}{\isacharbrackleft}{\kern0pt}{\isacharbrackright}{\kern0pt}{\isacharcomma}{\kern0pt}{\isacharparenleft}{\kern0pt}a\ {\isasymsharp}{\isacharquery}{\kern0pt}\ Unit{\isacharparenright}{\kern0pt}{\isacharhash}{\kern0pt}xs{\isacharparenright}{\kern0pt}\ {\isasymturnstile}{\isacharbraceleft}{\kern0pt}{\isacharbraceright}{\kern0pt}{\isasymrightarrow}\ {\isacharparenleft}{\kern0pt}{\isacharbrackleft}{\kern0pt}{\isacharbrackright}{\kern0pt}{\isacharcomma}{\kern0pt}xs{\isacharparenright}{\kern0pt}{\isachardoublequoteclose}\ {\isacharbar}{\kern0pt}\isanewline
\ \ paar{\isacharunderscore}{\kern0pt}cred{\isacharbrackleft}{\kern0pt}intro{\isacharbang}{\kern0pt}{\isacharbrackright}{\kern0pt}{\isacharcolon}{\kern0pt}\ \ \ \ {\isachardoublequoteopen}{\isacharparenleft}{\kern0pt}{\isacharbrackleft}{\kern0pt}{\isacharbrackright}{\kern0pt}{\isacharcomma}{\kern0pt}{\isacharparenleft}{\kern0pt}a\ {\isasymsharp}{\isacharquery}{\kern0pt}\ Paar\ t{\isadigit{1}}\ t{\isadigit{2}}{\isacharparenright}{\kern0pt}{\isacharhash}{\kern0pt}xs{\isacharparenright}{\kern0pt}\ {\isasymturnstile}{\isacharbraceleft}{\kern0pt}{\isacharbraceright}{\kern0pt}{\isasymrightarrow}\ {\isacharparenleft}{\kern0pt}{\isacharbrackleft}{\kern0pt}{\isacharbrackright}{\kern0pt}{\isacharcomma}{\kern0pt}{\isacharparenleft}{\kern0pt}a{\isasymsharp}{\isacharquery}{\kern0pt}t{\isadigit{1}}{\isacharparenright}{\kern0pt}{\isacharhash}{\kern0pt}{\isacharparenleft}{\kern0pt}a{\isasymsharp}{\isacharquery}{\kern0pt}t{\isadigit{2}}{\isacharparenright}{\kern0pt}{\isacharhash}{\kern0pt}xs{\isacharparenright}{\kern0pt}{\isachardoublequoteclose}\ {\isacharbar}{\kern0pt}\isanewline
\ \ func{\isacharunderscore}{\kern0pt}cred{\isacharbrackleft}{\kern0pt}intro{\isacharbang}{\kern0pt}{\isacharbrackright}{\kern0pt}{\isacharcolon}{\kern0pt}\ \ \ \ {\isachardoublequoteopen}{\isacharparenleft}{\kern0pt}{\isacharbrackleft}{\kern0pt}{\isacharbrackright}{\kern0pt}{\isacharcomma}{\kern0pt}{\isacharparenleft}{\kern0pt}a\ {\isasymsharp}{\isacharquery}{\kern0pt}\ Func\ F\ t{\isacharparenright}{\kern0pt}{\isacharhash}{\kern0pt}xs{\isacharparenright}{\kern0pt}\ {\isasymturnstile}{\isacharbraceleft}{\kern0pt}{\isacharbraceright}{\kern0pt}{\isasymrightarrow}\ {\isacharparenleft}{\kern0pt}{\isacharbrackleft}{\kern0pt}{\isacharbrackright}{\kern0pt}{\isacharcomma}{\kern0pt}{\isacharparenleft}{\kern0pt}a{\isasymsharp}{\isacharquery}{\kern0pt}t{\isacharparenright}{\kern0pt}{\isacharhash}{\kern0pt}xs{\isacharparenright}{\kern0pt}{\isachardoublequoteclose}\ {\isacharbar}{\kern0pt}\isanewline
\ \ abst{\isacharunderscore}{\kern0pt}aa{\isacharunderscore}{\kern0pt}cred{\isacharbrackleft}{\kern0pt}intro{\isacharbang}{\kern0pt}{\isacharbrackright}{\kern0pt}{\isacharcolon}{\kern0pt}\ {\isachardoublequoteopen}{\isacharparenleft}{\kern0pt}{\isacharbrackleft}{\kern0pt}{\isacharbrackright}{\kern0pt}{\isacharcomma}{\kern0pt}{\isacharparenleft}{\kern0pt}a\ {\isasymsharp}{\isacharquery}{\kern0pt}\ Abst\ a\ t{\isacharparenright}{\kern0pt}{\isacharhash}{\kern0pt}xs{\isacharparenright}{\kern0pt}\ {\isasymturnstile}{\isacharbraceleft}{\kern0pt}{\isacharbraceright}{\kern0pt}{\isasymrightarrow}\ {\isacharparenleft}{\kern0pt}{\isacharbrackleft}{\kern0pt}{\isacharbrackright}{\kern0pt}{\isacharcomma}{\kern0pt}xs{\isacharparenright}{\kern0pt}{\isachardoublequoteclose}\ {\isacharbar}{\kern0pt}\isanewline
\ \ abst{\isacharunderscore}{\kern0pt}ab{\isacharunderscore}{\kern0pt}cred{\isacharbrackleft}{\kern0pt}intro{\isacharbang}{\kern0pt}{\isacharbrackright}{\kern0pt}{\isacharcolon}{\kern0pt}\ {\isachardoublequoteopen}a{\isasymnoteq}b\ {\isasymLongrightarrow}{\isacharparenleft}{\kern0pt}{\isacharbrackleft}{\kern0pt}{\isacharbrackright}{\kern0pt}{\isacharcomma}{\kern0pt}{\isacharparenleft}{\kern0pt}a\ {\isasymsharp}{\isacharquery}{\kern0pt}\ Abst\ b\ t{\isacharparenright}{\kern0pt}{\isacharhash}{\kern0pt}xs{\isacharparenright}{\kern0pt}\ {\isasymturnstile}{\isacharbraceleft}{\kern0pt}{\isacharbraceright}{\kern0pt}{\isasymrightarrow}\ {\isacharparenleft}{\kern0pt}{\isacharbrackleft}{\kern0pt}{\isacharbrackright}{\kern0pt}{\isacharcomma}{\kern0pt}{\isacharparenleft}{\kern0pt}a{\isasymsharp}{\isacharquery}{\kern0pt}t{\isacharparenright}{\kern0pt}{\isacharhash}{\kern0pt}xs{\isacharparenright}{\kern0pt}{\isachardoublequoteclose}\ {\isacharbar}{\kern0pt}\isanewline
\ \ atom{\isacharunderscore}{\kern0pt}cred{\isacharbrackleft}{\kern0pt}intro{\isacharbang}{\kern0pt}{\isacharbrackright}{\kern0pt}{\isacharcolon}{\kern0pt}\ \ \ \ {\isachardoublequoteopen}a{\isasymnoteq}b\ {\isasymLongrightarrow}{\isacharparenleft}{\kern0pt}{\isacharbrackleft}{\kern0pt}{\isacharbrackright}{\kern0pt}{\isacharcomma}{\kern0pt}{\isacharparenleft}{\kern0pt}a\ {\isasymsharp}{\isacharquery}{\kern0pt}\ Atom\ b{\isacharparenright}{\kern0pt}{\isacharhash}{\kern0pt}xs{\isacharparenright}{\kern0pt}\ {\isasymturnstile}{\isacharbraceleft}{\kern0pt}{\isacharbraceright}{\kern0pt}{\isasymrightarrow}\ {\isacharparenleft}{\kern0pt}{\isacharbrackleft}{\kern0pt}{\isacharbrackright}{\kern0pt}{\isacharcomma}{\kern0pt}xs{\isacharparenright}{\kern0pt}{\isachardoublequoteclose}\ {\isacharbar}{\kern0pt}\isanewline
\ \ susp{\isacharunderscore}{\kern0pt}cred{\isacharbrackleft}{\kern0pt}intro{\isacharbang}{\kern0pt}{\isacharbrackright}{\kern0pt}{\isacharcolon}{\kern0pt}\ \ \ \ {\isachardoublequoteopen}{\isacharparenleft}{\kern0pt}{\isacharbrackleft}{\kern0pt}{\isacharbrackright}{\kern0pt}{\isacharcomma}{\kern0pt}{\isacharparenleft}{\kern0pt}a\ {\isasymsharp}{\isacharquery}{\kern0pt}\ Susp\ pi\ X{\isacharparenright}{\kern0pt}{\isacharhash}{\kern0pt}xs{\isacharparenright}{\kern0pt}\ {\isasymturnstile}\ {\isacharbraceleft}{\kern0pt}{\isacharparenleft}{\kern0pt}{\isacharparenleft}{\kern0pt}swapas\ {\isacharparenleft}{\kern0pt}rev\ pi{\isacharparenright}{\kern0pt}\ a{\isacharparenright}{\kern0pt}{\isacharcomma}{\kern0pt}X{\isacharparenright}{\kern0pt}{\isacharbraceright}{\kern0pt}{\isasymrightarrow}\ {\isacharparenleft}{\kern0pt}{\isacharbrackleft}{\kern0pt}{\isacharbrackright}{\kern0pt}{\isacharcomma}{\kern0pt}xs{\isacharparenright}{\kern0pt}{\isachardoublequoteclose}\isanewline
\isanewline
\isanewline
\isanewline
\isanewline
\isacommand{inductive}\isamarkupfalse%
\ red{\isacharunderscore}{\kern0pt}plus\ {\isacharcolon}{\kern0pt}{\isacharcolon}{\kern0pt}\ {\isachardoublequoteopen}problem{\isacharunderscore}{\kern0pt}type\ {\isasymRightarrow}\ unifier{\isacharunderscore}{\kern0pt}type\ {\isasymRightarrow}\ problem{\isacharunderscore}{\kern0pt}type\ {\isasymRightarrow}\ bool{\isachardoublequoteclose}\ {\isacharparenleft}{\kern0pt}{\isachardoublequoteopen}{\isacharunderscore}{\kern0pt}\ {\isasymTurnstile}\ {\isacharunderscore}{\kern0pt}\ {\isasymRightarrow}\ {\isacharunderscore}{\kern0pt}\ {\isachardoublequoteclose}\ {\isacharbrackleft}{\kern0pt}{\isadigit{8}}{\isadigit{0}}{\isacharcomma}{\kern0pt}{\isadigit{8}}{\isadigit{0}}{\isacharcomma}{\kern0pt}{\isadigit{8}}{\isadigit{0}}{\isacharbrackright}{\kern0pt}\ {\isadigit{8}}{\isadigit{0}}{\isacharparenright}{\kern0pt}\isanewline
\ \ \isakeyword{where}\isanewline
\ \ sred{\isacharunderscore}{\kern0pt}single{\isacharbrackleft}{\kern0pt}intro{\isacharbang}{\kern0pt}{\isacharbrackright}{\kern0pt}{\isacharcolon}{\kern0pt}\ {\isachardoublequoteopen}{\isasymlbrakk}P{\isadigit{1}}{\isasymturnstile}s{\isadigit{1}}\ {\isasymleadsto}\ P{\isadigit{2}}{\isasymrbrakk}\ {\isasymLongrightarrow}\ P{\isadigit{1}}{\isasymTurnstile}{\isacharparenleft}{\kern0pt}{\isacharbraceleft}{\kern0pt}{\isacharbraceright}{\kern0pt}{\isacharcomma}{\kern0pt}s{\isadigit{1}}{\isacharparenright}{\kern0pt}{\isasymRightarrow}P{\isadigit{2}}{\isachardoublequoteclose}\ {\isacharbar}{\kern0pt}\isanewline
\ \ cred{\isacharunderscore}{\kern0pt}single{\isacharbrackleft}{\kern0pt}intro{\isacharbang}{\kern0pt}{\isacharbrackright}{\kern0pt}{\isacharcolon}{\kern0pt}\ {\isachardoublequoteopen}{\isasymlbrakk}P{\isadigit{1}}{\isasymturnstile}nabla{\isadigit{1}}{\isasymrightarrow}P{\isadigit{2}}{\isasymrbrakk}\ {\isasymLongrightarrow}\ P{\isadigit{1}}\ {\isasymTurnstile}\ {\isacharparenleft}{\kern0pt}nabla{\isadigit{1}}{\isacharcomma}{\kern0pt}{\isacharbrackleft}{\kern0pt}{\isacharbrackright}{\kern0pt}{\isacharparenright}{\kern0pt}\ {\isasymRightarrow}\ P{\isadigit{2}}{\isachardoublequoteclose}\ {\isacharbar}{\kern0pt}\isanewline
\ \ sred{\isacharunderscore}{\kern0pt}step{\isacharbrackleft}{\kern0pt}intro{\isacharbang}{\kern0pt}{\isacharbrackright}{\kern0pt}{\isacharcolon}{\kern0pt}\ \ \ {\isachardoublequoteopen}{\isasymlbrakk}P{\isadigit{1}}{\isasymturnstile}s{\isadigit{1}}{\isasymleadsto}P{\isadigit{2}}{\isacharsemicolon}{\kern0pt}\ P{\isadigit{2}}{\isasymTurnstile}{\isacharparenleft}{\kern0pt}nabla{\isadigit{2}}{\isacharcomma}{\kern0pt}s{\isadigit{2}}{\isacharparenright}{\kern0pt}{\isasymRightarrow}P{\isadigit{3}}{\isasymrbrakk}\ {\isasymLongrightarrow}\ P{\isadigit{1}}{\isasymTurnstile}{\isacharparenleft}{\kern0pt}nabla{\isadigit{2}}{\isacharcomma}{\kern0pt}{\isacharparenleft}{\kern0pt}s{\isadigit{2}}{\isasymbullet}s{\isadigit{1}}{\isacharparenright}{\kern0pt}{\isacharparenright}{\kern0pt}{\isasymRightarrow}P{\isadigit{3}}{\isachardoublequoteclose}\ {\isacharbar}{\kern0pt}\isanewline
\ \ cred{\isacharunderscore}{\kern0pt}step{\isacharbrackleft}{\kern0pt}intro{\isacharbang}{\kern0pt}{\isacharbrackright}{\kern0pt}{\isacharcolon}{\kern0pt}\ \ \ {\isachardoublequoteopen}{\isasymlbrakk}P{\isadigit{1}}{\isasymturnstile}nabla{\isadigit{1}}{\isasymrightarrow}P{\isadigit{2}}{\isacharsemicolon}{\kern0pt}\ P{\isadigit{2}}{\isasymTurnstile}{\isacharparenleft}{\kern0pt}nabla{\isadigit{2}}{\isacharcomma}{\kern0pt}{\isacharbrackleft}{\kern0pt}{\isacharbrackright}{\kern0pt}{\isacharparenright}{\kern0pt}{\isasymRightarrow}P{\isadigit{3}}{\isasymrbrakk}\ {\isasymLongrightarrow}\ P{\isadigit{1}}{\isasymTurnstile}{\isacharparenleft}{\kern0pt}nabla{\isadigit{2}}{\isasymunion}nabla{\isadigit{1}}{\isacharcomma}{\kern0pt}{\isacharbrackleft}{\kern0pt}{\isacharbrackright}{\kern0pt}{\isacharparenright}{\kern0pt}{\isasymRightarrow}P{\isadigit{3}}{\isachardoublequoteclose}\isanewline
\isanewline
\isanewline
\isacommand{lemma}\isamarkupfalse%
\ mgu{\isacharunderscore}{\kern0pt}idem{\isacharcolon}{\kern0pt}\ \isanewline
\ \ \isakeyword{assumes}\ {\isachardoublequoteopen}{\isacharparenleft}{\kern0pt}nabla{\isadigit{1}}{\isacharcomma}{\kern0pt}s{\isadigit{1}}{\isacharparenright}{\kern0pt}\ {\isasymin}\ U\ P{\isachardoublequoteclose}\isanewline
\ \ \ \ \isakeyword{and}\ {\isachardoublequoteopen}{\isasymforall}{\isacharparenleft}{\kern0pt}nabla{\isadigit{2}}{\isacharcomma}{\kern0pt}s{\isadigit{2}}{\isacharparenright}{\kern0pt}\ {\isasymin}\ U\ P{\isachardot}{\kern0pt}\ nabla{\isadigit{2}}\ {\isasymTurnstile}{\isacharparenleft}{\kern0pt}subst\ s{\isadigit{2}}{\isacharparenright}{\kern0pt}\ nabla{\isadigit{1}}\ {\isasymand}\ \ nabla{\isadigit{2}}{\isasymTurnstile}subst{\isacharparenleft}{\kern0pt}s{\isadigit{2}}\ {\isasymbullet}\ s{\isadigit{1}}{\isacharparenright}{\kern0pt}{\isasymapprox}subst\ s{\isadigit{2}}{\isachardoublequoteclose}\isanewline
\ \ \isakeyword{shows}\ {\isachardoublequoteopen}mgu\ P\ {\isacharparenleft}{\kern0pt}nabla{\isadigit{1}}{\isacharcomma}{\kern0pt}s{\isadigit{1}}{\isacharparenright}{\kern0pt}\ {\isasymand}\ idem\ {\isacharparenleft}{\kern0pt}nabla{\isadigit{1}}{\isacharcomma}{\kern0pt}s{\isadigit{1}}{\isacharparenright}{\kern0pt}{\isachardoublequoteclose}\isanewline
%
\isadelimproof
\ \ %
\endisadelimproof
%
\isatagproof
\isacommand{using}\isamarkupfalse%
\ assms\ mgu{\isacharunderscore}{\kern0pt}def\ idem{\isacharunderscore}{\kern0pt}def\ \isacommand{by}\isamarkupfalse%
\ auto%
\endisatagproof
{\isafoldproof}%
%
\isadelimproof
\isanewline
%
\endisadelimproof
\isanewline
\isanewline
\isacommand{lemma}\isamarkupfalse%
\ problem{\isacharunderscore}{\kern0pt}subst{\isacharunderscore}{\kern0pt}comm{\isacharcolon}{\kern0pt}\ {\isachardoublequoteopen}{\isacharparenleft}{\kern0pt}{\isacharparenleft}{\kern0pt}nabla{\isacharcomma}{\kern0pt}s{\isadigit{2}}{\isacharparenright}{\kern0pt}\ {\isasymin}\ U\ {\isacharparenleft}{\kern0pt}apply{\isacharunderscore}{\kern0pt}subst\ s{\isadigit{1}}\ P{\isacharparenright}{\kern0pt}{\isacharparenright}{\kern0pt}\ {\isacharequal}{\kern0pt}\ {\isacharparenleft}{\kern0pt}{\isacharparenleft}{\kern0pt}nabla{\isacharcomma}{\kern0pt}{\isacharparenleft}{\kern0pt}s{\isadigit{2}}{\isasymbullet}s{\isadigit{1}}{\isacharparenright}{\kern0pt}{\isacharparenright}{\kern0pt}{\isasymin}U\ P{\isacharparenright}{\kern0pt}{\isachardoublequoteclose}\isanewline
%
\isadelimproof
\ \ %
\endisadelimproof
%
\isatagproof
\isacommand{using}\isamarkupfalse%
\ all{\isacharunderscore}{\kern0pt}solutions{\isacharunderscore}{\kern0pt}def\ apply{\isacharunderscore}{\kern0pt}subst{\isacharunderscore}{\kern0pt}def\ subst{\isacharunderscore}{\kern0pt}comp{\isacharunderscore}{\kern0pt}expand\ \isacommand{by}\isamarkupfalse%
\ auto%
\endisatagproof
{\isafoldproof}%
%
\isadelimproof
\isanewline
%
\endisadelimproof
\isanewline
\isanewline
\isanewline
\isacommand{lemma}\isamarkupfalse%
\ P{\isadigit{1}}{\isacharunderscore}{\kern0pt}to{\isacharunderscore}{\kern0pt}P{\isadigit{2}}{\isacharunderscore}{\kern0pt}sred{\isacharcolon}{\kern0pt}\isanewline
\ \ \isakeyword{assumes}\ {\isachardoublequoteopen}{\isacharparenleft}{\kern0pt}nabla{\isadigit{1}}{\isacharcomma}{\kern0pt}s{\isadigit{1}}{\isacharparenright}{\kern0pt}\ {\isasymin}\ U\ P{\isadigit{1}}{\isachardoublequoteclose}\ \isakeyword{and}\ {\isachardoublequoteopen}P{\isadigit{1}}\ {\isasymturnstile}s{\isadigit{2}}\ {\isasymleadsto}\ P{\isadigit{2}}{\isachardoublequoteclose}\isanewline
\ \ \isakeyword{shows}\ {\isachardoublequoteopen}{\isacharparenleft}{\kern0pt}nabla{\isadigit{1}}{\isacharcomma}{\kern0pt}s{\isadigit{1}}{\isacharparenright}{\kern0pt}\ {\isasymin}\ U\ P{\isadigit{2}}\ \ {\isasymand}\ {\isacharparenleft}{\kern0pt}nabla{\isadigit{1}}{\isasymTurnstile}subst\ {\isacharparenleft}{\kern0pt}s{\isadigit{1}}{\isasymbullet}s{\isadigit{2}}{\isacharparenright}{\kern0pt}{\isasymapprox}subst\ s{\isadigit{1}}{\isacharparenright}{\kern0pt}{\isachardoublequoteclose}\isanewline
%
\isadelimproof
\ \ %
\endisadelimproof
%
\isatagproof
\isacommand{using}\isamarkupfalse%
\ assms\isanewline
\isacommand{proof}\isamarkupfalse%
{\isacharparenleft}{\kern0pt}induction\ arbitrary{\isacharcolon}{\kern0pt}\ nabla{\isadigit{1}}\ s{\isadigit{1}}\ rule{\isacharcolon}{\kern0pt}\ s{\isacharunderscore}{\kern0pt}red{\isachardot}{\kern0pt}induct{\isacharbrackleft}{\kern0pt}OF\ {\isacartoucheopen}P{\isadigit{1}}\ {\isasymturnstile}\ s{\isadigit{2}}\ {\isasymleadsto}\ P{\isadigit{2}}{\isacartoucheclose}{\isacharbrackright}{\kern0pt}{\isacharparenright}{\kern0pt}\isanewline
\ \ \isacommand{case}\isamarkupfalse%
\ {\isacharparenleft}{\kern0pt}{\isadigit{1}}\ xs\ ys{\isacharparenright}{\kern0pt}\isanewline
\ \ \ \isacommand{then}\isamarkupfalse%
\ \isacommand{have}\isamarkupfalse%
\ subst{\isacharcolon}{\kern0pt}\ {\isachardoublequoteopen}nabla{\isadigit{1}}\ {\isasymTurnstile}\ subst\ {\isacharparenleft}{\kern0pt}s{\isadigit{1}}\ {\isasymbullet}\ {\isacharbrackleft}{\kern0pt}{\isacharbrackright}{\kern0pt}{\isacharparenright}{\kern0pt}\ {\isasymapprox}\ subst\ s{\isadigit{1}}{\isachardoublequoteclose}\ \isanewline
\ \ \ \ \ \isacommand{using}\isamarkupfalse%
\ equ{\isacharunderscore}{\kern0pt}refl\ subst{\isacharunderscore}{\kern0pt}equ{\isacharunderscore}{\kern0pt}def\ \isacommand{by}\isamarkupfalse%
\ simp\isanewline
\ \ \ \isacommand{from}\isamarkupfalse%
\ {\isadigit{1}}\isanewline
\ \ \ \isacommand{have}\isamarkupfalse%
\ eqs{\isacharcolon}{\kern0pt}\ {\isachardoublequoteopen}{\isasymforall}\ {\isacharparenleft}{\kern0pt}t{\isadigit{1}}{\isacharcomma}{\kern0pt}t{\isadigit{2}}{\isacharparenright}{\kern0pt}\ {\isasymin}\ set\ xs{\isachardot}{\kern0pt}\ nabla{\isadigit{1}}\ {\isasymturnstile}\ subst\ s{\isadigit{1}}\ t{\isadigit{1}}\ {\isasymapprox}\ subst\ s{\isadigit{1}}\ t{\isadigit{2}}{\isachardoublequoteclose}\ \isanewline
\ \ \ \ \ \isakeyword{and}\ fresh{\isacharcolon}{\kern0pt}\ {\isachardoublequoteopen}{\isasymforall}\ {\isacharparenleft}{\kern0pt}a{\isacharcomma}{\kern0pt}t{\isacharparenright}{\kern0pt}\ {\isasymin}\ set\ ys{\isachardot}{\kern0pt}\ nabla{\isadigit{1}}\ {\isasymturnstile}\ a\ {\isasymsharp}\ subst\ s{\isadigit{1}}\ t{\isachardoublequoteclose}\ \isanewline
\ \ \ \ \ \isacommand{using}\isamarkupfalse%
\ all{\isacharunderscore}{\kern0pt}solutions{\isacharunderscore}{\kern0pt}def\ \isacommand{by}\isamarkupfalse%
\ simp{\isacharunderscore}{\kern0pt}all\isanewline
\ \ \ \isacommand{with}\isamarkupfalse%
\ subst\ \isacommand{show}\isamarkupfalse%
\ {\isacharquery}{\kern0pt}case\isanewline
\ \ \ \ \ \isacommand{using}\isamarkupfalse%
\ all{\isacharunderscore}{\kern0pt}solutions{\isacharunderscore}{\kern0pt}def\ \isacommand{by}\isamarkupfalse%
\ auto\isanewline
\ \isacommand{next}\isamarkupfalse%
\isanewline
\ \ \ \isacommand{case}\isamarkupfalse%
\ {\isacharparenleft}{\kern0pt}{\isadigit{2}}\ t{\isadigit{1}}\ t{\isadigit{2}}\ t{\isadigit{3}}\ t{\isadigit{4}}\ xs\ ys{\isacharparenright}{\kern0pt}\isanewline
\ \ \ \isacommand{then}\isamarkupfalse%
\ \isacommand{have}\isamarkupfalse%
\ subst{\isacharcolon}{\kern0pt}\ {\isachardoublequoteopen}nabla{\isadigit{1}}\ {\isasymTurnstile}\ subst\ {\isacharparenleft}{\kern0pt}s{\isadigit{1}}{\isasymbullet}\ {\isacharbrackleft}{\kern0pt}{\isacharbrackright}{\kern0pt}{\isacharparenright}{\kern0pt}\ {\isasymapprox}\ subst\ s{\isadigit{1}}{\isachardoublequoteclose}\ \isanewline
\ \ \ \ \ \isacommand{using}\isamarkupfalse%
\ equ{\isacharunderscore}{\kern0pt}refl\ subst{\isacharunderscore}{\kern0pt}equ{\isacharunderscore}{\kern0pt}def\ \isacommand{by}\isamarkupfalse%
\ simp\isanewline
\isanewline
\ \ \ \isacommand{from}\isamarkupfalse%
\ {\isadigit{2}}\isanewline
\ \ \ \isacommand{have}\isamarkupfalse%
\ fresh{\isacharcolon}{\kern0pt}\ {\isachardoublequoteopen}{\isasymforall}\ {\isacharparenleft}{\kern0pt}a{\isacharcomma}{\kern0pt}t{\isacharparenright}{\kern0pt}\ {\isasymin}\ set\ ys{\isachardot}{\kern0pt}\ nabla{\isadigit{1}}\ {\isasymturnstile}\ a\ {\isasymsharp}\ subst\ s{\isadigit{1}}\ t{\isachardoublequoteclose}\isanewline
\ \ \ \ \ \isacommand{using}\isamarkupfalse%
\ all{\isacharunderscore}{\kern0pt}solutions{\isacharunderscore}{\kern0pt}def\ \isacommand{by}\isamarkupfalse%
\ simp\isanewline
\isanewline
\ \ \ \isacommand{from}\isamarkupfalse%
\ {\isadigit{2}}\isanewline
\ \ \ \isacommand{have}\isamarkupfalse%
\ i{\isacharcolon}{\kern0pt}\ {\isachardoublequoteopen}{\isasymforall}\ {\isacharparenleft}{\kern0pt}t{\isadigit{1}}{\isacharcomma}{\kern0pt}t{\isadigit{2}}{\isacharparenright}{\kern0pt}\ {\isasymin}\ set\ {\isacharparenleft}{\kern0pt}{\isacharparenleft}{\kern0pt}Paar\ t{\isadigit{1}}\ t{\isadigit{2}}{\isacharcomma}{\kern0pt}\ Paar\ t{\isadigit{3}}\ t{\isadigit{4}}{\isacharparenright}{\kern0pt}\ {\isacharhash}{\kern0pt}\ xs{\isacharparenright}{\kern0pt}{\isachardot}{\kern0pt}\ nabla{\isadigit{1}}\ {\isasymturnstile}\ subst\ s{\isadigit{1}}\ t{\isadigit{1}}\ {\isasymapprox}\ subst\ s{\isadigit{1}}\ t{\isadigit{2}}{\isachardoublequoteclose}\isanewline
\ \ \ \ \ \isacommand{using}\isamarkupfalse%
\ all{\isacharunderscore}{\kern0pt}solutions{\isacharunderscore}{\kern0pt}def\ \isacommand{by}\isamarkupfalse%
\ simp\isanewline
\ \ \ \isacommand{hence}\isamarkupfalse%
\ {\isachardoublequoteopen}nabla{\isadigit{1}}\ {\isasymturnstile}\ subst\ s{\isadigit{1}}\ t{\isadigit{1}}\ {\isasymapprox}\ subst\ s{\isadigit{1}}\ t{\isadigit{3}}{\isachardoublequoteclose}\ \isanewline
\ \ \ \ {\isachardoublequoteopen}nabla{\isadigit{1}}\ {\isasymturnstile}\ subst\ s{\isadigit{1}}\ t{\isadigit{2}}\ {\isasymapprox}\ subst\ s{\isadigit{1}}\ t{\isadigit{4}}{\isachardoublequoteclose}\isanewline
\ \ \ \ \isacommand{using}\isamarkupfalse%
\ Equ{\isacharunderscore}{\kern0pt}elims{\isacharparenleft}{\kern0pt}{\isadigit{4}}{\isacharparenright}{\kern0pt}{\isacharbrackleft}{\kern0pt}of\ nabla{\isadigit{1}}\ {\isacartoucheopen}{\isacharparenleft}{\kern0pt}subst\ s{\isadigit{1}}\ t{\isadigit{1}}{\isacharparenright}{\kern0pt}{\isacartoucheclose}\ \ {\isacartoucheopen}{\isacharparenleft}{\kern0pt}subst\ s{\isadigit{1}}\ t{\isadigit{2}}{\isacharparenright}{\kern0pt}{\isacartoucheclose}\ \ {\isacartoucheopen}{\isacharparenleft}{\kern0pt}subst\ s{\isadigit{1}}\ t{\isadigit{3}}{\isacharparenright}{\kern0pt}{\isacartoucheclose}\ \ {\isacartoucheopen}{\isacharparenleft}{\kern0pt}subst\ s{\isadigit{1}}\ t{\isadigit{4}}{\isacharparenright}{\kern0pt}{\isacartoucheclose}{\isacharbrackright}{\kern0pt}\ \isanewline
\ \ \ \ \ \ subst{\isacharunderscore}{\kern0pt}paar{\isacharbrackleft}{\kern0pt}of\ s{\isadigit{1}}{\isacharbrackright}{\kern0pt}\ \isacommand{by}\isamarkupfalse%
\ auto{\isacharplus}{\kern0pt}\isanewline
\ \ \ \ \isacommand{with}\isamarkupfalse%
\ i\ \isacommand{have}\isamarkupfalse%
\ eqs{\isacharcolon}{\kern0pt}\isanewline
\ \ \ \ {\isachardoublequoteopen}{\isasymforall}\ {\isacharparenleft}{\kern0pt}t{\isadigit{1}}{\isacharcomma}{\kern0pt}t{\isadigit{2}}{\isacharparenright}{\kern0pt}\ {\isasymin}\ set\ {\isacharparenleft}{\kern0pt}{\isacharparenleft}{\kern0pt}t{\isadigit{1}}{\isacharcomma}{\kern0pt}t{\isadigit{3}}{\isacharparenright}{\kern0pt}{\isacharhash}{\kern0pt}{\isacharparenleft}{\kern0pt}t{\isadigit{2}}{\isacharcomma}{\kern0pt}\ t{\isadigit{4}}{\isacharparenright}{\kern0pt}\ {\isacharhash}{\kern0pt}\ xs{\isacharparenright}{\kern0pt}{\isachardot}{\kern0pt}\ nabla{\isadigit{1}}\ {\isasymturnstile}\ subst\ s{\isadigit{1}}\ t{\isadigit{1}}\ {\isasymapprox}\ subst\ s{\isadigit{1}}\ t{\isadigit{2}}{\isachardoublequoteclose}\isanewline
\ \ \ \ \ \ \isacommand{by}\isamarkupfalse%
\ auto\isanewline
\isanewline
\ \ \ \ \isacommand{from}\isamarkupfalse%
\ fresh\ eqs\ subst\ \isanewline
\ \ \ \ \isacommand{show}\isamarkupfalse%
\ {\isacharquery}{\kern0pt}case\ \isacommand{using}\isamarkupfalse%
\ all{\isacharunderscore}{\kern0pt}solutions{\isacharunderscore}{\kern0pt}def\ \isacommand{by}\isamarkupfalse%
\ simp\isanewline
\ \isacommand{next}\isamarkupfalse%
\isanewline
\ \ \ \isacommand{case}\isamarkupfalse%
\ {\isacharparenleft}{\kern0pt}{\isadigit{3}}\ F\ t{\isadigit{1}}\ t{\isadigit{2}}\ xs\ ys{\isacharparenright}{\kern0pt}\isanewline
\ \ \ \isacommand{then}\isamarkupfalse%
\ \isacommand{have}\isamarkupfalse%
\ subst{\isacharcolon}{\kern0pt}\ {\isachardoublequoteopen}nabla{\isadigit{1}}\ {\isasymTurnstile}\ subst\ {\isacharparenleft}{\kern0pt}s{\isadigit{1}}\ {\isasymbullet}\ {\isacharbrackleft}{\kern0pt}{\isacharbrackright}{\kern0pt}{\isacharparenright}{\kern0pt}\ {\isasymapprox}\ subst\ s{\isadigit{1}}{\isachardoublequoteclose}\ \isanewline
\ \ \ \ \ \isacommand{using}\isamarkupfalse%
\ equ{\isacharunderscore}{\kern0pt}refl\ subst{\isacharunderscore}{\kern0pt}equ{\isacharunderscore}{\kern0pt}def\ \isacommand{by}\isamarkupfalse%
\ simp\isanewline
\isanewline
\ \ \ \isacommand{from}\isamarkupfalse%
\ {\isadigit{3}}\isanewline
\ \ \ \isacommand{have}\isamarkupfalse%
\ fresh{\isacharcolon}{\kern0pt}\ {\isachardoublequoteopen}{\isasymforall}\ {\isacharparenleft}{\kern0pt}a{\isacharcomma}{\kern0pt}t{\isacharparenright}{\kern0pt}\ {\isasymin}\ set\ ys{\isachardot}{\kern0pt}\ nabla{\isadigit{1}}\ {\isasymturnstile}\ a\ {\isasymsharp}\ subst\ s{\isadigit{1}}\ t{\isachardoublequoteclose}\isanewline
\ \ \ \ \ \isacommand{using}\isamarkupfalse%
\ all{\isacharunderscore}{\kern0pt}solutions{\isacharunderscore}{\kern0pt}def\ \isacommand{by}\isamarkupfalse%
\ simp\isanewline
\isanewline
\ \ \ \isacommand{from}\isamarkupfalse%
\ {\isadigit{3}}\isanewline
\ \ \ \isacommand{have}\isamarkupfalse%
\ i{\isacharcolon}{\kern0pt}\ {\isachardoublequoteopen}{\isasymforall}\ {\isacharparenleft}{\kern0pt}t{\isadigit{1}}{\isacharcomma}{\kern0pt}t{\isadigit{2}}{\isacharparenright}{\kern0pt}\ {\isasymin}\ set\ {\isacharparenleft}{\kern0pt}{\isacharparenleft}{\kern0pt}Func\ F\ t{\isadigit{1}}{\isacharcomma}{\kern0pt}\ Func\ F\ t{\isadigit{2}}{\isacharparenright}{\kern0pt}\ {\isacharhash}{\kern0pt}\ xs{\isacharparenright}{\kern0pt}{\isachardot}{\kern0pt}\ nabla{\isadigit{1}}\ {\isasymturnstile}\ subst\ s{\isadigit{1}}\ t{\isadigit{1}}\ {\isasymapprox}\ subst\ s{\isadigit{1}}\ t{\isadigit{2}}{\isachardoublequoteclose}\isanewline
\ \ \ \ \ \isacommand{using}\isamarkupfalse%
\ all{\isacharunderscore}{\kern0pt}solutions{\isacharunderscore}{\kern0pt}def\ \isacommand{by}\isamarkupfalse%
\ simp\isanewline
\ \ \ \isacommand{hence}\isamarkupfalse%
\ {\isachardoublequoteopen}nabla{\isadigit{1}}\ {\isasymturnstile}\ subst\ s{\isadigit{1}}\ t{\isadigit{1}}\ {\isasymapprox}\ subst\ s{\isadigit{1}}\ t{\isadigit{2}}{\isachardoublequoteclose}\ \isanewline
\ \ \ \ \ \isacommand{using}\isamarkupfalse%
\ Equ{\isacharunderscore}{\kern0pt}elims{\isacharparenleft}{\kern0pt}{\isadigit{5}}{\isacharparenright}{\kern0pt}{\isacharbrackleft}{\kern0pt}of\ nabla{\isadigit{1}}\ F\ {\isacartoucheopen}subst\ s{\isadigit{1}}\ t{\isadigit{1}}{\isacartoucheclose}\ {\isacartoucheopen}subst\ s{\isadigit{1}}\ t{\isadigit{2}}{\isacartoucheclose}{\isacharbrackright}{\kern0pt}\ \isanewline
\ \ \ \ \ \ subst{\isacharunderscore}{\kern0pt}func\ \isacommand{by}\isamarkupfalse%
\ auto\isanewline
\ \ \ \isacommand{with}\isamarkupfalse%
\ i\ \isacommand{have}\isamarkupfalse%
\ eqs{\isacharcolon}{\kern0pt}\ \isanewline
\ \ \ \ \ {\isachardoublequoteopen}{\isasymforall}\ {\isacharparenleft}{\kern0pt}t{\isadigit{1}}{\isacharcomma}{\kern0pt}t{\isadigit{2}}{\isacharparenright}{\kern0pt}\ {\isasymin}\ set\ {\isacharparenleft}{\kern0pt}{\isacharparenleft}{\kern0pt}t{\isadigit{1}}{\isacharcomma}{\kern0pt}\ t{\isadigit{2}}{\isacharparenright}{\kern0pt}\ {\isacharhash}{\kern0pt}\ xs{\isacharparenright}{\kern0pt}{\isachardot}{\kern0pt}\ nabla{\isadigit{1}}\ {\isasymturnstile}\ subst\ s{\isadigit{1}}\ t{\isadigit{1}}\ {\isasymapprox}\ subst\ s{\isadigit{1}}\ t{\isadigit{2}}{\isachardoublequoteclose}\isanewline
\ \ \ \ \ \isacommand{by}\isamarkupfalse%
\ auto\isanewline
\isanewline
\ \ \ \isacommand{from}\isamarkupfalse%
\ subst\ fresh\ eqs\isanewline
\ \ \ \isacommand{show}\isamarkupfalse%
\ {\isacharquery}{\kern0pt}case\ \isacommand{using}\isamarkupfalse%
\ all{\isacharunderscore}{\kern0pt}solutions{\isacharunderscore}{\kern0pt}def\ \isacommand{by}\isamarkupfalse%
\ simp\isanewline
\ \isacommand{next}\isamarkupfalse%
\isanewline
\ \ \ \isacommand{case}\isamarkupfalse%
\ {\isacharparenleft}{\kern0pt}{\isadigit{4}}\ a\ t{\isadigit{1}}\ t{\isadigit{2}}\ xs\ ys{\isacharparenright}{\kern0pt}\isanewline
\ \ \ \isacommand{then}\isamarkupfalse%
\ \isacommand{have}\isamarkupfalse%
\ subst{\isacharcolon}{\kern0pt}\ {\isachardoublequoteopen}nabla{\isadigit{1}}\ {\isasymTurnstile}\ subst\ {\isacharparenleft}{\kern0pt}s{\isadigit{1}}\ {\isasymbullet}\ {\isacharbrackleft}{\kern0pt}{\isacharbrackright}{\kern0pt}{\isacharparenright}{\kern0pt}\ {\isasymapprox}\ subst\ s{\isadigit{1}}{\isachardoublequoteclose}\ \isanewline
\ \ \ \ \ \isacommand{using}\isamarkupfalse%
\ equ{\isacharunderscore}{\kern0pt}refl\ subst{\isacharunderscore}{\kern0pt}equ{\isacharunderscore}{\kern0pt}def\ \isacommand{by}\isamarkupfalse%
\ simp\isanewline
\isanewline
\ \ \ \isacommand{from}\isamarkupfalse%
\ {\isadigit{4}}\isanewline
\ \ \ \isacommand{have}\isamarkupfalse%
\ fresh{\isacharcolon}{\kern0pt}\ {\isachardoublequoteopen}{\isasymforall}\ {\isacharparenleft}{\kern0pt}a{\isacharcomma}{\kern0pt}t{\isacharparenright}{\kern0pt}\ {\isasymin}\ set\ ys{\isachardot}{\kern0pt}\ nabla{\isadigit{1}}\ {\isasymturnstile}\ a\ {\isasymsharp}\ subst\ s{\isadigit{1}}\ t{\isachardoublequoteclose}\isanewline
\ \ \ \ \ \isacommand{using}\isamarkupfalse%
\ all{\isacharunderscore}{\kern0pt}solutions{\isacharunderscore}{\kern0pt}def\ \isacommand{by}\isamarkupfalse%
\ simp\isanewline
\isanewline
\ \ \ \isacommand{from}\isamarkupfalse%
\ {\isadigit{4}}\isanewline
\ \ \ \isacommand{have}\isamarkupfalse%
\ i{\isacharcolon}{\kern0pt}\ {\isachardoublequoteopen}{\isasymforall}\ {\isacharparenleft}{\kern0pt}t{\isadigit{1}}{\isacharcomma}{\kern0pt}t{\isadigit{2}}{\isacharparenright}{\kern0pt}\ {\isasymin}\ set\ {\isacharparenleft}{\kern0pt}{\isacharparenleft}{\kern0pt}Abst\ a\ t{\isadigit{1}}{\isacharcomma}{\kern0pt}\ Abst\ a\ t{\isadigit{2}}{\isacharparenright}{\kern0pt}\ {\isacharhash}{\kern0pt}\ xs{\isacharparenright}{\kern0pt}{\isachardot}{\kern0pt}\ nabla{\isadigit{1}}\ {\isasymturnstile}\ subst\ s{\isadigit{1}}\ t{\isadigit{1}}\ {\isasymapprox}\ subst\ s{\isadigit{1}}\ t{\isadigit{2}}{\isachardoublequoteclose}\isanewline
\ \ \ \ \ \isacommand{using}\isamarkupfalse%
\ all{\isacharunderscore}{\kern0pt}solutions{\isacharunderscore}{\kern0pt}def\ \isacommand{by}\isamarkupfalse%
\ simp\isanewline
\ \ \ \isacommand{hence}\isamarkupfalse%
\ {\isachardoublequoteopen}nabla{\isadigit{1}}\ {\isasymturnstile}\ subst\ s{\isadigit{1}}\ t{\isadigit{1}}\ {\isasymapprox}\ subst\ s{\isadigit{1}}\ t{\isadigit{2}}{\isachardoublequoteclose}\ \isanewline
\ \ \ \ \ \isacommand{using}\isamarkupfalse%
\ Equ{\isacharunderscore}{\kern0pt}elims{\isacharparenleft}{\kern0pt}{\isadigit{6}}{\isacharparenright}{\kern0pt}{\isacharbrackleft}{\kern0pt}of\ nabla{\isadigit{1}}\ a\ {\isacartoucheopen}subst\ s{\isadigit{1}}\ t{\isadigit{1}}{\isacartoucheclose}\ {\isacartoucheopen}subst\ s{\isadigit{1}}\ t{\isadigit{2}}{\isacartoucheclose}{\isacharbrackright}{\kern0pt}\ \isanewline
\ \ \ \ \ \ subst{\isacharunderscore}{\kern0pt}abst\ \isacommand{by}\isamarkupfalse%
\ auto\isanewline
\ \ \ \isacommand{with}\isamarkupfalse%
\ i\ \isacommand{have}\isamarkupfalse%
\ eqs{\isacharcolon}{\kern0pt}\ \isanewline
\ \ \ \ \ {\isachardoublequoteopen}{\isasymforall}\ {\isacharparenleft}{\kern0pt}t{\isadigit{1}}{\isacharcomma}{\kern0pt}t{\isadigit{2}}{\isacharparenright}{\kern0pt}\ {\isasymin}\ set\ {\isacharparenleft}{\kern0pt}{\isacharparenleft}{\kern0pt}t{\isadigit{1}}{\isacharcomma}{\kern0pt}\ t{\isadigit{2}}{\isacharparenright}{\kern0pt}\ {\isacharhash}{\kern0pt}\ xs{\isacharparenright}{\kern0pt}{\isachardot}{\kern0pt}\ nabla{\isadigit{1}}\ {\isasymturnstile}\ subst\ s{\isadigit{1}}\ t{\isadigit{1}}\ {\isasymapprox}\ subst\ s{\isadigit{1}}\ t{\isadigit{2}}{\isachardoublequoteclose}\isanewline
\ \ \ \ \ \isacommand{by}\isamarkupfalse%
\ auto\isanewline
\isanewline
\ \ \ \isacommand{from}\isamarkupfalse%
\ subst\ fresh\ eqs\isanewline
\ \ \ \isacommand{show}\isamarkupfalse%
\ {\isacharquery}{\kern0pt}case\ \isacommand{using}\isamarkupfalse%
\ all{\isacharunderscore}{\kern0pt}solutions{\isacharunderscore}{\kern0pt}def\ \isacommand{by}\isamarkupfalse%
\ simp\isanewline
\ \isacommand{next}\isamarkupfalse%
\isanewline
\ \ \ \isacommand{case}\isamarkupfalse%
\ {\isacharparenleft}{\kern0pt}{\isadigit{5}}\ a\ b\ t{\isadigit{1}}\ t{\isadigit{2}}\ xs\ ys{\isacharparenright}{\kern0pt}\isanewline
\ \ \ \isacommand{then}\isamarkupfalse%
\ \isacommand{have}\isamarkupfalse%
\ subst{\isacharcolon}{\kern0pt}\ {\isachardoublequoteopen}nabla{\isadigit{1}}\ {\isasymTurnstile}\ subst\ {\isacharparenleft}{\kern0pt}s{\isadigit{1}}\ {\isasymbullet}\ {\isacharbrackleft}{\kern0pt}{\isacharbrackright}{\kern0pt}{\isacharparenright}{\kern0pt}\ {\isasymapprox}\ subst\ s{\isadigit{1}}{\isachardoublequoteclose}\ \isanewline
\ \ \ \ \ \isacommand{using}\isamarkupfalse%
\ equ{\isacharunderscore}{\kern0pt}refl\ subst{\isacharunderscore}{\kern0pt}equ{\isacharunderscore}{\kern0pt}def\ \isacommand{by}\isamarkupfalse%
\ simp\isanewline
\isanewline
\ \ \ \isacommand{from}\isamarkupfalse%
\ {\isadigit{5}}\isanewline
\ \ \ \isacommand{have}\isamarkupfalse%
\ eqsP{\isadigit{1}}{\isacharcolon}{\kern0pt}\ {\isachardoublequoteopen}{\isasymforall}\ {\isacharparenleft}{\kern0pt}t{\isadigit{1}}{\isacharcomma}{\kern0pt}t{\isadigit{2}}{\isacharparenright}{\kern0pt}\ {\isasymin}\ set\ {\isacharparenleft}{\kern0pt}{\isacharparenleft}{\kern0pt}Abst\ a\ t{\isadigit{1}}{\isacharcomma}{\kern0pt}\ Abst\ b\ t{\isadigit{2}}{\isacharparenright}{\kern0pt}\ {\isacharhash}{\kern0pt}\ xs{\isacharparenright}{\kern0pt}{\isachardot}{\kern0pt}\ nabla{\isadigit{1}}\ {\isasymturnstile}\ subst\ s{\isadigit{1}}\ t{\isadigit{1}}\ {\isasymapprox}\ subst\ s{\isadigit{1}}\ t{\isadigit{2}}{\isachardoublequoteclose}\isanewline
\ \ \ \ \ \isacommand{using}\isamarkupfalse%
\ all{\isacharunderscore}{\kern0pt}solutions{\isacharunderscore}{\kern0pt}def\ \isacommand{by}\isamarkupfalse%
\ simp\isanewline
\ \ \ \isacommand{hence}\isamarkupfalse%
\ i{\isacharcolon}{\kern0pt}\ {\isachardoublequoteopen}nabla{\isadigit{1}}\ {\isasymturnstile}\ Abst\ a\ {\isacharparenleft}{\kern0pt}subst\ s{\isadigit{1}}\ t{\isadigit{1}}{\isacharparenright}{\kern0pt}\ {\isasymapprox}\ Abst\ b\ {\isacharparenleft}{\kern0pt}subst\ s{\isadigit{1}}\ t{\isadigit{2}}{\isacharparenright}{\kern0pt}{\isachardoublequoteclose}\isanewline
\ \ \ \ \ \isacommand{using}\isamarkupfalse%
\ subst{\isacharunderscore}{\kern0pt}abst{\isacharbrackleft}{\kern0pt}of\ s{\isadigit{1}}{\isacharbrackright}{\kern0pt}\ \isacommand{by}\isamarkupfalse%
\ simp\isanewline
\ \ \ \isanewline
\ \ \ \isacommand{from}\isamarkupfalse%
\ {\isadigit{5}}\isanewline
\ \ \ \isacommand{have}\isamarkupfalse%
\ ii{\isacharcolon}{\kern0pt}\ {\isachardoublequoteopen}{\isasymforall}\ {\isacharparenleft}{\kern0pt}a{\isacharcomma}{\kern0pt}t{\isacharparenright}{\kern0pt}\ {\isasymin}\ set\ ys{\isachardot}{\kern0pt}\ nabla{\isadigit{1}}\ {\isasymturnstile}\ a\ {\isasymsharp}\ subst\ s{\isadigit{1}}\ t{\isachardoublequoteclose}\isanewline
\ \ \ \ \ \isacommand{using}\isamarkupfalse%
\ all{\isacharunderscore}{\kern0pt}solutions{\isacharunderscore}{\kern0pt}def\ \isacommand{by}\isamarkupfalse%
\ auto\isanewline
\ \ \ \isacommand{from}\isamarkupfalse%
\ {\isadigit{5}}\ i\isanewline
\ \ \ \isacommand{have}\isamarkupfalse%
\ {\isachardoublequoteopen}nabla{\isadigit{1}}\ {\isasymturnstile}\ a\ {\isasymsharp}\ subst\ s{\isadigit{1}}\ t{\isadigit{2}}{\isachardoublequoteclose}\isanewline
\ \ \ \ \ \isacommand{using}\isamarkupfalse%
\ Equ{\isacharunderscore}{\kern0pt}elims{\isacharparenleft}{\kern0pt}{\isadigit{7}}{\isacharparenright}{\kern0pt}\ \isacommand{by}\isamarkupfalse%
\ blast\isanewline
\ \ \ \isacommand{with}\isamarkupfalse%
\ ii\ \isanewline
\ \ \ \isacommand{have}\isamarkupfalse%
\ fresh{\isacharcolon}{\kern0pt}\ {\isachardoublequoteopen}{\isasymforall}\ {\isacharparenleft}{\kern0pt}a{\isacharcomma}{\kern0pt}t{\isacharparenright}{\kern0pt}\ {\isasymin}\ set\ {\isacharparenleft}{\kern0pt}{\isacharparenleft}{\kern0pt}a{\isacharcomma}{\kern0pt}t{\isadigit{2}}{\isacharparenright}{\kern0pt}\ {\isacharhash}{\kern0pt}\ ys{\isacharparenright}{\kern0pt}{\isachardot}{\kern0pt}\ nabla{\isadigit{1}}\ {\isasymturnstile}\ a\ {\isasymsharp}\ subst\ s{\isadigit{1}}\ t{\isachardoublequoteclose}\isanewline
\ \ \ \ \ \isacommand{by}\isamarkupfalse%
\ simp\isanewline
\ \ \ \isanewline
\ \ \ \isacommand{from}\isamarkupfalse%
\ i\ {\isadigit{5}}\isanewline
\ \ \ \isacommand{have}\isamarkupfalse%
\ {\isachardoublequoteopen}nabla{\isadigit{1}}\ {\isasymturnstile}\ subst\ s{\isadigit{1}}\ t{\isadigit{1}}\ {\isasymapprox}\ subst\ s{\isadigit{1}}\ {\isacharparenleft}{\kern0pt}swap\ {\isacharbrackleft}{\kern0pt}{\isacharparenleft}{\kern0pt}a{\isacharcomma}{\kern0pt}\ b{\isacharparenright}{\kern0pt}{\isacharbrackright}{\kern0pt}\ t{\isadigit{2}}{\isacharparenright}{\kern0pt}{\isachardoublequoteclose}\ \isanewline
\ \ \ \ \ \isacommand{using}\isamarkupfalse%
\ {\isadigit{5}}\ Equ{\isacharunderscore}{\kern0pt}elims{\isacharparenleft}{\kern0pt}{\isadigit{7}}{\isacharparenright}{\kern0pt}{\isacharbrackleft}{\kern0pt}of\ nabla{\isadigit{1}}\ a\ {\isacartoucheopen}subst\ s{\isadigit{1}}\ t{\isadigit{1}}{\isacartoucheclose}\ b\ {\isacartoucheopen}subst\ s{\isadigit{1}}\ t{\isadigit{2}}{\isacartoucheclose}{\isacharbrackright}{\kern0pt}\ \ \isanewline
\ \ \ \ \ \ \ subst{\isacharunderscore}{\kern0pt}swap{\isacharunderscore}{\kern0pt}comm\ \isacommand{by}\isamarkupfalse%
\ simp\isanewline
\ \ \ \isacommand{with}\isamarkupfalse%
\ eqsP{\isadigit{1}}\ \isacommand{have}\isamarkupfalse%
\ eqs{\isacharcolon}{\kern0pt}\ \isanewline
\ \ \ \ \ {\isachardoublequoteopen}{\isasymforall}\ {\isacharparenleft}{\kern0pt}t{\isadigit{1}}{\isacharcomma}{\kern0pt}t{\isadigit{2}}{\isacharparenright}{\kern0pt}\ {\isasymin}\ set\ {\isacharparenleft}{\kern0pt}{\isacharparenleft}{\kern0pt}t{\isadigit{1}}{\isacharcomma}{\kern0pt}\ swap\ {\isacharbrackleft}{\kern0pt}{\isacharparenleft}{\kern0pt}a{\isacharcomma}{\kern0pt}b{\isacharparenright}{\kern0pt}{\isacharbrackright}{\kern0pt}\ t{\isadigit{2}}{\isacharparenright}{\kern0pt}\ {\isacharhash}{\kern0pt}\ xs{\isacharparenright}{\kern0pt}{\isachardot}{\kern0pt}\ nabla{\isadigit{1}}\ {\isasymturnstile}\ subst\ s{\isadigit{1}}\ t{\isadigit{1}}\ {\isasymapprox}\ subst\ s{\isadigit{1}}\ t{\isadigit{2}}{\isachardoublequoteclose}\isanewline
\ \ \ \ \ \isacommand{by}\isamarkupfalse%
\ auto\isanewline
\ \ \ \ \isanewline
\ \ \ \isacommand{from}\isamarkupfalse%
\ subst\ fresh\ eqs\ \isanewline
\ \ \ \isacommand{show}\isamarkupfalse%
\ {\isacharquery}{\kern0pt}case\ \isacommand{using}\isamarkupfalse%
\ all{\isacharunderscore}{\kern0pt}solutions{\isacharunderscore}{\kern0pt}def\ \isacommand{by}\isamarkupfalse%
\ simp\isanewline
\ \isacommand{next}\isamarkupfalse%
\isanewline
\ \ \ \isacommand{case}\isamarkupfalse%
\ {\isacharparenleft}{\kern0pt}{\isadigit{6}}\ a\ xs\ ys{\isacharparenright}{\kern0pt}\isanewline
\ \ \ \isacommand{then}\isamarkupfalse%
\ \isacommand{have}\isamarkupfalse%
\ subst{\isacharcolon}{\kern0pt}\ {\isachardoublequoteopen}nabla{\isadigit{1}}\ {\isasymTurnstile}\ subst\ {\isacharparenleft}{\kern0pt}s{\isadigit{1}}\ {\isasymbullet}\ {\isacharbrackleft}{\kern0pt}{\isacharbrackright}{\kern0pt}{\isacharparenright}{\kern0pt}\ {\isasymapprox}\ subst\ s{\isadigit{1}}{\isachardoublequoteclose}\ \isanewline
\ \ \ \ \ \isacommand{using}\isamarkupfalse%
\ equ{\isacharunderscore}{\kern0pt}refl\ subst{\isacharunderscore}{\kern0pt}equ{\isacharunderscore}{\kern0pt}def\ \isacommand{by}\isamarkupfalse%
\ simp\isanewline
\isanewline
\ \ \ \isacommand{from}\isamarkupfalse%
\ {\isadigit{6}}\isanewline
\ \ \ \isacommand{have}\isamarkupfalse%
\ fresh{\isacharcolon}{\kern0pt}\ {\isachardoublequoteopen}{\isasymforall}\ {\isacharparenleft}{\kern0pt}a{\isacharcomma}{\kern0pt}t{\isacharparenright}{\kern0pt}\ {\isasymin}\ set\ ys{\isachardot}{\kern0pt}\ nabla{\isadigit{1}}\ {\isasymturnstile}\ a\ {\isasymsharp}\ subst\ s{\isadigit{1}}\ t{\isachardoublequoteclose}\isanewline
\ \ \ \ \ \isacommand{using}\isamarkupfalse%
\ all{\isacharunderscore}{\kern0pt}solutions{\isacharunderscore}{\kern0pt}def\ \isacommand{by}\isamarkupfalse%
\ auto\isanewline
\isanewline
\ \ \ \isacommand{from}\isamarkupfalse%
\ {\isadigit{6}}\isanewline
\ \ \ \isacommand{have}\isamarkupfalse%
\ eqs{\isacharcolon}{\kern0pt}\ {\isachardoublequoteopen}{\isasymforall}\ {\isacharparenleft}{\kern0pt}t{\isadigit{1}}{\isacharcomma}{\kern0pt}t{\isadigit{2}}{\isacharparenright}{\kern0pt}\ {\isasymin}\ set\ xs{\isachardot}{\kern0pt}\ nabla{\isadigit{1}}\ {\isasymturnstile}\ subst\ s{\isadigit{1}}\ t{\isadigit{1}}\ {\isasymapprox}\ subst\ s{\isadigit{1}}\ t{\isadigit{2}}{\isachardoublequoteclose}\isanewline
\ \ \ \ \ \isacommand{using}\isamarkupfalse%
\ all{\isacharunderscore}{\kern0pt}solutions{\isacharunderscore}{\kern0pt}def\ \isacommand{by}\isamarkupfalse%
\ auto\isanewline
\ \ \isanewline
\ \ \ \isacommand{from}\isamarkupfalse%
\ subst\ fresh\ eqs\isanewline
\ \ \ \isacommand{show}\isamarkupfalse%
\ {\isacharquery}{\kern0pt}case\ \isacommand{using}\isamarkupfalse%
\ all{\isacharunderscore}{\kern0pt}solutions{\isacharunderscore}{\kern0pt}def\ \isacommand{by}\isamarkupfalse%
\ simp\isanewline
\ \isacommand{next}\isamarkupfalse%
\isanewline
\ \ \ \isacommand{case}\isamarkupfalse%
\ {\isacharparenleft}{\kern0pt}{\isadigit{7}}\ pi{\isadigit{1}}\ X\ pi{\isadigit{2}}\ xs\ ys{\isacharparenright}{\kern0pt}\isanewline
\ \ \ \isacommand{then}\isamarkupfalse%
\ \isacommand{have}\isamarkupfalse%
\ subst{\isacharcolon}{\kern0pt}\ {\isachardoublequoteopen}nabla{\isadigit{1}}\ {\isasymTurnstile}\ subst\ {\isacharparenleft}{\kern0pt}s{\isadigit{1}}\ {\isasymbullet}\ {\isacharbrackleft}{\kern0pt}{\isacharbrackright}{\kern0pt}{\isacharparenright}{\kern0pt}\ {\isasymapprox}\ subst\ s{\isadigit{1}}{\isachardoublequoteclose}\ \isanewline
\ \ \ \ \ \isacommand{using}\isamarkupfalse%
\ equ{\isacharunderscore}{\kern0pt}refl\ subst{\isacharunderscore}{\kern0pt}equ{\isacharunderscore}{\kern0pt}def\ \isacommand{by}\isamarkupfalse%
\ simp\isanewline
\isanewline
\ \ \ \isacommand{from}\isamarkupfalse%
\ {\isadigit{7}}\isanewline
\ \ \ \isacommand{have}\isamarkupfalse%
\ i{\isacharcolon}{\kern0pt}\ {\isachardoublequoteopen}{\isasymforall}\ {\isacharparenleft}{\kern0pt}t{\isadigit{1}}{\isacharcomma}{\kern0pt}t{\isadigit{2}}{\isacharparenright}{\kern0pt}\ {\isasymin}\ set\ {\isacharparenleft}{\kern0pt}{\isacharparenleft}{\kern0pt}Susp\ pi{\isadigit{1}}\ X{\isacharcomma}{\kern0pt}\ Susp\ pi{\isadigit{2}}\ X{\isacharparenright}{\kern0pt}\ {\isacharhash}{\kern0pt}\ xs{\isacharparenright}{\kern0pt}{\isachardot}{\kern0pt}\ nabla{\isadigit{1}}\ {\isasymturnstile}\ subst\ s{\isadigit{1}}\ t{\isadigit{1}}\ {\isasymapprox}\ subst\ s{\isadigit{1}}\ t{\isadigit{2}}{\isachardoublequoteclose}\isanewline
\ \ \ \ \ {\isachardoublequoteopen}{\isasymforall}\ {\isacharparenleft}{\kern0pt}a{\isacharcomma}{\kern0pt}t{\isacharparenright}{\kern0pt}\ {\isasymin}\ set\ ys{\isachardot}{\kern0pt}\ nabla{\isadigit{1}}\ {\isasymturnstile}\ a\ {\isasymsharp}\ subst\ s{\isadigit{1}}\ t{\isachardoublequoteclose}\isanewline
\ \ \ \ \ \isacommand{using}\isamarkupfalse%
\ all{\isacharunderscore}{\kern0pt}solutions{\isacharunderscore}{\kern0pt}def\ \isacommand{by}\isamarkupfalse%
\ simp{\isacharplus}{\kern0pt}\isanewline
\isanewline
\isanewline
\ \ \ \isacommand{from}\isamarkupfalse%
\ {\isadigit{7}}{\isacharparenleft}{\kern0pt}{\isadigit{1}}{\isacharparenright}{\kern0pt}\isanewline
\ \ \ \isacommand{have}\isamarkupfalse%
\ {\isachardoublequoteopen}nabla{\isadigit{1}}\ {\isasymturnstile}\ swap\ pi{\isadigit{1}}\ {\isacharparenleft}{\kern0pt}look{\isacharunderscore}{\kern0pt}up\ X\ s{\isadigit{1}}{\isacharparenright}{\kern0pt}\ {\isasymapprox}\ swap\ pi{\isadigit{2}}\ {\isacharparenleft}{\kern0pt}look{\isacharunderscore}{\kern0pt}up\ X\ s{\isadigit{1}}{\isacharparenright}{\kern0pt}{\isachardoublequoteclose}\isanewline
\ \ \ \ \ \isacommand{using}\isamarkupfalse%
\ all{\isacharunderscore}{\kern0pt}solutions{\isacharunderscore}{\kern0pt}def\ subst{\isacharunderscore}{\kern0pt}susp\ \isacommand{by}\isamarkupfalse%
\ simp\isanewline
\ \ \ \isacommand{hence}\isamarkupfalse%
\ {\isachardoublequoteopen}{\isasymforall}a{\isasymin}ds\ pi{\isadigit{1}}\ pi{\isadigit{2}}{\isachardot}{\kern0pt}\ \ nabla{\isadigit{1}}\ {\isasymturnstile}\ a\ {\isasymsharp}\ {\isacharparenleft}{\kern0pt}look{\isacharunderscore}{\kern0pt}up\ X\ s{\isadigit{1}}{\isacharparenright}{\kern0pt}{\isachardoublequoteclose}\isanewline
\ \ \ \ \ \isacommand{using}\isamarkupfalse%
\ equ{\isacharunderscore}{\kern0pt}pi{\isadigit{1}}{\isacharunderscore}{\kern0pt}pi{\isadigit{2}}{\isacharunderscore}{\kern0pt}dec\ \isacommand{by}\isamarkupfalse%
\ simp\isanewline
\ \ \ \isacommand{hence}\isamarkupfalse%
\ {\isachardoublequoteopen}{\isasymforall}a{\isasymin}\ set\ {\isacharparenleft}{\kern0pt}ds{\isacharunderscore}{\kern0pt}list\ pi{\isadigit{1}}\ pi{\isadigit{2}}{\isacharparenright}{\kern0pt}{\isachardot}{\kern0pt}\ nabla{\isadigit{1}}\ {\isasymturnstile}\ a\ {\isasymsharp}\ subst\ s{\isadigit{1}}\ {\isacharparenleft}{\kern0pt}Susp\ {\isacharbrackleft}{\kern0pt}{\isacharbrackright}{\kern0pt}\ X{\isacharparenright}{\kern0pt}{\isachardoublequoteclose}\isanewline
\ \ \ \ \ \isacommand{using}\isamarkupfalse%
\ subst{\isacharunderscore}{\kern0pt}susp\ ds{\isacharunderscore}{\kern0pt}list{\isacharunderscore}{\kern0pt}equ{\isacharunderscore}{\kern0pt}ds{\isacharbrackleft}{\kern0pt}of\ pi{\isadigit{1}}\ pi{\isadigit{2}}{\isacharbrackright}{\kern0pt}\ \isacommand{by}\isamarkupfalse%
\ simp\isanewline
\ \ \ \isacommand{hence}\isamarkupfalse%
\ {\isachardoublequoteopen}{\isasymforall}\ {\isacharparenleft}{\kern0pt}a{\isacharcomma}{\kern0pt}t{\isacharparenright}{\kern0pt}\ {\isasymin}\ set\ {\isacharparenleft}{\kern0pt}map\ {\isacharparenleft}{\kern0pt}{\isasymlambda}a{\isachardot}{\kern0pt}\ a\ {\isasymsharp}{\isacharquery}{\kern0pt}\ Susp\ {\isacharbrackleft}{\kern0pt}{\isacharbrackright}{\kern0pt}\ X{\isacharparenright}{\kern0pt}\ {\isacharparenleft}{\kern0pt}ds{\isacharunderscore}{\kern0pt}list\ pi{\isadigit{1}}\ pi{\isadigit{2}}{\isacharparenright}{\kern0pt}{\isacharparenright}{\kern0pt}{\isachardot}{\kern0pt}\ nabla{\isadigit{1}}\ {\isasymturnstile}\ a\ {\isasymsharp}\ subst\ s{\isadigit{1}}\ t{\isachardoublequoteclose}\isanewline
\ \ \ \ \isacommand{by}\isamarkupfalse%
\ {\isacharparenleft}{\kern0pt}auto\ split{\isacharcolon}{\kern0pt}\ prod{\isachardot}{\kern0pt}splits{\isacharparenright}{\kern0pt}\isanewline
\ \ \ \isacommand{with}\isamarkupfalse%
\ {\isadigit{7}}\ i{\isacharparenleft}{\kern0pt}{\isadigit{2}}{\isacharparenright}{\kern0pt}\isanewline
\ \ \ \isacommand{have}\isamarkupfalse%
\ fresh{\isacharcolon}{\kern0pt}\ \isanewline
\ \ \ \ \ {\isachardoublequoteopen}{\isasymforall}\ {\isacharparenleft}{\kern0pt}a{\isacharcomma}{\kern0pt}t{\isacharparenright}{\kern0pt}\ {\isasymin}\ set\ {\isacharparenleft}{\kern0pt}map\ {\isacharparenleft}{\kern0pt}{\isasymlambda}a{\isachardot}{\kern0pt}\ {\isacharparenleft}{\kern0pt}a{\isacharcomma}{\kern0pt}\ Susp\ {\isacharbrackleft}{\kern0pt}{\isacharbrackright}{\kern0pt}\ X{\isacharparenright}{\kern0pt}{\isacharparenright}{\kern0pt}\ {\isacharparenleft}{\kern0pt}ds{\isacharunderscore}{\kern0pt}list\ pi{\isadigit{1}}\ pi{\isadigit{2}}{\isacharparenright}{\kern0pt}\ {\isacharat}{\kern0pt}\ ys{\isacharparenright}{\kern0pt}{\isachardot}{\kern0pt}\ nabla{\isadigit{1}}\ {\isasymturnstile}\ a\ {\isasymsharp}\ subst\ s{\isadigit{1}}\ t{\isachardoublequoteclose}\isanewline
\ \ \ \ \isacommand{by}\isamarkupfalse%
\ auto\isanewline
\isanewline
\ \ \ \isacommand{from}\isamarkupfalse%
\ i{\isacharparenleft}{\kern0pt}{\isadigit{1}}{\isacharparenright}{\kern0pt}\isanewline
\ \ \ \isacommand{have}\isamarkupfalse%
\ eqs{\isacharcolon}{\kern0pt}\ \ \isanewline
\ \ \ \ {\isachardoublequoteopen}{\isasymforall}\ {\isacharparenleft}{\kern0pt}t{\isadigit{1}}{\isacharcomma}{\kern0pt}t{\isadigit{2}}{\isacharparenright}{\kern0pt}\ {\isasymin}\ set\ xs{\isachardot}{\kern0pt}\ nabla{\isadigit{1}}\ {\isasymturnstile}\ subst\ s{\isadigit{1}}\ t{\isadigit{1}}\ {\isasymapprox}\ subst\ s{\isadigit{1}}\ t{\isadigit{2}}{\isachardoublequoteclose}\isanewline
\ \ \ \ \ \isacommand{by}\isamarkupfalse%
\ simp\isanewline
\ \ \isanewline
\ \ \ \isacommand{from}\isamarkupfalse%
\ subst\ fresh\ eqs\isanewline
\ \ \ \isacommand{show}\isamarkupfalse%
\ {\isacharquery}{\kern0pt}case\ \isacommand{using}\isamarkupfalse%
\ all{\isacharunderscore}{\kern0pt}solutions{\isacharunderscore}{\kern0pt}def\ \isacommand{by}\isamarkupfalse%
\ simp\isanewline
\ \isacommand{next}\isamarkupfalse%
\isanewline
\ \ \ \isacommand{case}\isamarkupfalse%
\ {\isacharparenleft}{\kern0pt}{\isadigit{8}}\ X\ t{\isacharprime}{\kern0pt}\ pi\ xs\ ys{\isacharparenright}{\kern0pt}\isanewline
\ \ \ \isacommand{hence}\isamarkupfalse%
\ {\isachardoublequoteopen}nabla{\isadigit{1}}\ {\isasymturnstile}\ subst\ s{\isadigit{1}}\ {\isacharparenleft}{\kern0pt}Susp\ pi\ X{\isacharparenright}{\kern0pt}\ {\isasymapprox}\ subst\ s{\isadigit{1}}\ t{\isacharprime}{\kern0pt}{\isachardoublequoteclose}\isanewline
\ \ \ \ \ \isacommand{using}\isamarkupfalse%
\ all{\isacharunderscore}{\kern0pt}solutions{\isacharunderscore}{\kern0pt}def{\isacharbrackleft}{\kern0pt}of\ {\isacartoucheopen}{\isacharparenleft}{\kern0pt}{\isacharparenleft}{\kern0pt}Susp\ pi\ X{\isacharcomma}{\kern0pt}\ t{\isacharprime}{\kern0pt}{\isacharparenright}{\kern0pt}\ {\isacharhash}{\kern0pt}\ xs{\isacharcomma}{\kern0pt}\ ys{\isacharparenright}{\kern0pt}{\isacartoucheclose}{\isacharbrackright}{\kern0pt}\ \isacommand{by}\isamarkupfalse%
\ auto\isanewline
\ \ \ \isacommand{hence}\isamarkupfalse%
\ subst{\isacharcolon}{\kern0pt}\ {\isachardoublequoteopen}nabla{\isadigit{1}}\ {\isasymTurnstile}\ subst\ {\isacharparenleft}{\kern0pt}s{\isadigit{1}}\ {\isasymbullet}\ {\isacharbrackleft}{\kern0pt}{\isacharparenleft}{\kern0pt}X{\isacharcomma}{\kern0pt}\ swap\ {\isacharparenleft}{\kern0pt}rev\ pi{\isacharparenright}{\kern0pt}\ t{\isacharprime}{\kern0pt}{\isacharparenright}{\kern0pt}{\isacharbrackright}{\kern0pt}{\isacharparenright}{\kern0pt}\ {\isasymapprox}\ subst\ s{\isadigit{1}}{\isachardoublequoteclose}\isanewline
\ \ \ \ \ \isacommand{using}\isamarkupfalse%
\ unif{\isacharunderscore}{\kern0pt}{\isadigit{1}}\ \isacommand{by}\isamarkupfalse%
\ simp\isanewline
\isanewline
\ \ \ \isacommand{from}\isamarkupfalse%
\ {\isadigit{8}}\isanewline
\ \ \ \isacommand{have}\isamarkupfalse%
\ eqs{\isacharunderscore}{\kern0pt}old{\isacharcolon}{\kern0pt}\isanewline
\ \ \ \ \ {\isachardoublequoteopen}{\isasymforall}{\isacharparenleft}{\kern0pt}t{\isadigit{1}}{\isacharcomma}{\kern0pt}t{\isadigit{2}}{\isacharparenright}{\kern0pt}{\isasymin}set\ xs{\isachardot}{\kern0pt}\ nabla{\isadigit{1}}\ {\isasymturnstile}\ subst\ s{\isadigit{1}}\ t{\isadigit{1}}\ {\isasymapprox}\ subst\ s{\isadigit{1}}\ t{\isadigit{2}}{\isachardoublequoteclose}\isanewline
\ \ \ \ \ \isakeyword{and}\ fresh{\isacharunderscore}{\kern0pt}old{\isacharcolon}{\kern0pt}\isanewline
\ \ \ \ \ {\isachardoublequoteopen}{\isasymforall}{\isacharparenleft}{\kern0pt}a{\isacharcomma}{\kern0pt}t{\isacharparenright}{\kern0pt}{\isasymin}set\ ys{\isachardot}{\kern0pt}\ nabla{\isadigit{1}}\ {\isasymturnstile}\ a\ {\isasymsharp}\ subst\ s{\isadigit{1}}\ t{\isachardoublequoteclose}\isanewline
\ \ \ \ \ \isacommand{using}\isamarkupfalse%
\ all{\isacharunderscore}{\kern0pt}solutions{\isacharunderscore}{\kern0pt}def\ \isacommand{by}\isamarkupfalse%
\ auto\isanewline
\isanewline
\ \ \ \isacommand{from}\isamarkupfalse%
\ eqs{\isacharunderscore}{\kern0pt}old\ \isanewline
\ \ \ \isacommand{have}\isamarkupfalse%
\ eqs{\isacharcolon}{\kern0pt}\ {\isachardoublequoteopen}{\isasymforall}{\isacharparenleft}{\kern0pt}t{\isadigit{1}}{\isacharcomma}{\kern0pt}t{\isadigit{2}}{\isacharparenright}{\kern0pt}{\isasymin}set\ xs{\isachardot}{\kern0pt}\ \isanewline
\ \ \ \ nabla{\isadigit{1}}\ {\isasymturnstile}\ subst\ {\isacharparenleft}{\kern0pt}s{\isadigit{1}}\ {\isasymbullet}\ {\isacharbrackleft}{\kern0pt}{\isacharparenleft}{\kern0pt}X{\isacharcomma}{\kern0pt}\ swap\ {\isacharparenleft}{\kern0pt}rev\ pi{\isacharparenright}{\kern0pt}\ t{\isacharprime}{\kern0pt}{\isacharparenright}{\kern0pt}{\isacharbrackright}{\kern0pt}{\isacharparenright}{\kern0pt}\ t{\isadigit{1}}\ {\isasymapprox}\ \ subst\ {\isacharparenleft}{\kern0pt}s{\isadigit{1}}\ {\isasymbullet}\ {\isacharbrackleft}{\kern0pt}{\isacharparenleft}{\kern0pt}X{\isacharcomma}{\kern0pt}\ swap\ {\isacharparenleft}{\kern0pt}rev\ pi{\isacharparenright}{\kern0pt}\ t{\isacharprime}{\kern0pt}{\isacharparenright}{\kern0pt}{\isacharbrackright}{\kern0pt}{\isacharparenright}{\kern0pt}\ t{\isadigit{2}}{\isachardoublequoteclose}\isanewline
\ \ \ \ \ \isacommand{using}\isamarkupfalse%
\ unif{\isacharunderscore}{\kern0pt}{\isadigit{2}}a{\isacharbrackleft}{\kern0pt}OF\ subst{\isacharbrackright}{\kern0pt}\ \isacommand{by}\isamarkupfalse%
\ auto\isanewline
\isanewline
\ \ \ \isacommand{from}\isamarkupfalse%
\ fresh{\isacharunderscore}{\kern0pt}old\ \isanewline
\ \ \ \isacommand{have}\isamarkupfalse%
\ fresh{\isacharcolon}{\kern0pt}\ {\isachardoublequoteopen}{\isasymforall}{\isacharparenleft}{\kern0pt}a{\isacharcomma}{\kern0pt}t{\isacharparenright}{\kern0pt}{\isasymin}set\ ys{\isachardot}{\kern0pt}\ nabla{\isadigit{1}}\ {\isasymturnstile}\ a\ {\isasymsharp}\ subst\ {\isacharparenleft}{\kern0pt}s{\isadigit{1}}\ {\isasymbullet}\ {\isacharbrackleft}{\kern0pt}{\isacharparenleft}{\kern0pt}X{\isacharcomma}{\kern0pt}\ swap\ {\isacharparenleft}{\kern0pt}rev\ pi{\isacharparenright}{\kern0pt}\ t{\isacharprime}{\kern0pt}{\isacharparenright}{\kern0pt}{\isacharbrackright}{\kern0pt}{\isacharparenright}{\kern0pt}\ t{\isachardoublequoteclose}\isanewline
\ \ \ \ \ \isacommand{using}\isamarkupfalse%
\ unif{\isacharunderscore}{\kern0pt}{\isadigit{2}}b{\isacharbrackleft}{\kern0pt}OF\ subst{\isacharbrackright}{\kern0pt}\ \isacommand{by}\isamarkupfalse%
\ auto\isanewline
\isanewline
\ \ \ \isacommand{from}\isamarkupfalse%
\ eqs\ fresh\isanewline
\ \ \ \isacommand{have}\isamarkupfalse%
\ {\isachardoublequoteopen}{\isacharparenleft}{\kern0pt}nabla{\isadigit{1}}{\isacharcomma}{\kern0pt}\ {\isacharparenleft}{\kern0pt}s{\isadigit{1}}\ {\isasymbullet}\ {\isacharbrackleft}{\kern0pt}{\isacharparenleft}{\kern0pt}X{\isacharcomma}{\kern0pt}\ swap\ {\isacharparenleft}{\kern0pt}rev\ pi{\isacharparenright}{\kern0pt}\ t{\isacharprime}{\kern0pt}{\isacharparenright}{\kern0pt}{\isacharbrackright}{\kern0pt}{\isacharparenright}{\kern0pt}{\isacharparenright}{\kern0pt}\ {\isasymin}\ U\ {\isacharparenleft}{\kern0pt}xs{\isacharcomma}{\kern0pt}\ ys{\isacharparenright}{\kern0pt}{\isachardoublequoteclose}\isanewline
\ \ \ \ \ \isacommand{using}\isamarkupfalse%
\ all{\isacharunderscore}{\kern0pt}solutions{\isacharunderscore}{\kern0pt}def\ \isacommand{by}\isamarkupfalse%
\ simp\isanewline
\ \ \ \isacommand{hence}\isamarkupfalse%
\ U{\isacharcolon}{\kern0pt}\ {\isachardoublequoteopen}{\isacharparenleft}{\kern0pt}nabla{\isadigit{1}}{\isacharcomma}{\kern0pt}\ s{\isadigit{1}}{\isacharparenright}{\kern0pt}\ {\isasymin}\ U\ {\isacharparenleft}{\kern0pt}apply{\isacharunderscore}{\kern0pt}subst\ {\isacharbrackleft}{\kern0pt}{\isacharparenleft}{\kern0pt}X{\isacharcomma}{\kern0pt}\ swap\ {\isacharparenleft}{\kern0pt}rev\ pi{\isacharparenright}{\kern0pt}\ t{\isacharprime}{\kern0pt}{\isacharparenright}{\kern0pt}{\isacharbrackright}{\kern0pt}\ {\isacharparenleft}{\kern0pt}xs{\isacharcomma}{\kern0pt}\ ys{\isacharparenright}{\kern0pt}{\isacharparenright}{\kern0pt}{\isachardoublequoteclose}\isanewline
\ \ \ \ \ \isacommand{using}\isamarkupfalse%
\ problem{\isacharunderscore}{\kern0pt}subst{\isacharunderscore}{\kern0pt}comm\ \isacommand{by}\isamarkupfalse%
\ simp\isanewline
\ \ \ \isanewline
\ \ \ \isacommand{from}\isamarkupfalse%
\ subst\ U\isanewline
\ \ \ \isacommand{show}\isamarkupfalse%
\ {\isacharquery}{\kern0pt}case\ \isacommand{by}\isamarkupfalse%
\ simp\isanewline
\ \isacommand{next}\isamarkupfalse%
\isanewline
\ \ \ \isacommand{case}\isamarkupfalse%
\ {\isacharparenleft}{\kern0pt}{\isadigit{9}}\ X\ t{\isacharprime}{\kern0pt}\ pi\ xs\ ys{\isacharparenright}{\kern0pt}\isanewline
\ \ \ \isacommand{hence}\isamarkupfalse%
\ {\isachardoublequoteopen}nabla{\isadigit{1}}\ {\isasymturnstile}\ \ subst\ s{\isadigit{1}}\ t{\isacharprime}{\kern0pt}\ {\isasymapprox}\ subst\ s{\isadigit{1}}\ {\isacharparenleft}{\kern0pt}Susp\ pi\ X{\isacharparenright}{\kern0pt}{\isachardoublequoteclose}\isanewline
\ \ \ \ \ \isacommand{using}\isamarkupfalse%
\ all{\isacharunderscore}{\kern0pt}solutions{\isacharunderscore}{\kern0pt}def{\isacharbrackleft}{\kern0pt}of\ {\isacartoucheopen}{\isacharparenleft}{\kern0pt}{\isacharparenleft}{\kern0pt}t{\isacharprime}{\kern0pt}{\isacharcomma}{\kern0pt}\ Susp\ pi\ X{\isacharparenright}{\kern0pt}\ {\isacharhash}{\kern0pt}\ xs{\isacharcomma}{\kern0pt}\ ys{\isacharparenright}{\kern0pt}{\isacartoucheclose}{\isacharbrackright}{\kern0pt}\ \isacommand{by}\isamarkupfalse%
\ simp\isanewline
\ \ \ \isacommand{hence}\isamarkupfalse%
\ subst{\isacharcolon}{\kern0pt}\ {\isachardoublequoteopen}nabla{\isadigit{1}}\ {\isasymTurnstile}\ subst\ {\isacharparenleft}{\kern0pt}s{\isadigit{1}}\ {\isasymbullet}\ {\isacharbrackleft}{\kern0pt}{\isacharparenleft}{\kern0pt}X{\isacharcomma}{\kern0pt}\ swap\ {\isacharparenleft}{\kern0pt}rev\ pi{\isacharparenright}{\kern0pt}\ t{\isacharprime}{\kern0pt}{\isacharparenright}{\kern0pt}{\isacharbrackright}{\kern0pt}{\isacharparenright}{\kern0pt}\ {\isasymapprox}\ subst\ s{\isadigit{1}}\ {\isachardoublequoteclose}\isanewline
\ \ \ \ \ \isacommand{using}\isamarkupfalse%
\ unif{\isacharunderscore}{\kern0pt}{\isadigit{1}}{\isacharbrackleft}{\kern0pt}OF\ equ{\isacharunderscore}{\kern0pt}symm{\isacharbrackright}{\kern0pt}\ \isacommand{by}\isamarkupfalse%
\ blast\isanewline
\isanewline
\ \ \ \isacommand{from}\isamarkupfalse%
\ {\isadigit{9}}\isanewline
\ \ \ \isacommand{have}\isamarkupfalse%
\ eqs{\isacharunderscore}{\kern0pt}old{\isacharcolon}{\kern0pt}\isanewline
\ \ \ \ {\isachardoublequoteopen}{\isasymforall}{\isacharparenleft}{\kern0pt}t{\isadigit{1}}{\isacharcomma}{\kern0pt}t{\isadigit{2}}{\isacharparenright}{\kern0pt}{\isasymin}set\ xs{\isachardot}{\kern0pt}\ nabla{\isadigit{1}}\ {\isasymturnstile}\ subst\ s{\isadigit{1}}\ t{\isadigit{1}}\ {\isasymapprox}\ subst\ s{\isadigit{1}}\ t{\isadigit{2}}{\isachardoublequoteclose}\isanewline
\ \ \ \ \ \isakeyword{and}\ fresh{\isacharunderscore}{\kern0pt}old{\isacharcolon}{\kern0pt}\isanewline
\ \ \ \ {\isachardoublequoteopen}{\isasymforall}{\isacharparenleft}{\kern0pt}a{\isacharcomma}{\kern0pt}t{\isacharparenright}{\kern0pt}{\isasymin}\ set\ ys{\isachardot}{\kern0pt}\ nabla{\isadigit{1}}\ {\isasymturnstile}\ a\ {\isasymsharp}\ subst\ s{\isadigit{1}}\ t{\isachardoublequoteclose}\isanewline
\ \ \ \ \ \ \isacommand{using}\isamarkupfalse%
\ all{\isacharunderscore}{\kern0pt}solutions{\isacharunderscore}{\kern0pt}def\ \isacommand{by}\isamarkupfalse%
\ auto\isanewline
\isanewline
\ \ \ \isacommand{from}\isamarkupfalse%
\ eqs{\isacharunderscore}{\kern0pt}old\ \isanewline
\ \ \ \isacommand{have}\isamarkupfalse%
\ eqs{\isacharcolon}{\kern0pt}\ {\isachardoublequoteopen}{\isasymforall}{\isacharparenleft}{\kern0pt}t{\isadigit{1}}{\isacharcomma}{\kern0pt}t{\isadigit{2}}{\isacharparenright}{\kern0pt}{\isasymin}set\ xs{\isachardot}{\kern0pt}\ \isanewline
\ \ \ \ nabla{\isadigit{1}}\ {\isasymturnstile}\ subst\ {\isacharparenleft}{\kern0pt}s{\isadigit{1}}\ {\isasymbullet}\ {\isacharbrackleft}{\kern0pt}{\isacharparenleft}{\kern0pt}X{\isacharcomma}{\kern0pt}\ swap\ {\isacharparenleft}{\kern0pt}rev\ pi{\isacharparenright}{\kern0pt}\ t{\isacharprime}{\kern0pt}{\isacharparenright}{\kern0pt}{\isacharbrackright}{\kern0pt}{\isacharparenright}{\kern0pt}\ t{\isadigit{1}}\ {\isasymapprox}\ \ subst\ {\isacharparenleft}{\kern0pt}s{\isadigit{1}}\ {\isasymbullet}\ {\isacharbrackleft}{\kern0pt}{\isacharparenleft}{\kern0pt}X{\isacharcomma}{\kern0pt}\ swap\ {\isacharparenleft}{\kern0pt}rev\ pi{\isacharparenright}{\kern0pt}\ t{\isacharprime}{\kern0pt}{\isacharparenright}{\kern0pt}{\isacharbrackright}{\kern0pt}{\isacharparenright}{\kern0pt}\ t{\isadigit{2}}{\isachardoublequoteclose}\isanewline
\ \ \ \ \isacommand{using}\isamarkupfalse%
\ unif{\isacharunderscore}{\kern0pt}{\isadigit{2}}a{\isacharbrackleft}{\kern0pt}OF\ subst{\isacharbrackright}{\kern0pt}\ \isacommand{by}\isamarkupfalse%
\ auto\isanewline
\isanewline
\ \ \ \isacommand{from}\isamarkupfalse%
\ fresh{\isacharunderscore}{\kern0pt}old\ \isanewline
\ \ \ \isacommand{have}\isamarkupfalse%
\ fresh{\isacharcolon}{\kern0pt}\ {\isachardoublequoteopen}{\isasymforall}{\isacharparenleft}{\kern0pt}a{\isacharcomma}{\kern0pt}t{\isacharparenright}{\kern0pt}{\isasymin}set\ ys{\isachardot}{\kern0pt}\ nabla{\isadigit{1}}\ {\isasymturnstile}\ a\ {\isasymsharp}\ subst\ {\isacharparenleft}{\kern0pt}s{\isadigit{1}}\ {\isasymbullet}\ {\isacharbrackleft}{\kern0pt}{\isacharparenleft}{\kern0pt}X{\isacharcomma}{\kern0pt}\ swap\ {\isacharparenleft}{\kern0pt}rev\ pi{\isacharparenright}{\kern0pt}\ t{\isacharprime}{\kern0pt}{\isacharparenright}{\kern0pt}{\isacharbrackright}{\kern0pt}{\isacharparenright}{\kern0pt}\ t{\isachardoublequoteclose}\isanewline
\ \ \ \ \ \isacommand{using}\isamarkupfalse%
\ unif{\isacharunderscore}{\kern0pt}{\isadigit{2}}b{\isacharbrackleft}{\kern0pt}OF\ subst{\isacharbrackright}{\kern0pt}\ \isacommand{by}\isamarkupfalse%
\ auto\isanewline
\isanewline
\ \ \ \isacommand{from}\isamarkupfalse%
\ eqs\ fresh\isanewline
\ \ \ \isacommand{have}\isamarkupfalse%
\ {\isachardoublequoteopen}{\isacharparenleft}{\kern0pt}nabla{\isadigit{1}}{\isacharcomma}{\kern0pt}\ {\isacharparenleft}{\kern0pt}s{\isadigit{1}}\ {\isasymbullet}\ {\isacharbrackleft}{\kern0pt}{\isacharparenleft}{\kern0pt}X{\isacharcomma}{\kern0pt}\ swap\ {\isacharparenleft}{\kern0pt}rev\ pi{\isacharparenright}{\kern0pt}\ t{\isacharprime}{\kern0pt}{\isacharparenright}{\kern0pt}{\isacharbrackright}{\kern0pt}{\isacharparenright}{\kern0pt}{\isacharparenright}{\kern0pt}\ {\isasymin}\ U\ {\isacharparenleft}{\kern0pt}xs{\isacharcomma}{\kern0pt}\ ys{\isacharparenright}{\kern0pt}{\isachardoublequoteclose}\isanewline
\ \ \ \ \ \isacommand{using}\isamarkupfalse%
\ all{\isacharunderscore}{\kern0pt}solutions{\isacharunderscore}{\kern0pt}def\ \isacommand{by}\isamarkupfalse%
\ simp\isanewline
\ \ \ \isacommand{hence}\isamarkupfalse%
\ U{\isacharcolon}{\kern0pt}\ {\isachardoublequoteopen}{\isacharparenleft}{\kern0pt}nabla{\isadigit{1}}{\isacharcomma}{\kern0pt}\ s{\isadigit{1}}{\isacharparenright}{\kern0pt}\ {\isasymin}\ U\ {\isacharparenleft}{\kern0pt}apply{\isacharunderscore}{\kern0pt}subst\ {\isacharbrackleft}{\kern0pt}{\isacharparenleft}{\kern0pt}X{\isacharcomma}{\kern0pt}\ swap\ {\isacharparenleft}{\kern0pt}rev\ pi{\isacharparenright}{\kern0pt}\ t{\isacharprime}{\kern0pt}{\isacharparenright}{\kern0pt}{\isacharbrackright}{\kern0pt}\ {\isacharparenleft}{\kern0pt}xs{\isacharcomma}{\kern0pt}\ ys{\isacharparenright}{\kern0pt}{\isacharparenright}{\kern0pt}{\isachardoublequoteclose}\isanewline
\ \ \ \ \ \isacommand{using}\isamarkupfalse%
\ problem{\isacharunderscore}{\kern0pt}subst{\isacharunderscore}{\kern0pt}comm\ \isacommand{by}\isamarkupfalse%
\ simp\isanewline
\isanewline
\ \ \ \isacommand{from}\isamarkupfalse%
\ subst\ U\isanewline
\ \ \isacommand{show}\isamarkupfalse%
\ {\isacharquery}{\kern0pt}case\ \isacommand{by}\isamarkupfalse%
\ simp\isanewline
\isacommand{qed}\isamarkupfalse%
%
\endisatagproof
{\isafoldproof}%
%
\isadelimproof
\isanewline
%
\endisadelimproof
\isanewline
\isanewline
\isanewline
\isacommand{lemma}\isamarkupfalse%
\ P{\isadigit{1}}{\isacharunderscore}{\kern0pt}from{\isacharunderscore}{\kern0pt}P{\isadigit{2}}{\isacharunderscore}{\kern0pt}sred{\isacharcolon}{\kern0pt}\ \isanewline
\ \ \isakeyword{assumes}\ {\isachardoublequoteopen}{\isacharparenleft}{\kern0pt}nabla{\isadigit{1}}{\isacharcomma}{\kern0pt}s{\isadigit{1}}{\isacharparenright}{\kern0pt}{\isasymin}\ U\ P{\isadigit{2}}{\isachardoublequoteclose}\ \isakeyword{and}\ {\isachardoublequoteopen}P{\isadigit{1}}{\isasymturnstile}s{\isadigit{2}}{\isasymleadsto}P{\isadigit{2}}{\isachardoublequoteclose}\isanewline
\ \ \isakeyword{shows}\ {\isachardoublequoteopen}{\isacharparenleft}{\kern0pt}nabla{\isadigit{1}}{\isacharcomma}{\kern0pt}s{\isadigit{1}}{\isasymbullet}s{\isadigit{2}}{\isacharparenright}{\kern0pt}{\isasymin}\ U\ P{\isadigit{1}}{\isachardoublequoteclose}\isanewline
%
\isadelimproof
\ %
\endisadelimproof
%
\isatagproof
\isacommand{using}\isamarkupfalse%
\ assms\isanewline
\isacommand{proof}\isamarkupfalse%
{\isacharparenleft}{\kern0pt}induction\ arbitrary{\isacharcolon}{\kern0pt}\ nabla{\isadigit{1}}\ s{\isadigit{1}}\ rule{\isacharcolon}{\kern0pt}\ s{\isacharunderscore}{\kern0pt}red{\isachardot}{\kern0pt}induct{\isacharbrackleft}{\kern0pt}OF\ {\isacartoucheopen}P{\isadigit{1}}\ {\isasymturnstile}\ s{\isadigit{2}}\ {\isasymleadsto}\ P{\isadigit{2}}{\isacartoucheclose}{\isacharbrackright}{\kern0pt}{\isacharparenright}{\kern0pt}\isanewline
\ \ \isacommand{case}\isamarkupfalse%
\ {\isacharparenleft}{\kern0pt}{\isadigit{1}}\ xs\ ys{\isacharparenright}{\kern0pt}\isanewline
\ \ \isacommand{hence}\isamarkupfalse%
\ {\isachardoublequoteopen}{\isasymforall}\ {\isacharparenleft}{\kern0pt}t{\isadigit{1}}{\isacharcomma}{\kern0pt}t{\isadigit{2}}{\isacharparenright}{\kern0pt}\ {\isasymin}\ set\ xs{\isachardot}{\kern0pt}\ nabla{\isadigit{1}}\ {\isasymturnstile}\ subst\ {\isacharparenleft}{\kern0pt}s{\isadigit{1}}\ {\isasymbullet}\ {\isacharbrackleft}{\kern0pt}{\isacharbrackright}{\kern0pt}{\isacharparenright}{\kern0pt}\ t{\isadigit{1}}\ {\isasymapprox}\ subst\ {\isacharparenleft}{\kern0pt}s{\isadigit{1}}\ {\isasymbullet}\ {\isacharbrackleft}{\kern0pt}{\isacharbrackright}{\kern0pt}{\isacharparenright}{\kern0pt}\ t{\isadigit{2}}{\isachardoublequoteclose}\ \isanewline
\ \ \ \ \ \isakeyword{and}\ fresh{\isacharcolon}{\kern0pt}\ {\isachardoublequoteopen}{\isasymforall}\ {\isacharparenleft}{\kern0pt}a{\isacharcomma}{\kern0pt}t{\isacharparenright}{\kern0pt}\ {\isasymin}\ set\ ys{\isachardot}{\kern0pt}\ nabla{\isadigit{1}}\ {\isasymturnstile}\ a\ {\isasymsharp}\ subst\ {\isacharparenleft}{\kern0pt}s{\isadigit{1}}\ {\isasymbullet}\ {\isacharbrackleft}{\kern0pt}{\isacharbrackright}{\kern0pt}{\isacharparenright}{\kern0pt}\ t{\isachardoublequoteclose}\isanewline
\ \ \ \ \isacommand{using}\isamarkupfalse%
\ all{\isacharunderscore}{\kern0pt}solutions{\isacharunderscore}{\kern0pt}def\ \isacommand{by}\isamarkupfalse%
\ simp{\isacharunderscore}{\kern0pt}all\isanewline
\ \ \isacommand{hence}\isamarkupfalse%
\ eqs{\isacharcolon}{\kern0pt}\ \isanewline
\ \ \ \ {\isachardoublequoteopen}{\isasymforall}\ {\isacharparenleft}{\kern0pt}t{\isadigit{1}}{\isacharcomma}{\kern0pt}t{\isadigit{2}}{\isacharparenright}{\kern0pt}\ {\isasymin}\ set\ {\isacharparenleft}{\kern0pt}{\isacharparenleft}{\kern0pt}Unit{\isacharcomma}{\kern0pt}\ Unit{\isacharparenright}{\kern0pt}\ {\isacharhash}{\kern0pt}\ xs{\isacharparenright}{\kern0pt}{\isachardot}{\kern0pt}\ nabla{\isadigit{1}}\ {\isasymturnstile}\ subst\ {\isacharparenleft}{\kern0pt}s{\isadigit{1}}\ {\isasymbullet}\ {\isacharbrackleft}{\kern0pt}{\isacharbrackright}{\kern0pt}{\isacharparenright}{\kern0pt}\ t{\isadigit{1}}\ {\isasymapprox}\ subst\ {\isacharparenleft}{\kern0pt}s{\isadigit{1}}\ {\isasymbullet}\ {\isacharbrackleft}{\kern0pt}{\isacharbrackright}{\kern0pt}{\isacharparenright}{\kern0pt}\ t{\isadigit{2}}{\isachardoublequoteclose}\isanewline
\ \ \ \ \isacommand{using}\isamarkupfalse%
\ subst{\isacharunderscore}{\kern0pt}unit{\isacharbrackleft}{\kern0pt}of\ {\isacartoucheopen}{\isacharparenleft}{\kern0pt}s{\isadigit{1}}\ {\isasymbullet}\ {\isacharbrackleft}{\kern0pt}{\isacharbrackright}{\kern0pt}{\isacharparenright}{\kern0pt}{\isacartoucheclose}{\isacharbrackright}{\kern0pt}\ equ{\isacharunderscore}{\kern0pt}refl{\isacharbrackleft}{\kern0pt}of\ nabla{\isadigit{1}}\ Unit{\isacharbrackright}{\kern0pt}\ \isacommand{by}\isamarkupfalse%
\ auto\isanewline
\ \ \isacommand{from}\isamarkupfalse%
\ fresh\ eqs\ \isacommand{show}\isamarkupfalse%
\ {\isacharquery}{\kern0pt}case\isanewline
\ \ \ \ \isacommand{using}\isamarkupfalse%
\ all{\isacharunderscore}{\kern0pt}solutions{\isacharunderscore}{\kern0pt}def\ \isacommand{by}\isamarkupfalse%
\ simp\isanewline
\isacommand{next}\isamarkupfalse%
\isanewline
\ \ \isacommand{case}\isamarkupfalse%
\ {\isacharparenleft}{\kern0pt}{\isadigit{2}}\ t{\isadigit{1}}\ t{\isadigit{2}}\ t{\isadigit{3}}\ t{\isadigit{4}}\ xs\ ys{\isacharparenright}{\kern0pt}\isanewline
\ \ \ \isacommand{hence}\isamarkupfalse%
\ eqs{\isacharunderscore}{\kern0pt}old{\isacharcolon}{\kern0pt}\ {\isachardoublequoteopen}{\isasymforall}\ {\isacharparenleft}{\kern0pt}t{\isadigit{1}}{\isacharcomma}{\kern0pt}t{\isadigit{2}}{\isacharparenright}{\kern0pt}\ {\isasymin}\ set\ {\isacharparenleft}{\kern0pt}{\isacharparenleft}{\kern0pt}t{\isadigit{1}}{\isacharcomma}{\kern0pt}\ t{\isadigit{3}}{\isacharparenright}{\kern0pt}\ {\isacharhash}{\kern0pt}\ {\isacharparenleft}{\kern0pt}t{\isadigit{2}}{\isacharcomma}{\kern0pt}\ t{\isadigit{4}}{\isacharparenright}{\kern0pt}\ {\isacharhash}{\kern0pt}\ xs{\isacharparenright}{\kern0pt}{\isachardot}{\kern0pt}\ nabla{\isadigit{1}}\ {\isasymturnstile}\ subst\ {\isacharparenleft}{\kern0pt}s{\isadigit{1}}\ {\isasymbullet}\ {\isacharbrackleft}{\kern0pt}{\isacharbrackright}{\kern0pt}{\isacharparenright}{\kern0pt}\ t{\isadigit{1}}\ {\isasymapprox}\ subst\ {\isacharparenleft}{\kern0pt}s{\isadigit{1}}\ {\isasymbullet}\ {\isacharbrackleft}{\kern0pt}{\isacharbrackright}{\kern0pt}{\isacharparenright}{\kern0pt}\ t{\isadigit{2}}{\isachardoublequoteclose}\ \isanewline
\ \ \ \ \ \isakeyword{and}\ fresh{\isacharcolon}{\kern0pt}\ {\isachardoublequoteopen}{\isasymforall}\ {\isacharparenleft}{\kern0pt}a{\isacharcomma}{\kern0pt}t{\isacharparenright}{\kern0pt}\ {\isasymin}\ set\ ys{\isachardot}{\kern0pt}\ nabla{\isadigit{1}}\ {\isasymturnstile}\ a\ {\isasymsharp}\ subst\ {\isacharparenleft}{\kern0pt}s{\isadigit{1}}\ {\isasymbullet}\ {\isacharbrackleft}{\kern0pt}{\isacharbrackright}{\kern0pt}{\isacharparenright}{\kern0pt}\ t{\isachardoublequoteclose}\isanewline
\ \ \ \ \ \isacommand{using}\isamarkupfalse%
\ all{\isacharunderscore}{\kern0pt}solutions{\isacharunderscore}{\kern0pt}def\ \isacommand{by}\isamarkupfalse%
\ simp{\isacharunderscore}{\kern0pt}all\isanewline
\ \ \ \isacommand{hence}\isamarkupfalse%
\ {\isachardoublequoteopen}nabla{\isadigit{1}}\ {\isasymturnstile}\ subst\ {\isacharparenleft}{\kern0pt}s{\isadigit{1}}\ {\isasymbullet}\ {\isacharbrackleft}{\kern0pt}{\isacharbrackright}{\kern0pt}{\isacharparenright}{\kern0pt}\ t{\isadigit{1}}\ {\isasymapprox}\ subst\ {\isacharparenleft}{\kern0pt}s{\isadigit{1}}\ {\isasymbullet}\ {\isacharbrackleft}{\kern0pt}{\isacharbrackright}{\kern0pt}{\isacharparenright}{\kern0pt}\ t{\isadigit{3}}{\isachardoublequoteclose}\ \isanewline
\ \ \ \ \ {\isachardoublequoteopen}nabla{\isadigit{1}}\ {\isasymturnstile}\ subst\ {\isacharparenleft}{\kern0pt}s{\isadigit{1}}\ {\isasymbullet}\ {\isacharbrackleft}{\kern0pt}{\isacharbrackright}{\kern0pt}{\isacharparenright}{\kern0pt}\ t{\isadigit{2}}\ {\isasymapprox}\ subst\ {\isacharparenleft}{\kern0pt}s{\isadigit{1}}\ {\isasymbullet}\ {\isacharbrackleft}{\kern0pt}{\isacharbrackright}{\kern0pt}{\isacharparenright}{\kern0pt}\ t{\isadigit{4}}{\isachardoublequoteclose}\ \isacommand{by}\isamarkupfalse%
\ simp{\isacharunderscore}{\kern0pt}all\isanewline
\ \ \ \isacommand{hence}\isamarkupfalse%
\ {\isachardoublequoteopen}nabla{\isadigit{1}}\ {\isasymturnstile}\ subst\ {\isacharparenleft}{\kern0pt}s{\isadigit{1}}\ {\isasymbullet}\ {\isacharbrackleft}{\kern0pt}{\isacharbrackright}{\kern0pt}{\isacharparenright}{\kern0pt}\ {\isacharparenleft}{\kern0pt}Paar\ t{\isadigit{1}}\ t{\isadigit{2}}{\isacharparenright}{\kern0pt}\ {\isasymapprox}\ subst\ {\isacharparenleft}{\kern0pt}s{\isadigit{1}}\ {\isasymbullet}\ {\isacharbrackleft}{\kern0pt}{\isacharbrackright}{\kern0pt}{\isacharparenright}{\kern0pt}\ {\isacharparenleft}{\kern0pt}Paar\ t{\isadigit{3}}\ t{\isadigit{4}}{\isacharparenright}{\kern0pt}{\isachardoublequoteclose}\isanewline
\ \ \ \ \ \isacommand{using}\isamarkupfalse%
\ equ{\isacharunderscore}{\kern0pt}paar\ \isacommand{by}\isamarkupfalse%
\ simp\isanewline
\ \ \ \isacommand{with}\isamarkupfalse%
\ eqs{\isacharunderscore}{\kern0pt}old\ \isacommand{have}\isamarkupfalse%
\ eqs{\isacharcolon}{\kern0pt}\isanewline
\ \ \ \ \ {\isachardoublequoteopen}{\isasymforall}\ {\isacharparenleft}{\kern0pt}t{\isadigit{1}}{\isacharcomma}{\kern0pt}t{\isadigit{2}}{\isacharparenright}{\kern0pt}\ {\isasymin}\ set\ {\isacharparenleft}{\kern0pt}{\isacharparenleft}{\kern0pt}Paar\ t{\isadigit{1}}\ t{\isadigit{2}}{\isacharcomma}{\kern0pt}\ Paar\ t{\isadigit{3}}\ t{\isadigit{4}}{\isacharparenright}{\kern0pt}\ {\isacharhash}{\kern0pt}\ xs{\isacharparenright}{\kern0pt}{\isachardot}{\kern0pt}\ nabla{\isadigit{1}}\ {\isasymturnstile}\ subst\ {\isacharparenleft}{\kern0pt}s{\isadigit{1}}\ {\isasymbullet}\ {\isacharbrackleft}{\kern0pt}{\isacharbrackright}{\kern0pt}{\isacharparenright}{\kern0pt}\ t{\isadigit{1}}\ {\isasymapprox}\ subst\ {\isacharparenleft}{\kern0pt}s{\isadigit{1}}\ {\isasymbullet}\ {\isacharbrackleft}{\kern0pt}{\isacharbrackright}{\kern0pt}{\isacharparenright}{\kern0pt}\ t{\isadigit{2}}{\isachardoublequoteclose}\isanewline
\ \ \ \ \ \isacommand{by}\isamarkupfalse%
\ auto\isanewline
\ \ \ \isacommand{from}\isamarkupfalse%
\ fresh\ eqs\ \isacommand{show}\isamarkupfalse%
\ {\isacharquery}{\kern0pt}case\isanewline
\ \ \ \ \isacommand{using}\isamarkupfalse%
\ all{\isacharunderscore}{\kern0pt}solutions{\isacharunderscore}{\kern0pt}def\ \isacommand{by}\isamarkupfalse%
\ simp\isanewline
\isacommand{next}\isamarkupfalse%
\isanewline
\ \ \isacommand{case}\isamarkupfalse%
\ {\isacharparenleft}{\kern0pt}{\isadigit{3}}\ f\ t{\isadigit{1}}\ t{\isadigit{2}}\ xs\ ys{\isacharparenright}{\kern0pt}\isanewline
\ \ \isacommand{hence}\isamarkupfalse%
\ eqs{\isacharunderscore}{\kern0pt}old{\isacharcolon}{\kern0pt}\ {\isachardoublequoteopen}{\isasymforall}\ {\isacharparenleft}{\kern0pt}t{\isadigit{1}}{\isacharcomma}{\kern0pt}t{\isadigit{2}}{\isacharparenright}{\kern0pt}\ {\isasymin}\ set\ {\isacharparenleft}{\kern0pt}{\isacharparenleft}{\kern0pt}t{\isadigit{1}}{\isacharcomma}{\kern0pt}\ t{\isadigit{2}}{\isacharparenright}{\kern0pt}\ {\isacharhash}{\kern0pt}\ xs{\isacharparenright}{\kern0pt}{\isachardot}{\kern0pt}\ nabla{\isadigit{1}}\ {\isasymturnstile}\ subst\ {\isacharparenleft}{\kern0pt}s{\isadigit{1}}\ {\isasymbullet}\ {\isacharbrackleft}{\kern0pt}{\isacharbrackright}{\kern0pt}{\isacharparenright}{\kern0pt}\ t{\isadigit{1}}\ {\isasymapprox}\ subst\ {\isacharparenleft}{\kern0pt}s{\isadigit{1}}\ {\isasymbullet}\ {\isacharbrackleft}{\kern0pt}{\isacharbrackright}{\kern0pt}{\isacharparenright}{\kern0pt}\ t{\isadigit{2}}{\isachardoublequoteclose}\ \isanewline
\ \ \ \ \ \isakeyword{and}\ fresh{\isacharcolon}{\kern0pt}\ {\isachardoublequoteopen}{\isasymforall}\ {\isacharparenleft}{\kern0pt}a{\isacharcomma}{\kern0pt}t{\isacharparenright}{\kern0pt}\ {\isasymin}\ set\ ys{\isachardot}{\kern0pt}\ nabla{\isadigit{1}}\ {\isasymturnstile}\ a\ {\isasymsharp}\ subst\ {\isacharparenleft}{\kern0pt}s{\isadigit{1}}\ {\isasymbullet}\ {\isacharbrackleft}{\kern0pt}{\isacharbrackright}{\kern0pt}{\isacharparenright}{\kern0pt}\ t{\isachardoublequoteclose}\isanewline
\ \ \ \ \isacommand{using}\isamarkupfalse%
\ all{\isacharunderscore}{\kern0pt}solutions{\isacharunderscore}{\kern0pt}def\ \isacommand{by}\isamarkupfalse%
\ simp{\isacharunderscore}{\kern0pt}all\isanewline
\ \ \isacommand{hence}\isamarkupfalse%
\ {\isachardoublequoteopen}nabla{\isadigit{1}}\ {\isasymturnstile}\ subst\ {\isacharparenleft}{\kern0pt}s{\isadigit{1}}\ {\isasymbullet}\ {\isacharbrackleft}{\kern0pt}{\isacharbrackright}{\kern0pt}{\isacharparenright}{\kern0pt}\ {\isacharparenleft}{\kern0pt}Func\ f\ t{\isadigit{1}}{\isacharparenright}{\kern0pt}\ {\isasymapprox}\ subst\ {\isacharparenleft}{\kern0pt}s{\isadigit{1}}\ {\isasymbullet}\ {\isacharbrackleft}{\kern0pt}{\isacharbrackright}{\kern0pt}{\isacharparenright}{\kern0pt}\ {\isacharparenleft}{\kern0pt}Func\ f\ t{\isadigit{2}}{\isacharparenright}{\kern0pt}{\isachardoublequoteclose}\isanewline
\ \ \ \ \isacommand{using}\isamarkupfalse%
\ equ{\isacharunderscore}{\kern0pt}func\ \isacommand{by}\isamarkupfalse%
\ simp\isanewline
\ \ \isacommand{with}\isamarkupfalse%
\ eqs{\isacharunderscore}{\kern0pt}old\ \isacommand{have}\isamarkupfalse%
\ eqs{\isacharcolon}{\kern0pt}\isanewline
\ \ \ \ \ {\isachardoublequoteopen}{\isasymforall}\ {\isacharparenleft}{\kern0pt}t{\isadigit{1}}{\isacharcomma}{\kern0pt}t{\isadigit{2}}{\isacharparenright}{\kern0pt}\ {\isasymin}\ set\ {\isacharparenleft}{\kern0pt}{\isacharparenleft}{\kern0pt}Func\ f\ t{\isadigit{1}}{\isacharcomma}{\kern0pt}\ Func\ f\ t{\isadigit{2}}{\isacharparenright}{\kern0pt}\ {\isacharhash}{\kern0pt}\ xs{\isacharparenright}{\kern0pt}{\isachardot}{\kern0pt}\ nabla{\isadigit{1}}\ {\isasymturnstile}\ subst\ {\isacharparenleft}{\kern0pt}s{\isadigit{1}}\ {\isasymbullet}\ {\isacharbrackleft}{\kern0pt}{\isacharbrackright}{\kern0pt}{\isacharparenright}{\kern0pt}\ t{\isadigit{1}}\ {\isasymapprox}\ subst\ {\isacharparenleft}{\kern0pt}s{\isadigit{1}}\ {\isasymbullet}\ {\isacharbrackleft}{\kern0pt}{\isacharbrackright}{\kern0pt}{\isacharparenright}{\kern0pt}\ t{\isadigit{2}}{\isachardoublequoteclose}\isanewline
\ \ \ \ \ \isacommand{by}\isamarkupfalse%
\ auto\isanewline
\ \ \isacommand{from}\isamarkupfalse%
\ fresh\ eqs\ \isacommand{show}\isamarkupfalse%
\ {\isacharquery}{\kern0pt}case\isanewline
\ \ \ \ \isacommand{using}\isamarkupfalse%
\ all{\isacharunderscore}{\kern0pt}solutions{\isacharunderscore}{\kern0pt}def\ \isacommand{by}\isamarkupfalse%
\ simp\isanewline
\isacommand{next}\isamarkupfalse%
\isanewline
\ \ \isacommand{case}\isamarkupfalse%
\ {\isacharparenleft}{\kern0pt}{\isadigit{4}}\ a\ t{\isadigit{1}}\ t{\isadigit{2}}\ xs\ ys{\isacharparenright}{\kern0pt}\isanewline
\ \ \isacommand{hence}\isamarkupfalse%
\ eqs{\isacharunderscore}{\kern0pt}old{\isacharcolon}{\kern0pt}\ {\isachardoublequoteopen}{\isasymforall}\ {\isacharparenleft}{\kern0pt}t{\isadigit{1}}{\isacharcomma}{\kern0pt}t{\isadigit{2}}{\isacharparenright}{\kern0pt}\ {\isasymin}\ set\ {\isacharparenleft}{\kern0pt}{\isacharparenleft}{\kern0pt}t{\isadigit{1}}{\isacharcomma}{\kern0pt}\ t{\isadigit{2}}{\isacharparenright}{\kern0pt}\ {\isacharhash}{\kern0pt}\ xs{\isacharparenright}{\kern0pt}{\isachardot}{\kern0pt}\ nabla{\isadigit{1}}\ {\isasymturnstile}\ subst\ {\isacharparenleft}{\kern0pt}s{\isadigit{1}}\ {\isasymbullet}\ {\isacharbrackleft}{\kern0pt}{\isacharbrackright}{\kern0pt}{\isacharparenright}{\kern0pt}\ t{\isadigit{1}}\ {\isasymapprox}\ subst\ {\isacharparenleft}{\kern0pt}s{\isadigit{1}}\ {\isasymbullet}\ {\isacharbrackleft}{\kern0pt}{\isacharbrackright}{\kern0pt}{\isacharparenright}{\kern0pt}\ t{\isadigit{2}}{\isachardoublequoteclose}\ \isanewline
\ \ \ \ \ \isakeyword{and}\ fresh{\isacharcolon}{\kern0pt}\ {\isachardoublequoteopen}{\isasymforall}\ {\isacharparenleft}{\kern0pt}a{\isacharcomma}{\kern0pt}t{\isacharparenright}{\kern0pt}\ {\isasymin}\ set\ ys{\isachardot}{\kern0pt}\ nabla{\isadigit{1}}\ {\isasymturnstile}\ a\ {\isasymsharp}\ subst\ {\isacharparenleft}{\kern0pt}s{\isadigit{1}}\ {\isasymbullet}\ {\isacharbrackleft}{\kern0pt}{\isacharbrackright}{\kern0pt}{\isacharparenright}{\kern0pt}\ t{\isachardoublequoteclose}\isanewline
\ \ \ \ \isacommand{using}\isamarkupfalse%
\ all{\isacharunderscore}{\kern0pt}solutions{\isacharunderscore}{\kern0pt}def\ \isacommand{by}\isamarkupfalse%
\ simp{\isacharunderscore}{\kern0pt}all\isanewline
\ \ \isacommand{hence}\isamarkupfalse%
\ {\isachardoublequoteopen}nabla{\isadigit{1}}\ {\isasymturnstile}\ subst\ {\isacharparenleft}{\kern0pt}s{\isadigit{1}}\ {\isasymbullet}\ {\isacharbrackleft}{\kern0pt}{\isacharbrackright}{\kern0pt}{\isacharparenright}{\kern0pt}\ {\isacharparenleft}{\kern0pt}Abst\ a\ t{\isadigit{1}}{\isacharparenright}{\kern0pt}\ {\isasymapprox}\ subst\ {\isacharparenleft}{\kern0pt}s{\isadigit{1}}\ {\isasymbullet}\ {\isacharbrackleft}{\kern0pt}{\isacharbrackright}{\kern0pt}{\isacharparenright}{\kern0pt}\ {\isacharparenleft}{\kern0pt}Abst\ a\ t{\isadigit{2}}{\isacharparenright}{\kern0pt}{\isachardoublequoteclose}\isanewline
\ \ \ \ \isacommand{using}\isamarkupfalse%
\ equ{\isacharunderscore}{\kern0pt}abst{\isacharunderscore}{\kern0pt}aa\ \isacommand{by}\isamarkupfalse%
\ simp\isanewline
\ \ \isacommand{with}\isamarkupfalse%
\ eqs{\isacharunderscore}{\kern0pt}old\ \isacommand{have}\isamarkupfalse%
\ eqs{\isacharcolon}{\kern0pt}\isanewline
\ \ \ \ \ {\isachardoublequoteopen}{\isasymforall}\ {\isacharparenleft}{\kern0pt}t{\isadigit{1}}{\isacharcomma}{\kern0pt}t{\isadigit{2}}{\isacharparenright}{\kern0pt}\ {\isasymin}\ set\ {\isacharparenleft}{\kern0pt}{\isacharparenleft}{\kern0pt}Abst\ a\ t{\isadigit{1}}{\isacharcomma}{\kern0pt}\ Abst\ a\ t{\isadigit{2}}{\isacharparenright}{\kern0pt}\ {\isacharhash}{\kern0pt}\ xs{\isacharparenright}{\kern0pt}{\isachardot}{\kern0pt}\ nabla{\isadigit{1}}\ {\isasymturnstile}\ subst\ {\isacharparenleft}{\kern0pt}s{\isadigit{1}}\ {\isasymbullet}\ {\isacharbrackleft}{\kern0pt}{\isacharbrackright}{\kern0pt}{\isacharparenright}{\kern0pt}\ t{\isadigit{1}}\ {\isasymapprox}\ subst\ {\isacharparenleft}{\kern0pt}s{\isadigit{1}}\ {\isasymbullet}\ {\isacharbrackleft}{\kern0pt}{\isacharbrackright}{\kern0pt}{\isacharparenright}{\kern0pt}\ t{\isadigit{2}}{\isachardoublequoteclose}\isanewline
\ \ \ \ \ \isacommand{by}\isamarkupfalse%
\ auto\isanewline
\ \ \isacommand{from}\isamarkupfalse%
\ fresh\ eqs\ \isacommand{show}\isamarkupfalse%
\ {\isacharquery}{\kern0pt}case\isanewline
\ \ \ \ \isacommand{using}\isamarkupfalse%
\ all{\isacharunderscore}{\kern0pt}solutions{\isacharunderscore}{\kern0pt}def\ \isacommand{by}\isamarkupfalse%
\ simp\isanewline
\isacommand{next}\isamarkupfalse%
\isanewline
\ \ \isacommand{case}\isamarkupfalse%
\ {\isacharparenleft}{\kern0pt}{\isadigit{5}}\ a\ b\ t{\isadigit{1}}\ t{\isadigit{2}}\ xs\ ys{\isacharparenright}{\kern0pt}\isanewline
\ \ \isacommand{hence}\isamarkupfalse%
\ eqs{\isacharunderscore}{\kern0pt}old{\isacharcolon}{\kern0pt}\ {\isachardoublequoteopen}{\isasymforall}\ {\isacharparenleft}{\kern0pt}t{\isadigit{1}}{\isacharcomma}{\kern0pt}t{\isadigit{2}}{\isacharparenright}{\kern0pt}\ {\isasymin}\ set\ {\isacharparenleft}{\kern0pt}{\isacharparenleft}{\kern0pt}t{\isadigit{1}}{\isacharcomma}{\kern0pt}\ swap\ {\isacharbrackleft}{\kern0pt}{\isacharparenleft}{\kern0pt}a{\isacharcomma}{\kern0pt}\ b{\isacharparenright}{\kern0pt}{\isacharbrackright}{\kern0pt}\ t{\isadigit{2}}{\isacharparenright}{\kern0pt}\ {\isacharhash}{\kern0pt}\ xs{\isacharparenright}{\kern0pt}{\isachardot}{\kern0pt}\ nabla{\isadigit{1}}\ {\isasymturnstile}\ subst\ {\isacharparenleft}{\kern0pt}s{\isadigit{1}}\ {\isasymbullet}\ {\isacharbrackleft}{\kern0pt}{\isacharbrackright}{\kern0pt}{\isacharparenright}{\kern0pt}\ t{\isadigit{1}}\ {\isasymapprox}\ subst\ {\isacharparenleft}{\kern0pt}s{\isadigit{1}}\ {\isasymbullet}\ {\isacharbrackleft}{\kern0pt}{\isacharbrackright}{\kern0pt}{\isacharparenright}{\kern0pt}\ t{\isadigit{2}}{\isachardoublequoteclose}\ \isanewline
\ \ \ \ \isakeyword{and}\ fresh{\isacharcolon}{\kern0pt}\ {\isachardoublequoteopen}{\isasymforall}\ {\isacharparenleft}{\kern0pt}a{\isacharcomma}{\kern0pt}t{\isacharparenright}{\kern0pt}\ {\isasymin}\ set\ {\isacharparenleft}{\kern0pt}{\isacharparenleft}{\kern0pt}a{\isacharcomma}{\kern0pt}\ t{\isadigit{2}}{\isacharparenright}{\kern0pt}\ {\isacharhash}{\kern0pt}\ ys{\isacharparenright}{\kern0pt}{\isachardot}{\kern0pt}\ nabla{\isadigit{1}}\ {\isasymturnstile}\ a\ {\isasymsharp}\ subst\ {\isacharparenleft}{\kern0pt}s{\isadigit{1}}\ {\isasymbullet}\ {\isacharbrackleft}{\kern0pt}{\isacharbrackright}{\kern0pt}{\isacharparenright}{\kern0pt}\ t{\isachardoublequoteclose}\isanewline
\ \ \ \ \isacommand{using}\isamarkupfalse%
\ all{\isacharunderscore}{\kern0pt}solutions{\isacharunderscore}{\kern0pt}def\ \isacommand{by}\isamarkupfalse%
\ simp{\isacharunderscore}{\kern0pt}all\isanewline
\ \ \isacommand{hence}\isamarkupfalse%
\ {\isachardoublequoteopen}nabla{\isadigit{1}}\ {\isasymturnstile}\ subst\ {\isacharparenleft}{\kern0pt}s{\isadigit{1}}\ {\isasymbullet}\ {\isacharbrackleft}{\kern0pt}{\isacharbrackright}{\kern0pt}{\isacharparenright}{\kern0pt}\ t{\isadigit{1}}\ {\isasymapprox}\ swap\ {\isacharbrackleft}{\kern0pt}{\isacharparenleft}{\kern0pt}a{\isacharcomma}{\kern0pt}\ b{\isacharparenright}{\kern0pt}{\isacharbrackright}{\kern0pt}\ {\isacharparenleft}{\kern0pt}subst\ {\isacharparenleft}{\kern0pt}s{\isadigit{1}}\ {\isasymbullet}\ {\isacharbrackleft}{\kern0pt}{\isacharbrackright}{\kern0pt}{\isacharparenright}{\kern0pt}\ t{\isadigit{2}}{\isacharparenright}{\kern0pt}{\isachardoublequoteclose}\isanewline
\ \ \ \ \ \ {\isachardoublequoteopen}nabla{\isadigit{1}}\ {\isasymturnstile}\ a\ {\isasymsharp}\ subst\ {\isacharparenleft}{\kern0pt}s{\isadigit{1}}\ {\isasymbullet}\ {\isacharbrackleft}{\kern0pt}{\isacharbrackright}{\kern0pt}{\isacharparenright}{\kern0pt}\ t{\isadigit{2}}{\isachardoublequoteclose}\ \isanewline
\ \ \ \ \isacommand{using}\isamarkupfalse%
\ subst{\isacharunderscore}{\kern0pt}swap{\isacharunderscore}{\kern0pt}comm{\isacharbrackleft}{\kern0pt}of\ {\isacartoucheopen}{\isacharparenleft}{\kern0pt}s{\isadigit{1}}\ {\isasymbullet}\ {\isacharbrackleft}{\kern0pt}{\isacharbrackright}{\kern0pt}{\isacharparenright}{\kern0pt}{\isacartoucheclose}\ {\isacartoucheopen}{\isacharbrackleft}{\kern0pt}{\isacharparenleft}{\kern0pt}a{\isacharcomma}{\kern0pt}\ b{\isacharparenright}{\kern0pt}{\isacharbrackright}{\kern0pt}{\isacartoucheclose}\ t{\isadigit{2}}{\isacharbrackright}{\kern0pt}\ \isacommand{by}\isamarkupfalse%
\ simp{\isacharunderscore}{\kern0pt}all\isanewline
\ \ \isacommand{hence}\isamarkupfalse%
\ {\isachardoublequoteopen}nabla{\isadigit{1}}\ {\isasymturnstile}\ subst\ {\isacharparenleft}{\kern0pt}s{\isadigit{1}}\ {\isasymbullet}\ {\isacharbrackleft}{\kern0pt}{\isacharbrackright}{\kern0pt}{\isacharparenright}{\kern0pt}\ {\isacharparenleft}{\kern0pt}Abst\ a\ t{\isadigit{1}}{\isacharparenright}{\kern0pt}\ {\isasymapprox}\ subst\ {\isacharparenleft}{\kern0pt}s{\isadigit{1}}\ {\isasymbullet}\ {\isacharbrackleft}{\kern0pt}{\isacharbrackright}{\kern0pt}{\isacharparenright}{\kern0pt}\ {\isacharparenleft}{\kern0pt}Abst\ b\ t{\isadigit{2}}{\isacharparenright}{\kern0pt}{\isachardoublequoteclose}\isanewline
\ \ \ \ \isacommand{using}\isamarkupfalse%
\ equ{\isacharunderscore}{\kern0pt}abst{\isacharunderscore}{\kern0pt}ab{\isacharbrackleft}{\kern0pt}OF\ {\isadigit{5}}{\isacharparenleft}{\kern0pt}{\isadigit{1}}{\isacharparenright}{\kern0pt}\ {\isacartoucheopen}nabla{\isadigit{1}}\ {\isasymturnstile}\ a\ {\isasymsharp}\ subst\ {\isacharparenleft}{\kern0pt}s{\isadigit{1}}\ {\isasymbullet}\ {\isacharbrackleft}{\kern0pt}{\isacharbrackright}{\kern0pt}{\isacharparenright}{\kern0pt}\ t{\isadigit{2}}{\isacartoucheclose}{\isacharbrackright}{\kern0pt}\ subst{\isacharunderscore}{\kern0pt}abst\ \isacommand{by}\isamarkupfalse%
\ simp\isanewline
\ \ \isacommand{with}\isamarkupfalse%
\ eqs{\isacharunderscore}{\kern0pt}old\ \isacommand{have}\isamarkupfalse%
\ eqs{\isacharcolon}{\kern0pt}\isanewline
\ \ \ \ \ {\isachardoublequoteopen}{\isasymforall}\ {\isacharparenleft}{\kern0pt}t{\isadigit{1}}{\isacharcomma}{\kern0pt}t{\isadigit{2}}{\isacharparenright}{\kern0pt}\ {\isasymin}\ set\ {\isacharparenleft}{\kern0pt}{\isacharparenleft}{\kern0pt}Abst\ a\ t{\isadigit{1}}{\isacharcomma}{\kern0pt}\ Abst\ b\ t{\isadigit{2}}{\isacharparenright}{\kern0pt}\ {\isacharhash}{\kern0pt}\ xs{\isacharparenright}{\kern0pt}{\isachardot}{\kern0pt}\ nabla{\isadigit{1}}\ {\isasymturnstile}\ subst\ {\isacharparenleft}{\kern0pt}s{\isadigit{1}}\ {\isasymbullet}\ {\isacharbrackleft}{\kern0pt}{\isacharbrackright}{\kern0pt}{\isacharparenright}{\kern0pt}\ t{\isadigit{1}}\ {\isasymapprox}\ subst\ {\isacharparenleft}{\kern0pt}s{\isadigit{1}}\ {\isasymbullet}\ {\isacharbrackleft}{\kern0pt}{\isacharbrackright}{\kern0pt}{\isacharparenright}{\kern0pt}\ t{\isadigit{2}}{\isachardoublequoteclose}\isanewline
\ \ \ \ \ \isacommand{by}\isamarkupfalse%
\ auto\isanewline
\ \ \isacommand{from}\isamarkupfalse%
\ fresh\ eqs\ \isacommand{show}\isamarkupfalse%
\ {\isacharquery}{\kern0pt}case\isanewline
\ \ \ \ \isacommand{using}\isamarkupfalse%
\ all{\isacharunderscore}{\kern0pt}solutions{\isacharunderscore}{\kern0pt}def\ \isacommand{by}\isamarkupfalse%
\ simp\isanewline
\isacommand{next}\isamarkupfalse%
\isanewline
\ \ \isacommand{case}\isamarkupfalse%
\ {\isacharparenleft}{\kern0pt}{\isadigit{6}}\ a\ xs\ ys{\isacharparenright}{\kern0pt}\isanewline
\ \ \isacommand{hence}\isamarkupfalse%
\ {\isachardoublequoteopen}{\isasymforall}\ {\isacharparenleft}{\kern0pt}t{\isadigit{1}}{\isacharcomma}{\kern0pt}t{\isadigit{2}}{\isacharparenright}{\kern0pt}\ {\isasymin}\ set\ xs{\isachardot}{\kern0pt}\ nabla{\isadigit{1}}\ {\isasymturnstile}\ subst\ {\isacharparenleft}{\kern0pt}s{\isadigit{1}}\ {\isasymbullet}\ {\isacharbrackleft}{\kern0pt}{\isacharbrackright}{\kern0pt}{\isacharparenright}{\kern0pt}\ t{\isadigit{1}}\ {\isasymapprox}\ subst\ {\isacharparenleft}{\kern0pt}s{\isadigit{1}}\ {\isasymbullet}\ {\isacharbrackleft}{\kern0pt}{\isacharbrackright}{\kern0pt}{\isacharparenright}{\kern0pt}\ t{\isadigit{2}}{\isachardoublequoteclose}\ \isanewline
\ \ \ \ \ \isakeyword{and}\ fresh{\isacharcolon}{\kern0pt}\ {\isachardoublequoteopen}{\isasymforall}\ {\isacharparenleft}{\kern0pt}a{\isacharcomma}{\kern0pt}t{\isacharparenright}{\kern0pt}\ {\isasymin}\ set\ ys{\isachardot}{\kern0pt}\ nabla{\isadigit{1}}\ {\isasymturnstile}\ a\ {\isasymsharp}\ subst\ {\isacharparenleft}{\kern0pt}s{\isadigit{1}}\ {\isasymbullet}\ {\isacharbrackleft}{\kern0pt}{\isacharbrackright}{\kern0pt}{\isacharparenright}{\kern0pt}\ t{\isachardoublequoteclose}\isanewline
\ \ \ \ \isacommand{using}\isamarkupfalse%
\ all{\isacharunderscore}{\kern0pt}solutions{\isacharunderscore}{\kern0pt}def\ \isacommand{by}\isamarkupfalse%
\ simp{\isacharunderscore}{\kern0pt}all\isanewline
\ \ \isacommand{hence}\isamarkupfalse%
\ eqs{\isacharcolon}{\kern0pt}\ \isanewline
\ \ \ \ {\isachardoublequoteopen}{\isasymforall}\ {\isacharparenleft}{\kern0pt}t{\isadigit{1}}{\isacharcomma}{\kern0pt}t{\isadigit{2}}{\isacharparenright}{\kern0pt}\ {\isasymin}\ set\ {\isacharparenleft}{\kern0pt}{\isacharparenleft}{\kern0pt}Atom\ a{\isacharcomma}{\kern0pt}\ Atom\ a{\isacharparenright}{\kern0pt}\ {\isacharhash}{\kern0pt}\ xs{\isacharparenright}{\kern0pt}{\isachardot}{\kern0pt}\ nabla{\isadigit{1}}\ {\isasymturnstile}\ subst\ {\isacharparenleft}{\kern0pt}s{\isadigit{1}}\ {\isasymbullet}\ {\isacharbrackleft}{\kern0pt}{\isacharbrackright}{\kern0pt}{\isacharparenright}{\kern0pt}\ t{\isadigit{1}}\ {\isasymapprox}\ subst\ {\isacharparenleft}{\kern0pt}s{\isadigit{1}}\ {\isasymbullet}\ {\isacharbrackleft}{\kern0pt}{\isacharbrackright}{\kern0pt}{\isacharparenright}{\kern0pt}\ t{\isadigit{2}}{\isachardoublequoteclose}\isanewline
\ \ \ \ \isacommand{using}\isamarkupfalse%
\ subst{\isacharunderscore}{\kern0pt}unit{\isacharbrackleft}{\kern0pt}of\ {\isacartoucheopen}{\isacharparenleft}{\kern0pt}s{\isadigit{1}}\ {\isasymbullet}\ {\isacharbrackleft}{\kern0pt}{\isacharbrackright}{\kern0pt}{\isacharparenright}{\kern0pt}{\isacartoucheclose}{\isacharbrackright}{\kern0pt}\ equ{\isacharunderscore}{\kern0pt}refl{\isacharbrackleft}{\kern0pt}of\ nabla{\isadigit{1}}\ Unit{\isacharbrackright}{\kern0pt}\ \isacommand{by}\isamarkupfalse%
\ auto\isanewline
\ \ \isacommand{from}\isamarkupfalse%
\ fresh\ eqs\ \isacommand{show}\isamarkupfalse%
\ {\isacharquery}{\kern0pt}case\isanewline
\ \ \ \ \isacommand{using}\isamarkupfalse%
\ all{\isacharunderscore}{\kern0pt}solutions{\isacharunderscore}{\kern0pt}def\ \isacommand{by}\isamarkupfalse%
\ simp\isanewline
\isacommand{next}\isamarkupfalse%
\isanewline
\ \ \isacommand{case}\isamarkupfalse%
\ {\isacharparenleft}{\kern0pt}{\isadigit{7}}\ pi{\isadigit{1}}\ X\ pi{\isadigit{2}}\ xs\ ys{\isacharparenright}{\kern0pt}\isanewline
\ \ \isacommand{hence}\isamarkupfalse%
\ eqs{\isacharunderscore}{\kern0pt}old{\isacharcolon}{\kern0pt}\ {\isachardoublequoteopen}{\isasymforall}\ {\isacharparenleft}{\kern0pt}t{\isadigit{1}}{\isacharcomma}{\kern0pt}t{\isadigit{2}}{\isacharparenright}{\kern0pt}\ {\isasymin}\ set\ xs{\isachardot}{\kern0pt}\ nabla{\isadigit{1}}\ {\isasymturnstile}\ subst\ {\isacharparenleft}{\kern0pt}s{\isadigit{1}}\ {\isasymbullet}\ {\isacharbrackleft}{\kern0pt}{\isacharbrackright}{\kern0pt}{\isacharparenright}{\kern0pt}\ t{\isadigit{1}}\ {\isasymapprox}\ subst\ {\isacharparenleft}{\kern0pt}s{\isadigit{1}}\ {\isasymbullet}\ {\isacharbrackleft}{\kern0pt}{\isacharbrackright}{\kern0pt}{\isacharparenright}{\kern0pt}\ t{\isadigit{2}}{\isachardoublequoteclose}\ \isanewline
\ \ \ \ \ \isakeyword{and}\ fresh{\isacharunderscore}{\kern0pt}old{\isacharcolon}{\kern0pt}\ \isanewline
\ \ \ \ \ {\isachardoublequoteopen}{\isasymforall}\ {\isacharparenleft}{\kern0pt}a{\isacharcomma}{\kern0pt}t{\isacharparenright}{\kern0pt}\ {\isasymin}\ set\ {\isacharparenleft}{\kern0pt}map\ {\isacharparenleft}{\kern0pt}{\isasymlambda}a{\isachardot}{\kern0pt}\ {\isacharparenleft}{\kern0pt}a{\isacharcomma}{\kern0pt}\ Susp\ {\isacharbrackleft}{\kern0pt}{\isacharbrackright}{\kern0pt}\ X{\isacharparenright}{\kern0pt}{\isacharparenright}{\kern0pt}\ {\isacharparenleft}{\kern0pt}ds{\isacharunderscore}{\kern0pt}list\ pi{\isadigit{1}}\ pi{\isadigit{2}}{\isacharparenright}{\kern0pt}\ {\isacharat}{\kern0pt}\ ys{\isacharparenright}{\kern0pt}{\isachardot}{\kern0pt}\ nabla{\isadigit{1}}\ {\isasymturnstile}\ a\ {\isasymsharp}\ subst\ {\isacharparenleft}{\kern0pt}s{\isadigit{1}}\ {\isasymbullet}\ {\isacharbrackleft}{\kern0pt}{\isacharbrackright}{\kern0pt}{\isacharparenright}{\kern0pt}\ t{\isachardoublequoteclose}\isanewline
\ \ \ \ \isacommand{using}\isamarkupfalse%
\ all{\isacharunderscore}{\kern0pt}solutions{\isacharunderscore}{\kern0pt}def\ \isacommand{by}\isamarkupfalse%
\ simp{\isacharunderscore}{\kern0pt}all\isanewline
\ \ \isacommand{hence}\isamarkupfalse%
\ {\isachardoublequoteopen}{\isasymforall}c{\isasymin}\ set\ {\isacharparenleft}{\kern0pt}ds{\isacharunderscore}{\kern0pt}list\ pi{\isadigit{1}}\ pi{\isadigit{2}}{\isacharparenright}{\kern0pt}{\isachardot}{\kern0pt}\ nabla{\isadigit{1}}\ {\isasymturnstile}\ c\ {\isasymsharp}\ subst\ {\isacharparenleft}{\kern0pt}s{\isadigit{1}}\ {\isasymbullet}\ {\isacharbrackleft}{\kern0pt}{\isacharbrackright}{\kern0pt}{\isacharparenright}{\kern0pt}\ {\isacharparenleft}{\kern0pt}Susp\ {\isacharbrackleft}{\kern0pt}{\isacharbrackright}{\kern0pt}\ X{\isacharparenright}{\kern0pt}{\isachardoublequoteclose}\isanewline
\ \ \ \ \isacommand{by}\isamarkupfalse%
\ {\isacharparenleft}{\kern0pt}auto\ split{\isacharcolon}{\kern0pt}\ prod{\isachardot}{\kern0pt}splits{\isacharparenright}{\kern0pt}\isanewline
\ \ \isacommand{hence}\isamarkupfalse%
\ {\isachardoublequoteopen}{\isasymforall}c\ {\isasymin}\ ds\ pi{\isadigit{1}}\ pi{\isadigit{2}}{\isachardot}{\kern0pt}\ nabla{\isadigit{1}}\ {\isasymturnstile}\ c\ {\isasymsharp}\ subst\ {\isacharparenleft}{\kern0pt}s{\isadigit{1}}\ {\isasymbullet}\ {\isacharbrackleft}{\kern0pt}{\isacharbrackright}{\kern0pt}{\isacharparenright}{\kern0pt}\ {\isacharparenleft}{\kern0pt}Susp\ {\isacharbrackleft}{\kern0pt}{\isacharbrackright}{\kern0pt}\ X{\isacharparenright}{\kern0pt}{\isachardoublequoteclose}\isanewline
\ \ \ \ \isacommand{using}\isamarkupfalse%
\ ds{\isacharunderscore}{\kern0pt}list{\isacharunderscore}{\kern0pt}equ{\isacharunderscore}{\kern0pt}ds\ \isacommand{by}\isamarkupfalse%
\ simp\isanewline
\ \ \isacommand{hence}\isamarkupfalse%
\ {\isachardoublequoteopen}{\isasymforall}c\ {\isasymin}\ ds\ pi{\isadigit{1}}\ pi{\isadigit{2}}{\isachardot}{\kern0pt}\ nabla{\isadigit{1}}\ {\isasymturnstile}\ c\ {\isasymsharp}\ {\isacharparenleft}{\kern0pt}look{\isacharunderscore}{\kern0pt}up\ X\ {\isacharparenleft}{\kern0pt}s{\isadigit{1}}\ {\isasymbullet}\ {\isacharbrackleft}{\kern0pt}{\isacharbrackright}{\kern0pt}{\isacharparenright}{\kern0pt}{\isacharparenright}{\kern0pt}{\isachardoublequoteclose}\isanewline
\ \ \ \ \isacommand{using}\isamarkupfalse%
\ subst{\isacharunderscore}{\kern0pt}susp{\isacharbrackleft}{\kern0pt}of\ {\isacartoucheopen}{\isacharparenleft}{\kern0pt}s{\isadigit{1}}\ {\isasymbullet}\ {\isacharbrackleft}{\kern0pt}{\isacharbrackright}{\kern0pt}{\isacharparenright}{\kern0pt}{\isacartoucheclose}\ {\isacartoucheopen}{\isacharbrackleft}{\kern0pt}{\isacharbrackright}{\kern0pt}{\isacartoucheclose}\ X{\isacharbrackright}{\kern0pt}\ swap{\isacharunderscore}{\kern0pt}id\ \isacommand{by}\isamarkupfalse%
\ simp\isanewline
\ \ \isacommand{hence}\isamarkupfalse%
\ {\isachardoublequoteopen}nabla{\isadigit{1}}\ {\isasymturnstile}\ subst\ {\isacharparenleft}{\kern0pt}s{\isadigit{1}}\ {\isasymbullet}\ {\isacharbrackleft}{\kern0pt}{\isacharbrackright}{\kern0pt}{\isacharparenright}{\kern0pt}\ {\isacharparenleft}{\kern0pt}Susp\ pi{\isadigit{1}}\ X{\isacharparenright}{\kern0pt}\ {\isasymapprox}\ subst\ {\isacharparenleft}{\kern0pt}s{\isadigit{1}}\ {\isasymbullet}\ {\isacharbrackleft}{\kern0pt}{\isacharbrackright}{\kern0pt}{\isacharparenright}{\kern0pt}\ {\isacharparenleft}{\kern0pt}Susp\ pi{\isadigit{2}}\ X{\isacharparenright}{\kern0pt}{\isachardoublequoteclose}\isanewline
\ \ \ \ \isacommand{using}\isamarkupfalse%
\ \ equ{\isacharunderscore}{\kern0pt}equiv{\isacharunderscore}{\kern0pt}pi\ subst{\isacharunderscore}{\kern0pt}susp{\isacharbrackleft}{\kern0pt}of\ {\isacartoucheopen}{\isacharparenleft}{\kern0pt}s{\isadigit{1}}\ {\isasymbullet}\ {\isacharbrackleft}{\kern0pt}{\isacharbrackright}{\kern0pt}{\isacharparenright}{\kern0pt}{\isacartoucheclose}\ {\isacharunderscore}{\kern0pt}\ X{\isacharbrackright}{\kern0pt}\ \isacommand{by}\isamarkupfalse%
\ auto\isanewline
\ \ \isacommand{with}\isamarkupfalse%
\ eqs{\isacharunderscore}{\kern0pt}old\ fresh{\isacharunderscore}{\kern0pt}old\ \isacommand{have}\isamarkupfalse%
\isanewline
\ \ eqs{\isacharcolon}{\kern0pt}\ {\isachardoublequoteopen}{\isasymforall}\ {\isacharparenleft}{\kern0pt}t{\isadigit{1}}{\isacharcomma}{\kern0pt}t{\isadigit{2}}{\isacharparenright}{\kern0pt}\ {\isasymin}\ set\ {\isacharparenleft}{\kern0pt}{\isacharparenleft}{\kern0pt}Susp\ pi{\isadigit{1}}\ X{\isacharcomma}{\kern0pt}\ Susp\ pi{\isadigit{2}}\ X{\isacharparenright}{\kern0pt}{\isacharhash}{\kern0pt}xs{\isacharparenright}{\kern0pt}{\isachardot}{\kern0pt}\ \isanewline
\ \ \ \ nabla{\isadigit{1}}\ {\isasymturnstile}\ subst\ {\isacharparenleft}{\kern0pt}s{\isadigit{1}}\ {\isasymbullet}\ {\isacharbrackleft}{\kern0pt}{\isacharbrackright}{\kern0pt}{\isacharparenright}{\kern0pt}\ t{\isadigit{1}}\ {\isasymapprox}\ subst\ {\isacharparenleft}{\kern0pt}s{\isadigit{1}}\ {\isasymbullet}\ {\isacharbrackleft}{\kern0pt}{\isacharbrackright}{\kern0pt}{\isacharparenright}{\kern0pt}\ t{\isadigit{2}}{\isachardoublequoteclose}\isanewline
\ \ \isakeyword{and}\ fresh{\isacharcolon}{\kern0pt}\ {\isachardoublequoteopen}{\isasymforall}\ {\isacharparenleft}{\kern0pt}a{\isacharcomma}{\kern0pt}t{\isacharparenright}{\kern0pt}\ {\isasymin}\ set\ ys{\isachardot}{\kern0pt}\ nabla{\isadigit{1}}\ {\isasymturnstile}\ a\ {\isasymsharp}\ subst\ {\isacharparenleft}{\kern0pt}s{\isadigit{1}}\ {\isasymbullet}\ {\isacharbrackleft}{\kern0pt}{\isacharbrackright}{\kern0pt}{\isacharparenright}{\kern0pt}\ t{\isachardoublequoteclose}\isanewline
\ \ \ \ \isacommand{by}\isamarkupfalse%
\ simp{\isacharunderscore}{\kern0pt}all\isanewline
\ \ \isacommand{then}\isamarkupfalse%
\ \isacommand{show}\isamarkupfalse%
\ {\isacharquery}{\kern0pt}case\ \isanewline
\ \ \ \ \isacommand{using}\isamarkupfalse%
\ all{\isacharunderscore}{\kern0pt}solutions{\isacharunderscore}{\kern0pt}def\ \isacommand{by}\isamarkupfalse%
\ simp\isanewline
\isacommand{next}\isamarkupfalse%
\isanewline
\ \ \isacommand{case}\isamarkupfalse%
\ {\isacharparenleft}{\kern0pt}{\isadigit{8}}\ X\ t{\isacharprime}{\kern0pt}\ pi\ xs\ ys{\isacharparenright}{\kern0pt}\isanewline
\ \ \isacommand{hence}\isamarkupfalse%
\ {\isachardoublequoteopen}{\isacharparenleft}{\kern0pt}nabla{\isadigit{1}}{\isacharcomma}{\kern0pt}\ s{\isadigit{1}}\ {\isasymbullet}\ {\isacharbrackleft}{\kern0pt}{\isacharparenleft}{\kern0pt}X{\isacharcomma}{\kern0pt}\ swap\ {\isacharparenleft}{\kern0pt}rev\ pi{\isacharparenright}{\kern0pt}\ t{\isacharprime}{\kern0pt}{\isacharparenright}{\kern0pt}{\isacharbrackright}{\kern0pt}{\isacharparenright}{\kern0pt}\ {\isasymin}\ U\ {\isacharparenleft}{\kern0pt}xs{\isacharcomma}{\kern0pt}\ ys{\isacharparenright}{\kern0pt}{\isachardoublequoteclose}\isanewline
\ \ \ \ \isacommand{using}\isamarkupfalse%
\ problem{\isacharunderscore}{\kern0pt}subst{\isacharunderscore}{\kern0pt}comm\ \isacommand{by}\isamarkupfalse%
\ simp\isanewline
\ \ \isacommand{hence}\isamarkupfalse%
\ eqs{\isacharunderscore}{\kern0pt}old{\isacharcolon}{\kern0pt}\ \isanewline
\ \ \ \ {\isachardoublequoteopen}{\isasymforall}\ {\isacharparenleft}{\kern0pt}t{\isadigit{1}}{\isacharcomma}{\kern0pt}t{\isadigit{2}}{\isacharparenright}{\kern0pt}\ {\isasymin}\ set\ xs{\isachardot}{\kern0pt}\ \isanewline
\ \ \ \ nabla{\isadigit{1}}\ {\isasymturnstile}\ subst\ {\isacharparenleft}{\kern0pt}s{\isadigit{1}}\ {\isasymbullet}\ {\isacharbrackleft}{\kern0pt}{\isacharparenleft}{\kern0pt}X{\isacharcomma}{\kern0pt}\ swap\ {\isacharparenleft}{\kern0pt}rev\ pi{\isacharparenright}{\kern0pt}\ t{\isacharprime}{\kern0pt}{\isacharparenright}{\kern0pt}{\isacharbrackright}{\kern0pt}{\isacharparenright}{\kern0pt}\ t{\isadigit{1}}\ {\isasymapprox}\ subst\ {\isacharparenleft}{\kern0pt}s{\isadigit{1}}\ {\isasymbullet}\ {\isacharbrackleft}{\kern0pt}{\isacharparenleft}{\kern0pt}X{\isacharcomma}{\kern0pt}\ swap\ {\isacharparenleft}{\kern0pt}rev\ pi{\isacharparenright}{\kern0pt}\ t{\isacharprime}{\kern0pt}{\isacharparenright}{\kern0pt}{\isacharbrackright}{\kern0pt}{\isacharparenright}{\kern0pt}\ t{\isadigit{2}}{\isachardoublequoteclose}\isanewline
\ \ \isakeyword{and}\ fresh{\isacharcolon}{\kern0pt}\ {\isachardoublequoteopen}{\isasymforall}\ {\isacharparenleft}{\kern0pt}a{\isacharcomma}{\kern0pt}t{\isacharparenright}{\kern0pt}\ {\isasymin}\ set\ ys{\isachardot}{\kern0pt}\ nabla{\isadigit{1}}\ {\isasymturnstile}\ a\ {\isasymsharp}\ subst\ {\isacharparenleft}{\kern0pt}s{\isadigit{1}}\ {\isasymbullet}\ {\isacharbrackleft}{\kern0pt}{\isacharparenleft}{\kern0pt}X{\isacharcomma}{\kern0pt}\ swap\ {\isacharparenleft}{\kern0pt}rev\ pi{\isacharparenright}{\kern0pt}\ t{\isacharprime}{\kern0pt}{\isacharparenright}{\kern0pt}{\isacharbrackright}{\kern0pt}{\isacharparenright}{\kern0pt}\ t{\isachardoublequoteclose}\isanewline
\ \ \ \ \isacommand{using}\isamarkupfalse%
\ all{\isacharunderscore}{\kern0pt}solutions{\isacharunderscore}{\kern0pt}def\ \isacommand{by}\isamarkupfalse%
\ simp{\isacharunderscore}{\kern0pt}all\isanewline
\isanewline
\ \ \isacommand{have}\isamarkupfalse%
\ {\isachardoublequoteopen}subst\ {\isacharparenleft}{\kern0pt}s{\isadigit{1}}\ {\isasymbullet}\ {\isacharbrackleft}{\kern0pt}{\isacharparenleft}{\kern0pt}X{\isacharcomma}{\kern0pt}\ swap\ {\isacharparenleft}{\kern0pt}rev\ pi{\isacharparenright}{\kern0pt}\ t{\isacharprime}{\kern0pt}{\isacharparenright}{\kern0pt}{\isacharbrackright}{\kern0pt}{\isacharparenright}{\kern0pt}\ {\isacharparenleft}{\kern0pt}Susp\ pi\ X{\isacharparenright}{\kern0pt}\ {\isacharequal}{\kern0pt}\ subst\ s{\isadigit{1}}\ {\isacharparenleft}{\kern0pt}swap\ pi\ {\isacharparenleft}{\kern0pt}swap\ {\isacharparenleft}{\kern0pt}rev\ pi{\isacharparenright}{\kern0pt}\ t{\isacharprime}{\kern0pt}{\isacharparenright}{\kern0pt}{\isacharparenright}{\kern0pt}{\isachardoublequoteclose}\isanewline
\ \ \ \ \isacommand{using}\isamarkupfalse%
\ subst{\isacharunderscore}{\kern0pt}susp\ subst{\isacharunderscore}{\kern0pt}swap{\isacharunderscore}{\kern0pt}comm\ subst{\isacharunderscore}{\kern0pt}comp{\isacharunderscore}{\kern0pt}expand\ eq{\isacharunderscore}{\kern0pt}fst{\isacharunderscore}{\kern0pt}iff\ look{\isacharunderscore}{\kern0pt}up{\isachardot}{\kern0pt}simps{\isacharparenleft}{\kern0pt}{\isadigit{2}}{\isacharparenright}{\kern0pt}\ snd{\isacharunderscore}{\kern0pt}eqD\ \isacommand{by}\isamarkupfalse%
\ metis\isanewline
\ \ \isacommand{also}\isamarkupfalse%
\ \isacommand{have}\isamarkupfalse%
\ {\isachardoublequoteopen}{\isachardot}{\kern0pt}{\isachardot}{\kern0pt}{\isachardot}{\kern0pt}\ {\isacharequal}{\kern0pt}\ swap\ pi\ {\isacharparenleft}{\kern0pt}swap\ {\isacharparenleft}{\kern0pt}rev\ pi{\isacharparenright}{\kern0pt}\ {\isacharparenleft}{\kern0pt}subst\ s{\isadigit{1}}\ t{\isacharprime}{\kern0pt}{\isacharparenright}{\kern0pt}{\isacharparenright}{\kern0pt}{\isachardoublequoteclose}\isanewline
\ \ \ \ \isacommand{using}\isamarkupfalse%
\ subst{\isacharunderscore}{\kern0pt}swap{\isacharunderscore}{\kern0pt}comm\ \isacommand{by}\isamarkupfalse%
\ simp\isanewline
\ \ \isacommand{hence}\isamarkupfalse%
\ {\isachardoublequoteopen}nabla{\isadigit{1}}\ {\isasymturnstile}\ subst\ {\isacharparenleft}{\kern0pt}s{\isadigit{1}}\ {\isasymbullet}\ {\isacharbrackleft}{\kern0pt}{\isacharparenleft}{\kern0pt}X{\isacharcomma}{\kern0pt}\ swap\ {\isacharparenleft}{\kern0pt}rev\ pi{\isacharparenright}{\kern0pt}\ t{\isacharprime}{\kern0pt}{\isacharparenright}{\kern0pt}{\isacharbrackright}{\kern0pt}{\isacharparenright}{\kern0pt}\ {\isacharparenleft}{\kern0pt}Susp\ pi\ X{\isacharparenright}{\kern0pt}\ {\isasymapprox}\ subst\ s{\isadigit{1}}\ t{\isacharprime}{\kern0pt}{\isachardoublequoteclose}\isanewline
\ \ \ \ \isacommand{using}\isamarkupfalse%
\ equ{\isacharunderscore}{\kern0pt}refl\ calculation\ rev{\isacharunderscore}{\kern0pt}pi{\isacharunderscore}{\kern0pt}pi{\isacharunderscore}{\kern0pt}equ{\isacharbrackleft}{\kern0pt}of\ nabla{\isadigit{1}}\ pi\ {\isacartoucheopen}subst\ s{\isadigit{1}}\ t{\isacharprime}{\kern0pt}{\isacartoucheclose}{\isacharbrackright}{\kern0pt}\isanewline
\ \ \ \ \ \ equ{\isacharunderscore}{\kern0pt}trans{\isacharbrackleft}{\kern0pt}of\ nabla{\isadigit{1}}\ {\isacartoucheopen}subst\ {\isacharparenleft}{\kern0pt}s{\isadigit{1}}\ {\isasymbullet}\ {\isacharbrackleft}{\kern0pt}{\isacharparenleft}{\kern0pt}X{\isacharcomma}{\kern0pt}\ swap\ {\isacharparenleft}{\kern0pt}rev\ pi{\isacharparenright}{\kern0pt}\ t{\isacharprime}{\kern0pt}{\isacharparenright}{\kern0pt}{\isacharbrackright}{\kern0pt}{\isacharparenright}{\kern0pt}\ {\isacharparenleft}{\kern0pt}Susp\ pi\ X{\isacharparenright}{\kern0pt}{\isacartoucheclose}\ \isanewline
\ \ \ \ \ \ \ \ {\isacartoucheopen}swap\ {\isacharparenleft}{\kern0pt}rev\ pi{\isacharparenright}{\kern0pt}\ {\isacharparenleft}{\kern0pt}swap\ pi\ {\isacharparenleft}{\kern0pt}subst\ s{\isadigit{1}}\ t{\isacharprime}{\kern0pt}{\isacharparenright}{\kern0pt}{\isacharparenright}{\kern0pt}{\isacartoucheclose}\ {\isacartoucheopen}subst\ s{\isadigit{1}}\ t{\isacharprime}{\kern0pt}{\isacartoucheclose}{\isacharbrackright}{\kern0pt}\isanewline
\ \ \ \ \ \ \ swap{\isacharunderscore}{\kern0pt}inv{\isacharunderscore}{\kern0pt}side\ \isacommand{by}\isamarkupfalse%
\ auto\isanewline
\isanewline
\ \ \isacommand{with}\isamarkupfalse%
\ {\isadigit{8}}\ eqs{\isacharunderscore}{\kern0pt}old\isanewline
\ \ \isacommand{have}\isamarkupfalse%
\ eqs{\isacharcolon}{\kern0pt}\ {\isachardoublequoteopen}{\isasymforall}\ {\isacharparenleft}{\kern0pt}t{\isadigit{1}}{\isacharcomma}{\kern0pt}t{\isadigit{2}}{\isacharparenright}{\kern0pt}\ {\isasymin}\ set\ {\isacharparenleft}{\kern0pt}{\isacharparenleft}{\kern0pt}Susp\ pi\ X{\isacharcomma}{\kern0pt}\ t{\isacharprime}{\kern0pt}{\isacharparenright}{\kern0pt}{\isacharhash}{\kern0pt}xs{\isacharparenright}{\kern0pt}{\isachardot}{\kern0pt}\ \isanewline
\ \ \ \ nabla{\isadigit{1}}\ {\isasymturnstile}\ subst\ {\isacharparenleft}{\kern0pt}s{\isadigit{1}}\ {\isasymbullet}\ {\isacharbrackleft}{\kern0pt}{\isacharparenleft}{\kern0pt}X{\isacharcomma}{\kern0pt}\ swap\ {\isacharparenleft}{\kern0pt}rev\ pi{\isacharparenright}{\kern0pt}\ t{\isacharprime}{\kern0pt}{\isacharparenright}{\kern0pt}{\isacharbrackright}{\kern0pt}{\isacharparenright}{\kern0pt}\ t{\isadigit{1}}\ {\isasymapprox}\ subst\ {\isacharparenleft}{\kern0pt}s{\isadigit{1}}\ {\isasymbullet}\ {\isacharbrackleft}{\kern0pt}{\isacharparenleft}{\kern0pt}X{\isacharcomma}{\kern0pt}\ swap\ {\isacharparenleft}{\kern0pt}rev\ pi{\isacharparenright}{\kern0pt}\ t{\isacharprime}{\kern0pt}{\isacharparenright}{\kern0pt}{\isacharbrackright}{\kern0pt}{\isacharparenright}{\kern0pt}\ t{\isadigit{2}}{\isachardoublequoteclose}\isanewline
\ \ \ \ \isacommand{using}\isamarkupfalse%
\ subst{\isacharunderscore}{\kern0pt}comp{\isacharunderscore}{\kern0pt}expand\ subst{\isacharunderscore}{\kern0pt}not{\isacharunderscore}{\kern0pt}occurs\ \isacommand{by}\isamarkupfalse%
\ simp\isanewline
\isanewline
\ \ \isacommand{from}\isamarkupfalse%
\ fresh\ eqs\ \isacommand{show}\isamarkupfalse%
\ {\isacharquery}{\kern0pt}case\isanewline
\ \ \ \ \isacommand{using}\isamarkupfalse%
\ all{\isacharunderscore}{\kern0pt}solutions{\isacharunderscore}{\kern0pt}def\ \isacommand{by}\isamarkupfalse%
\ simp\isanewline
\isacommand{next}\isamarkupfalse%
\isanewline
\ \ \isacommand{case}\isamarkupfalse%
\ {\isacharparenleft}{\kern0pt}{\isadigit{9}}\ X\ t{\isacharprime}{\kern0pt}\ pi\ xs\ ys{\isacharparenright}{\kern0pt}\isanewline
\ \ \isacommand{hence}\isamarkupfalse%
\ {\isachardoublequoteopen}{\isacharparenleft}{\kern0pt}nabla{\isadigit{1}}{\isacharcomma}{\kern0pt}\ s{\isadigit{1}}\ {\isasymbullet}\ {\isacharbrackleft}{\kern0pt}{\isacharparenleft}{\kern0pt}X{\isacharcomma}{\kern0pt}\ swap\ {\isacharparenleft}{\kern0pt}rev\ pi{\isacharparenright}{\kern0pt}\ t{\isacharprime}{\kern0pt}{\isacharparenright}{\kern0pt}{\isacharbrackright}{\kern0pt}{\isacharparenright}{\kern0pt}\ {\isasymin}\ U\ {\isacharparenleft}{\kern0pt}xs{\isacharcomma}{\kern0pt}\ ys{\isacharparenright}{\kern0pt}{\isachardoublequoteclose}\isanewline
\ \ \ \ \isacommand{using}\isamarkupfalse%
\ problem{\isacharunderscore}{\kern0pt}subst{\isacharunderscore}{\kern0pt}comm\ \isacommand{by}\isamarkupfalse%
\ simp\isanewline
\ \ \isacommand{hence}\isamarkupfalse%
\ eqs{\isacharunderscore}{\kern0pt}old{\isacharcolon}{\kern0pt}\ \isanewline
\ \ \ \ {\isachardoublequoteopen}{\isasymforall}\ {\isacharparenleft}{\kern0pt}t{\isadigit{1}}{\isacharcomma}{\kern0pt}t{\isadigit{2}}{\isacharparenright}{\kern0pt}\ {\isasymin}\ set\ xs{\isachardot}{\kern0pt}\ \isanewline
\ \ \ \ nabla{\isadigit{1}}\ {\isasymturnstile}\ subst\ {\isacharparenleft}{\kern0pt}s{\isadigit{1}}\ {\isasymbullet}\ {\isacharbrackleft}{\kern0pt}{\isacharparenleft}{\kern0pt}X{\isacharcomma}{\kern0pt}\ swap\ {\isacharparenleft}{\kern0pt}rev\ pi{\isacharparenright}{\kern0pt}\ t{\isacharprime}{\kern0pt}{\isacharparenright}{\kern0pt}{\isacharbrackright}{\kern0pt}{\isacharparenright}{\kern0pt}\ t{\isadigit{1}}\ {\isasymapprox}\ subst\ {\isacharparenleft}{\kern0pt}s{\isadigit{1}}\ {\isasymbullet}\ {\isacharbrackleft}{\kern0pt}{\isacharparenleft}{\kern0pt}X{\isacharcomma}{\kern0pt}\ swap\ {\isacharparenleft}{\kern0pt}rev\ pi{\isacharparenright}{\kern0pt}\ t{\isacharprime}{\kern0pt}{\isacharparenright}{\kern0pt}{\isacharbrackright}{\kern0pt}{\isacharparenright}{\kern0pt}\ t{\isadigit{2}}{\isachardoublequoteclose}\isanewline
\ \ \isakeyword{and}\ fresh{\isacharcolon}{\kern0pt}\ {\isachardoublequoteopen}{\isasymforall}\ {\isacharparenleft}{\kern0pt}a{\isacharcomma}{\kern0pt}t{\isacharparenright}{\kern0pt}\ {\isasymin}\ set\ ys{\isachardot}{\kern0pt}\ nabla{\isadigit{1}}\ {\isasymturnstile}\ a\ {\isasymsharp}\ subst\ {\isacharparenleft}{\kern0pt}s{\isadigit{1}}\ {\isasymbullet}\ {\isacharbrackleft}{\kern0pt}{\isacharparenleft}{\kern0pt}X{\isacharcomma}{\kern0pt}\ swap\ {\isacharparenleft}{\kern0pt}rev\ pi{\isacharparenright}{\kern0pt}\ t{\isacharprime}{\kern0pt}{\isacharparenright}{\kern0pt}{\isacharbrackright}{\kern0pt}{\isacharparenright}{\kern0pt}\ t{\isachardoublequoteclose}\isanewline
\ \ \ \ \isacommand{using}\isamarkupfalse%
\ all{\isacharunderscore}{\kern0pt}solutions{\isacharunderscore}{\kern0pt}def\ \isacommand{by}\isamarkupfalse%
\ simp{\isacharunderscore}{\kern0pt}all\isanewline
\isanewline
\ \ \isacommand{have}\isamarkupfalse%
\ {\isachardoublequoteopen}subst\ {\isacharparenleft}{\kern0pt}s{\isadigit{1}}\ {\isasymbullet}\ {\isacharbrackleft}{\kern0pt}{\isacharparenleft}{\kern0pt}X{\isacharcomma}{\kern0pt}\ swap\ {\isacharparenleft}{\kern0pt}rev\ pi{\isacharparenright}{\kern0pt}\ t{\isacharprime}{\kern0pt}{\isacharparenright}{\kern0pt}{\isacharbrackright}{\kern0pt}{\isacharparenright}{\kern0pt}\ {\isacharparenleft}{\kern0pt}Susp\ pi\ X{\isacharparenright}{\kern0pt}\ {\isacharequal}{\kern0pt}\ subst\ s{\isadigit{1}}\ {\isacharparenleft}{\kern0pt}swap\ pi\ {\isacharparenleft}{\kern0pt}swap\ {\isacharparenleft}{\kern0pt}rev\ pi{\isacharparenright}{\kern0pt}\ t{\isacharprime}{\kern0pt}{\isacharparenright}{\kern0pt}{\isacharparenright}{\kern0pt}{\isachardoublequoteclose}\isanewline
\ \ \ \ \isacommand{using}\isamarkupfalse%
\ subst{\isacharunderscore}{\kern0pt}susp\ subst{\isacharunderscore}{\kern0pt}swap{\isacharunderscore}{\kern0pt}comm\ subst{\isacharunderscore}{\kern0pt}comp{\isacharunderscore}{\kern0pt}expand\ eq{\isacharunderscore}{\kern0pt}fst{\isacharunderscore}{\kern0pt}iff\ look{\isacharunderscore}{\kern0pt}up{\isachardot}{\kern0pt}simps{\isacharparenleft}{\kern0pt}{\isadigit{2}}{\isacharparenright}{\kern0pt}\ snd{\isacharunderscore}{\kern0pt}eqD\ \isacommand{by}\isamarkupfalse%
\ metis\isanewline
\ \ \isacommand{also}\isamarkupfalse%
\ \isacommand{have}\isamarkupfalse%
\ {\isachardoublequoteopen}{\isachardot}{\kern0pt}{\isachardot}{\kern0pt}{\isachardot}{\kern0pt}\ {\isacharequal}{\kern0pt}\ swap\ pi\ {\isacharparenleft}{\kern0pt}swap\ {\isacharparenleft}{\kern0pt}rev\ pi{\isacharparenright}{\kern0pt}\ {\isacharparenleft}{\kern0pt}subst\ s{\isadigit{1}}\ t{\isacharprime}{\kern0pt}{\isacharparenright}{\kern0pt}{\isacharparenright}{\kern0pt}{\isachardoublequoteclose}\isanewline
\ \ \ \ \isacommand{using}\isamarkupfalse%
\ subst{\isacharunderscore}{\kern0pt}swap{\isacharunderscore}{\kern0pt}comm\ \isacommand{by}\isamarkupfalse%
\ simp\isanewline
\ \ \isacommand{hence}\isamarkupfalse%
\ {\isachardoublequoteopen}nabla{\isadigit{1}}\ {\isasymturnstile}\ subst\ {\isacharparenleft}{\kern0pt}s{\isadigit{1}}\ {\isasymbullet}\ {\isacharbrackleft}{\kern0pt}{\isacharparenleft}{\kern0pt}X{\isacharcomma}{\kern0pt}\ swap\ {\isacharparenleft}{\kern0pt}rev\ pi{\isacharparenright}{\kern0pt}\ t{\isacharprime}{\kern0pt}{\isacharparenright}{\kern0pt}{\isacharbrackright}{\kern0pt}{\isacharparenright}{\kern0pt}\ {\isacharparenleft}{\kern0pt}Susp\ pi\ X{\isacharparenright}{\kern0pt}\ {\isasymapprox}\ subst\ s{\isadigit{1}}\ t{\isacharprime}{\kern0pt}{\isachardoublequoteclose}\isanewline
\ \ \ \ \isacommand{using}\isamarkupfalse%
\ equ{\isacharunderscore}{\kern0pt}refl\ calculation\ rev{\isacharunderscore}{\kern0pt}pi{\isacharunderscore}{\kern0pt}pi{\isacharunderscore}{\kern0pt}equ{\isacharbrackleft}{\kern0pt}of\ nabla{\isadigit{1}}\ pi\ {\isacartoucheopen}subst\ s{\isadigit{1}}\ t{\isacharprime}{\kern0pt}{\isacartoucheclose}{\isacharbrackright}{\kern0pt}\isanewline
\ \ \ \ \ \ equ{\isacharunderscore}{\kern0pt}trans{\isacharbrackleft}{\kern0pt}of\ nabla{\isadigit{1}}\ {\isacartoucheopen}subst\ {\isacharparenleft}{\kern0pt}s{\isadigit{1}}\ {\isasymbullet}\ {\isacharbrackleft}{\kern0pt}{\isacharparenleft}{\kern0pt}X{\isacharcomma}{\kern0pt}\ swap\ {\isacharparenleft}{\kern0pt}rev\ pi{\isacharparenright}{\kern0pt}\ t{\isacharprime}{\kern0pt}{\isacharparenright}{\kern0pt}{\isacharbrackright}{\kern0pt}{\isacharparenright}{\kern0pt}\ {\isacharparenleft}{\kern0pt}Susp\ pi\ X{\isacharparenright}{\kern0pt}{\isacartoucheclose}\ \isanewline
\ \ \ \ \ \ \ \ {\isacartoucheopen}swap\ {\isacharparenleft}{\kern0pt}rev\ pi{\isacharparenright}{\kern0pt}\ {\isacharparenleft}{\kern0pt}swap\ pi\ {\isacharparenleft}{\kern0pt}subst\ s{\isadigit{1}}\ t{\isacharprime}{\kern0pt}{\isacharparenright}{\kern0pt}{\isacharparenright}{\kern0pt}{\isacartoucheclose}\ {\isacartoucheopen}subst\ s{\isadigit{1}}\ t{\isacharprime}{\kern0pt}{\isacartoucheclose}{\isacharbrackright}{\kern0pt}\isanewline
\ \ \ \ \ \ \ swap{\isacharunderscore}{\kern0pt}inv{\isacharunderscore}{\kern0pt}side\ \isacommand{by}\isamarkupfalse%
\ auto\isanewline
\ \ \isacommand{hence}\isamarkupfalse%
\ {\isachardoublequoteopen}nabla{\isadigit{1}}\ {\isasymturnstile}\ subst\ s{\isadigit{1}}\ t{\isacharprime}{\kern0pt}\ {\isasymapprox}\ subst\ {\isacharparenleft}{\kern0pt}s{\isadigit{1}}\ {\isasymbullet}\ {\isacharbrackleft}{\kern0pt}{\isacharparenleft}{\kern0pt}X{\isacharcomma}{\kern0pt}\ swap\ {\isacharparenleft}{\kern0pt}rev\ pi{\isacharparenright}{\kern0pt}\ t{\isacharprime}{\kern0pt}{\isacharparenright}{\kern0pt}{\isacharbrackright}{\kern0pt}{\isacharparenright}{\kern0pt}\ {\isacharparenleft}{\kern0pt}Susp\ pi\ X{\isacharparenright}{\kern0pt}{\isachardoublequoteclose}\isanewline
\ \ \ \ \isacommand{using}\isamarkupfalse%
\ equ{\isacharunderscore}{\kern0pt}symm\ \isacommand{by}\isamarkupfalse%
\ simp\isanewline
\isanewline
\ \ \isacommand{with}\isamarkupfalse%
\ {\isadigit{9}}\ eqs{\isacharunderscore}{\kern0pt}old\isanewline
\ \ \isacommand{have}\isamarkupfalse%
\ eqs{\isacharcolon}{\kern0pt}\ {\isachardoublequoteopen}{\isasymforall}\ {\isacharparenleft}{\kern0pt}t{\isadigit{1}}{\isacharcomma}{\kern0pt}t{\isadigit{2}}{\isacharparenright}{\kern0pt}\ {\isasymin}\ set\ {\isacharparenleft}{\kern0pt}{\isacharparenleft}{\kern0pt}t{\isacharprime}{\kern0pt}{\isacharcomma}{\kern0pt}\ Susp\ pi\ X{\isacharparenright}{\kern0pt}{\isacharhash}{\kern0pt}xs{\isacharparenright}{\kern0pt}{\isachardot}{\kern0pt}\ \isanewline
\ \ \ \ nabla{\isadigit{1}}\ {\isasymturnstile}\ subst\ {\isacharparenleft}{\kern0pt}s{\isadigit{1}}\ {\isasymbullet}\ {\isacharbrackleft}{\kern0pt}{\isacharparenleft}{\kern0pt}X{\isacharcomma}{\kern0pt}\ swap\ {\isacharparenleft}{\kern0pt}rev\ pi{\isacharparenright}{\kern0pt}\ t{\isacharprime}{\kern0pt}{\isacharparenright}{\kern0pt}{\isacharbrackright}{\kern0pt}{\isacharparenright}{\kern0pt}\ t{\isadigit{1}}\ {\isasymapprox}\ subst\ {\isacharparenleft}{\kern0pt}s{\isadigit{1}}\ {\isasymbullet}\ {\isacharbrackleft}{\kern0pt}{\isacharparenleft}{\kern0pt}X{\isacharcomma}{\kern0pt}\ swap\ {\isacharparenleft}{\kern0pt}rev\ pi{\isacharparenright}{\kern0pt}\ t{\isacharprime}{\kern0pt}{\isacharparenright}{\kern0pt}{\isacharbrackright}{\kern0pt}{\isacharparenright}{\kern0pt}\ t{\isadigit{2}}{\isachardoublequoteclose}\isanewline
\ \ \ \ \isacommand{using}\isamarkupfalse%
\ subst{\isacharunderscore}{\kern0pt}comp{\isacharunderscore}{\kern0pt}expand\ subst{\isacharunderscore}{\kern0pt}not{\isacharunderscore}{\kern0pt}occurs\ \isacommand{by}\isamarkupfalse%
\ simp\isanewline
\isanewline
\ \ \isacommand{from}\isamarkupfalse%
\ fresh\ eqs\ \isacommand{show}\isamarkupfalse%
\ {\isacharquery}{\kern0pt}case\isanewline
\ \ \ \ \isacommand{using}\isamarkupfalse%
\ all{\isacharunderscore}{\kern0pt}solutions{\isacharunderscore}{\kern0pt}def\ \isacommand{by}\isamarkupfalse%
\ simp\isanewline
\isacommand{qed}\isamarkupfalse%
%
\endisatagproof
{\isafoldproof}%
%
\isadelimproof
\isanewline
%
\endisadelimproof
\isanewline
\isanewline
\isacommand{lemma}\isamarkupfalse%
\ P{\isadigit{1}}{\isacharunderscore}{\kern0pt}to{\isacharunderscore}{\kern0pt}P{\isadigit{2}}{\isacharunderscore}{\kern0pt}cred{\isacharcolon}{\kern0pt}\ \isanewline
\ \ \isakeyword{assumes}\ {\isachardoublequoteopen}{\isacharparenleft}{\kern0pt}nabla{\isadigit{1}}{\isacharcomma}{\kern0pt}s{\isadigit{1}}{\isacharparenright}{\kern0pt}{\isasymin}\ U\ P{\isadigit{1}}{\isachardoublequoteclose}\isanewline
\ \ \isakeyword{and}\ {\isachardoublequoteopen}P{\isadigit{1}}\ {\isasymturnstile}nabla{\isadigit{2}}\ {\isasymrightarrow}\ P{\isadigit{2}}{\isachardoublequoteclose}\isanewline
\isakeyword{shows}\ {\isachardoublequoteopen}{\isacharparenleft}{\kern0pt}nabla{\isadigit{1}}{\isacharcomma}{\kern0pt}s{\isadigit{1}}{\isacharparenright}{\kern0pt}{\isasymin}\ U\ P{\isadigit{2}}\ \ {\isasymand}\ {\isacharparenleft}{\kern0pt}nabla{\isadigit{1}}{\isasymTurnstile}{\isacharparenleft}{\kern0pt}subst\ s{\isadigit{1}}{\isacharparenright}{\kern0pt}\ nabla{\isadigit{2}}{\isacharparenright}{\kern0pt}{\isachardoublequoteclose}\isanewline
%
\isadelimproof
\ \ %
\endisadelimproof
%
\isatagproof
\isacommand{using}\isamarkupfalse%
\ assms\isanewline
\isacommand{proof}\isamarkupfalse%
{\isacharparenleft}{\kern0pt}induction\ arbitrary{\isacharcolon}{\kern0pt}\ nabla{\isadigit{1}}\ s{\isadigit{1}}\ rule{\isacharcolon}{\kern0pt}\ c{\isacharunderscore}{\kern0pt}red{\isachardot}{\kern0pt}induct{\isacharbrackleft}{\kern0pt}OF\ {\isacartoucheopen}P{\isadigit{1}}\ {\isasymturnstile}nabla{\isadigit{2}}{\isasymrightarrow}\ P{\isadigit{2}}{\isacartoucheclose}{\isacharbrackright}{\kern0pt}{\isacharparenright}{\kern0pt}\isanewline
\ \ \isacommand{case}\isamarkupfalse%
\ {\isacharparenleft}{\kern0pt}{\isadigit{2}}\ a\ t{\isadigit{1}}\ t{\isadigit{2}}\ ys{\isacharparenright}{\kern0pt}\isanewline
\ \ \isacommand{hence}\isamarkupfalse%
\ fresh{\isacharunderscore}{\kern0pt}old{\isacharcolon}{\kern0pt}\ {\isachardoublequoteopen}{\isasymforall}\ {\isacharparenleft}{\kern0pt}a{\isacharcomma}{\kern0pt}t{\isacharparenright}{\kern0pt}\ {\isasymin}\ set\ {\isacharparenleft}{\kern0pt}{\isacharparenleft}{\kern0pt}a{\isacharcomma}{\kern0pt}\ Paar\ t{\isadigit{1}}\ t{\isadigit{2}}{\isacharparenright}{\kern0pt}\ {\isacharhash}{\kern0pt}\ ys{\isacharparenright}{\kern0pt}{\isachardot}{\kern0pt}\ nabla{\isadigit{1}}\ {\isasymturnstile}\ a\ {\isasymsharp}\ subst\ s{\isadigit{1}}\ t{\isachardoublequoteclose}\isanewline
\ \ \ \ \isakeyword{and}\ eqs{\isacharcolon}{\kern0pt}\ {\isachardoublequoteopen}{\isasymforall}\ {\isacharparenleft}{\kern0pt}t{\isadigit{1}}{\isacharcomma}{\kern0pt}t{\isadigit{2}}{\isacharparenright}{\kern0pt}\ {\isasymin}\ set\ {\isacharbrackleft}{\kern0pt}{\isacharbrackright}{\kern0pt}{\isachardot}{\kern0pt}\ nabla{\isadigit{1}}\ {\isasymturnstile}\ subst\ s{\isadigit{1}}\ t{\isadigit{1}}\ {\isasymapprox}\ subst\ s{\isadigit{1}}\ t{\isadigit{2}}{\isachardoublequoteclose}\isanewline
\ \ \ \ \isacommand{using}\isamarkupfalse%
\ all{\isacharunderscore}{\kern0pt}solutions{\isacharunderscore}{\kern0pt}def\ \isacommand{by}\isamarkupfalse%
\ simp{\isacharunderscore}{\kern0pt}all\isanewline
\ \ \isacommand{hence}\isamarkupfalse%
\ {\isachardoublequoteopen}nabla{\isadigit{1}}\ {\isasymturnstile}\ a\ {\isasymsharp}\ subst\ s{\isadigit{1}}\ {\isacharparenleft}{\kern0pt}Paar\ t{\isadigit{1}}\ t{\isadigit{2}}{\isacharparenright}{\kern0pt}{\isachardoublequoteclose}\ \isanewline
\ \ \ \ \isacommand{by}\isamarkupfalse%
\ simp\isanewline
\ \ \isacommand{hence}\isamarkupfalse%
\ {\isachardoublequoteopen}nabla{\isadigit{1}}\ {\isasymturnstile}\ a\ {\isasymsharp}\ subst\ s{\isadigit{1}}\ t{\isadigit{1}}{\isachardoublequoteclose}\ {\isachardoublequoteopen}nabla{\isadigit{1}}\ {\isasymturnstile}\ a\ {\isasymsharp}\ subst\ s{\isadigit{1}}\ t{\isadigit{2}}{\isachardoublequoteclose}\isanewline
\ \ \ \ \isacommand{using}\isamarkupfalse%
\ Fresh{\isacharunderscore}{\kern0pt}elims{\isacharparenleft}{\kern0pt}{\isadigit{5}}{\isacharparenright}{\kern0pt}{\isacharbrackleft}{\kern0pt}of\ nabla{\isadigit{1}}\ a\ {\isacartoucheopen}subst\ s{\isadigit{1}}\ t{\isadigit{1}}{\isacartoucheclose}\isanewline
\ \ \ \ \ \ \ \ {\isacartoucheopen}subst\ s{\isadigit{1}}\ t{\isadigit{2}}{\isacartoucheclose}{\isacharbrackright}{\kern0pt}\ \isacommand{by}\isamarkupfalse%
\ auto\isanewline
\ \ \isacommand{with}\isamarkupfalse%
\ fresh{\isacharunderscore}{\kern0pt}old\ \isacommand{have}\isamarkupfalse%
\isanewline
\ \ \ fresh{\isacharcolon}{\kern0pt}\ {\isachardoublequoteopen}{\isasymforall}\ {\isacharparenleft}{\kern0pt}a{\isacharcomma}{\kern0pt}t{\isacharparenright}{\kern0pt}\ {\isasymin}\ set\ {\isacharparenleft}{\kern0pt}{\isacharparenleft}{\kern0pt}a{\isacharcomma}{\kern0pt}\ t{\isadigit{1}}{\isacharparenright}{\kern0pt}\ {\isacharhash}{\kern0pt}\ {\isacharparenleft}{\kern0pt}a{\isacharcomma}{\kern0pt}\ t{\isadigit{2}}{\isacharparenright}{\kern0pt}\ {\isacharhash}{\kern0pt}\ ys{\isacharparenright}{\kern0pt}{\isachardot}{\kern0pt}\ nabla{\isadigit{1}}\ {\isasymturnstile}\ a\ {\isasymsharp}\ subst\ s{\isadigit{1}}\ t{\isachardoublequoteclose}\isanewline
\ \ \ \ \isacommand{by}\isamarkupfalse%
\ auto\isanewline
\ \ \isacommand{from}\isamarkupfalse%
\ fresh\ eqs\isanewline
\ \ \isacommand{show}\isamarkupfalse%
\ {\isacharquery}{\kern0pt}case\isanewline
\ \ \ \ \isacommand{using}\isamarkupfalse%
\ all{\isacharunderscore}{\kern0pt}solutions{\isacharunderscore}{\kern0pt}def\ ext{\isacharunderscore}{\kern0pt}subst{\isacharunderscore}{\kern0pt}def\ \isacommand{by}\isamarkupfalse%
\ simp\isanewline
\isacommand{next}\isamarkupfalse%
\isanewline
\ \ \isacommand{case}\isamarkupfalse%
\ {\isacharparenleft}{\kern0pt}{\isadigit{3}}\ a\ f\ t\ ys{\isacharparenright}{\kern0pt}\isanewline
\ \ \isacommand{hence}\isamarkupfalse%
\ fresh{\isacharunderscore}{\kern0pt}old{\isacharcolon}{\kern0pt}\ {\isachardoublequoteopen}{\isasymforall}\ {\isacharparenleft}{\kern0pt}a{\isacharcomma}{\kern0pt}t{\isacharparenright}{\kern0pt}\ {\isasymin}\ set\ {\isacharparenleft}{\kern0pt}{\isacharparenleft}{\kern0pt}a{\isacharcomma}{\kern0pt}\ Func\ f\ t{\isacharparenright}{\kern0pt}\ {\isacharhash}{\kern0pt}\ ys{\isacharparenright}{\kern0pt}{\isachardot}{\kern0pt}\ nabla{\isadigit{1}}\ {\isasymturnstile}\ a\ {\isasymsharp}\ subst\ s{\isadigit{1}}\ t{\isachardoublequoteclose}\isanewline
\ \ \ \ \isakeyword{and}\ eqs{\isacharcolon}{\kern0pt}\ {\isachardoublequoteopen}{\isasymforall}\ {\isacharparenleft}{\kern0pt}t{\isadigit{1}}{\isacharcomma}{\kern0pt}t{\isadigit{2}}{\isacharparenright}{\kern0pt}\ {\isasymin}\ set\ {\isacharbrackleft}{\kern0pt}{\isacharbrackright}{\kern0pt}{\isachardot}{\kern0pt}\ nabla{\isadigit{1}}\ {\isasymturnstile}\ subst\ s{\isadigit{1}}\ t{\isadigit{1}}\ {\isasymapprox}\ subst\ s{\isadigit{1}}\ t{\isadigit{2}}{\isachardoublequoteclose}\isanewline
\ \ \ \ \isacommand{using}\isamarkupfalse%
\ all{\isacharunderscore}{\kern0pt}solutions{\isacharunderscore}{\kern0pt}def\ \isacommand{by}\isamarkupfalse%
\ simp{\isacharunderscore}{\kern0pt}all\isanewline
\ \ \isacommand{hence}\isamarkupfalse%
\ {\isachardoublequoteopen}nabla{\isadigit{1}}\ {\isasymturnstile}\ a\ {\isasymsharp}\ subst\ s{\isadigit{1}}\ {\isacharparenleft}{\kern0pt}Func\ f\ t{\isacharparenright}{\kern0pt}{\isachardoublequoteclose}\ \isanewline
\ \ \ \ \isacommand{by}\isamarkupfalse%
\ simp\isanewline
\ \ \isacommand{hence}\isamarkupfalse%
\ {\isachardoublequoteopen}nabla{\isadigit{1}}\ {\isasymturnstile}\ a\ {\isasymsharp}\ subst\ s{\isadigit{1}}\ t{\isachardoublequoteclose}\ \isanewline
\ \ \ \ \isacommand{using}\isamarkupfalse%
\ Fresh{\isacharunderscore}{\kern0pt}elims{\isacharparenleft}{\kern0pt}{\isadigit{6}}{\isacharparenright}{\kern0pt}\ \isacommand{by}\isamarkupfalse%
\ auto\isanewline
\ \ \isacommand{with}\isamarkupfalse%
\ fresh{\isacharunderscore}{\kern0pt}old\ \isacommand{have}\isamarkupfalse%
\isanewline
\ \ \ fresh{\isacharcolon}{\kern0pt}\ {\isachardoublequoteopen}{\isasymforall}\ {\isacharparenleft}{\kern0pt}a{\isacharcomma}{\kern0pt}t{\isacharparenright}{\kern0pt}\ {\isasymin}\ set\ {\isacharparenleft}{\kern0pt}{\isacharparenleft}{\kern0pt}a{\isacharcomma}{\kern0pt}\ t{\isacharparenright}{\kern0pt}\ {\isacharhash}{\kern0pt}\ ys{\isacharparenright}{\kern0pt}{\isachardot}{\kern0pt}\ nabla{\isadigit{1}}\ {\isasymturnstile}\ a\ {\isasymsharp}\ subst\ s{\isadigit{1}}\ t{\isachardoublequoteclose}\isanewline
\ \ \ \ \isacommand{by}\isamarkupfalse%
\ auto\isanewline
\ \ \isacommand{from}\isamarkupfalse%
\ fresh\ eqs\isanewline
\ \ \isacommand{show}\isamarkupfalse%
\ {\isacharquery}{\kern0pt}case\isanewline
\ \ \ \ \isacommand{using}\isamarkupfalse%
\ all{\isacharunderscore}{\kern0pt}solutions{\isacharunderscore}{\kern0pt}def\ ext{\isacharunderscore}{\kern0pt}subst{\isacharunderscore}{\kern0pt}def\ \isacommand{by}\isamarkupfalse%
\ simp\isanewline
\isacommand{next}\isamarkupfalse%
\isanewline
\ \ \isacommand{case}\isamarkupfalse%
\ {\isacharparenleft}{\kern0pt}{\isadigit{5}}\ a\ b\ t\ ys{\isacharparenright}{\kern0pt}\isanewline
\ \ \isacommand{hence}\isamarkupfalse%
\ fresh{\isacharunderscore}{\kern0pt}old{\isacharcolon}{\kern0pt}\ {\isachardoublequoteopen}{\isasymforall}\ {\isacharparenleft}{\kern0pt}a{\isacharcomma}{\kern0pt}t{\isacharprime}{\kern0pt}{\isacharparenright}{\kern0pt}\ {\isasymin}\ set\ {\isacharparenleft}{\kern0pt}{\isacharparenleft}{\kern0pt}a{\isacharcomma}{\kern0pt}\ Abst\ b\ t{\isacharparenright}{\kern0pt}\ {\isacharhash}{\kern0pt}\ ys{\isacharparenright}{\kern0pt}{\isachardot}{\kern0pt}\ nabla{\isadigit{1}}\ {\isasymturnstile}\ a\ {\isasymsharp}\ subst\ s{\isadigit{1}}\ t{\isacharprime}{\kern0pt}{\isachardoublequoteclose}\isanewline
\ \ \ \ \isakeyword{and}\ eqs{\isacharcolon}{\kern0pt}\ {\isachardoublequoteopen}{\isasymforall}\ {\isacharparenleft}{\kern0pt}t{\isadigit{1}}{\isacharcomma}{\kern0pt}t{\isadigit{2}}{\isacharparenright}{\kern0pt}\ {\isasymin}\ set\ {\isacharbrackleft}{\kern0pt}{\isacharbrackright}{\kern0pt}{\isachardot}{\kern0pt}\ nabla{\isadigit{1}}\ {\isasymturnstile}\ subst\ s{\isadigit{1}}\ t{\isadigit{1}}\ {\isasymapprox}\ subst\ s{\isadigit{1}}\ t{\isadigit{2}}{\isachardoublequoteclose}\isanewline
\ \ \ \ \isacommand{using}\isamarkupfalse%
\ all{\isacharunderscore}{\kern0pt}solutions{\isacharunderscore}{\kern0pt}def\ \isacommand{by}\isamarkupfalse%
\ simp{\isacharunderscore}{\kern0pt}all\isanewline
\ \ \isacommand{hence}\isamarkupfalse%
\ {\isachardoublequoteopen}nabla{\isadigit{1}}\ {\isasymturnstile}\ a\ {\isasymsharp}\ subst\ s{\isadigit{1}}\ {\isacharparenleft}{\kern0pt}Abst\ b\ t{\isacharparenright}{\kern0pt}{\isachardoublequoteclose}\isanewline
\ \ \ \ \isacommand{by}\isamarkupfalse%
\ simp\isanewline
\ \ \isacommand{hence}\isamarkupfalse%
\ {\isachardoublequoteopen}nabla{\isadigit{1}}\ {\isasymturnstile}\ a\ {\isasymsharp}\ subst\ s{\isadigit{1}}\ t{\isachardoublequoteclose}\isanewline
\ \ \ \ \isacommand{using}\isamarkupfalse%
\ {\isadigit{5}}\ Fresh{\isacharunderscore}{\kern0pt}elims{\isacharparenleft}{\kern0pt}{\isadigit{1}}{\isacharparenright}{\kern0pt}{\isacharbrackleft}{\kern0pt}of\ nabla{\isadigit{1}}\ a\ b\ {\isacartoucheopen}subst\ s{\isadigit{1}}\ t{\isacartoucheclose}{\isacharbrackright}{\kern0pt}\ \isacommand{by}\isamarkupfalse%
\ auto\isanewline
\ \ \isacommand{with}\isamarkupfalse%
\ fresh{\isacharunderscore}{\kern0pt}old\ \isacommand{have}\isamarkupfalse%
\isanewline
\ \ fresh{\isacharcolon}{\kern0pt}\ {\isachardoublequoteopen}{\isasymforall}\ {\isacharparenleft}{\kern0pt}a{\isacharcomma}{\kern0pt}t{\isacharprime}{\kern0pt}{\isacharparenright}{\kern0pt}\ {\isasymin}\ set\ {\isacharparenleft}{\kern0pt}{\isacharparenleft}{\kern0pt}a{\isacharcomma}{\kern0pt}\ t{\isacharparenright}{\kern0pt}\ {\isacharhash}{\kern0pt}\ ys{\isacharparenright}{\kern0pt}{\isachardot}{\kern0pt}\ nabla{\isadigit{1}}\ {\isasymturnstile}\ a\ {\isasymsharp}\ subst\ s{\isadigit{1}}\ t{\isacharprime}{\kern0pt}{\isachardoublequoteclose}\isanewline
\ \ \ \ \isacommand{by}\isamarkupfalse%
\ auto\isanewline
\ \ \isacommand{from}\isamarkupfalse%
\ fresh\ eqs\isanewline
\ \ \isacommand{show}\isamarkupfalse%
\ {\isacharquery}{\kern0pt}case\isanewline
\ \ \ \ \isacommand{using}\isamarkupfalse%
\ all{\isacharunderscore}{\kern0pt}solutions{\isacharunderscore}{\kern0pt}def\ ext{\isacharunderscore}{\kern0pt}subst{\isacharunderscore}{\kern0pt}def\ \isacommand{by}\isamarkupfalse%
\ simp\isanewline
\isacommand{next}\isamarkupfalse%
\isanewline
\ \ \isacommand{case}\isamarkupfalse%
\ {\isacharparenleft}{\kern0pt}{\isadigit{7}}\ b\ pi\ X\ ys{\isacharparenright}{\kern0pt}\isanewline
\ \ \isacommand{hence}\isamarkupfalse%
\ fresh{\isacharunderscore}{\kern0pt}old{\isacharcolon}{\kern0pt}\ {\isachardoublequoteopen}{\isasymforall}\ {\isacharparenleft}{\kern0pt}a{\isacharcomma}{\kern0pt}t{\isacharparenright}{\kern0pt}\ {\isasymin}\ set\ {\isacharparenleft}{\kern0pt}{\isacharparenleft}{\kern0pt}b{\isacharcomma}{\kern0pt}\ Susp\ pi\ X{\isacharparenright}{\kern0pt}\ {\isacharhash}{\kern0pt}\ ys{\isacharparenright}{\kern0pt}{\isachardot}{\kern0pt}\ nabla{\isadigit{1}}\ {\isasymturnstile}\ a\ {\isasymsharp}\ subst\ s{\isadigit{1}}\ t{\isachardoublequoteclose}\isanewline
\ \ \ \ \isakeyword{and}\ eqs{\isacharcolon}{\kern0pt}\ {\isachardoublequoteopen}{\isasymforall}\ {\isacharparenleft}{\kern0pt}t{\isadigit{1}}{\isacharcomma}{\kern0pt}t{\isadigit{2}}{\isacharparenright}{\kern0pt}\ {\isasymin}\ set\ {\isacharbrackleft}{\kern0pt}{\isacharbrackright}{\kern0pt}{\isachardot}{\kern0pt}\ nabla{\isadigit{1}}\ {\isasymturnstile}\ subst\ s{\isadigit{1}}\ t{\isadigit{1}}\ {\isasymapprox}\ subst\ s{\isadigit{1}}\ t{\isadigit{2}}{\isachardoublequoteclose}\isanewline
\ \ \ \ \isacommand{using}\isamarkupfalse%
\ all{\isacharunderscore}{\kern0pt}solutions{\isacharunderscore}{\kern0pt}def\ \isacommand{by}\isamarkupfalse%
\ simp{\isacharunderscore}{\kern0pt}all\isanewline
\ \ \isacommand{hence}\isamarkupfalse%
\ fresh{\isacharcolon}{\kern0pt}\ {\isachardoublequoteopen}{\isasymforall}\ {\isacharparenleft}{\kern0pt}a{\isacharcomma}{\kern0pt}t{\isacharparenright}{\kern0pt}\ {\isasymin}\ set\ ys{\isachardot}{\kern0pt}\ nabla{\isadigit{1}}\ {\isasymturnstile}\ a\ {\isasymsharp}\ subst\ s{\isadigit{1}}\ t{\isachardoublequoteclose}\ \isacommand{by}\isamarkupfalse%
\ simp\isanewline
\ \ \isacommand{from}\isamarkupfalse%
\ fresh{\isacharunderscore}{\kern0pt}old\ \isacommand{have}\isamarkupfalse%
\isanewline
\ \ {\isachardoublequoteopen}nabla{\isadigit{1}}\ {\isasymturnstile}\ b\ {\isasymsharp}\ subst\ s{\isadigit{1}}\ {\isacharparenleft}{\kern0pt}Susp\ pi\ X{\isacharparenright}{\kern0pt}{\isachardoublequoteclose}\ \isacommand{by}\isamarkupfalse%
\ simp\isanewline
\ \ \isacommand{hence}\isamarkupfalse%
\ {\isachardoublequoteopen}nabla{\isadigit{1}}\ {\isasymturnstile}\ swapas\ {\isacharparenleft}{\kern0pt}rev\ pi{\isacharparenright}{\kern0pt}\ b\ {\isasymsharp}\ subst\ s{\isadigit{1}}\ {\isacharparenleft}{\kern0pt}Susp\ {\isacharbrackleft}{\kern0pt}{\isacharbrackright}{\kern0pt}\ X{\isacharparenright}{\kern0pt}{\isachardoublequoteclose}\isanewline
\ \ \ \ \isacommand{using}\isamarkupfalse%
\ fresh{\isacharunderscore}{\kern0pt}swap{\isacharunderscore}{\kern0pt}left\ subst{\isacharunderscore}{\kern0pt}susp\ \isacommand{by}\isamarkupfalse%
\ simp\isanewline
\ \ \isacommand{hence}\isamarkupfalse%
\ ext{\isacharcolon}{\kern0pt}\ {\isachardoublequoteopen}nabla{\isadigit{1}}\ {\isasymTurnstile}\ {\isacharparenleft}{\kern0pt}subst\ s{\isadigit{1}}{\isacharparenright}{\kern0pt}\ {\isacharbraceleft}{\kern0pt}{\isacharparenleft}{\kern0pt}swapas\ {\isacharparenleft}{\kern0pt}rev\ pi{\isacharparenright}{\kern0pt}\ b{\isacharcomma}{\kern0pt}\ X{\isacharparenright}{\kern0pt}{\isacharbraceright}{\kern0pt}{\isachardoublequoteclose}\isanewline
\ \ \ \ \isacommand{using}\isamarkupfalse%
\ ext{\isacharunderscore}{\kern0pt}subst{\isacharunderscore}{\kern0pt}def\ \isacommand{by}\isamarkupfalse%
\ simp\isanewline
\ \ \isacommand{from}\isamarkupfalse%
\ fresh\ eqs\ ext\isanewline
\ \ \ \ \isacommand{show}\isamarkupfalse%
\ {\isacharquery}{\kern0pt}case\ \isacommand{using}\isamarkupfalse%
\ all{\isacharunderscore}{\kern0pt}solutions{\isacharunderscore}{\kern0pt}def\ \isacommand{by}\isamarkupfalse%
\ simp\isanewline
\isacommand{qed}\isamarkupfalse%
\ {\isacharparenleft}{\kern0pt}auto\ simp\ add{\isacharcolon}{\kern0pt}\ all{\isacharunderscore}{\kern0pt}solutions{\isacharunderscore}{\kern0pt}def\ ext{\isacharunderscore}{\kern0pt}subst{\isacharunderscore}{\kern0pt}def{\isacharparenright}{\kern0pt}%
\endisatagproof
{\isafoldproof}%
%
\isadelimproof
\isanewline
%
\endisadelimproof
\isanewline
\isanewline
\isacommand{lemma}\isamarkupfalse%
\ P{\isadigit{1}}{\isacharunderscore}{\kern0pt}from{\isacharunderscore}{\kern0pt}P{\isadigit{2}}{\isacharunderscore}{\kern0pt}cred{\isacharcolon}{\kern0pt}\isanewline
\ \ \isakeyword{assumes}\ {\isachardoublequoteopen}{\isacharparenleft}{\kern0pt}nabla{\isadigit{1}}{\isacharcomma}{\kern0pt}s{\isadigit{1}}{\isacharparenright}{\kern0pt}{\isasymin}U\ P{\isadigit{2}}{\isachardoublequoteclose}\isanewline
\ \ \ \ \isakeyword{and}\ {\isachardoublequoteopen}P{\isadigit{1}}\ {\isasymturnstile}nabla{\isadigit{2}}{\isasymrightarrow}\ P{\isadigit{2}}{\isachardoublequoteclose}\ \isakeyword{and}\ {\isachardoublequoteopen}nabla{\isadigit{3}}{\isasymTurnstile}{\isacharparenleft}{\kern0pt}subst\ s{\isadigit{1}}{\isacharparenright}{\kern0pt}\ nabla{\isadigit{2}}{\isachardoublequoteclose}\isanewline
\ \ \isakeyword{shows}\ {\isachardoublequoteopen}{\isacharparenleft}{\kern0pt}nabla{\isadigit{1}}{\isasymunion}nabla{\isadigit{3}}{\isacharcomma}{\kern0pt}s{\isadigit{1}}{\isacharparenright}{\kern0pt}{\isasymin}\ U\ P{\isadigit{1}}{\isachardoublequoteclose}\ \isanewline
%
\isadelimproof
\ \ %
\endisadelimproof
%
\isatagproof
\isacommand{using}\isamarkupfalse%
\ assms\isanewline
\isacommand{proof}\isamarkupfalse%
{\isacharparenleft}{\kern0pt}induction\ arbitrary{\isacharcolon}{\kern0pt}\ nabla{\isadigit{1}}\ s{\isadigit{1}}\ rule{\isacharcolon}{\kern0pt}\ c{\isacharunderscore}{\kern0pt}red{\isachardot}{\kern0pt}induct{\isacharbrackleft}{\kern0pt}OF\ {\isacartoucheopen}P{\isadigit{1}}\ {\isasymturnstile}nabla{\isadigit{2}}{\isasymrightarrow}\ P{\isadigit{2}}{\isacartoucheclose}{\isacharbrackright}{\kern0pt}{\isacharparenright}{\kern0pt}\isanewline
\ \ \isacommand{case}\isamarkupfalse%
\ {\isacharparenleft}{\kern0pt}{\isadigit{2}}\ a\ t{\isadigit{1}}\ t{\isadigit{2}}\ xs{\isacharparenright}{\kern0pt}\isanewline
\ \ \isacommand{hence}\isamarkupfalse%
\ fresh{\isacharunderscore}{\kern0pt}old{\isacharcolon}{\kern0pt}\isanewline
\ \ \ \ {\isachardoublequoteopen}{\isasymforall}{\isacharparenleft}{\kern0pt}b{\isacharcomma}{\kern0pt}t{\isacharparenright}{\kern0pt}{\isasymin}set\ {\isacharparenleft}{\kern0pt}{\isacharparenleft}{\kern0pt}a{\isacharcomma}{\kern0pt}t{\isadigit{1}}{\isacharparenright}{\kern0pt}{\isacharhash}{\kern0pt}{\isacharparenleft}{\kern0pt}a{\isacharcomma}{\kern0pt}t{\isadigit{2}}{\isacharparenright}{\kern0pt}{\isacharhash}{\kern0pt}xs{\isacharparenright}{\kern0pt}{\isachardot}{\kern0pt}\ nabla{\isadigit{1}}\ {\isasymturnstile}\ b\ {\isasymsharp}\ subst\ s{\isadigit{1}}\ t{\isachardoublequoteclose}\isanewline
\ \ \ \ \isakeyword{and}\ {\isachardoublequoteopen}{\isasymforall}\ {\isacharparenleft}{\kern0pt}t{\isadigit{1}}{\isacharcomma}{\kern0pt}\ t{\isadigit{2}}{\isacharparenright}{\kern0pt}\ {\isasymin}\ set\ {\isacharbrackleft}{\kern0pt}{\isacharbrackright}{\kern0pt}{\isachardot}{\kern0pt}\ nabla{\isadigit{1}}\ {\isasymturnstile}\ subst\ s{\isadigit{1}}\ t{\isadigit{1}}\ {\isasymapprox}\ subst\ s{\isadigit{1}}\ t{\isadigit{2}}{\isachardoublequoteclose}\isanewline
\ \ \ \ \isacommand{using}\isamarkupfalse%
\ all{\isacharunderscore}{\kern0pt}solutions{\isacharunderscore}{\kern0pt}def\ \isacommand{by}\isamarkupfalse%
\ auto\isanewline
\ \ \isacommand{hence}\isamarkupfalse%
\ weak{\isadigit{1}}{\isacharcolon}{\kern0pt}\ {\isachardoublequoteopen}{\isasymforall}{\isacharparenleft}{\kern0pt}b{\isacharcomma}{\kern0pt}t{\isacharparenright}{\kern0pt}{\isasymin}set\ xs{\isachardot}{\kern0pt}\ {\isacharparenleft}{\kern0pt}nabla{\isadigit{1}}\ {\isasymunion}\ nabla{\isadigit{3}}{\isacharparenright}{\kern0pt}\ {\isasymturnstile}\ b\ {\isasymsharp}\ subst\ s{\isadigit{1}}\ t{\isachardoublequoteclose}\isanewline
\ \ \ \ \ \ \ \ {\isachardoublequoteopen}{\isasymforall}\ {\isacharparenleft}{\kern0pt}t{\isadigit{1}}{\isacharcomma}{\kern0pt}\ t{\isadigit{2}}{\isacharparenright}{\kern0pt}\ {\isasymin}\ set\ {\isacharbrackleft}{\kern0pt}{\isacharbrackright}{\kern0pt}{\isachardot}{\kern0pt}\ {\isacharparenleft}{\kern0pt}nabla{\isadigit{1}}\ {\isasymunion}\ nabla{\isadigit{3}}{\isacharparenright}{\kern0pt}\ {\isasymturnstile}\ subst\ s{\isadigit{1}}\ t{\isadigit{1}}\ {\isasymapprox}\ subst\ s{\isadigit{1}}\ t{\isadigit{2}}{\isachardoublequoteclose}\isanewline
\ \ \ \ \isacommand{using}\isamarkupfalse%
\ fresh{\isacharunderscore}{\kern0pt}weak\ \isacommand{by}\isamarkupfalse%
\ auto\isanewline
\ \ \isacommand{from}\isamarkupfalse%
\ fresh{\isacharunderscore}{\kern0pt}old\ \isacommand{have}\isamarkupfalse%
\ {\isachardoublequoteopen}nabla{\isadigit{1}}\ {\isasymturnstile}\ a\ {\isasymsharp}\ subst\ s{\isadigit{1}}\ t{\isadigit{1}}{\isachardoublequoteclose}\ {\isachardoublequoteopen}nabla{\isadigit{1}}\ {\isasymturnstile}\ a\ {\isasymsharp}\ subst\ s{\isadigit{1}}\ t{\isadigit{2}}{\isachardoublequoteclose}\isanewline
\ \ \ \ \isacommand{by}\isamarkupfalse%
\ simp{\isacharplus}{\kern0pt}\isanewline
\ \ \isacommand{hence}\isamarkupfalse%
\ {\isachardoublequoteopen}nabla{\isadigit{1}}\ {\isasymturnstile}\ a\ {\isasymsharp}\ subst\ s{\isadigit{1}}\ {\isacharparenleft}{\kern0pt}Paar\ t{\isadigit{1}}\ t{\isadigit{2}}{\isacharparenright}{\kern0pt}{\isachardoublequoteclose}\isanewline
\ \ \ \ \isacommand{using}\isamarkupfalse%
\ fresh{\isacharunderscore}{\kern0pt}paar\ subst{\isacharunderscore}{\kern0pt}paar\ \isacommand{by}\isamarkupfalse%
\ auto\isanewline
\ \ \isacommand{hence}\isamarkupfalse%
\ {\isachardoublequoteopen}{\isacharparenleft}{\kern0pt}nabla{\isadigit{1}}\ {\isasymunion}\ nabla{\isadigit{3}}{\isacharparenright}{\kern0pt}\ {\isasymturnstile}\ a\ {\isasymsharp}\ subst\ s{\isadigit{1}}\ {\isacharparenleft}{\kern0pt}Paar\ t{\isadigit{1}}\ t{\isadigit{2}}{\isacharparenright}{\kern0pt}{\isachardoublequoteclose}\isanewline
\ \ \ \ \isacommand{using}\isamarkupfalse%
\ fresh{\isacharunderscore}{\kern0pt}weak\ \isacommand{by}\isamarkupfalse%
\ auto\isanewline
\ \ \isacommand{with}\isamarkupfalse%
\ weak{\isadigit{1}}\isanewline
\ \ \isacommand{have}\isamarkupfalse%
\ fresh{\isacharcolon}{\kern0pt}\ {\isachardoublequoteopen}{\isasymforall}{\isacharparenleft}{\kern0pt}b{\isacharcomma}{\kern0pt}t{\isacharparenright}{\kern0pt}{\isasymin}set\ {\isacharparenleft}{\kern0pt}{\isacharparenleft}{\kern0pt}a{\isacharcomma}{\kern0pt}Paar\ t{\isadigit{1}}\ t{\isadigit{2}}{\isacharparenright}{\kern0pt}{\isacharhash}{\kern0pt}xs{\isacharparenright}{\kern0pt}{\isachardot}{\kern0pt}\ {\isacharparenleft}{\kern0pt}nabla{\isadigit{1}}\ {\isasymunion}\ nabla{\isadigit{3}}{\isacharparenright}{\kern0pt}\ {\isasymturnstile}\ b\ {\isasymsharp}\ subst\ s{\isadigit{1}}\ t{\isachardoublequoteclose}\isanewline
\ \ \ \ \isacommand{by}\isamarkupfalse%
\ simp\isanewline
\ \ \isacommand{with}\isamarkupfalse%
\ weak{\isadigit{1}}{\isacharparenleft}{\kern0pt}{\isadigit{2}}{\isacharparenright}{\kern0pt}\ \isacommand{show}\isamarkupfalse%
\ {\isacharquery}{\kern0pt}case\ \isanewline
\ \ \ \ \isacommand{using}\isamarkupfalse%
\ all{\isacharunderscore}{\kern0pt}solutions{\isacharunderscore}{\kern0pt}def\ ext{\isacharunderscore}{\kern0pt}subst{\isacharunderscore}{\kern0pt}def\ \isacommand{by}\isamarkupfalse%
\ simp\isanewline
\isacommand{next}\isamarkupfalse%
\isanewline
\ \ \isacommand{case}\isamarkupfalse%
\ {\isacharparenleft}{\kern0pt}{\isadigit{3}}\ a\ f\ t{\isacharprime}{\kern0pt}\ xs{\isacharparenright}{\kern0pt}\isanewline
\ \ \isacommand{hence}\isamarkupfalse%
\ fresh{\isacharunderscore}{\kern0pt}old{\isacharcolon}{\kern0pt}\isanewline
\ \ \ \ {\isachardoublequoteopen}{\isasymforall}{\isacharparenleft}{\kern0pt}b{\isacharcomma}{\kern0pt}t{\isacharparenright}{\kern0pt}{\isasymin}set\ {\isacharparenleft}{\kern0pt}{\isacharparenleft}{\kern0pt}a{\isacharcomma}{\kern0pt}t{\isacharprime}{\kern0pt}{\isacharparenright}{\kern0pt}{\isacharhash}{\kern0pt}xs{\isacharparenright}{\kern0pt}{\isachardot}{\kern0pt}\ nabla{\isadigit{1}}\ {\isasymturnstile}\ b\ {\isasymsharp}\ subst\ s{\isadigit{1}}\ t{\isachardoublequoteclose}\isanewline
\ \ \ \ \isakeyword{and}\ {\isachardoublequoteopen}{\isasymforall}\ {\isacharparenleft}{\kern0pt}t{\isadigit{1}}{\isacharcomma}{\kern0pt}\ t{\isadigit{2}}{\isacharparenright}{\kern0pt}\ {\isasymin}\ set\ {\isacharbrackleft}{\kern0pt}{\isacharbrackright}{\kern0pt}{\isachardot}{\kern0pt}\ nabla{\isadigit{1}}\ {\isasymturnstile}\ subst\ s{\isadigit{1}}\ t{\isadigit{1}}\ {\isasymapprox}\ subst\ s{\isadigit{1}}\ t{\isadigit{2}}{\isachardoublequoteclose}\isanewline
\ \ \ \ \isacommand{using}\isamarkupfalse%
\ all{\isacharunderscore}{\kern0pt}solutions{\isacharunderscore}{\kern0pt}def\ \isacommand{by}\isamarkupfalse%
\ auto\isanewline
\ \ \isacommand{hence}\isamarkupfalse%
\ weak{\isadigit{1}}{\isacharcolon}{\kern0pt}\ {\isachardoublequoteopen}{\isasymforall}{\isacharparenleft}{\kern0pt}b{\isacharcomma}{\kern0pt}t{\isacharparenright}{\kern0pt}{\isasymin}set\ xs{\isachardot}{\kern0pt}\ {\isacharparenleft}{\kern0pt}nabla{\isadigit{1}}\ {\isasymunion}\ nabla{\isadigit{3}}{\isacharparenright}{\kern0pt}\ {\isasymturnstile}\ b\ {\isasymsharp}\ subst\ s{\isadigit{1}}\ t{\isachardoublequoteclose}\isanewline
\ \ \ \ \ \ \ \ {\isachardoublequoteopen}{\isasymforall}\ {\isacharparenleft}{\kern0pt}t{\isadigit{1}}{\isacharcomma}{\kern0pt}\ t{\isadigit{2}}{\isacharparenright}{\kern0pt}\ {\isasymin}\ set\ {\isacharbrackleft}{\kern0pt}{\isacharbrackright}{\kern0pt}{\isachardot}{\kern0pt}\ {\isacharparenleft}{\kern0pt}nabla{\isadigit{1}}\ {\isasymunion}\ nabla{\isadigit{3}}{\isacharparenright}{\kern0pt}\ {\isasymturnstile}\ subst\ s{\isadigit{1}}\ t{\isadigit{1}}\ {\isasymapprox}\ subst\ s{\isadigit{1}}\ t{\isadigit{2}}{\isachardoublequoteclose}\isanewline
\ \ \ \ \isacommand{using}\isamarkupfalse%
\ fresh{\isacharunderscore}{\kern0pt}weak\ \isacommand{by}\isamarkupfalse%
\ auto\isanewline
\ \ \isacommand{from}\isamarkupfalse%
\ fresh{\isacharunderscore}{\kern0pt}old\ \isacommand{have}\isamarkupfalse%
\ {\isachardoublequoteopen}nabla{\isadigit{1}}\ {\isasymturnstile}\ a\ {\isasymsharp}\ subst\ s{\isadigit{1}}\ t{\isacharprime}{\kern0pt}{\isachardoublequoteclose}\ \isanewline
\ \ \ \ \isacommand{by}\isamarkupfalse%
\ simp\isanewline
\ \ \isacommand{hence}\isamarkupfalse%
\ {\isachardoublequoteopen}nabla{\isadigit{1}}\ {\isasymturnstile}\ a\ {\isasymsharp}\ subst\ s{\isadigit{1}}\ {\isacharparenleft}{\kern0pt}Func\ f\ t{\isacharprime}{\kern0pt}{\isacharparenright}{\kern0pt}{\isachardoublequoteclose}\isanewline
\ \ \ \ \isacommand{using}\isamarkupfalse%
\ fresh{\isacharunderscore}{\kern0pt}paar\ subst{\isacharunderscore}{\kern0pt}paar\ \isacommand{by}\isamarkupfalse%
\ auto\isanewline
\ \ \isacommand{hence}\isamarkupfalse%
\ {\isachardoublequoteopen}{\isacharparenleft}{\kern0pt}nabla{\isadigit{1}}\ {\isasymunion}\ nabla{\isadigit{3}}{\isacharparenright}{\kern0pt}\ {\isasymturnstile}\ a\ {\isasymsharp}\ subst\ s{\isadigit{1}}\ {\isacharparenleft}{\kern0pt}Func\ f\ t{\isacharprime}{\kern0pt}{\isacharparenright}{\kern0pt}{\isachardoublequoteclose}\isanewline
\ \ \ \ \isacommand{using}\isamarkupfalse%
\ fresh{\isacharunderscore}{\kern0pt}weak\ \isacommand{by}\isamarkupfalse%
\ auto\isanewline
\ \ \isacommand{with}\isamarkupfalse%
\ weak{\isadigit{1}}\isanewline
\ \ \isacommand{have}\isamarkupfalse%
\ fresh{\isacharcolon}{\kern0pt}\ {\isachardoublequoteopen}{\isasymforall}{\isacharparenleft}{\kern0pt}b{\isacharcomma}{\kern0pt}t{\isacharparenright}{\kern0pt}{\isasymin}set\ {\isacharparenleft}{\kern0pt}{\isacharparenleft}{\kern0pt}a{\isacharcomma}{\kern0pt}\ Func\ f\ t{\isacharprime}{\kern0pt}{\isacharparenright}{\kern0pt}{\isacharhash}{\kern0pt}xs{\isacharparenright}{\kern0pt}{\isachardot}{\kern0pt}\ {\isacharparenleft}{\kern0pt}nabla{\isadigit{1}}\ {\isasymunion}\ nabla{\isadigit{3}}{\isacharparenright}{\kern0pt}\ {\isasymturnstile}\ b\ {\isasymsharp}\ subst\ s{\isadigit{1}}\ t{\isachardoublequoteclose}\isanewline
\ \ \ \ \isacommand{by}\isamarkupfalse%
\ simp\isanewline
\ \ \isacommand{with}\isamarkupfalse%
\ weak{\isadigit{1}}{\isacharparenleft}{\kern0pt}{\isadigit{2}}{\isacharparenright}{\kern0pt}\ \isacommand{show}\isamarkupfalse%
\ {\isacharquery}{\kern0pt}case\ \isanewline
\ \ \ \ \isacommand{using}\isamarkupfalse%
\ all{\isacharunderscore}{\kern0pt}solutions{\isacharunderscore}{\kern0pt}def\ ext{\isacharunderscore}{\kern0pt}subst{\isacharunderscore}{\kern0pt}def\ \isacommand{by}\isamarkupfalse%
\ simp\isanewline
\isacommand{next}\isamarkupfalse%
\isanewline
\ \ \isacommand{case}\isamarkupfalse%
\ {\isacharparenleft}{\kern0pt}{\isadigit{5}}\ a\ b\ t{\isacharprime}{\kern0pt}\ xs{\isacharparenright}{\kern0pt}\isanewline
\ \ \ \isacommand{hence}\isamarkupfalse%
\ fresh{\isacharunderscore}{\kern0pt}old{\isacharcolon}{\kern0pt}\isanewline
\ \ \ \ {\isachardoublequoteopen}{\isasymforall}{\isacharparenleft}{\kern0pt}c{\isacharcomma}{\kern0pt}t{\isacharparenright}{\kern0pt}{\isasymin}set\ {\isacharparenleft}{\kern0pt}{\isacharparenleft}{\kern0pt}a{\isacharcomma}{\kern0pt}t{\isacharprime}{\kern0pt}{\isacharparenright}{\kern0pt}{\isacharhash}{\kern0pt}xs{\isacharparenright}{\kern0pt}{\isachardot}{\kern0pt}\ nabla{\isadigit{1}}\ {\isasymturnstile}\ c\ {\isasymsharp}\ subst\ s{\isadigit{1}}\ t{\isachardoublequoteclose}\isanewline
\ \ \ \ \isakeyword{and}\ {\isachardoublequoteopen}{\isasymforall}\ {\isacharparenleft}{\kern0pt}t{\isadigit{1}}{\isacharcomma}{\kern0pt}\ t{\isadigit{2}}{\isacharparenright}{\kern0pt}\ {\isasymin}\ set\ {\isacharbrackleft}{\kern0pt}{\isacharbrackright}{\kern0pt}{\isachardot}{\kern0pt}\ nabla{\isadigit{1}}\ {\isasymturnstile}\ subst\ s{\isadigit{1}}\ t{\isadigit{1}}\ {\isasymapprox}\ subst\ s{\isadigit{1}}\ t{\isadigit{2}}{\isachardoublequoteclose}\isanewline
\ \ \ \ \ \isacommand{using}\isamarkupfalse%
\ all{\isacharunderscore}{\kern0pt}solutions{\isacharunderscore}{\kern0pt}def\ \isacommand{by}\isamarkupfalse%
\ auto\isanewline
\ \ \ \isacommand{hence}\isamarkupfalse%
\ weak{\isadigit{1}}{\isacharcolon}{\kern0pt}\ {\isachardoublequoteopen}{\isasymforall}{\isacharparenleft}{\kern0pt}c{\isacharcomma}{\kern0pt}t{\isacharparenright}{\kern0pt}{\isasymin}set\ xs{\isachardot}{\kern0pt}\ {\isacharparenleft}{\kern0pt}nabla{\isadigit{1}}\ {\isasymunion}\ nabla{\isadigit{3}}{\isacharparenright}{\kern0pt}\ {\isasymturnstile}\ c\ {\isasymsharp}\ subst\ s{\isadigit{1}}\ t{\isachardoublequoteclose}\isanewline
\ \ \ \ \ \ {\isachardoublequoteopen}{\isasymforall}\ {\isacharparenleft}{\kern0pt}t{\isadigit{1}}{\isacharcomma}{\kern0pt}\ t{\isadigit{2}}{\isacharparenright}{\kern0pt}\ {\isasymin}\ set\ {\isacharbrackleft}{\kern0pt}{\isacharbrackright}{\kern0pt}{\isachardot}{\kern0pt}\ {\isacharparenleft}{\kern0pt}nabla{\isadigit{1}}\ {\isasymunion}\ nabla{\isadigit{3}}{\isacharparenright}{\kern0pt}\ {\isasymturnstile}\ subst\ s{\isadigit{1}}\ t{\isadigit{1}}\ {\isasymapprox}\ subst\ s{\isadigit{1}}\ t{\isadigit{2}}{\isachardoublequoteclose}\isanewline
\ \ \ \ \ \isacommand{using}\isamarkupfalse%
\ fresh{\isacharunderscore}{\kern0pt}weak\ \isacommand{by}\isamarkupfalse%
\ auto\isanewline
\ \ \ \isacommand{from}\isamarkupfalse%
\ fresh{\isacharunderscore}{\kern0pt}old\ \isacommand{have}\isamarkupfalse%
\ {\isachardoublequoteopen}nabla{\isadigit{1}}\ {\isasymturnstile}\ a\ {\isasymsharp}\ subst\ s{\isadigit{1}}\ t{\isacharprime}{\kern0pt}{\isachardoublequoteclose}\ \isanewline
\ \ \ \ \ \isacommand{by}\isamarkupfalse%
\ simp\isanewline
\ \ \ \isacommand{hence}\isamarkupfalse%
\ {\isachardoublequoteopen}nabla{\isadigit{1}}\ {\isasymturnstile}\ a\ {\isasymsharp}\ subst\ s{\isadigit{1}}\ {\isacharparenleft}{\kern0pt}Abst\ b\ t{\isacharprime}{\kern0pt}{\isacharparenright}{\kern0pt}{\isachardoublequoteclose}\isanewline
\ \ \ \isacommand{using}\isamarkupfalse%
\ fresh{\isacharunderscore}{\kern0pt}abst{\isacharunderscore}{\kern0pt}ab{\isacharbrackleft}{\kern0pt}OF\ {\isacartoucheopen}nabla{\isadigit{1}}\ {\isasymturnstile}\ a\ {\isasymsharp}\ subst\ s{\isadigit{1}}\ t{\isacharprime}{\kern0pt}{\isacartoucheclose}\ {\isadigit{5}}{\isacharparenleft}{\kern0pt}{\isadigit{1}}{\isacharparenright}{\kern0pt}{\isacharbrackright}{\kern0pt}\ subst{\isacharunderscore}{\kern0pt}abst\ \isacommand{by}\isamarkupfalse%
\ auto\isanewline
\ \ \ \isacommand{hence}\isamarkupfalse%
\ {\isachardoublequoteopen}{\isacharparenleft}{\kern0pt}nabla{\isadigit{1}}\ {\isasymunion}\ nabla{\isadigit{3}}{\isacharparenright}{\kern0pt}\ {\isasymturnstile}\ a\ {\isasymsharp}\ subst\ s{\isadigit{1}}\ {\isacharparenleft}{\kern0pt}Abst\ b\ t{\isacharprime}{\kern0pt}{\isacharparenright}{\kern0pt}{\isachardoublequoteclose}\isanewline
\ \ \ \ \ \isacommand{using}\isamarkupfalse%
\ fresh{\isacharunderscore}{\kern0pt}weak\ \isacommand{by}\isamarkupfalse%
\ auto\isanewline
\ \ \ \isacommand{with}\isamarkupfalse%
\ weak{\isadigit{1}}\isanewline
\ \ \ \isacommand{have}\isamarkupfalse%
\ fresh{\isacharcolon}{\kern0pt}\ {\isachardoublequoteopen}{\isasymforall}{\isacharparenleft}{\kern0pt}b{\isacharcomma}{\kern0pt}t{\isacharparenright}{\kern0pt}{\isasymin}set\ {\isacharparenleft}{\kern0pt}{\isacharparenleft}{\kern0pt}a{\isacharcomma}{\kern0pt}\ Abst\ b\ t{\isacharprime}{\kern0pt}{\isacharparenright}{\kern0pt}{\isacharhash}{\kern0pt}xs{\isacharparenright}{\kern0pt}{\isachardot}{\kern0pt}\ {\isacharparenleft}{\kern0pt}nabla{\isadigit{1}}\ {\isasymunion}\ nabla{\isadigit{3}}{\isacharparenright}{\kern0pt}\ {\isasymturnstile}\ b\ {\isasymsharp}\ subst\ s{\isadigit{1}}\ t{\isachardoublequoteclose}\isanewline
\ \ \ \ \ \isacommand{by}\isamarkupfalse%
\ simp\isanewline
\ \ \ \isacommand{with}\isamarkupfalse%
\ weak{\isadigit{1}}{\isacharparenleft}{\kern0pt}{\isadigit{2}}{\isacharparenright}{\kern0pt}\ \isacommand{show}\isamarkupfalse%
\ {\isacharquery}{\kern0pt}case\ \isanewline
\ \ \ \ \ \isacommand{using}\isamarkupfalse%
\ all{\isacharunderscore}{\kern0pt}solutions{\isacharunderscore}{\kern0pt}def\ ext{\isacharunderscore}{\kern0pt}subst{\isacharunderscore}{\kern0pt}def\ \isacommand{by}\isamarkupfalse%
\ simp\isanewline
\isacommand{next}\isamarkupfalse%
\isanewline
\ \ \isacommand{case}\isamarkupfalse%
\ {\isacharparenleft}{\kern0pt}{\isadigit{7}}\ b\ pi\ X\ xs{\isacharparenright}{\kern0pt}\isanewline
\ \ \isacommand{hence}\isamarkupfalse%
\ fresh{\isacharunderscore}{\kern0pt}old{\isacharcolon}{\kern0pt}\ {\isachardoublequoteopen}{\isasymforall}\ {\isacharparenleft}{\kern0pt}a{\isacharcomma}{\kern0pt}t{\isacharparenright}{\kern0pt}\ {\isasymin}\ set\ xs{\isachardot}{\kern0pt}\ nabla{\isadigit{1}}\ {\isasymturnstile}\ a\ {\isasymsharp}\ subst\ s{\isadigit{1}}\ t{\isachardoublequoteclose}\isanewline
\ \ \ \ \isakeyword{and}\ eqs{\isacharcolon}{\kern0pt}\ {\isachardoublequoteopen}{\isasymforall}\ {\isacharparenleft}{\kern0pt}t{\isadigit{1}}{\isacharcomma}{\kern0pt}t{\isadigit{2}}{\isacharparenright}{\kern0pt}\ {\isasymin}\ set\ {\isacharbrackleft}{\kern0pt}{\isacharbrackright}{\kern0pt}{\isachardot}{\kern0pt}\ nabla{\isadigit{1}}\ {\isasymturnstile}\ subst\ s{\isadigit{1}}\ t{\isadigit{1}}\ {\isasymapprox}\ subst\ s{\isadigit{1}}\ t{\isadigit{2}}{\isachardoublequoteclose}\isanewline
\ \ \ \ \isacommand{using}\isamarkupfalse%
\ all{\isacharunderscore}{\kern0pt}solutions{\isacharunderscore}{\kern0pt}def\ \isacommand{by}\isamarkupfalse%
\ simp{\isacharunderscore}{\kern0pt}all\isanewline
\ \ \isacommand{hence}\isamarkupfalse%
\ weak{\isadigit{1}}{\isacharcolon}{\kern0pt}\ {\isachardoublequoteopen}{\isasymforall}{\isacharparenleft}{\kern0pt}a{\isacharcomma}{\kern0pt}t{\isacharparenright}{\kern0pt}{\isasymin}set\ xs{\isachardot}{\kern0pt}\ {\isacharparenleft}{\kern0pt}nabla{\isadigit{1}}\ {\isasymunion}\ nabla{\isadigit{3}}{\isacharparenright}{\kern0pt}\ {\isasymturnstile}\ a\ {\isasymsharp}\ subst\ s{\isadigit{1}}\ t{\isachardoublequoteclose}\isanewline
\ \ \ \ \ \ {\isachardoublequoteopen}{\isasymforall}\ {\isacharparenleft}{\kern0pt}t{\isadigit{1}}{\isacharcomma}{\kern0pt}\ t{\isadigit{2}}{\isacharparenright}{\kern0pt}\ {\isasymin}\ set\ {\isacharbrackleft}{\kern0pt}{\isacharbrackright}{\kern0pt}{\isachardot}{\kern0pt}\ {\isacharparenleft}{\kern0pt}nabla{\isadigit{1}}\ {\isasymunion}\ nabla{\isadigit{3}}{\isacharparenright}{\kern0pt}\ {\isasymturnstile}\ subst\ s{\isadigit{1}}\ t{\isadigit{1}}\ {\isasymapprox}\ subst\ s{\isadigit{1}}\ t{\isadigit{2}}{\isachardoublequoteclose}\isanewline
\ \ \ \ \ \isacommand{using}\isamarkupfalse%
\ fresh{\isacharunderscore}{\kern0pt}weak\ \isacommand{by}\isamarkupfalse%
\ auto\isanewline
\ \ \isacommand{from}\isamarkupfalse%
\ {\isadigit{7}}\ \isacommand{have}\isamarkupfalse%
\ {\isachardoublequoteopen}nabla{\isadigit{3}}\ {\isasymturnstile}\ swapas\ {\isacharparenleft}{\kern0pt}rev\ pi{\isacharparenright}{\kern0pt}\ b\ {\isasymsharp}\ subst\ s{\isadigit{1}}\ {\isacharparenleft}{\kern0pt}Susp\ {\isacharbrackleft}{\kern0pt}{\isacharbrackright}{\kern0pt}\ X{\isacharparenright}{\kern0pt}{\isachardoublequoteclose}\isanewline
\ \ \ \ \isacommand{using}\isamarkupfalse%
\ ext{\isacharunderscore}{\kern0pt}subst{\isacharunderscore}{\kern0pt}def\ \isacommand{by}\isamarkupfalse%
\ auto\isanewline
\ \ \isacommand{hence}\isamarkupfalse%
\ {\isachardoublequoteopen}nabla{\isadigit{3}}\ {\isasymturnstile}\ b\ {\isasymsharp}\ subst\ s{\isadigit{1}}\ {\isacharparenleft}{\kern0pt}Susp\ pi\ X{\isacharparenright}{\kern0pt}{\isachardoublequoteclose}\isanewline
\ \ \ \ \isacommand{using}\isamarkupfalse%
\ fresh{\isacharunderscore}{\kern0pt}swap{\isacharunderscore}{\kern0pt}right\ subst{\isacharunderscore}{\kern0pt}susp\ \isacommand{by}\isamarkupfalse%
\ auto\isanewline
\ \ \isacommand{hence}\isamarkupfalse%
\ weak{\isadigit{2}}{\isacharcolon}{\kern0pt}\ {\isachardoublequoteopen}{\isacharparenleft}{\kern0pt}nabla{\isadigit{1}}\ {\isasymunion}\ nabla{\isadigit{3}}{\isacharparenright}{\kern0pt}\ {\isasymturnstile}\ b\ {\isasymsharp}\ subst\ s{\isadigit{1}}\ {\isacharparenleft}{\kern0pt}Susp\ pi\ X{\isacharparenright}{\kern0pt}{\isachardoublequoteclose}\isanewline
\ \ \ \ \isacommand{using}\isamarkupfalse%
\ fresh{\isacharunderscore}{\kern0pt}weak{\isacharbrackleft}{\kern0pt}of\ nabla{\isadigit{3}}\ b\ {\isacharunderscore}{\kern0pt}\ nabla{\isadigit{1}}{\isacharbrackright}{\kern0pt}\ Un{\isacharunderscore}{\kern0pt}commute{\isacharbrackleft}{\kern0pt}of\ nabla{\isadigit{3}}\ nabla{\isadigit{1}}{\isacharbrackright}{\kern0pt}\ \isacommand{by}\isamarkupfalse%
\ auto\isanewline
\ \ \isacommand{with}\isamarkupfalse%
\ weak{\isadigit{1}}{\isacharparenleft}{\kern0pt}{\isadigit{1}}{\isacharparenright}{\kern0pt}\ \isanewline
\ \ \isacommand{have}\isamarkupfalse%
\ {\isachardoublequoteopen}{\isasymforall}{\isacharparenleft}{\kern0pt}a{\isacharcomma}{\kern0pt}t{\isacharparenright}{\kern0pt}{\isasymin}set\ {\isacharparenleft}{\kern0pt}{\isacharparenleft}{\kern0pt}b{\isacharcomma}{\kern0pt}Susp\ pi\ X{\isacharparenright}{\kern0pt}{\isacharhash}{\kern0pt}xs{\isacharparenright}{\kern0pt}{\isachardot}{\kern0pt}\ {\isacharparenleft}{\kern0pt}nabla{\isadigit{1}}\ {\isasymunion}\ nabla{\isadigit{3}}{\isacharparenright}{\kern0pt}\ {\isasymturnstile}\ a\ {\isasymsharp}\ subst\ s{\isadigit{1}}\ t{\isachardoublequoteclose}\isanewline
\ \ \ \ \isacommand{by}\isamarkupfalse%
\ simp\isanewline
\ \ \isacommand{with}\isamarkupfalse%
\ weak{\isadigit{2}}\ \isacommand{show}\isamarkupfalse%
\ {\isacharquery}{\kern0pt}case\ \isanewline
\ \ \ \ \isacommand{using}\isamarkupfalse%
\ ext{\isacharunderscore}{\kern0pt}subst{\isacharunderscore}{\kern0pt}def\ all{\isacharunderscore}{\kern0pt}solutions{\isacharunderscore}{\kern0pt}def\ \isacommand{by}\isamarkupfalse%
\ simp\isanewline
\isacommand{qed}\isamarkupfalse%
\ {\isacharparenleft}{\kern0pt}auto\ simp\ add{\isacharcolon}{\kern0pt}\ ext{\isacharunderscore}{\kern0pt}subst{\isacharunderscore}{\kern0pt}def\ all{\isacharunderscore}{\kern0pt}solutions{\isacharunderscore}{\kern0pt}def\ fresh{\isacharunderscore}{\kern0pt}weak{\isacharparenright}{\kern0pt}%
\endisatagproof
{\isafoldproof}%
%
\isadelimproof
\isanewline
%
\endisadelimproof
\isanewline
\isanewline
\isanewline
\isanewline
\isacommand{lemma}\isamarkupfalse%
\ P{\isadigit{1}}{\isacharunderscore}{\kern0pt}to{\isacharunderscore}{\kern0pt}P{\isadigit{2}}{\isacharunderscore}{\kern0pt}red{\isacharunderscore}{\kern0pt}plus{\isadigit{1}}{\isacharcolon}{\kern0pt}\ \isanewline
\ \ \isakeyword{assumes}\ {\isachardoublequoteopen}P{\isadigit{1}}\ {\isasymTurnstile}\ {\isacharparenleft}{\kern0pt}nabla{\isacharcomma}{\kern0pt}s{\isacharparenright}{\kern0pt}\ {\isasymRightarrow}\ P{\isadigit{2}}{\isachardoublequoteclose}\isanewline
\ \ \isakeyword{and}\ {\isachardoublequoteopen}{\isacharparenleft}{\kern0pt}nabla{\isadigit{1}}{\isacharcomma}{\kern0pt}s{\isadigit{1}}{\isacharparenright}{\kern0pt}\ {\isasymin}\ U\ P{\isadigit{1}}{\isachardoublequoteclose}\isanewline
\ \ \isakeyword{shows}\ {\isachardoublequoteopen}{\isacharparenleft}{\kern0pt}nabla{\isadigit{1}}{\isacharcomma}{\kern0pt}s{\isadigit{1}}{\isacharparenright}{\kern0pt}\ {\isasymin}\ U\ P{\isadigit{2}}{\isachardoublequoteclose}\isanewline
%
\isadelimproof
\ \ %
\endisadelimproof
%
\isatagproof
\isacommand{using}\isamarkupfalse%
\ assms\isanewline
\isacommand{proof}\isamarkupfalse%
{\isacharparenleft}{\kern0pt}induction\ P{\isadigit{1}}{\isasymequiv}{\isachardoublequoteopen}P{\isadigit{1}}{\isachardoublequoteclose}\ nabla{\isasymequiv}{\isachardoublequoteopen}{\isacharparenleft}{\kern0pt}nabla{\isacharcomma}{\kern0pt}\ s{\isacharparenright}{\kern0pt}{\isachardoublequoteclose}\ P{\isadigit{2}}{\isasymequiv}{\isachardoublequoteopen}P{\isadigit{2}}{\isachardoublequoteclose}\ arbitrary{\isacharcolon}{\kern0pt}\ s\ nabla{\isadigit{1}}\ nabla\ s{\isadigit{1}}\ P{\isadigit{1}}\ P{\isadigit{2}}\ rule{\isacharcolon}{\kern0pt}\ red{\isacharunderscore}{\kern0pt}plus{\isachardot}{\kern0pt}induct{\isacharparenright}{\kern0pt}\isanewline
\ \ \isacommand{case}\isamarkupfalse%
\ {\isacharparenleft}{\kern0pt}sred{\isacharunderscore}{\kern0pt}single\ P{\isadigit{1}}\ s\ P{\isadigit{2}}\ nabla{\isadigit{1}}\ s{\isadigit{1}}{\isacharparenright}{\kern0pt}\isanewline
\ \ \isacommand{have}\isamarkupfalse%
\ {\isachardoublequoteopen}P{\isadigit{1}}\ {\isasymturnstile}\ s\ {\isasymleadsto}\ P{\isadigit{2}}{\isachardoublequoteclose}\ {\isachardoublequoteopen}{\isacharparenleft}{\kern0pt}nabla{\isadigit{1}}{\isacharcomma}{\kern0pt}\ s{\isadigit{1}}{\isacharparenright}{\kern0pt}\ {\isasymin}\ U\ P{\isadigit{1}}{\isachardoublequoteclose}\ \isacommand{by}\isamarkupfalse%
\ fact{\isacharplus}{\kern0pt}\isanewline
\ \ \isacommand{then}\isamarkupfalse%
\ \isacommand{show}\isamarkupfalse%
\ {\isachardoublequoteopen}{\isacharparenleft}{\kern0pt}nabla{\isadigit{1}}{\isacharcomma}{\kern0pt}\ s{\isadigit{1}}{\isacharparenright}{\kern0pt}\ {\isasymin}\ U\ P{\isadigit{2}}{\isachardoublequoteclose}\isanewline
\ \ \ \ \isacommand{using}\isamarkupfalse%
\ P{\isadigit{1}}{\isacharunderscore}{\kern0pt}to{\isacharunderscore}{\kern0pt}P{\isadigit{2}}{\isacharunderscore}{\kern0pt}sred\ \isacommand{by}\isamarkupfalse%
\ blast\isanewline
\isacommand{next}\isamarkupfalse%
\isanewline
\ \ \isacommand{case}\isamarkupfalse%
\ {\isacharparenleft}{\kern0pt}cred{\isacharunderscore}{\kern0pt}single\ P{\isadigit{1}}\ nabla{\isadigit{1}}\ P{\isadigit{2}}\ nabla{\isadigit{2}}\ s{\isadigit{1}}{\isacharparenright}{\kern0pt}\isanewline
\ \ \isacommand{have}\isamarkupfalse%
\ {\isachardoublequoteopen}P{\isadigit{1}}\ {\isasymturnstile}\ nabla{\isadigit{1}}\ {\isasymrightarrow}P{\isadigit{2}}{\isachardoublequoteclose}\ {\isachardoublequoteopen}{\isacharparenleft}{\kern0pt}nabla{\isadigit{2}}{\isacharcomma}{\kern0pt}\ s{\isadigit{1}}{\isacharparenright}{\kern0pt}\ {\isasymin}\ U\ P{\isadigit{1}}{\isachardoublequoteclose}\ \isacommand{by}\isamarkupfalse%
\ fact{\isacharplus}{\kern0pt}\isanewline
\ \ \isacommand{then}\isamarkupfalse%
\ \isacommand{show}\isamarkupfalse%
\ {\isachardoublequoteopen}{\isacharparenleft}{\kern0pt}nabla{\isadigit{2}}{\isacharcomma}{\kern0pt}\ s{\isadigit{1}}{\isacharparenright}{\kern0pt}\ {\isasymin}\ U\ P{\isadigit{2}}{\isachardoublequoteclose}\isanewline
\ \ \ \ \isacommand{using}\isamarkupfalse%
\ P{\isadigit{1}}{\isacharunderscore}{\kern0pt}to{\isacharunderscore}{\kern0pt}P{\isadigit{2}}{\isacharunderscore}{\kern0pt}cred\ \isacommand{by}\isamarkupfalse%
\ blast\isanewline
\isacommand{next}\isamarkupfalse%
\isanewline
\ \ \isacommand{case}\isamarkupfalse%
\ {\isacharparenleft}{\kern0pt}sred{\isacharunderscore}{\kern0pt}step\ P{\isadigit{1}}\ s\ P{\isadigit{2}}\ nabla{\isadigit{2}}\ s{\isadigit{2}}\ P{\isadigit{3}}\ nabla{\isadigit{1}}\ s{\isadigit{1}}{\isacharparenright}{\kern0pt}\isanewline
\ \ \isacommand{have}\isamarkupfalse%
\ {\isachardoublequoteopen}P{\isadigit{1}}\ {\isasymturnstile}\ s\ {\isasymleadsto}\ P{\isadigit{2}}{\isachardoublequoteclose}\ \isakeyword{and}\ {\isachardoublequoteopen}{\isacharparenleft}{\kern0pt}nabla{\isadigit{1}}{\isacharcomma}{\kern0pt}\ s{\isadigit{1}}{\isacharparenright}{\kern0pt}\ {\isasymin}\ U\ P{\isadigit{1}}{\isachardoublequoteclose}\ \isacommand{by}\isamarkupfalse%
\ fact{\isacharplus}{\kern0pt}\isanewline
\ \ \isacommand{then}\isamarkupfalse%
\ \isacommand{have}\isamarkupfalse%
\ {\isachardoublequoteopen}{\isacharparenleft}{\kern0pt}nabla{\isadigit{1}}{\isacharcomma}{\kern0pt}\ s{\isadigit{1}}{\isacharparenright}{\kern0pt}\ {\isasymin}\ U\ P{\isadigit{2}}{\isachardoublequoteclose}\ \isacommand{using}\isamarkupfalse%
\ P{\isadigit{1}}{\isacharunderscore}{\kern0pt}to{\isacharunderscore}{\kern0pt}P{\isadigit{2}}{\isacharunderscore}{\kern0pt}sred\ \isacommand{by}\isamarkupfalse%
\ blast\ \isanewline
\ \ \isacommand{moreover}\isamarkupfalse%
\ \isanewline
\ \ \isacommand{have}\isamarkupfalse%
\ IH{\isacharcolon}{\kern0pt}\ {\isachardoublequoteopen}{\isacharparenleft}{\kern0pt}nabla{\isadigit{1}}{\isacharcomma}{\kern0pt}\ s{\isadigit{1}}{\isacharparenright}{\kern0pt}\ {\isasymin}\ U\ P{\isadigit{2}}\ {\isasymLongrightarrow}\ {\isacharparenleft}{\kern0pt}nabla{\isadigit{1}}{\isacharcomma}{\kern0pt}\ s{\isadigit{1}}{\isacharparenright}{\kern0pt}\ {\isasymin}\ U\ P{\isadigit{3}}{\isachardoublequoteclose}\ \isacommand{by}\isamarkupfalse%
\ fact\isanewline
\ \ \isacommand{ultimately}\isamarkupfalse%
\ \isanewline
\ \ \isacommand{show}\isamarkupfalse%
\ {\isachardoublequoteopen}{\isacharparenleft}{\kern0pt}nabla{\isadigit{1}}{\isacharcomma}{\kern0pt}\ s{\isadigit{1}}{\isacharparenright}{\kern0pt}\ {\isasymin}\ U\ P{\isadigit{3}}{\isachardoublequoteclose}\ \isacommand{by}\isamarkupfalse%
\ simp\ \isanewline
\isacommand{next}\isamarkupfalse%
\isanewline
\ \ \isacommand{case}\isamarkupfalse%
\ {\isacharparenleft}{\kern0pt}cred{\isacharunderscore}{\kern0pt}step\ P{\isadigit{1}}\ nabla{\isadigit{1}}\ P{\isadigit{2}}\ nabla{\isadigit{2}}\ P{\isadigit{3}}\ nabla{\isadigit{3}}\ s{\isadigit{1}}{\isacharparenright}{\kern0pt}\isanewline
\ \ \isacommand{have}\isamarkupfalse%
\ {\isachardoublequoteopen}P{\isadigit{1}}\ {\isasymturnstile}\ nabla{\isadigit{1}}\ {\isasymrightarrow}\ P{\isadigit{2}}{\isachardoublequoteclose}\ \isakeyword{and}\ {\isachardoublequoteopen}{\isacharparenleft}{\kern0pt}nabla{\isadigit{3}}{\isacharcomma}{\kern0pt}\ s{\isadigit{1}}{\isacharparenright}{\kern0pt}\ {\isasymin}\ U\ P{\isadigit{1}}{\isachardoublequoteclose}\ \isacommand{by}\isamarkupfalse%
\ fact{\isacharplus}{\kern0pt}\isanewline
\ \ \isacommand{then}\isamarkupfalse%
\ \isacommand{have}\isamarkupfalse%
\ {\isachardoublequoteopen}{\isacharparenleft}{\kern0pt}nabla{\isadigit{3}}{\isacharcomma}{\kern0pt}\ s{\isadigit{1}}{\isacharparenright}{\kern0pt}\ {\isasymin}\ U\ P{\isadigit{2}}{\isachardoublequoteclose}\ \isacommand{using}\isamarkupfalse%
\ P{\isadigit{1}}{\isacharunderscore}{\kern0pt}to{\isacharunderscore}{\kern0pt}P{\isadigit{2}}{\isacharunderscore}{\kern0pt}cred\ \isacommand{by}\isamarkupfalse%
\ blast\ \isanewline
\ \ \isacommand{moreover}\isamarkupfalse%
\ \isanewline
\ \ \isacommand{have}\isamarkupfalse%
\ IH{\isacharcolon}{\kern0pt}\ {\isachardoublequoteopen}{\isacharparenleft}{\kern0pt}nabla{\isadigit{3}}{\isacharcomma}{\kern0pt}\ s{\isadigit{1}}{\isacharparenright}{\kern0pt}\ {\isasymin}\ U\ P{\isadigit{2}}\ {\isasymLongrightarrow}\ {\isacharparenleft}{\kern0pt}nabla{\isadigit{3}}{\isacharcomma}{\kern0pt}\ s{\isadigit{1}}{\isacharparenright}{\kern0pt}\ {\isasymin}\ U\ P{\isadigit{3}}{\isachardoublequoteclose}\ \isacommand{by}\isamarkupfalse%
\ fact\isanewline
\ \ \isacommand{ultimately}\isamarkupfalse%
\ \isacommand{show}\isamarkupfalse%
\ {\isachardoublequoteopen}{\isacharparenleft}{\kern0pt}nabla{\isadigit{3}}{\isacharcomma}{\kern0pt}\ s{\isadigit{1}}{\isacharparenright}{\kern0pt}\ {\isasymin}\ U\ P{\isadigit{3}}{\isachardoublequoteclose}\ \isacommand{by}\isamarkupfalse%
\ simp\isanewline
\isacommand{qed}\isamarkupfalse%
%
\endisatagproof
{\isafoldproof}%
%
\isadelimproof
\isanewline
%
\endisadelimproof
\isanewline
\isacommand{lemma}\isamarkupfalse%
\ P{\isadigit{1}}{\isacharunderscore}{\kern0pt}to{\isacharunderscore}{\kern0pt}P{\isadigit{2}}{\isacharunderscore}{\kern0pt}red{\isacharunderscore}{\kern0pt}plus{\isadigit{3}}{\isacharcolon}{\kern0pt}\ \isanewline
\ \ \isakeyword{assumes}\ {\isachardoublequoteopen}P{\isadigit{1}}\ {\isasymTurnstile}\ {\isacharparenleft}{\kern0pt}nabla{\isacharcomma}{\kern0pt}s{\isacharparenright}{\kern0pt}\ {\isasymRightarrow}\ P{\isadigit{2}}{\isachardoublequoteclose}\isanewline
\ \ \isakeyword{and}\ {\isachardoublequoteopen}{\isacharparenleft}{\kern0pt}nabla{\isadigit{1}}{\isacharcomma}{\kern0pt}s{\isadigit{1}}{\isacharparenright}{\kern0pt}{\isasymin}\ U\ P{\isadigit{1}}{\isachardoublequoteclose}\isanewline
\ \ \isakeyword{shows}\ {\isachardoublequoteopen}nabla{\isadigit{1}}{\isasymTurnstile}{\isacharparenleft}{\kern0pt}subst\ s{\isadigit{1}}{\isacharparenright}{\kern0pt}\ nabla{\isachardoublequoteclose}\isanewline
%
\isadelimproof
\ \ %
\endisadelimproof
%
\isatagproof
\isacommand{using}\isamarkupfalse%
\ assms\isanewline
\isacommand{proof}\isamarkupfalse%
{\isacharparenleft}{\kern0pt}induction\ P{\isadigit{1}}{\isasymequiv}{\isachardoublequoteopen}P{\isadigit{1}}{\isachardoublequoteclose}\ nabla{\isasymequiv}{\isachardoublequoteopen}{\isacharparenleft}{\kern0pt}nabla{\isacharcomma}{\kern0pt}\ s{\isacharparenright}{\kern0pt}{\isachardoublequoteclose}\ P{\isadigit{2}}{\isasymequiv}{\isachardoublequoteopen}P{\isadigit{2}}{\isachardoublequoteclose}\ arbitrary{\isacharcolon}{\kern0pt}\ s\ nabla{\isadigit{1}}\ nabla\ s{\isadigit{1}}\ P{\isadigit{1}}\ P{\isadigit{2}}\ rule{\isacharcolon}{\kern0pt}\ red{\isacharunderscore}{\kern0pt}plus{\isachardot}{\kern0pt}induct{\isacharparenright}{\kern0pt}\isanewline
\ \ \isacommand{case}\isamarkupfalse%
\ {\isacharparenleft}{\kern0pt}sred{\isacharunderscore}{\kern0pt}single\ P{\isadigit{1}}\ s\ P{\isadigit{2}}\ nabla{\isadigit{1}}\ s{\isadigit{1}}{\isacharparenright}{\kern0pt}\isanewline
\ \ \isacommand{then}\isamarkupfalse%
\ \isacommand{show}\isamarkupfalse%
\ {\isachardoublequoteopen}nabla{\isadigit{1}}\ {\isasymTurnstile}\ {\isacharparenleft}{\kern0pt}subst\ s{\isadigit{1}}{\isacharparenright}{\kern0pt}\ {\isacharbraceleft}{\kern0pt}{\isacharbraceright}{\kern0pt}{\isachardoublequoteclose}\ \isacommand{by}\isamarkupfalse%
\ {\isacharparenleft}{\kern0pt}simp\ add{\isacharcolon}{\kern0pt}\ ext{\isacharunderscore}{\kern0pt}subst{\isacharunderscore}{\kern0pt}def{\isacharparenright}{\kern0pt}\ \isanewline
\isacommand{next}\isamarkupfalse%
\isanewline
\ \ \isacommand{case}\isamarkupfalse%
\ {\isacharparenleft}{\kern0pt}cred{\isacharunderscore}{\kern0pt}single\ P{\isadigit{1}}\ nabla{\isadigit{1}}\ P{\isadigit{2}}\ nabla{\isadigit{2}}\ s{\isadigit{1}}{\isacharparenright}{\kern0pt}\isanewline
\ \ \isacommand{have}\isamarkupfalse%
\ {\isachardoublequoteopen}P{\isadigit{1}}\ \ {\isasymturnstile}\ nabla{\isadigit{1}}\ {\isasymrightarrow}\ P{\isadigit{2}}{\isachardoublequoteclose}\ {\isachardoublequoteopen}{\isacharparenleft}{\kern0pt}nabla{\isadigit{2}}{\isacharcomma}{\kern0pt}\ s{\isadigit{1}}{\isacharparenright}{\kern0pt}\ {\isasymin}\ U\ P{\isadigit{1}}{\isachardoublequoteclose}\ \isacommand{by}\isamarkupfalse%
\ fact{\isacharplus}{\kern0pt}\isanewline
\ \ \isacommand{then}\isamarkupfalse%
\ \isacommand{show}\isamarkupfalse%
\ {\isachardoublequoteopen}nabla{\isadigit{2}}\ {\isasymTurnstile}\ {\isacharparenleft}{\kern0pt}subst\ s{\isadigit{1}}{\isacharparenright}{\kern0pt}\ nabla{\isadigit{1}}{\isachardoublequoteclose}\ \isacommand{using}\isamarkupfalse%
\ P{\isadigit{1}}{\isacharunderscore}{\kern0pt}to{\isacharunderscore}{\kern0pt}P{\isadigit{2}}{\isacharunderscore}{\kern0pt}cred\ \isacommand{by}\isamarkupfalse%
\ blast\isanewline
\isacommand{next}\isamarkupfalse%
\isanewline
\ \ \isacommand{case}\isamarkupfalse%
\ {\isacharparenleft}{\kern0pt}sred{\isacharunderscore}{\kern0pt}step\ P{\isadigit{1}}\ s\ P{\isadigit{2}}\ nabla{\isadigit{2}}\ s{\isadigit{2}}\ P{\isadigit{3}}\ nabla{\isadigit{1}}\ s{\isadigit{1}}{\isacharparenright}{\kern0pt}\isanewline
\ \ \isacommand{have}\isamarkupfalse%
\ {\isachardoublequoteopen}P{\isadigit{1}}\ \ {\isasymturnstile}\ s\ {\isasymleadsto}\ P{\isadigit{2}}{\isachardoublequoteclose}\ {\isachardoublequoteopen}{\isacharparenleft}{\kern0pt}nabla{\isadigit{1}}{\isacharcomma}{\kern0pt}\ s{\isadigit{1}}{\isacharparenright}{\kern0pt}{\isasymin}\ U\ P{\isadigit{1}}{\isachardoublequoteclose}\ \isacommand{by}\isamarkupfalse%
\ fact{\isacharplus}{\kern0pt}\ \isanewline
\ \ \isacommand{then}\isamarkupfalse%
\ \isacommand{have}\isamarkupfalse%
\ {\isachardoublequoteopen}{\isacharparenleft}{\kern0pt}nabla{\isadigit{1}}{\isacharcomma}{\kern0pt}\ s{\isadigit{1}}{\isacharparenright}{\kern0pt}{\isasymin}\ U\ P{\isadigit{2}}{\isachardoublequoteclose}\ \isacommand{using}\isamarkupfalse%
\ P{\isadigit{1}}{\isacharunderscore}{\kern0pt}to{\isacharunderscore}{\kern0pt}P{\isadigit{2}}{\isacharunderscore}{\kern0pt}sred\ \isacommand{by}\isamarkupfalse%
\ blast\ \isanewline
\ \ \isacommand{moreover}\isamarkupfalse%
\ \isacommand{have}\isamarkupfalse%
\ IH{\isacharcolon}{\kern0pt}\ {\isachardoublequoteopen}{\isacharparenleft}{\kern0pt}nabla{\isadigit{1}}{\isacharcomma}{\kern0pt}\ s{\isadigit{1}}{\isacharparenright}{\kern0pt}{\isasymin}\ U\ P{\isadigit{2}}\ {\isasymLongrightarrow}\ \ nabla{\isadigit{1}}\ {\isasymTurnstile}\ {\isacharparenleft}{\kern0pt}subst\ s{\isadigit{1}}{\isacharparenright}{\kern0pt}\ nabla{\isadigit{2}}{\isachardoublequoteclose}\ \isacommand{by}\isamarkupfalse%
\ fact\ \isanewline
\ \ \isacommand{ultimately}\isamarkupfalse%
\ \isacommand{show}\isamarkupfalse%
\ {\isachardoublequoteopen}nabla{\isadigit{1}}\ {\isasymTurnstile}\ {\isacharparenleft}{\kern0pt}subst\ s{\isadigit{1}}{\isacharparenright}{\kern0pt}\ nabla{\isadigit{2}}{\isachardoublequoteclose}\ \isacommand{by}\isamarkupfalse%
\ simp\ \isanewline
\isacommand{next}\isamarkupfalse%
\isanewline
\ \ \isacommand{case}\isamarkupfalse%
\ {\isacharparenleft}{\kern0pt}cred{\isacharunderscore}{\kern0pt}step\ P{\isadigit{1}}\ nabla{\isadigit{1}}\ P{\isadigit{2}}\ nabla{\isadigit{2}}\ P{\isadigit{3}}\ nabla{\isadigit{3}}\ s{\isadigit{1}}{\isacharparenright}{\kern0pt}\isanewline
\ \ \isacommand{have}\isamarkupfalse%
\ IH{\isacharcolon}{\kern0pt}\ {\isachardoublequoteopen}{\isacharparenleft}{\kern0pt}nabla{\isadigit{3}}{\isacharcomma}{\kern0pt}\ s{\isadigit{1}}{\isacharparenright}{\kern0pt}\ {\isasymin}\ U\ P{\isadigit{2}}\ {\isasymLongrightarrow}\ nabla{\isadigit{3}}\ {\isasymTurnstile}\ {\isacharparenleft}{\kern0pt}subst\ s{\isadigit{1}}{\isacharparenright}{\kern0pt}\ nabla{\isadigit{2}}{\isachardoublequoteclose}\ \isacommand{by}\isamarkupfalse%
\ fact\isanewline
\ \ \isacommand{have}\isamarkupfalse%
\ as{\isacharcolon}{\kern0pt}\ {\isachardoublequoteopen}P{\isadigit{1}}\ {\isasymturnstile}\ nabla{\isadigit{1}}\ {\isasymrightarrow}\ P{\isadigit{2}}{\isachardoublequoteclose}\ {\isachardoublequoteopen}{\isacharparenleft}{\kern0pt}nabla{\isadigit{3}}{\isacharcomma}{\kern0pt}\ s{\isadigit{1}}{\isacharparenright}{\kern0pt}\ {\isasymin}\ U\ P{\isadigit{1}}{\isachardoublequoteclose}\ \isacommand{by}\isamarkupfalse%
\ fact{\isacharplus}{\kern0pt}\ \isanewline
\ \ \isanewline
\ \ \isacommand{from}\isamarkupfalse%
\ as\ \isacommand{have}\isamarkupfalse%
\ {\isachardoublequoteopen}nabla{\isadigit{3}}\ {\isasymTurnstile}\ {\isacharparenleft}{\kern0pt}subst\ s{\isadigit{1}}{\isacharparenright}{\kern0pt}\ nabla{\isadigit{1}}{\isachardoublequoteclose}\ \isacommand{using}\isamarkupfalse%
\ P{\isadigit{1}}{\isacharunderscore}{\kern0pt}to{\isacharunderscore}{\kern0pt}P{\isadigit{2}}{\isacharunderscore}{\kern0pt}cred\ \isacommand{by}\isamarkupfalse%
\ blast\ \ \ \ \isanewline
\ \ \isacommand{moreover}\isamarkupfalse%
\ \isanewline
\ \ \isacommand{from}\isamarkupfalse%
\ as\ \isacommand{have}\isamarkupfalse%
\ {\isachardoublequoteopen}{\isacharparenleft}{\kern0pt}nabla{\isadigit{3}}{\isacharcomma}{\kern0pt}\ s{\isadigit{1}}{\isacharparenright}{\kern0pt}\ {\isasymin}\ U\ P{\isadigit{2}}{\isachardoublequoteclose}\ \isacommand{using}\isamarkupfalse%
\ P{\isadigit{1}}{\isacharunderscore}{\kern0pt}to{\isacharunderscore}{\kern0pt}P{\isadigit{2}}{\isacharunderscore}{\kern0pt}cred\ \isacommand{by}\isamarkupfalse%
\ blast\ \isanewline
\ \ \isacommand{then}\isamarkupfalse%
\ \isacommand{have}\isamarkupfalse%
\ {\isachardoublequoteopen}nabla{\isadigit{3}}\ {\isasymTurnstile}\ {\isacharparenleft}{\kern0pt}subst\ s{\isadigit{1}}{\isacharparenright}{\kern0pt}\ nabla{\isadigit{2}}{\isachardoublequoteclose}\ \isacommand{using}\isamarkupfalse%
\ IH\ \isacommand{by}\isamarkupfalse%
\ blast\isanewline
\ \ \isacommand{ultimately}\isamarkupfalse%
\ \isacommand{show}\isamarkupfalse%
\ {\isachardoublequoteopen}nabla{\isadigit{3}}\ {\isasymTurnstile}\ {\isacharparenleft}{\kern0pt}subst\ s{\isadigit{1}}{\isacharparenright}{\kern0pt}\ {\isacharparenleft}{\kern0pt}nabla{\isadigit{2}}\ {\isasymunion}\ nabla{\isadigit{1}}{\isacharparenright}{\kern0pt}{\isachardoublequoteclose}\isanewline
\ \ \ \ \isacommand{by}\isamarkupfalse%
{\isacharparenleft}{\kern0pt}auto\ simp\ add{\isacharcolon}{\kern0pt}\ ext{\isacharunderscore}{\kern0pt}subst{\isacharunderscore}{\kern0pt}def{\isacharparenright}{\kern0pt}\isanewline
\isacommand{qed}\isamarkupfalse%
%
\endisatagproof
{\isafoldproof}%
%
\isadelimproof
\isanewline
%
\endisadelimproof
\isanewline
\isacommand{lemma}\isamarkupfalse%
\ P{\isadigit{1}}{\isacharunderscore}{\kern0pt}to{\isacharunderscore}{\kern0pt}P{\isadigit{2}}{\isacharunderscore}{\kern0pt}red{\isacharunderscore}{\kern0pt}plus{\isadigit{2}}{\isacharcolon}{\kern0pt}\isanewline
\ \ \isakeyword{assumes}\ {\isachardoublequoteopen}P{\isadigit{1}}\ {\isasymTurnstile}\ {\isacharparenleft}{\kern0pt}nabla{\isacharcomma}{\kern0pt}\ s{\isacharparenright}{\kern0pt}\ {\isasymRightarrow}\ P{\isadigit{2}}{\isachardoublequoteclose}\isanewline
\ \ \isakeyword{and}\ {\isachardoublequoteopen}{\isacharparenleft}{\kern0pt}nabla{\isadigit{1}}{\isacharcomma}{\kern0pt}\ s{\isadigit{1}}{\isacharparenright}{\kern0pt}\ {\isasymin}\ U\ P{\isadigit{1}}{\isachardoublequoteclose}\isanewline
\isakeyword{shows}\ {\isachardoublequoteopen}nabla{\isadigit{1}}\ {\isasymTurnstile}\ subst\ {\isacharparenleft}{\kern0pt}s{\isadigit{1}}\ {\isasymbullet}\ s{\isacharparenright}{\kern0pt}\ {\isasymapprox}\ subst\ s{\isadigit{1}}{\isachardoublequoteclose}\isanewline
%
\isadelimproof
\ \ %
\endisadelimproof
%
\isatagproof
\isacommand{using}\isamarkupfalse%
\ assms\isanewline
\isacommand{proof}\isamarkupfalse%
{\isacharparenleft}{\kern0pt}induction\ P{\isadigit{1}}{\isasymequiv}{\isachardoublequoteopen}P{\isadigit{1}}{\isachardoublequoteclose}\ nablas{\isasymequiv}{\isachardoublequoteopen}{\isacharparenleft}{\kern0pt}nabla{\isacharcomma}{\kern0pt}\ s{\isacharparenright}{\kern0pt}{\isachardoublequoteclose}\ P{\isadigit{2}}{\isasymequiv}{\isachardoublequoteopen}P{\isadigit{2}}{\isachardoublequoteclose}\ arbitrary{\isacharcolon}{\kern0pt}\ s\ nabla{\isadigit{1}}\ nabla\ s{\isadigit{1}}\ P{\isadigit{1}}\ P{\isadigit{2}}\ rule{\isacharcolon}{\kern0pt}\ red{\isacharunderscore}{\kern0pt}plus{\isachardot}{\kern0pt}induct{\isacharparenright}{\kern0pt}\isanewline
\ \ \isacommand{case}\isamarkupfalse%
\ {\isacharparenleft}{\kern0pt}sred{\isacharunderscore}{\kern0pt}single\ P{\isadigit{1}}\ s\ P{\isadigit{2}}\ nabla{\isadigit{1}}\ s{\isadigit{1}}{\isacharparenright}{\kern0pt}\isanewline
\ \ \isacommand{have}\isamarkupfalse%
\ {\isachardoublequoteopen}P{\isadigit{1}}\ {\isasymturnstile}s\ {\isasymleadsto}P{\isadigit{2}}{\isachardoublequoteclose}\ {\isachardoublequoteopen}{\isacharparenleft}{\kern0pt}nabla{\isadigit{1}}{\isacharcomma}{\kern0pt}\ s{\isadigit{1}}{\isacharparenright}{\kern0pt}\ {\isasymin}\ U\ P{\isadigit{1}}{\isachardoublequoteclose}\ \isacommand{by}\isamarkupfalse%
\ fact{\isacharplus}{\kern0pt}\isanewline
\ \ \isacommand{then}\isamarkupfalse%
\ \isacommand{show}\isamarkupfalse%
\ {\isachardoublequoteopen}nabla{\isadigit{1}}\ {\isasymTurnstile}\ subst\ {\isacharparenleft}{\kern0pt}s{\isadigit{1}}\ {\isasymbullet}\ s{\isacharparenright}{\kern0pt}\ {\isasymapprox}subst\ s{\isadigit{1}}{\isachardoublequoteclose}\isanewline
\ \ \ \ \isacommand{using}\isamarkupfalse%
\ P{\isadigit{1}}{\isacharunderscore}{\kern0pt}to{\isacharunderscore}{\kern0pt}P{\isadigit{2}}{\isacharunderscore}{\kern0pt}sred\ \isacommand{by}\isamarkupfalse%
\ blast\isanewline
\isacommand{next}\isamarkupfalse%
\isanewline
\ \ \isacommand{case}\isamarkupfalse%
\ {\isacharparenleft}{\kern0pt}cred{\isacharunderscore}{\kern0pt}single\ P{\isadigit{1}}\ nabla{\isadigit{1}}\ P{\isadigit{2}}\ nabla{\isadigit{2}}\ s{\isadigit{1}}{\isacharparenright}{\kern0pt}\isanewline
\ \ \isacommand{show}\isamarkupfalse%
\ {\isachardoublequoteopen}nabla{\isadigit{2}}\ {\isasymTurnstile}\ subst\ {\isacharparenleft}{\kern0pt}s{\isadigit{1}}\ {\isasymbullet}\ {\isacharbrackleft}{\kern0pt}{\isacharbrackright}{\kern0pt}{\isacharparenright}{\kern0pt}\ {\isasymapprox}subst\ s{\isadigit{1}}{\isachardoublequoteclose}\isanewline
\ \ \ \ \isacommand{by}\isamarkupfalse%
\ {\isacharparenleft}{\kern0pt}simp\ add{\isacharcolon}{\kern0pt}\ subst{\isacharunderscore}{\kern0pt}equ{\isacharunderscore}{\kern0pt}refl{\isacharparenright}{\kern0pt}\ \isanewline
\isacommand{next}\isamarkupfalse%
\isanewline
\ \ \isacommand{case}\isamarkupfalse%
\ {\isacharparenleft}{\kern0pt}sred{\isacharunderscore}{\kern0pt}step\ P{\isadigit{1}}\ s\ P{\isadigit{2}}\ nabla{\isadigit{2}}\ s{\isadigit{2}}\ P{\isadigit{3}}\ nabla{\isadigit{1}}\ s{\isadigit{1}}{\isacharparenright}{\kern0pt}\isanewline
\ \ \isacommand{have}\isamarkupfalse%
\ IH{\isacharcolon}{\kern0pt}\ {\isachardoublequoteopen}{\isacharparenleft}{\kern0pt}nabla{\isadigit{1}}{\isacharcomma}{\kern0pt}\ s{\isadigit{1}}{\isacharparenright}{\kern0pt}\ {\isasymin}\ U\ P{\isadigit{2}}\ {\isasymLongrightarrow}\ \ nabla{\isadigit{1}}\ {\isasymTurnstile}\ subst\ {\isacharparenleft}{\kern0pt}s{\isadigit{1}}\ {\isasymbullet}\ s{\isadigit{2}}{\isacharparenright}{\kern0pt}\ {\isasymapprox}subst\ s{\isadigit{1}}{\isachardoublequoteclose}\ \isacommand{by}\isamarkupfalse%
\ fact\isanewline
\ \ \isacommand{have}\isamarkupfalse%
\ as{\isacharcolon}{\kern0pt}\ {\isachardoublequoteopen}P{\isadigit{1}}\ {\isasymturnstile}s\ {\isasymleadsto}P{\isadigit{2}}{\isachardoublequoteclose}\ {\isachardoublequoteopen}{\isacharparenleft}{\kern0pt}nabla{\isadigit{1}}{\isacharcomma}{\kern0pt}\ s{\isadigit{1}}{\isacharparenright}{\kern0pt}\ {\isasymin}\ U\ P{\isadigit{1}}{\isachardoublequoteclose}\ \isacommand{by}\isamarkupfalse%
\ fact{\isacharplus}{\kern0pt}\isanewline
\isanewline
\ \ \isacommand{from}\isamarkupfalse%
\ as\ \isacommand{have}\isamarkupfalse%
\ {\isachardoublequoteopen}{\isacharparenleft}{\kern0pt}nabla{\isadigit{1}}{\isacharcomma}{\kern0pt}\ s{\isadigit{1}}{\isacharparenright}{\kern0pt}\ {\isasymin}\ U\ P{\isadigit{2}}{\isachardoublequoteclose}\ \isacommand{using}\isamarkupfalse%
\ P{\isadigit{1}}{\isacharunderscore}{\kern0pt}to{\isacharunderscore}{\kern0pt}P{\isadigit{2}}{\isacharunderscore}{\kern0pt}sred\ \isacommand{by}\isamarkupfalse%
\ blast\ \isanewline
\ \ \isacommand{with}\isamarkupfalse%
\ IH\ \isacommand{have}\isamarkupfalse%
\ {\isachardoublequoteopen}nabla{\isadigit{1}}\ {\isasymTurnstile}\ subst\ {\isacharparenleft}{\kern0pt}s{\isadigit{1}}\ {\isasymbullet}\ s{\isadigit{2}}{\isacharparenright}{\kern0pt}\ {\isasymapprox}\ subst\ s{\isadigit{1}}{\isachardoublequoteclose}\ \isacommand{by}\isamarkupfalse%
\ blast\isanewline
\ \ \isacommand{then}\isamarkupfalse%
\ \isacommand{have}\isamarkupfalse%
\ {\isachardoublequoteopen}nabla{\isadigit{1}}\ {\isasymTurnstile}\ subst\ {\isacharparenleft}{\kern0pt}{\isacharparenleft}{\kern0pt}s{\isadigit{1}}\ {\isasymbullet}\ s{\isadigit{2}}{\isacharparenright}{\kern0pt}\ {\isasymbullet}\ s{\isacharparenright}{\kern0pt}\ {\isasymapprox}\ subst\ {\isacharparenleft}{\kern0pt}s{\isadigit{1}}\ {\isasymbullet}\ s{\isacharparenright}{\kern0pt}{\isachardoublequoteclose}\isanewline
\ \ \ \ \isacommand{by}\isamarkupfalse%
\ {\isacharparenleft}{\kern0pt}simp\ add{\isacharcolon}{\kern0pt}\ subst{\isacharunderscore}{\kern0pt}cancel{\isacharunderscore}{\kern0pt}right{\isacharparenright}{\kern0pt}\ \ \isanewline
\ \ \isacommand{moreover}\isamarkupfalse%
\isanewline
\ \ \isacommand{from}\isamarkupfalse%
\ as\ \isacommand{have}\isamarkupfalse%
\ {\isachardoublequoteopen}nabla{\isadigit{1}}\ {\isasymTurnstile}\ subst\ {\isacharparenleft}{\kern0pt}s{\isadigit{1}}\ {\isasymbullet}\ s{\isacharparenright}{\kern0pt}\ {\isasymapprox}subst\ s{\isadigit{1}}{\isachardoublequoteclose}\isanewline
\ \ \ \ \isacommand{using}\isamarkupfalse%
\ P{\isadigit{1}}{\isacharunderscore}{\kern0pt}to{\isacharunderscore}{\kern0pt}P{\isadigit{2}}{\isacharunderscore}{\kern0pt}sred\ \isacommand{by}\isamarkupfalse%
\ blast\ \isanewline
\ \ \isacommand{ultimately}\isamarkupfalse%
\isanewline
\ \ \isacommand{show}\isamarkupfalse%
\ {\isachardoublequoteopen}nabla{\isadigit{1}}\ {\isasymTurnstile}\ subst\ {\isacharparenleft}{\kern0pt}s{\isadigit{1}}\ {\isasymbullet}\ {\isacharparenleft}{\kern0pt}s{\isadigit{2}}\ {\isasymbullet}\ s{\isacharparenright}{\kern0pt}{\isacharparenright}{\kern0pt}\ {\isasymapprox}\ subst\ s{\isadigit{1}}{\isachardoublequoteclose}\isanewline
\ \ \ \ \isacommand{using}\isamarkupfalse%
\ subst{\isacharunderscore}{\kern0pt}assoc\ subst{\isacharunderscore}{\kern0pt}trans\ \isacommand{by}\isamarkupfalse%
\ metis\isanewline
\isacommand{next}\isamarkupfalse%
\isanewline
\ \ \isacommand{case}\isamarkupfalse%
\ {\isacharparenleft}{\kern0pt}cred{\isacharunderscore}{\kern0pt}step\ P{\isadigit{1}}\ nabla{\isadigit{1}}\ P{\isadigit{2}}\ nabla{\isadigit{2}}\ P{\isadigit{3}}{\isacharparenright}{\kern0pt}\isanewline
\ \ \isacommand{show}\isamarkupfalse%
\ {\isachardoublequoteopen}nabla{\isadigit{1}}\ {\isasymTurnstile}\ subst\ {\isacharparenleft}{\kern0pt}s{\isadigit{1}}\ {\isasymbullet}\ {\isacharbrackleft}{\kern0pt}{\isacharbrackright}{\kern0pt}{\isacharparenright}{\kern0pt}\ {\isasymapprox}\ subst\ s{\isadigit{1}}{\isachardoublequoteclose}\isanewline
\ \ \ \ \isacommand{by}\isamarkupfalse%
\ {\isacharparenleft}{\kern0pt}simp\ add{\isacharcolon}{\kern0pt}\ subst{\isacharunderscore}{\kern0pt}equ{\isacharunderscore}{\kern0pt}refl{\isacharparenright}{\kern0pt}\isanewline
\isacommand{qed}\isamarkupfalse%
%
\endisatagproof
{\isafoldproof}%
%
\isadelimproof
\isanewline
%
\endisadelimproof
\isanewline
\isacommand{lemma}\isamarkupfalse%
\ P{\isadigit{1}}{\isacharunderscore}{\kern0pt}from{\isacharunderscore}{\kern0pt}P{\isadigit{2}}{\isacharunderscore}{\kern0pt}red{\isacharunderscore}{\kern0pt}plus{\isacharcolon}{\kern0pt}\isanewline
\ \ \isakeyword{assumes}\ {\isachardoublequoteopen}P{\isadigit{1}}\ {\isasymTurnstile}{\isacharparenleft}{\kern0pt}nabla{\isacharcomma}{\kern0pt}s{\isacharparenright}{\kern0pt}{\isasymRightarrow}P{\isadigit{2}}{\isachardoublequoteclose}\isanewline
\ \ \ \ \isakeyword{and}\ {\isachardoublequoteopen}{\isacharparenleft}{\kern0pt}nabla{\isadigit{1}}{\isacharcomma}{\kern0pt}s{\isadigit{1}}{\isacharparenright}{\kern0pt}{\isasymin}\ U\ P{\isadigit{2}}{\isachardoublequoteclose}\ \isakeyword{and}\ {\isachardoublequoteopen}nabla{\isadigit{3}}{\isasymTurnstile}{\isacharparenleft}{\kern0pt}subst\ s{\isadigit{1}}{\isacharparenright}{\kern0pt}{\isacharparenleft}{\kern0pt}nabla{\isacharparenright}{\kern0pt}{\isachardoublequoteclose}\isanewline
\ \ \isakeyword{shows}\ {\isachardoublequoteopen}{\isacharparenleft}{\kern0pt}nabla{\isadigit{1}}{\isasymunion}nabla{\isadigit{3}}{\isacharcomma}{\kern0pt}{\isacharparenleft}{\kern0pt}s{\isadigit{1}}\ {\isasymbullet}\ s{\isacharparenright}{\kern0pt}{\isacharparenright}{\kern0pt}{\isasymin}\ U\ P{\isadigit{1}}{\isachardoublequoteclose}\isanewline
%
\isadelimproof
\ \ %
\endisadelimproof
%
\isatagproof
\isacommand{using}\isamarkupfalse%
\ assms\isanewline
\isacommand{proof}\isamarkupfalse%
{\isacharparenleft}{\kern0pt}induction\ P{\isadigit{1}}\ {\isasymequiv}\ {\isachardoublequoteopen}P{\isadigit{1}}{\isachardoublequoteclose}\ nablas{\isasymequiv}{\isachardoublequoteopen}{\isacharparenleft}{\kern0pt}nabla{\isacharcomma}{\kern0pt}\ s{\isacharparenright}{\kern0pt}{\isachardoublequoteclose}\ P{\isadigit{2}}{\isasymequiv}{\isachardoublequoteopen}P{\isadigit{2}}{\isachardoublequoteclose}\ arbitrary{\isacharcolon}{\kern0pt}\ s\ nabla{\isadigit{1}}\ nabla\ nabla{\isadigit{3}}\ s{\isadigit{1}}\ P{\isadigit{1}}\ P{\isadigit{2}}\ rule{\isacharcolon}{\kern0pt}\ red{\isacharunderscore}{\kern0pt}plus{\isachardot}{\kern0pt}induct{\isacharparenright}{\kern0pt}\isanewline
\ \ \isacommand{case}\isamarkupfalse%
\ {\isacharparenleft}{\kern0pt}sred{\isacharunderscore}{\kern0pt}single\ P{\isadigit{1}}\ s\ P{\isadigit{2}}\ nabla{\isadigit{1}}\ nabla{\isadigit{3}}{\isacharparenright}{\kern0pt}\isanewline
\ \ \isacommand{have}\isamarkupfalse%
\ {\isachardoublequoteopen}P{\isadigit{1}}\ {\isasymturnstile}\ s\ {\isasymleadsto}\ P{\isadigit{2}}{\isachardoublequoteclose}\ {\isachardoublequoteopen}{\isacharparenleft}{\kern0pt}nabla{\isadigit{1}}{\isacharcomma}{\kern0pt}\ s{\isadigit{1}}{\isacharparenright}{\kern0pt}\ {\isasymin}\ U\ P{\isadigit{2}}{\isachardoublequoteclose}\ \isacommand{by}\isamarkupfalse%
\ fact{\isacharplus}{\kern0pt}\isanewline
\ \ \isacommand{then}\isamarkupfalse%
\ \isacommand{have}\isamarkupfalse%
\ {\isachardoublequoteopen}{\isacharparenleft}{\kern0pt}nabla{\isadigit{1}}{\isacharcomma}{\kern0pt}\ s{\isadigit{1}}\ {\isasymbullet}\ s{\isacharparenright}{\kern0pt}\ {\isasymin}\ U\ P{\isadigit{1}}{\isachardoublequoteclose}\isanewline
\ \ \ \ \isacommand{using}\isamarkupfalse%
\ P{\isadigit{1}}{\isacharunderscore}{\kern0pt}from{\isacharunderscore}{\kern0pt}P{\isadigit{2}}{\isacharunderscore}{\kern0pt}sred\ \isacommand{by}\isamarkupfalse%
\ simp\isanewline
\ \ \isacommand{then}\isamarkupfalse%
\ \isacommand{show}\isamarkupfalse%
\ {\isachardoublequoteopen}{\isacharparenleft}{\kern0pt}nabla{\isadigit{1}}\ {\isasymunion}\ nabla{\isadigit{3}}{\isacharcomma}{\kern0pt}\ s{\isadigit{1}}\ {\isasymbullet}\ s{\isacharparenright}{\kern0pt}\ {\isasymin}\ U\ P{\isadigit{1}}{\isachardoublequoteclose}\isanewline
\ \ \ \ \isacommand{using}\isamarkupfalse%
\ P{\isadigit{1}}{\isacharunderscore}{\kern0pt}from{\isacharunderscore}{\kern0pt}P{\isadigit{2}}{\isacharunderscore}{\kern0pt}sred\ solution{\isacharunderscore}{\kern0pt}weak\ \isacommand{by}\isamarkupfalse%
\ simp\isanewline
\isacommand{next}\isamarkupfalse%
\isanewline
\ \ \isacommand{case}\isamarkupfalse%
\ {\isacharparenleft}{\kern0pt}cred{\isacharunderscore}{\kern0pt}single\ P{\isadigit{1}}\ nabla\ P{\isadigit{2}}\ nabla{\isadigit{1}}\ nabla{\isadigit{3}}{\isacharparenright}{\kern0pt}\isanewline
\ \ \isacommand{have}\isamarkupfalse%
\ {\isachardoublequoteopen}P{\isadigit{1}}\ {\isasymturnstile}\ nabla\ {\isasymrightarrow}\ P{\isadigit{2}}{\isachardoublequoteclose}\isanewline
\ \ \ \ \ \ {\isachardoublequoteopen}{\isacharparenleft}{\kern0pt}nabla{\isadigit{1}}{\isacharcomma}{\kern0pt}\ s{\isadigit{1}}{\isacharparenright}{\kern0pt}\ {\isasymin}\ U\ P{\isadigit{2}}{\isachardoublequoteclose}\ \isanewline
\ \ \ \ \ \ {\isachardoublequoteopen}nabla{\isadigit{3}}\ {\isasymTurnstile}\ {\isacharparenleft}{\kern0pt}subst\ s{\isadigit{1}}{\isacharparenright}{\kern0pt}\ nabla{\isachardoublequoteclose}\ \isacommand{by}\isamarkupfalse%
\ fact{\isacharplus}{\kern0pt}\isanewline
\ \ \isacommand{hence}\isamarkupfalse%
\ {\isachardoublequoteopen}{\isacharparenleft}{\kern0pt}nabla{\isadigit{1}}\ {\isasymunion}\ nabla{\isadigit{3}}{\isacharcomma}{\kern0pt}\ s{\isadigit{1}}{\isacharparenright}{\kern0pt}\ {\isasymin}\ U\ P{\isadigit{1}}{\isachardoublequoteclose}\isanewline
\ \ \ \ \isacommand{using}\isamarkupfalse%
\ P{\isadigit{1}}{\isacharunderscore}{\kern0pt}from{\isacharunderscore}{\kern0pt}P{\isadigit{2}}{\isacharunderscore}{\kern0pt}cred\ \isacommand{by}\isamarkupfalse%
\ simp\isanewline
\ \ \isacommand{then}\isamarkupfalse%
\ \isacommand{show}\isamarkupfalse%
\ {\isachardoublequoteopen}{\isacharparenleft}{\kern0pt}nabla{\isadigit{1}}\ {\isasymunion}\ nabla{\isadigit{3}}{\isacharcomma}{\kern0pt}\ s{\isadigit{1}}\ {\isasymbullet}\ {\isacharbrackleft}{\kern0pt}{\isacharbrackright}{\kern0pt}{\isacharparenright}{\kern0pt}\ {\isasymin}\ U\ P{\isadigit{1}}{\isachardoublequoteclose}\isanewline
\ \ \ \ \isacommand{using}\isamarkupfalse%
\ solution{\isacharunderscore}{\kern0pt}comp{\isacharunderscore}{\kern0pt}id\ \isacommand{by}\isamarkupfalse%
\ auto\isanewline
\isacommand{next}\isamarkupfalse%
\isanewline
\ \ \isacommand{case}\isamarkupfalse%
\ {\isacharparenleft}{\kern0pt}sred{\isacharunderscore}{\kern0pt}step\ P{\isadigit{1}}\ s\ P{\isadigit{2}}\ nabla{\isadigit{2}}\ s{\isadigit{2}}\ P{\isadigit{3}}\ nabla{\isadigit{1}}\ nabla{\isadigit{3}}{\isacharparenright}{\kern0pt}\isanewline
\ \ \isacommand{have}\isamarkupfalse%
\ IH{\isacharcolon}{\kern0pt}\ {\isachardoublequoteopen}{\isacharparenleft}{\kern0pt}nabla{\isadigit{1}}{\isacharcomma}{\kern0pt}\ s{\isadigit{1}}{\isacharparenright}{\kern0pt}\ {\isasymin}\ U\ P{\isadigit{3}}\ {\isasymLongrightarrow}\isanewline
\ \ \ \ \ nabla{\isadigit{3}}\ {\isasymTurnstile}\ {\isacharparenleft}{\kern0pt}subst\ s{\isadigit{1}}{\isacharparenright}{\kern0pt}\ nabla{\isadigit{2}}\ \ {\isasymLongrightarrow}\ {\isacharparenleft}{\kern0pt}nabla{\isadigit{1}}\ {\isasymunion}\ nabla{\isadigit{3}}{\isacharcomma}{\kern0pt}\ s{\isadigit{1}}\ {\isasymbullet}\ s{\isadigit{2}}{\isacharparenright}{\kern0pt}\ {\isasymin}\ U\ P{\isadigit{2}}{\isachardoublequoteclose}\ \isacommand{by}\isamarkupfalse%
\ fact{\isacharplus}{\kern0pt}\isanewline
\ \ \isacommand{have}\isamarkupfalse%
\ as{\isacharcolon}{\kern0pt}\ {\isachardoublequoteopen}{\isacharparenleft}{\kern0pt}nabla{\isadigit{1}}{\isacharcomma}{\kern0pt}\ s{\isadigit{1}}{\isacharparenright}{\kern0pt}\ {\isasymin}\ U\ P{\isadigit{3}}{\isachardoublequoteclose}\ \isanewline
\ \ \ \ {\isachardoublequoteopen}nabla{\isadigit{3}}\ {\isasymTurnstile}\ {\isacharparenleft}{\kern0pt}subst\ s{\isadigit{1}}{\isacharparenright}{\kern0pt}\ nabla{\isadigit{2}}{\isachardoublequoteclose}\ \isanewline
\ \ \ \ {\isachardoublequoteopen}P{\isadigit{1}}\ {\isasymturnstile}\ s\ {\isasymleadsto}\ P{\isadigit{2}}{\isachardoublequoteclose}\ \isacommand{by}\isamarkupfalse%
\ fact{\isacharplus}{\kern0pt}\isanewline
\ \ \isacommand{hence}\isamarkupfalse%
\ {\isachardoublequoteopen}{\isacharparenleft}{\kern0pt}nabla{\isadigit{1}}\ {\isasymunion}\ nabla{\isadigit{3}}{\isacharcomma}{\kern0pt}\ s{\isadigit{1}}\ {\isasymbullet}\ s{\isadigit{2}}{\isacharparenright}{\kern0pt}\ {\isasymin}\ U\ P{\isadigit{2}}{\isachardoublequoteclose}\isanewline
\ \ \ \ \isacommand{using}\isamarkupfalse%
\ IH\ \isacommand{by}\isamarkupfalse%
\ simp\isanewline
\ \ \isacommand{hence}\isamarkupfalse%
\ {\isachardoublequoteopen}{\isacharparenleft}{\kern0pt}nabla{\isadigit{1}}\ {\isasymunion}\ nabla{\isadigit{3}}{\isacharcomma}{\kern0pt}\ {\isacharparenleft}{\kern0pt}s{\isadigit{1}}\ {\isasymbullet}\ s{\isadigit{2}}{\isacharparenright}{\kern0pt}\ {\isasymbullet}\ s{\isacharparenright}{\kern0pt}\ {\isasymin}\ U\ P{\isadigit{1}}{\isachardoublequoteclose}\isanewline
\ \ \ \ \isacommand{using}\isamarkupfalse%
\ as{\isacharparenleft}{\kern0pt}{\isadigit{3}}{\isacharparenright}{\kern0pt}\isanewline
\ \ \ \ \isacommand{by}\isamarkupfalse%
\ {\isacharparenleft}{\kern0pt}simp\ add{\isacharcolon}{\kern0pt}\ P{\isadigit{1}}{\isacharunderscore}{\kern0pt}from{\isacharunderscore}{\kern0pt}P{\isadigit{2}}{\isacharunderscore}{\kern0pt}sred{\isacharparenright}{\kern0pt}\isanewline
\ \ \isacommand{then}\isamarkupfalse%
\ \isacommand{show}\isamarkupfalse%
\ {\isachardoublequoteopen}{\isacharparenleft}{\kern0pt}nabla{\isadigit{1}}\ {\isasymunion}\ nabla{\isadigit{3}}{\isacharcomma}{\kern0pt}\ s{\isadigit{1}}\ {\isasymbullet}\ {\isacharparenleft}{\kern0pt}s{\isadigit{2}}\ {\isasymbullet}\ s{\isacharparenright}{\kern0pt}{\isacharparenright}{\kern0pt}\ {\isasymin}\ U\ P{\isadigit{1}}{\isachardoublequoteclose}\isanewline
\ \ \ \ \isacommand{using}\isamarkupfalse%
\ solutions{\isacharunderscore}{\kern0pt}subst{\isacharunderscore}{\kern0pt}assoc\ \isacommand{by}\isamarkupfalse%
\ simp\isanewline
\isacommand{next}\isamarkupfalse%
\isanewline
\ \ \isacommand{case}\isamarkupfalse%
\ {\isacharparenleft}{\kern0pt}cred{\isacharunderscore}{\kern0pt}step\ P{\isadigit{1}}\ nabla\ P{\isadigit{2}}\ nabla{\isadigit{2}}\ P{\isadigit{3}}{\isacharparenright}{\kern0pt}\isanewline
\ \ \isacommand{have}\isamarkupfalse%
\ IH{\isacharcolon}{\kern0pt}\ {\isachardoublequoteopen}{\isacharparenleft}{\kern0pt}nabla{\isadigit{1}}{\isacharcomma}{\kern0pt}\ s{\isadigit{1}}{\isacharparenright}{\kern0pt}\ {\isasymin}\ U\ P{\isadigit{3}}\ {\isasymLongrightarrow}\isanewline
\ \ \ \ \ nabla{\isadigit{3}}\ {\isasymTurnstile}\ {\isacharparenleft}{\kern0pt}subst\ s{\isadigit{1}}{\isacharparenright}{\kern0pt}\ nabla{\isadigit{2}}\ \ {\isasymLongrightarrow}\ {\isacharparenleft}{\kern0pt}nabla{\isadigit{1}}\ {\isasymunion}\ nabla{\isadigit{3}}{\isacharcomma}{\kern0pt}\ s{\isadigit{1}}\ {\isasymbullet}\ {\isacharbrackleft}{\kern0pt}{\isacharbrackright}{\kern0pt}{\isacharparenright}{\kern0pt}\ {\isasymin}\ U\ P{\isadigit{2}}{\isachardoublequoteclose}\ \isacommand{by}\isamarkupfalse%
\ fact\isanewline
\ \ \isacommand{have}\isamarkupfalse%
\ as{\isacharcolon}{\kern0pt}\ {\isachardoublequoteopen}P{\isadigit{1}}\ {\isasymturnstile}\ nabla\ {\isasymrightarrow}\ P{\isadigit{2}}{\isachardoublequoteclose}\isanewline
\ \ \ \ \ \ \ \ \ \ \ {\isachardoublequoteopen}{\isacharparenleft}{\kern0pt}nabla{\isadigit{1}}{\isacharcomma}{\kern0pt}\ s{\isadigit{1}}{\isacharparenright}{\kern0pt}\ {\isasymin}\ U\ P{\isadigit{3}}{\isachardoublequoteclose}\isanewline
\ \ \ \ \ \ \ \ \ \ \ {\isachardoublequoteopen}nabla{\isadigit{3}}\ {\isasymTurnstile}\ {\isacharparenleft}{\kern0pt}subst\ s{\isadigit{1}}{\isacharparenright}{\kern0pt}\ {\isacharparenleft}{\kern0pt}nabla{\isadigit{2}}\ {\isasymunion}\ nabla{\isacharparenright}{\kern0pt}{\isachardoublequoteclose}\ \isacommand{by}\isamarkupfalse%
\ fact{\isacharplus}{\kern0pt}\isanewline
\ \ \isacommand{have}\isamarkupfalse%
\ ext{\isacharunderscore}{\kern0pt}substs{\isacharcolon}{\kern0pt}\ {\isachardoublequoteopen}nabla{\isadigit{3}}\ {\isasymTurnstile}\ {\isacharparenleft}{\kern0pt}subst\ s{\isadigit{1}}{\isacharparenright}{\kern0pt}\ {\isacharparenleft}{\kern0pt}nabla{\isadigit{2}}{\isacharparenright}{\kern0pt}{\isachardoublequoteclose}\ \isanewline
\ \ \ \ {\isachardoublequoteopen}nabla{\isadigit{3}}\ {\isasymTurnstile}\ {\isacharparenleft}{\kern0pt}subst\ s{\isadigit{1}}{\isacharparenright}{\kern0pt}\ nabla{\isachardoublequoteclose}\isanewline
\ \ \ \ \isacommand{using}\isamarkupfalse%
\ ext{\isacharunderscore}{\kern0pt}subst{\isacharunderscore}{\kern0pt}strong{\isacharbrackleft}{\kern0pt}OF\ as{\isacharparenleft}{\kern0pt}{\isadigit{3}}{\isacharparenright}{\kern0pt}{\isacharbrackright}{\kern0pt}\ \isacommand{by}\isamarkupfalse%
\ simp{\isacharplus}{\kern0pt}\isanewline
\ \ \isacommand{hence}\isamarkupfalse%
\ a{\isacharcolon}{\kern0pt}\ {\isachardoublequoteopen}{\isacharparenleft}{\kern0pt}nabla{\isadigit{1}}\ {\isasymunion}\ nabla{\isadigit{3}}{\isacharcomma}{\kern0pt}\ s{\isadigit{1}}\ {\isasymbullet}\ {\isacharbrackleft}{\kern0pt}{\isacharbrackright}{\kern0pt}{\isacharparenright}{\kern0pt}\ {\isasymin}\ U\ P{\isadigit{2}}{\isachardoublequoteclose}\isanewline
\ \ \ \ \isacommand{using}\isamarkupfalse%
\ IH\ as{\isacharparenleft}{\kern0pt}{\isadigit{2}}{\isacharparenright}{\kern0pt}\ \isacommand{by}\isamarkupfalse%
\ simp\isanewline
\ \ \isacommand{hence}\isamarkupfalse%
\ {\isachardoublequoteopen}{\isacharparenleft}{\kern0pt}nabla{\isadigit{1}}\ {\isasymunion}\ nabla{\isadigit{3}}\ {\isasymunion}\ nabla{\isadigit{3}}{\isacharcomma}{\kern0pt}\ s{\isadigit{1}}\ {\isasymbullet}\ {\isacharbrackleft}{\kern0pt}{\isacharbrackright}{\kern0pt}{\isacharparenright}{\kern0pt}\ {\isasymin}\ U\ P{\isadigit{1}}{\isachardoublequoteclose}\isanewline
\ \ \ \ \isacommand{using}\isamarkupfalse%
\ as{\isacharparenleft}{\kern0pt}{\isadigit{1}}{\isacharparenright}{\kern0pt}\ ext{\isacharunderscore}{\kern0pt}substs{\isacharparenleft}{\kern0pt}{\isadigit{2}}{\isacharparenright}{\kern0pt}\ P{\isadigit{1}}{\isacharunderscore}{\kern0pt}from{\isacharunderscore}{\kern0pt}P{\isadigit{2}}{\isacharunderscore}{\kern0pt}cred\ solution{\isacharunderscore}{\kern0pt}comp{\isacharunderscore}{\kern0pt}id\ \isacommand{by}\isamarkupfalse%
\ blast\isanewline
\ \ \isacommand{then}\isamarkupfalse%
\ \isacommand{show}\isamarkupfalse%
\ {\isachardoublequoteopen}{\isacharparenleft}{\kern0pt}nabla{\isadigit{1}}\ {\isasymunion}\ nabla{\isadigit{3}}{\isacharcomma}{\kern0pt}\ s{\isadigit{1}}\ {\isasymbullet}\ {\isacharbrackleft}{\kern0pt}{\isacharbrackright}{\kern0pt}{\isacharparenright}{\kern0pt}\ {\isasymin}\ U\ P{\isadigit{1}}{\isachardoublequoteclose}\isanewline
\ \ \ \ \isacommand{unfolding}\isamarkupfalse%
\ all{\isacharunderscore}{\kern0pt}solutions{\isacharunderscore}{\kern0pt}def\ \isacommand{using}\isamarkupfalse%
\ Un{\isacharunderscore}{\kern0pt}absorb\ \isacommand{by}\isamarkupfalse%
\ simp\isanewline
\isacommand{qed}\isamarkupfalse%
%
\endisatagproof
{\isafoldproof}%
%
\isadelimproof
\isanewline
%
\endisadelimproof
\isanewline
\isacommand{lemma}\isamarkupfalse%
\ mgu{\isacharcolon}{\kern0pt}\ \isanewline
\ \ \isakeyword{assumes}\ {\isachardoublequoteopen}P\ {\isasymTurnstile}{\isacharparenleft}{\kern0pt}nabla{\isacharcomma}{\kern0pt}s{\isacharparenright}{\kern0pt}{\isasymRightarrow}{\isacharparenleft}{\kern0pt}{\isacharbrackleft}{\kern0pt}{\isacharbrackright}{\kern0pt}{\isacharcomma}{\kern0pt}{\isacharbrackleft}{\kern0pt}{\isacharbrackright}{\kern0pt}{\isacharparenright}{\kern0pt}{\isachardoublequoteclose}\isanewline
\ \ \isakeyword{shows}\ {\isachardoublequoteopen}mgu\ P\ {\isacharparenleft}{\kern0pt}nabla{\isacharcomma}{\kern0pt}s{\isacharparenright}{\kern0pt}\ {\isasymand}\ idem\ {\isacharparenleft}{\kern0pt}nabla{\isacharcomma}{\kern0pt}s{\isacharparenright}{\kern0pt}{\isachardoublequoteclose}\isanewline
%
\isadelimproof
%
\endisadelimproof
%
\isatagproof
\isacommand{proof}\isamarkupfalse%
{\isacharparenleft}{\kern0pt}rule\ mgu{\isacharunderscore}{\kern0pt}idem{\isacharparenright}{\kern0pt}\isanewline
\ \ \isacommand{have}\isamarkupfalse%
\ i{\isacharcolon}{\kern0pt}\ {\isachardoublequoteopen}{\isacharparenleft}{\kern0pt}{\isacharbraceleft}{\kern0pt}{\isacharbraceright}{\kern0pt}{\isacharcomma}{\kern0pt}{\isacharbrackleft}{\kern0pt}{\isacharbrackright}{\kern0pt}{\isacharparenright}{\kern0pt}\ {\isasymin}\ U\ {\isacharparenleft}{\kern0pt}{\isacharbrackleft}{\kern0pt}{\isacharbrackright}{\kern0pt}{\isacharcomma}{\kern0pt}{\isacharbrackleft}{\kern0pt}{\isacharbrackright}{\kern0pt}{\isacharparenright}{\kern0pt}{\isachardoublequoteclose}\isanewline
\ \ \ \ \isacommand{unfolding}\isamarkupfalse%
\ all{\isacharunderscore}{\kern0pt}solutions{\isacharunderscore}{\kern0pt}def\ \isacommand{by}\isamarkupfalse%
\ simp\isanewline
\ \ \isacommand{then}\isamarkupfalse%
\ \isacommand{show}\isamarkupfalse%
\ {\isachardoublequoteopen}{\isacharparenleft}{\kern0pt}nabla{\isacharcomma}{\kern0pt}\ s{\isacharparenright}{\kern0pt}\ {\isasymin}\ U\ P{\isachardoublequoteclose}\isanewline
\ \ \ \ \isacommand{using}\isamarkupfalse%
\ P{\isadigit{1}}{\isacharunderscore}{\kern0pt}from{\isacharunderscore}{\kern0pt}P{\isadigit{2}}{\isacharunderscore}{\kern0pt}red{\isacharunderscore}{\kern0pt}plus{\isacharbrackleft}{\kern0pt}OF\ assms\ i{\isacharbrackright}{\kern0pt}\ \isanewline
\ \ \ \ \ \ ext{\isacharunderscore}{\kern0pt}subst{\isacharunderscore}{\kern0pt}id\ \ solution{\isacharunderscore}{\kern0pt}comp{\isacharunderscore}{\kern0pt}id{\isacharparenleft}{\kern0pt}{\isadigit{2}}{\isacharparenright}{\kern0pt}\ \isacommand{by}\isamarkupfalse%
\ simp\isanewline
\ \ \ \ \ \ \ \ \ \ \ \ \ \ \ \ \ \ \ \ \ \ \ \ \ \ \ \ \isanewline
\ \ \isacommand{show}\isamarkupfalse%
\ {\isachardoublequoteopen}{\isasymforall}{\isacharparenleft}{\kern0pt}nabla{\isadigit{2}}{\isacharcomma}{\kern0pt}\ s{\isadigit{2}}{\isacharparenright}{\kern0pt}{\isasymin}\ U\ P{\isachardot}{\kern0pt}\ \ nabla{\isadigit{2}}\ {\isasymTurnstile}\ {\isacharparenleft}{\kern0pt}subst\ s{\isadigit{2}}{\isacharparenright}{\kern0pt}\ nabla\ {\isasymand}\ \ \isanewline
\ \ \ \ \ \ \ \ \ \ \ \ \ \ \ \ \ \ \ \ \ \ \ nabla{\isadigit{2}}\ {\isasymTurnstile}\ subst\ {\isacharparenleft}{\kern0pt}s{\isadigit{2}}\ {\isasymbullet}\ s{\isacharparenright}{\kern0pt}\ {\isasymapprox}\ subst\ s{\isadigit{2}}{\isachardoublequoteclose}\isanewline
\ \ \isacommand{proof}\isamarkupfalse%
\isanewline
\ \ \ \ \isacommand{fix}\isamarkupfalse%
\ x\isanewline
\ \ \ \ \isacommand{assume}\isamarkupfalse%
\ {\isachardoublequoteopen}x\ {\isasymin}\ U\ P{\isachardoublequoteclose}\isanewline
\ \ \ \ \isacommand{then}\isamarkupfalse%
\ \isacommand{show}\isamarkupfalse%
\ {\isachardoublequoteopen}case\ x\ of\ {\isacharparenleft}{\kern0pt}nabla{\isadigit{2}}{\isacharcomma}{\kern0pt}\ s{\isadigit{2}}{\isacharparenright}{\kern0pt}\ {\isasymRightarrow}\ \isanewline
\ \ \ \ \ \ \ \ \ \ \ \ \ \ \ nabla{\isadigit{2}}\ {\isasymTurnstile}\ {\isacharparenleft}{\kern0pt}subst\ s{\isadigit{2}}{\isacharparenright}{\kern0pt}\ nabla\ {\isasymand}\ \ \isanewline
\ \ \ \ \ \ \ \ \ \ \ \ \ \ \ \ \ \ \ \ \ \ \ nabla{\isadigit{2}}\ {\isasymTurnstile}\ subst\ {\isacharparenleft}{\kern0pt}s{\isadigit{2}}\ {\isasymbullet}\ s{\isacharparenright}{\kern0pt}\ {\isasymapprox}\ subst\ s{\isadigit{2}}{\isachardoublequoteclose}\isanewline
\ \ \ \ \isacommand{proof}\isamarkupfalse%
\ {\isacharparenleft}{\kern0pt}cases\ x{\isacharparenright}{\kern0pt}\isanewline
\ \ \ \ \ \ \isacommand{case}\isamarkupfalse%
\ {\isacharparenleft}{\kern0pt}Pair\ nabla{\isadigit{2}}\ s{\isadigit{2}}{\isacharparenright}{\kern0pt}\isanewline
\ \ \ \ \ \ \isacommand{hence}\isamarkupfalse%
\ {\isachardoublequoteopen}{\isacharparenleft}{\kern0pt}nabla{\isadigit{2}}{\isacharcomma}{\kern0pt}s{\isadigit{2}}{\isacharparenright}{\kern0pt}\ {\isasymin}\ U\ P{\isachardoublequoteclose}\isanewline
\ \ \ \ \ \ \ \ \isacommand{using}\isamarkupfalse%
\ {\isacartoucheopen}x\ {\isasymin}\ U\ P{\isacartoucheclose}\ \isacommand{by}\isamarkupfalse%
\ auto\isanewline
\ \ \ \ \ \ \isacommand{hence}\isamarkupfalse%
\ {\isachardoublequoteopen}nabla{\isadigit{2}}\ {\isasymTurnstile}\ {\isacharparenleft}{\kern0pt}subst\ s{\isadigit{2}}{\isacharparenright}{\kern0pt}\ nabla\ {\isasymand}\ nabla{\isadigit{2}}\ {\isasymTurnstile}\ subst\ {\isacharparenleft}{\kern0pt}s{\isadigit{2}}\ {\isasymbullet}\ s{\isacharparenright}{\kern0pt}\ {\isasymapprox}\ subst\ s{\isadigit{2}}{\isachardoublequoteclose}\isanewline
\ \ \ \ \ \ \ \ \isacommand{using}\isamarkupfalse%
\ assms\ P{\isadigit{1}}{\isacharunderscore}{\kern0pt}to{\isacharunderscore}{\kern0pt}P{\isadigit{2}}{\isacharunderscore}{\kern0pt}red{\isacharunderscore}{\kern0pt}plus{\isadigit{2}}\ P{\isadigit{1}}{\isacharunderscore}{\kern0pt}to{\isacharunderscore}{\kern0pt}P{\isadigit{2}}{\isacharunderscore}{\kern0pt}red{\isacharunderscore}{\kern0pt}plus{\isadigit{3}}\ \isacommand{by}\isamarkupfalse%
\ auto{\isacharplus}{\kern0pt}\isanewline
\ \ \ \ \ \ \isacommand{with}\isamarkupfalse%
\ Pair\ \isacommand{show}\isamarkupfalse%
\ {\isacharquery}{\kern0pt}thesis\ \isacommand{by}\isamarkupfalse%
\ simp\isanewline
\ \ \ \ \isacommand{qed}\isamarkupfalse%
\isanewline
\ \ \isacommand{qed}\isamarkupfalse%
\isanewline
\isacommand{qed}\isamarkupfalse%
%
\endisatagproof
{\isafoldproof}%
%
\isadelimproof
\isanewline
%
\endisadelimproof
\isanewline
%
\isadelimtheory
\isanewline
%
\endisadelimtheory
%
\isatagtheory
\isacommand{end}\isamarkupfalse%
%
\endisatagtheory
{\isafoldtheory}%
%
\isadelimtheory
%
\endisadelimtheory
%
\end{isabellebody}%
\endinput
%:%file=~/nominal_unification_Isabelle/nominal_alpha_unification/Mgu.thy%:%
%:%10=1%:%
%:%11=1%:%
%:%12=2%:%
%:%13=3%:%
%:%14=4%:%
%:%15=5%:%
%:%20=5%:%
%:%23=6%:%
%:%24=7%:%
%:%25=8%:%
%:%26=9%:%
%:%27=9%:%
%:%28=10%:%
%:%29=11%:%
%:%30=11%:%
%:%31=12%:%
%:%32=13%:%
%:%33=14%:%
%:%34=15%:%
%:%35=16%:%
%:%36=17%:%
%:%37=18%:%
%:%38=19%:%
%:%39=19%:%
%:%40=20%:%
%:%41=21%:%
%:%42=21%:%
%:%43=22%:%
%:%44=23%:%
%:%45=23%:%
%:%46=24%:%
%:%47=25%:%
%:%49=27%:%
%:%50=28%:%
%:%51=29%:%
%:%52=30%:%
%:%53=31%:%
%:%54=32%:%
%:%55=32%:%
%:%56=33%:%
%:%57=33%:%
%:%58=34%:%
%:%59=34%:%
%:%60=35%:%
%:%61=36%:%
%:%62=37%:%
%:%63=37%:%
%:%64=38%:%
%:%65=39%:%
%:%66=40%:%
%:%67=41%:%
%:%68=42%:%
%:%69=42%:%
%:%70=43%:%
%:%71=44%:%
%:%72=45%:%
%:%73=46%:%
%:%74=47%:%
%:%75=47%:%
%:%76=48%:%
%:%77=49%:%
%:%78=50%:%
%:%79=51%:%
%:%80=52%:%
%:%81=53%:%
%:%82=53%:%
%:%83=54%:%
%:%85=56%:%
%:%86=57%:%
%:%87=58%:%
%:%88=59%:%
%:%89=60%:%
%:%90=61%:%
%:%91=61%:%
%:%92=62%:%
%:%93=63%:%
%:%94=64%:%
%:%95=65%:%
%:%96=66%:%
%:%97=67%:%
%:%98=68%:%
%:%99=68%:%
%:%100=69%:%
%:%101=70%:%
%:%102=71%:%
%:%103=71%:%
%:%104=72%:%
%:%105=73%:%
%:%106=74%:%
%:%107=75%:%
%:%108=75%:%
%:%109=76%:%
%:%110=77%:%
%:%111=78%:%
%:%112=79%:%
%:%113=80%:%
%:%114=81%:%
%:%115=82%:%
%:%116=82%:%
%:%117=83%:%
%:%118=84%:%
%:%119=85%:%
%:%120=86%:%
%:%121=87%:%
%:%122=88%:%
%:%123=89%:%
%:%124=90%:%
%:%125=91%:%
%:%126=92%:%
%:%127=93%:%
%:%128=94%:%
%:%129=95%:%
%:%130=96%:%
%:%131=97%:%
%:%132=98%:%
%:%133=98%:%
%:%134=99%:%
%:%135=100%:%
%:%136=101%:%
%:%137=102%:%
%:%138=103%:%
%:%139=104%:%
%:%140=105%:%
%:%141=106%:%
%:%142=107%:%
%:%143=108%:%
%:%144=109%:%
%:%145=110%:%
%:%146=111%:%
%:%147=112%:%
%:%148=112%:%
%:%149=113%:%
%:%150=114%:%
%:%153=115%:%
%:%157=115%:%
%:%158=115%:%
%:%159=116%:%
%:%160=116%:%
%:%161=116%:%
%:%162=116%:%
%:%167=116%:%
%:%170=117%:%
%:%171=118%:%
%:%172=118%:%
%:%173=119%:%
%:%174=120%:%
%:%177=121%:%
%:%181=121%:%
%:%182=121%:%
%:%183=121%:%
%:%188=121%:%
%:%191=122%:%
%:%192=123%:%
%:%193=123%:%
%:%194=124%:%
%:%197=125%:%
%:%201=125%:%
%:%202=125%:%
%:%203=125%:%
%:%204=125%:%
%:%209=125%:%
%:%212=126%:%
%:%213=127%:%
%:%214=128%:%
%:%215=129%:%
%:%216=129%:%
%:%217=130%:%
%:%218=131%:%
%:%219=132%:%
%:%220=133%:%
%:%221=134%:%
%:%222=135%:%
%:%223=136%:%
%:%224=137%:%
%:%225=138%:%
%:%226=139%:%
%:%227=140%:%
%:%228=141%:%
%:%229=142%:%
%:%230=142%:%
%:%231=143%:%
%:%232=144%:%
%:%233=145%:%
%:%234=146%:%
%:%235=147%:%
%:%236=148%:%
%:%237=149%:%
%:%238=150%:%
%:%239=150%:%
%:%240=151%:%
%:%241=152%:%
%:%242=153%:%
%:%245=154%:%
%:%249=154%:%
%:%250=154%:%
%:%251=154%:%
%:%256=154%:%
%:%259=155%:%
%:%260=156%:%
%:%261=157%:%
%:%262=157%:%
%:%265=158%:%
%:%269=158%:%
%:%270=158%:%
%:%271=158%:%
%:%276=158%:%
%:%279=159%:%
%:%280=160%:%
%:%281=161%:%
%:%282=162%:%
%:%283=162%:%
%:%284=163%:%
%:%285=164%:%
%:%288=165%:%
%:%292=165%:%
%:%293=165%:%
%:%294=166%:%
%:%295=166%:%
%:%296=167%:%
%:%297=167%:%
%:%298=168%:%
%:%299=168%:%
%:%300=168%:%
%:%301=169%:%
%:%302=169%:%
%:%303=169%:%
%:%304=170%:%
%:%305=170%:%
%:%306=171%:%
%:%307=171%:%
%:%308=172%:%
%:%309=173%:%
%:%310=173%:%
%:%311=173%:%
%:%312=174%:%
%:%313=174%:%
%:%314=174%:%
%:%315=175%:%
%:%316=175%:%
%:%317=175%:%
%:%318=176%:%
%:%319=176%:%
%:%320=177%:%
%:%321=177%:%
%:%322=178%:%
%:%323=178%:%
%:%324=178%:%
%:%325=179%:%
%:%326=179%:%
%:%327=179%:%
%:%328=180%:%
%:%329=181%:%
%:%330=181%:%
%:%331=182%:%
%:%332=182%:%
%:%333=183%:%
%:%334=183%:%
%:%335=183%:%
%:%336=184%:%
%:%337=185%:%
%:%338=185%:%
%:%339=186%:%
%:%340=186%:%
%:%341=187%:%
%:%342=187%:%
%:%343=187%:%
%:%344=188%:%
%:%345=188%:%
%:%346=189%:%
%:%347=190%:%
%:%348=190%:%
%:%349=191%:%
%:%350=191%:%
%:%351=192%:%
%:%352=192%:%
%:%353=192%:%
%:%354=193%:%
%:%355=194%:%
%:%356=194%:%
%:%357=195%:%
%:%358=196%:%
%:%359=196%:%
%:%360=197%:%
%:%361=197%:%
%:%362=197%:%
%:%363=197%:%
%:%364=198%:%
%:%365=198%:%
%:%366=199%:%
%:%367=199%:%
%:%368=200%:%
%:%369=200%:%
%:%370=200%:%
%:%371=201%:%
%:%372=201%:%
%:%373=201%:%
%:%374=202%:%
%:%375=203%:%
%:%376=203%:%
%:%377=204%:%
%:%378=204%:%
%:%379=205%:%
%:%380=205%:%
%:%381=205%:%
%:%382=206%:%
%:%383=207%:%
%:%384=207%:%
%:%385=208%:%
%:%386=208%:%
%:%387=209%:%
%:%388=209%:%
%:%389=209%:%
%:%390=210%:%
%:%391=210%:%
%:%392=211%:%
%:%393=211%:%
%:%394=212%:%
%:%395=212%:%
%:%396=213%:%
%:%397=213%:%
%:%398=213%:%
%:%399=214%:%
%:%400=215%:%
%:%401=215%:%
%:%402=216%:%
%:%403=217%:%
%:%404=217%:%
%:%405=218%:%
%:%406=218%:%
%:%407=218%:%
%:%408=218%:%
%:%409=219%:%
%:%410=219%:%
%:%411=220%:%
%:%412=220%:%
%:%413=221%:%
%:%414=221%:%
%:%415=221%:%
%:%416=222%:%
%:%417=222%:%
%:%418=222%:%
%:%419=223%:%
%:%420=224%:%
%:%421=224%:%
%:%422=225%:%
%:%423=225%:%
%:%424=226%:%
%:%425=226%:%
%:%426=226%:%
%:%427=227%:%
%:%428=228%:%
%:%429=228%:%
%:%430=229%:%
%:%431=229%:%
%:%432=230%:%
%:%433=230%:%
%:%434=230%:%
%:%435=231%:%
%:%436=231%:%
%:%437=232%:%
%:%438=232%:%
%:%439=233%:%
%:%440=233%:%
%:%441=234%:%
%:%442=234%:%
%:%443=234%:%
%:%444=235%:%
%:%445=236%:%
%:%446=236%:%
%:%447=237%:%
%:%448=238%:%
%:%449=238%:%
%:%450=239%:%
%:%451=239%:%
%:%452=239%:%
%:%453=239%:%
%:%454=240%:%
%:%455=240%:%
%:%456=241%:%
%:%457=241%:%
%:%458=242%:%
%:%459=242%:%
%:%460=242%:%
%:%461=243%:%
%:%462=243%:%
%:%463=243%:%
%:%464=244%:%
%:%465=245%:%
%:%466=245%:%
%:%467=246%:%
%:%468=246%:%
%:%469=247%:%
%:%470=247%:%
%:%471=247%:%
%:%472=248%:%
%:%473=248%:%
%:%474=249%:%
%:%475=249%:%
%:%476=249%:%
%:%477=250%:%
%:%478=251%:%
%:%479=251%:%
%:%480=252%:%
%:%481=252%:%
%:%482=253%:%
%:%483=253%:%
%:%484=253%:%
%:%485=254%:%
%:%486=254%:%
%:%487=255%:%
%:%488=255%:%
%:%489=256%:%
%:%490=256%:%
%:%491=256%:%
%:%492=257%:%
%:%493=257%:%
%:%494=258%:%
%:%495=258%:%
%:%496=259%:%
%:%497=259%:%
%:%498=260%:%
%:%499=261%:%
%:%500=261%:%
%:%501=262%:%
%:%502=262%:%
%:%503=263%:%
%:%504=263%:%
%:%505=264%:%
%:%506=264%:%
%:%507=265%:%
%:%508=265%:%
%:%509=265%:%
%:%510=266%:%
%:%511=267%:%
%:%512=267%:%
%:%513=268%:%
%:%514=269%:%
%:%515=269%:%
%:%516=270%:%
%:%517=270%:%
%:%518=270%:%
%:%519=270%:%
%:%520=271%:%
%:%521=271%:%
%:%522=272%:%
%:%523=272%:%
%:%524=273%:%
%:%525=273%:%
%:%526=273%:%
%:%527=274%:%
%:%528=274%:%
%:%529=274%:%
%:%530=275%:%
%:%531=276%:%
%:%532=276%:%
%:%533=277%:%
%:%534=277%:%
%:%535=278%:%
%:%536=278%:%
%:%537=278%:%
%:%538=279%:%
%:%539=280%:%
%:%540=280%:%
%:%541=281%:%
%:%542=281%:%
%:%543=282%:%
%:%544=282%:%
%:%545=282%:%
%:%546=283%:%
%:%547=284%:%
%:%548=284%:%
%:%549=285%:%
%:%550=285%:%
%:%551=285%:%
%:%552=285%:%
%:%553=286%:%
%:%554=286%:%
%:%555=287%:%
%:%556=287%:%
%:%557=288%:%
%:%558=288%:%
%:%559=288%:%
%:%560=289%:%
%:%561=289%:%
%:%562=289%:%
%:%563=290%:%
%:%564=291%:%
%:%565=291%:%
%:%566=292%:%
%:%567=292%:%
%:%568=293%:%
%:%569=294%:%
%:%570=294%:%
%:%571=294%:%
%:%572=295%:%
%:%573=296%:%
%:%574=297%:%
%:%575=297%:%
%:%576=298%:%
%:%577=298%:%
%:%578=299%:%
%:%579=299%:%
%:%580=299%:%
%:%581=300%:%
%:%582=300%:%
%:%583=301%:%
%:%584=301%:%
%:%585=301%:%
%:%586=302%:%
%:%587=302%:%
%:%588=303%:%
%:%589=303%:%
%:%590=303%:%
%:%591=304%:%
%:%592=304%:%
%:%593=305%:%
%:%594=305%:%
%:%595=306%:%
%:%596=306%:%
%:%597=307%:%
%:%598=307%:%
%:%599=308%:%
%:%600=309%:%
%:%601=309%:%
%:%602=310%:%
%:%603=311%:%
%:%604=311%:%
%:%605=312%:%
%:%606=312%:%
%:%607=313%:%
%:%608=314%:%
%:%609=314%:%
%:%610=315%:%
%:%611=316%:%
%:%612=316%:%
%:%613=317%:%
%:%614=317%:%
%:%615=317%:%
%:%616=317%:%
%:%617=318%:%
%:%618=318%:%
%:%619=319%:%
%:%620=319%:%
%:%621=320%:%
%:%622=320%:%
%:%623=321%:%
%:%624=321%:%
%:%625=321%:%
%:%626=322%:%
%:%627=322%:%
%:%628=323%:%
%:%629=323%:%
%:%630=323%:%
%:%631=324%:%
%:%632=325%:%
%:%633=325%:%
%:%634=326%:%
%:%635=326%:%
%:%636=327%:%
%:%637=328%:%
%:%638=329%:%
%:%639=330%:%
%:%640=330%:%
%:%641=330%:%
%:%642=331%:%
%:%643=332%:%
%:%644=332%:%
%:%645=333%:%
%:%646=333%:%
%:%647=334%:%
%:%648=335%:%
%:%649=335%:%
%:%650=335%:%
%:%651=336%:%
%:%652=337%:%
%:%653=337%:%
%:%654=338%:%
%:%655=338%:%
%:%656=339%:%
%:%657=339%:%
%:%658=339%:%
%:%659=340%:%
%:%660=341%:%
%:%661=341%:%
%:%662=342%:%
%:%663=342%:%
%:%664=343%:%
%:%665=343%:%
%:%666=343%:%
%:%667=344%:%
%:%668=344%:%
%:%669=345%:%
%:%670=345%:%
%:%671=345%:%
%:%672=346%:%
%:%673=347%:%
%:%674=347%:%
%:%675=348%:%
%:%676=348%:%
%:%677=348%:%
%:%678=349%:%
%:%679=349%:%
%:%680=350%:%
%:%681=350%:%
%:%682=351%:%
%:%683=351%:%
%:%684=352%:%
%:%685=352%:%
%:%686=352%:%
%:%687=353%:%
%:%688=353%:%
%:%689=354%:%
%:%690=354%:%
%:%691=354%:%
%:%692=355%:%
%:%693=356%:%
%:%694=356%:%
%:%695=357%:%
%:%696=357%:%
%:%697=358%:%
%:%698=359%:%
%:%699=360%:%
%:%700=361%:%
%:%701=361%:%
%:%702=361%:%
%:%703=362%:%
%:%704=363%:%
%:%705=363%:%
%:%706=364%:%
%:%707=364%:%
%:%708=365%:%
%:%709=366%:%
%:%710=366%:%
%:%711=366%:%
%:%712=367%:%
%:%713=368%:%
%:%714=368%:%
%:%715=369%:%
%:%716=369%:%
%:%717=370%:%
%:%718=370%:%
%:%719=370%:%
%:%720=371%:%
%:%721=372%:%
%:%722=372%:%
%:%723=373%:%
%:%724=373%:%
%:%725=374%:%
%:%726=374%:%
%:%727=374%:%
%:%728=375%:%
%:%729=375%:%
%:%730=376%:%
%:%731=376%:%
%:%732=376%:%
%:%733=377%:%
%:%734=378%:%
%:%735=378%:%
%:%736=379%:%
%:%737=379%:%
%:%738=379%:%
%:%739=380%:%
%:%745=380%:%
%:%748=381%:%
%:%749=382%:%
%:%750=383%:%
%:%751=384%:%
%:%752=384%:%
%:%753=385%:%
%:%754=386%:%
%:%757=387%:%
%:%761=387%:%
%:%762=387%:%
%:%763=388%:%
%:%764=388%:%
%:%765=389%:%
%:%766=389%:%
%:%767=390%:%
%:%768=390%:%
%:%769=391%:%
%:%770=392%:%
%:%771=392%:%
%:%772=392%:%
%:%773=393%:%
%:%774=393%:%
%:%775=394%:%
%:%776=395%:%
%:%777=395%:%
%:%778=395%:%
%:%779=396%:%
%:%780=396%:%
%:%781=396%:%
%:%782=397%:%
%:%783=397%:%
%:%784=397%:%
%:%785=398%:%
%:%786=398%:%
%:%787=399%:%
%:%788=399%:%
%:%789=400%:%
%:%790=400%:%
%:%791=401%:%
%:%792=402%:%
%:%793=402%:%
%:%794=402%:%
%:%795=403%:%
%:%796=403%:%
%:%797=404%:%
%:%798=404%:%
%:%799=405%:%
%:%800=405%:%
%:%801=406%:%
%:%802=406%:%
%:%803=406%:%
%:%804=407%:%
%:%805=407%:%
%:%806=407%:%
%:%807=408%:%
%:%808=409%:%
%:%809=409%:%
%:%810=410%:%
%:%811=410%:%
%:%812=410%:%
%:%813=411%:%
%:%814=411%:%
%:%815=411%:%
%:%816=412%:%
%:%817=412%:%
%:%818=413%:%
%:%819=413%:%
%:%820=414%:%
%:%821=414%:%
%:%822=415%:%
%:%823=416%:%
%:%824=416%:%
%:%825=416%:%
%:%826=417%:%
%:%827=417%:%
%:%828=418%:%
%:%829=418%:%
%:%830=418%:%
%:%831=419%:%
%:%832=419%:%
%:%833=419%:%
%:%834=420%:%
%:%835=421%:%
%:%836=421%:%
%:%837=422%:%
%:%838=422%:%
%:%839=422%:%
%:%840=423%:%
%:%841=423%:%
%:%842=423%:%
%:%843=424%:%
%:%844=424%:%
%:%845=425%:%
%:%846=425%:%
%:%847=426%:%
%:%848=426%:%
%:%849=427%:%
%:%850=428%:%
%:%851=428%:%
%:%852=428%:%
%:%853=429%:%
%:%854=429%:%
%:%855=430%:%
%:%856=430%:%
%:%857=430%:%
%:%858=431%:%
%:%859=431%:%
%:%860=431%:%
%:%861=432%:%
%:%862=433%:%
%:%863=433%:%
%:%864=434%:%
%:%865=434%:%
%:%866=434%:%
%:%867=435%:%
%:%868=435%:%
%:%869=435%:%
%:%870=436%:%
%:%871=436%:%
%:%872=437%:%
%:%873=437%:%
%:%874=438%:%
%:%875=438%:%
%:%876=439%:%
%:%877=440%:%
%:%878=440%:%
%:%879=440%:%
%:%880=441%:%
%:%881=441%:%
%:%882=442%:%
%:%883=443%:%
%:%884=443%:%
%:%885=443%:%
%:%886=444%:%
%:%887=444%:%
%:%888=445%:%
%:%889=445%:%
%:%890=445%:%
%:%891=446%:%
%:%892=446%:%
%:%893=446%:%
%:%894=447%:%
%:%895=448%:%
%:%896=448%:%
%:%897=449%:%
%:%898=449%:%
%:%899=449%:%
%:%900=450%:%
%:%901=450%:%
%:%902=450%:%
%:%903=451%:%
%:%904=451%:%
%:%905=452%:%
%:%906=452%:%
%:%907=453%:%
%:%908=453%:%
%:%909=454%:%
%:%910=455%:%
%:%911=455%:%
%:%912=455%:%
%:%913=456%:%
%:%914=456%:%
%:%915=457%:%
%:%916=458%:%
%:%917=458%:%
%:%918=458%:%
%:%919=459%:%
%:%920=459%:%
%:%921=459%:%
%:%922=460%:%
%:%923=460%:%
%:%924=460%:%
%:%925=461%:%
%:%926=461%:%
%:%927=462%:%
%:%928=462%:%
%:%929=463%:%
%:%930=463%:%
%:%931=464%:%
%:%932=465%:%
%:%933=466%:%
%:%934=466%:%
%:%935=466%:%
%:%936=467%:%
%:%937=467%:%
%:%938=468%:%
%:%939=468%:%
%:%940=469%:%
%:%941=469%:%
%:%942=470%:%
%:%943=470%:%
%:%944=470%:%
%:%945=471%:%
%:%946=471%:%
%:%947=472%:%
%:%948=472%:%
%:%949=472%:%
%:%950=473%:%
%:%951=473%:%
%:%952=474%:%
%:%953=474%:%
%:%954=474%:%
%:%955=475%:%
%:%956=475%:%
%:%957=475%:%
%:%958=476%:%
%:%959=477%:%
%:%960=478%:%
%:%961=479%:%
%:%962=479%:%
%:%963=480%:%
%:%964=480%:%
%:%965=480%:%
%:%966=481%:%
%:%967=481%:%
%:%968=481%:%
%:%969=482%:%
%:%970=482%:%
%:%971=483%:%
%:%972=483%:%
%:%973=484%:%
%:%974=484%:%
%:%975=485%:%
%:%976=485%:%
%:%977=485%:%
%:%978=486%:%
%:%979=486%:%
%:%980=487%:%
%:%981=488%:%
%:%982=489%:%
%:%983=490%:%
%:%984=490%:%
%:%985=490%:%
%:%986=491%:%
%:%987=492%:%
%:%988=492%:%
%:%989=493%:%
%:%990=493%:%
%:%991=493%:%
%:%992=494%:%
%:%993=494%:%
%:%994=494%:%
%:%995=495%:%
%:%996=495%:%
%:%997=495%:%
%:%998=496%:%
%:%999=496%:%
%:%1000=497%:%
%:%1001=497%:%
%:%1002=498%:%
%:%1003=499%:%
%:%1004=500%:%
%:%1005=500%:%
%:%1006=501%:%
%:%1007=502%:%
%:%1008=502%:%
%:%1009=503%:%
%:%1010=503%:%
%:%1011=504%:%
%:%1012=505%:%
%:%1013=505%:%
%:%1014=505%:%
%:%1015=506%:%
%:%1016=507%:%
%:%1017=507%:%
%:%1018=507%:%
%:%1019=508%:%
%:%1020=508%:%
%:%1021=508%:%
%:%1022=509%:%
%:%1023=509%:%
%:%1024=510%:%
%:%1025=510%:%
%:%1026=511%:%
%:%1027=511%:%
%:%1028=512%:%
%:%1029=512%:%
%:%1030=512%:%
%:%1031=513%:%
%:%1032=513%:%
%:%1033=514%:%
%:%1034=515%:%
%:%1035=516%:%
%:%1036=517%:%
%:%1037=517%:%
%:%1038=517%:%
%:%1039=518%:%
%:%1040=519%:%
%:%1041=519%:%
%:%1042=520%:%
%:%1043=520%:%
%:%1044=520%:%
%:%1045=521%:%
%:%1046=521%:%
%:%1047=521%:%
%:%1048=522%:%
%:%1049=522%:%
%:%1050=522%:%
%:%1051=523%:%
%:%1052=523%:%
%:%1053=524%:%
%:%1054=524%:%
%:%1055=525%:%
%:%1056=526%:%
%:%1057=527%:%
%:%1058=527%:%
%:%1059=528%:%
%:%1060=528%:%
%:%1061=529%:%
%:%1062=529%:%
%:%1063=529%:%
%:%1064=530%:%
%:%1065=531%:%
%:%1066=531%:%
%:%1067=532%:%
%:%1068=532%:%
%:%1069=533%:%
%:%1070=534%:%
%:%1071=534%:%
%:%1072=534%:%
%:%1073=535%:%
%:%1074=536%:%
%:%1075=536%:%
%:%1076=536%:%
%:%1077=537%:%
%:%1078=537%:%
%:%1079=537%:%
%:%1080=538%:%
%:%1086=538%:%
%:%1089=539%:%
%:%1090=540%:%
%:%1091=541%:%
%:%1092=541%:%
%:%1093=542%:%
%:%1094=543%:%
%:%1095=544%:%
%:%1098=545%:%
%:%1102=545%:%
%:%1103=545%:%
%:%1104=546%:%
%:%1105=546%:%
%:%1106=547%:%
%:%1107=547%:%
%:%1108=548%:%
%:%1109=548%:%
%:%1110=549%:%
%:%1111=550%:%
%:%1112=550%:%
%:%1113=550%:%
%:%1114=551%:%
%:%1115=551%:%
%:%1116=552%:%
%:%1117=552%:%
%:%1118=553%:%
%:%1119=553%:%
%:%1120=554%:%
%:%1121=554%:%
%:%1122=555%:%
%:%1123=555%:%
%:%1124=556%:%
%:%1125=556%:%
%:%1126=556%:%
%:%1127=557%:%
%:%1128=558%:%
%:%1129=558%:%
%:%1130=559%:%
%:%1131=559%:%
%:%1132=560%:%
%:%1133=560%:%
%:%1134=561%:%
%:%1135=561%:%
%:%1136=561%:%
%:%1137=562%:%
%:%1138=562%:%
%:%1139=563%:%
%:%1140=563%:%
%:%1141=564%:%
%:%1142=564%:%
%:%1143=565%:%
%:%1144=566%:%
%:%1145=566%:%
%:%1146=566%:%
%:%1147=567%:%
%:%1148=567%:%
%:%1149=568%:%
%:%1150=568%:%
%:%1151=569%:%
%:%1152=569%:%
%:%1153=570%:%
%:%1154=570%:%
%:%1155=570%:%
%:%1156=571%:%
%:%1157=571%:%
%:%1158=571%:%
%:%1159=572%:%
%:%1160=573%:%
%:%1161=573%:%
%:%1162=574%:%
%:%1163=574%:%
%:%1164=575%:%
%:%1165=575%:%
%:%1166=576%:%
%:%1167=576%:%
%:%1168=576%:%
%:%1169=577%:%
%:%1170=577%:%
%:%1171=578%:%
%:%1172=578%:%
%:%1173=579%:%
%:%1174=579%:%
%:%1175=580%:%
%:%1176=581%:%
%:%1177=581%:%
%:%1178=581%:%
%:%1179=582%:%
%:%1180=582%:%
%:%1181=583%:%
%:%1182=583%:%
%:%1183=584%:%
%:%1184=584%:%
%:%1185=585%:%
%:%1186=585%:%
%:%1187=585%:%
%:%1188=586%:%
%:%1189=586%:%
%:%1190=586%:%
%:%1191=587%:%
%:%1192=588%:%
%:%1193=588%:%
%:%1194=589%:%
%:%1195=589%:%
%:%1196=590%:%
%:%1197=590%:%
%:%1198=591%:%
%:%1199=591%:%
%:%1200=591%:%
%:%1201=592%:%
%:%1202=592%:%
%:%1203=593%:%
%:%1204=593%:%
%:%1205=594%:%
%:%1206=594%:%
%:%1207=595%:%
%:%1208=596%:%
%:%1209=596%:%
%:%1210=596%:%
%:%1211=597%:%
%:%1212=597%:%
%:%1213=597%:%
%:%1214=598%:%
%:%1215=598%:%
%:%1216=598%:%
%:%1217=599%:%
%:%1218=599%:%
%:%1219=600%:%
%:%1220=600%:%
%:%1221=601%:%
%:%1222=601%:%
%:%1223=601%:%
%:%1224=602%:%
%:%1225=602%:%
%:%1226=603%:%
%:%1227=603%:%
%:%1228=603%:%
%:%1229=604%:%
%:%1230=604%:%
%:%1231=605%:%
%:%1232=605%:%
%:%1233=605%:%
%:%1234=605%:%
%:%1235=606%:%
%:%1236=606%:%
%:%1241=606%:%
%:%1244=607%:%
%:%1245=608%:%
%:%1246=609%:%
%:%1247=609%:%
%:%1248=610%:%
%:%1249=611%:%
%:%1250=612%:%
%:%1253=613%:%
%:%1257=613%:%
%:%1258=613%:%
%:%1259=614%:%
%:%1260=614%:%
%:%1261=615%:%
%:%1262=615%:%
%:%1263=616%:%
%:%1264=616%:%
%:%1265=617%:%
%:%1266=618%:%
%:%1267=619%:%
%:%1268=619%:%
%:%1269=619%:%
%:%1270=620%:%
%:%1271=620%:%
%:%1272=621%:%
%:%1273=622%:%
%:%1274=622%:%
%:%1275=622%:%
%:%1276=623%:%
%:%1277=623%:%
%:%1278=623%:%
%:%1279=624%:%
%:%1280=624%:%
%:%1281=625%:%
%:%1282=625%:%
%:%1283=626%:%
%:%1284=626%:%
%:%1285=626%:%
%:%1286=627%:%
%:%1287=627%:%
%:%1288=628%:%
%:%1289=628%:%
%:%1290=628%:%
%:%1291=629%:%
%:%1292=629%:%
%:%1293=630%:%
%:%1294=630%:%
%:%1295=631%:%
%:%1296=631%:%
%:%1297=632%:%
%:%1298=632%:%
%:%1299=632%:%
%:%1300=633%:%
%:%1301=633%:%
%:%1302=633%:%
%:%1303=634%:%
%:%1304=634%:%
%:%1305=635%:%
%:%1306=635%:%
%:%1307=636%:%
%:%1308=636%:%
%:%1309=637%:%
%:%1310=638%:%
%:%1311=639%:%
%:%1312=639%:%
%:%1313=639%:%
%:%1314=640%:%
%:%1315=640%:%
%:%1316=641%:%
%:%1317=642%:%
%:%1318=642%:%
%:%1319=642%:%
%:%1320=643%:%
%:%1321=643%:%
%:%1322=643%:%
%:%1323=644%:%
%:%1324=644%:%
%:%1325=645%:%
%:%1326=645%:%
%:%1327=646%:%
%:%1328=646%:%
%:%1329=646%:%
%:%1330=647%:%
%:%1331=647%:%
%:%1332=648%:%
%:%1333=648%:%
%:%1334=648%:%
%:%1335=649%:%
%:%1336=649%:%
%:%1337=650%:%
%:%1338=650%:%
%:%1339=651%:%
%:%1340=651%:%
%:%1341=652%:%
%:%1342=652%:%
%:%1343=652%:%
%:%1344=653%:%
%:%1345=653%:%
%:%1346=653%:%
%:%1347=654%:%
%:%1348=654%:%
%:%1349=655%:%
%:%1350=655%:%
%:%1351=656%:%
%:%1352=656%:%
%:%1353=657%:%
%:%1354=658%:%
%:%1355=659%:%
%:%1356=659%:%
%:%1357=659%:%
%:%1358=660%:%
%:%1359=660%:%
%:%1360=661%:%
%:%1361=662%:%
%:%1362=662%:%
%:%1363=662%:%
%:%1364=663%:%
%:%1365=663%:%
%:%1366=663%:%
%:%1367=664%:%
%:%1368=664%:%
%:%1369=665%:%
%:%1370=665%:%
%:%1371=666%:%
%:%1372=666%:%
%:%1373=666%:%
%:%1374=667%:%
%:%1375=667%:%
%:%1376=668%:%
%:%1377=668%:%
%:%1378=668%:%
%:%1379=669%:%
%:%1380=669%:%
%:%1381=670%:%
%:%1382=670%:%
%:%1383=671%:%
%:%1384=671%:%
%:%1385=672%:%
%:%1386=672%:%
%:%1387=672%:%
%:%1388=673%:%
%:%1389=673%:%
%:%1390=673%:%
%:%1391=674%:%
%:%1392=674%:%
%:%1393=675%:%
%:%1394=675%:%
%:%1395=676%:%
%:%1396=676%:%
%:%1397=677%:%
%:%1398=678%:%
%:%1399=678%:%
%:%1400=678%:%
%:%1401=679%:%
%:%1402=679%:%
%:%1403=680%:%
%:%1404=681%:%
%:%1405=681%:%
%:%1406=681%:%
%:%1407=682%:%
%:%1408=682%:%
%:%1409=682%:%
%:%1410=683%:%
%:%1411=683%:%
%:%1412=683%:%
%:%1413=684%:%
%:%1414=684%:%
%:%1415=685%:%
%:%1416=685%:%
%:%1417=685%:%
%:%1418=686%:%
%:%1419=686%:%
%:%1420=687%:%
%:%1421=687%:%
%:%1422=687%:%
%:%1423=688%:%
%:%1424=688%:%
%:%1425=689%:%
%:%1426=689%:%
%:%1427=690%:%
%:%1428=690%:%
%:%1429=691%:%
%:%1430=691%:%
%:%1431=691%:%
%:%1432=692%:%
%:%1433=692%:%
%:%1434=692%:%
%:%1435=693%:%
%:%1436=693%:%
%:%1441=693%:%
%:%1444=694%:%
%:%1445=718%:%
%:%1446=719%:%
%:%1447=720%:%
%:%1448=721%:%
%:%1449=721%:%
%:%1450=722%:%
%:%1451=723%:%
%:%1452=724%:%
%:%1455=725%:%
%:%1459=725%:%
%:%1460=725%:%
%:%1461=726%:%
%:%1462=726%:%
%:%1463=727%:%
%:%1464=727%:%
%:%1465=728%:%
%:%1466=728%:%
%:%1467=728%:%
%:%1468=729%:%
%:%1469=729%:%
%:%1470=729%:%
%:%1471=730%:%
%:%1472=730%:%
%:%1473=730%:%
%:%1474=731%:%
%:%1475=731%:%
%:%1476=732%:%
%:%1477=732%:%
%:%1478=733%:%
%:%1479=733%:%
%:%1480=733%:%
%:%1481=734%:%
%:%1482=734%:%
%:%1483=734%:%
%:%1484=735%:%
%:%1485=735%:%
%:%1486=735%:%
%:%1487=736%:%
%:%1488=736%:%
%:%1489=737%:%
%:%1490=737%:%
%:%1491=738%:%
%:%1492=738%:%
%:%1493=738%:%
%:%1494=739%:%
%:%1495=739%:%
%:%1496=739%:%
%:%1497=739%:%
%:%1498=739%:%
%:%1499=740%:%
%:%1500=740%:%
%:%1501=741%:%
%:%1502=741%:%
%:%1503=741%:%
%:%1504=742%:%
%:%1505=742%:%
%:%1506=743%:%
%:%1507=743%:%
%:%1508=743%:%
%:%1509=744%:%
%:%1510=744%:%
%:%1511=745%:%
%:%1512=745%:%
%:%1513=746%:%
%:%1514=746%:%
%:%1515=746%:%
%:%1516=747%:%
%:%1517=747%:%
%:%1518=747%:%
%:%1519=747%:%
%:%1520=747%:%
%:%1521=748%:%
%:%1522=748%:%
%:%1523=749%:%
%:%1524=749%:%
%:%1525=749%:%
%:%1526=750%:%
%:%1527=750%:%
%:%1528=750%:%
%:%1529=750%:%
%:%1530=751%:%
%:%1536=751%:%
%:%1539=752%:%
%:%1540=753%:%
%:%1541=753%:%
%:%1542=754%:%
%:%1543=755%:%
%:%1544=756%:%
%:%1547=757%:%
%:%1551=757%:%
%:%1552=757%:%
%:%1553=758%:%
%:%1554=758%:%
%:%1555=759%:%
%:%1556=759%:%
%:%1557=760%:%
%:%1558=760%:%
%:%1559=760%:%
%:%1560=760%:%
%:%1561=761%:%
%:%1562=761%:%
%:%1563=762%:%
%:%1564=762%:%
%:%1565=763%:%
%:%1566=763%:%
%:%1567=763%:%
%:%1568=764%:%
%:%1569=764%:%
%:%1570=764%:%
%:%1571=764%:%
%:%1572=764%:%
%:%1573=765%:%
%:%1574=765%:%
%:%1575=766%:%
%:%1576=766%:%
%:%1577=767%:%
%:%1578=767%:%
%:%1579=767%:%
%:%1580=768%:%
%:%1581=768%:%
%:%1582=768%:%
%:%1583=768%:%
%:%1584=768%:%
%:%1585=769%:%
%:%1586=769%:%
%:%1587=769%:%
%:%1588=769%:%
%:%1589=770%:%
%:%1590=770%:%
%:%1591=770%:%
%:%1592=770%:%
%:%1593=771%:%
%:%1594=771%:%
%:%1595=772%:%
%:%1596=772%:%
%:%1597=773%:%
%:%1598=773%:%
%:%1599=773%:%
%:%1600=774%:%
%:%1601=774%:%
%:%1602=774%:%
%:%1603=775%:%
%:%1604=776%:%
%:%1605=776%:%
%:%1606=776%:%
%:%1607=776%:%
%:%1608=776%:%
%:%1609=777%:%
%:%1610=777%:%
%:%1611=778%:%
%:%1612=778%:%
%:%1613=778%:%
%:%1614=778%:%
%:%1615=778%:%
%:%1616=779%:%
%:%1617=779%:%
%:%1618=779%:%
%:%1619=779%:%
%:%1620=779%:%
%:%1621=780%:%
%:%1622=780%:%
%:%1623=780%:%
%:%1624=781%:%
%:%1625=781%:%
%:%1626=782%:%
%:%1632=782%:%
%:%1635=783%:%
%:%1636=784%:%
%:%1637=784%:%
%:%1638=785%:%
%:%1639=786%:%
%:%1640=787%:%
%:%1643=788%:%
%:%1647=788%:%
%:%1648=788%:%
%:%1649=789%:%
%:%1650=789%:%
%:%1651=790%:%
%:%1652=790%:%
%:%1653=791%:%
%:%1654=791%:%
%:%1655=791%:%
%:%1656=792%:%
%:%1657=792%:%
%:%1658=792%:%
%:%1659=793%:%
%:%1660=793%:%
%:%1661=793%:%
%:%1662=794%:%
%:%1663=794%:%
%:%1664=795%:%
%:%1665=795%:%
%:%1666=796%:%
%:%1667=796%:%
%:%1668=797%:%
%:%1669=797%:%
%:%1670=798%:%
%:%1671=798%:%
%:%1672=799%:%
%:%1673=799%:%
%:%1674=800%:%
%:%1675=800%:%
%:%1676=800%:%
%:%1677=801%:%
%:%1678=801%:%
%:%1679=801%:%
%:%1680=802%:%
%:%1681=803%:%
%:%1682=803%:%
%:%1683=803%:%
%:%1684=803%:%
%:%1685=803%:%
%:%1686=804%:%
%:%1687=804%:%
%:%1688=804%:%
%:%1689=804%:%
%:%1690=805%:%
%:%1691=805%:%
%:%1692=805%:%
%:%1693=806%:%
%:%1694=806%:%
%:%1695=807%:%
%:%1696=807%:%
%:%1697=808%:%
%:%1698=808%:%
%:%1699=808%:%
%:%1700=809%:%
%:%1701=809%:%
%:%1702=809%:%
%:%1703=810%:%
%:%1704=810%:%
%:%1705=811%:%
%:%1706=811%:%
%:%1707=812%:%
%:%1708=812%:%
%:%1709=812%:%
%:%1710=813%:%
%:%1711=813%:%
%:%1712=814%:%
%:%1713=814%:%
%:%1714=815%:%
%:%1715=815%:%
%:%1716=816%:%
%:%1717=816%:%
%:%1718=817%:%
%:%1724=817%:%
%:%1727=818%:%
%:%1728=819%:%
%:%1729=819%:%
%:%1730=820%:%
%:%1731=821%:%
%:%1732=822%:%
%:%1735=823%:%
%:%1739=823%:%
%:%1740=823%:%
%:%1741=824%:%
%:%1742=824%:%
%:%1743=825%:%
%:%1744=825%:%
%:%1745=826%:%
%:%1746=826%:%
%:%1747=826%:%
%:%1748=827%:%
%:%1749=827%:%
%:%1750=827%:%
%:%1751=828%:%
%:%1752=828%:%
%:%1753=828%:%
%:%1754=829%:%
%:%1755=829%:%
%:%1756=829%:%
%:%1757=830%:%
%:%1758=830%:%
%:%1759=830%:%
%:%1760=831%:%
%:%1761=831%:%
%:%1762=832%:%
%:%1763=832%:%
%:%1764=833%:%
%:%1765=833%:%
%:%1766=834%:%
%:%1767=835%:%
%:%1768=835%:%
%:%1769=836%:%
%:%1770=836%:%
%:%1771=837%:%
%:%1772=837%:%
%:%1773=837%:%
%:%1774=838%:%
%:%1775=838%:%
%:%1776=838%:%
%:%1777=839%:%
%:%1778=839%:%
%:%1779=839%:%
%:%1780=840%:%
%:%1781=840%:%
%:%1782=841%:%
%:%1783=841%:%
%:%1784=842%:%
%:%1785=842%:%
%:%1786=843%:%
%:%1787=843%:%
%:%1788=844%:%
%:%1789=844%:%
%:%1790=845%:%
%:%1791=846%:%
%:%1792=846%:%
%:%1793=847%:%
%:%1794=847%:%
%:%1795=848%:%
%:%1796=848%:%
%:%1797=848%:%
%:%1798=849%:%
%:%1799=849%:%
%:%1800=850%:%
%:%1801=850%:%
%:%1802=851%:%
%:%1803=851%:%
%:%1804=852%:%
%:%1805=852%:%
%:%1806=852%:%
%:%1807=853%:%
%:%1808=853%:%
%:%1809=853%:%
%:%1810=854%:%
%:%1811=854%:%
%:%1812=855%:%
%:%1813=855%:%
%:%1814=856%:%
%:%1815=856%:%
%:%1816=857%:%
%:%1817=857%:%
%:%1818=858%:%
%:%1819=858%:%
%:%1820=859%:%
%:%1821=860%:%
%:%1822=860%:%
%:%1823=861%:%
%:%1824=861%:%
%:%1825=862%:%
%:%1826=863%:%
%:%1827=863%:%
%:%1828=863%:%
%:%1829=864%:%
%:%1830=864%:%
%:%1831=865%:%
%:%1832=865%:%
%:%1833=865%:%
%:%1834=866%:%
%:%1835=866%:%
%:%1836=867%:%
%:%1837=867%:%
%:%1838=867%:%
%:%1839=868%:%
%:%1840=868%:%
%:%1841=868%:%
%:%1842=869%:%
%:%1843=869%:%
%:%1844=869%:%
%:%1845=869%:%
%:%1846=870%:%
%:%1852=870%:%
%:%1855=871%:%
%:%1856=872%:%
%:%1857=872%:%
%:%1858=873%:%
%:%1859=874%:%
%:%1866=875%:%
%:%1867=875%:%
%:%1868=876%:%
%:%1869=876%:%
%:%1870=877%:%
%:%1871=877%:%
%:%1872=877%:%
%:%1873=878%:%
%:%1874=878%:%
%:%1875=878%:%
%:%1876=879%:%
%:%1877=879%:%
%:%1878=880%:%
%:%1879=880%:%
%:%1880=881%:%
%:%1881=882%:%
%:%1882=882%:%
%:%1883=883%:%
%:%1884=884%:%
%:%1885=884%:%
%:%1886=885%:%
%:%1887=885%:%
%:%1888=886%:%
%:%1889=886%:%
%:%1890=887%:%
%:%1891=887%:%
%:%1892=887%:%
%:%1894=889%:%
%:%1895=890%:%
%:%1896=890%:%
%:%1897=891%:%
%:%1898=891%:%
%:%1899=892%:%
%:%1900=892%:%
%:%1901=893%:%
%:%1902=893%:%
%:%1903=893%:%
%:%1904=894%:%
%:%1905=894%:%
%:%1906=895%:%
%:%1907=895%:%
%:%1908=895%:%
%:%1909=896%:%
%:%1910=896%:%
%:%1911=896%:%
%:%1912=896%:%
%:%1913=897%:%
%:%1914=897%:%
%:%1915=898%:%
%:%1916=898%:%
%:%1917=899%:%
%:%1923=899%:%
%:%1926=900%:%
%:%1929=901%:%
%:%1934=902%:%



% optional bibliography
\addcontentsline{toc}{section}{References}
\bibliographystyle{abbrv}
\bibliography{root}



%
\end{document}

%%% Local Variables:
%%% mode: latex
%%% TeX-master: t
%%% End:
\endinput
%:%file=~/nominal_unification_Isabelle/nominal_alpha_unification/document/root.tex%:%
